\documentclass[12pt,fontset=fandol]{ctexart}

\makeatletter
\def\input@path{{../}{../../}}
\makeatother

\usepackage{researchnotes}

\title{全国大学生数学竞赛数学分析与高等数学试题选编}
\author{Jacobi Lie 原题解及方法稿\quad 课程选编与补充解析}
\date{2026年10月10日}

\begin{document}
\maketitle

\section*{编写说明、来源与阅读路线}
\addcontentsline{toc}{section}{编写说明、来源与阅读路线}
本书面向本科数学分析与高等数学课程，以全国大学生数学竞赛历年题目为线索，组织计算、证明及方法研习。数学专业组部分选自第1至17届初赛及决赛，并包含第13届补赛、第14届两次补赛；非数学专业组另列。数学类试卷中的纯高等代数、抽象代数、概率论及专门微分几何题不在选编范围。涉及常微分方程、实变量积分与逼近的交叉题予以保留，并在阅读时区分基础微积分与进阶分析。这里的“选编”不表示各组别历年全部试卷均已收齐。

主要可编辑题解来源为 Jacobi Lie 的\href{https://github.com/lzxcr/cmc-mathematics}{《全国大学生数学竞赛题解》与《数学竞赛的方法与结构》}，使用2026年10月10日下载的源码快照，版本标识为 \href{https://github.com/lzxcr/cmc-mathematics/tree/64916fa1fdbc3e444a652eca204d5669058ee321}{固定源码版本（64916fa）}。原作者的序言、解答、评注、方法说明及编排采用\href{https://creativecommons.org/licenses/by-sa/4.0/}{CC BY-SA 4.0（署名—相同方式共享4.0）}许可。本书对相关内容作课程筛选、章节重组、TeX 转写及说明补充；这些改编内容沿用同一许可。引用试题及第三方材料的权利仍归原权利人，该许可不改变它们原有的权利状态。本文不是组委会审定的官方题解。

数学专业组的题面沿用开放题解的规范化表述，部分题补明了定义域、正则性、退化情形或约定；不能据此把措辞差异解释为官方试卷修订。题后保留源码编号与可核查的逐题链接。初赛编号通常对应原大题；决赛中填空题可能已拆分，高低年级共同题已合并，因此“来源序号”指所引题解的整理编号，而不是强称原卷大题号。有关拆分方式见书末来源与编目说明。未逐项核对原卷扫描件的题目，其来源层级为公开题解转录。

本选编实际收录194道真题，其中数学专业组157道、非数学专业组37道。每道题与其解答相邻编排，同一题只计一次；专题方法中的例题、推广和反例不计为真题。

每道真题后紧接解答，建议先独立作答，再阅读题后的推导与说明。解答补充方法选择、关键计算及所用结论的条件，原来源有评注者保留评注；专题方法编则集中说明可迁移的证明工具。高等数学读者可先学习极限、导数、定积分、多元微积分、级数与初等微分方程；数学专业读者继续研习一致收敛、反常积分、一致估计和逼近。含有更强工具的讨论作为进阶阅读，不将其暗示为所有参赛组别的必考内容。

未特别声明时，变量取实数，$\log$ 与 $\ln$ 均指自然对数，三角函数使用弧度制。$C^k$ 表示有连续的 $k$ 阶导数；$O$、$o$ 和 $\sim$ 的含义在方法编中说明。常微分方程的通解须同时说明定义区间；分离变量时另查被除去的零解。积分换序、积分号下求导和极限交换须依据绝对可积、非负性、一致收敛或明确的控制条件。成功编译只能验证源文件和交叉引用，不能替代逐题数学审稿。
\tableofcontents
\clearpage


\part{数学专业组真题与详解}

\section{初赛：第 1 届（2009 年）}

本节选入4道课程相关题；来源序号保留原题解编排，未选入题目不重新补号。

\begin{example}[来源序号4]\label{ex:pre01-math:4}\label{pr:pre01-math:4}

设 $f_n\in C^\infty([a,b])$ 逐点收敛，且存在与 $n,x$ 无关的常数 $M$，使
\[
|f_n'(x)|\le M\qquad(x\in[a,b]).
\]

（1） 证明 $f_n$ 一致收敛。

（2） 极限函数是否必在 $[a,b]$ 上处处可导？

\par\href{https://github.com/lzxcr/cmc-mathematics/blob/64916fa1fdbc3e444a652eca204d5669058ee321/src/problems/preliminary/01/04.md}{题目与来源条目}

\end{example}

\begin{solve}\label{ans:pre01-math:4}

导数的一致控制意味着所有 $f_n$ 都是 $M$-Lipschitz 函数；逐点取极限后，$f=\lim f_n$ 也具有这一性质。

这里先把“导数有界”转化为可直接使用的差值估计。对任意 $x,y\in[a,b]$，中值定理给出 $|f_n(x)-f_n(y)|\le M|x-y|$；固定 $x,y$ 后让 $n\to\infty$，便得 $|f(x)-f(y)|\le M|x-y|$。所以下面可以同时控制 $f_n$ 与 $f$ 在相邻两点的变化。逐点收敛本来允许收敛速度随点变化，而有限网格上的有限多个下标可以取最大值，这正是得到统一下标的原因。

给定 $\varepsilon>0$，在 $[a,b]$ 中取有限个点 $x_1,\ldots,x_N$，使每个 $x$ 到某个 $x_j$ 的距离小于 $\varepsilon/[3(M+1)]$。当 $n$ 足够大时，所有这些点同时满足 $|f_n(x_j)-f(x_j)|<\varepsilon/3$，从而
\[
|f_n(x)-f(x)|
\le M|x-x_j|+|f_n(x_j)-f(x_j)|+M|x-x_j|<\varepsilon.
\]
这就证明了一致收敛。

极限未必处处可导。若 $a<b$，取 $c\in(a,b)$，令
\[
f_n(x)=\sqrt{(x-c)^2+\frac1n}.
\]
则 $|f_n'|\le1$，但极限 $|x-c|$ 在 $c$ 处不可导。

反例中的导数为 $f_n'(x)=(x-c)/\sqrt{(x-c)^2+1/n}$，其绝对值不超过 $1$；而且 $0\le f_n(x)-|x-c|\le n^{-1/2}$，确实一致趋近 $|x-c|$。极限函数在 $c$ 处的左、右导数分别为 $-1$ 和 $1$，从而明确否定了第二问。

\end{solve}

\begin{remark}[方法与条件]有限网格把逐点控制提升为全局控制，是紧空间上等度连续函数族的基本思想，也是 Arzelà--Ascoli 定理的核心。题中的无穷次可微远强于所需条件：统一的连续性模量已经足够。\end{remark}

\begin{example}[来源序号5]\label{ex:pre01-math:5}\label{pr:pre01-math:5}

设
\[
a_n=\int_0^{\pi/2}t\left|\frac{\sin(nt)}{\sin t}\right|^3\,dt.
\]
证明正项级数 $\sum_{n=1}^\infty a_n^{-1}$ 发散。

\par\href{https://github.com/lzxcr/cmc-mathematics/blob/64916fa1fdbc3e444a652eca204d5669058ee321/src/problems/preliminary/01/05.md}{题目与来源条目}

\end{example}

\begin{solve}\label{ans:pre01-math:5}

只需证明 $a_n=O(n)$。比值在原点附近至多为 $n$，离开原点后至多为 $O(t^{-1})$，两种估计在尺度 $t\asymp n^{-1}$ 相接。

由 $|\sin(nt)|\le n\sin t$ 及 $\sin t\ge2t/\pi$，有
\[
\left|\frac{\sin(nt)}{\sin t}\right|
\le\min\left\{n,\frac{\pi}{2t}\right\},\qquad 0<t\le\frac\pi2.
\]
第一条界可只用实变量证明：$|\sin((n+1)t)|\le |\sin(nt)|+\sin t$，从 $n=1$ 出发归纳即得 $|\sin(nt)|\le n\sin t$。第二条界使用 $|\sin(nt)|\le1$ 和正弦函数在 $[0,\pi/2]$ 上的凹性。若全区间只用第一条界，就只能得到 $O(n^3)$，不足以判定倒数级数；分段的目的正是让两种估计各自在合适的区域发挥作用。

在 $t=1/n$ 处分开积分，得到
\[
0<a_n\le n^3\int_0^{1/n}t\,dt
+\frac{\pi^3}{8}\int_{1/n}^{\pi/2}\frac{dt}{t^2}
\le\left(\frac12+\frac{\pi^3}{8}\right)n.
\]
故 $a_n^{-1}\ge c/n$，与调和级数比较即得结论。

\end{solve}

\begin{remark}[方法与条件]第一条三角估计也可由
\[
\frac{\sin(nt)}{\sin t}
=\sum_{j=0}^{n-1}e^{i(n-1-2j)t}
\]
及三角不等式直接得到。积分分段的尺度来自两种界的平衡；这一原则见\hyperref[app:estimate-scaling]{方法卷：尺度与分区估计}。\end{remark}

\begin{example}[来源序号6]\label{ex:pre01-math:6}\label{pr:pre01-math:6}

设 $f$ 在闭单位圆盘 $D$ 上二次连续可微，且
\[
\Delta f=f_{xx}+f_{yy}=x^2y^2.
\]
计算
\[
I=\iint_D\left(
\frac{x}{\sqrt{x^2+y^2}}f_x+
\frac{y}{\sqrt{x^2+y^2}}f_y
\right)\,dx\,dy.
\]
原点处的被积函数可任意赋值。

\par\href{https://github.com/lzxcr/cmc-mathematics/blob/64916fa1fdbc3e444a652eca204d5669058ee321/src/problems/preliminary/01/06.md}{题目与来源条目}

\end{example}

\begin{solve}\label{ans:pre01-math:6}

被积函数是径向导数 $\partial_r f$。把圆盘按半径分层，再对每个小圆盘使用散度定理，便可把未知的一阶导数换成已知的 Laplace 算子。记 $D_r=\{x^2+y^2\le r^2\}$，则
\[
\begin{aligned}
I
&=\int_0^1\left(\int_{\partial D_r}\partial_\nu f\,ds\right)dr\\
&=\int_0^1\left(\iint_{D_r}x^2y^2\,dx\,dy\right)dr.
\end{aligned}
\]
具体地，令 $x=r\cos\theta,y=r\sin\theta$，链式法则给出 $\partial_r f=f_x\cos\theta+f_y\sin\theta$，面积元为 $r\,dr\,d\theta$。在圆周上外法向量也是 $(\cos\theta,\sin\theta)$，弧长元为 $r\,d\theta$，所以每层角积分恰好等于边界通量。$f\in C^2$ 保证散度定理适用；原点是面积为零的集合，赋值不影响积分。

而
\[
\iint_{D_r}x^2y^2\,dx\,dy
=\int_0^r\rho^5\,d\rho
\int_0^{2\pi}\cos^2\theta\sin^2\theta\,d\theta
=\frac{\pi r^6}{24}.
\]
因此
\[
\boxed{I=\frac\pi{168}.}
\]

最后的角积分可展开为 $\cos^2\theta\sin^2\theta=(1-\cos4\theta)/8$，其在 $[0,2\pi]$ 上的积分为 $\pi/4$；径向积分是 $r^6/6$。再对 $r$ 从 $0$ 积到 $1$，即得 $(\pi/24)\cdot(1/7)=\pi/168$。

\end{solve}

\begin{remark}[方法与条件]$f$ 本身并不由 $\Delta f$ 唯一决定，但任意调和函数的每层边界通量均为零，所以这里的积分仍被唯一确定。更一般地，径向加权的通量积分可以转化为对 $\Delta f$ 的径向加权积分。\end{remark}

\begin{example}[来源序号7]\label{ex:pre01-math:7}\label{pr:pre01-math:7}

设 $f\in C([0,1])$，且在 $(0,1)$ 内二阶可导。连接 $(0,f(0))$ 与 $(1,f(1))$ 的直线还经过曲线 $y=f(x)$ 上的一点 $(c,f(c))$，其中 $0<c<1$。证明存在 $\xi\in(0,1)$，使 $f''(\xi)=0$。

\par\href{https://github.com/lzxcr/cmc-mathematics/blob/64916fa1fdbc3e444a652eca204d5669058ee321/src/problems/preliminary/01/07.md}{题目与来源条目}

\end{example}

\begin{solve}\label{ans:pre01-math:7}

减去弦所在的仿射函数，令
\[
g(x)=f(x)-\bigl((1-x)f(0)+xf(1)\bigr).
\]
则 $g(0)=g(c)=g(1)=0$。分别在 $[0,c]$ 与 $[c,1]$ 上应用 Rolle 定理，得到 $u\in(0,c)$、$v\in(c,1)$，满足 $g'(u)=g'(v)=0$。再对 $g'$ 在 $[u,v]$ 上应用 Rolle 定理，得
\[
f''(\xi)=g''(\xi)=0
\]
对某个 $\xi\in(u,v)$ 成立。

辅助函数的选择来自题设的共线条件：弦在横坐标 $x$ 处的纵坐标为 $(1-x)f(0)+xf(1)$，因此第三个交点恰好给出 $g(c)=0$。第一次应用 Rolle 定理所需的闭区间连续性由 $f\in C([0,1])$ 保证；第二次应用时，$u,v$ 都位于 $(0,1)$ 内，而 $g'$ 在这里可导，故在 $[u,v]$ 上连续。这样三处零点依次产生两处一阶导数零点和一处二阶导数零点。

\end{solve}

\begin{remark}[方法与条件]减去插值多项式后反复使用 Rolle 定理，是插值余项公式的基本机制。二阶导数无需连续；二阶可导已保证应用 Rolle 定理所需的一阶导数连续性。\end{remark}

\section{初赛：第 2 届（2010 年）}

本节选入5道课程相关题；来源序号保留原题解编排，未选入题目不重新补号。

\begin{example}[来源序号1]\label{ex:pre02-math:1}\label{pr:pre02-math:1}

设 $a\in\mathbb R$，$0<\varepsilon<1$，定义
\[
x_0=a,\qquad x_{n+1}=a+\varepsilon\sin x_n.
\]
证明 $x_n$ 收敛，且其极限是方程 $x-\varepsilon\sin x=a$ 的唯一实根。

\par\href{https://github.com/lzxcr/cmc-mathematics/blob/64916fa1fdbc3e444a652eca204d5669058ee321/src/problems/preliminary/02/01.md}{题目与来源条目}

\end{example}

\begin{solve}\label{ans:pre02-math:1}

这是压缩映射的不动点迭代。令 $T(x)=a+\varepsilon\sin x$，则
\[
|T(x)-T(y)|\le\varepsilon|x-y|.
\]
因此 $|x_{n+1}-x_n|\le\varepsilon^n|x_1-x_0|$，对任意 $m>n$，
\[
|x_m-x_n|\le\frac{\varepsilon^n}{1-\varepsilon}|x_1-x_0|.
\]
上述估计可逐项核查。中值定理给出 $|\sin x-\sin y|\le|x-y|$，故相邻差每递推一步至少缩小到原来的 $\varepsilon$ 倍。再由
\[
|x_m-x_n|\le\sum_{k=n}^{m-1}|x_{k+1}-x_k|
\le |x_1-x_0|\sum_{k=n}^{m-1}\varepsilon^k
\le |x_1-x_0|\frac{\varepsilon^n}{1-\varepsilon},
\]
得到与 $m$ 无关、随 $n$ 趋零的界。这满足 Cauchy 收敛准则，而实数的完备性保证极限确实存在。

故 $x_n$ 为 Cauchy 列，收敛于某个 $\xi$。对递推式取极限，得 $\xi=T(\xi)$。

若还有不动点 $\eta$，则 $|\xi-\eta|\le\varepsilon|\xi-\eta|$，由 $\varepsilon<1$ 得 $\eta=\xi$。

\end{solve}

\begin{remark}[方法与条件]任意实初值都给出同一个极限，并有误差界
\[
|x_n-\xi|\le\frac{\varepsilon^n}{1-\varepsilon}|x_1-x_0|.
\]
压缩性同时保证存在、唯一与收敛速度，是一种完整的求解机制。\end{remark}

\begin{example}[来源序号3]\label{ex:pre02-math:3}\label{pr:pre02-math:3}

设 $D\subset\mathbb R^2$ 是凸区域，$f:D\to\mathbb R$ 为凸函数，即
\[
f(tu+(1-t)v)\le tf(u)+(1-t)f(v)
\qquad(u,v\in D,\ 0<t<1).
\]
判断 $f$ 是否必在 $D$ 上连续。

\par\href{https://github.com/lzxcr/cmc-mathematics/blob/64916fa1fdbc3e444a652eca204d5669058ee321/src/problems/preliminary/02/03.md}{题目与来源条目}

\end{example}

\begin{solve}\label{ans:pre02-math:3}

按“区域”指开连通集的通常约定，结论成立。固定 $x_0\in D$，取闭正方形
\[
Q=x_0+[-r,r]^2\subset D.
\]
设 $M$ 为四个顶点处函数值的最大值。每个正方形内的点都是顶点的凸组合，所以 $f\le M$ 于 $Q$。

这里 $\|h\|_\infty=\max(|h_1|,|h_2|)$。正方形内一点的两个坐标分别作区间端点的凸组合，再把两次组合相乘，就把该点写成四个顶点的凸组合；两点凸性反复使用即得到 $f\le M$。因此此处的局部上界不依赖任何尚未证明的连续性。

当 $0<\|h\|_\infty\le r$ 时，记 $t=\|h\|_\infty/r$。沿着从 $x_0$ 指向正方形边界的线段使用凸性，得
\[
f(x_0+h)-f(x_0)\le t\bigl(M-f(x_0)\bigr).
\]
对 $-h$ 应用同一估计，再由
\[
2f(x_0)\le f(x_0+h)+f(x_0-h),
\]
得到反向界。因此
\[
|f(x_0+h)-f(x_0)|
\le\frac{M-f(x_0)}r\|h\|_\infty\longrightarrow0.
\]

为看清上下界怎样配合，令 $z=x_0+h/t\in Q$，则 $x_0+h=(1-t)x_0+tz$，凸性直接给出上界。另一边由中点不等式得 $f(x_0+h)-f(x_0)\ge-[f(x_0-h)-f(x_0)]\ge-t(M-f(x_0))$。两边合起来得到绝对值估计，且 $M\ge f(x_0)$。任意 $x_0\in D$ 都成立，故 $f$ 在整个开区域内连续。

\end{solve}

\begin{remark}[方法与条件]有限维空间中，有限值凸函数在定义域的相对内部自动连续，并且局部 Lipschitz。开集条件不能直接丢掉：在 $[0,1]^2$ 上令 $f(x,y)=\mathbf1_{\{x=0\}}$，它是凸函数，却在整条边 $x=0$ 上不连续。若把题中的“区域”理解为可能包含边界的凸集，就必须作这一修正。\end{remark}

\begin{example}[来源序号4]\label{ex:pre02-math:4}\label{pr:pre02-math:4}

设 $f$ 在 $[0,1]$ 上黎曼可积，在 $x=1$ 处左可导，且 $f(1)=0$、$f'(1)=a$。证明
\[
\lim_{n\to\infty}n^2\int_0^1x^nf(x)\,dx=-a.
\]

\par\href{https://github.com/lzxcr/cmc-mathematics/blob/64916fa1fdbc3e444a652eca204d5669058ee321/src/problems/preliminary/02/04.md}{题目与来源条目}

\end{example}

\begin{solve}\label{ans:pre02-math:4}

权重 $x^n$ 把积分集中到端点 $1$，因此只应保留 $f$ 在该处的一阶行为。令
\[
g(x)=\frac{f(x)}{1-x}\quad(x<1),\qquad g(1)=-a.
\]
则 $g$ 有界，且 $g(x)\to-a$ 当 $x\to1^-$。

权重
\[
d\mu_n(x)=(n+1)(n+2)x^n(1-x)\,dx
\]
的总质量为 $1$，而对每个 $\delta>0$，
\[
\mu_n([0,1-\delta])\le(n+1)(n+2)(1-\delta)^n\longrightarrow0.
\]
归一化常数来自
\[
\int_0^1x^n(1-x)\,dx=\frac1{n+1}-\frac1{n+2}=\frac1{(n+1)(n+2)}.
\]
又因 $f(1)=0$，左导数定义给出 $f(x)/(1-x)\to-a$。在端点附近这个商因有有限极限而有界；远离端点时，$f$ 的有界性与 $1-x\ge\delta$ 给出有界性，所以 $g$ 在全区间有界。

把积分分为 $[0,1-\delta]$ 与 $[1-\delta,1]$，分别使用 $g$ 的有界性及端点极限，便得 $\int g\,d\mu_n\to-a$。于是
\[
n^2\int_0^1x^nf(x)\,dx
=\frac{n^2}{(n+1)(n+2)}\int g\,d\mu_n\longrightarrow-a.
\]

分段论证可写成明确的误差界。给定 $\eta>0$，先选 $\delta$ 使 $|g(x)+a|<\eta$ 于 $[1-\delta,1]$，再令 $K=\sup_{[0,1]}|g+a|$，则
\[
\left|\int g\,d\mu_n+a\right|
\le K\mu_n([0,1-\delta])+\eta\mu_n([1-\delta,1])
\le K\mu_n([0,1-\delta])+\eta.
\]
先令 $n\to\infty$，再令 $\eta\downarrow0$，得到所需极限。

\end{solve}

\begin{remark}[方法与条件]这里的 $\mu_n$ 是趋向端点点质量的概率分布。相同论证适用于更高阶端点展开，见\hyperref[app:integration-local]{方法卷：积分的局部化}。\end{remark}

\begin{example}[来源序号7]\label{ex:pre02-math:7}\label{pr:pre02-math:7}

设 $f$ 在 $[0,1]$ 上黎曼可积，且 $0\le f\le1$。证明对每个 $\varepsilon>0$，存在仅取 $0,1$ 的有限分段常值函数 $g$，使任意 $[\alpha,\beta]\subset[0,1]$ 都满足
\[
\left|\int_\alpha^\beta(f-g)\,dx\right|<\varepsilon.
\]

\par\href{https://github.com/lzxcr/cmc-mathematics/blob/64916fa1fdbc3e444a652eca204d5669058ee321/src/problems/preliminary/02/07.md}{题目与来源条目}

\end{example}

\begin{solve}\label{ans:pre02-math:7}

在每个小区间内保留总质量，再把它集中成取值为 $1$ 的一小段。取整数 $N>2/\varepsilon$，对 $0\le j<N$ 令
\[
m_j=\int_{j/N}^{(j+1)/N}f(x)\,dx\in[0,1/N].
\]
定义 $g$ 在 $[j/N,j/N+m_j)$ 上等于 $1$，在该小区间的其余部分等于 $0$。于是每个完整小区间都满足
\[
\int_{j/N}^{(j+1)/N}(f-g)\,dx=0.
\]
任意 $[\alpha,\beta]$ 除去完整小区间后，至多剩下两段总长不超过 $2/N$ 的边角区间。由 $|f-g|\le1$，
\[
\left|\int_\alpha^\beta(f-g)\,dx\right|\le\frac2N<\varepsilon.
\]

构造中 $m_j$ 是面积，也正是函数 $g=1$ 所占小段的长度，因此 $\int_{j/N}^{(j+1)/N}g=m_j=\int_{j/N}^{(j+1)/N}f$。任意积分区间的中间整格全部抵消，只有包含两个端点的格子可能有误差；若两端落在同一格，误差甚至不超过 $1/N$。各分点处怎样指定 $g$ 的值不影响积分，分点总数有限，故 $g$ 确为题目要求的有限分段常值函数。

\end{solve}

\begin{remark}[方法与条件]逼近的是所有区间上的积分，而不是函数值本身。例如 $f\equiv1/2$ 时，任何这样的 $g$ 都满足 $\|f-g\|_{L^1}=1/2$。细尺度上的 $0,1$ 振荡仍能在积分意义下逼近任意介于 $0$ 与 $1$ 之间的密度。\end{remark}

\begin{example}[来源序号8]\label{ex:pre02-math:8}\label{pr:pre02-math:8}

设 $\varphi:(0,\infty)\to(0,\infty)$ 严格递减且连续，满足 $\varphi(t)\to\infty$ 当 $t\to0^+$，并且
\[
\int_0^\infty\varphi(t)\,dt
=\int_0^\infty\varphi^{-1}(t)\,dt=a<\infty.
\]
证明
\[
\int_0^\infty\varphi(t)^2\,dt
+\int_0^\infty\varphi^{-1}(t)^2\,dt
\ge\frac12a^{3/2}.
\]

\par\href{https://github.com/lzxcr/cmc-mathematics/blob/64916fa1fdbc3e444a652eca204d5669058ee321/src/problems/preliminary/02/08.md}{题目与来源条目}

\end{example}

\begin{solve}\label{ans:pre02-math:8}

函数与反函数描述同一个图下区域
\[
\Omega=\{(x,y):x>0,\ 0<y<\varphi(x)\}.
\]
它的面积为 $a$，水平截面的长度是 $\varphi^{-1}(y)$。由非负函数的 Fubini 定理，题中左端等于
\[
2\iint_\Omega(x+y)\,dx\,dy.
\]
由于 $\varphi$ 正且递减，其在无穷远的极限存在；积分有限迫使该极限为 $0$。结合零端极限和连续性，$\varphi$ 的值域是 $(0,\infty)$，反函数在所需范围内有定义。对固定 $y>0$，$y<\varphi(x)$ 等价于 $x<\varphi^{-1}(y)$，因而
\[
\int_0^\infty\varphi(x)^2\,dx=2\iint_\Omega y\,dx\,dy,
\qquad
\int_0^\infty\varphi^{-1}(y)^2\,dy=2\iint_\Omega x\,dx\,dy.
\]
这些等式对非负积分成立，允许取值为无穷大。

若该积分无穷大，结论显然。否则令 $r=\sqrt{2a}$，取同面积三角形
\[
T=\{x,y>0:x+y<r\}.
\]
$T$ 内的权重 $x+y$ 不超过 $r$，外部的权重不小于 $r$，故
\[
\iint_\Omega(x+y)-\iint_T(x+y)
=\iint_{\mathbb R_+^2}(x+y-r)(\mathbf1_\Omega-\mathbf1_T)\ge0.
\]
因此实际上有更强的估计
\[
\begin{aligned}
\int_0^\infty\varphi^2+\int_0^\infty(\varphi^{-1})^2
&\ge2\iint_T(x+y)\,dx\,dy\\
&=\frac{2r^3}{3}
=\frac{4\sqrt2}{3}a^{3/2}
>\frac12a^{3/2}.
\end{aligned}
\]

比较式中的被积函数为什么非负，可以分区查看：在 $\Omega\setminus T$ 上两个因子均非负，在 $T\setminus\Omega$ 上两个因子均非正，其他地方第二个因子为零。减去常数 $r$ 不改变积分差，因为两区域面积相等。最后以 $s=x+y$ 分层，$\iint_T(x+y)\,dx\,dy=\int_0^r s\cdot s\,ds=r^3/3$，于是得到所列常数。

\end{solve}

\begin{remark}[方法与条件]固定质量后，优先占据权重最小的区域，就使加权积分最小。这是一种基本的重排原则。此处把反函数积分转成同一个区域的另一方向截面，几何结构便直接显现；两次面积积分相等也由此自然成立。\end{remark}

\section{初赛：第 3 届（2011 年）}

本节选入4道课程相关题；来源序号保留原题解编排，未选入题目不重新补号。

\begin{example}[来源序号2]\label{ex:pre03-math:2}\label{pr:pre03-math:2}

设 $f_1,\ldots,f_n$ 是 $[0,1]$ 上的非负连续函数。证明存在 $\xi\in[0,1]$，使
\[
\prod_{k=1}^nf_k(\xi)
\le\prod_{k=1}^n\int_0^1f_k(x)\,dx.
\]

\par\href{https://github.com/lzxcr/cmc-mathematics/blob/64916fa1fdbc3e444a652eca204d5669058ee321/src/problems/preliminary/03/02.md}{题目与来源条目}

\end{example}

\begin{solve}\label{ans:pre03-math:2}

先取乘积的 $n$ 次方根，使每个因子恰好适用 Hölder 不等式。令
\[
h(x)=\left(\prod_{k=1}^nf_k(x)\right)^{1/n}.
\]
由 Hölder 不等式，
\[
\int_0^1h(x)\,dx
\le\prod_{k=1}^n\left(\int_0^1f_k(x)\,dx\right)^{1/n}.
\]
连续函数 $h$ 在区间上必有一点不超过其平均值。取这样的 $\xi$，再把所得不等式取 $n$ 次方即可。

当 $n=1$ 时结论就是积分平均值的基本性质。$n\ge2$ 时，将 Hölder 不等式用于 $n$ 个函数 $f_k^{1/n}$，各指数都取 $n$，其倒数和为 $1$，便得到上面的乘积积分界。另一方面，$h$ 连续且非负，在闭区间取得最小值 $h(\xi)$，从而 $h(\xi)\le\int_0^1h$。把这两个不等式相接并取 $n$ 次方即可；即使某个积分为零也不必除以它，因此退化情形同样包含在证明中。

\end{solve}

\begin{remark}[方法与条件]若使用一般区间 $[a,b]$，右端应写成各函数平均值的乘积，即每个积分都除以 $b-a$。本题的自然对象是概率测度下的平均值，而非未经归一化的积分。\end{remark}

\begin{example}[来源序号4]\label{ex:pre03-math:4}\label{pr:pre03-math:4}

设 $A,B,C$ 是三角形的三个内角。求
\[
3\sin A+4\sin B+18\sin C
\]
的最大值。

\par\href{https://github.com/lzxcr/cmc-mathematics/blob/64916fa1fdbc3e444a652eca204d5669058ee321/src/problems/preliminary/03/04.md}{题目与来源条目}

\end{example}

\begin{solve}\label{ans:pre03-math:4}

固定 $C$ 后，用 $B=\pi-C-A$ 合并前两项，再用 Cauchy--Schwarz 不等式：
\[
\begin{aligned}
3\sin A+4\sin B
&=(3+4\cos C)\sin A+4\sin C\cos A\\
&\le\sqrt{25+24\cos C}.
\end{aligned}
\]
故令 $c=\cos C$，只需考察
\[
H(c)=\sqrt{25+24c}+18\sqrt{1-c^2},\qquad -1<c<1.
\]
这里先对所有合法三角形给出统一上界，再检验上界能否实现；固定某个 $C$ 时 Cauchy--Schwarz 的等号不一定都可行，这不影响求上界。二阶导数为
\[
H''(c)=-\frac{144}{(25+24c)^{3/2}}-\frac{18}{(1-c^2)^{3/2}}<0.
\]
因此 $H'$ 严格递减，只要找到一个零点，该点左侧函数递增、右侧函数递减，便确定了全局最大值。

$H$ 严格凹，而且
\[
H'(c)=\frac{12}{\sqrt{25+24c}}-\frac{18c}{\sqrt{1-c^2}},
\qquad H'(1/8)=0.
\]
所以其唯一最大点为 $c=1/8$，最大值为
\[
H(1/8)=2\sqrt7+\frac{27\sqrt7}{4}
=\boxed{\frac{35\sqrt7}{4}}.
\]

取锐角 $A,B,C$ 满足
\[
\cos A=\frac34,\qquad \cos B=\frac9{16},\qquad \cos C=\frac18.
\]
此时 $\cos(A+B)=-1/8$，故 $A+B+C=\pi$，且前述 Cauchy--Schwarz 不等式取等号，最大值确能达到。

\end{solve}

\begin{remark}[方法与条件]合并同一角的正弦、余弦先消去一个变量；剩下的函数具有凹性，验证一个驻点就足够。凹性同时提供全局最优性与唯一性，免去了平方方程后逐一排除多余根的过程。\end{remark}

\begin{example}[来源序号5]\label{ex:pre03-math:5}\label{pr:pre03-math:5}

证明：对任意实数 $\alpha$，存在仅取 $-1,1$ 的数列 $(a_n)$，使
\[
\lim_{n\to\infty}
\left(\sum_{k=1}^n\sqrt{n+a_k}-n^{3/2}\right)=\alpha.
\]

\par\href{https://github.com/lzxcr/cmc-mathematics/blob/64916fa1fdbc3e444a652eca204d5669058ee321/src/problems/preliminary/03/05.md}{题目与来源条目}

\end{example}

\begin{solve}\label{ans:pre03-math:5}

Taylor 展开把目标转化为控制符号的部分和。对 $a_k=\pm1$ 一致地有
\[
\sqrt{n+a_k}=\sqrt n+\frac{a_k}{2\sqrt n}+O(n^{-3/2}),
\]
故若 $S_n=\sum_{k=1}^na_k$，则
\[
\sum_{k=1}^n\sqrt{n+a_k}-n^{3/2}
=\frac{S_n}{2\sqrt n}+O(n^{-1/2}).
\]
展开的余项必须对所有 $k\le n$ 使用同一个界。因为 $a_k/n$ 只取 $\pm1/n$，函数 $\sqrt{1+u}$ 的二阶导数在 $|u|\le1/2$ 上有界，所以单项误差至多为 $Cn^{-3/2}$。求和后误差至多为 $Cn^{-1/2}\to0$，不会因项数增加而破坏所求极限。

因此只需构造 $S_n=2\alpha\sqrt n+O(1)$。

设 $p_n$ 为前 $n$ 项中 $1$ 的个数，应有 $S_n=2p_n-n$。自然地取
\[
p_n=\left\lfloor\frac n2+\alpha\sqrt n\right\rfloor
\qquad(n\ge N),
\]
其中 $N$ 足够大，使 $0\le p_N\le N$，且
\[
0<\frac12+\frac{\alpha}{\sqrt{n+1}+\sqrt n}<1
\qquad(n\ge N).
\]
这保证 $p_{n+1}-p_n\in\{0,1\}$。先任取含 $p_N$ 个 $1$ 的前 $N$ 项，再令
\[
a_{n+1}=2(p_{n+1}-p_n)-1.
\]
所得数列满足 $S_n=2p_n-n=2\alpha\sqrt n+O(1)$，代回即得结论。

记 $r_n=n/2+\alpha\sqrt n$。上面选取的 $N$ 保证 $0<r_{n+1}-r_n<1$，故相邻整数部分只能相同或相差 $1$；这恰好意味着每一步能够追加一个负号或正号。又 $-1<\lfloor r_n\rfloor-r_n\le0$，从而 $|S_n-2\alpha\sqrt n|<2$，于是 $S_n/(2\sqrt n)\to\alpha$。有限个初始符号只负责让构造从合法计数开始，不另产生渐近误差。

\end{solve}

\begin{remark}[方法与条件]符号和既有整数约束，又有与 $n$ 同奇偶的约束；改为记录正号个数后，两种约束同时由取整处理。相同办法可以跟踪许多相邻增量小于 $1$ 的缓慢变化目标。\end{remark}

\begin{example}[来源序号7]\label{ex:pre03-math:7}\label{pr:pre03-math:7}

设 $F:[0,\infty)\to\mathbb R$ 单调递减，且 $F(t)\to0$ 当 $t\to\infty$。已知
\[
\lim_{n\to\infty}\int_0^\infty F(t)\sin\frac tn\,dt=0.
\]
证明：

（1） $tF(t)\to0$ 当 $t\to\infty$；

（2） $\displaystyle\lim_{s\to0}\int_0^\infty F(t)\sin(st)\,dt=0$。

\par\href{https://github.com/lzxcr/cmc-mathematics/blob/64916fa1fdbc3e444a652eca204d5669058ee321/src/problems/preliminary/03/07.md}{题目与来源条目}

\end{example}

\begin{solve}\label{ans:pre03-math:7}

记 $S(s)=\int_0^\infty F(t)\sin(st)\,dt$。由单调性及极限条件，$F\ge0$；对 $s\ne0$，积分由 Dirichlet 判别法收敛。

\textbf{（1）} 单调性使相邻的正、负半波配对后非负。换元 $t=nu$，并在第一对半波内只保留 $0\le u\le\pi/2$，得到
\[
\begin{aligned}
S(1/n)
&=n\sum_{k=0}^\infty\int_0^\pi
\bigl[F(n(2k\pi+u))-F(n((2k+2)\pi-u))\bigr]\sin u\,du\\
&\ge n\bigl[F(\pi n/2)-F(3\pi n/2)\bigr]\ge0.
\end{aligned}
\]
配对时将负半波的变量反向写成 $2\pi-u$，才得到同一个正权重 $\sin u$。在第一对半波的 $0\le u\le\pi/2$ 上，$nu\le\pi n/2$ 且 $n(2\pi-u)\ge3\pi n/2$，由递减性得括号至少为 $F(\pi n/2)-F(3\pi n/2)$；同时 $\int_0^{\pi/2}\sin u\,du=1$。余下配对项均非负，可以舍去。此配对只取完整周期的部分积分，再利用已知收敛性取极限，并没有对条件收敛的积分任意重排。

令 $c=\pi/2$、$\varepsilon_n=\sup_{m\ge n}S(1/m)$，则 $\varepsilon_n\to0$。将上述差分估计用于 $n,3n,3^2n,\ldots$，再望远镜求和：
\[
\begin{aligned}
nF(cn)
&=\sum_{k=0}^\infty n\bigl[F(c3^kn)-F(c3^{k+1}n)\bigr]\\
&\le\sum_{k=0}^\infty\frac{\varepsilon_n}{3^k}
=\frac32\varepsilon_n\longrightarrow0.
\end{aligned}
\]
对于 $cn\le t<c(n+1)$，有 $0\le tF(t)\le c(n+1)F(cn)$，所以 $tF(t)\to0$。

\textbf{（2）} 只需考虑 $s>0$，因为 $S$ 为奇函数。取 $R=1/s$。在 $[0,R]$ 上用 $|\sin(st)|\le st$；在尾部用单调函数的 Dirichlet 余项估计，得到
\[
|S(s)|\le\frac1R\int_0^R tF(t)\,dt+2RF(R).
\]
第一项是趋零函数 $tF(t)$ 的 Cesàro 平均，第二项也由 （1） 趋于零，故 $S(s)\to0$。

尾项界的依据是：$\int_R^u\sin(st)\,dt$ 的绝对值不超过 $2/s$，对递减的 $F$ 使用分部积分的单调函数形式，便有 $|\int_R^\infty F(t)\sin(st)\,dt|\le2F(R)/s$。而若 $h(t)=tF(t)\to0$，固定足够大的 $A$ 后，$R^{-1}\int_0^Rh\le R^{-1}\int_0^A h+\sup_{t\ge A}h(t)$，先令 $R\to\infty$ 再令 $A\to\infty$，就证明了这里使用的平均值极限。

\end{solve}

\begin{remark}[方法与条件]这是一个 Tauber 型论证：变换在一列趋零频率上的信息，经由单调性，反推出原函数的尾部衰减，再得到全部趋零频率上的极限。正负半波配对提供非负性，几何尺度上的望远镜求和负责把局部差分控制提升为尾部控制。\end{remark}

\section{初赛：第 4 届（2012 年）}

本节选入4道课程相关题；来源序号保留原题解编排，未选入题目不重新补号。

\begin{example}[来源序号2]\label{ex:pre04-math:2}\label{pr:pre04-math:2}

设 $\Gamma$ 为抛物线，点 $P$ 位于抛物线与焦点同侧的开区域内。过 $P$ 的直线 $L$ 与 $\Gamma$ 围成有界区域时，记其面积为 $A(L)$。证明：$A(L)$ 最小，当且仅当 $P$ 是 $L$ 被 $\Gamma$ 截出的弦的中点。

\par\href{https://github.com/lzxcr/cmc-mathematics/blob/64916fa1fdbc3e444a652eca204d5669058ee321/src/problems/preliminary/04/02.md}{题目与来源条目}

\end{example}

\begin{solve}\label{ans:pre04-math:2}

可逆仿射变换保持中点，并将所有面积乘以同一正常数，所以可先将抛物线化为 $y=x^2$。若 $P=(x_0,y_0)$，则 $h=y_0-x_0^2>0$。再作保持该抛物线的剪切平移
\[
u=x-x_0,\qquad v=y-2x_0x+x_0^2,
\]
便将 $P$ 化为 $(0,h)$。

能围出有界区域的直线均可写成 $v=mu+h$。它与抛物线间的竖直间隙为
\[
mu+h-u^2=R^2-(u-m/2)^2,
\qquad R^2=h+\frac{m^2}{4}.
\]
因此
\[
A(m)=\int_{-R}^{R}(R^2-t^2)\,dt
=\frac43\left(h+\frac{m^2}{4}\right)^{3/2},
\]
唯一最小值在 $m=0$ 处取得。

两交点的 $u$ 坐标之和为 $m$，故弦中点的 $u$ 坐标为 $m/2$；由于 $P$ 与中点都在 $L$ 上，二者相同恰等价于 $m=0$。结论得证。

剪切变换的线性部分行列式为 $1$，且将 $y=x^2$ 化为 $v=u^2$，所以它确实同时保持面积和抛物线。交点满足 $u^2-mu-h=0$，根为 $m/2\pm R$；换元 $t=u-m/2$ 后积分范围才成为 $[-R,R]$。由于 $h>0$，每个有限斜率都给出两个交点，而竖直线只能与抛物线相交一次，不能围成题设区域。这样考察 $m\in\mathbb R$ 已穷尽所有候选直线。

\end{solve}

\begin{remark}[方法与条件]剪切把给定点移到抛物线的对称轴上，同时保留抛物线与面积；最优弦的中点条件因此成为显式的对称性。这里使用的是仿射几何中的面积与中点，而非角度或长度。\end{remark}

\begin{example}[来源序号3]\label{ex:pre04-math:3}\label{pr:pre04-math:3}

设 $f\in C^1([0,\infty))$，$f(0)>0$，且 $f'\ge0$。若
\[
\int_0^\infty\frac{dx}{f(x)+f'(x)}<\infty,
\]
证明
\[
\int_0^\infty\frac{dx}{f(x)}<\infty.
\]

\par\href{https://github.com/lzxcr/cmc-mathematics/blob/64916fa1fdbc3e444a652eca204d5669058ee321/src/problems/preliminary/04/03.md}{题目与来源条目}

\end{example}

\begin{solve}\label{ans:pre04-math:3}

比较两者的差，并把它控制为一个可积的全导数。由于 $f\ge f(0)>0$、$f'\ge0$，
\[
0\le\frac1f-\frac1{f+f'}
=\frac{f'}{f(f+f')}
\le\frac{f'}{f^2}
=-\left(\frac1f\right)'.
\]
在 $[0,R]$ 上积分，得到
\[
\int_0^R\frac{dx}{f(x)}
\le\int_0^R\frac{dx}{f(x)+f'(x)}
+\frac1{f(0)}-\frac1{f(R)}
\le\int_0^\infty\frac{dx}{f(x)+f'(x)}+\frac1{f(0)}.
\]
令 $R\to\infty$ 即得结论。

仅由 $1/(f+f')\le1/f$，只能知道已知收敛积分较小，不能推出较大的积分收敛，所以要额外估计二者之差。最后令 $J(R)=\int_0^R1/f$，因为被积函数正，$J(R)$ 随 $R$ 单调递增；上式又给出不依赖 $R$ 的有限上界，故 $J(R)$ 存在有限极限。这才是由有限区间估计推出反常积分收敛的完整依据。

\end{solve}

\begin{remark}[方法与条件]题设方向的点态比较本身不足以证明收敛；两项的差却有一个统一可积的上界。识别全导数使“分母增加一个导数”所带来的总改变量被端点值控制。\end{remark}

\begin{example}[来源序号5]\label{ex:pre04-math:5}\label{pr:pre04-math:5}

设 $n$ 为正整数，且
\[
\frac{(1+x)^n}{(1-x)^3}=\sum_{i=0}^\infty a_ix^i,
\qquad |x|<1.
\]
求 $\sum_{i=0}^{n-1}a_i$。

\par\href{https://github.com/lzxcr/cmc-mathematics/blob/64916fa1fdbc3e444a652eca204d5669058ee321/src/problems/preliminary/04/05.md}{题目与来源条目}

\end{example}

\begin{solve}\label{ans:pre04-math:5}

求系数的部分和，等价于先给生成函数乘上 $(1-x)^{-1}$。用 $[x^m]$ 表示取 $x^m$ 的系数，则
\[
S=\sum_{i=0}^{n-1}a_i
=[x^{n-1}]\frac{(1+x)^n}{(1-x)^4}
=\sum_{k=0}^n\binom nk\binom{k+2}{3}.
\]
最后一式由二项式展开并将指标 $k$ 换为 $n-k$ 得到，$k=0$ 项为零。

为展开系数提取，先注意 $(1-x)^{-1}=\sum_{j\ge0}x^j$，所以其与 $\sum_{i\ge0}a_ix^i$ 的 $x^{n-1}$ 系数正是所求部分和。又
\[
(1-x)^{-4}=\sum_{j\ge0}\binom{j+3}{3}x^j,
\]
于是该系数为 $\sum_{k=0}^{n-1}\binom nk\binom{n-k+2}{3}$。令新指标等于 $n-k$ 并利用 $\binom n{n-k}=\binom nk$，即可得到上面使用的求和式。

利用
\[
\binom{k+2}{3}=\binom k3+2\binom k2+\binom k1,
\qquad
\sum_{k=0}^n\binom nk\binom kr=\binom nr2^{n-r},
\]
可得
\[
\begin{aligned}
S
&=\binom n3 2^{n-3}+2\binom n2 2^{n-2}+n2^{n-1}\\
&=\boxed{\frac{n(n+2)(n+7)}3\,2^{n-4}}.
\end{aligned}
\]
这里约定 $r>n$ 时 $\binom nr=0$，故公式对所有正整数 $n$ 都成立。

\end{solve}

\begin{remark}[方法与条件]第二个恒等式先选出 $r$ 个必选元素，再任取其余元素即可解释。生成函数将累加化为乘法，二项式基底又将多项式权重化为简单的组合计数。\end{remark}

\begin{example}[来源序号6]\label{ex:pre04-math:6}\label{pr:pre04-math:6}

设 $f:[0,1]\to\mathbb R$ 可微，满足
\[
f(0)=f(1),\qquad \int_0^1f(x)\,dx=0,
\qquad f'(x)\ne1\quad(x\in[0,1]).
\]
证明对每个正整数 $n$，
\[
\left|\sum_{k=0}^{n-1}f\left(\frac kn\right)\right|<\frac12.
\]

\par\href{https://github.com/lzxcr/cmc-mathematics/blob/64916fa1fdbc3e444a652eca204d5669058ee321/src/problems/preliminary/04/06.md}{题目与来源条目}

\end{example}

\begin{solve}\label{ans:pre04-math:6}

由 Rolle 定理，某处有 $f'=0$。导数具有 Darboux 性质，又处处不取 $1$，故 $f'<1$。于是
\[
g(x)=f(x)-x
\]
严格递减，且
\[
\int_0^1g(x)\,dx=-\frac12,\qquad g(0)-g(1)=1.
\]
详细说，先由 $f(0)=f(1)$ 得某个 $\eta$ 满足 $f'(\eta)=0$；如果另有一点 $f'>1$，导数的介值性将迫使中间某处 $f'=1$，与题设矛盾。因此处处 $f'<1$，中值定理进一步给出 $g$ 严格递减。对每一小区间 $[(k-1)/n,k/n]$，内部的函数值严格介于两个端点值之间，积分后得到严格的左右和不等式，不能只写成非严格比较。

严格递减函数的左、右端点 Riemann 和严格夹住积分：
\[
\frac1n\sum_{k=1}^{n}g(k/n)
< -\frac12
<\frac1n\sum_{k=0}^{n-1}g(k/n).
\]
两端的差恰为 $1/n$，所以
\[
-\frac n2<\sum_{k=0}^{n-1}g(k/n)<-\frac n2+1.
\]
加上 $\sum_{k=0}^{n-1}k/n=(n-1)/2$，得到
\[
-\frac12<\sum_{k=0}^{n-1}f(k/n)<\frac12.
\]

\end{solve}

\begin{remark}[方法与条件]导数未必连续，但 Darboux 定理仍保证它不能越过未曾取到的值。减去线性函数后，题目化为单调函数左右和之间的面积估计；端点条件精确确定了两种和的差。\end{remark}

\section{初赛：第 5 届（2013 年）}

本节选入3道课程相关题；来源序号保留原题解编排，未选入题目不重新补号。

\begin{example}[来源序号3]\label{ex:pre05-math:3}\label{pr:pre05-math:3}

设 $f\in C^2([0,a])$，$f'(0)=1$、$f''(0)\ne0$，且 $0<f(x)<x$ 对 $0<x<a$ 成立。令 $x_{n+1}=f(x_n)$，其中 $x_1\in(0,a)$。

（1）证明 $x_n$ 收敛并求极限；（2）判断 $nx_n$ 是否收敛，若收敛，求其极限。

\par\href{https://github.com/lzxcr/cmc-mathematics/blob/64916fa1fdbc3e444a652eca204d5669058ee321/src/problems/preliminary/05/03.md}{题目与来源条目}

\end{example}

\begin{solve}\label{ans:pre05-math:3}

（1）$0<x_{n+1}<x_n\le x_1<a$，故 $x_n\to\ell\ge0$。连续性给 $f(\ell)=\ell$，而区间 $(0,a)$ 内没有不动点，故 $\ell=0$；同时由题设与连续性知 $f(0)=0$。

（2）Taylor 展开为
\[
f(x)=x+\frac{f''(0)}2x^2+o(x^2).
\]
由 $f(x)<x$ 及 $f''(0)\ne0$ 知 $f''(0)<0$。因此
\[
\frac1{x_{n+1}}-\frac1{x_n}
=\frac{x_n-f(x_n)}{x_nf(x_n)}
\longrightarrow-\frac{f''(0)}2.
\]
令 $z_n=1/x_n$。上面的差分极限利用了 $f(x_n)/x_n\to f'(0)=1$，因此分母 $x_nf(x_n)$ 与 $x_n^2$ 等价。若记 $L=-f''(0)/2>0$，则
\[
\frac{z_n}{n}=\frac{z_1}{n}+\frac1n\sum_{j=1}^{n-1}(z_{j+1}-z_j)\longrightarrow L.
\]
这里用的是收敛数列的算术平均仍趋于同一极限；因为 $L>0$，最后才可以取倒数。

对差分求平均，得 $1/(nx_n)\to-f''(0)/2$，从而
\[
\boxed{\displaystyle nx_n\longrightarrow-\frac2{f''(0)}}.
\]

\end{solve}

\begin{remark}[方法与条件]不动点的线性项等于恒等映射时，二次项决定收敛速度。取倒数将 $x_n-x_{n+1}$ 的二阶衰减变为近似常量的增量；一般阶数的处理见\hyperref[app:estimate-recursion]{递推与渐近尺度}。\end{remark}

\begin{example}[来源序号4]\label{ex:pre05-math:4}\label{pr:pre05-math:4}

设 $a>1$，可微函数 $f:(0,\infty)\to(0,\infty)$。证明存在 $x_n\to\infty$，使 $f'(x_n)<f(ax_n)$ 对每个 $n$ 成立。

\par\href{https://github.com/lzxcr/cmc-mathematics/blob/64916fa1fdbc3e444a652eca204d5669058ee321/src/problems/preliminary/05/04.md}{题目与来源条目}

\end{example}

\begin{solve}\label{ans:pre05-math:4}

若结论不成立，则有 $M>0$，使所有 $x\ge M$ 均满足 $f'(x)\ge f(ax)>0$，故 $f$ 在 $[M,\infty)$ 上递增。

取 $x$ 使 $x/a\ge M$。在 $[x/a,x]$ 上用中值定理，存在 $\xi\in(x/a,x)$，满足
\[
f(x)-f(x/a)=(x-x/a)f'(\xi)
\ge(1-a^{-1})x f(a\xi)
\ge(1-a^{-1})x f(x).
\]
左端严格小于 $f(x)$，所以 $(1-a^{-1})x<1$，这与 $x$ 可任意大矛盾。因此所需不等式成立的点集无上界，从中选取 $x_n\ge n$ 即可。

反设的逻辑是：若不等式成立的点集无上界，就能每次选一个不小于 $n$ 的点，直接得到所需数列。因此否定结论意味着从某个位置开始不等式处处失败。中值点满足 $a\xi>x$，而这些点全部落在递增区间内，所以 $f(a\xi)\ge f(x)$。再用 $f(x/a)>0$ 得到严格上界，除以正数 $f(x)$ 后便出现不可能对任意大 $x$ 成立的线性约束。

\end{solve}

\begin{remark}[方法与条件]反证假设先给出单调性，再使整个区间内的导数都受同一个端点函数值控制。区间长度随 $x$ 增长，由此产生矛盾，无须假定 $f'$ 连续或可积分。\end{remark}

\begin{example}[来源序号5]\label{ex:pre05-math:5}\label{pr:pre05-math:5}

设 $f:[-1,1]\to\mathbb R$ 为偶函数，且在 $[0,1]$ 上单调递增；$g:[-1,1]\to\mathbb R$ 为凸函数，即
\[
g(tx+(1-t)y)\le tg(x)+(1-t)g(y)
\]
对 $x,y\in[-1,1]$、$0<t<1$ 成立。证明
\[
2\int_{-1}^1f(x)g(x)\,\mathrm dx
\ge\left(\int_{-1}^1f(x)\,\mathrm dx\right)
\left(\int_{-1}^1g(x)\,\mathrm dx\right).
\]

\par\href{https://github.com/lzxcr/cmc-mathematics/blob/64916fa1fdbc3e444a652eca204d5669058ee321/src/problems/preliminary/05/05.md}{题目与来源条目}

\end{example}

\begin{solve}\label{ans:pre05-math:5}

令 $h(x)=[g(x)+g(-x)]/2$。它是偶凸函数，因而在 $[0,1]$ 上递增：对 $0\le x\le y$，$x$ 是 $y,-y$ 的凸组合，故 $h(x)\le h(y)$。

于是对 $x,y\in[0,1]$ 有 $(f(x)-f(y))(h(x)-h(y))\ge0$。在单位正方形积分得到
\[
\int_0^1fh\ge\left(\int_0^1f\right)\left(\int_0^1h\right).
\]
二重积分展开的四项分别是 $\int fh$、$\int fh$、$-(\int f)(\int h)$ 和 $-(\int f)(\int h)$，故非负性恰好给出上述不等式。对 $y>0$，还可明确写出 $x=\frac{y+x}{2y}y+\frac{y-x}{2y}(-y)$，由凸性与偶性得到 $h(x)\le h(y)$；$y=0$ 时无需论证。对称化的意义在于，$g$ 未必单调，但与偶函数 $f$ 积分时，$g$ 的奇部自动抵消。

再由偶性，$\int_{-1}^1fg=2\int_0^1fh$、$\int_{-1}^1f=2\int_0^1f$、$\int_{-1}^1g=2\int_0^1h$，即得结论。上述函数均有界且 Riemann 可积，故积分合法。

\end{solve}

\begin{remark}[方法与条件]先对称化，便能把凸性转换为同向单调性；剩下的是非负协方差恒等式。相关的凸性与积分不等式见\hyperref[app:estimate-convexity]{凸性方法}。\end{remark}

\section{初赛：第 6 届（2014 年）}

本节选入3道课程相关题；来源序号保留原题解编排，未选入题目不重新补号。

\begin{example}[来源序号2]\label{ex:pre06-math:2}\label{pr:pre06-math:2}

设 $f\in C([0,1])$ 非负且严格递增。

（1）证明每个正整数 $n$ 都对应唯一 $x_n\in[0,1]$，使
\[
f(x_n)^n=\int_0^1f(x)^n\,\mathrm dx.
\]
（2）证明 $x_n\to1$。

\par\href{https://github.com/lzxcr/cmc-mathematics/blob/64916fa1fdbc3e444a652eca204d5669058ee321/src/problems/preliminary/06/02.md}{题目与来源条目}

\end{example}

\begin{solve}\label{ans:pre06-math:2}

（1）$f^n$ 连续且严格递增，其积分严格介于端点值之间，故由介值定理存在唯一的 $x_n\in(0,1)$。

（2）任取 $0<\varepsilon<1$，有
\[
\varepsilon^{1/n}f(1-\varepsilon)
\le\left(\int_0^1f^n\right)^{1/n}=f(x_n)\le f(1).
\]
先令 $n\to\infty$，再令 $\varepsilon\downarrow0$，得 $f(x_n)\to f(1)$。严格单调连续函数的逆函数连续，故 $x_n\to1$。

第二问的下界只保留末端长度为 $\varepsilon$ 的区间：$\int_0^1 f^n\ge\int_{1-\varepsilon}^1f^n\ge\varepsilon f(1-\varepsilon)^n$，取 $n$ 次方根即得。不使用反函数也能收尾：若 $x_n$ 不趋于 $1$，就有某个 $\eta>0$ 及无穷多个 $n$ 满足 $x_n\le1-\eta$，于是 $f(x_n)\le f(1-\eta)<f(1)$，与已证极限矛盾。

\end{solve}

\begin{remark}[方法与条件]这里先证明 $L^n$ 范数趋向最大值，再由函数的可逆性恢复取值点。更一般地，在有限测度空间上，$\|f\|_p\to\|f\|_\infty$；“最大值附近有正测度”是本题下界的关键。\end{remark}

\begin{example}[来源序号4]\label{ex:pre06-math:4}\label{pr:pre06-math:4}

设 $f$ 在 $\mathbb R$ 上二阶可微，且 $f,f',f''$ 均严格为正。若存在 $a,b>0$，使
\[
f''(x)\le af(x)+bf'(x),\qquad x\in\mathbb R,
\]
证明：（1）$\lim_{x\to-\infty}f'(x)=0$；（2）存在常数 $c$，使 $f'\le cf$；（3）求由 $a,b$ 保证的最小普适常数 $c$。

\par\href{https://github.com/lzxcr/cmc-mathematics/blob/64916fa1fdbc3e444a652eca204d5669058ee321/src/problems/preliminary/06/04.md}{题目与来源条目}

\end{example}

\begin{solve}\label{ans:pre06-math:4}

（1）$f,f'$ 均递增，故在 $(-\infty,0]$ 上有界，并各有有限的左端极限。中值定理给出
\[
0<f'(x)\le f(x+1)-f(x)\longrightarrow0.
\]

（2）令
\[
c=\frac{b+\sqrt{b^2+4a}}2,\qquad q=c-b>0.
\]
由 $c^2-bc-a=0$，对 $h=f'-cf$ 有 $h'\le-qh$，故 $e^{qx}h(x)$ 递减。它在 $x\to-\infty$ 时趋于 $0$，因为 $f,f'$ 在左半轴有界。因此 $h\le0$，即 $f'\le cf$。

（3）函数 $f(x)=e^{cx}$ 满足全部条件并处处取等，所以此 $c$ 不能减小。

积分因子的常数来自消去 $f$ 项：$h'=f''-cf'\le af+(b-c)f'=af-qf'$，而 $qc=c(c-b)=a$，所以右端正是 $-q(f'-cf)=-qh$。于是 $(e^{qx}h)'\le0$。对任意 $t<x$，$e^{qx}h(x)\le e^{qt}h(t)$；让 $t\to-\infty$ 后右端为零，便得 $h(x)\le0$。这种从左端极限控制函数符号的论证不要求 $f''$ 连续。

\end{solve}

\begin{remark}[方法与条件]二阶微分不等式通过特征方程分解成一阶单调量。第三问求的是对全部合条件函数同时有效的常数；对某个固定函数，其最小常数为 $\sup_x f'(x)/f(x)$，可能更小。\end{remark}

\begin{example}[来源序号6]\label{ex:pre06-math:6}\label{pr:pre06-math:6}

设 $0<\alpha<1$，正数列 $a_n$ 满足
\[
\liminf_{n\to\infty}n^\alpha\left(\frac{a_n}{a_{n+1}}-1\right)
=\lambda\in(0,\infty).
\]
证明对每个 $k>0$ 均有 $n^ka_n\to0$。

\par\href{https://github.com/lzxcr/cmc-mathematics/blob/64916fa1fdbc3e444a652eca204d5669058ee321/src/problems/preliminary/06/06.md}{题目与来源条目}

\end{example}

\begin{solve}\label{ans:pre06-math:6}

取 $0<c<\lambda$，则充分大时 $a_n/a_{n+1}\ge1+cn^{-\alpha}$。再用 $\log(1+u)\ge u/2$（$0\le u\le1$），得到
\[
\log a_n\le C-\frac c2\sum_{j=N}^{n-1}j^{-\alpha}
\le C'-c'n^{1-\alpha}
\]
其中 $c'>0$。于是
\[
0<n^ka_n\le e^{C'}\exp\{k\log n-c'n^{1-\alpha}\}\longrightarrow0.
\]

由下极限等于正数 $\lambda$，任取 $c<\lambda$ 后，所需下界对全部充分大的下标成立，而不仅仅对子列成立。将 $\log a_{j+1}\le\log a_j-\log(1+cj^{-\alpha})$ 从 $N$ 累加到 $n-1$，即得到对数估计；再用 $\sum_{j=N}^{n-1}j^{-\alpha}\ge\int_N^n t^{-\alpha}\,dt=(n^{1-\alpha}-N^{1-\alpha})/(1-\alpha)$。由于 $\log n/n^{1-\alpha}\to0$，任何固定幂次 $n^k$ 都无法抵消这个指数衰减。

\end{solve}

\begin{remark}[方法与条件]取对数将乘法衰减变成可累积的加法衰减。证明实际给出比任意幂更快的上界 $a_n\le C e^{-c'n^{1-\alpha}}$；只需题设下极限为正，不必要求它有限。\end{remark}

\section{初赛：第 7 届（2015 年）}

本节选入3道课程相关题；来源序号保留原题解编排，未选入题目不重新补号。

\begin{example}[来源序号4]\label{ex:pre07-math:4}\label{pr:pre07-math:4}

正数列满足
\[
a_{n+1}=a_n+\frac n{a_n},\qquad a_1>0.
\]
证明极限 $\lim_{n\to\infty}n(a_n-n)$ 存在。

\par\href{https://github.com/lzxcr/cmc-mathematics/blob/64916fa1fdbc3e444a652eca204d5669058ee321/src/problems/preliminary/07/04.md}{题目与来源条目}

\end{example}

\begin{solve}\label{ans:pre07-math:4}

由 $a_2=a_1+a_1^{-1}\ge2$，令 $d_n=a_n-n$，则对 $n\ge2$，
\[
d_{n+1}=d_n\left(1-\frac1{n+d_n}\right).
\]
故 $0\le d_{n+1}\le d_n\le d_2$。若 $d_2=0$，此后 $a_n=n$，所求极限为零。

若 $d_2>0$，令 $c_n=nd_n>0$，则
\[
\frac{c_{n+1}}{c_n}
=1+\frac{d_n-1}{n(n+d_n)}=1+u_n,
\qquad |u_n|\le\frac{d_2+1}{n^2}.
\]
因此 $\sum_n|u_n|<\infty$，且每个因子 $1+u_n$ 为正。充分大时 $|\log(1+u_n)|\le2|u_n|$，所以 $\log c_n$ 收敛到有限值，$c_n$ 收敛到一个正数。

先减去 $n$，是因为递推增量 $n/a_n$ 若趋近 $1$，主项就应当是 $n$；这里不预先假定此猜测正确，而直接计算出误差的精确递推。接着乘以 $n$ 后，一阶的 $1/n$ 变化互相抵消，只留下可求和的 $O(n^{-2})$ 相对变化。具体地，$\log c_n=\log c_2+\sum_{j=2}^{n-1}\log(1+u_j)$；尾部绝对收敛、前面有限项均有定义，故其极限为有限实数。指数函数的连续性给出 $c_n$ 的正极限，并非只得到有界性。

\end{solve}

\begin{remark}[方法与条件]减去精确解 $a_n=n$，再乘预期尺度 $n$，就得到偏离 $1$ 仅为可求和小量的比值。无需求出依赖 $a_1$ 的极限常数；证明还说明除 $a_1=1$ 外，极限严格为正。\end{remark}

\begin{example}[来源序号5]\label{ex:pre07-math:5}\label{pr:pre07-math:5}

设 $f$ 为 $[0,\infty)$ 上的有界连续函数，$h$ 连续且
\[
\int_0^\infty|h(t)|\,\mathrm dt=a<1.
\]
定义 $g_0=f$ 及
\[
g_n(x)=f(x)+\int_0^xh(t)g_{n-1}(t)\,\mathrm dt,\qquad n\ge1.
\]
证明 $g_n$ 收敛到连续函数，并求极限函数。

\par\href{https://github.com/lzxcr/cmc-mathematics/blob/64916fa1fdbc3e444a652eca204d5669058ee321/src/problems/preliminary/07/05.md}{题目与来源条目}

\end{example}

\begin{solve}\label{ans:pre07-math:5}

在有界连续函数空间上定义 $Tu(x)=\int_0^xh(t)u(t)\,\mathrm dt$，则 $\|T\|\le a<1$，而
\[
g_n=\sum_{j=0}^nT^jf.
\]
所以该级数在整个半轴上一致收敛，极限 $g$ 连续，且满足 $g=f+Tg$。

这里取一致范数 $\|u\|_\infty=\sup_{x\ge0}|u(x)|$，则 $|Tu(x)|\le a\|u\|_\infty$。归纳展开递推式得到有限和表示，并有 $\|T^jf\|_\infty\le a^j\|f\|_\infty$；几何级数收敛保证函数项级数一致收敛。又 $\|Tg_n-Tg\|_\infty\le a\|g_n-g\|_\infty\to0$，故对递推取极限是合法的。

令 $H(x)=\int_0^xh(t)\,\mathrm dt$、$u=g-f$。虽然 $f$ 未必可微，但 $u=Tg$ 可微，并满足 $u'=h(f+u)$、$u(0)=0$。乘积分因子 $e^{-H}$ 并积分，得
\[
\boxed{g(x)=f(x)+\int_0^xe^{H(x)-H(t)}h(t)f(t)\,\mathrm dt}.
\]

\end{solve}

\begin{remark}[方法与条件]Volterra 积分的次序结构实际上更强：$\|T^j\|\le a^j/j!$，因为迭代积分区域是有序单纯形。因此只要 $\int_0^\infty|h|<\infty$，即使不小于 $1$，同一迭代仍一致收敛。一般的算子范数几何级数只是这里较粗的充分条件。\end{remark}

\begin{example}[来源序号6]\label{ex:pre07-math:6}\label{pr:pre07-math:6}

设连续函数 $f:\mathbb R\to\mathbb R$ 有下界或有上界，且存在 $a>0$，使
\[
f(x)+a\int_{x-1}^xf(t)\,\mathrm dt
\]
与 $x$ 无关。证明 $f$ 必为常数。

\par\href{https://github.com/lzxcr/cmc-mathematics/blob/64916fa1fdbc3e444a652eca204d5669058ee321/src/problems/preliminary/07/06.md}{题目与来源条目}

\end{example}

\begin{solve}\label{ans:pre07-math:6}

将 $f$ 换成 $-f$ 可交换上下界，故只需考虑有下界的情形。令 $m=\inf f$、$g=f-m\ge0$，则 $\inf g=0$，且
\[
g(x)+a\int_{x-1}^xg(t)\,\mathrm dt=C\ge0.
\]
积分项可微，所以 $g$ 可微；求导得 $g'+ag=ag(x-1)\ge0$，故 $e^{ax}g(x)$ 递增。对 $x-1\le t\le x$，因此有 $g(t)\le e^{a(x-t)}g(x)$，代回得到
\[
C\le g(x)+ag(x)\int_{x-1}^xe^{a(x-t)}\,\mathrm dt=e^ag(x).
\]
取下确界即得 $C=0$，再由非负性知 $g\equiv0$，所以 $f\equiv m$。

下确界不必在某点取到，但定义保证存在点列 $x_j$ 使 $g(x_j)\to0$。因此将 $C\le e^ag(x_j)$ 取极限，仍能推出 $C\le0$，结合 $C\ge0$ 得 $C=0$。原恒等式的两项都非负，故每个 $x$ 都有 $0\le g(x)\le C=0$。求导也没有额外假设：连续函数的变限积分可微，而恒等式把 $g$ 表示成常数减去该积分，所以 $g$ 的可微性是由题设推出来的。

\end{solve}

\begin{remark}[方法与条件]关键是先减去下确界，使延迟方程保留正性；积分因子随即把局部平均与单点值比较。下确界不必在某处取得。\end{remark}

\section{初赛：第 8 届（2016 年）}

本节选入3道课程相关题；来源序号保留原题解编排，未选入题目不重新补号。

\begin{example}[来源序号4]\label{ex:pre08-math:4}\label{pr:pre08-math:4}

设 $f_0,f_1$ 为 $[0,1]$ 上的正连续函数，满足 $\int_0^1f_0\le\int_0^1f_1$。递推定义
\[
f_{n+1}(x)=\frac{2f_n(x)^2}{f_n(x)+f_{n-1}(x)},\qquad n\ge1.
\]
证明 $a_n=\int_0^1f_n(x)\,\mathrm dx$（$n\ge1$）单调递增且收敛。

\par\href{https://github.com/lzxcr/cmc-mathematics/blob/64916fa1fdbc3e444a652eca204d5669058ee321/src/problems/preliminary/08/04.md}{题目与来源条目}

\end{example}

\begin{solve}\label{ans:pre08-math:4}

Cauchy--Schwarz 不等式给出
\[
a_{n+1}=2\int_0^1\frac{f_n^2}{f_n+f_{n-1}}
\ge\frac{2a_n^2}{a_n+a_{n-1}}.
\]
Cauchy--Schwarz 的具体应用是
\[
\left(\int_0^1 f_n\right)^2
=\left(\int_0^1\frac{f_n}{\sqrt{f_n+f_{n-1}}}\sqrt{f_n+f_{n-1}}\right)^2
\le\left(\int_0^1\frac{f_n^2}{f_n+f_{n-1}}\right)(a_n+a_{n-1}).
\]
分母始终为正，故所有运算成立。若 $a_n\ge a_{n-1}$，则 $2a_n^2/(a_n+a_{n-1})\ge a_n$，这就完成了单调性的归纳步骤。

由 $a_1\ge a_0$ 归纳可得 $a_{n+1}\ge a_n$。

为证明有界，逐点令 $u_n=f_{n-1}/f_n$，则 $u_{n+1}=(1+u_n)/2$，所以
\[
\frac{f_{n+1}}{f_n}
=\frac1{1+2^{-n}(f_0/f_1-1)}
\le\frac1{1-2^{-n}}.
\]
由 $-\log(1-t)\le2t$（$0\le t\le1/2$），
\[
f_n\le f_1\prod_{j=1}^{n-1}(1-2^{-j})^{-1}\le e^2 f_1.
\]
因此 $a_n\le e^2a_1$，单调有界，必收敛。

\end{solve}

\begin{remark}[方法与条件]积分的单调性由凸性控制，逐点的大小则由比值递推控制。题设只比较两个初始积分，不要求 $f_1\ge f_0$ 逐点成立；把这两种顺序混用会使证明失效。\end{remark}

\begin{example}[来源序号5]\label{ex:pre08-math:5}\label{pr:pre08-math:5}

设 $\alpha>1$。证明不存在定义于 $[0,\infty)$ 的正可导函数 $f$，使
\[
f'(x)\ge f(x)^\alpha,\qquad x\ge0.
\]

\par\href{https://github.com/lzxcr/cmc-mathematics/blob/64916fa1fdbc3e444a652eca204d5669058ee321/src/problems/preliminary/08/05.md}{题目与来源条目}

\end{example}

\begin{solve}\label{ans:pre08-math:5}

令 $g=f^{1-\alpha}>0$，则
\[
g'=(1-\alpha)f^{-\alpha}f'\le-(\alpha-1).
\]
由中值定理，$g(x)\le g(0)-(\alpha-1)x$，右端在 $x>g(0)/(\alpha-1)$ 时为负，与 $g>0$ 矛盾。

变换 $g=f^{1-\alpha}$ 的目的是抵消右端的 $f^\alpha$：链式法则产生因子 $f^{-\alpha}$，相乘后成为常数。由于 $1-\alpha<0$，乘上这个系数时不等号方向必须反转。对 $g(x)-g(0)$ 用中值定理便能得到线性上界，不需要假定 $f'$ 连续或对导数作不加说明的积分。矛盾说明正解不能在整个半轴上存在；例如等式方程的正解也会在有限位置爆破。

\end{solve}

\begin{remark}[方法与条件]这种换元把超线性增长转为匀速下降，实际给出有限时间爆破的上界 $f(0)^{1-\alpha}/(\alpha-1)$。一般形式见\hyperref[app:dynamics-blowup]{微分方程的增长与爆破}。\end{remark}

\begin{example}[来源序号6]\label{ex:pre08-math:6}\label{pr:pre08-math:6}

设 $f,g$ 为 $[0,1]$ 上的单调递增函数，满足 $0\le f,g\le1$，且 $\int_0^1f=\int_0^1g$。证明
\[
\int_0^1|f(x)-g(x)|\,\mathrm dx\le\frac12.
\]

\par\href{https://github.com/lzxcr/cmc-mathematics/blob/64916fa1fdbc3e444a652eca204d5669058ee321/src/problems/preliminary/08/06.md}{题目与来源条目}

\end{example}

\begin{solve}\label{ans:pre08-math:6}

记共同积分为 $m$。同向单调性给出
\[
\int_0^1fg-m^2
=\frac12\int_0^1\!\int_0^1(f(x)-f(y))(g(x)-g(y))\,\mathrm dx\,\mathrm dy\ge0.
\]
又因 $0\le f,g\le1$，有 $\min(f,g)\ge fg$。因此
\[
\int_0^1|f-g|=2m-2\int_0^1\min(f,g)
\le2m(1-m)\le\frac12.
\]

点态恒等式 $|u-v|=u+v-2\min(u,v)$ 将绝对值积分转化为交叠部分的下界。若 $u\le v$，则 $uv\le u$，另一种大小关系同理，所以 $\min(u,v)\ge uv$。最后 $2m(1-m)=\tfrac12-2(m-\tfrac12)^2\le\tfrac12$。常数确实可取到：令 $f\equiv1/2$，$g$ 在左半区间为 $0$、右半区间为 $1$，二者同向单调、积分均为 $1/2$，而差的绝对值积分恰为 $1/2$。

\end{solve}

\begin{remark}[方法与条件]常数 $1/2$ 最优：取 $f\equiv1/2$、$g=1_{[1/2,1]}$ 即可取等。证明把绝对差变为重叠量，再由单调性给出重叠的下界。\end{remark}

\section{初赛：第 9 届（2017 年）}

本节选入3道课程相关题；来源序号保留原题解编排，未选入题目不重新补号。

\begin{example}[来源序号2]\label{ex:pre09-math:2}\label{pr:pre09-math:2}

设 $a_n$ 递增且 $a_1>1$。证明级数
\[
\sum_{n=1}^\infty\frac{a_{n+1}-a_n}{a_n\log a_{n+1}}
\]
收敛当且仅当 $a_n$ 有界。若把分母中的 $a_n$ 换成 $a_{n+1}$，此结论是否仍成立？

\par\href{https://github.com/lzxcr/cmc-mathematics/blob/64916fa1fdbc3e444a652eca204d5669058ee321/src/problems/preliminary/09/02.md}{题目与来源条目}

\end{example}

\begin{solve}\label{ans:pre09-math:2}

若 $a_n\le M$，每一项不超过 $(a_{n+1}-a_n)/(a_1\log a_1)$，由望远镜求和知级数收敛。

反过来，设级数收敛，令 $b_n=\log a_n>0$。由 $e^t-1\ge t$，
\[
\frac{a_{n+1}-a_n}{a_n\log a_{n+1}}
\ge\frac{b_{n+1}-b_n}{b_{n+1}}=:r_n\ge0.
\]
因为 $a_{n+1}/a_n=e^{b_{n+1}-b_n}$，原级数的项正是 $[e^{b_{n+1}-b_n}-1]/b_{n+1}$。又 $r_n=1-b_n/b_{n+1}<1$，所以 $\log(b_{n+1}/b_n)=-\log(1-r_n)$ 有定义。对 $0\le r\le1/2$，$-\log(1-r)=\int_0^r(1-t)^{-1}\,dt\le2r$，这说明为何级数收敛能控制对数的总增量。

故 $\sum r_n<\infty$，充分大时 $r_n\le1/2$，于是
\[
\log b_{n+1}-\log b_n=-\log(1-r_n)\le2r_n.
\]
求和得 $\log b_n$ 有界，因而 $a_n$ 有界。

替换分母后结论不成立。取 $a_n=e^{2^n}$，虽然 $a_n\to\infty$，但
\[
0\le\frac{a_{n+1}-a_n}{a_{n+1}\log a_{n+1}}
\le\frac1{2^{n+1}},
\]
相应级数仍收敛。

\end{solve}

\begin{remark}[方法与条件]这里不能随意将左右端点的分母互换：当相邻项增长很快时，两者并不同阶。证明先取对数，再把级数收敛转为乘法增长有界。\end{remark}

\begin{example}[来源序号5]\label{ex:pre09-math:5}\label{pr:pre09-math:5}

设 $f(x)=\arctan x$，$A$ 为常数。若有限极限
\[
B=\lim_{n\to\infty}\left(\sum_{k=1}^nf(k/n)-An\right)
\]
存在，求 $A,B$。

\par\href{https://github.com/lzxcr/cmc-mathematics/blob/64916fa1fdbc3e444a652eca204d5669058ee321/src/problems/preliminary/09/05.md}{题目与来源条目}

\end{example}

\begin{solve}\label{ans:pre09-math:5}

将括号内除以 $n$，由 Riemann 和知
\[
A=\int_0^1\arctan x\,\mathrm dx=\frac\pi4-\frac12\log2.
\]
在第 $k$ 个小区间以右端点作 Taylor 展开。因 $f''$ 有界，逐段积分并求和得
\[
\sum_{k=1}^nf(k/n)-n\int_0^1f
=\frac1{2n}\sum_{k=1}^nf'(k/n)+O(n^{-1}).
\]
把每格计算写开：设 $b=k/n$，对 $b-1/n\le x\le b$ 有 $f(x)=f(b)+f'(b)(x-b)+O(n^{-2})$，其中余项常数可对所有格子统一选取。积分后
\[
\int_{b-1/n}^bf(x)\,dx=\frac{f(b)}n-\frac{f'(b)}{2n^2}+O(n^{-3}).
\]
求和再乘以 $n$，总误差为 $O(n^{-1})$，于是得到上式。主项 $A$ 的积分则由分部积分 $\int\arctan x\,dx=x\arctan x-\tfrac12\log(1+x^2)$ 求得。

右端趋于 $\frac12\int_0^1f'=[f(1)-f(0)]/2=\pi/8$。所以
\[
\boxed{A=\frac\pi4-\frac12\log2,\qquad B=\frac\pi8}.
\]

\end{solve}

\begin{remark}[方法与条件]第一项来自面积，下一项来自端点偏差。这是 Euler--Maclaurin 公式的首个修正项，但本题用一次 Taylor 展开已足够，见\hyperref[app:integration-riemann]{Riemann 和的误差}。\end{remark}

\begin{example}[来源序号6]\label{ex:pre09-math:6}\label{pr:pre09-math:6}

设 $f(x)=1-x^2+x^3$（$0\le x\le1$）。计算并说明
\[
\lim_{n\to\infty}
\frac{\int_0^1f(x)^n\log(x+2)\,\mathrm dx}{\int_0^1f(x)^n\,\mathrm dx}.
\]

\par\href{https://github.com/lzxcr/cmc-mathematics/blob/64916fa1fdbc3e444a652eca204d5669058ee321/src/problems/preliminary/09/06.md}{题目与来源条目}

\end{example}

\begin{solve}\label{ans:pre09-math:6}

$f$ 在 $0,1$ 都取最大值 $1$，但两处下降速度不同。记分母为 $I_n$。当 $0\le x\le n^{-1/2}$ 时，$f(x)\ge1-1/n$，所以对 $n\ge2$，
\[
I_n\ge n^{-1/2}(1-1/n)^n\ge e^{-2}n^{-1/2}.
\]
另一方面，对任意固定 $0<\delta<1$，由 $1-f(x)=x^2(1-x)$ 及 $1-u\le e^{-u}$，
\[
\int_\delta^1f(x)^n\,\mathrm dx
\le\int_\delta^1e^{-n\delta^2(1-x)}\,\mathrm dx
\le\frac1{n\delta^2}.
\]
因此，以密度 $f^n/I_n$ 定义的概率质量全部集中到 $0$：区间 $[\delta,1]$ 的质量趋零。利用 $\log(x+2)$ 在 $0$ 连续且整体有界，先固定 $\delta$，再令 $\delta\downarrow0$，即得所求极限为 $\boxed{\log2}$。

为直接核对比值极限，令 $w_n=f^n/I_n$，则 $\int_0^1w_n=1$，且 $\int_\delta^1w_n\le e^2/(\delta^2\sqrt n)$。于是
\[
\left|\int_0^1w_n\log(x+2)\,dx-\log2\right|
\le\sup_{0\le x\le\delta}|\log(x+2)-\log2|
+\log(3/2)\frac{e^2}{\delta^2\sqrt n}.
\]
先固定 $\delta$ 令 $n$ 趋于无穷，再缩小 $\delta$ 即得零。虽然两个端点都使 $f=1$，但不能把它们的函数值简单平均，必须比较各端点附近的积分质量。

\end{solve}

\begin{remark}[方法与条件]最大值相同不意味着两峰贡献相同：在 $0$ 附近 $1-f\sim x^2$，质量尺度为 $n^{-1/2}$；在 $1$ 附近 $1-f\sim1-x$，质量尺度为 $n^{-1}$。更平坦的峰主导归一化积分，见\hyperref[app:integration-local]{局部化与尺度}。\end{remark}

\section{初赛：第 10 届（2018 年）}

本节选入3道课程相关题；来源序号保留原题解编排，未选入题目不重新补号。

\begin{example}[来源序号4]\label{ex:pre10-math:4}\label{pr:pre10-math:4}

设 $f\in C^1([0,1])$，在 $0$ 处有任意阶导数且 $f^{(n)}(0)=0$（$n\ge0$）。若存在 $C>0$，使
\[
|xf'(x)|\le C|f(x)|,\qquad 0\le x\le1,
\]
证明：（1）$\lim_{x\downarrow0}f(x)/x^n=0$ 对每个非负整数 $n$ 成立；（2）$f\equiv0$。

\par\href{https://github.com/lzxcr/cmc-mathematics/blob/64916fa1fdbc3e444a652eca204d5669058ee321/src/problems/preliminary/10/04.md}{题目与来源条目}

\end{example}

\begin{solve}\label{ans:pre10-math:4}

（1）对任意 $n$，由 $0$ 处的 Taylor 公式及全部低阶导数为零，$f(x)=o(x^n)$。

（2）在 $x>0$ 时计算
\[
\left(\frac{f(x)^2}{x^{2C}}\right)'
=\frac{2xf(x)f'(x)-2Cf(x)^2}{x^{2C+1}}\le0.
\]
取整数 $n>C$，由第一问知 $f(x)^2/x^{2C}\to0$。这个函数非负且递减，故对任意 $x>0$，其值不超过左端极限 $0$，只能恒为零。再用 $f(0)=0$，得结论。

使用平方是为了避免在 $f$ 的零点处求 $|f|$ 的导数，也无需预先证明 $f$ 不变号。分子满足 $xf f'\le |f|\,|xf'|\le Cf^2$，因而导数非正。若 $n>C$，则 $f(x)^2/x^{2C}=(f(x)/x^n)^2x^{2(n-C)}\to0$。对固定 $x>0$ 和任意 $0<t<x$，递减性给出 $0\le f(x)^2/x^{2C}\le f(t)^2/t^{2C}$；让 $t\downarrow0$，得到逐点为零。

\end{solve}

\begin{remark}[方法与条件]全部导数在一点为零，单独并不能推出函数恒零；关键是微分不等式限制了离开零点后的增长速度。平方消除了绝对值在零点处的可微性问题，幂权则消除了系数 $C/x$。\end{remark}

\begin{example}[来源序号5]\label{ex:pre10-math:5}\label{pr:pre10-math:5}

设 $a_n>0$，$\sum_{n=1}^\infty b_n$ 绝对收敛，且对 $n\ge2$ 有
\[
\frac{a_n}{a_{n+1}}\le1+\frac1n+\frac1{n\log n}+b_n.
\]
证明：（1）$\displaystyle\frac{a_n}{a_{n+1}}<\frac{n+1}{n}\frac{\log(n+1)}{\log n}+b_n$；（2）$\sum_{n=1}^\infty a_n$ 发散。

\par\href{https://github.com/lzxcr/cmc-mathematics/blob/64916fa1fdbc3e444a652eca204d5669058ee321/src/problems/preliminary/10/05.md}{题目与来源条目}

\end{example}

\begin{solve}\label{ans:pre10-math:5}

（1）由 $\log(1+1/n)>1/(n+1)$，
\[
\frac{n+1}{n}\frac{\log(n+1)}{\log n}
=1+\frac1n+\frac{n+1}{n\log n}\log(1+1/n)
>1+\frac1n+\frac1{n\log n},
\]
故第一问成立。

（2）令 $c_n=n\log n\,a_n$（$n\ge2$），第一问给出
\[
\frac{c_n}{c_{n+1}}
<1+\frac{n\log n}{(n+1)\log(n+1)}b_n
\le1+|b_n|.
\]
相乘并用 $1+u\le e^u$，得到
\[
c_n\ge c_2\exp\!\left(-\sum_{j=2}^\infty|b_j|\right)=:c>0.
\]
所以 $a_n\ge c/(n\log n)$，由比较判别法知级数发散。

第一问所用严格不等式可由积分得到：$\log(1+1/n)=\int_n^{n+1}t^{-1}\,dt>1/(n+1)$。第二问的乘积方向也应注意：$c_j/c_{j+1}\le1+|b_j|$ 等价于 $c_{j+1}\ge c_j/(1+|b_j|)$，连续相乘才给出正的统一下界。最后 $\int_2^R dx/(x\log x)=\log\log R-\log\log2\to\infty$，由积分判别法知比较级数发散。

\end{solve}

\begin{remark}[方法与条件]临界比较列 $1/(n\log n)$ 正是前两个比值修正项对应的尺度；绝对可和的扰动只改变一个正的乘法常数，不改变发散性。\end{remark}

\begin{example}[来源序号6]\label{ex:pre10-math:6}\label{pr:pre10-math:6}

设可微函数 $f:\mathbb R\to(0,\infty)$ 满足
\[
|f'(x)-f'(y)|\le|x-y|^\alpha,\qquad x,y\in\mathbb R,
\]
其中 $0<\alpha\le1$。证明对每个 $x\in\mathbb R$，
\[
|f'(x)|^{(\alpha+1)/\alpha}<\frac{\alpha+1}{\alpha}f(x).
\]

\par\href{https://github.com/lzxcr/cmc-mathematics/blob/64916fa1fdbc3e444a652eca204d5669058ee321/src/problems/preliminary/10/06.md}{题目与来源条目}

\end{example}

\begin{solve}\label{ans:pre10-math:6}

导数的 Hölder 条件给出对任意实数 $h$ 的估计
\[
f(x+h)\le f(x)+f'(x)h+\frac{|h|^{\alpha+1}}{\alpha+1}.
\]
这是对 $f'(x+t)-f'(x)$ 积分的直接结果。

更明确地，导数满足 Hölder 连续，故可以用微积分基本定理写成
\[
f(x+h)-f(x)-f'(x)h
=h\int_0^1[f'(x+th)-f'(x)]\,dt.
\]
余项绝对值不超过 $|h|^{\alpha+1}\int_0^1t^\alpha\,dt=|h|^{\alpha+1}/(\alpha+1)$。沿导数的反方向走距离 $r$ 时，线性项为 $-|f'(x)|r$；令其与余项之和对 $r$ 的导数为零，即 $r^\alpha=|f'(x)|$，就得到下面的步长选择。

若 $f'(x)=0$，结论由 $f(x)>0$ 立即成立。否则取使右端最小的步长
\[
h=-\operatorname{sgn}(f'(x))|f'(x)|^{1/\alpha}.
\]
由 $f(x+h)>0$，得到
\[
0<f(x)-\frac\alpha{\alpha+1}|f'(x)|^{(\alpha+1)/\alpha},
\]
即所求。

\end{solve}

\begin{remark}[方法与条件]正性把函数允许下降的幅度限制为 $f(x)$，导数的连续模则控制线性近似的余项。选择使两者平衡的步长，就是\hyperref[app:estimate-interpolation]{尺度选择与插值估计}；$\alpha=1$ 时得到 $|f'|^2<2f$。\end{remark}

\section{初赛：第 11 届（2019 年）}

本节选入3道课程相关题；来源序号保留原题解编排，未选入题目不重新补号。

\begin{example}[来源序号2]\label{ex:ch10-11th-A:2}\label{pr:ch10-11th-A:2}

设 $\alpha>0$，$f(x)$ 在 $[0,1]$ 上非负，有二阶导函数，$f(0)=0$，且在 $[0,1]$ 上不恒为零。求证：存在 $\xi\in(0,1)$ 使得
\[
\xi f''(\xi)+(\alpha+1)f'(\xi)>\alpha f(\xi).
\]

\par\href{https://github.com/lzxcr/cmc-mathematics/blob/64916fa1fdbc3e444a652eca204d5669058ee321/src/problems/preliminary/11/02.md}{题目与来源条目}

\end{example}

\begin{solve}\label{ans:ch10-11th-A:2}

反设 $xf''+(\alpha+1)f'\leq\alpha f$ 在 $(0,1)$ 上恒成立。左端是 $xf'+\alpha f$ 的导数，因而先积分一次可把二阶条件降到一阶。令 $F(x)=\int_0^x f(t)\,\mathrm dt$，则
\[
H(x)=xf'(x)+\alpha f(x)-\alpha F(x)
\]
满足 $H'\leq0$、$H(0)=0$，故 $xf'+\alpha f\leq\alpha F$。乘以 $x^{\alpha-1}$，从 $0$ 积分至 $x$，并利用 $F$ 非负递增，得到
\[
x^\alpha f(x)
\leq\alpha\int_0^x t^{\alpha-1}F(t)\,\mathrm dt
\leq x^\alpha F(x).
\]
因此 $F'=f\leq F$，从而 $(e^{-x}F)'\leq0$。但 $F(0)=0$、$F\geq0$，只能有 $F\equiv0$，继而 $f\equiv0$，与题设矛盾。

降阶后再乘 $x^{\alpha-1}$，是因为 $(x^\alpha f)'=x^{\alpha-1}(xf'+\alpha f)$。严格处理零端时，可先在 $[\varepsilon,x]$ 上积分，再让 $\varepsilon\downarrow0$；$f$ 连续有界使 $\varepsilon^\alpha f(\varepsilon)\to0$，而 $\alpha>0$ 保证 $t^{\alpha-1}$ 在零端可积。又 $0\le F(t)\le F(x)$，故积分至多为 $F(x)x^\alpha$，从而最终得到 $F'\le F$。

\end{solve}

\begin{remark}[方法与条件]这个证明是积分比较原理：先把二阶微分不等式降为 $f\leq\int_0^x f$，再用积分因子排除从零初值长出的非负解。后一结论就是 Gronwall 不等式的零初值情形。\end{remark}

\begin{example}[来源序号5]\label{ex:ch10-11th-A:5}\label{pr:ch10-11th-A:5}

设 $\delta>0$，$\alpha\in(0,1)$，实数列 $\left\{x_{n}\right\}$ 满足
\[
x_{n+1}=x_{n}\left(1-\frac{h_{n}}{n^{\alpha}}\right)+\frac{1}{n^{\alpha+\delta}},\qquad n\geq 1,
\]
其中 $\left\{h_{n}\right\}$ 有正的上下界。证明：$\left\{n^{\delta}x_{n}\right\}$ 有界。

\par\href{https://github.com/lzxcr/cmc-mathematics/blob/64916fa1fdbc3e444a652eca204d5669058ee321/src/problems/preliminary/11/05.md}{题目与来源条目}

\end{example}

\begin{solve}\label{ans:ch10-11th-A:5}

驱动项 $n^{-\alpha-\delta}$ 与衰减系数 $n^{-\alpha}$ 的比值为 $n^{-\delta}$，这提示按题目中的尺度归一化。记 $0<c\leq h_n\leq C$，令 $y_n=n^\delta x_n$。递推式化为
\[
y_{n+1}=\left(1+\frac1n\right)^\delta
\left(1-\frac{h_n}{n^\alpha}\right)y_n
+\left(1+\frac1n\right)^\delta\frac1{n^\alpha}.
\]
因为 $\alpha<1$，有 $n^{-1}=o(n^{-\alpha})$。因此对充分大的 $n$，第一项系数非负且不超过 $1-c/(2n^\alpha)$，第二项不超过 $2/n^\alpha$。于是
\[
|y_{n+1}|\leq\left(1-\frac{c}{2n^\alpha}\right)|y_n|+\frac2{n^\alpha}.
\]
取起始下标 $N$ 足够大及 $K=\max\{|y_N|,4/c\}$。若 $|y_n|\leq K$，上式便给出 $|y_{n+1}|\leq K$。归纳可知尾项有界，再加有限个首项，得到所需结论。

系数估计可以展开为 $(1+1/n)^\delta(1-h_n n^{-\alpha})=1-h_n n^{-\alpha}+O(n^{-1})$，余项常数因 $h_n\le C$ 可统一选择。由于 $n^{-1}/n^{-\alpha}\to0$，可使余项不超过 $c/(2n^\alpha)$；同时选 $N$ 使 $h_n/n^\alpha\le1$，保证系数非负。归纳界中 $-(c/2)K+2\le0$ 正是选择 $K\ge4/c$ 的理由，有限个首项即使变号也不影响尾部论证。

\end{solve}

\begin{remark}[方法与条件]缩放 $y_n=n^\delta x_n$ 后，递推成为一个稳定的仿射平均：衰减量与驱动量同为 $n^{-\alpha}$，而缩放引入的扰动只有 $O(1/n)$。条件 $\alpha<1$ 正是保证前者控制后者。

这类递推中如何选尺度、建立不变区间并进一步求渐近，见\hyperref[app:estimate-recursion]{方法卷：递推中的平衡尺度}。\end{remark}

\begin{example}[来源序号6]\label{ex:ch10-11th-A:6}\label{pr:ch10-11th-A:6}

设 $f(x)=\frac{1}{1+e^{x}}$。
（1）证明 $f(x)$ 是 $[0,+\infty)$ 上的凸函数。进一步，证明当 $x,y\geq 0$ 时成立
\[
f(x)+f(y)\leq f(0)+f(x+y).
\]
（2）设 $n\geq 3$，试确定集合
\[
E\equiv\left\{\sum_{k=1}^{n}f\left(x_{k}\right)\ \middle|\ \sum_{k=1}^{n}x_{k}=0,\ x_{1},\dots,x_{n}\in\mathbb{R}\right\}.
\]

\par\href{https://github.com/lzxcr/cmc-mathematics/blob/64916fa1fdbc3e444a652eca204d5669058ee321/src/problems/preliminary/11/06.md}{题目与来源条目}

\end{example}

\begin{solve}\label{ans:ch10-11th-A:6}

(i) 计算得
\[
f''(x)=\frac{e^x(e^x-1)}{(1+e^x)^3}\geq0\qquad(x\geq0).
\]
所以 $f$ 凸。若 $s=x+y>0$，分别将 $x,y$ 写成端点 $0,s$ 的凸组合并相加，即得 $f(x)+f(y)\leq f(0)+f(s)$；$s=0$ 时为等式。

(ii) 答案为 $E=(1,n-1)$。先证下界。若所有 $x_k=0$，和为 $n/2>1$；否则有正项也有负项。若至少有两个负项，它们的函数值都大于 $1/2$，故总和大于 $1$。若只有一个负项 $-T$，则任取正项 $a$ 都有 $0<a\leq T$，从而
\[
f(-T)+f(a)\geq f(-T)+f(T)=1.
\]
其余 $n-2$ 项的函数值都为正，故总和仍严格大于 $1$。再用 $f(t)+f(-t)=1$，对 $(-x_k)$ 应用下界，便得总和严格小于 $n-1$。

为证明两界之间没有遗漏，只需沿一条满足约束的连续路径取值。取
\[
(x_1,\ldots,x_n)=(t,\ldots,t,-(n-1)t),
\]
当 $t\to+\infty$ 时和趋于 $1$，当 $t\to-\infty$ 时趋于 $n-1$。沿这条路径，函数和关于 $t$ 连续；由介值定理，它已经取到 $(1,n-1)$ 的每个值。结合严格上下界，故 $E=(1,n-1)$。

第一问的两组权重分别为 $x/s$ 与 $y/s$，把 $f(x)\le(y/s)f(0)+(x/s)f(s)$ 和对应的 $f(y)$ 不等式相加，端点前的系数都变成 $1$。第二问中 $f(-t)=1-f(t)$ 直接由分母同乘 $e^t$ 得到。对于任意给定 $z\in(1,n-1)$，连续路径两端的极限保证存在有限的 $t_-<t_+$ 使路径值分别大于、小于 $z$；在这个有限区间上用介值定理，便可取到 $z$，无需证明路径单调。

\end{solve}

\begin{remark}[方法与条件]第一问将固定和的两个非负变量向端点集中，这是凸性的直接作用。第二问则主要使用 $f(-t)=1-f(t)$ 及单调性：先排除两个端点，再由一条连接两种极限配置的路径得到全部中间值。两问所利用的结构不同，无须强行把值域问题化成高维极值计算。\end{remark}

\section{初赛：第 12 届（2020 年）}

本节选入3道课程相关题；来源序号保留原题解编排，未选入题目不重新补号。

\begin{example}[来源序号2]\label{ex:ch09-12th-A:2}\label{pr:ch09-12th-A:2}

求极限
\[
I=\lim_{n\to\infty}\frac{\ln n}{\ln\left(1^{2020}+2^{2020}+\cdots+n^{2020}\right)}.
\]

\par\href{https://github.com/lzxcr/cmc-mathematics/blob/64916fa1fdbc3e444a652eca204d5669058ee321/src/problems/preliminary/12/02.md}{题目与来源条目}

\end{example}

\begin{solve}\label{ans:ch09-12th-A:2}

对数会消去乘法常数，真正需要辨认的是和式的增长幂次。先提出每项的 $n^{2020}$ 和求和带来的另一因子 $n$，由 Riemann 和，
\[
\frac1{n^{2021}}\sum_{k=1}^n k^{2020}
=\frac1n\sum_{k=1}^n\left(\frac kn\right)^{2020}
\longrightarrow\int_0^1t^{2020}\,\mathrm dt=\frac1{2021}.
\]
因此分母为 $2021\ln n-\ln2021+o(1)$，所求极限为 $1/2021$。

记上面的归一化和为 $c_n$，则 $c_n\to1/2021>0$，故 $\ln c_n\to-\ln2021$。因此
\[
\frac{\ln n}{\ln\sum_{k=1}^nk^{2020}}
=\frac1{2021+(\ln c_n)/\ln n}\longrightarrow\frac1{2021}.
\]
这里必须先证明 $c_n$ 趋于正数，才能对其取对数并断言这部分只贡献有界项；仅知道原和趋于无穷大不足以判断两个对数的比值。

\end{solve}

\begin{example}[来源序号5]\label{ex:ch09-12th-A:5}\label{pr:ch09-12th-A:5}

设 $\varphi$ 是 $\mathbb{R}$ 上严格单调增加的连续函数，$\psi$ 是 $\varphi$ 的反函数，实数列 $\left\{x_{n}\right\}$ 满足
\[
x_{n+2}=\psi\left(\left(1-\frac{1}{\sqrt{n}}\right)\varphi\left(x_{n}\right)+\frac{1}{\sqrt{n}}\varphi\left(x_{n+1}\right)\right),\qquad n\geq 2.
\]
证明：$\left\{x_{n}\right\}$ 收敛或举例说明 $\left\{x_{n}\right\}$ 有可能发散。

\par\href{https://github.com/lzxcr/cmc-mathematics/blob/64916fa1fdbc3e444a652eca204d5669058ee321/src/problems/preliminary/12/05.md}{题目与来源条目}

\end{example}

\begin{solve}\label{ans:ch09-12th-A:5}

数列必收敛。递推中的非线性只来自坐标变换 $\varphi$，先在 $\varphi$ 坐标中观察，它就是相邻两项的凸平均。令 $y_n=\varphi(x_n)$，则
\[
y_{n+2}=\left(1-\frac1{\sqrt n}\right)y_n
+\frac1{\sqrt n}y_{n+1}\qquad(n\geq2).
\]
因此以 $y_n,y_{n+1}$ 为端点的闭区间逐次嵌套，且其长度满足
\[
|y_{n+2}-y_{n+1}|
=\left(1-\frac1{\sqrt n}\right)|y_{n+1}-y_n|.
\]
于是长度不超过
\[
|y_3-y_2|\exp\left(-\sum_{j=2}^n j^{-1/2}\right)\longrightarrow0.
\]
嵌套区间缩为一点，故 $y_n\to L$。所有 $y_n$ 都在端点为 $\varphi(x_2),\varphi(x_3)$ 的紧区间内，故 $L$ 属于 $\varphi$ 的值域。严格单调连续函数的反函数连续，得到 $x_n\to\psi(L)$。

令 $I_n=[\min(y_n,y_{n+1}),\max(y_n,y_{n+1})]$。因为新项是两端点的凸组合，$I_{n+1}\subset I_n$；其左端点递增、右端点递减，且都位于 $I_2$ 中，故分别收敛。长度趋零使两个端点极限相等，夹在其中的 $y_n$ 就收敛。特别要注意，单有 $|y_{n+1}-y_n|\to0$ 并不足以推出收敛，本题还利用了区间嵌套结构。

\end{solve}

\begin{example}[来源序号6]\label{ex:ch09-12th-A:6}\label{pr:ch09-12th-A:6}

对于有界区间 $[a,b]$ 的划分
\[
P:a=x_{0}<x_{1}<\cdots<x_{n+1}=b,
\]
其范数定义为
\[
\|P\|=\max_{0\leq k\leq n}\left(x_{k+1}-x_{k}\right).
\]
现设 $[a,b]$ 上函数 $f$ 满足 Lipschitz 条件，即存在常数 $M>0$，使得对任何 $x,y\in[a,b]$，成立 $|f(x)-f(y)|\leq M|x-y|$。定义
\[
s(f;P)\equiv\sum_{k=0}^{n}\sqrt{\left|x_{k+1}-x_{k}\right|^{2}+\left|f\left(x_{k+1}\right)-f\left(x_{k}\right)\right|^{2}}.
\]
若 $\lim\limits_{\|P\|\to 0^{+}}s(f;P)$ 存在，则称曲线 $y=f(x)$ 可求长。记 $P_{n}$ 为 $[a,b]$ 的 $2^{n}$ 等分。证明：
（1）$\lim\limits_{n\to\infty}s\left(f;P_{n}\right)$ 存在；
（2）曲线 $y=f(x)$ 可求长。

\par\href{https://github.com/lzxcr/cmc-mathematics/blob/64916fa1fdbc3e444a652eca204d5669058ee321/src/problems/preliminary/12/06.md}{题目与来源条目}

\end{example}

\begin{solve}\label{ans:ch09-12th-A:6}

记 $C=\sqrt{1+M^2}$。任意划分都满足
\[
s(f;P)\leq C(b-a),
\]
且加细划分不会减少折线长。

（1）二分划分逐次加细，因此 $s(f;P_n)$ 单调有界，极限存在。

（2）仅有二分划分的极限还不能控制任意划分。先固定一条接近最优的折线，再证明任意细划分加入其有限个顶点时只付出很小的长度误差。令 $L=\sup_P s(f;P)<\infty$。任给 $\varepsilon>0$，取固定划分 $Q$ 使 $s(f;Q)>L-\varepsilon$，设其内部分点数为 $N$。将 $Q$ 的分点加入任意划分 $P$，只影响至多 $N$ 个长度不超过 $\|P\|$ 的区间。在每个受影响区间上，加细后的折线长不超过该区间长度乘以 $C$，故
\[
0\leq s(f;P\cup Q)-s(f;P)\leq CN\|P\|.
\]
结合 $s(f;P\cup Q)\geq s(f;Q)$，得到
\[
L-\varepsilon-CN\|P\|<s(f;P)\leq L.
\]
先令 $\|P\|\to0$，再令 $\varepsilon\downarrow0$，即得 $s(f;P)\to L$。

任意一段弦长不超过 $C(x_{k+1}-x_k)$，求和得到统一上界；插入中间点时，三角不等式保证“两段之和不小于原一段”，所以加细的长度单调性也有直接依据。第二问中 $Q$ 必须先固定，之后其分点数 $N$ 才是与任意划分 $P$ 无关的常数。若要求总误差小于 $2\varepsilon$，取 $\|P\|<\varepsilon/[C(N+1)]$ 即可，这给出了任意划分极限的完整量词次序。

\end{solve}

\begin{remark}[方法与条件]折线长度的上确界给出不依赖参数的弧长定义；本题证明任意细划分都能逼近它。导数积分公式需要进一步的绝对连续性，不能从“几乎处处有导数”直接推出。弧长、绝对连续参数及变分之间的联系，见方法卷：从折线长到速度积分。\end{remark}

\section{初赛：第 13 届（2021 年）}

本节选入3道课程相关题；来源序号保留原题解编排，未选入题目不重新补号。

\begin{example}[来源序号2]\label{ex:ch07-13th-A:2}\label{pr:ch07-13th-A:2}

设 $B\subset\mathbb R^{n}(n\geq 2)$ 是单位开球，函数 $u,v$ 在 $\bar{B}$ 上连续，在 $B$ 内二阶连续可导，满足
\[
\begin{cases}-\Delta u-\left(1-u^{2}-v^{2}\right)u=0, & x\in B,\\ -\Delta v-\left(1-u^{2}-v^{2}\right)v=0, & x\in B,\\ u(x)=v(x)=0, & x\in\partial B,\end{cases}
\]
其中 $x=\left(x_{1},x_{2},\ldots,x_{n}\right)$，
\[
\Delta u=\frac{\partial^{2}u}{\partial x_{1}^{2}}+\frac{\partial^{2}u}{\partial x_{2}^{2}}+\cdots+\frac{\partial^{2}u}{\partial x_{n}^{2}},
\]
$\partial B$ 表示 $B$ 的边界。证明：
\[
u^{2}(x)+v^{2}(x)\leq 1\quad(\forall x\in\bar{B}).
\]

\par\href{https://github.com/lzxcr/cmc-mathematics/blob/64916fa1fdbc3e444a652eca204d5669058ee321/src/problems/preliminary/13/02.md}{题目与来源条目}

\end{example}

\begin{solve}\label{ans:ch07-13th-A:2}

两条方程通过同一个系数 $1-u^2-v^2$ 耦合；将它们合成为模平方，导数交叉项便变成非负的能量项。令 $w=u^2+v^2$。由方程，
\[
\Delta w=2(|\nabla u|^2+|\nabla v|^2)+2(w-1)w.
\]
若 $\max_{\bar B}w>1$，因边界上 $w=0$，最大值在内点取得；该点应有 $\Delta w\leq0$，但上式右端严格为正，矛盾。因此 $u^2+v^2\leq1$。

对每个坐标求两次导数，有 $\partial_{ii}(u^2)=2(\partial_i u)^2+2u\partial_{ii}u$，对 $v$ 同理；将方程改写成 $\Delta u=(w-1)u$、$\Delta v=(w-1)v$，相加即得公式。闭球紧且 $w$ 连续，最大值确能取得。若内点 $x_*$ 为最大值点，将 $w$ 限制在经过它的各条坐标直线上，一元极值必要条件给出 $\partial_{ii}w(x_*)\le0$，相加即 $\Delta w(x_*)\le0$。

\end{solve}

\begin{remark}[方法与条件]同一证明对任意有界区域和任意分量数都成立：若向量函数 $U$ 满足 $-\Delta U=(1-|U|^2)U$ 且边界上 $|U|\leq1$，则整个区域内 $|U|\leq1$。关键是对标量 $|U|^2$ 使用最大值原理。\end{remark}

\begin{example}[来源序号5]\label{ex:ch07-13th-A:5}\label{pr:ch07-13th-A:5}

设 $\alpha>1$，证明：
（1）\[
\int_{0}^{+\infty}\mathrm{d}x\int_{0}^{+\infty}e^{-t^{\alpha}x}\sin x\,\mathrm{d}t=\int_{0}^{+\infty}\mathrm{d}t\int_{0}^{+\infty}e^{-t^{\alpha}x}\sin x\,\mathrm{d}x.
\]
（2）计算
\[
\int_{0}^{+\infty}\sin x^{3}\,\mathrm{d}x\cdot\int_{0}^{+\infty}\sin x^{\frac{3}{2}}\,\mathrm{d}x.
\]

\par\href{https://github.com/lzxcr/cmc-mathematics/blob/64916fa1fdbc3e444a652eca204d5669058ee321/src/problems/preliminary/13/05.md}{题目与来源条目}

\end{example}

\begin{solve}\label{ans:ch07-13th-A:5}

（1）记 $p=1/\alpha\in(0,1)$。先对 $t$ 积分，
\[
\int_0^\infty e^{-t^\alpha x}\,\mathrm dt
=\frac{\Gamma(p)}\alpha x^{-p}.
\]
因此左边的外层积分在零点可积，在无穷远由 Dirichlet 判别法收敛。

直接换序的障碍在无穷远处：$x^{-p}\sin x$ 只条件收敛。引入阻尼 $\varepsilon>0$ 后，二重积分绝对收敛，可以用 Fubini 定理：
\[
\frac{\Gamma(p)}\alpha\int_0^\infty e^{-\varepsilon x}x^{-p}\sin x\,\mathrm dx
=\int_0^\infty\frac{\mathrm dt}{1+(t^\alpha+\varepsilon)^2}.
\]
阻尼后的绝对可积性可以直接检查：先对 $t$ 积分，绝对值二重积分化为常数倍的 $\int_0^\infty e^{-\varepsilon x}x^{-p}|\sin x|\,dx$；零端用 $|\sin x|\le x$，无穷端用指数衰减，均可积。内层有理表达式来自 $\int_0^\infty e^{-bx}\sin x\,dx=1/(1+b^2)$（$b>0$），可由两次分部积分求出。

右端被 $(1+t^{2\alpha})^{-1}$ 控制，故可令 $\varepsilon\downarrow0$。左端也回到原反常积分：在有限区间用控制收敛，在 $[R,\infty)$ 上分部积分给出对 $\varepsilon\geq0$ 一致的尾估计 $O(R^{-p})$。于是两边都等于 $\int_0^\infty(1+t^{2\alpha})^{-1}\,\mathrm dt$，换序成立。

（2）对 $0<s<1$，Gamma 积分的阻尼形式为
\[
\int_0^\infty x^{s-1}e^{-(\varepsilon-i)x}\,\mathrm dx
=\Gamma(s)(\varepsilon-i)^{-s}.
\]
这里复幂取 $-\pi/2<\arg(\varepsilon-i)<0$ 的分支；该公式由正实参数的 Gamma 积分向右半平面解析延拓得到。同样用 Dirichlet 尾估计令 $\varepsilon\downarrow0$，取虚部得
\[
\int_0^\infty x^{s-1}\sin x\,\mathrm dx
=\Gamma(s)\sin\frac{\pi s}{2}.
\]
对 $r>1$ 作代换 $u=x^r$，得到
\[
\int_0^\infty\sin(x^r)\,\mathrm dx
=\frac1r\Gamma(1/r)\sin\frac{\pi}{2r}.
\]
取 $r=3,3/2$，用 $\Gamma(1/3)\Gamma(2/3)=2\pi/\sqrt3$，所求乘积为
\[
\frac{\Gamma(1/3)}6\cdot
\frac{\sqrt3\,\Gamma(2/3)}3=\frac\pi9.
\]

第二问中代换后的指数分别为 $s=1/3$ 与 $s=2/3$，都在 $(0,1)$ 内，正是零端可积且无穷端满足 Dirichlet 判别的范围。所用 Gamma 乘积也可由 Beta 积分核查：$\Gamma(1/3)\Gamma(2/3)=\int_0^\infty t^{-2/3}/(1+t)\,dt=3\int_0^\infty du/(1+u^3)=2\pi/\sqrt3$。最后一个积分用 $1+u^3=(u+1)(u^2-u+1)$ 作部分分式分解即可计算；它说明这里的特殊函数常数最终归结为实变量有理积分。

\end{solve}

\begin{remark}[方法与条件]原二重积分不绝对可积，因此不能直接套用 Fubini 定理；加阻尼后换序，再用一致尾估计去掉阻尼，是振荡积分的一种通用处理方法。注意右边的累次积分已经是正函数 $(1+t^{2\alpha})^{-1}$ 的积分，它自身绝对收敛。

绝对可积、阻尼换序与一致尾部控制的各自作用，见\hyperref[app:integration-exchange]{方法卷：积分交换与尾部估计}。\end{remark}

\begin{example}[来源序号6]\label{ex:ch07-13th-A:6}\label{pr:ch07-13th-A:6}

设 $f,g$ 为 $\mathbb{R}$ 上的非负连续可微函数，满足：$\forall x\in\mathbb{R}$，成立
\[
f^{\prime}(x)\geq 6+f(x)-f^{2}(x),\qquad g^{\prime}(x)\leq 6+g(x)-g^{2}(x).
\]
证明：
（1）$\forall\varepsilon\in(0,1)$ 以及 $x\in\mathbb{R}$，存在 $\xi\in(-\infty,x)$ 使得 $f(\xi)\geq 3-\varepsilon$；
（2）$\forall x\in\mathbb{R}$，成立 $f(x)\geq 3$；
（3）$\forall x\in\mathbb{R}$，存在 $\eta\in(-\infty,x)$ 使得 $g(\eta)\leq 3$；
（4）$\forall x\in\mathbb{R}$，成立 $g(x)\leq 3$。

\par\href{https://github.com/lzxcr/cmc-mathematics/blob/64916fa1fdbc3e444a652eca204d5669058ee321/src/problems/preliminary/13/06.md}{题目与来源条目}

\end{example}

\begin{solve}\label{ans:ch07-13th-A:6}

右端二次式的两个零点为 $3,-2$。取到这两个点的差之比，可以把二次项约掉，同时非负性保证分母不为零。对非负函数 $u$，令
\[
W_u(x)=\frac{u(x)-3}{u(x)+2},\qquad -\frac32\leq W_u(x)<1.
\]
因为 $6+u-u^2=-(u-3)(u+2)$，所以题设等价地给出
\[
W_f'+5W_f\geq0,\qquad W_g'+5W_g\leq0.
\]
乘以积分因子 $e^{5x}$。对任意 $t<x$，有
\[
W_f(x)\geq e^{-5(x-t)}W_f(t)
\geq-\frac32e^{-5(x-t)},
\]
\[
W_g(x)\leq e^{-5(x-t)}W_g(t)
\leq e^{-5(x-t)}.
\]
令 $t\to-\infty$，得到 $W_f(x)\geq0$、$W_g(x)\leq0$，即 $f(x)\geq3$、$g(x)\leq3$。这证明了 （2）、（4）；（1）、（3）随即成立，甚至任取 $\xi<x$、$\eta<x$ 都满足相应要求。

商的求导给出 $W_u'=5u'/(u+2)^2$，而 $5W_u=5(u-3)(u+2)/(u+2)^2$，所以题设的二次微分不等式恰好化为线性不等式。关键是 $u\ge0$ 提供了 $-3/2\le W_u<1$ 的双侧常数界，乘上 $e^{-5(x-t)}$ 后在 $t\to-\infty$ 时都消失。如果函数只定义在有限左端起始的区间，就不能使用这一步，结论也不会由同一论证自动成立。

\end{solve}

\begin{remark}[方法与条件]代换 $(u-3)/(u+2)$ 利用了 Riccati 多项式的两个平衡点 $3,-2$，把非线性微分不等式化为线性不等式。全实轴上的非负性使新变量有统一界，再沿时间向 $-\infty$ 推移即可排除越过平衡点的情形；仅在半直线上给出条件时，结论一般不成立。\end{remark}

\section{初赛：第 13 届补赛（2021 年）}

本节选入3道课程相关题；来源序号保留原题解编排，未选入题目不重新补号。

\begin{remark}[来源说明]因疫情影响，部分赛区或学校延期举行的统一补赛试卷。\end{remark}

\begin{example}[来源序号2]\label{ex:ch08-13th-A-makeup:2}\label{pr:ch08-13th-A-makeup:2}

计算积分
\[
\int_{0}^{+\infty}\frac{x-x^{2}+x^{3}-x^{4}+\cdots-x^{2018}}{(1+x)^{2021}}\,\mathrm{d}x.
\]

\par\href{https://github.com/lzxcr/cmc-mathematics/blob/64916fa1fdbc3e444a652eca204d5669058ee321/src/problems/preliminary/13-makeup/02.md}{题目与来源条目}

\end{example}

\begin{solve}\label{ans:ch08-13th-A-makeup:2}

分母在 $x\mapsto1/x$ 下只改变一个幂次，这提示先寻找积分之间的配对关系，无需逐项求值。记 $I_k=\int_0^\infty x^k(1+x)^{-2021}\,\mathrm dx$（$1\leq k\leq2018$）；这些积分都绝对收敛。代换 $x=1/t$ 给出
\[
I_k=I_{2019-k}.
\]
而 $k$ 与 $2019-k$ 的奇偶性相反，所以分子中的相应两项积分相消，原积分等于 $0$。

换元时三个因子必须同时变换：$x^k=t^{-k}$，$(1+x)^{-2021}=t^{2021}(1+t)^{-2021}$，以及 $dx=-t^{-2}dt$，合起来的幂次是 $2019-k$。该配对将 $1,\ldots,2018$ 两两对应，且无整数固定点，因为 $2k=2019$ 无整数解。零端各项为 $O(x^k)$，无穷端为 $O(x^{k-2021})$，最高幂也只有 $x^{-3}$，因此有限项相消不存在反常积分分别发散的问题。

\end{solve}

\begin{remark}[方法与条件]倒数代换把 Beta 积分的两个参数互换：对 $-1<k<m-1$，
\[
\int_0^\infty\frac{x^k}{(1+x)^m}\,\mathrm dx
=B(k+1,m-k-1)=B(m-k-1,k+1).
\]
本题的零值由该对称性直接决定，无需分别计算各项。\end{remark}

\begin{example}[来源序号5]\label{ex:ch08-13th-A-makeup:5}\label{pr:ch08-13th-A-makeup:5}

设 $a_{k}\in(0,1)$，$1\leq k\leq 2021$，且
\[
a_{n+2021}=\left(a_{n}+a_{n+1}+\cdots+a_{n+2020}\right)^{1/2022},
\]
其中 $n=1,2,\cdots$，每一步均取正根。证明：$\lim_{n\to+\infty}a_{n}$ 存在。

\par\href{https://github.com/lzxcr/cmc-mathematics/blob/64916fa1fdbc3e444a652eca204d5669058ee321/src/problems/preliminary/13-makeup/05.md}{题目与来源条目}

\end{example}

\begin{solve}\label{ans:ch08-13th-A-makeup:5}

正的常数解满足 $X^{2022}=2021X$，故先用这个平衡值归一化。令 $X=2021^{1/2021}$、$x_n=a_n/X$，则正根递推化为
\[
x_{n+2021}^{2022}=\frac1{2021}\sum_{j=0}^{2020}x_{n+j}.
\]
取 $s=\min_{1\leq j\leq2021}x_j>0$。因初始项均小于 $1$，递推归纳给出 $s\leq x_n<1$：若窗口内各项在 $[s,1)$ 中，其均值也在该区间，开 $2022$ 次方后仍在 $[s,1)$。

记 $\ell=\liminf x_n\in[s,1]$。在递推式中取下极限，有限和的下极限不小于各下极限之和，故
\[
\ell^{2022}\geq\ell.
\]
由 $\ell>0$ 得 $\ell\geq1$，结合 $x_n<1$ 可知 $x_n\to1$。因此 $a_n\to2021^{1/2021}$。

取下极限这一步也可以不用运算法则而直接证明。任给 $0<\varepsilon<\ell$，全部充分大的 $k$ 都有 $x_k\ge\ell-\varepsilon$，故每个足够靠后的窗口均满足 $x_{n+2021}^{2022}\ge\ell-\varepsilon$。沿着趋于下极限的子列取极限，再让 $\varepsilon\downarrow0$，得到 $\ell^{2022}\ge\ell$。若 $0<\ell<1$，则 $\ell^{2022}<\ell$，矛盾；所以必须 $\ell=1$，而整体上界也是 $1$，上下极限相同。

\end{solve}

\begin{remark}[方法与条件]同一论证还适用于任意正初值：归一化后用初始项与 $1$ 的最小值、最大值夹住全序列，再分别对下、上极限使用 $\ell^{2022}\geq\ell$、$L^{2022}\leq L$，仍得 $a_n\to2021^{1/2021}$。

若将正根递推放宽为偶次幂等式，允许后续项取负值，则可能不收敛。取 $r\in(1/2021,1)$ 使
\[
2020r-(2021r)^{1/2022}-r^{2022}=0;
\]
左端在两端点异号，故这样的 $r$ 存在。令 $s=(2021r)^{1/2022}$，周期重复由 $2021$ 个 $r$ 和一个 $-s$ 组成的块。由 $s^{2022}=2021r$ 与 $r^{2022}=2020r-s$，所得数列满足偶次幂等式与全部初始条件，却不收敛。\end{remark}

\begin{example}[来源序号6]\label{ex:ch08-13th-A-makeup:6}\label{pr:ch08-13th-A-makeup:6}

（1）证明：
\[
\lim_{\alpha\to 0^{+}}\sum_{n=1}^{\infty}\frac{\cos\left(n+\frac{1}{2}\right)}{n^{1+\alpha}}=\sum_{n=1}^{\infty}\frac{\cos\left(n+\frac{1}{2}\right)}{n};
\]
（2）计算
\[
\lim_{\alpha\to 0^{+}}\sum_{n=1}^{\infty}\frac{\sin n}{n^{\alpha}},
\]
并说明理由。

\par\href{https://github.com/lzxcr/cmc-mathematics/blob/64916fa1fdbc3e444a652eca204d5669058ee321/src/problems/preliminary/13-makeup/06.md}{题目与来源条目}

\end{example}

\begin{solve}\label{ans:ch08-13th-A-makeup:6}

（1）$c_n=\cos(n+1/2)$ 的部分和一致有界，因为
\[
\sum_{n=1}^N c_n=\frac{\sin(N+1)-\sin1}{2\sin(1/2)}.
\]
分部求和给出对 $\alpha\geq0$ 一致的尾估计
\[
\left|\sum_{n=N}^M\frac{c_n}{n^{1+\alpha}}\right|\leq\frac C N.
\]
因此级数关于 $\alpha\geq0$ 一致收敛，可在 $\alpha\downarrow0$ 时逐项取极限，得到所需等式。

（2）此时不能指望一致尾估计延续到 $\alpha=0$，因为极限级数的通项不趋零。改将振荡项写成离散差分，让差分转移到缓慢变化的权重上。对每个 $\alpha>0$，Dirichlet 判别法保证级数收敛。用
\[
2\sin(1/2)\sin n=\cos(n-1/2)-\cos(n+1/2)
\]
作分部求和，得到
\[
\sum_{n=1}^\infty\frac{\sin n}{n^\alpha}
=\frac12\cot\frac12
-\frac1{2\sin(1/2)}\sum_{n=1}^\infty d_n(\alpha)c_n,
\quad d_n(\alpha)=n^{-\alpha}-(n+1)^{-\alpha}.
\]
由 $t^{-\alpha}$ 的凸性，$d_n(\alpha)$ 非负、递减至零。再次分部求和，利用 $c_n$ 的有界部分和，
\[
\left|\sum_{n=1}^\infty d_n(\alpha)c_n\right|
\leq C d_1(\alpha)=C(1-2^{-\alpha})\longrightarrow0.
\]
所以所求极限为 $\tfrac12\cot(1/2)$。

两次分部求和使用的是同一个事实：若 $C_j=\sum_{n=1}^jc_n$ 且 $|C_j|\le K$，$v_n\downarrow0$，则
\[
\sum_{n=1}^Mv_nc_n=C_Mv_M+\sum_{n=1}^{M-1}C_n(v_n-v_{n+1}),
\quad
\left|\sum_{n=1}^\infty v_nc_n\right|\le Kv_1.
\]
第一问从第 $N$ 项起应用它，尾部由 $N^{-1-\alpha}\le N^{-1}$ 统一控制；第二问则令 $v_n=d_n(\alpha)$，得到由首个差分控制的整个余项。这样既说明参数极限为何存在，也解释为什么不能把第二问直接代入 $\alpha=0$。

\end{solve}

\begin{remark}[方法与条件]第二问的逐项极限是 $\sin n$，而 $\sum\sin n$ 的通项不趋于零，普通意义下不收敛。因此这里得到的是一种加权求和值。对固定 $0<\theta<2\pi$，同一证明给出
\[
\lim_{\alpha\downarrow0}\sum_{n=1}^\infty
\frac{\sin(n\theta)}{n^\alpha}=\frac12\cot\frac\theta2.
\]
它也等于几何阻尼的 Abel 极限 $\lim_{r\uparrow1}\operatorname{Im}[re^{i\theta}/(1-re^{i\theta})]$。

分部求和与振荡积分分部积分的平行结构，见\hyperref[app:integration-exchange]{方法卷：积分交换与一致尾部}。\end{remark}

\section{初赛：第 14 届（2022 年）}

本节选入3道课程相关题；来源序号保留原题解编排，未选入题目不重新补号。

\begin{example}[来源序号2]\label{ex:ch04-14th-A:2}\label{pr:ch04-14th-A:2}

设 $\lim_{n\rightarrow+\infty}\frac{a_{n}}{n^{2}}=a$，$\lim_{n\rightarrow+\infty}\frac{b_{n}}{n^{2}}=b$。证明极限 $\lim_{n\rightarrow+\infty}\frac{1}{n^{5}}\sum_{k=0}^{n}a_{k}b_{n-k}$ 存在并求其值。

\par\href{https://github.com/lzxcr/cmc-mathematics/blob/64916fa1fdbc3e444a652eca204d5669058ee321/src/problems/preliminary/14/02.md}{题目与来源条目}

\end{example}

\begin{solve}\label{ans:ch04-14th-A:2}

卷积中的典型项有 $n^4$ 的尺度，又有约 $n$ 项，因而 $n^5$ 归一化会留下区间 $[0,1]$ 上的积分。为使两端的小下标也受到控制，写 $a_k=k^2(a+\varepsilon_k)$、$b_k=k^2(b+\eta_k)$（$k\geq1$），其中 $\varepsilon_k,\eta_k\to0$ 且有界。于是
\[
\frac1{n^5}\sum_{k=1}^{n-1}a_kb_{n-k}
=\frac{ab}{n}\sum_{k=1}^{n-1}
 \left(\frac kn\right)^2\left(1-\frac kn\right)^2+o(1).
\]
确实，展开后的误差绝对值不超过
\[
\frac Cn\sum_{k=1}^{n-1}(|\varepsilon_k|+|\eta_k|)\longrightarrow0
\]
（收敛数列的 Cesàro 平均仍收敛）。两端项 $(a_0b_n+a_nb_0)/n^5=O(n^{-3})$ 也趋于零。故由 Riemann 和，所求极限为
\[
ab\int_0^1t^2(1-t)^2\,\mathrm dt=\frac{ab}{30}.
\]

展开误差时出现 $a\eta_{n-k}+b\varepsilon_k+\varepsilon_k\eta_{n-k}$。由于两误差列有界，其绝对值不超过常数倍的 $|\varepsilon_k|+|\eta_{n-k}|$；再用 $0\le(k/n)^2(1-k/n)^2\le1$，便得所列平均值上界。指标变换 $j=n-k$ 使 $\eta$ 的求和恢复通常顺序。最后 $\int_0^1t^2(1-t)^2\,dt=1/3-1/2+1/5=1/30$，即使 $a$ 或 $b$ 为零，同样的误差估计仍然成立。

\end{solve}

\begin{remark}[方法与条件]更一般地，若 $a_n/n^r\to a$、$b_n/n^s\to b$，其中 $r,s\geq0$，则同样的 Riemann 和论证给出
\[
\frac{\sum_{k=0}^n a_kb_{n-k}}{n^{r+s+1}}
\longrightarrow ab\int_0^1t^r(1-t)^s\,\mathrm dt
=ab\,B(r+1,s+1).
\]
卷积把两个幂次相加再增加一阶，常数由 Beta 积分决定。

离散归一化和式与连续极限的关系，见\hyperref[app:integration-riemann]{方法卷：Riemann 和与测度}。\end{remark}

\begin{example}[来源序号5]\label{ex:ch04-14th-A:5}\label{pr:ch04-14th-A:5}

设 $\lim_{n\rightarrow+\infty}\beta_{n}=0$，函数 $f$ 在 $[-1,2]$ 上有界，在 $[0,1]$ 上 Riemann 可积。证明：

\[
\lim_{n\rightarrow+\infty}\frac{1}{n}\sum_{k=1}^{n}f\left(\frac{k}{n}+\beta_{n}\right)=\int_{0}^{1}f(x)\mathrm{d}x.
\]

\par\href{https://github.com/lzxcr/cmc-mathematics/blob/64916fa1fdbc3e444a652eca204d5669058ee321/src/problems/preliminary/14/05.md}{题目与来源条目}

\end{example}

\begin{solve}\label{ans:ch04-14th-A:5}

$f$ 未必连续，故不能逐项比较 $f(k/n+\beta_n)$ 与 $f(k/n)$。可比较的是每段区间获得的采样点比例，再用 Riemann 可积性的上下阶梯夹逼。记 $t_{n,k}=k/n+\beta_n$。落在 $[0,1]$ 外的采样点至多有 $n|\beta_n|+2$ 个，其贡献为 $o(1)$，因为 $f$ 在 $[-1,2]$ 上有界。

对任意固定区间 $I\subset[0,1]$，等距网格的计数满足
\[
\frac1n\#\{1\leq k\leq n:t_{n,k}\in I\}\longrightarrow |I|.
\]
这是因为平移前后的截断区间长度相差至多 $2|\beta_n|$，且网格计数与“长度乘以 $n$”至多相差常数。因此结论先对阶梯函数成立。

由 Riemann 可积性，任给 $\varepsilon>0$，可在 $[0,1]$ 上取阶梯函数 $g\leq f\leq h$ 且 $\int_0^1(h-g)<\varepsilon$。将内部采样和夹在 $g,h$ 的采样和之间，再令 $n\to\infty$、$\varepsilon\downarrow0$，即得所求极限。

上下阶梯函数可由同一个固定划分上 $f$ 的上、下确界构造，二者积分差就是 Darboux 上下和之差。采样恰落在有限个分点时，可把这些点单独处理：每个固定分点至多对应一个采样点，总贡献为 $O(1/n)$。外部采样的绝对贡献不超过 $\|f\|_\infty(|\beta_n|+2/n)$，故无需假设 $f$ 在 $[0,1]$ 以外可积或连续。所有估计都从充分大的 $n$ 开始，此时采样点已落在 $[-1,2]$ 内。

\end{solve}

\begin{remark}[方法与条件]本题与\hyperref[pr:ch01-17th-A:6]{具有渐近密度的离散取样} 都先证明区间计数，再扩展测试函数。统一观点见\hyperref[app:integration-riemann]{方法卷：Riemann 和与测度}。\end{remark}

\begin{example}[来源序号6]\label{ex:ch04-14th-A:6}\label{pr:ch04-14th-A:6}

设 $f$ 在 $[0,+\infty)$ 的任意闭区间上 Riemann 可积。对于 $x\geq 0$，定义 $F(x)=\int_{0}^{x}t^{\alpha}f(t+x)\mathrm{d}t$。

（1）若 $\alpha\in(-1,0)$ 且 $\lim_{x\rightarrow+\infty}f(x)=A$，证明：$F$ 在 $[0,+\infty)$ 上一致连续。

（2）若 $\alpha\in(0,1)$，$f$ 以 $T>0$ 为周期，$\int_{0}^{3}f(t)\mathrm{d}t=2022$。证明：$F$ 在 $[0,+\infty)$ 上非一致连续。

\par\href{https://github.com/lzxcr/cmc-mathematics/blob/64916fa1fdbc3e444a652eca204d5669058ee321/src/problems/preliminary/14/06.md}{题目与来源条目}

\end{example}

\begin{solve}\label{ans:ch04-14th-A:6}

（1）不对可能不连续的 $f$ 求导，而把平移转移到积分核。题设保证 $M=\sup_{t\geq0}|f(t)|<\infty$。设 $y=x+h$，$h>0$，引入
\[
B=\int_h^{x+h}t^\alpha f(t+x)\,\mathrm dt
=\int_0^x(t+h)^\alpha f(t+y)\,\mathrm dt.
\]
分别比较 $B$ 与 $F(x),F(y)$，利用 $t^\alpha$ 递减，得到
\[
|F(y)-F(x)|\leq M\left[
\int_0^h t^\alpha\,\mathrm dt
+2\int_x^{x+h}t^\alpha\,\mathrm dt
+\int_0^x\bigl(t^\alpha-(t+h)^\alpha\bigr)\,\mathrm dt\right]
\leq\frac{4M}{1+\alpha}h^{1+\alpha}.
\]
故 $F$ 实际上全局 Hölder 连续。这一结论只需 $f$ 有界。

（2）令 $m=T^{-1}\int_0^T f$，$G(s)=\int_0^s(f(t)-m)\,\mathrm dt$。则 $G$ 连续、有界且以 $T$ 为周期。分部积分给出
\[
F(x)=\frac{m}{1+\alpha}x^{1+\alpha}
+x^\alpha G(2x)
-\alpha\int_0^x t^{\alpha-1}G(t+x)\,\mathrm dt.
\]
若 $m\neq0$，后两项为 $O(x^\alpha)$，故 $F$ 超线性增长。一致连续函数在半直线上至多线性增长（将长区间分成固定短区间并累加增量），所以 $F$ 不一致连续。

若 $m=0$，超线性均值项消失，必须改查振荡的幅度。记最后一项为 $R(x)$。因 $\alpha-1\in(-1,0)$，（1）的估计说明 $R$ 一致连续。另一方面，$G(3)=2022\neq G(0)$，所以 $G$ 非常值。任给 $\delta>0$，可取 $X\geq0$、$0<h<\delta$ 使 $G(2X+2h)\neq G(2X)$。利用周期性，
\[
\left|(X+jT+h)^\alpha G(2X+2h)
 -(X+jT)^\alpha G(2X)\right|\longrightarrow\infty.
\]
因此 $x^\alpha G(2x)$ 不一致连续；与一致连续的 $R$ 相加后，$F$ 仍不一致连续。

第一问中核的平移误差满足
\[
\int_0^x[t^\alpha-(t+h)^\alpha]\,dt
=\int_0^h t^\alpha\,dt-\int_x^{x+h}t^\alpha\,dt
\le\frac{h^{1+\alpha}}{1+\alpha},
\]
这也说明奇点没有造成不可控误差。第二问若 $G$ 非常值，则它不可能在每个长度小于任意给定 $2\delta$ 的区间上都恒定，否则用相邻小区间连接任意两点就会推出它处处相等。选定使差值 $d=G(2X+2h)-G(2X)\ne0$ 的两点后，振荡项的增量为 $(X+jT)^\alpha d+o(1)$，故绝对值趋于无穷；一致连续的 $R$ 在固定的足够小步长下增量统一有界，不能抵消这一增长。

\end{solve}

\begin{remark}[方法与条件]第 （1）问展示了积分核的平移正则化：即使 $f$ 不连续，弱奇异核 $t^\alpha$ 的积分仍获得 $1+\alpha$ 阶 Hölder 控制。第 （2）问的两种障碍不同：非零均值产生超线性主项；零均值时，周期振荡被 $x^\alpha$ 放大。条件 $\int_0^3f\neq0$ 保证后一种情况下原函数不是常数。\end{remark}

\section{初赛：第 14 届补赛（2022 年）}

本节选入3道课程相关题；来源序号保留原题解编排，未选入题目不重新补号。

\begin{remark}[来源说明]考试形式：闭卷；考试时间：150 分钟；满分：100 分。\end{remark}

\begin{example}[来源序号2]\label{ex:ch05-14th-A-makeup:2}\label{pr:ch05-14th-A-makeup:2}

设 $f$ 是 $[0,1]$ 上的凸函数，求证：

\[
\int_0^1 t(1-t)f(t) \mathrm{d}t \leq \frac{1}{3} \int_0^1 \left(t^3+(1-t)^3\right)f(t) \mathrm{d}t.
\]

\par\href{https://github.com/lzxcr/cmc-mathematics/blob/64916fa1fdbc3e444a652eca204d5669058ee321/src/problems/preliminary/14-makeup/02.md}{题目与来源条目}

\end{example}

\begin{solve}\label{ans:ch05-14th-A-makeup:2}

希望用凸性比较 $f(t)$ 与两侧点的函数值，因而选一对以 $t$ 为重心的点；让它们分别扫过 $[0,t]$ 和 $[t,1]$，便会产生题中的积分权重。对 $0<t<1$、$0\leq s\leq1$，恒等式
\[
t=(1-t)(ts)+t(1-s+ts)
\]
和凸性给出 $f(t)\leq(1-t)f(ts)+tf(1-s+ts)$。对 $s$ 积分并乘以 $t(1-t)$，得到
\[
t(1-t)f(t)\leq(1-t)^2\int_0^t f(u)\,\mathrm du
+t^2\int_t^1f(u)\,\mathrm du.
\]
再对 $t$ 积分、交换三角形区域上的积分次序，便得
\[
\int_0^1t(1-t)f(t)\,\mathrm dt
\leq\frac13\int_0^1\bigl((1-u)^3+u^3\bigr)f(u)\,\mathrm du.
\]
有限凸函数在闭区间上有界、内部连续，故这里的积分与换序均合法；端点可能的跳跃不影响积分。

两次换元分别给出 $\int_0^1f(ts)\,ds=t^{-1}\int_0^tf(u)\,du$ 和 $\int_0^1f(1-s+ts)\,ds=(1-t)^{-1}\int_t^1f(u)\,du$，所以乘上 $t(1-t)$ 后恰出现平方权重。换序后，第一个三角形的内积分为 $\int_u^1(1-t)^2\,dt=(1-u)^3/3$，第二个为 $\int_0^ut^2\,dt=u^3/3$。这逐项产生了题目右端的三次权重，并不需要 $f$ 非负。

\end{solve}

\begin{remark}[方法与条件]两边权重的总质量与一阶矩相同，因此仿射函数对差值没有贡献。若 $f\in C^2$，两次分部积分可把差值写为
\[
\int_0^1\left[t(1-t)-\frac{t^3+(1-t)^3}{3}\right]f(t)\,\mathrm dt
=-\frac16\int_0^1t^2(1-t)^2f''(t)\,\mathrm dt.
\]
右端将不等式归结为“非负曲率与非正核的配对”。一般凸函数未必有二阶导数，主解的凸组合论证无需这项额外正则性；用折线与 $(t-s)_+$ 统一两种证明的方式，见\hyperref[app:estimate-convexity]{方法卷：凸序与积分核}。\end{remark}

\begin{example}[来源序号5]\label{ex:ch05-14th-A-makeup:5}\label{pr:ch05-14th-A-makeup:5}

设 $x_0,y_0 \in \mathbb{R}$，对于 $n \geq 0$，定义

\[
\begin{cases}
x_{n+1} = \frac{1}{x_n^2+x_n y_n+2 y_n^2+1}, \\
y_{n+1} = \frac{1}{2 x_n^2+x_n y_n+y_n^2+1}.
\end{cases}
\]

证明：（1）$\lim_{n \to +\infty} \left(x_n-y_n\right) = 0$。（2）点列 $\{x_n\}$ 收敛。

\par\href{https://github.com/lzxcr/cmc-mathematics/blob/64916fa1fdbc3e444a652eca204d5669058ee321/src/problems/preliminary/14-makeup/05.md}{题目与来源条目}

\end{example}

\begin{solve}\label{ans:ch05-14th-A-makeup:5}

（1）交换 $x,y$ 会交换两条递推式，故先看反对称量 $x-y$，检验两分量是否趋于同步。两个分母均为正定二次型加 $1$，故 $0<x_n,y_n\leq1$（$n\geq1$）。记 $s=x_n+y_n$，$d=x_n-y_n$，通分可得
\[
|x_{n+1}-y_{n+1}|
\leq\frac{s}{1+3x_n^2+2x_ny_n+3y_n^2}|d|
=\frac{s}{1+2s^2+d^2}|d|
\leq\frac12|d|.
\]
最后一步由 $1+2s^2\geq2s$ 得到。因此差值几何衰减至零。

（2）由 （1），$x_{n+1}=T(x_n)+o(1)$，其中 $T(t)=1/(1+4t^2)$。设 $L=\limsup x_n$、$\ell=\liminf x_n$。对连续递减函数，下极限处对应上极限，反之亦然（可由逼近上、下极限的子列及夹逼得到）。因此
\[
L=\frac1{1+4\ell^2},\qquad
\ell=\frac1{1+4L^2}.
\]
相减得 $(L-\ell)(1-4L\ell)=0$。若 $L\neq\ell$，则 $4L\ell=1$，代回得 $L+\ell=1$，从而 $(L-\ell)^2=(L+\ell)^2-4L\ell=0$，矛盾。因此数列收敛。

极限满足 $4L^3+L=1$，左端严格递增且 $L=1/2$ 是解，故 $x_n,y_n\to1/2$。

分母正定可由 $x^2+xy+2y^2=(x+y/2)^2+7y^2/4$ 及交换 $x,y$ 后的式子确认。进入 $[0,1]^2$ 后，分母至多为 $5$，故从第二步起还有 $x_n,y_n\ge1/5$，上、下极限均为正。由 $L(1+4\ell^2)=1$、$\ell(1+4L^2)=1$ 相减得到 $(L-\ell)(1-4L\ell)=0$；若后一因子为零，第一式变成 $L+\ell=1$，再结合乘积 $L\ell=1/4$ 就只能有 $L=\ell=1/2$。

\end{solve}

\begin{remark}[方法与条件]先证明两个分量同步，再研究一维映射的可能二周期，是耦合递推的有效路径。本题在极限点的 Jacobian 特征值为 $-1$ 和 $1/4$，不能直接套用严格压缩原理；上下极限方程恰好排除了非平凡二周期。\end{remark}

\begin{example}[来源序号6]\label{ex:ch05-14th-A-makeup:6}\label{pr:ch05-14th-A-makeup:6}

设 $\psi(s) = \int_0^{+\infty} \frac{\ln(1+sx)}{x(1+x)} \mathrm{d}x$。证明：

（1）$\psi(1) = \frac{\pi^2}{6}$。

（2）对于 $s > 0$，有 $\psi(s) + \psi\left(\frac{1}{s}\right) = \frac{\pi^2}{3} + \frac{1}{2} \ln^2 s$。

\par\href{https://github.com/lzxcr/cmc-mathematics/blob/64916fa1fdbc3e444a652eca204d5669058ee321/src/problems/preliminary/14-makeup/06.md}{题目与来源条目}

\end{example}

\begin{solve}\label{ans:ch05-14th-A-makeup:6}

对数含参数而分母有理，先把对数写成关于参数的积分，便可在非负情形下合法换序，得到更简单的有理积分。对 $s>0$，使用 $\ln(1+sx)=\int_0^s x/(1+tx)\,\mathrm dt$ 及非负积分的换序，得到
\[
\psi(s)=\int_0^s\!\mathrm dt\int_0^\infty
\frac{\mathrm dx}{(1+x)(1+tx)}
=\int_0^s\frac{\ln t}{t-1}\,\mathrm dt.
\]
有理积分可先取有限上限 $R$：当 $t\ne1$ 时，
\[
\int_0^R\frac{dx}{(1+x)(1+tx)}
=\frac{\ln(1+R)-\ln(1+tR)}{1-t}
\longrightarrow\frac{\ln t}{t-1}.
\]
$t=1$ 时直接积分得 $1$，也等于右端的连续极限。参数积分在零端仅有 $-\ln t$ 型奇点，因而有限；这同时证实了原积分收敛。

因此
\[
\psi(1)=\int_0^1\frac{-\ln t}{1-t}\,\mathrm dt
=\sum_{k=0}^\infty\int_0^1(-\ln t)t^k\,\mathrm dt
=\sum_{k=1}^\infty\frac1{k^2}=\frac{\pi^2}{6}.
\]
又 $\psi'(s)=\ln s/(s-1)$，其中 $s=1$ 处取连续值。令 $H(s)=\psi(s)+\psi(1/s)$，链式法则给出
\[
H'(s)=\frac{\ln s}{s-1}-\frac{\ln s}{s(s-1)}
=\frac{\ln s}{s}.
\]
从 $1$ 积分至 $s$，并用 $H(1)=\pi^2/3$，即得
\[
\psi(s)+\psi(1/s)=\frac{\pi^2}{3}+\frac12\ln^2s.
\]

\end{solve}

\begin{remark}[方法与条件]将 $\psi(s)=\pi^2/6-\operatorname{Li}_2(1-s)$ 代入恒等式，得到
\[
\operatorname{Li}_2(1-s)+\operatorname{Li}_2(1-1/s)
=-\frac12\ln^2s\qquad(s>0).
\]
这是一条实变量二重对数恒等式。其定义与积分表示见\hyperref[pr:ch02-16th-A:2]{相关参数积分题}；两题由代换 $x\mapsto x^2$ 直接联系。\end{remark}

\section{初赛：第 14 届第二次补赛（2022 年）}

本节选入3道课程相关题；来源序号保留原题解编排，未选入题目不重新补号。

\begin{example}[来源序号2]\label{ex:ch06-14th-A-makeup2:2}\label{pr:ch06-14th-A-makeup2:2}

设 $\Gamma=\{(x_n)_{n\geq1}:x_n\in\{0,2\}\}$，并令
\[
\Pi(x)=\sum_{n=1}^\infty\frac{x_n}{3^n}.
\]
对 $x\in\Gamma$ 定义
\[
f(x)=\limsup_{n\to\infty}\frac{x_1+\cdots+x_n}{n}.
\]
证明：（1）$\Pi$ 是单射；（2）$\Pi(\Gamma)$ 的每一点都是其聚点；（3）$f(\Gamma)=[0,2]$。

\par\href{https://github.com/lzxcr/cmc-mathematics/blob/64916fa1fdbc3e444a652eca204d5669058ee321/src/problems/preliminary/14-makeup2/02.md}{题目与来源条目}

\end{example}

\begin{solve}\label{ans:ch06-14th-A-makeup2:2}

（1）若 $x,y$ 首次在第 $k$ 位不同，则
\[
|\Pi(x)-\Pi(y)|\geq\frac2{3^k}-\sum_{j>k}\frac2{3^j}
=\frac1{3^k}>0.
\]
故 $\Pi$ 为单射。

（2）将 $x$ 的第 $k$ 位翻转得到 $x^{(k)}$。则 $\Pi(x^{(k)})\neq\Pi(x)$，且距离为 $2/3^k\to0$，故每点都是聚点。

（3）每个均值均在 $[0,2]$，故 $f(\Gamma)\subseteq[0,2]$。反之，给定 $\alpha\in[0,2]$，让前 $n$ 项之和以有界误差跟踪 $n\alpha$，即可使均值趋于 $\alpha$。具体地，令 $p=\alpha/2$，取
\[
x_n=2\bigl(\lfloor np\rfloor-\lfloor(n-1)p\rfloor\bigr)\in\{0,2\}.
\]
则 $\sum_{j=1}^n x_j=2\lfloor np\rfloor$，所以其均值趋于 $2p=\alpha$，从而 $f(x)=\alpha$。

第一问估计的是“首个不同数字的贡献减去所有后续数字可能抵消的最大量”，而后面的几何级数之和只有 $3^{-k}$，小于首项的 $2\cdot3^{-k}$，所以三进制表示在限定数字集 $\{0,2\}$ 中没有歧义。第三问中 $0\le p\le1$ 保证相邻两个整数部分相差 $0$ 或 $1$，故构造合法；且 $|2\lfloor np\rfloor/n-2p|<2/n$，实际上均值有普通极限，因而其上极限也等于所求值。

\end{solve}

\begin{remark}[方法与条件]像集 $\Pi(\Gamma)$ 是 Cantor 三分集。乘积空间 $\{0,2\}^{\mathbb N}$ 紧，编码 $\Pi$ 连续且单射，所以它与 Cantor 集同胚。

上极限使 $f$ 对每个编码都有定义，即使其均值不收敛。例如交替连接 $0$ 块与 $2$ 块，使每块长度与之前总长度之比趋于无穷，则均值的下极限为 $0$、上极限为 $2$，故 $f(x)=2$。若记 $S=\{n:x_n=2\}$，则 $f(x)$ 正是 $S$ 的上自然密度的两倍。有关上、下极限的运算，见\hyperref[app:limit-limsup]{方法卷：上下极限与子列}。

第（3）问的取整构造给出更强的结论：对每个 $\alpha\in[0,2]$，都有编码满足 $|\sum_{j=1}^n x_j-n\alpha|<2$，从而均值本身收敛到 $\alpha$。\end{remark}

\begin{example}[来源序号5]\label{ex:ch06-14th-A-makeup2:5}\label{pr:ch06-14th-A-makeup2:5}

设 $f$ 与 $g$ 在 $[a,b]$ 上可导，且对任何 $x\in[a,b]$，$g^{\prime}(x)\neq 0$。又
\[
\int_{a}^{b}f(x)\mathrm{d}x=\int_{a}^{b}f(x)g(x)\mathrm{d}x=0.
\]
证明：存在 $\xi\in(a,b)$ 使得 $f^{\prime}(\xi)=0$。

\par\href{https://github.com/lzxcr/cmc-mathematics/blob/64916fa1fdbc3e444a652eca204d5669058ee321/src/problems/preliminary/14-makeup2/05.md}{题目与来源条目}

\end{example}

\begin{solve}\label{ans:ch06-14th-A-makeup2:5}

第一个积分条件让原函数在两端为零；第二个条件需要再制造一个函数，使其导数包含这个原函数。这样可逐次增加零点并使用 Rolle 定理。设 $F(x)=\int_a^x f(t)\,\mathrm dt$，并令
\[
H(x)=g(x)F(x)-\int_a^x f(t)g(t)\,\mathrm dt.
\]
两个积分条件给出 $F(a)=F(b)=H(a)=H(b)=0$，且 $H'=g'F$。由 Rolle 定理，存在 $c\in(a,b)$ 使 $H'(c)=0$；因 $g'(c)\neq0$，故 $F(c)=0$。

$F$ 在 $a,c,b$ 三点为零，连续应用 Rolle 定理，先得到 $f=F'$ 的两个内点零点，再得到它们之间的 $f'$ 零点，即为所求。

将求导写开便可看到辅助函数的作用：$H'=g'F+gF'-fg=g'F$，其中 $F'=f$ 使后两项完全抵消。得到 $F(c)=0$ 后，分别在 $[a,c]$、$[c,b]$ 上用 Rolle 定理，存在 $u<c<v$ 使 $f(u)=f(v)=0$；再在 $[u,v]$ 上对 $f$ 使用 Rolle 定理，得到 $f'(\xi)=0$。题设可导保证所需连续性，整个论证无需假设 $g'$ 连续。

\end{solve}

\begin{remark}[方法与条件]证明避开了对 $g'$ 的积分，只需其点态存在且不为零。相关的普遍原理是正交性迫使振荡：连续非零函数若与 $1,x,\ldots,x^{m-1}$ 都正交，则至少变号 $m$ 次，否则可用这些变号点构造次数小于 $m$ 的多项式，使其与原函数的乘积不变号，从而与正交性矛盾。\end{remark}

\begin{example}[来源序号6]\label{ex:ch06-14th-A-makeup2:6}\label{pr:ch06-14th-A-makeup2:6}

设 $A=\left\{\sqrt{\frac{m}{n}}\ \middle|\ m,n\in\mathbb{Z}^{+}\right\}\setminus\mathbb{Q}$，$x_{0}=\sum_{n=0}^{\infty}\frac{1}{10^{n!}}$。对于 $x>0$，定义
\[
f(x)=\begin{cases}0, & x\text{ 为无理数},\\ \frac{1}{q^{\alpha}}, & x=\frac{p}{q}\text{ 为既约分数}.\end{cases}
\]
证明：（1）对任何 $x\in A$，存在常数 $M_{x}>0$，使得对任何既约分数 $\frac{p}{q}$ 都有 $\left|x-\frac{p}{q}\right|\geq\frac{M_{x}}{q^{2}}$；（2）$f$ 在 $A$ 中的每个点可微的充要条件是 $\alpha>2$；（3）对任何 $\alpha$，函数 $f$ 在 $x_{0}$ 处均不可微。

\par\href{https://github.com/lzxcr/cmc-mathematics/blob/64916fa1fdbc3e444a652eca204d5669058ee321/src/problems/preliminary/14-makeup2/06.md}{题目与来源条目}

\end{example}

\begin{solve}\label{ans:ch06-14th-A-makeup2:6}

（1）固定 $x=\sqrt{m/n}\notin\mathbb Q$。若 $|x-p/q|<1$，则 $|x+p/q|\leq2x+1$，且非零整数 $mq^2-np^2$ 的绝对值至少为 $1$。于是
\[
\left|x-\frac pq\right|
=\frac{|mq^2-np^2|}{nq^2|x+p/q|}
\geq\frac1{n(2x+1)q^2}.
\]
若 $|x-p/q|\geq1$，结论也成立。可取 $M_x=\min\{1,[n(2x+1)]^{-1}\}$。

（2）沿无理数趋近时差商恒为零，因此若导数存在，只能等于零；问题转为比较有理点函数值 $q^{-\alpha}$ 与其到 $x$ 的距离。若 $\alpha>2$，有理数 $t=p/q\to x$ 时，其既约分母 $q\to\infty$，故
\[
\left|\frac{f(t)-f(x)}{t-x}\right|
\leq M_x^{-1}q^{2-\alpha}\longrightarrow0.
\]
无理数 $t$ 对应的差商恒为零，所以 $f'(x)=0$。

分母为何必须趋于无穷，也需说明：在 $x$ 的某个有界邻域内，分母不超过固定整数 $Q$ 的既约分数只有有限个；它们都不等于无理数 $x$，所以与 $x$ 的距离有正的最小值。一个趋于 $x$ 的有理数列最终会避开这些点，任意 $Q$ 都如此，故分母趋于无穷。

若 $\alpha\leq2$，Dirichlet 逼近定理给出无穷多个趋于 $x$ 的既约分数 $p/q$，满足 $|x-p/q|<q^{-2}$。沿此列差商绝对值大于 $q^{2-\alpha}\geq1$；沿无理数趋近则为零，故导数不存在。

（3）$x_0$ 的十进制展开不最终周期，因而无理。对 $N\geq2$，记
\[
t_N=\sum_{j=0}^N10^{-j!}=p_N/q_N.
\]
分子末位为 $1$，所以既约分母为 $q_N=10^{N!}$；并且
\[
0<x_0-t_N\leq2\,10^{-(N+1)!}.
\]
于是对任意固定实数 $\alpha$，
\[
\left|\frac{f(t_N)-f(x_0)}{t_N-x_0}\right|
\geq\frac12\,10^{(N+1-\alpha)N!}\longrightarrow\infty.
\]
故 $f$ 在 $x_0$ 处不可微。

Dirichlet 逼近可由抽屉原理理解：将 $0,x,2x,\ldots,Qx$ 的小数部分分入 $Q$ 个等长小区间，可找到 $1\le q\le Q$ 及整数 $p$ 使 $|qx-p|\le1/Q$，于是 $|x-p/q|\le1/(qQ)\le1/q^2$。约分只会使所需上界放宽；随着 $Q\to\infty$，无理性迫使出现无穷多个不同的逼近分数。第三问的阶乘位数则使误差远小于任意固定的分母幂次，这正是它与二次无理数的区别。

\end{solve}

\begin{remark}[方法与条件]对无理数 $x$，定义无理性指数 $\mu(x)$ 为 $|x-p/q|<q^{-\nu}$ 有无穷多既约解的指数 $\nu$ 的上确界。若 $\alpha>\mu(x)$，上述差商趋于零；若 $\alpha<\mu(x)$，则不可微。临界情形需额外信息。本题二次无理数的指数为 $2$，而 $x_0$ 的指数无穷大。

Roth 定理把“指数为 $2$”推广到全部代数无理数，但不给出 （1）中统一的 $M_x/q^2$ 下界；后者是二次无理数的更强性质。\end{remark}

\section{初赛：第 15 届（2023 年）}

本节选入3道课程相关题；来源序号保留原题解编排，未选入题目不重新补号。

\begin{example}[来源序号2]\label{ex:ch03-15th-A:2}\label{pr:ch03-15th-A:2}

设 $f(x,y,z) = x^2 + \left(y^2 + z^2\right)(1-x)^3$。

（1）计算 $f$ 的驻点。

（2）求 $f$ 在 $\Sigma$ 上的最小值，其中 $\Sigma$ 是 $\left\{(x,y,z) \mid |x| \leq 2, y^2 + z^2 \leq 4\right\}$ 的边界。

（3）求 $f$ 在椭球 $x^2 + \frac{y^2}{2} + \frac{z^2}{3} \leq 1$ 上的最小值。

\par\href{https://github.com/lzxcr/cmc-mathematics/blob/64916fa1fdbc3e444a652eca204d5669058ee321/src/problems/preliminary/15/02.md}{题目与来源条目}

\end{example}

\begin{solve}\label{ans:ch03-15th-A:2}

（1）驻点方程为
\[
2x-3(y^2+z^2)(1-x)^2=0,\qquad
 y(1-x)^3=z(1-x)^3=0.
\]
$x=1$ 不满足第一式，故 $y=z=0$，继而 $x=0$。唯一驻点为 $(0,0,0)$。

（2）先观察 $(1-x)^3$ 的符号，再利用 $y^2+z^2\leq4$；这可直接给出全域下界，免去对各片边界分别求驻点。实际上，在整个实心圆柱内都有 $f\geq0$：当 $x\leq1$ 时各项非负；当 $1\leq x\leq2$ 时，
\[
f\geq x^2-4(x-1)^3
\geq x^2-4(x-1)^2=(2-x)(3x-2)\geq0.
\]
在边界圆周 $x=2,\ y^2+z^2=4$ 上取到 $0$。故所求最小值为 $0$。

（3）椭球内 $-1\leq x\leq1$，所以 $f\geq0$；原点属于椭球且 $f(0,0,0)=0$。最小值仍为 $0$。

第二问中当 $x>1$ 时，$(1-x)^3$ 为负，故将 $y^2+z^2$ 换成最大值 $4$ 才得到下界；不等号方向不能沿用正系数情形。又 $0\le x-1\le1$，所以 $(x-1)^3\le(x-1)^2$，这解释了随后的代数估计。虽然原点使 $f=0$，它不在圆柱边界上，因此必须另行指出边界圆周上的取等点；第三问的约束是实心椭球，原点才是合法取等点。

\end{solve}

\begin{example}[来源序号5]\label{ex:ch03-15th-A:5}\label{pr:ch03-15th-A:5}

设 $E \subset \mathbb{R}^n$ 非空有界，$c \in \mathbb{R}^n$，$c$ 非零。用 $\operatorname{diam} E = \sup_{x,y \in E} \|x - y\|$ 表示 $E$ 的直径，记 $E + c = \{x + c \mid x \in E\}$。证明：$\operatorname{diam} E < \operatorname{diam}(E \cup (E + c))$。

\par\href{https://github.com/lzxcr/cmc-mathematics/blob/64916fa1fdbc3e444a652eca204d5669058ee321/src/problems/preliminary/15/05.md}{题目与来源条目}

\end{example}

\begin{solve}\label{ans:ch03-15th-A:5}

记 $D=\operatorname{diam}(E\cup(E+c))$。任取 $x,y\in E$，取正、反两种跨集合配对，便得到位移为 $x-y+c$ 和 $x-y-c$ 的两段；它们的距离都不超过 $D$。将平方相加可消去未知符号的交叉项。由平行四边形恒等式，
\[
D^2\geq\frac{\|x-y+c\|^2+\|x-y-c\|^2}{2}
=\|x-y\|^2+\|c\|^2.
\]
对 $x,y$ 取上确界，得到更强的估计
\[
D\geq\sqrt{(\operatorname{diam}E)^2+\|c\|^2}
>\operatorname{diam}E.
\]

两种跨集合的距离分别是 $\|(x+c)-y\|$ 和 $\|x-(y+c)\|$，两端确实都属于并集。令 $d=\operatorname{diam}E$，即使 $E$ 不闭、直径不由某对点取到，仍可选取 $x_j,y_j\in E$ 使 $\|x_j-y_j\|\to d$；将平方估计逐对应用并取极限，得到 $D^2\ge d^2+\|c\|^2$。因 $E$ 有界，$d,D$ 有限；因 $c\ne0$，最后的平方根严格大于 $d$。

\end{solve}

\begin{remark}[方法与条件]这个证明逐字适用于任意实或复 Hilbert 空间，既不需要有限维，也不需要直径能够取到。定量下界来自内积范数的平行四边形恒等式。\end{remark}

\begin{example}[来源序号6]\label{ex:ch03-15th-A:6}\label{pr:ch03-15th-A:6}

设 $a = \sqrt[3]{3}$，$x_1 = a$，$x_{n+1} = a^{x_n}\ (n = 1, 2, \ldots)$。证明：数列 $\left\{x_n\right\}_{n=1}^{\infty}$ 极限存在，但不是 3。

\par\href{https://github.com/lzxcr/cmc-mathematics/blob/64916fa1fdbc3e444a652eca204d5669058ee321/src/problems/preliminary/15/06.md}{题目与来源条目}

\end{example}

\begin{solve}\label{ans:ch03-15th-A:6}

$3$ 虽是不动点，却未必是从给定初值出发所到达的不动点。只要找到小于 $3$ 的不变上界并证明单调性，就能同时解决两个要求。函数 $h(t)=\ln t/t$ 在 $(0,e)$ 上递增，在 $(e,\infty)$ 上递减，故 $\ln3/3<1/e$。于是
\[
a=3^{1/3}<e^{1/e}<e,\qquad a^e<e.
\]
由于 $t\mapsto a^t$ 递增，$x_2>x_1$ 蕴含整个数列递增；又 $x_1<e$ 和 $a^e<e$ 保证 $x_n<e$ 对所有 $n$ 成立。因此 $x_n$ 收敛，且其极限 $L\leq e<3$。

首先 $a>1$ 且 $x_1=a>1$，所以 $x_2=a^a>a=x_1$。若 $x_n>x_{n-1}$，指数函数的严格递增性给出 $x_{n+1}>x_n$，完成单调归纳。另一方面 $h'(t)=(1-\ln t)/t^2$，$3>e$ 使 $h(3)<h(e)=1/e$，即 $e\ln a<1$，所以 $a^e<e$。若仅由极限方程 $L=a^L$ 认出 $L=3$ 是一根，并不能断言这就是迭代极限；不变区间提供了排除该根的依据。

\end{solve}

\begin{remark}[方法与条件]不动点方程 $L=a^L$ 等价于 $\ln L/L=\ln3/3$，它有两个正根：一个在 $(1,e)$ 内，另一个是 $3$。本题迭代收敛到较小根。对 $T(t)=a^t$，在不动点有 $T'(L)=\ln L$，所以较小根是吸引点，而 $3$ 因 $\ln3>1$ 为排斥点。\end{remark}

\section{初赛：第 16 届（2024 年）}

本节选入3道课程相关题；来源序号保留原题解编排，未选入题目不重新补号。

\begin{example}[来源序号2]\label{ex:ch02-16th-A:2}\label{pr:ch02-16th-A:2}

设 $\varphi(t)=\int_{0}^{+\infty} \frac{\ln(1+tx^{2})}{x(1+x^{2})} \mathrm{d} x, t \geq 0$，求 $\varphi(1), \varphi(2)$。

\par\href{https://github.com/lzxcr/cmc-mathematics/blob/64916fa1fdbc3e444a652eca204d5669058ee321/src/problems/preliminary/16/02.md}{题目与来源条目}

\end{example}

\begin{solve}\label{ans:ch02-16th-A:2}

分母与对数都含 $x^2$，先统一变量，再把参数从对数中取出，可以把计算归结为有理函数积分。作代换 $u=x^2$，再用 $\ln(1+tu)=\int_0^t u/(1+su)\,\mathrm ds$。被积函数非负，Tonelli 定理允许交换积分：
\[
\varphi(t)=\frac12\int_0^t\!\mathrm ds
\int_0^\infty\frac{\mathrm du}{(1+u)(1+su)}
=\frac12\int_0^t\frac{\ln s}{s-1}\,\mathrm ds.
\]
其中 $s=1$ 处取连续延拓值 $1$，$s=0$ 处的对数奇点可积。因此
\[
\varphi(1)=\frac12\sum_{k=0}^\infty
\int_0^1(-\ln s)s^k\,\mathrm ds
=\frac12\sum_{k=1}^\infty\frac1{k^2}
=\frac{\pi^2}{12}.
\]
又令 $r=s-1$，
\[
\varphi(2)-\varphi(1)
=\frac12\int_0^1\frac{\ln(1+r)}r\,\mathrm dr
=\frac12\sum_{k=1}^\infty\frac{(-1)^{k+1}}{k^2}
=\frac{\pi^2}{24}.
\]
故 $\varphi(2)=\pi^2/8$。后一项积分可先在 $[0,\rho]$ 上展开，再令 $\rho\uparrow1$；积分后的级数绝对收敛。

换元 $u=x^2$ 时 $dx/x=du/(2u)$，这正是整段计算中因子 $1/2$ 的来源。求交错平方倒数和时，将偶数项从全部平方倒数和中扣除两次：
\[
\sum_{k\ge1}\frac{(-1)^{k+1}}{k^2}
=\sum_{k\ge1}\frac1{k^2}-2\sum_{j\ge1}\frac1{(2j)^2}
=\frac12\cdot\frac{\pi^2}{6}=\frac{\pi^2}{12}.
\]
因此 $\varphi(2)-\varphi(1)=\pi^2/24$，加上 $\pi^2/12$ 才得到 $\pi^2/8$。

\end{solve}

\begin{remark}[方法与条件]本题与\hyperref[pr:ch05-14th-A-makeup:6]{对数积分反演恒等式}中的函数 $\psi$ 满足 $\varphi(t)=\psi(t)/2$。用实变量定义
\[
\operatorname{Li}_2(z)=-\int_0^z\frac{\ln(1-u)}u\,\mathrm du\qquad(z\leq1),
\]
便可统一写成 $\varphi(t)=\tfrac12[\pi^2/6-\operatorname{Li}_2(1-t)]$。幂级数 $\sum z^k/k^2$ 只在其收敛范围内用于表示该函数，不能直接用于任意负实数。\end{remark}

\begin{example}[来源序号5]\label{ex:ch02-16th-A:5}\label{pr:ch02-16th-A:5}

讨论级数 $\sum_{n=2}^{\infty} \frac{(-1)^{n}}{n+(-1)^{[\sqrt{n}]}}$，$\sum_{n=2}^{\infty} \frac{(-1)^{n}}{\sqrt{n}+(-1)^{[\sqrt{n}]}}$ 的收敛性，其中 $[x]$ 表示 $x$ 的整数部分。

\par\href{https://github.com/lzxcr/cmc-mathematics/blob/64916fa1fdbc3e444a652eca204d5669058ee321/src/problems/preliminary/16/05.md}{题目与来源条目}

\end{example}

\begin{solve}\label{ans:ch02-16th-A:5}

分母扰动并不单调，不能直接用交错级数判别法。先展开成标准交错主项与余项；只有尚不绝对可和的部分才需要分析其振荡。记 $\varepsilon_n=(-1)^{\lfloor\sqrt n\rfloor}$。第一项有展开
\[
\frac{(-1)^n}{n+\varepsilon_n}
=\frac{(-1)^n}{n}+O(n^{-2}),
\]
故第一个级数收敛。

第二项有展开
\[
\frac{(-1)^n}{\sqrt n+\varepsilon_n}
=\frac{(-1)^n}{\sqrt n}
-\frac{(-1)^n\varepsilon_n}{n}+O(n^{-3/2}).
\]
只需考察中间项。令 $S_N=\sum_{n=2}^N(-1)^n\varepsilon_n$。在每个平方数区间 $[j^2,(j+1)^2)$ 上，$\varepsilon_n$ 不变，交错和的绝对值至多为 $1$；因此 $S_N=O(\sqrt N)$。分部求和给出
\[
\sum_{n=2}^N\frac{(-1)^n\varepsilon_n}{n}
=\frac{S_N}{N}+\sum_{n=2}^{N-1}\frac{S_n}{n(n+1)},
\]
右端有极限，因为 $S_N/N\to0$ 且后一项绝对收敛。故第二个级数也收敛。

两者绝对值的通项分别与 $1/n$、$1/\sqrt n$ 等价，绝对值级数均发散。因此两个级数均为条件收敛。

第二个展开的余项可用恒等式 $1/(1+u)=1-u+u^2/(1+u)$ 验证，取 $u=\varepsilon_n/\sqrt n$ 后余项为 $O(n^{-3/2})$，常数对充分大的 $n$ 统一成立。平方数区间虽越来越长，但内部仍是符号交替，每个完整区间和末端不完整区间的贡献都至多为 $1$，区间个数只有 $O(\sqrt N)$。于是 $|S_n|/[n(n+1)]=O(n^{-3/2})$ 可求和，这补足了单纯交错判别无法处理的扰动项。

\end{solve}

\begin{example}[来源序号6]\label{ex:ch02-16th-A:6}\label{pr:ch02-16th-A:6}

（1）设 $f_{1}(t)=\frac{t+3}{2}, f_{2}(t)=\frac{t+6}{3}, \{n_{k}\}$ 为取值于 $\{1,2\}$ 的整数列。令 $F_{1}(t)=f_{n_{1}}(t), F_{k+1}(t)=F_{k}[f_{n_{k+1}}(t)] (k \geq 1)$。证明：对任何 $x \in \mathbb{R}$，极限 $\lim_{k \to +\infty} F_{k}(x)$ 存在且与 $x$ 无关。

（2）若题 （1）中的 $f_{1}, f_{2}$ 改为 $f_{1}(t)=t-\arctan t, f_{2}(t)=2 \arctan t-t$，结论如何？

\par\href{https://github.com/lzxcr/cmc-mathematics/blob/64916fa1fdbc3e444a652eca204d5669058ee321/src/problems/preliminary/16/06.md}{题目与来源条目}

\end{example}

\begin{solve}\label{ans:ch02-16th-A:6}

（1）两个仿射映射有共同不动点 $3$，且
\[
|f_i(t)-3|\leq\tfrac12|t-3|\qquad(i=1,2).
\]
任意次序的 $k$ 重复合都满足 $|F_k(x)-3|\leq2^{-k}|x-3|$，故极限为 $3$。

（2）结论仍成立，极限为 $0$。两映射均为奇函数，且对 $t\neq0$ 有
\[
|t-\arctan t|<|t|,\qquad |2\arctan t-t|<|t|.
\]
当 $x=0$ 时结论显然。对于 $x\neq0$，复合次序随 $k$ 改变，不宜把 $F_k(x)$ 当作固定方向的单一迭代；应给出对任意复合次序都有效的估计。固定 $x$ 和 $\varepsilon\in(0,|x|)$。连续性保证存在 $q<1$，使得
\[
|f_i(t)|\leq q|t|\qquad
(\varepsilon\leq|t|\leq|x|,\ i=1,2).
\]
从内向外计算任意 $k$ 重复合：若中间值曾进入 $[-\varepsilon,\varepsilon]$，以后不会离开；否则每一步都按比例 $q$ 缩小。因此
\[
|F_k(x)|\leq\max\{\varepsilon,q^k|x|\}.
\]
先令 $k\to\infty$，再令 $\varepsilon\downarrow0$，即得结论。

对 $t>0$，$0<\arctan t<t$，所以 $0<t-\arctan t<t$，且 $-t<2\arctan t-t<t$；由奇性可推广到负数。固定环带 $\varepsilon\le|t|\le|x|$ 后，两个连续比值 $|f_i(t)|/|t|$ 在紧集上取得最大值，且每点都严格小于 $1$，所以其共同最大值 $q$ 也严格小于 $1$。这种统一性是任意复合次序下仍能收敛的关键，单点的严格缩小本身还不够。

\end{solve}

\begin{remark}[方法与条件]第二问只需有限族连续映射具有共同不动点，且严格减小到该点的距离；在每个避开不动点的紧环带上，便有共同的收缩比例。它也提醒我们区分收敛与收敛速度：$t-\arctan t=O(t^3)$ 在原点附近收敛很快，而 $2\arctan t-t=t-\tfrac23t^3+O(t^5)$ 的导数在原点等于 $1$。\end{remark}

\section{初赛：第 17 届（2025 年）}

本节选入3道课程相关题；来源序号保留原题解编排，未选入题目不重新补号。

\begin{remark}[来源说明]考试形式：闭卷；考试时间：150 分钟；满分：100 分。\end{remark}

\begin{example}[来源序号2]\label{ex:ch01-17th-A:2}\label{pr:ch01-17th-A:2}

设 $N \geq 1$，$S$ 是包含 $N$ 个整数的集合，满足如下的加性唯一性条件：若 $n_k \in S (1 \leq k \leq 4)$ 且 $n_1 + n_2 = n_3 + n_4$，则必有 $n_1 = n_3$ 或 $n_1 = n_4$。令 $f(x) := \sum_{n \in S} e^{2\pi i nx}$。

（1）计算 $\int_0^1 |f(x)|^2 \mathrm{d}x$，$\int_0^1 |f(x)|^4 \mathrm{d}x$。

（2）证明：$\int_0^1 |f(x)| \mathrm{d}x \geq \frac{\sqrt{N}}{2}$。

\par\href{https://github.com/lzxcr/cmc-mathematics/blob/64916fa1fdbc3e444a652eca204d5669058ee321/src/problems/preliminary/17/02.md}{题目与来源条目}

\end{example}

\begin{solve}\label{ans:ch01-17th-A:2}

（1）利用整数频率的正交性，积分只保留零频率项：
\[
\int_0^1|f|^2\,\mathrm dx=N,\qquad
\int_0^1|f|^4\,\mathrm dx
=\#\{(a,b,c,d)\in S^4:a+b=c+d\}.
\]
题设保证后一集合只有 $(c,d)=(a,b)$ 或 $(b,a)$ 两类；两类各有 $N^2$ 个元素，重合部分为 $a=b=c=d$ 的 $N$ 个元素。因此
\[
\int_0^1|f|^4\,\mathrm dx=2N^2-N.
\]

（2）二次矩已经固定，四次矩又受到加性唯一性控制；若一次矩过小，这两条信息就会冲突。把二次矩插在一次与四次矩之间，Hölder 不等式给出
\[
N=\int_0^1 |f|^{2/3}|f|^{4/3}\,\mathrm dx
\leq\left(\int_0^1|f|\,\mathrm dx\right)^{2/3}
     \left(\int_0^1|f|^4\,\mathrm dx\right)^{1/3}.
\]
从而得到比题设更强的下界
\[
\int_0^1|f|\,\mathrm dx
\geq\frac{N^{3/2}}{\sqrt{2N^2-N}}
>\sqrt{\frac N2}>\frac{\sqrt N}{2}.
\]

正交性具体指：整数 $m\ne0$ 时 $\int_0^1e^{2\pi imx}\,dx=0$，$m=0$ 时积分为 $1$。将 $|f|^4$ 展开为四重有限和后，每项频率为 $a+b-c-d$，因此积分变为有序四元组计数，交换次序不涉及收敛问题。第二问中 Hölder 的两个指数取 $3/2$ 与 $3$，相应幂次恰为 $|f|$ 与 $|f|^4$；不等式两边取 $3/2$ 次方再除以正的四次矩平方根，即得一次矩下界。

\end{solve}

\begin{remark}[方法与条件]满足题设的集合称为 Sidon 集，其中“和的表示唯一”是忽略两加数次序而言。一般有限整数集的加性能量
\[
E(S)=\#\{a+b=c+d:a,b,c,d\in S\}
\]
总满足 $E(S)\geq2N^2-N$；同一插值论证给出 $\|f\|_1\geq N^{3/2}/\sqrt{E(S)}$。因此，加法关系越少，指数和的 $L^1$ 范数越不易因抵消而变小。

二次矩和四次矩把频率的计数关系转成函数范数；相关正交性与加性能量的系统解释，见\hyperref[app:integration-fourier]{方法卷：Fourier 正交与能量}。\end{remark}

\begin{example}[来源序号5]\label{ex:ch01-17th-A:5}\label{pr:ch01-17th-A:5}

设 $f \in C[0,1] \cap C^2(0,1)$，$\sup_{x \in (0,1)} f''(x) = 2$。证明：存在唯一的二次多项式 $P(x) = x^2 + bx + c$，使得 $P(0) - f(0) = P(1) - f(1) = 0$，且下列结论之一成立：

（1）$f(x) > P(x) (\forall x \in (0,1))$；

（2）$f(x) = P(x) (\forall x \in [0,1])$。

\par\href{https://github.com/lzxcr/cmc-mathematics/blob/64916fa1fdbc3e444a652eca204d5669058ee321/src/problems/preliminary/17/05.md}{题目与来源条目}

\end{example}

\begin{solve}\label{ans:ch01-17th-A:5}

端点条件唯一确定
\[
P(x)=x^2+\bigl(f(1)-f(0)-1\bigr)x+f(0).
\]
减去端点插值多项式后，边界值变成零，二阶导数的上界变成凹性条件。令 $h=f-P$，则 $h(0)=h(1)=0$ 且 $h''=f''-2\leq0$，所以 $h$ 是凹函数，$h\geq0$。

若 $h$ 不恒为零，存在 $\xi\in(0,1)$ 使 $h(\xi)>0$。将 $x\in(0,\xi)$ 写成 $0,\xi$ 的凸组合，将 $x\in(\xi,1)$ 写成 $\xi,1$ 的凸组合，凹性分别给出
\[
h(x)\geq\frac{x}{\xi}h(\xi)>0,
\qquad
h(x)\geq\frac{1-x}{1-\xi}h(\xi)>0.
\]
故 $f>P$ 在整个开区间成立；否则 $f=P$ 在闭区间恒成立。

上确界等于 $2$ 只说明每点 $f''\le2$，不要求某点真的取到 $2$；证明也只用到了这条上界。由 $h''\le0$ 可知 $h'$ 单调不增，进而得到开区间内的凹性；利用端点连续性将凹性延伸到闭区间。端点连线为零，所以 $h(x)\ge(1-x)h(0)+xh(1)=0$。一旦某处严格为正，连接该点与左右端点的两条弦就给出全开区间的严格正下界，因此不会出现“有些内点相等、有些内点严格大于”的第三种情况。

\end{solve}

\begin{remark}[方法与条件]证明只需要 $f''\leq2$，无需上确界恰为 $2$。更一般地，若 $f''\leq M$，令 $P$ 是端点与 $f$ 相同且二次项系数为 $M/2$ 的多项式，同样有 $f\geq P$，并具有“恒等或内部严格大于”的二择一性质。这是一维强最大值原理的简单形式。\end{remark}

\begin{example}[来源序号6]\label{ex:ch01-17th-A:6}\label{pr:ch01-17th-A:6}

设 $\mathcal{A}$ 为正整数集 $\mathbb{Z}_+$ 的子集，已知 $\lim_{n \rightarrow +\infty} \frac{S_n(0,1)}{n} = p$ 存在，其中对于 $0 \leq a \leq b \leq 1$，$S_n(a,b)$ 表示集合 $\{k \in \mathcal{A} \mid an \leq k \leq bn\}$ 的元素个数。

（1）若 $0 < a < b < 1$，证明极限 $\lim_{n \rightarrow +\infty} \frac{S_n(a,b)}{n}$ 存在并求其值。

（2）若 $f$ 在 $[0,1]$ 上 Riemann 可积，证明：$\lim_{n \rightarrow +\infty} \frac{1}{n} \sum_{k \leq n, k \in \mathcal{A}} f\left(\frac{k}{n}\right) = p \int_0^1 f(x) \mathrm{d}x$。

\par\href{https://github.com/lzxcr/cmc-mathematics/blob/64916fa1fdbc3e444a652eca204d5669058ee321/src/problems/preliminary/17/06.md}{题目与来源条目}

\end{example}

\begin{solve}\label{ans:ch01-17th-A:6}

（1）记 $N(t)=\#\{k\in\mathcal A:k\leq t\}$。题设即 $N(t)=pt+o(t)$，而端点至多影响一个计数，因此
\[
S_n(a,b)=N(bn)-N(an)+O(1)=p(b-a)n+o(n).
\]
所求极限为 $p(b-a)$。同式对 $0\leq a\leq b\leq1$ 也成立。

（2）第一问给出了每个区间内的渐近点密度；阶梯函数只是有限个区间指示函数的线性组合，因此先在这一类上验证，再用上下逼近即可。令 $L_n(f)=n^{-1}\sum_{k\leq n,\,k\in\mathcal A}f(k/n)$。由 （1）及线性性，对阶梯函数 $s$ 有 $L_n(s)\to p\int_0^1s$；有限个端点的取值不影响极限。

若 $f$ Riemann 可积，任给 $\varepsilon>0$，可取阶梯函数 $g\leq f\leq h$，使 $\int_0^1(h-g)<\varepsilon$。由 $L_n$ 的正性，
\[
p\int_0^1g\leq\liminf_nL_n(f)
\leq\limsup_nL_n(f)\leq p\int_0^1h.
\]
令 $\varepsilon\downarrow0$，即得 $L_n(f)\to p\int_0^1f$。

整数下标的密度极限可以推广到实数 $t$，因为 $N(t)=N(\lfloor t\rfloor)$ 且 $\lfloor t\rfloor/t\to1$。区间计数在左端是否包含整数只带来至多 $1$ 的误差，除以 $n$ 后消失。第二问中 $0\le p\le1$，而 $L_n$ 的正性指 $u\le v$ 蕴含 $L_n(u)\le L_n(v)$，不要求 $f$ 本身非负。上下极限的间距至多为 $p\int(h-g)\le\varepsilon$，同时目标积分也夹在相同两界之间，故夹逼得到所述极限。

\end{solve}

\begin{remark}[方法与条件]令 $\mu_n=n^{-1}\sum_{k\leq n,\,k\in\mathcal A}\delta_{k/n}$，则结论包含有限测度的弱收敛 $\mu_n\Rightarrow p\,\mathrm dx$。Riemann 可积函数的间断点集为零测集，上下阶梯函数的积分逼近把连续测试函数的结论延伸到了这里。

不能任意换成有界可测函数：取 $\mathcal A=\mathbb Z_+$、$f=\mathbf1_{\mathbb Q\cap[0,1]}$，所有离散平均都等于 $1$，而 Lebesgue 积分等于 $0$。

区间计数如何延伸为测度收敛，以及为何仍须控制间断点集，见\hyperref[app:integration-riemann]{方法卷：Riemann 和与测度}。\end{remark}

\section{决赛：第 1 届（2010 年）}

本节选入7道课程相关题；来源序号保留原题解编排，未选入题目不重新补号。

\begin{remark}[来源说明]本届数学类决赛不分高、低年级组。\end{remark}

\begin{remark}[来源说明]原卷第一大题的四个填空小题分别列为第1至4题；原卷第二至八大题依次列为第5至11题。\end{remark}

\begin{example}[来源序号1]\label{ex:fin01-all:1}\label{pr:fin01-all:1}

设 $\beta>\alpha>0$，求积分
\[
\int_0^\infty\frac{e^{-\alpha x^2}-e^{-\beta x^2}}{x^2}\,dx.
\]

\par\href{https://github.com/lzxcr/cmc-mathematics/blob/64916fa1fdbc3e444a652eca204d5669058ee321/src/problems/finals/01/01.md}{题目与来源条目}

\end{example}

\begin{solve}\label{ans:fin01-all:1}

把差商写成参数积分，再用 Tonelli 定理交换非负积分：
\[
\begin{aligned}
\int_0^\infty\frac{e^{-\alpha x^2}-e^{-\beta x^2}}{x^2}\,dx
&=\int_\alpha^\beta\int_0^\infty e^{-t x^2}\,dx\,dt\\
&=\frac{\sqrt\pi}{2}\int_\alpha^\beta t^{-1/2}\,dt
=\sqrt\pi(\sqrt\beta-\sqrt\alpha).
\end{aligned}
\]

参数恒等式来自 $\partial_t e^{-tx^2}=-x^2e^{-tx^2}$，即 $e^{-\alpha x^2}-e^{-\beta x^2}=x^2\int_\alpha^\beta e^{-tx^2}\,dt$。因为 $\beta>\alpha>0$，二重积分的被积函数非负，交换积分不需事先知道其有限；换序后的计算给出有限值，也就证明原反常积分收敛。内积分用 $u=\sqrt t\,x$ 化为 $t^{-1/2}\int_0^\infty e^{-u^2}\,du=\sqrt\pi/(2\sqrt t)$。

\end{solve}

\begin{remark}[方法与条件]参数积分同时消除了原点处的表面奇性。交换积分的依据是非负性，见\hyperref[app:integration-exchange]{方法卷：极限交换}。\end{remark}

\begin{example}[来源序号2]\label{ex:fin01-all:2}\label{pr:fin01-all:2}

设 $k>0$。若方程 $kx+x^{-2}=1$ 在 $(0,\infty)$ 中恰有一个实根，求 $k$。

\par\href{https://github.com/lzxcr/cmc-mathematics/blob/64916fa1fdbc3e444a652eca204d5669058ee321/src/problems/finals/01/02.md}{题目与来源条目}

\end{example}

\begin{solve}\label{ans:fin01-all:2}

方程等价于 $k=h(x):=x^{-1}-x^{-3}$，且正根必有 $x>1$。函数 $h$ 在 $(1,\sqrt3)$ 上严格递增，在 $(\sqrt3,\infty)$ 上严格递减，两端极限均为 $0$。故恰有一根当且仅当水平线经过其最高点，即
\[
k=h(\sqrt3)=\frac{2}{3\sqrt3}.
\]

由 $kx=1-x^{-2}>0$ 可先排除 $0<x\le1$。求导得 $h'(x)=(3-x^2)/x^4$，所以 $\sqrt3$ 是唯一极大值点。当 $0<k<h(\sqrt3)$ 时，连续性及两段严格单调性保证左右各有一根；当 $k>h(\sqrt3)$ 时没有根；等于最大值时只有极大值点这一根。因此“恰有一个根”确定的是相切情形，而不能随意解某个代数方程后挑选根。

\end{solve}

\begin{example}[来源序号3]\label{ex:fin01-all:3}\label{pr:fin01-all:3}

设 $f\in C[a,b]$，右导数 $f'_+(a)$ 存在且非零。对每个 $x\in(a,b)$，取 $\xi_x\in[a,x]$ 使
\[
\int_a^x f(t)\,dt=(x-a)f(\xi_x).
\]
求 $\displaystyle\lim_{x\downarrow a}\frac{\xi_x-a}{x-a}$。

\par\href{https://github.com/lzxcr/cmc-mathematics/blob/64916fa1fdbc3e444a652eca204d5669058ee321/src/problems/finals/01/03.md}{题目与来源条目}

\end{example}

\begin{solve}\label{ans:fin01-all:3}

记 $c=f'_+(a)\ne0$、$h=x-a$。一阶展开给出
\[
\frac1h\int_a^{a+h}f(t)\,dt=f(a)+\frac c2h+o(h),
\qquad
f(\xi_x)=f(a)+c(\xi_x-a)+o(h).
\]
第二式利用了 $0\le\xi_x-a\le h$，故余项对中值点的选择一致。比较两式即得所求极限为 $1/2$。

右可导意味着 $f(a+t)=f(a)+ct+t\eta(t)$，其中 $\eta(t)\to0$。积分余项满足 $|h^{-1}\int_0^h t\eta(t)\,dt|\le(h/2)\sup_{0<t\le h}|\eta(t)|=o(h)$；在中值点处，余项也至多为 $h\sup_{0<t\le h}|\eta(t)|$。相等后除以非零的 $ch$，得到 $(\xi_x-a)/h=1/2+o(1)$。这说明即使中值点不唯一，任何允许的选择都有相同极限。

\end{solve}

\begin{example}[来源序号5]\label{ex:fin01-all:5}\label{pr:fin01-all:5}

设 $f:(-1,1)\to\mathbb R$ 在 $0$ 处可导，且 $f(0)=0$。证明
\[
\lim_{n\to\infty}\sum_{k=1}^n f\!\left(\frac{k}{n^2}\right)=\frac{f'(0)}2.
\]

\par\href{https://github.com/lzxcr/cmc-mathematics/blob/64916fa1fdbc3e444a652eca204d5669058ee321/src/problems/finals/01/05.md}{题目与来源条目}

\end{example}

\begin{solve}\label{ans:fin01-all:5}

写 $f(t)=f'(0)t+t\varepsilon(t)$，其中 $\varepsilon(t)\to0$。由于所有采样点都落在 $(0,1/n]$，
\[
\left|\sum_{k=1}^n f\!\left(\frac{k}{n^2}\right)
-f'(0)\frac{n+1}{2n}\right|
\le\frac{n+1}{2n}\sup_{0<t\le1/n}|\varepsilon(t)|\longrightarrow0.
\]

这里对 $n\ge2$ 考察和式即可，所有自变量都在定义域 $(-1,1)$ 内。线性主项的和为 $f'(0)n^{-2}\sum_{k=1}^nk=f'(0)(n+1)/(2n)$。余项若仅说“每项趋零”，不足以应对项数也趋于无穷的情形；上述估计利用全部采样点同时缩到零点附近，使余项系数的上确界趋零，而总权重 $(n+1)/(2n)$ 始终有界。

\end{solve}

\begin{example}[来源序号6]\label{ex:fin01-all:6}\label{pr:fin01-all:6}

设 $f:[0,\infty)\to\mathbb R$ 一致连续，且对每个固定的 $x\ge0$，均有 $f(x+n)\to0$。证明函数列 $f_n(x)=f(x+n)$ 在 $[0,1]$ 上一致收敛于 $0$。

\par\href{https://github.com/lzxcr/cmc-mathematics/blob/64916fa1fdbc3e444a652eca204d5669058ee321/src/problems/finals/01/06.md}{题目与来源条目}

\end{example}

\begin{solve}\label{ans:fin01-all:6}

一致连续性使 $\{f_n\}$ 等度连续。给定 $\varepsilon>0$，可取 $[0,1]$ 的有限网 $x_1,\ldots,x_m$，使每个 $x$ 均有某个 $x_j$ 满足
\[
|f_n(x)-f_n(x_j)|<\varepsilon/2\qquad(\text{对所有 }n).
\]
有限多个 $f_n(x_j)$ 同时趋于 $0$，故充分大的 $n$ 满足 $\|f_n\|_\infty<\varepsilon$。

一致连续性给出同一个 $\delta>0$，使任意 $s,t\ge0$ 且 $|s-t|<\delta$ 时都有 $|f(s)-f(t)|<\varepsilon/2$。把 $[0,1]$ 分成长度小于 $\delta$ 的有限网格，取其端点为 $x_j$；平移后距离仍是 $|x-x_j|$，所以估计对全部 $n$ 成立。对每个固定 $x_j$ 用题设找到 $N_j$，再取 $N=\max_jN_j$，便有 $|f(x+n)|\le|f(x+n)-f(x_j+n)|+|f(x_j+n)|<\varepsilon$，其中 $N$ 与 $x$ 无关。

\end{solve}

\begin{remark}[方法与条件]这是一条一般的紧致性原则：紧空间上，等度连续函数列若逐点收敛到连续函数，便一致收敛。这里还推出更强的结论 $f(t)\to0$（$t\to\infty$），因为可写 $t=n+x$，$x\in[0,1)$。\end{remark}

\begin{example}[来源序号7]\label{ex:fin01-all:7}\label{pr:fin01-all:7}

设 $D=\{(x,y):x^2+y^2<1\}$。函数 $f,g$ 在 $D$ 上连续且具有二阶偏导数，$g$ 有界，并满足
\[
f(x,y)\longrightarrow+\infty\quad(x^2+y^2\to1),
\qquad \Delta f=e^f,\quad \Delta g\ge e^g.
\]
证明 $f\ge g$ 在 $D$ 上处处成立。

\par\href{https://github.com/lzxcr/cmc-mathematics/blob/64916fa1fdbc3e444a652eca204d5669058ee321/src/problems/finals/01/07.md}{题目与来源条目}

\end{example}

\begin{solve}\label{ans:fin01-all:7}

令 $u=g-f$。边界附近有 $u<0$；若某处 $u>0$，则 $u$ 在内部一点 $p$ 达到正的最大值。因此
\[
0\ge\Delta u(p)\ge e^{g(p)}-e^{f(p)}>0,
\]
矛盾。

最大值的存在可在紧圆盘上严格实现。设 $g\le M$，由边界爆破条件，可选 $r<1$ 使 $f>M+1$ 在 $r<\sqrt{x^2+y^2}<1$ 上成立，故那里 $u<-1$。若另有一点 $u>0$，它必在闭圆盘 $\overline D_r$ 内；$u$ 连续，所以其正的最大值在此紧集取得，且取得点不在边界。沿两条坐标方向使用二阶导数的极大值必要条件，得 $u_{xx},u_{yy}\le0$；另一方面指数函数严格递增，$g(p)>f(p)$ 迫使拉普拉斯差为正，形成矛盾。

\end{solve}

\begin{remark}[方法与条件]这是半线性方程的比较原理：在极值点比较二阶导数，再利用非线性项的单调性。将 $e^t$ 换成任意严格递增函数，证明仍然成立。\end{remark}

\begin{example}[来源序号8]\label{ex:fin01-all:8}\label{pr:fin01-all:8}

设 $R_\varepsilon=[0,1-\varepsilon]^2$，并定义
\[
I=\lim_{\varepsilon\downarrow0}\iint_{R_\varepsilon}\frac{dx\,dy}{1-xy}.
\]

（1）证明 $I=\sum_{n=1}^\infty n^{-2}$。
（2）利用变换 $u=(x+y)/2$、$v=(y-x)/2$ 求 $I$，从而计算上述级数。

\par\href{https://github.com/lzxcr/cmc-mathematics/blob/64916fa1fdbc3e444a652eca204d5669058ee321/src/problems/finals/01/08.md}{题目与来源条目}

\end{example}

\begin{solve}\label{ans:fin01-all:8}

非负项级数可以逐项积分，由单调收敛定理，
\[
I=\iint_{[0,1]^2}\sum_{m=0}^\infty x^my^m\,dx\,dy
=\sum_{m=0}^\infty\frac1{(m+1)^2}.
\]
在指定变换下，$dx\,dy=2\,du\,dv$，积分域为
$0\le u\le1$、$|v|\le\min(u,1-u)$。利用关于 $v$ 的对称性并先对 $v$ 积分，得
\[
I=4\int_0^1\frac1{\sqrt{1-u^2}}
\arctan\frac{\min(u,1-u)}{\sqrt{1-u^2}}\,du.
\]
再令 $u=\sin\theta$。当 $0\le\theta\le\pi/6$ 时反正切为 $\theta$；当 $\pi/6\le\theta\le\pi/2$ 时反正切为 $\pi/4-\theta/2$。于是
\[
I=4\left[\int_0^{\pi/6}\theta\,d\theta
+\int_{\pi/6}^{\pi/2}\left(\frac\pi4-\frac\theta2\right)d\theta\right]
=\frac{\pi^2}{6}.
\]

逆变换为 $x=u-v,y=u+v$，故 $xy=u^2-v^2$、Jacobian 绝对值为 $2$。正方形的四条边给出 $|v|\le u$ 和 $|v|\le1-u$，两种上界在 $u=1/2$ 交换，对应 $\theta=\pi/6$。后一段使用半角恒等式 $(1-\sin\theta)/\cos\theta=\tan(\pi/4-\theta/2)$，其角度落在反正切主值范围内。整个积分最初可按非负广义积分处理，最终计算有限，因而同时证明了二重积分和级数都收敛。

\end{solve}

\begin{remark}[方法与条件]级数先变成非负二重积分，几何换元再把积分化为一次函数的面积。这里无须预先知道级数和，也无须在奇点附近单独估计。\end{remark}

\section{决赛：第 2 届（2011 年）}

本节选入3道课程相关题；来源序号保留原题解编排，未选入题目不重新补号。

\begin{remark}[来源说明]本届数学类决赛不分高、低年级组。\end{remark}

\begin{example}[来源序号2]\label{ex:fin02-all:2}\label{pr:fin02-all:2}

设 $0<f(x)<1$（$x\ge0$），且积分 $\int_0^\infty f(x)\,dx$ 与 $\int_0^\infty xf(x)\,dx$ 均收敛。证明
\[
\int_0^\infty xf(x)\,dx>
\frac12\left(\int_0^\infty f(x)\,dx\right)^2.
\]

\par\href{https://github.com/lzxcr/cmc-mathematics/blob/64916fa1fdbc3e444a652eca204d5669058ee321/src/problems/finals/02/02.md}{题目与来源条目}

\end{example}

\begin{solve}\label{ans:fin02-all:2}

记 $m=\int_0^\infty f>0$，与等质量的函数 $g=\boldsymbol1_{[0,m]}$ 比较。由于 $\int(f-g)=0$，
\[
\int_0^\infty xf(x)\,dx-\frac{m^2}{2}
=\int_0^\infty(x-m)(f(x)-g(x))\,dx>0.
\]
最后的被积函数非负，并在 $(0,m)$ 上严格为正。

函数 $g$ 把总质量 $m$ 全部放到最靠近原点、密度为 $1$ 的区间内，其一阶矩为 $\int_0^m x\,dx=m^2/2$。在 $0<x<m$ 上，$x-m<0$ 且 $f-1<0$，乘积严格为正；在 $x>m$ 上，$x-m>0$ 且 $f>0$，乘积也严格为正。被积函数可积，且在一个正长度区间上处处为正，所以积分严格为正。这解释了为何题设的严格密度界带来严格不等式，而不仅是非严格下界。

\end{solve}

\begin{remark}[方法与条件]在 $0\le f\le1$、总质量固定时，将质量尽量向左放置，使递增权重的积分最小。这是重排与“浴缸原理”的最基本形式；允许端点取值后，等号仅由 $\boldsymbol1_{[0,m]}$（几乎处处意义下）取得。\end{remark}

\begin{example}[来源序号3]\label{ex:fin02-all:3}\label{pr:fin02-all:3}

设级数 $\sum_{n=1}^\infty na_n$ 收敛。证明
\[
t_n=\sum_{k=1}^\infty ka_{n+k}
\]
有意义，且 $t_n\to0$。

\par\href{https://github.com/lzxcr/cmc-mathematics/blob/64916fa1fdbc3e444a652eca204d5669058ee321/src/problems/finals/02/03.md}{题目与来源条目}

\end{example}

\begin{solve}\label{ans:fin02-all:3}

令 $R_j=\sum_{m=j}^\infty ma_m$，则 $R_j\to0$。固定 $n$，记 $w_k=k/(n+k)$。分部求和并令上限趋于无穷，得
\[
t_n=w_1R_{n+1}+\sum_{k=2}^\infty(w_k-w_{k-1})R_{n+k}.
\]
这里 $w_k\uparrow1$，右端绝对收敛，系数均非负且总和为 $1$；有限分部求和中的边界项 $w_mR_{n+m+1}$ 趋于 $0$。因此
\[
|t_n|\le\sup_{j\ge n+1}|R_j|\longrightarrow0.
\]

把 $ka_{n+k}$ 写成 $w_k(R_{n+k}-R_{n+k+1})$，有限求和得到
\[
\sum_{k=1}^mka_{n+k}=w_1R_{n+1}+\sum_{k=2}^m(w_k-w_{k-1})R_{n+k}-w_mR_{n+m+1}.
\]
其中 $w_k-w_{k-1}=n/[(n+k)(n+k-1)]\ge0$，加上首个系数后总和趋于 $1$。$R_j$ 有界保证右侧级数绝对收敛，$R_j\to0$ 保证边界项消失。这先证明了每个 $t_n$ 有意义，再用同一表示证明 $t_n\to0$，没有把尚未存在的无穷和直接换序。

\end{solve}

\begin{remark}[方法与条件]分部求和把看似增长的权重变成了原级数尾和的凸组合；只需收敛，无须绝对收敛。\end{remark}

\begin{example}[来源序号5]\label{ex:fin02-all:5}\label{pr:fin02-all:5}

设连续函数 $f:\mathbb R\to\mathbb R$ 满足
\[
M:=\sup_{x,y\in\mathbb R}|f(x+y)-f(x)-f(y)|<\infty.
\]
证明存在 $a\in\mathbb R$，使 $\sup_x|f(x)-ax|<\infty$。

\par\href{https://github.com/lzxcr/cmc-mathematics/blob/64916fa1fdbc3e444a652eca204d5669058ee321/src/problems/finals/02/05.md}{题目与来源条目}

\end{example}

\begin{solve}\label{ans:fin02-all:5}

令 $g_m(x)=2^{-m}f(2^mx)$。由近似可加性，
\[
|g_{m+1}(x)-g_m(x)|\le\frac{M}{2^{m+1}}.
\]
故 $g_m$ 在整个实轴上一致收敛于连续函数 $g$，且 $|g-f|\le M$。又有
\[
|g_m(x+y)-g_m(x)-g_m(y)|\le2^{-m}M,
\]
取极限得 $g(x+y)=g(x)+g(y)$。连续的可加函数必有 $g(x)=g(1)x$，取 $a=g(1)$ 即可，且得到准确估计
\[
\sup_x|f(x)-ax|\le M.
\]

一致收敛来自尾差估计 $|g_r(x)-g_m(x)|\le\sum_{j=m}^{r-1}M2^{-j-1}\le M2^{-m}$，它对全部实数 $x$ 同时成立。由可加性有 $g(0)=0$、$g(-x)=-g(x)$，并通过整数倍与等分得到 $g(q)=qg(1)$ 对每个有理数 $q$ 成立。对任意实数 $x$ 取有理数列 $q_j\to x$，连续性再给出 $g(x)=xg(1)$，这补足了“连续可加必为线性”的论证。

\end{solve}

\begin{remark}[方法与条件]这称为可加方程的 Hyers 稳定性：一致有界的方程误差，对应一致有界的函数误差。缩放使误差形成可求和的几何级数；所得斜率也唯一，因为两条不同斜率的直线不可能相差有界。\end{remark}

\section{决赛：第 3 届（2012 年）}

本节选入4道课程相关题；来源序号保留原题解编排，未选入题目不重新补号。

\begin{remark}[来源说明]本届数学类决赛不分高、低年级组。\end{remark}

\begin{example}[来源序号2]\label{ex:fin03-all:2}\label{pr:fin03-all:2}

设 $f$ 在 $[a,b]$ 上有二阶导数，且 $f''$ 黎曼可积。证明对每个 $x\in[a,b]$，
\[
f(x)=f(a)+f'(a)(x-a)+\int_a^x(x-t)f''(t)\,dt.
\]

\par\href{https://github.com/lzxcr/cmc-mathematics/blob/64916fa1fdbc3e444a652eca204d5669058ee321/src/problems/finals/03/02.md}{题目与来源条目}

\end{example}

\begin{solve}\label{ans:fin03-all:2}

对固定的 $x$ 分部积分：
\[
\begin{aligned}
\int_a^x(x-t)f''(t)\,dt
&=\bigl[(x-t)f'(t)\bigr]_a^x+\int_a^x f'(t)\,dt\\
&=-(x-a)f'(a)+f(x)-f(a).
\end{aligned}
\]
这些等式只需要导数可积的 Newton--Leibniz 公式，不要求 $f''$ 连续。

可将分部积分理解为对乘积 $q(t)=(x-t)f'(t)$ 使用 Newton--Leibniz 公式。其导数为 $q'(t)=-f'(t)+(x-t)f''(t)$；$f'$ 因可导而连续，$f''$ 按题设黎曼可积，所以 $q'$ 黎曼可积，公式适用。端点 $t=x$ 的乘积为零，$t=a$ 的乘积为 $(x-a)f'(a)$，产生前面的负号。移项后即为带积分余项的一阶 Taylor 公式，$x=a$ 时各项直接相等。

\end{solve}

\begin{remark}[方法与条件]这是 Taylor 公式的积分余项。核 $(x-t)$ 来自两次积分；更高阶情形的核为 $(x-t)^{m-1}/(m-1)!$。见\hyperref[app:estimate-taylor]{方法卷：Taylor 余项}。\end{remark}

\begin{example}[来源序号3]\label{ex:fin03-all:3}\label{pr:fin03-all:3}

设 $0<k_0<k_1<\cdots<k_n$ 为整数，$A_1,\ldots,A_n\in\mathbb R$。求
\[
f(x)=\sin(k_0x)+\sum_{j=1}^n A_j\sin(k_jx)
\]
在 $[0,2\pi)$ 上不同零点个数的最小值，并证明。

\par\href{https://github.com/lzxcr/cmc-mathematics/blob/64916fa1fdbc3e444a652eca204d5669058ee321/src/problems/finals/03/03.md}{题目与来源条目}

\end{example}

\begin{solve}\label{ans:fin03-all:3}

最小值为 $2k_0$，在所有 $A_j=0$ 时达到。下面证明下界。

非零三角多项式在一周期内只有有限个零点，且变号点的个数为偶数，记为 $2m$。若 $m<k_0$，设这些变号点为 $x_1,\ldots,x_{2m}$，并令
\[
p(x)=\prod_{j=1}^{2m}\sin\frac{x-x_j}{2}
\]
（$m=0$ 时令 $p=1$）。这是次数不超过 $m$ 的三角多项式。它与 $f$ 在同样的位置变号，故适当调整整体符号后，$f(x)p(x)\ge0$ 且不恒为零，从而 $\int_0^{2\pi}fp>0$。

另一方面，$f$ 的所有频率都至少为 $k_0>m$，三角函数的正交性给出 $\int_0^{2\pi}fp=0$，矛盾。因此 $f$ 至少有 $2k_0$ 个变号点，也就至少有这么多零点。

零点有限可由代数化说明：令 $z=e^{ix}$，用 $\sin(kx)=(z^k-z^{-k})/(2i)$ 表示，再乘以最高频率对应的 $z$ 的幂，得到非零普通多项式，因此单位圆上的根只有有限个。$p$ 的 $2m$ 个半角因子相乘后，各指数是 $2m$ 个 $\pm1/2$ 之和，故都是绝对值不超过 $m$ 的整数，确实是周期为 $2\pi$ 的低阶三角多项式。每经过一个变号点，$f$ 与 $p$ 同时变号，乘积符号保持不变；其他零点不改变符号。乘积连续且某处为正，所以其积分严格为正。

\end{solve}

\begin{remark}[方法与条件]这是 Sturm--Hurwitz 振荡定理的有限 Fourier 级数版本：最低非零频率决定至少发生多少次振荡。证明将“零点太少”转化为一个低频的同号试验函数，再用正交性排除。\end{remark}

\begin{example}[来源序号4]\label{ex:fin03-all:4}\label{pr:fin03-all:4}

设正数列 $(a_n)$ 满足
\[
\liminf_{n\to\infty}a_n=1,\qquad
\limsup_{n\to\infty}a_n<\infty,\qquad
\lim_{n\to\infty}(a_1\cdots a_n)^{1/n}=1.
\]
证明 $\displaystyle\lim_{n\to\infty}\frac1n\sum_{k=1}^n a_k=1$。

\par\href{https://github.com/lzxcr/cmc-mathematics/blob/64916fa1fdbc3e444a652eca204d5669058ee321/src/problems/finals/03/04.md}{题目与来源条目}

\end{example}

\begin{solve}\label{ans:fin03-all:4}

令 $b_n=\log a_n$。题设给出 $\liminf b_n=0$，且
\[
\frac1n\sum_{k=1}^n b_k\longrightarrow0.
\]
记 $b_n^-:=\max(-b_n,0)$，则 $b_n^-\to0$，因而其 Cesàro 均值也趋于 $0$。由 $|b_n|=b_n+2b_n^-$，得
\[
\frac1n\sum_{k=1}^n|b_k|\longrightarrow0.
\]
又因 $a_n$ 最终处于某个紧区间 $[1/2,M]$，在该区间上有
$|a_n-1|\le C|\log a_n|$。所以 $n^{-1}\sum_{k=1}^n|a_k-1|\to0$，结论随之成立。

由 $\liminf b_n=0$，给定 $\varepsilon>0$ 后，最终 $b_n\ge-\varepsilon$，即 $0\le b_n^-\le\varepsilon$，所以负部趋零。对数平均又趋零，迫使正部的平均也趋零，排除了正负大幅抵消的可能。最后对 $e^t$ 使用中值定理：当 $a_n\in[1/2,M]$ 时，$|e^{b_n}-1|\le\max(1,M)|b_n|$。有限个早期项除以 $n$ 后消失，故这个最终成立的估计足以控制全体算术平均。

\end{solve}

\begin{remark}[方法与条件]几何平均控制的是对数的平均。下极限条件排除了持续的负偏差，再由有界性把对数尺度上的控制转回原尺度；实际上证明了平均绝对偏差也趋于零。\end{remark}

\begin{example}[来源序号7]\label{ex:fin03-all:7}\label{pr:fin03-all:7}

设 $F,G:[0,\infty)\to[0,\infty)$ 单调递减，且
\[
x(F(x)+G(x))\longrightarrow0\qquad(x\to\infty).
\]

（1）证明对每个 $\varepsilon>0$，
\[
\lim_{x\to\infty}\int_\varepsilon^\infty xF(xt)\cos t\,dt=0.
\]
（2）若进一步有
\[
\lim_{n\to\infty}\int_0^\infty(F(t)-G(t))\cos(t/n)\,dt=0,
\]
证明
\[
\lim_{s\to0}\int_0^\infty(F(t)-G(t))\cos(st)\,dt=0,
\]
其中极限中的积分取 $s\ne0$。

\par\href{https://github.com/lzxcr/cmc-mathematics/blob/64916fa1fdbc3e444a652eca204d5669058ee321/src/problems/finals/03/07.md}{题目与来源条目}

\end{example}

\begin{solve}\label{ans:fin03-all:7}

对任意趋于零的非负递减函数 $h$，Dirichlet 估计给出
\[
\left|\int_R^\infty h(t)\cos(st)\,dt\right|
\le\frac{2h(R)}{|s|}\qquad(s\ne0).
\]
这也保证了本题各振荡积分的收敛。将 $h(t)=xF(xt)$、$s=1$、$R=\varepsilon$ 代入，第一问的积分绝对值至多为
$2xF(x\varepsilon)\to0$。

第二问记 $H=F-G$、$K=F+G$，以及
$J(s)=\int_0^\infty H(t)\cos(st)\,dt$。由偶性只须考虑 $s\downarrow0$。取 $n=\lfloor1/s\rfloor$，则
$|s-1/n|\le1/[n(n+1)]$。在 $[0,n]$ 上，用余弦函数的 Lipschitz 性，得
\[
\left|\int_0^nH(t)\bigl(\cos(st)-\cos(t/n)\bigr)\,dt\right|
\le\frac1{n(n+1)}\int_0^n tK(t)\,dt\longrightarrow0,
\]
因为 $tK(t)$ 有界。在 $[n,\infty)$ 上，上述 Dirichlet 估计将两条尾项之和控制为
\[
2\bigl(s^{-1}+n\bigr)K(n)\le(4n+2)K(n)\longrightarrow0.
\]
故 $J(s)-J(1/n)\to0$。结合题设 $J(1/n)\to0$ 即得结论。

选择 $n=\lfloor1/s\rfloor$ 意味着 $n\le1/s<n+1$，所以 $1/(n+1)<s\le1/n$，从而频率差至多为 $1/[n(n+1)]$。在近端，$|\cos u-\cos v|\le|u-v|$ 将频率差转化为一个小的乘子；在远端，分别对非负单调的 $F,G$ 应用 Dirichlet 估计，再相加，而不是误把 $F-G$ 当作单调函数。$tK(t)\to0$ 还保证它在全半轴有界，故近端误差至多为常数乘 $n/[n(n+1)]\to0$。

\end{solve}

\begin{remark}[方法与条件]这里不要求 $F-G$ 绝对可积。靠近原点的频率由离散序列补成连续参数，关键是把积分分成有限区间和振荡尾项：前者比较频率，后者利用单调性带来的抵消。\end{remark}

\section{决赛：第 4 届（2013 年）}

本节选入4道课程相关题；来源序号保留原题解编排，未选入题目不重新补号。

\begin{remark}[来源说明]本届数学类决赛不分高、低年级组。\end{remark}

\begin{example}[来源序号1]\label{ex:fin04-all:1}\label{pr:fin04-all:1}

设 $A>0$。直线 $\ell$ 与双曲线右支 $x^2-y^2=2$、$x>0$ 围成有限区域，其面积为 $A$。证明：

（1）所有这样的弦的中点轨迹为一条双曲线的右支。
（2）每条弦所在的直线都是该中点轨迹的切线。

\par\href{https://github.com/lzxcr/cmc-mathematics/blob/64916fa1fdbc3e444a652eca204d5669058ee321/src/problems/finals/04/01.md}{题目与来源条目}

\end{example}

\begin{solve}\label{ans:fin04-all:1}

将双曲线参数化为 $r(t)=\sqrt2(\cosh t,\sinh t)$，设两交点的参数为 $m-h,m+h$，其中 $h>0$。线性变换
\[
\begin{pmatrix}x\\y\end{pmatrix}\longmapsto
\begin{pmatrix}\cosh m&-\sinh m\\-\sinh m&\cosh m\end{pmatrix}
\begin{pmatrix}x\\y\end{pmatrix}
\]
保持双曲线与面积，把两端点送到 $r(-h),r(h)$。故所围面积只依赖 $h$，且
\[
A=2\int_0^{\sqrt2\sinh h}\left(\sqrt2\cosh h-\sqrt{2+y^2}\right)dy
=\sinh(2h)-2h.
\]
右式在 $h>0$ 上从 $0$ 严格增至无穷，所以 $A$ 唯一确定 $h$。

弦中点为
\[
M=\sqrt2\cosh h\,(\cosh m,\sinh m).
\]
因此其轨迹为 $x^2-y^2=2\cosh^2h$ 的右支。弦所在直线的方程为
\[
x\cosh m-y\sinh m=\sqrt2\cosh h,
\]
恰为此轨迹在 $M$ 处的切线。

双曲函数满足 $\cosh^2m-\sinh^2m=1$，所以所用矩阵行列式为 $1$，面积不变；加法公式则说明它把参数 $t$ 变成 $t-m$。面积函数的导数为 $2\cosh(2h)-2>0$，故固定正面积确实唯一确定端点参数差 $2h$，而中心参数 $m$ 可任取实数，因此所求轨迹是整个右支。对于 $x^2-y^2=R^2$，在 $(x_0,y_0)$ 处的切线为 $x_0x-y_0y=R^2$；代入 $R=\sqrt2\cosh h$ 和中点坐标后，正好得到弦的方程。

\end{solve}

\begin{remark}[方法与条件]双曲函数参数把“弦的位置”和“弦张开的程度”分离为 $m,h$；保持面积的双曲旋转消去了位置参数。固定面积因而固定 $h$，中点轨迹自然成为与原曲线同心、同渐近线的双曲线。\end{remark}

\begin{example}[来源序号2]\label{ex:fin04-all:2}\label{pr:fin04-all:2}

设 $f:[a,b]\to[a,b]$ 满足：存在 $0<L<1$，对任意 $x\ne y$ 有
$|f(x)-f(y)|<L|x-y|$。从任意 $x_1\in[a,b]$ 出发，定义
\[
x_{n+1}=\frac{x_n+f(x_n)}2.
\]
证明 $x_n$ 收敛，且其极限为 $f$ 的不动点。

\par\href{https://github.com/lzxcr/cmc-mathematics/blob/64916fa1fdbc3e444a652eca204d5669058ee321/src/problems/finals/04/02.md}{题目与来源条目}

\end{example}

\begin{solve}\label{ans:fin04-all:2}

映射 $T=(I+f)/2$ 将完备空间 $[a,b]$ 映入自身，且
\[
|T(x)-T(y)|\le\frac{1+L}{2}|x-y|.
\]
因此 $T$ 是压缩映射。由压缩映射原理，$x_n=T^{n-1}(x_1)$ 收敛到唯一的不动点 $x_*$；而 $T(x_*)=x_*$ 等价于 $f(x_*)=x_*$。

也可直接展开压缩映射原理在本题中的证明。记 $q=(1+L)/2<1$，则 $|x_{n+1}-x_n|\le q^{n-1}|x_2-x_1|$；对 $m>n$ 累加相邻差，得 $|x_m-x_n|\le q^{n-1}|x_2-x_1|/(1-q)$。因此数列为 Cauchy 列，在闭区间内收敛。$f$ 的 Lipschitz 条件保证连续性，对递推取极限得 $x_*=(x_*+f(x_*))/2$；如果还有另一不动点，两点距离不超过自身的 $L$ 倍，只能为零。

\end{solve}

\begin{remark}[方法与条件]更一般地，松弛迭代 $T_\theta=(1-\theta)I+\theta f$ 在 $0<\theta\le1$ 时仍为压缩映射，压缩常数不超过 $1-\theta(1-L)$。本题的平均操作就是 $\theta=1/2$。\end{remark}

\begin{example}[来源序号4]\label{ex:fin04-all:4}\label{pr:fin04-all:4}

设 $f:(a,b)\to\mathbb R$ 可导。若存在 $x_0$ 的邻域，使其上每个 $x\ne x_0$ 都满足
\[
f(x)>f(x_0)+f'(x_0)(x-x_0),
\]
则称 $x_0$ 为严格下支撑点；反向严格不等式定义严格上支撑点。证明：若 $f$ 不是仿射函数，则至少存在一种这样的点。

\par\href{https://github.com/lzxcr/cmc-mathematics/blob/64916fa1fdbc3e444a652eca204d5669058ee321/src/problems/finals/04/04.md}{题目与来源条目}

\end{example}

\begin{solve}\label{ans:fin04-all:4}

由于 $f$ 非仿射，可取 $u<c<v$，使 $f$ 的图像在这三点不共线。令 $\ell$ 为连接两端点的仿射函数，并设 $h=f-\ell$。不妨将 $f$ 换成 $-f$，使 $h(c)>0$，而 $h(u)=h(v)=0$。

取充分小的 $\varepsilon>0$，使函数
$H(x)=h(x)+\varepsilon(x-c)^2$ 在两个端点的值都小于 $H(c)$。其最大值因而在某个 $\xi\in(u,v)$ 达到，且 $H'(\xi)=0$。展开 $H(x)\le H(\xi)$，得到
\[
f(x)\le f(\xi)+f'(\xi)(x-\xi)-\varepsilon(x-\xi)^2
\qquad(u\le x\le v).
\]
因此 $\xi$ 是严格上支撑点。若先前将 $f$ 换成了 $-f$，则得到原函数的严格下支撑点。

可以明确选取 $0<\varepsilon<h(c)/\max\{(u-c)^2,(v-c)^2\}$，以保证最大值不落在两个端点。又 $H'(\xi)=0$ 给出 $h'(\xi)=-2\varepsilon(\xi-c)$，于是
\[
h(x)-h(\xi)\le\varepsilon[(\xi-c)^2-(x-c)^2]
=h'(\xi)(x-\xi)-\varepsilon(x-\xi)^2.
\]
加回仿射函数后得到所列切线不等式。只要 $x\ne\xi$，最后的负平方项就严格为负，即使 $H$ 的最大值不是唯一，也仍能得到严格支撑点。

\end{solve}

\begin{remark}[方法与条件]这里采用“严格下支撑点”“严格上支撑点”的术语，以免与函数在整个区间上的凹凸性混淆。“非仿射”明确排除所有 $f(x)=ax+b$，包括常数函数；这类函数处处与其切线重合，不能满足严格支撑不等式。证明的核心是用微小二次扰动把非严格接触变成严格接触，不需要二阶导数。\end{remark}

\begin{example}[来源序号6]\label{ex:fin04-all:6}\label{pr:fin04-all:6}

设 $n$ 为非负整数，$\mathcal P_n$ 为所有 $n$ 次首一实系数多项式组成的集合。证明
\[
\inf_{\substack{b\in\mathbb R,\ a>0\\P\in\mathcal P_n}}
\frac1{a^{n+1}}\int_b^{b+a}|P(x)|\,dx>0.
\]

\par\href{https://github.com/lzxcr/cmc-mathematics/blob/64916fa1fdbc3e444a652eca204d5669058ee321/src/problems/finals/04/06.md}{题目与来源条目}

\end{example}

\begin{solve}\label{ans:fin04-all:6}

令 $Q(t)=a^{-n}P(b+at)$，则 $Q$ 仍为 $n$ 次首一多项式，且
\[
\frac1{a^{n+1}}\int_b^{b+a}|P(x)|\,dx=\int_0^1|Q(t)|\,dt.
\]
现在只需处理固定区间。在次数不超过 $n$ 的多项式空间上，
$\|Q\|_1=\int_0^1|Q(t)|\,dt$ 是范数，而最高次项系数 $\ell(Q)$ 是线性泛函。有限维空间上的线性泛函连续，故存在仅依赖 $n$ 的 $C_n>0$，使
\[
|\ell(Q)|\le C_n\|Q\|_1.
\]
对首一多项式有 $\ell(Q)=1$，因此所求下界至少为 $C_n^{-1}>0$。

有限维范数结论在这里可直接证明。对系数向量 $c=(c_0,\ldots,c_n)$ 的欧氏单位球面，函数 $J(c)=\int_0^1|\sum_{j=0}^nc_jt^j|\,dt$ 连续：三角不等式把其变化控制在 $\sum_j|c_j-d_j|/(j+1)$ 之内。非零多项式不可能在区间上恒为零，所以 $J(c)>0$；紧性保证最小值 $m_n>0$。按系数长度缩放得 $\|Q\|_1\ge m_n\|c\|_2\ge m_n|c_n|$，首一时 $c_n=1$，于是显式获得所需的正下界。

\end{solve}

\begin{remark}[方法与条件]上述有限维结论也可直接由紧致性证明：在系数向量的单位球面上，连续函数 $Q\mapsto\int_0^1|Q|$ 处处为正，所以有正的最小值。题目中的 $a^{n+1}$ 正是次数为 $n$ 的多项式在长度缩放下的自然量纲；先缩放，再用有限维紧致性，免去了逐次差分估计。\end{remark}

\section{决赛：第 5 届（2014 年）}

本节选入4道课程相关题；来源序号保留原题解编排，未选入题目不重新补号。

\begin{remark}[来源说明]原卷第1至4题为必做题，第5至11题任选两题；本选编按课程范围选入其中相关题目。高、低年级重复题已合并，并保留两卷原题号。\end{remark}

\begin{example}[来源序号3]\label{ex:fin05-hi:3}\label{pr:fin05-hi:3}

设 $f:[0,\infty)\to[0,\infty)$ 可导，$f(0)=0$，且 $f'(x)\le1/2$。若 $M=\int_0^\infty f(x)\,dx<\infty$，证明对任意 $\alpha>1$，积分 $\int_0^\infty f(x)^\alpha\,dx$ 收敛，并且
\[
\int_0^\infty f(x)^\alpha\,dx\le M^{(\alpha+1)/2}.
\]

\par\href{https://github.com/lzxcr/cmc-mathematics/blob/64916fa1fdbc3e444a652eca204d5669058ee321/src/problems/finals/05/03.md}{题目与来源条目}

\end{example}

\begin{solve}\label{ans:fin05-hi:3}

令 $F(t)=\int_0^t f(x)\,dx$。由于
\[
(F-f^2)'=f(1-2f')\ge0,\qquad F(0)-f(0)^2=0,
\]
有 $f(t)^2\le F(t)\le M$。于是已经得到
$\int f^\alpha\le M^{(\alpha-1)/2}\int f=M^{(\alpha+1)/2}$。

事实上保留 $F(t)$ 可以得到更强的估计：
\[
\int_0^T f^\alpha
\le\int_0^T fF^{(\alpha-1)/2}
=\frac{2}{\alpha+1}F(T)^{(\alpha+1)/2}.
\]
令 $T\to\infty$，便得收敛性以及系数为 $2/(\alpha+1)$ 的加强结论。

$F-f^2$ 可导且导数非负，所以由中值定理在半轴上递增，无需对可能不连续的 $f'$ 作积分。又 $f\ge0$ 且 $\alpha>1$，因此从 $f^2\le M$ 得 $f^{\alpha-1}\le M^{(\alpha-1)/2}$，乘以 $f$ 后可与已知可积函数比较。加强式中用 $F'=f$，故 $fF^{(\alpha-1)/2}$ 是 $2F^{(\alpha+1)/2}/(\alpha+1)$ 的导数；在 $F=0$ 处这个等式仍成立，因为幂指数大于 $1$。

\end{solve}

\begin{remark}[方法与条件]导数的单侧上界限制了函数“长到某个高度”之前必须积累的面积。辅助量 $F-f^2$ 把这一几何事实变成单调性；不需要对下降速度作任何限制。\end{remark}

\begin{example}[来源序号4]\label{ex:fin05-hi:4}\label{pr:fin05-hi:4}

对所有根均为实数的多项式 $p$，以 $d(p)$ 表示其最大根与最小根之差。设 $n\ge2$，求使每个这样的 $n$ 次多项式均满足
\[
d(p')\ge C\,d(p)
\]
的最大常数 $C$。

\par\href{https://github.com/lzxcr/cmc-mathematics/blob/64916fa1fdbc3e444a652eca204d5669058ee321/src/problems/finals/05/04.md}{题目与来源条目}

\end{example}

\begin{solve}\label{ans:fin05-hi:4}

答案是 $C=\sqrt{(n-2)/n}$。$n=2$ 时结论显然，以下设 $n>2$。

设 $p$ 首一，根为 $r_1\le\cdots\le r_n$。令
$D=\operatorname{diag}(r_1,\ldots,r_n)$、$e=(1,\ldots,1)^{\mathsf T}/\sqrt n$，并把 $D$ 正交压缩到 $e^\perp$，记所得对称算子为 $T$。分块行列式公式给出
\[
\det(tI_{n-1}-T)
=\det(tI_n-D)\,e^{\mathsf T}(tI_n-D)^{-1}e
=\frac1n p'(t).
\]
这里先在 $t\ne r_i$ 时计算，再按多项式恒等式延拓。因此 $p'$ 的根就是 $T$ 的特征值。

压缩算子指先用 $D$ 作用，再正交投影回 $e^\perp$。取以 $e$ 为最后一个向量的标准正交基，则 $tI-D$ 的左上块是 $tI_{n-1}-T$；逆矩阵右下角元素等于该块行列式除以全行列式，这就是上面使用的分块公式。另一方面，$e^{\mathsf T}(tI-D)^{-1}e=n^{-1}\sum_i(t-r_i)^{-1}$，乘以 $p(t)=\prod_i(t-r_i)$ 后正好是 $p'(t)/n$。重根情形通过多项式恒等式同样包含在内。

以 $e_i$ 表示标准基。在 $e^\perp$ 中取两个正交单位向量
\[
u=\frac{e_1-e_n}{\sqrt2},\qquad
v=\sqrt{\frac{n}{2(n-2)}}
\left(e_1+e_n-\frac2n\sum_{j=1}^n e_j\right).
\]
有
\[
u^{\mathsf T}Tv=u^{\mathsf T}Dv
=\frac{r_1-r_n}{2}\sqrt{\frac{n-2}{n}}.
\]
单位向量 $(u+v)/\sqrt2$ 与 $(u-v)/\sqrt2$ 的 Rayleigh 商之差为 $2u^{\mathsf T}Tv$，而任意两个 Rayleigh 商之差的绝对值不超过谱直径。故
\[
d(p')\ge2|u^{\mathsf T}Tv|=\sqrt{\frac{n-2}{n}}\,d(p).
\]
最后取 $p(x)=x(x-1)(x-1/2)^{n-2}$。令 $z=x-1/2$，则
\[
p'(x)=z^{n-3}\left(nz^2-\frac{n-2}{4}\right),
\]
从而 $d(p')=\sqrt{(n-2)/n}$、$d(p)=1$，证明常数最优。

对实对称算子，任意单位向量在一组标准正交特征向量下展开，Rayleigh 商 $z^{\mathsf T}Tz$ 是特征值以坐标平方为权重的凸组合，所以必介于最小、最大特征值之间。这解释了为何两个试探向量的商之差能控制导数根的直径。取等多项式在平移后为 $(z^2-1/4)z^{n-2}$，求导即得所列二次因子，其两个端根为 $\pm\tfrac12\sqrt{(n-2)/n}$，中间可能出现的零根不改变直径。

\end{solve}

\begin{remark}[方法与条件]实根多项式求导可以理解为实对称算子的正交压缩。Rolle 定理的零点交错随之成为谱交错，而本题的定量界来自两个明确试验向量的 Rayleigh 商。这一联系把零点几何纳入了统一的谱理论。\end{remark}

\begin{example}[来源序号5]\label{ex:fin05-hi:5}\label{pr:fin05-hi:5}

设 $f(x,y)$ 在 $[a,b]\times\mathbb R$ 上关于 $y$ 单调递减。可导函数 $y,z$ 满足
\[
y'(x)=f(x,y(x)),\qquad z'(x)\le f(x,z(x)),\qquad z(a)\le y(a).
\]
证明 $z(x)\le y(x)$ 对所有 $x\in[a,b]$ 成立。

\par\href{https://github.com/lzxcr/cmc-mathematics/blob/64916fa1fdbc3e444a652eca204d5669058ee321/src/problems/finals/05/05.md}{题目与来源条目}

\end{example}

\begin{solve}\label{ans:fin05-hi:5}

令 $w=z-y$。在开集 $\{w>0\}$ 上，由单调性有
\[
w'\le f(x,z)-f(x,y)\le0.
\]
若该开集非空，取其一个区间分支，设左端点为 $c$。连续性与初值给出 $w(c)=0$，但 $w$ 在该分支上不增，因此不可能为正。矛盾。

若存在 $x_*>a$ 使 $w(x_*)>0$，考察包含它的正值区间，取左端点 $c\ge a$。当 $c>a$ 时由连续性与区间的极大性得 $w(c)=0$；当 $c=a$ 时，初值给 $w(a)\le0$，右侧正值又迫使 $w(a)\ge0$，仍为零。在区间中任取 $c<t<x_*$，由中值定理及 $w'\le0$ 得 $w(x_*)\le w(t)$；令 $t\downarrow c$，得到 $w(x_*)\le0$，矛盾。这避免了只凭“第一次相交时的导数方向”作不充分判断。

\end{solve}

\begin{remark}[方法与条件]这是标量微分方程的比较原理。只在可能违反结论的区域比较导数，便可排除解曲线向上穿越；无须求出解，也不要求关于 $y$ 的 Lipschitz 条件。\end{remark}

\begin{example}[来源序号12，高年级组选做]\label{ex:fin05-hi:11}\label{pr:fin05-hi:11}

定义多项式列
\[
T_0(x)=1,\qquad T_1(x)=x,\qquad
T_{n+1}(x)=2xT_n(x)-T_{n-1}(x)\quad(n\ge1).
\]

（1）证明 $T_n(x)=\cos(n\arccos x)$（$-1\le x\le1$）。
（2）证明 $(T_n)$ 在内积
\[
\langle f,g\rangle=\int_{-1}^1\frac{f(x)g(x)}{\sqrt{1-x^2}}\,dx
\]
下两两正交。
（3）设 $n\ge1$，$P$ 为 $n$ 次首一实多项式。证明
\[
\|P\|_\infty\ge2^{1-n},\qquad
\|P\|_\infty:=\max_{[-1,1]}|P|,
\]
且等号当且仅当 $P=2^{1-n}T_n$ 时成立。

\par\href{https://github.com/lzxcr/cmc-mathematics/blob/64916fa1fdbc3e444a652eca204d5669058ee321/src/problems/finals/05/12.md}{题目与来源条目}

\end{example}

\begin{solve}\label{ans:fin05-hi:11}

余弦恒等式
$2\cos\theta\cos(n\theta)=\cos((n+1)\theta)+\cos((n-1)\theta)$
与初值相同，故递推即得第一问。令 $x=\cos\theta$，又有
\[
\langle T_n,T_m\rangle=\int_0^\pi\cos(n\theta)\cos(m\theta)\,d\theta=0
\qquad(n\ne m).
\]

递推表明 $T_n$ 的首项系数为 $2^{n-1}$，且在节点 $x_j=\cos(j\pi/n)$（$0\le j\le n$）上满足 $T_n(x_j)=(-1)^j$。记
\[
w_j=\frac1{\left|\prod_{k\ne j}(x_j-x_k)\right|}>0.
\]
Lagrange 插值公式中比较最高次项，得
\[
1=\sum_{j=0}^n(-1)^jw_jP(x_j)
\le\|P\|_\infty\sum_{j=0}^n w_j.
\]
把 $P$ 换成首一多项式 $2^{1-n}T_n$，知 $\sum_jw_j=2^{n-1}$，故得到所求下界。

若等号成立，由所有 $w_j>0$，必有 $P(x_j)=2^{1-n}(-1)^j$ 对每个 $j$ 成立。这给出 $n+1$ 个插值条件，唯一确定 $P=2^{1-n}T_n$。反向显然成立。

插值节点按 $x_0>x_1>\cdots>x_n$ 排列，故 $\prod_{k\ne j}(x_j-x_k)$ 中恰有 $j$ 个负因子，其倒数为 $(-1)^jw_j$。Lagrange 基多项式的最高次项系数正是这个倒数，因此比较首项系数得到求和恒等式。又对每个 $j$ 都有 $(-1)^jP(x_j)\le\|P\|_\infty$，若加权求和取等且每个权重为正，就必须每一项都取等。最后两个次数不超过 $n$ 的多项式在 $n+1$ 个不同点相同，其差只能是零多项式，从而得到唯一性。

\end{solve}

\begin{remark}[方法与条件]Chebyshev 多项式一面来自圆周上的 Fourier 正交性，一面解决一致范数下的最优逼近问题。极值点的交替符号决定最优性，正的插值权重进一步给出唯一性，避免了仅靠非严格变号计数可能留下的空隙。\end{remark}

\section{决赛：第 6 届（2015 年）}

本节选入5道课程相关题；来源序号保留原题解编排，未选入题目不重新补号。

\begin{example}[来源序号2]\label{ex:fin06-lo:2}\label{pr:fin06-lo:2}

（填空题）级数 $\sum_{n=1}^{\infty}\frac{n}{3^n}$ 的和为 $\underline{\qquad\qquad}$。

\par\href{https://github.com/lzxcr/cmc-mathematics/blob/64916fa1fdbc3e444a652eca204d5669058ee321/src/problems/finals/06/02.md}{题目与来源条目}

\end{example}

\begin{solve}\label{ans:fin06-lo:2}

答案为 $\frac34$。对几何级数在收敛圆 $|x|<1$ 内逐项求导，得
\[
\sum_{n\ge1}nx^n=x\frac{\mathrm d}{\mathrm dx}\frac1{1-x}
=\frac{x}{(1-x)^2}.
\]
取 $x=1/3$ 即得 $\sum_{n\ge1}n/3^n=3/4$。

幂级数逐项求导在任意 $|x|\le r<1$ 上合法，因为导数项的绝对值被 $nr^{n-1}$ 控制，而其级数收敛。也可用纯求和方法核查：记 $S=\sum_{n\ge1}n/3^n$，则 $S/3=\sum_{n\ge2}(n-1)/3^n$，相减得 $2S/3=\sum_{n\ge1}3^{-n}=1/2$，所以 $S=3/4$；这些相减操作因原级数绝对收敛而成立。

\end{solve}

\begin{example}[来源序号3]\label{ex:fin06-lo:3}\label{pr:fin06-lo:3}

（填空题）计算 $I=\iint_{x^2+y^2+z^2=1}\left(x^2+2y^2+3z^2\right)\mathrm{d}S=\underline{\qquad\qquad}$。

\par\href{https://github.com/lzxcr/cmc-mathematics/blob/64916fa1fdbc3e444a652eca204d5669058ee321/src/problems/finals/06/03.md}{题目与来源条目}

\end{example}

\begin{solve}\label{ans:fin06-lo:3}

答案为 $8\pi$。单位球面及其面积元在坐标交换下不变，故
\[
\iint_Sx^2\,\mathrm dS=\iint_Sy^2\,\mathrm dS
=\iint_Sz^2\,\mathrm dS
=\frac13\iint_S(x^2+y^2+z^2)\,\mathrm dS=\frac{4\pi}{3}.
\]
所求积分因而为 $(1+2+3)\cdot4\pi/3=8\pi$。

坐标交换是保持球面与面积的正交变换，因此三项积分相等；这里相等的是积分，不能说球面上逐点有 $x^2=y^2=z^2$。单位球面上逐点成立的是 $x^2+y^2+z^2=1$，积分后总和等于球面积 $4\pi$，所以每项为 $4\pi/3$。若要从参数公式核对面积，取球坐标面积元 $\sin\theta\,d\theta\,d\varphi$，在 $0\le\theta\le\pi$、$0\le\varphi\le2\pi$ 上积分即得 $4\pi$。

\end{solve}

\begin{example}[来源序号7]\label{ex:fin06-lo:7}\label{pr:fin06-lo:7}

设
\[
f(x)=\begin{cases}x\sin\dfrac{1}{x},&x\ne0,\\0,&x=0,\end{cases}
\]
求最大常数 $\alpha$，使
\[
\sup_{x\ne y}\frac{|f(x)-f(y)|}{|x-y|^{\alpha}}<+\infty .
\]

\par\href{https://github.com/lzxcr/cmc-mathematics/blob/64916fa1fdbc3e444a652eca204d5669058ee321/src/problems/finals/06/07.md}{题目与来源条目}

\end{example}

\begin{solve}\label{ans:fin06-lo:7}

最大值为 $\alpha=1/2$。因为 $f$ 是偶函数且 $\bigl||x|-|y|\bigr|\le|x-y|$，证明上界只需考虑 $0\le x<y$。记 $d=y-x$；由 $|f|\le1$，$d\ge1$ 时已有 $|f(y)-f(x)|\le2\sqrt d$。

设 $0<d<1$。在距零点约 $r$ 处，振幅约为 $r$，而导数给出的增量界约为 $d/r$；平衡二者，取 $r=\sqrt d$。

若 $x\le r$，则 $y=x+d\le2r$，直接用振幅控制：
\[
|f(y)-f(x)|\le x+y\le3r=3\sqrt d.
\]
若 $x>r$，区间 $[x,y]$ 避开零点。由 $|f'(t)|\le1+1/t$，
\[
|f(y)-f(x)|
\le\int_x^y|f'(t)|\,\mathrm dt
\le d\left(1+\frac1x\right)
\le d+\frac d r\le2\sqrt d.
\]

故 $f$ 在全实轴上满足统一的 $1/2$ 阶 Hölder 估计。

再取
\[
x_n=\frac1{n\pi},\qquad y_n=\frac1{(n+1/2)\pi}.
\]
此时 $|f(x_n)-f(y_n)|=y_n\sim1/(\pi n)$，而 $|x_n-y_n|\sim1/(2\pi n^2)$，故
\[
\frac{|f(x_n)-f(y_n)|}{|x_n-y_n|^\alpha}
\sim2^\alpha\pi^{\alpha-1}n^{2\alpha-1}.
\]
对任意 $\alpha>1/2$，此比值趋于无穷，证明了最优性。

全局有界性由 $|\sin u|\le|u|$ 给出 $|x\sin(1/x)|\le1$；近零时另有 $|f(x)|\le|x|$，这两个界负责不同区域。反例点列让一个相位恰为正弦的零点，另一个相位恰为峰值，所以函数值差约为 $1/n$，而横坐标距离只有 $1/n^2$。若要更高的 Hölder 阶，就要求纵向变化比这一平方根尺度更小，与这组实际振荡冲突。

\end{solve}

\begin{remark}[方法与条件]这里只比较 $f(x)$ 与 $f(0)$ 时有 $|f(x)-f(0)|\le|x|$；但本题要求对任意两点统一估计。距零点约 $x$ 处，一次显著振荡的距离约为 $x^2$，函数值却改变约 $x$，因而最优指数降为 $1/2$。

这一证明可推广到 $x^a\sin(x^{-b})$：在 $[0,1]$ 上，$a,b>0$ 时的最优一致指数为 $\min\{1,a/(b+1)\}$，见\hyperref[app:estimate-holder]{振幅与局部导数的两尺度估计}。若一般函数绝对连续且导数属于 $L^p$，也可直接对导数积分使用 Hölder 不等式，见\hyperref[app:estimate-derivative]{导数积分与连续模}。本题的 $f'$ 在零点附近绝对不可积，故不能直接使用这一全区间方法；解答中的积分只在避开零点的区间上进行。

若允许 $\alpha$ 取任意实数，则全实轴上的有界性恰在 $0\le\alpha\le1/2$ 时成立。负指数不能由有界函数自动得到：固定 $x=0$、令 $y\to\infty$，由 $f(y)\to1$ 可知该比值在 $\alpha<0$ 时仍趋于无穷。\end{remark}

\begin{example}[来源序号8]\label{ex:fin06-lo:8}\label{pr:fin06-lo:8}

$a(t)$，$f(t)$ 为实连续函数，$\forall t\in\mathbb{R}$，有
\[
f(t)>0,\qquad a(t)\ge1,\qquad \int_0^{\infty}f(t)\,\mathrm{d}t=+\infty .
\]
已知 $x(t)$ 满足 $x''(t)+a(t)f(x(t))\le0$，$\forall t\in\mathbb{R}$。求证：$x(t)$ 在 $[0,+\infty)$ 上有上界。

\par\href{https://github.com/lzxcr/cmc-mathematics/blob/64916fa1fdbc3e444a652eca204d5669058ee321/src/problems/finals/06/08.md}{题目与来源条目}

\end{example}

\begin{solve}\label{ans:fin06-lo:8}

令 $F(s)=\int_0^sf(u)\,\mathrm du$，则 $F$ 严格递增且 $F(s)\to\infty$（$s\to\infty$）。由题设 $x''<0$，故 $x'$ 严格递减。

若 $x$ 在正半轴无上界，则必有 $x'(t)>0$ 对所有 $t\ge0$ 成立：一旦某处 $x'\le0$，此后 $x$ 就单调不增。因而此时 $x(t)\to\infty$。但在这条递增轨道上，能量
\[
E(t)=\frac12x'(t)^2+F(x(t))
\]
满足
\[
E'(t)=x'(t)\bigl(x''(t)+f(x(t))\bigr)
\le-(a(t)-1)x'(t)f(x(t))\le0.
\]
于是 $F(x(t))\le E(t)\le E(0)$，与 $F(x(t))\to\infty$ 矛盾。故 $x$ 在 $[0,\infty)$ 上有上界。

若在某个 $t_0$ 有 $x'(t_0)\le0$，由于 $x'$ 严格递减，$x$ 在其后不增，而 $[0,t_0]$ 上连续函数有界，故全半轴有上界。因此在反设下确可乘以正数 $x'$，保持不等号方向。能量论证只用到 $F(s)\to\infty$ 当 $s\to+\infty$，并不要求 $F$ 在全部实数上非负；因为动能 $x'^2/2\ge0$，仍然有 $F(x(t))\le E(t)$，足以排除轨道向正无穷增长。

\end{solve}

\begin{remark}[方法与条件]核心是沿假设存在的递增轨道使用\hyperref[app:dynamics-energy]{能量}：无穷远的势能壁垒阻止轨道逃向正无穷。若 $a(t)\ge a_0>0$，改用 $x'^2/2+a_0F(x)$ 即可。

条件 $\int_0^\infty f=\infty$ 不能直接删去。取
\[
a\equiv1,\qquad f(s)=\frac{e^s}{(1+e^s)^3},
\qquad x(t)+e^{x(t)}=t.
\]
映射 $s\mapsto s+e^s$ 是全实轴上的光滑严格递增双射，故 $x$ 在全实轴上有定义，且
\[
x'=\frac1{1+e^x},\qquad
x''=-\frac{e^x}{(1+e^x)^3}=-f(x).
\]
这里 $\int_0^\infty f(s)\,\mathrm ds=1/8<\infty$，而 $x(t)\to\infty$。所有定义域与正性条件均保留。另一方面，取 $f\equiv1$、$a\equiv1$、$x=-t^2$，可见原结论不能加强为同时有下界。\end{remark}

\begin{example}[来源序号9，低年级组]\label{ex:fin06-lo:9}\label{pr:fin06-lo:9}

设 $f(x)$ 在区间 $[0,1]$ 上连续可导，且 $f(0)=f(1)=0$。求证：
\[
\left[\int_0^1 xf(x)\,\mathrm{d}x\right]^2\le\frac{1}{45}\int_0^1\left(f'(x)\right)^2\mathrm{d}x,
\]
等号当且仅当 $f(x)=A\left(x-x^3\right)$ 时成立，其中 $A$ 是常数。

\par\href{https://github.com/lzxcr/cmc-mathematics/blob/64916fa1fdbc3e444a652eca204d5669058ee321/src/problems/finals/06/09.md}{题目与来源条目}

\end{example}

\begin{solve}\label{ans:fin06-lo:9}

由边界条件，$\int_0^1f'=0$。分部积分后，可自由给积分核加常数，于是
\[
\int_0^1xf(x)\,\mathrm dx
=-\frac12\int_0^1x^2f'(x)\,\mathrm dx
=\int_0^1\frac{1-3x^2}{6}f'(x)\,\mathrm dx.
\]
Cauchy 不等式给出
\[
\left(\int_0^1xf\right)^2
\le\left(\int_0^1\left(\frac{1-3x^2}{6}\right)^2\mathrm dx\right)
\left(\int_0^1(f')^2\mathrm dx\right)
=\frac1{45}\int_0^1(f')^2\mathrm dx.
\]
等号当且仅当 $f'$ 与 $1-3x^2$ 成比例；由连续性及 $f(0)=0$，这等价于 $f(x)=A(x-x^3)$。这些函数也满足 $f(1)=0$，故等号条件完整，常数 $1/45$ 最佳。

核中加入的常数 $1/6$ 是为了得到最小的平方积分。对任意常数 $c$，$\int xf=\int(c-x^2/2)f'$，而
\[
\int_0^1(c-x^2/2)^2\,dx=c^2-\frac c3+\frac1{20}
=(c-1/6)^2+\frac1{45}.
\]
因此选 $c=1/6$ 正好得到最优 Cauchy 界。等号时 $f'=A(1-3x^2)$，积分给 $f=A(x-x^3)+B$，由 $f(0)=0$ 消去 $B$；任何非零 $A$ 都取得等号，故常数不能减小。

\end{solve}

\begin{remark}[方法与条件]加上常数 $1/6$，正是把核 $-x^2/2$ 正交投影到常数函数的正交补，使其平方积分最小。等价地，令 $u=(x-x^3)/6$，则 $-u''=x$、$u(0)=u(1)=0$，而 $\int xf=\int u'f'$：最优常数就是这个代表元的能量 $\int(u')^2=1/45$。这一方法可系统处理带零边界条件的\hyperref[app:inequality-energy]{积分型最优常数问题}。\end{remark}

\section{决赛：第 7 届（2016 年）}

本节选入6道课程相关题；来源序号保留原题解编排，未选入题目不重新补号。

\begin{example}[来源序号2]\label{ex:fin07-hi:2}\label{pr:fin07-hi:2}

（填空题）令 $a_{n}=\int_{0}^{\frac{\pi}{4}}\tan^{n}x\,\mathrm{d}x$。若 $\sum_{n=1}^{\infty}a_{n}^{p}$ 收敛，则 $p$ 取值范围 $\underline{\qquad\qquad}$。

\par\href{https://github.com/lzxcr/cmc-mathematics/blob/64916fa1fdbc3e444a652eca204d5669058ee321/src/problems/finals/07/02.md}{题目与来源条目}

\end{example}

\begin{solve}\label{ans:fin07-hi:2}

答案为 $p>1$。换元 $t=\tan x$ 得
\[
\frac{1}{2(n+1)}\leq a_n=\int_0^1\frac{t^n}{1+t^2}\,\mathrm{d}t\leq\frac{1}{n+1}.
\]
故 $p>0$ 时与 $\sum n^{-p}$ 同敛散；$p\leq0$ 时通项不趋于零。

换元后 $dx=dt/(1+t^2)$，且 $t\in[0,1]$，所以 $1/2\le(1+t^2)^{-1}\le1$，积分即可得到双侧界。当 $p>0$ 时两边取 $p$ 次方不改变方向，从而 $2^{-p}(n+1)^{-p}\le a_n^p\le(n+1)^{-p}$，可用正项级数比较判别法。若 $p=0$，通项等于 $1$；若 $p<0$，因 $a_n\to0^+$，通项反而趋于无穷。于是恰好只有 $p>1$ 收敛。

\end{solve}

\begin{example}[来源序号3]\label{ex:fin07-hi:3}\label{pr:fin07-hi:3}

（填空题）设 $D:x^{2}+2y^{2}\leq 2x+4y$，则积分
\[
I=\iint_{D}(x+y)\,\mathrm{d}x\mathrm{d}y=\underline{\hspace{2cm}}.
\]

\par\href{https://github.com/lzxcr/cmc-mathematics/blob/64916fa1fdbc3e444a652eca204d5669058ee321/src/problems/finals/07/03.md}{题目与来源条目}

\end{example}

\begin{solve}\label{ans:fin07-hi:3}

答案为 $3\sqrt2\pi$。平移 $u=x-1,v=y-1$ 后，区域为 $u^2+2v^2\leq3$，被积函数为 $2+u+v$。中心对称使后两项积分为零；椭圆面积为 $3\pi/\sqrt2$，故 $I=2\cdot3\pi/\sqrt2$。

完成平方得 $(x-1)^2+2(y-1)^2\le3$，所以平移后的半轴长分别是 $\sqrt3$ 与 $\sqrt{3/2}$，面积为二者乘积再乘 $\pi$。若记区域为 $E$，变换 $(u,v)\mapsto(-u,-v)$ 保持区域及面积元，同时使 $u+v$ 变号，所以 $\iint_E(u+v)=0$。被积函数的常数部分为 $2$，因此积分等于中心点处的函数值乘以椭圆面积。

\end{solve}

\begin{example}[来源序号7]\label{ex:fin07-hi:7}\label{pr:fin07-hi:7}

设单位圆 $\Gamma$ 的外切 $n$ 边形 $A_{1}A_{2}\cdots A_{n}$ 各边与 $\Gamma$ 分别切于 $B_{1},B_{2},\cdots,B_{n}$。令 $P_{A},P_{B}$ 分别表示多边形 $A_{1}A_{2}\cdots A_{n}$ 与 $B_{1}B_{2}\cdots B_{n}$ 的周长。求证：
\[
P_{A}^{\frac{1}{3}}P_{B}^{\frac{2}{3}}>2\pi.
\]

\par\href{https://github.com/lzxcr/cmc-mathematics/blob/64916fa1fdbc3e444a652eca204d5669058ee321/src/problems/finals/07/07.md}{题目与来源条目}

\end{example}

\begin{solve}\label{ans:fin07-hi:7}

设相邻切点所对圆心角的一半为 $\alpha_i\in(0,\pi/2)$，则
\[
\sum_i\alpha_i=\pi,\qquad P_A=2\sum_i\tan\alpha_i,\qquad P_B=2\sum_i\sin\alpha_i.
\]
目标中的指数 $1/3,2/3$ 提示先对每个角证明 $\tan^{1/3}x\sin^{2/3}x>x$，再用 Hölder 将各项组合。等价地，先建立点态估计 $\sin x>x\cos^{1/3}x$。令 $g(x)=\sin x/\cos^{1/3}x-x$，则 $g(0)=0$，且对三个正数 $1,\cos^2x,\cos^2x$ 使用算术---几何平均不等式，
\[
g'(x)=\frac{1+2\cos^2x}{3\cos^{4/3}x}-1>0\qquad(0<x<\pi/2).
\]
再用 Hölder 不等式，得到
\[
P_A^{1/3}P_B^{2/3}\geq2\sum_i(\tan\alpha_i)^{1/3}(\sin\alpha_i)^{2/3}>2\sum_i\alpha_i=2\pi.
\]

连接圆心与相邻切点、切线交点，可得到直角三角形：单位半径对应半弦长 $\sin\alpha_i$，每段切线长为 $\tan\alpha_i$，所以两种周长公式成立。点态严格性来自算术—几何平均的等号条件：$1,\cos^2x,\cos^2x$ 在 $0<x<\pi/2$ 时不全相等。最后 Hölder 使用指数 $3$ 与 $3/2$，即 $\sum a_i^{1/3}b_i^{2/3}\le(\sum a_i)^{1/3}(\sum b_i)^{2/3}$；有限和中每个点态不等式都严格，所以最终严格大于 $2\pi$。

\end{solve}

\begin{remark}[方法与条件]常数 $2\pi$ 最佳：取外切正 $n$ 边形，令 $n\to\infty$，两种周长都趋于 $2\pi$。半径为 $r$ 时，右端相应为 $2\pi r$。\end{remark}

\begin{example}[来源序号8，低年级组]\label{ex:fin07-lo:8}\label{pr:fin07-lo:8}

设 $a(x),f(x)$ 为 $\mathbb{R}$ 上的连续函数，且对任意 $x\in\mathbb{R}$ 有 $a(x)>0$。已知
\[
\int_{0}^{\infty}a(x)\,\mathrm{d}x=+\infty,\quad \lim_{x\to+\infty}\frac{f(x)}{a(x)}=0,\quad y'(x)+a(x)y(x)=f(x),\ x\in\mathbb{R},
\]
求证：$\lim_{x\to+\infty}y(x)=0$。

\par\href{https://github.com/lzxcr/cmc-mathematics/blob/64916fa1fdbc3e444a652eca204d5669058ee321/src/problems/finals/07/08.md}{题目与来源条目}

\end{example}

\begin{solve}\label{ans:fin07-lo:8}

积分因子把初值的影响与强迫项分开；题设的两个极限分别控制这两部分。令 $A(x)=\int_0^xa(t)\,\mathrm{d}t$，则 $A(x)\to\infty$。给定 $\varepsilon>0$，取 $T$ 使 $|f(t)|\leq\varepsilon a(t)$（$t\geq T$）。常数变易公式给出
\[
|y(x)|\leq e^{A(T)-A(x)}|y(T)|+\varepsilon\int_T^xa(t)e^{A(t)-A(x)}\,\mathrm{d}t
=e^{A(T)-A(x)}|y(T)|+\varepsilon(1-e^{A(T)-A(x)}).
\]
故 $\limsup_{x\to\infty}|y(x)|\leq\varepsilon$。令 $\varepsilon\downarrow0$ 即得结论。

具体地，方程乘以 $e^{A(x)}$ 后变为 $(e^Ay)'=e^Af$，在 $[T,x]$ 上积分再除以 $e^{A(x)}$ 得到常数变易公式。由于 $A'=a$，积分核满足 $\int_T^xa(t)e^{A(t)-A(x)}dt=[e^{A(t)-A(x)}]_T^x$，于是得到精确的 $1-e^{A(T)-A(x)}$。先固定由 $\varepsilon$ 决定的 $T$，再让 $x\to\infty$，初值项被无限累积的阻尼消去；最后才让 $\varepsilon\to0$。

\end{solve}

\begin{remark}[方法与条件]若 $f(x)/a(x)\to L$，对 $y-L$ 应用同一证明，便得 $y(x)\to L$。而 $\int_0^\infty a=\infty$ 不能省略：取 $a(x)=e^{-x}$、$f=0$，非零齐次解在正无穷处趋于非零常数。积分因子及这种“遗忘初值”的估计见\hyperref[app:dynamics-factor]{方法卷：积分因子与误差传播}。\end{remark}

\begin{example}[来源序号10，低年级组]\label{ex:fin07-lo:9}\label{pr:fin07-lo:9}

设 $f(x)$ 是定义在 $\mathbb{R}$ 上的连续函数，且满足方程
\[
xf(x)=2\int_{\frac{x}{2}}^{x}f(t)\,\mathrm{d}t+\frac{x^{2}}{4}.
\]
求 $f(x)$。

\par\href{https://github.com/lzxcr/cmc-mathematics/blob/64916fa1fdbc3e444a652eca204d5669058ee321/src/problems/finals/07/10.md}{题目与来源条目}

\end{example}

\begin{solve}\label{ans:fin07-lo:9}

原式把 $f(x)$ 与较小尺度上的平均联系起来。先试仿射函数，再减去已确定的部分：令 $g(x)=f(x)-x-f(0)$，则 $g(0)=0$，且
\[
g(x)=2\int_{1/2}^1g(tx)\,\mathrm{d}t.
\]
固定 $b>0$。若 $M=\max_{[0,b]}|g|>0$，取首次达到 $M$ 的点 $x_*>0$。若 $g(x_*)=\pm M$，连续函数在一段区间上的平均达到其最大值或最小值，只能在整段恒取该极值。因此平均恒等式迫使 $g$ 在 $[x_*/2,x_*]$ 恒等于 $g(x_*)$，与首次达到的定义矛盾。因此 $g=0$ 于 $[0,b]$。对函数 $g(-x)$ 同理，故
\[
f(x)=x+C\qquad(C\in\mathbb R).
\]
直接代回可验证这些函数均满足原方程。

仿射部分可以直接检验：$2\int_{x/2}^x(t+C)\,dt+x^2/4=x^2+Cx=x(x+C)$。减去后对 $x\ne0$ 令积分变量为 $tx$，即得到区间长度为 $1/2$、总权重为 $1$ 的平均公式，$x=0$ 时也成立。若平均值等于最大值，但区间某点低于最大值，连续性就会给出一段正长度的严格低值区间，使平均值严格偏小，矛盾；最小值同理。因此首次取极值点之前的半段也要取相同极值，才能产生所需矛盾。

\end{solve}

\begin{remark}[方法与条件]最大值原理适用于更一般的缩放平均：若连续函数 $g$ 满足 $g(0)=0$、$g(x)=\int g(tx)\,\mathrm d\mu(t)$，其中 $\mu$ 是 $(0,1)$ 上的概率测度，则同一“首次达到最大模”的论证给出 $g\equiv0$。关键在于所有被平均的点都更靠近原点。\end{remark}

\begin{example}[来源序号13，高年级组选做]\label{ex:fin07-hi:11}\label{pr:fin07-hi:11}

实系数多项式 $p(x)$ 的模 1 范数定义为：
\[
\|p\|_{1}:=\int_{0}^{1}\left|p(x)\right|\mathrm{d}x.
\]

（1）求二次实系数多项式 $p(x)$ 使得 $p(x)\leq x^{3}$，对任意 $x\in[0,1]$ 成立，且 $\left\|x^{3}-p(x)\right\|_{1}$ 达到最小。

（2）求三次实系数多项式 $p(x)$ 使得 $p(x)\leq x^{4}$，对任意 $x\in[0,1]$ 成立，且 $\left\|x^{4}-p(x)\right\|_{1}$ 达到最小。

\par\href{https://github.com/lzxcr/cmc-mathematics/blob/64916fa1fdbc3e444a652eca204d5669058ee321/src/problems/finals/07/13.md}{题目与来源条目}

\end{example}

\begin{solve}\label{ans:fin07-hi:11}

约束 $p\leq x^d$ 使误差积分等于 $1/(d+1)-\int_0^1p$，所以只需最大化 $\int p$。正权求积公式把逐点上界转成积分上界；再构造恰在求积节点接触 $x^d$ 的多项式，即可证明最优。

（1）对 $1,x,x^2$ 检查矩条件，得到对任意次数不超过 $2$ 的多项式 $p$，
\[
\int_0^1p(x)\,\mathrm{d}x=\frac14p(0)+\frac34p(2/3)\leq\frac34(2/3)^3=\frac29.
\]
取 $q(x)=x(x-2/3)^2\geq0$，则 $p_*(x)=x^3-q(x)=\frac43x^2-\frac49x$ 达到上界，最小误差为 $1/4-2/9=1/36$。等号要求 $q(0)=q(2/3)=0$；内点 $2/3$ 必为偶重根，因此这一最优多项式唯一。

（2）令 $r_\pm=(1\pm1/\sqrt3)/2$。对 $1,x,x^2,x^3$ 直接检查，得到两点 Gauss 求积公式：
\[
\int_0^1p(x)\,\mathrm{d}x=\frac12p(r_-)+\frac12p(r_+)\leq\frac12r_-^4+\frac12r_+^4.
\]
取 $q(x)=(x^2-x+1/6)^2$，则
\[
p_*(x)=x^4-q(x)=2x^3-\frac43x^2+\frac13x-\frac1{36}
\]
达到上界，最小误差为 $\int_0^1q=1/180$。等号迫使 $r_-,r_+$ 均为 $q$ 的二重根，所以最优解同样唯一。

第一问中权重 $1/4,3/4$ 对单项式 $1,x,x^2$ 分别给出 $1,1/2,1/3$，与积分相同，所以对任意二次以下多项式都精确。第二问中 $r_-+r_+=1$、$r_-r_+=1/6$，因此两节点的二次和为 $2/3$、三次和为 $1/2$，乘以 $1/2$ 后也恰为相应积分。取等时误差多项式在内部节点非负且为零，因而该零点不可能是奇重根；再结合总次数与首项系数为 $1$，就唯一确定了所列因式分解。

\end{solve}

\begin{remark}[方法与条件]最优余项在内部接触点必须有偶重零点，否则会在该点两侧变号；这解释了平方因子的来源。求积公式的正权重则使这些接触点给出全局上界。一般的插值与积分误差观点见\hyperref[app:estimate-interpolation]{方法卷：插值}。\end{remark}

\section{决赛：第 8 届（2017 年）}

本节选入5道课程相关题；来源序号保留原题解编排，未选入题目不重新补号。

\begin{remark}[来源说明]原卷第 1--4 大题必答，原卷第 5--10 大题任选三题。\end{remark}

\begin{example}[来源序号2]\label{ex:fin08-hi:2}\label{pr:fin08-hi:2}

（填空题）设 $a$ 为实数，关于 $x$ 的方程
\[
3x^{4}-8x^{3}-30x^{2}+72x+a=0
\]
有虚根的充分必要条件是 $a$ 满足 $\underline{\qquad\qquad}$。

\par\href{https://github.com/lzxcr/cmc-mathematics/blob/64916fa1fdbc3e444a652eca204d5669058ee321/src/problems/finals/08/02.md}{题目与来源条目}

\end{example}

\begin{solve}\label{ans:fin08-hi:2}

答案为 $a<-37$ 或 $a>27$。记左端多项式为 $P$，则
\[
P'(x)=12(x+2)(x-1)(x-3),\qquad(P(-2),P(1),P(3))=(a-152,a+37,a-27).
\]
由四个单调区间可知，四根均为实数（计重数）当且仅当两处极小值非正、极大值非负，即 $-37\leq a\leq27$。

导数的符号依次为负、正、负、正，所以 $-2,3$ 是极小值点，$1$ 是极大值点，且两端 $P(x)\to+\infty$。若极大值非负而两个极小值非正，每个单调区间都提供一根，端点相切时按重数计数。条件合起来是 $a\le152$、$a\le27$、$a\ge-37$，其中第一条被第二条包含。边界值 $a=-37,27$ 对应实重根，仍没有非实根；因此所求条件必须取开区间 $a<-37$ 或 $a>27$。

\end{solve}

\begin{example}[来源序号3]\label{ex:fin08-hi:3}\label{pr:fin08-hi:3}

（填空题）计算曲面积分
\[
I=\iint_{S}\frac{ax\,\mathrm{d}y\mathrm{d}z+\left(z+a\right)^{2}\mathrm{d}x\mathrm{d}y}{\sqrt{x^{2}+y^{2}+z^{2}}}\quad (a>0 \text{ 为常数}),
\]
其中 $S:z=-\sqrt{a^{2}-x^{2}-y^{2}}$，取上侧，则 $I=$ $\underline{\qquad\qquad}$。

\par\href{https://github.com/lzxcr/cmc-mathematics/blob/64916fa1fdbc3e444a652eca204d5669058ee321/src/problems/finals/08/03.md}{题目与来源条目}

\end{example}

\begin{solve}\label{ans:fin08-hi:3}

答案为 $-\pi a^3/2$。球面上分母恒为 $a$，可先把被积式看作多项式向量场 $F=(ax,0,(z+a)^2)$ 的通量，再除以 $a$。设 $\Omega$ 为下半球体；$S$ 的上侧是其内侧，补上取下侧的赤道圆盘 $D$ 后，闭曲面整体向内。因此
\[
aI+\iint_DF\cdot\nu\,\mathrm dS=-\iiint_\Omega(3a+2z)\,\mathrm dV.
\]
圆盘通量为 $-\pi a^4$，而 $|\Omega|=2\pi a^3/3$、$\int_\Omega z\,\mathrm dV=-\pi a^4/4$。代入得 $aI=-\pi a^4/2$。

圆盘 $D$ 位于 $z=0$，向内法向量为 $(0,0,-1)$，所以 $F\cdot\nu=-a^2$，乘面积 $\pi a^2$ 得 $-\pi a^4$。球体的水平截面积为 $\pi(a^2-z^2)$，从而
\[
\int_\Omega z\,dV=\pi\int_{-a}^0z(a^2-z^2)\,dz=-\frac{\pi a^4}{4}.
\]
散度积分是 $3a(2\pi a^3/3)+2(-\pi a^4/4)=3\pi a^4/2$，取向内负号后再减去圆盘通量，得到 $aI=-\pi a^4/2$。

\end{solve}

\begin{example}[来源序号7]\label{ex:fin08-hi:7}\label{pr:fin08-hi:7}

对 $\mathbb{R}$ 上无穷次可微的（复值）函数 $\varphi(x)$，称 $\varphi\in\mathcal{S}$，如果 $\forall m,k\geq 0$ 成立 $\sup_{x\in\mathbb{R}}\left|x^{m}\varphi^{(k)}(x)\right|<+\infty$。若 $f\in\mathcal{S}$，可定义 $\hat{f}(x)=\int_{R}f(y)e^{-2\pi ixy}\,\mathrm{d}y$（$\forall x\in\mathbb{R}$）。证明：$\hat{f}\in\mathcal{S}$，且 $f(x)=\int_{R}\hat{f}(y)e^{2\pi ixy}\,\mathrm{d}y$（$\forall x\in\mathbb{R}$）。

\par\href{https://github.com/lzxcr/cmc-mathematics/blob/64916fa1fdbc3e444a652eca204d5669058ee321/src/problems/finals/08/07.md}{题目与来源条目}

\end{example}

\begin{solve}\label{ans:fin08-hi:7}

先用微分与乘法的对偶性证明速降，再用 Gaussian 近似恒等算子证明反演。

积分号下求导与分部积分均由速降性保证，对任意非负整数 $m,k$，
\[
\xi^m\widehat f^{(k)}(\xi)=\frac1{(2\pi i)^m}\int_{\mathbb R}
\frac{\mathrm d^m}{\mathrm dy^m}\bigl[(-2\pi iy)^kf(y)\bigr]e^{-2\pi i\xi y}\,\mathrm dy.
\]
右端的被积系数是有限个“多项式乘 $f$ 的导数”之和，仍然速降，故其 $L^1$ 范数有限且与 $\xi$ 无关。因此 $\widehat f\in\mathcal S$。

速降如何保证这些运算，可作如下估计。任取非负整数 $k$，由定义可使 $|y^kf(y)|$ 在 $|y|\ge1$ 时被 $C|y|^{-2}$ 控制，故可积；这为每次对频率求导提供可积控制函数。分部积分涉及的各阶导数同样满足这一尾界，边界项趋于零。要证明频率导数连续，也使用相同控制函数和控制收敛定理，因而不仅存在各阶导数，而且它们连续。

令 $\varepsilon>0$。Gaussian 积分恒等式
\[
\int_{\mathbb R}e^{-\pi\varepsilon\xi^2}e^{2\pi it\xi}\,\mathrm d\xi
=\varepsilon^{-1/2}e^{-\pi t^2/\varepsilon}
\]
可由对 $t$ 求导、分部积分及 $t=0$ 时的 Gaussian 积分直接得到。由 Fubini 定理，
\[
\int_{\mathbb R}\widehat f(\xi)e^{-\pi\varepsilon\xi^2}e^{2\pi ix\xi}\,\mathrm d\xi
=\int_{\mathbb R}f(x-\sqrt\varepsilon u)e^{-\pi u^2}\,\mathrm du.
\]
令 $\varepsilon\downarrow0$：左端由 $\widehat f\in L^1$ 使用控制收敛，右端由 $f$ 连续有界及 $\int e^{-\pi u^2}=1$ 使用控制收敛，即得
\[
f(x)=\int_{\mathbb R}\widehat f(\xi)e^{2\pi ix\xi}\,\mathrm d\xi.
\]

Gaussian 恒等式也可详细推出。记其左端为 $K_\varepsilon(t)$，对 $t$ 求导，再利用 $(e^{-\pi\varepsilon\xi^2})'=-2\pi\varepsilon\xi e^{-\pi\varepsilon\xi^2}$ 分部积分，得 $K_\varepsilon'(t)=-(2\pi t/\varepsilon)K_\varepsilon(t)$。结合 $K_\varepsilon(0)=\varepsilon^{-1/2}$，解此一阶方程即得所列核。引入该核后的二重积分绝对值不超过 $\|f\|_1\int e^{-\pi\varepsilon\xi^2}d\xi<\infty$，故 Fubini 换序合法。

\end{solve}

\begin{remark}[方法与条件]Gaussian 因子的作用是让交换积分始终发生在绝对可积的范围内；交换后，它又成为总质量为 $1$、逐渐集中到原点的核。因此反演由“频率端截去远处”转成“空间端局部平均”。这一联系见\hyperref[app:integration-fourier]{方法卷：Fourier 变换与正则化}。\end{remark}

\begin{example}[来源序号8，低年级组]\label{ex:fin08-lo:8}\label{pr:fin08-lo:8}

设 $n>1$ 为正整数，令
\[
S_{n}=\left(\frac{1}{n}\right)^{n}+\left(\frac{2}{n}\right)^{n}+\cdots+\left(\frac{n-1}{n}\right)^{n}.
\]
（1）证明：数列 $S_{n}$ 单调增加且有界，从而极限 $\lim_{n\to\infty}S_{n}$ 存在；
（2）求极限 $\lim_{n\to\infty}S_{n}$。

\par\href{https://github.com/lzxcr/cmc-mathematics/blob/64916fa1fdbc3e444a652eca204d5669058ee321/src/problems/finals/08/08.md}{题目与来源条目}

\end{example}

\begin{solve}\label{ans:fin08-lo:8}

幂次 $n$ 很大时，远离 $1$ 的项迅速衰减，主要贡献来自右端附近。因此把指标改成距右端的网格数，定义 $b_{n,k}=(1-k/n)^n$（$1\leq k<n$），其余 $k\geq n$ 令 $b_{n,k}=0$，则 $S_n=\sum_{k\geq1}b_{n,k}$。

固定 $k$，对 $n>k$，算术---几何平均不等式给出
\[
1-\frac{k}{n+1}=\frac{n(1-k/n)+1}{n+1}>\left(1-\frac{k}{n}\right)^{n/(n+1)}.
\]
所以 $b_{n+1,k}>b_{n,k}$，从而 $S_n$ 严格递增。同时
\[
0\leq b_{n,k}\leq e^{-k},\qquad b_{n,k}\longrightarrow e^{-k}.
\]
可和上界既证明有界性，也允许对计数测度使用控制收敛定理，得到
\[
\lim_{n\to\infty}S_n=\sum_{k=1}^\infty e^{-k}=\frac1{e-1}.
\]

也可以用有限头部加一致尾部避免抽象的控制收敛表述。给定 $\varepsilon>0$，先取 $K$ 使 $\sum_{k>K}e^{-k}<\varepsilon$；对于有限个 $k\le K$，$b_{n,k}\to e^{-k}$ 可以逐项取极限，其和的误差最终小于 $\varepsilon$。剩余两条尾和都被 $\varepsilon$ 控制，因此整体误差最终小于 $3\varepsilon$。特别是原和中靠近 $k=n$ 的项数在变化，只有这个统一尾界才允许把极限移入无限和。

\end{solve}

\begin{remark}[方法与条件]将幂次换成 $an$（$a>0$），同样以 $e^{-ak}$ 控制各项，极限为 $1/(e^a-1)$。这里有效区域的宽度与网格宽度都是 $1/n$，所以保留下来的是离散几何级数，而不是通常的区间积分。交换无穷和与极限的依据见\hyperref[app:integration-exchange]{方法卷：积分与极限交换}。\end{remark}

\begin{example}[来源序号10，低年级组]\label{ex:fin08-lo:9}\label{pr:fin08-lo:9}

求证：常微分方程 $\frac{\mathrm{d}y}{\mathrm{d}x}=-y^{3}+\sin x$，$x\in[0,2\pi]$ 有唯一的满足 $y(0)=y(2\pi)$ 的解。

\par\href{https://github.com/lzxcr/cmc-mathematics/blob/64916fa1fdbc3e444a652eca204d5669058ee321/src/problems/finals/08/10.md}{题目与来源条目}

\end{example}

\begin{solve}\label{ans:fin08-lo:9}

周期边界条件要求初值经过 $2\pi$ 的演化回到自身，因此研究返回映射的不动点。区间 $[-2,2]$ 在正向流下不变：在上、下端点分别有 $-y^3+\sin x<0$、$>0$。因此初值 $r\in[-2,2]$ 对应的解存在于 $[0,2\pi]$，且 $P(r)=y(2\pi;r)$ 连续并把 $[-2,2]$ 映入自身。于是 $P(-2)+2\geq0$、$P(2)-2\leq0$，对连续函数 $P(r)-r$ 应用介值定理，得到不动点。

若 $y,z$ 都满足周期边界条件，则差 $w=y-z$ 满足
\[
w'=-(y^2+yz+z^2)w,\qquad
w(2\pi)=w(0)\exp\left[-\int_0^{2\pi}(y^2+yz+z^2)\,\mathrm dx\right].
\]
若 $w(0)\ne0$，则初始时 $y,z$ 不同时为零；连续性及 $y^2+yz+z^2\geq\frac12(y^2+z^2)$ 保证该积分严格为正，指数因子严格小于 $1$，与 $w(2\pi)=w(0)$ 矛盾。因此解唯一。

右端 $-y^3+\sin x$ 连续，并对 $y$ 局部 Lipschitz，因为其偏导数为 $-3y^2$。因此初值问题局部存在唯一，且连续依赖初值；不变区间又保证解保持有界，按延拓定理可覆盖整个 $[0,2\pi]$。唯一性部分的积分因子公式同时表明：若两解某时相等，则差恒为零；若初值不同，差始终不为零，且经过一周期后其绝对值严格缩小，无法同时满足相同的周期边界条件。

\end{solve}

\begin{remark}[方法与条件]存在性只需返回映射把一个闭区间映入自身；唯一性则来自两解之间距离在一周期内严格缩小，并不需要先得到一个统一的压缩常数。两种论证各用到不同的结构，见\hyperref[app:dynamics-flow]{方法卷：流映射与周期解}。\end{remark}

\section{决赛：第 9 届（2018 年）}

本节选入4道课程相关题；来源序号保留原题解编排，未选入题目不重新补号。

\begin{remark}[来源说明]原卷第 1--4 大题必答，原卷第 5--10 大题任选三题。\end{remark}

\begin{example}[来源序号2]\label{ex:fin09-hi:2}\label{pr:fin09-hi:2}

（填空题）$\lim_{x \to 0} \frac{\ln(1+\tan x)-\ln(1+\sin x)}{x^{3}}=$ $\underline{\qquad\qquad}$。

\par\href{https://github.com/lzxcr/cmc-mathematics/blob/64916fa1fdbc3e444a652eca204d5669058ee321/src/problems/finals/09/02.md}{题目与来源条目}

\end{example}

\begin{solve}\label{ans:fin09-hi:2}

答案为 $1/2$。由对数的中值定理，分子为 $(\tan x-\sin x)/(1+\theta_x)$，其中 $\theta_x\to0$；又 $\tan x-\sin x=\sin x(1-\cos x)/\cos x\sim x^3/2$。

这里中值定理作用于函数 $u\mapsto\ln(1+u)$，两个端点是 $\sin x$ 与 $\tan x$；在 $x$ 足够接近零时二者均大于 $-1$，所以导数 $1/(1+u)$ 有定义。由夹逼 $\theta_x\to0$，再逐因子计算
\[
\frac{\ln(1+\tan x)-\ln(1+\sin x)}{x^3}
=\frac1{(1+\theta_x)\cos x}\frac{\sin x}{x}\frac{1-\cos x}{x^2}
\longrightarrow1\cdot1\cdot\frac12.
\]
这避免了在两个相近对数相减时直接作不合法的一阶等价替换。

\end{solve}

\begin{example}[来源序号3]\label{ex:fin09-hi:3}\label{pr:fin09-hi:3}

（填空题）设 $\Gamma$ 为空间曲线
\[
\begin{cases} x=\pi \sin\left(t/2\right), \\ y=t-\sin t, \\ z=\sin 2t \end{cases}
\]
从 $t=0$ 到 $t=\pi$ 的一段，则第二型曲线积分
\[
\int_{\Gamma} e^{\sin x}\left(\cos x \cos y\, \mathrm{d} x-\sin y\, \mathrm{d} y\right)+\cos z\, \mathrm{d} z
\]
为 $\underline{\qquad\qquad}$。

\par\href{https://github.com/lzxcr/cmc-mathematics/blob/64916fa1fdbc3e444a652eca204d5669058ee321/src/problems/finals/09/03.md}{题目与来源条目}

\end{example}

\begin{solve}\label{ans:fin09-hi:3}

答案为 $-2$。被积微分形式为 $\mathrm d(e^{\sin x}\cos y+\sin z)$，起、终点分别为 $(0,0,0)$ 与 $(\pi,\pi,0)$，代入势函数之差即可。

令 $U(x,y,z)=e^{\sin x}\cos y+\sin z$，直接求偏导得 $U_x=e^{\sin x}\cos x\cos y$、$U_y=-e^{\sin x}\sin y$、$U_z=\cos z$，与三个微分系数逐项一致。因此沿参数曲线有 $dU/dt=U_xx'+U_yy'+U_zz'$，积分只取端点差。起点的势为 $1$，终点的势为 $-1$，所以积分为 $-1-1=-2$，也与给定方向一致。

\end{solve}

\begin{example}[来源序号7]\label{ex:fin09-hi:7}\label{pr:fin09-hi:7}

设 $f(x)$ 在 $[0, 1]$ 上有二阶连续导函数，且 $f(0)f(1) \geq 0$。求证：
\[
\int_0^1 \left|f'(x)\right| \mathrm{d}x \leq 2 \int_0^1 |f(x)| \mathrm{d}x + \int_0^1 \left|f''(x)\right| \mathrm{d}x.
\]

\par\href{https://github.com/lzxcr/cmc-mathematics/blob/64916fa1fdbc3e444a652eca204d5669058ee321/src/problems/finals/09/07.md}{题目与来源条目}

\end{example}

\begin{solve}\label{ans:fin09-hi:7}

二阶导数的积分控制一阶导数的振幅；剩下只需用 $f$ 的积分控制一阶导数的最低水平。令 $m=\min_{[0,1]}|f'|$。由微积分基本定理，
\[
\int_0^1|f'|\leq\max|f'|\leq m+\int_0^1|f''|.
\]
只需证明 $m\leq2\int_0^1|f|$。若 $m=0$，无需再证。若 $m>0$，则 $f$ 严格单调，端点同号保证 $f$ 在整个区间不变号。通过替换 $f$ 为 $-f$ 或 $f(1-x)$，可归约为 $f\geq0,f'\geq m$。于是 $f(x)\geq mx$，积分得到 $\int_0^1|f|\geq m/2$。

若 $x_0$ 使 $|f'(x_0)|=m$，则任意 $x$ 都有 $|f'(x)|\le m+|\int_{x_0}^xf''(t)dt|\le m+\int_0^1|f''|$，从而得到最开始的振幅界。$m>0$ 时连续的 $f'$ 不取零，故符号恒定；严格单调函数的值介于端点值之间，而两端同号或有一端为零，便保证 $f$ 不变号。取反号和反向变量都保持题中三个绝对值积分不变，所以可以归约到 $f\ge0$ 且 $f'\ge m$ 后积分。

\end{solve}

\begin{remark}[方法与条件]端点同号条件确实参与了证明：取 $f(x)=x-1/2$，则左端为 $1$、右端为 $1/2$，不等式不成立。常数 $2$ 也是最佳的，因为 $f(x)=x$ 时等号成立。\end{remark}

\begin{example}[来源序号10，低年级组]\label{ex:fin09-lo:9}\label{pr:fin09-lo:9}

有界连续函数 $g(t): \mathbb{R} \rightarrow \mathbb{R}$ 满足 $1 < g(t) < 2$。$x(t), t \in \mathbb{R}$ 是方程 $\ddot{x}(t) = g(t)x$ 的单调正解。求证：存在常数 $C_2 > C_1 > 0$ 满足
\[
C_1 x(t) < |\dot{x}(t)| < C_2 x(t), \quad t \in \mathbb{R}.
\]

\par\href{https://github.com/lzxcr/cmc-mathematics/blob/64916fa1fdbc3e444a652eca204d5669058ee321/src/problems/finals/09/10.md}{题目与来源条目}

\end{example}

\begin{solve}\label{ans:fin09-lo:9}

反转时间可归约为 $x$ 单调递增。由于 $x''=gx>0$，有 $x'>0$。当 $t\to-\infty$ 时，$x,x'$ 均有非负极限，而且二者都为零：$x'$ 的极限若为正，向左积分会使 $x<0$；$x$ 的极限若为正，则 $x''>x$ 使 $x'$ 向左变负。

比较 $x'^{\,2}$ 与 $x^2$ 时，它们的导数都含 $xx'$，所以将方程乘以 $x'>0$。从负无穷积分，得到
\[
\frac{x'(t)^2}{2}=\int_{-\infty}^tg(s)x(s)x'(s)\,\mathrm ds.
\]
由 $1<g<2$ 及 $x(-\infty)=0$，
\[
\frac{x(t)^2}{2}<\frac{x'(t)^2}{2}<x(t)^2.
\]
所以可取 $C_1=1,C_2=\sqrt2$，且两个不等号都严格。

左端极限为零可以量化：若 $x'(-\infty)=L>0$，则 $x(t)\le x(0)+Lt$ 对 $t<0$ 成立，足够负时违背正性；若 $x(-\infty)=\ell>0$，则 $x''(t)>\ell$，向左积分同样使 $x'(t)$ 最终为负。先在 $[T,t]$ 上积分 $x'x''=gxx'$，再令 $T\to-\infty$，左端趋于 $x'(t)^2/2$。由于 $xx'>0$，且 $g-1$、$2-g$ 连续严格为正，在任一非空有限子区间的积分均为正，所以两端比较都是严格的。

\end{solve}

\begin{remark}[方法与条件]若 $0<a\leq g\leq b$，同一能量比较给出 $\sqrt a\,x\leq|x'|\leq\sqrt b\,x$；常系数指数解表明通常不能改为严格不等号。等价地，正解的对数导数 $x'/x$ 被夹在两个常数之间。能量与对数导数的联系见\hyperref[app:dynamics-energy]{方法卷：能量方法}。\end{remark}

\section{决赛：第 10 届（2019 年）}

本节选入6道课程相关题；来源序号保留原题解编排，未选入题目不重新补号。

\begin{remark}[来源说明]原卷第 1--4 大题必答，原卷第 5--10 大题任选三题。\end{remark}

\begin{example}[来源序号2]\label{ex:fin10-hi:2}\label{pr:fin10-hi:2}

（填空题）设 $y(x) \in C^{1}[0,1)$ 满足 $y(x) \in [0,\pi]$ 及
\[
x = \begin{cases} \frac{\sin y(x)}{y(x)}, & y \in (0,\pi], \\ 1, & y=0, \end{cases}
\]
则 $y'(0) = $ $\underline{\qquad\qquad}$。

\par\href{https://github.com/lzxcr/cmc-mathematics/blob/64916fa1fdbc3e444a652eca204d5669058ee321/src/problems/finals/10/02.md}{题目与来源条目}

\end{example}

\begin{solve}\label{ans:fin10-hi:2}

答案为 $-\pi$。$x=0$ 时 $y=\pi$，且 $\left.(\sin y/y)'\right|_{y=\pi}=-1/\pi$，由反函数求导公式可得。

在 $0\le y\le\pi$ 中，$\sin y/y$ 的连续延拓值在 $y=0$ 是 $1$，而唯一使其为零的点是 $y=\pi$，所以 $y(0)=\pi$。对 $x=\sin y(x)/y(x)$ 求导，得到
\[
1=\frac{y\cos y-\sin y}{y^2}\,y'(x).
\]
代入 $x=0,y=\pi$ 得 $1=(-1/\pi)y'(0)$，即 $y'(0)=-\pi$。这里 $x=0$ 是定义区间端点，所求是与 $C^1[0,1)$ 约定一致的右导数。

\end{solve}

\begin{example}[来源序号3]\label{ex:fin10-hi:3}\label{pr:fin10-hi:3}

（填空题）设 $f(x) = \int_{x}^{+\infty} e^{-t^{2}} \mathrm{d} t$，则 $\int_{0}^{+\infty} xf(x) \mathrm{d} x = $ $\underline{\qquad\qquad}$。

\par\href{https://github.com/lzxcr/cmc-mathematics/blob/64916fa1fdbc3e444a652eca204d5669058ee321/src/problems/finals/10/03.md}{题目与来源条目}

\end{example}

\begin{solve}\label{ans:fin10-hi:3}

答案为 $\sqrt\pi/8$。Tonelli 定理给出
\[
\int_0^\infty x\int_x^\infty e^{-t^2}\,\mathrm dt\,\mathrm dx
=\frac12\int_0^\infty t^2e^{-t^2}\,\mathrm dt=\frac{\sqrt\pi}{8}.
\]

积分区域原来写成 $0\le x\le t<\infty$，换序后固定 $t$，$x$ 从 $0$ 积到 $t$，故 $\int_0^t x\,dx=t^2/2$。再利用 $d(e^{-t^2})=-2te^{-t^2}dt$，
\[
\int_0^\infty t^2e^{-t^2}dt
=-\frac12[te^{-t^2}]_0^\infty+\frac12\int_0^\infty e^{-t^2}dt=\frac{\sqrt\pi}{4}.
\]
端点项为零，外面的 $1/2$ 再给出 $\sqrt\pi/8$。非负性允许先换序，而算出的有限结果确认原积分收敛。

\end{solve}

\begin{example}[来源序号7]\label{ex:fin10-hi:7}\label{pr:fin10-hi:7}

（满分 20 分）设 $\alpha > 0$，$f(x) \in C^{1}[0,1]$，且对任何非负整数 $n$，$f^{(n)}(0)$ 均存在且为零。进一步，存在常数 $C > 0$ 使得
\[
\left| x^{\alpha} f'(x) \right| \leq C | f(x) | \quad (\forall x \in [0,1]).
\]
证明：

（1）若 $\alpha = 1$，则在 $[0,1]$ 上 $f(x) \equiv 0$；

（2）若 $\alpha > 1$，举例说明在 $[0,1]$ 上 $f(x) \equiv 0$ 可以不成立。

\par\href{https://github.com/lzxcr/cmc-mathematics/blob/64916fa1fdbc3e444a652eca204d5669058ee321/src/problems/finals/10/07.md}{题目与来源条目}

\end{example}

\begin{solve}\label{ans:fin10-hi:7}

（1）不直接除以 $f$，以免零点造成缺口；改对 $f^2$ 使用积分因子。取整数 $N>C$。零点处的平坦性及 Peano 形式的 Taylor 公式给出 $f(x)=o(x^N)$，故 $x^{-2C}f(x)^2\to0$。而
\[
\frac{\mathrm d}{\mathrm dx}\bigl(x^{-2C}f(x)^2\bigr)
=2x^{-2C-1}\bigl(xff'-Cf^2\bigr)\leq0.
\]
这个非负函数单调不增，且左端极限为零，因此恒为零，即 $f\equiv0$。

（2）取
\[
f(x)=\begin{cases}\exp(-x^{1-\alpha}),&x>0,\\0,&x=0.\end{cases}
\]
它在零点无穷阶平坦，且 $x^\alpha f'(x)=(\alpha-1)f(x)$，满足全部条件，但不恒为零。

第二问设 $\beta=\alpha-1>0$，则 $f(x)=e^{-x^{-\beta}}$。在 $x>0$ 每求导一次，都得到有限个“负幂 $x^{-r}$ 乘 $e^{-x^{-\beta}}$”的项；令 $u=x^{-\beta}$ 后，其绝对值形如 $u^{r/\beta}e^{-u}\to0$。这些项除以 $x$ 后仍趋零，故可逐阶用导数定义把各阶导数连续延伸到零点，并证明其值全为零。又直接求导得 $f'=\beta x^{-\alpha}f$，所以常数取 $C=\alpha-1$ 即满足要求。

\end{solve}

\begin{remark}[方法与条件]临界性来自积分因子：$\int_x^1t^{-\alpha}\,\mathrm dt$ 在 $\alpha=1$ 时按 $|\log x|$ 增长，只容许幂次量级的非零衰减；$\alpha>1$ 时按 $x^{1-\alpha}$ 增长，容许平坦的指数解。推导见\hyperref[app:dynamics-factor]{方法卷：积分因子}。\end{remark}

\begin{example}[来源序号8，低年级组]\label{ex:fin10-lo:8}\label{pr:fin10-lo:8}

（满分 15 分）设 $c\in(0,1)$，$x_{1}\in(0,1)$ 且 $x_{1}\neq c\left(1-x_{1}^{2}\right)$，$x_{n+1}=c\left(1-x_{n}^{2}\right)\ (n\geq1)$。证明：$\left\{x_{n}\right\}$ 收敛当且仅当 $c\in\left(0,\frac{\sqrt{3}}{2}\right]$。

\par\href{https://github.com/lzxcr/cmc-mathematics/blob/64916fa1fdbc3e444a652eca204d5669058ee321/src/problems/finals/10/08.md}{题目与来源条目}

\end{example}

\begin{solve}\label{ans:fin10-lo:8}

令 $F(x)=c(1-x^2)$，则 $F$ 把 $(0,1)$ 映入自身且严格递减，$F^2$ 递增。比较 $x_3$ 与 $x_1$ 后，$F^2$ 的递增性会把这一大小关系逐步传递到所有奇数项；偶数项因 $F$ 递减而沿相反方向单调。因此两个子列都单调有界，设其极限为 $u,v$，则
\[
u=c(1-v^2),\qquad v=c(1-u^2).
\]
若 $u\ne v$，相减得 $u+v=1/c$，故 $u,v$ 是
\[
ct^2-t+1/c-c=0
\]
的两个不同实根，要求判别式 $4c^2-3>0$。所以 $0<c\leq\sqrt3/2$ 时必有 $u=v$，整个数列收敛。

反之，唯一不动点为 $r=(\sqrt{1+4c^2}-1)/(2c)$。当 $c>\sqrt3/2$ 时，$|F'(r)|=\sqrt{1+4c^2}-1>1$，故在 $r$ 的某邻域中，$|F(x)-r|>|x-r|$。初值不是不动点，且 $F$ 单射保证迭代永不落到 $r$，所以数列不可能趋于 $r$，因而不收敛。

结合题设 $c\in(0,1)$，所求条件即 $c\in(0,\sqrt3/2]$。

必要性还可用精确差式说明排斥性：$F(x)-F(r)=-c(x+r)(x-r)$，所以 $|F(x)-r|/|x-r|=c(x+r)\to2cr>1$ 当 $x\to r$。因此存在 $q>1$ 和邻域，使其中非零误差每步至少乘以 $q$。若数列趋于 $r$，最终永远留在该邻域，误差反而至少几何增长，矛盾。在临界值 $c=\sqrt3/2$ 时导数绝对值等于 $1$，上述排斥论证不适用，但前面的二周期判别式为零仍保证两个子列极限相同。

\end{solve}

\begin{remark}[方法与条件]超过临界值时，两子列的不同极限就是
\[
u,v=\frac{1\pm\sqrt{4c^2-3}}{2c},
\]
它们构成二周期。临界值处 $F'(r)=-1$，一阶局部判据不作结论，但二周期方程仍能决定全局收敛。参见\hyperref[app:dynamics-iteration]{方法卷：递减映射与二周期}。\end{remark}

\begin{example}[来源序号10，低年级组]\label{ex:fin10-lo:9}\label{pr:fin10-lo:9}

（满分 15 分）已知 $a,b,c\in C(\mathbb{R},\mathbb{R})$，方程
\[
\frac{\mathrm{d}y}{\mathrm{d}x}=a(x)y^{2}+b(x)y+c(x)
\]
只有有限个 $2\pi$ 周期解。求它的 $2\pi$ 周期解个数的最大值。

\par\href{https://github.com/lzxcr/cmc-mathematics/blob/64916fa1fdbc3e444a652eca204d5669058ee321/src/problems/finals/10/10.md}{题目与来源条目}

\end{example}

\begin{solve}\label{ans:fin10-lo:9}

最大值为 $2$。例如 $y'=y(y+1)$ 仅有两个周期解 $y\equiv0,-1$，因为其他解严格单调。

证明上界时，把两个已知解变成常数解 $0,1$，二次方程便化为可分离变量的形式。若已有不同周期解 $y_1,y_2$，由初值唯一性，$d=y_1-y_2$ 处处非零。变换 $y=y_2+dz$ 将方程化为
\[
z'=a(x)d(x)z(z-1).
\]
若还存在第三个周期解 $z_*$，则 $z_*$ 处处不取 $0,1$，从而
\[
F(x):=\int_0^xa(t)d(t)\,\mathrm dt
=\log\left|\frac{z_*(x)-1}{z_*(x)}\right|-\log\left|\frac{z_*(0)-1}{z_*(0)}\right|
\]
是周期函数。此时对每个 $C<0$，
\[
z_C(x)=\frac1{1-Ce^{F(x)}}
\]
都是全局周期解，且分母始终为正，产生无穷多个不同的 $y$，与有限性矛盾。故有限时至多两个。

把 $y=y_2+dz$ 求导，并减去 $y_2'=ay_2^2+by_2+c$，再用 $d'=[a(y_1+y_2)+b]d$，线性项抵消后就得到 $z'=ad(z^2-z)$。此外
\[
\frac{d}{dx}\log\left|\frac{z-1}{z}\right|=\frac{z'}{z(z-1)}=ad.
\]
若第三个解周期，则对数表达式也是周期的，因此 $F$ 周期，无需另行假设原方程系数周期。选择 $C<0$ 使 $1-Ce^F>1$，保证构造的解处处有定义；不同 $C$ 给出不同初值，故确为无穷多个不同周期解。

\end{solve}

\begin{remark}[方法与条件]Riccati 方程通过分式线性变换与线性系统相联系。两个已知解固定了“坐标轴”，第三个周期解又迫使时间因子周期，从而产生整族周期解；这比单纯计数零点更能解释上界 $2$ 的来源。返回映射观点见\hyperref[app:dynamics-flow]{方法卷：流映射与周期解}。\end{remark}

\begin{example}[来源序号13，高年级组选做]\label{ex:fin10-hi:11}\label{pr:fin10-hi:11}

（满分 10 分）设 $\delta: a = x_0 < x_1 < \dots < x_n = b$ 是 $[a, b]$ 区间的一个剖分。用 $S[a, b]$ 表示满足下列条件的分片实系数多项式全体构成的集合：对任意 $s(x) \in S[a, b]$，

（1）$s(x) \mid_{[x_i, x_{i+1}]}$ 是三次多项式，$i = 0, 1, \dots, n-1$；

（2）$s(x)$ 在区间 $[a, b]$ 上二阶连续可导。

证明：

1) 对区间 $[a,b]$ 的任意实函数 $f(x)$，存在唯一的 $s(x)\in S[a,b]$ 满足
\[
s\left(x_{i}\right)=f\left(x_{i}\right),\ i=0,1,\dots,n, \qquad s^{\prime\prime}(a)=s^{\prime\prime}(b)=0;
\]

2) 设 $f\in C^{2}[a,b]$，则对满足 1) 的函数 $s(x)$ 有
\[
\int_{a}^{b}\left(s''(x)\right)^{2}\mathrm{d}x\leq \int_{a}^{b}\left(f''(x)\right)^{2}\mathrm{d}x,
\]
且等号成立当且仅当 $s(x)=f(x)$。

\par\href{https://github.com/lzxcr/cmc-mathematics/blob/64916fa1fdbc3e444a652eca204d5669058ee321/src/problems/finals/10/13.md}{题目与来源条目}

\end{example}

\begin{solve}\label{ans:fin10-hi:11}

（1）三次多项式的二阶导数为一次函数，所以以节点处二阶导数为未知量，拼接条件最简单。记 $h_i=x_{i+1}-x_i$、$f_i=f(x_i)$、$M_i=s''(x_i)$，自然边界条件为 $M_0=M_n=0$。各小区间上的样条必为
\[
\begin{aligned}
s_i(x)={}&\frac{M_i(x_{i+1}-x)^3+M_{i+1}(x-x_i)^3}{6h_i}\\
&+\left(f_i-\frac{M_ih_i^2}{6}\right)\frac{x_{i+1}-x}{h_i}
+\left(f_{i+1}-\frac{M_{i+1}h_i^2}{6}\right)\frac{x-x_i}{h_i}.
\end{aligned}
\]
这已保证函数值和二阶导数连续；一阶导数连续恰等价于
\[
h_{i-1}M_{i-1}+2(h_{i-1}+h_i)M_i+h_iM_{i+1}
=6\left(\frac{f_{i+1}-f_i}{h_i}-\frac{f_i-f_{i-1}}{h_{i-1}}\right)
\]
（$1\leq i<n$）。系数矩阵严格对角占优，故可逆，得到唯一的 $M_i$ 和 $s$。当 $n=1$ 时直接得到端点间的直线。

（2）令 $g=f-s$，则所有节点上 $g=0$。逐段分部积分两次，由 $s_i^{(3)}$ 为常数、$s''$ 连续以及自然边界条件，
\[
\int_a^bs''g''=\sum_i[s_i''g']_{x_i}^{x_{i+1}}
-\sum_i s_i^{(3)}[g]_{x_i}^{x_{i+1}}=0.
\]
故有精确的能量分解
\[
\int_a^b(f'')^2=\int_a^b(s'')^2+\int_a^b(g'')^2.
\]
等号要求 $g''=0$，而线性函数 $g$ 在两端点都为零，故 $g=0$。

这里分片三次多项式按样条的通常约定包含次数低于三次的情形。矩阵可逆性可直接验证：若齐次系统有非零向量，取绝对值最大的分量 $M_i$，则该行给出 $2(h_{i-1}+h_i)|M_i|\le(h_{i-1}+h_i)|M_i|$，矛盾。能量证明中的边界项在内部节点两两相消，是因为 $s''$ 与 $g'$ 都连续；外端由 $s''(a)=s''(b)=0$ 消失，含 $[g]$ 的项则因每个节点 $g=0$ 消失。于是交叉项为零，展开 $(s''+g'')^2$ 就得到精确分解。

\end{solve}

\begin{remark}[方法与条件]能量恒等式表明，样条是仿射插值空间中关于二阶导数内积的正交极小元。竞争函数无需满足自然边界条件；端点处 $s''=0$ 正是消去边界项的条件。投影方法及能量分解见\hyperref[app:integration-hilbert]{方法卷：Hilbert 空间方法}。\end{remark}

\section{决赛：第 11 届（2021 年）}

本节选入4道课程相关题；来源序号保留原题解编排，未选入题目不重新补号。

\begin{remark}[来源说明]原卷第 1--4 大题必答，原卷第 5--10 大题任选三题。\end{remark}

\begin{example}[来源序号1]\label{ex:fin11-high:1}\label{pr:fin11-high:1}

（填空题）$\lim_{n \to +\infty} \left( \prod_{k=n}^{3n-1} k \right)^{\frac{1}{2n}} \sin \frac{1}{n} = $ $\underline{\qquad\qquad}$。

\par\href{https://github.com/lzxcr/cmc-mathematics/blob/64916fa1fdbc3e444a652eca204d5669058ee321/src/problems/finals/11/01.md}{题目与来源条目}

\end{example}

\begin{solve}\label{ans:fin11-high:1}

答案为 $3\sqrt3/e$。取对数并分离 $\log n$，
\[
\frac1{2n}\sum_{k=n}^{3n-1}\log(k/n)\longrightarrow\frac12\int_1^3\log t\,\mathrm dt=\frac32\log3-1,
\]
再用 $n\sin(1/n)\to1$ 即得。

乘积共有 $2n$ 项，逐项提出 $n$ 后，原式成为
\[
n\sin(1/n)\exp\left\{\frac1{2n}\sum_{k=n}^{3n-1}\ln(k/n)\right\}.
\]
网格步长为 $1/n$，所以上述指数是 $[1,3]$ 上左端点黎曼和的一半。由 $\int_1^3\ln t\,dt=[t\ln t-t]_1^3=3\ln3-2$，指数极限为 $\tfrac32\ln3-1$；取指数得到 $3^{3/2}/e$，再乘趋于 $1$ 的正弦因子即为答案。

\end{solve}

\begin{example}[来源序号2]\label{ex:fin11-high:2}\label{pr:fin11-high:2}

（填空题）已知 $f$ 在区间 $(-1,3)$ 内有二阶连续导数，$f(0) = 12, f(2) = 2f'(2) + 8$，则 $\int_{0}^{1} xf''(2x) \mathrm{d} x = $ $\underline{\qquad\qquad}$。

\par\href{https://github.com/lzxcr/cmc-mathematics/blob/64916fa1fdbc3e444a652eca204d5669058ee321/src/problems/finals/11/02.md}{题目与来源条目}

\end{example}

\begin{solve}\label{ans:fin11-high:2}

答案为 $1$。换元 $t=2x$ 并分部积分，
\[
\int_0^1xf''(2x)\,\mathrm dx=\frac14[2f'(2)-f(2)+f(0)]=1.
\]

换元后 $x=t/2,dx=dt/2$，故两个因子合起来是 $1/4$：$\int_0^1xf''(2x)dx=\tfrac14\int_0^2tf''(t)dt$。分部积分给出 $\int_0^2tf''=[tf']_0^2-\int_0^2f'=2f'(2)-f(2)+f(0)$。题设关系恰好消去未知端点导数：$2f'(2)-f(2)=-8$，加上 $f(0)=12$ 得 $4$，再除以 $4$ 即得 $1$。

\end{solve}

\begin{example}[来源序号7]\label{ex:fin11-high:7}\label{pr:fin11-high:7}

（1）证明：函数方程 $x^{3} - 3x = t$ 存在三个在闭区间 $[-2,2]$ 上连续，在开区间 $(-2,2)$ 内连续可微的解 $x = \varphi_{1}(t), x = \varphi_{2}(t), x = \varphi_{3}(t)$ 满足：
$\varphi_{1}(-t) = -\varphi_{3}(t), \varphi_{2}(-t) = -\varphi_{2}(t), |t| \leq 2$。

（2）若 $f$ 是 $[-2,2]$ 上的连续偶函数，证明：
$\int_{1}^{2} f\left(x^{3} - 3x\right) \mathrm{d} x = \int_{0}^{1} f\left(x^{3} - 3x\right) \mathrm{d} x$。

\par\href{https://github.com/lzxcr/cmc-mathematics/blob/64916fa1fdbc3e444a652eca204d5669058ee321/src/problems/finals/11/07.md}{题目与来源条目}

\end{example}

\begin{solve}\label{ans:fin11-high:7}

（1）令 $g(x)=x^3-3x$。它在 $[-2,-1]$、$[-1,1]$、$[1,2]$ 上分别严格递增、递减、递增，且每个分支的值域均为 $[-2,2]$。取三支反函数为 $\varphi_1,\varphi_2,\varphi_3$，反函数定理给出所需连续性和内部可微性。由 $g$ 为奇函数，立即得到
\[
\varphi_1(-t)=-\varphi_3(t),\qquad\varphi_2(-t)=-\varphi_2(t).
\]

（2）不必显式求出三支反函数；只要找它们之间的关系，换元后各支的 Jacobian 就能相消。Vieta 公式给出 $\varphi_1+\varphi_2+\varphi_3=0$，从而内部有 $\sum_j\varphi_j'=0$。在三支上分别换元，注意中间一支反向，得到
\[
\int_{-2}^{-1}f(g(x))\,\mathrm dx-\int_{-1}^{1}f(g(x))\,\mathrm dx+\int_1^2f(g(x))\,\mathrm dx
=\int_{-2}^2f(t)\sum_{j=1}^3\varphi_j'(t)\,\mathrm dt=0.
\]
端点积分按广义积分理解，因每支单调且总变差有限而绝对收敛。又 $f\circ g$ 为偶函数，上式便是所求等式的两倍。

三支的端点值依次是 $\varphi_1(-2)=-2,\varphi_1(2)=-1$，$\varphi_2(-2)=1,\varphi_2(2)=-1$，以及 $\varphi_3(-2)=1,\varphi_3(2)=2$。内部有 $\varphi_j'=1/[3(\varphi_j^2-1)]$，中间支为负、两侧支为正，所以换元时中间积分带负号。端点处导数可能无界，但 $\int_{-2}^2|\varphi_j'|$ 等于该支的区间长度，有限；再乘有界的 $f$ 后仍绝对可积。因此可先在截断参数区间积分，再取端点极限，保证根和恒等式的积分使用合法。

\end{solve}

\begin{remark}[方法与条件]三支反函数是同一个三次方程的三个根，根之和固定，所以它们导数的带符号和为零。这解释了积分恒等式为何不需要求根公式。$x=2\cos\theta$ 与三倍角公式提供另一种参数化，但不参与相消机制。\end{remark}

\begin{example}[来源序号10，低年级组]\label{ex:fin11-low:9}\label{pr:fin11-low:9}

设 $m, n$ 为正整数。 证明：当参数 $k \neq 0$ 时，微分方程
\[
y^{\prime}(x) = ky^{2n}(x) + x^{2m-1}
\]
的所有解都不是全局解（全局解即指定义在 $(-\infty,+\infty)$ 上的解）。

\par\href{https://github.com/lzxcr/cmc-mathematics/blob/64916fa1fdbc3e444a652eca204d5669058ee321/src/problems/finals/11/10.md}{题目与来源条目}

\end{example}

\begin{solve}\label{ans:fin11-low:9}

若 $k<0$，将 $y(x)$ 换成 $y(-x)$，便把系数变成 $-k>0$ 而保留强迫项 $x^{2m-1}$，故只须考虑 $k>0$。

假设存在全局解。对 $x\geq0$，$y'(x)\geq x^{2m-1}$，故 $y(x)\geq y(0)+x^{2m}/(2m)$，从某个 $T$ 起为正。对正解段，选择导数能抵消 $y^{2n}$ 的负幂作为新变量，得到
\[
\frac{\mathrm d}{\mathrm dx}y(x)^{1-2n}=(1-2n)y(x)^{-2n}y'(x)\leq-(2n-1)k.
\]
右端为负常数，积分会使正数 $y(x)^{1-2n}$ 在有限时间变为非正，矛盾。因此不存在全局解。

若 $z(x)=y(-x)$，则 $z'(x)=-ky(-x)^{2n}+x^{2m-1}$，因为强迫项次数为奇数，确实只改变 $k$ 的符号。在 $k>0$ 的情形中，选 $T$ 使 $y(T)>0$，随后 $y'>0$ 保证 $y$ 一直为正。负幂估计积分为 $y(x)^{1-2n}\le y(T)^{1-2n}-(2n-1)k(x-T)$，右端在一个明确的有限时刻变为零，此后不可能继续保持正的实值。因而反设的全局解必不存在。

\end{solve}

\begin{remark}[方法与条件]更一般地，若正解满足 $y'\geq ky^p$、$k>0,p>1$，则从 $y(T)>0$ 开始，其可存在的向右时间不超过 $y(T)^{1-p}/[k(p-1)]$。本题的外力负责使解最终进入正半轴，超线性项负责有限时爆破。见\hyperref[app:dynamics-blowup]{方法卷：爆破与等待解}。\end{remark}

\section{决赛：第 12 届（2021 年）}

本节选入6道课程相关题；来源序号保留原题解编排，未选入题目不重新补号。

\begin{example}[来源序号1]\label{ex:fin12-high:1}\label{pr:fin12-high:1}

（填空题）设 $\Omega:(x-2)^{2}+(y-3)^{2}+(z-4)^{2}\leq1$，则积分
\[
\iiint_{\Omega}\left(x^{2}+2y^{2}+3z^{2}\right)\,\mathrm{d}x\,\mathrm{d}y\,\mathrm{d}z=
\]
$\underline{\qquad\qquad}$。

\par\href{https://github.com/lzxcr/cmc-mathematics/blob/64916fa1fdbc3e444a652eca204d5669058ee321/src/problems/finals/12/01.md}{题目与来源条目}

\end{example}

\begin{solve}\label{ans:fin12-high:1}

答案为 $1424\pi/15$。平移到单位球 $B$ 后奇次项积分为零，且
\[
|B|=\frac{4\pi}{3},\qquad\iiint_Bu^2\,\mathrm dV=\iiint_Bv^2\,\mathrm dV=\iiint_Bw^2\,\mathrm dV=\frac{4\pi}{15}.
\]
故积分为 $70\cdot4\pi/3+6\cdot4\pi/15=1424\pi/15$。

令 $u=x-2,v=y-3,w=z-4$，展开被积函数为 $u^2+2v^2+3w^2+4u+12v+24w+70$。球的中心对称性使三个一次项积分为零。为求二次矩，利用三个坐标积分相等，再用球坐标计算 $\iiint_B(u^2+v^2+w^2)dV=4\pi\int_0^1r^4dr=4\pi/5$，所以每个坐标平方积分为 $4\pi/15$。将常数体积项和二次项分别求和，得到 $(1400+24)\pi/15$。

\end{solve}

\begin{example}[来源序号2]\label{ex:fin12-high:2}\label{pr:fin12-high:2}

（填空题）设 $x_{n}=\sum_{k=1}^{n}\frac{e^{k^{2}}}{k},y_{n}=\int_{0}^{n}e^{x^{2}}\,\mathrm{d}x$，则 $\lim_{n\to\infty}\frac{x_{n}}{y_{n}}=$ $\underline{\qquad\qquad}$。

\par\href{https://github.com/lzxcr/cmc-mathematics/blob/64916fa1fdbc3e444a652eca204d5669058ee321/src/problems/finals/12/02.md}{题目与来源条目}

\end{example}

\begin{solve}\label{ans:fin12-high:2}

答案为 $2$。和式由最后一项支配：前 $n-1$ 项总和至多为 $ne^{(n-1)^2}=o(e^{n^2}/n)$。另一方面，对连续变量 $t\to\infty$ 使用 L'Hôpital 法则，因 $[e^{t^2}/(2t)]'=e^{t^2}(1-1/(2t^2))$，有 $\int_0^te^{x^2}\,\mathrm dx\sim e^{t^2}/(2t)$。取 $t=n$，故比值趋于 $2$。

离散和中舍去末项后的相对误差至多为
\[
\frac{ne^{(n-1)^2}}{e^{n^2}/n}=n^2e^{-2n+1}\longrightarrow0,
\]
所以 $x_n/(e^{n^2}/n)\to1$。对积分的比较，分子和 $e^{t^2}/(2t)$ 都趋于无穷，后者导数在充分大的 $t$ 为正，导数比趋于 $1$，满足 L'Hôpital 法则的条件。连续积分主要来自端点附近宽度约为 $1/n$ 的区域，与离散求和只有一个末项的结构不同，故出现因子 $2$。

\end{solve}

\begin{example}[来源序号7]\label{ex:fin12-high:7}\label{pr:fin12-high:7}

称实函数 $f$ 满足条件 $(P)$：若 $f$ 在 $[0,1]$ 上非负连续，
$f(1)>f(0)=0,\int_{0}^{1}\frac{1}{f(x)}\mathrm{d}x=+\infty$，
且对任何 $x_{1},x_{2}\in[0,1]$ 成立
\[
f\left(\frac{x_{1}+x_{2}}{2}\right)\geq\frac{f\left(x_{1}\right)+f\left(x_{2}\right)}{2}.
\]

（1）令 $c>0$，对于 $f_{1}(x)=cx$ 和 $f_{2}(x)=\sqrt{x}$，分别验证 $f_{1},f_{2}$ 是否满足条件 $(P)$，并计算
$\lim_{x\to 0^{+}}\left(f_{1}(x)-xf_{1}^{\prime}(x)\right)^{m}e^{f_{1}^{\prime}(x)}$ 和 $\lim_{x\to 0^{+}}\left(f_{2}(x)-xf_{2}^{\prime}(x)\right)^{m}e^{f_{2}^{\prime}(x)}$。

（2）证明：$\forall m\geq 1$，存在满足条件 $(P)$ 的函数 $f$ 以及趋于零的正数列 $\left\{x_{n}\right\}$，使得 $f$ 在每一点 $x_{n}$ 可导，且
\[
\lim_{n\to+\infty}\left(f\left(x_{n}\right)-x_{n}f^{\prime}\left(x_{n}\right)\right)^{m}e^{f^{\prime}\left(x_{n}\right)}=+\infty.
\]

\par\href{https://github.com/lzxcr/cmc-mathematics/blob/64916fa1fdbc3e444a652eca204d5669058ee321/src/problems/finals/12/07.md}{题目与来源条目}

\end{example}

\begin{solve}\label{ans:fin12-high:7}

（1）$f_1(x)=cx$ 满足全部条件，且 $f_1-xf_1'=0$，故第一个极限为 $0$。$f_2(x)=\sqrt x$ 虽非负、连续、凹，但 $\int_0^1x^{-1/2}\,\mathrm dx=2<\infty$，不满足 $(P)$；对任何 $m\geq1$，
\[
(f_2-xf_2')^me^{f_2'}=\left(\frac{\sqrt x}{2}\right)^me^{1/(2\sqrt x)}\longrightarrow\infty.
\]

（2）$\sqrt x$ 已给出足够大的斜率，却使倒数积分收敛。把它替换成稀疏节点间的弦线，会压低函数、增大倒数积分，同时保留大斜率。为使第 $n$ 段至少贡献 $1/n$，取 $a_1=1$，递推定义
\[
a_{n+1}=a_n\exp\left(-\frac2{n\sqrt{a_n}}\right),\qquad n\geq1.
\]
于是 $a_n\downarrow0$（因为 $-\log a_n\geq2\sum_{j<n}1/j$）。令 $f(0)=0$，在 $[a_{n+1},a_n]$ 上线性连接 $(a_{n+1},\sqrt{a_{n+1}})$ 与 $(a_n,\sqrt{a_n})$。各段斜率为
\[
k_n=\frac1{\sqrt{a_n}+\sqrt{a_{n+1}}}.
\]
这些斜率随 $x$ 增大而下降，故 $f$ 凹；又 $0\leq f(x)\leq\sqrt x$，保证零点处连续，且 $f(1)=1$。每段积分满足
\[
\int_{a_{n+1}}^{a_n}\frac{\mathrm dx}{f(x)}
=\frac1{2k_n}\log\frac{a_n}{a_{n+1}}\geq\frac1n,
\]
故 $\int_0^1\mathrm dx/f(x)=\infty$。

取 $x_n=(a_n+a_{n+1})/2$。在该点 $f'=k_n\geq1/(2\sqrt{a_n})$，而切线截距满足
\[
f(x_n)-x_nf'(x_n)=\frac{\sqrt{a_na_{n+1}}}{\sqrt{a_n}+\sqrt{a_{n+1}}}
\geq\frac{\sqrt{a_n}}2\exp\left(-\frac1{n\sqrt{a_n}}\right).
\]
因此目标量的对数至少为
\[
\left(\frac12-\frac mn\right)\frac1{\sqrt{a_n}}+\frac m2\log a_n-m\log2\longrightarrow\infty.
\]
最后一步用了 $a_n\to0$ 且 $a_n^{-1/2}$ 的增长快于 $|\log a_n|$。这个同一个 $f$ 实际上对每个固定 $m\geq1$ 都有效。

折线段可明确写成 $f(x)=k_nx+b_n$，其中 $b_n=\sqrt{a_na_{n+1}}/(\sqrt{a_n}+\sqrt{a_{n+1}})>0$。因此该段倒数积分为 $k_n^{-1}\log[f(a_n)/f(a_{n+1})]=(2k_n)^{-1}\log(a_n/a_{n+1})$，代入递推关系后至少为 $1/n$。中点 $x_n$ 位于开线段内，故导数确实等于 $k_n$，且 $f(x_n)-x_nf'(x_n)=b_n$。最后对固定 $m$ 取 $n\ge4m$，对数下界中的 $1/2-m/n\ge1/4$，于是主项至少为 $1/(4\sqrt{a_n})$，必压过负的对数项。

\end{solve}

\begin{remark}[方法与条件]构造的参数并非猜出：先令每段倒数积分不小于调和级数的一项，再检查切线截距的损失为 $\exp[-m/(n\sqrt{a_n})]$，最终被斜率带来的 $\exp[1/(2\sqrt{a_n})]$ 压过。因此一个函数可同时适用于每个固定的 $m$；关键是先规定必须积累的量，再安排尺度。\end{remark}

\begin{example}[来源序号8，低年级组]\label{ex:fin12-low:8}\label{pr:fin12-low:8}

设 $\alpha, \beta, \alpha_{1}, \alpha_{2}, \beta_{1}, \beta_{2}$ 和 $A$ 均为实数。 回答以下问题：

（1）\[
\lim_{n \to \infty} \sin(n\alpha + \beta) = A
\]
成立的充要条件是什么？

（2）\[
\lim_{n \to \infty} \left(\sin(n\alpha_1 + \beta_1) + \sin(n\alpha_2 + \beta_2)\right) = 0
\]
成立的充要条件是什么？

\par\href{https://github.com/lzxcr/cmc-mathematics/blob/64916fa1fdbc3e444a652eca204d5669058ee321/src/problems/finals/12/08.md}{题目与来源条目}

\end{example}

\begin{solve}\label{ans:fin12-low:8}

使用有限频率的线性独立性：若互不相同的单位复数 $z_1,\ldots,z_s$ 满足 $\sum_jc_jz_j^n\to0$，则所有 $c_j=0$。证明只需乘以 $\overline{z_k}^{\,n}$ 后作 Cesàro 平均；非同频项的几何和平均趋于零，剩下 $c_k$。

（1）将正弦写成复指数并合并相同频率，得到充要条件：
\[
\alpha\equiv0\pmod{2\pi},\quad A=\sin\beta;
\quad\text{或}\quad
\alpha\equiv\pi\pmod{2\pi},\quad\sin\beta=0,\quad A=0.
\]

（2）以下所有同余均模 $2\pi$。充要条件为下列四种情形之一：

- $\alpha_1,\alpha_2\in\pi\mathbb Z$ 且 $\alpha_1\equiv\alpha_2$，并且 $\sin\beta_1+\sin\beta_2=0$。
- $\alpha_1,\alpha_2\in\pi\mathbb Z$ 且 $\alpha_1\not\equiv\alpha_2$，并且 $\sin\beta_1=\sin\beta_2=0$。
- $\alpha_1\notin\pi\mathbb Z$、$\alpha_2\equiv\alpha_1$，并且 $\beta_1-\beta_2\equiv\pi$。
- $\alpha_1\notin\pi\mathbb Z$、$\alpha_2\equiv-\alpha_1$，并且 $\beta_1+\beta_2\equiv0$。

为验证穷尽性，将和式写成
\[
\frac1{2i}\sum_{j=1}^2\left[e^{i\beta_j}(e^{i\alpha_j})^n-e^{-i\beta_j}(e^{-i\alpha_j})^n\right].
\]
当两频率属于 $\{1,-1\}$ 时，分别消去常数和交错项，给出前两种情形；否则非实频率必须成对相同，给出后两种相位条件。上述引理保证无其他可能。每种条件下和式都恒为零，故充分性也成立。

频率引理中的平均公式为 $N^{-1}\sum_{n=1}^N(z_j\overline z_k)^n$：当 $j=k$ 时等于 $1$，否则其分子是有界的几何和，除以 $N$ 后趋零。另一方面，原序列趋零且所乘因子模为 $1$，其平均也趋零，所以留下的 $c_k$ 必为零。第一问对 $\sin(n\alpha+\beta)-A$ 使用引理，包含额外的常数频率 $1$。第二问在同频率情形的系数条件是 $e^{i\beta_1}+e^{i\beta_2}=0$，在反频率情形则是 $e^{i\beta_1}-e^{-i\beta_2}=0$，分别化为所列两条相位同余。

\end{solve}

\begin{remark}[方法与条件]Cesàro 平均在这里相当于取离散 Fourier 系数：同频项保留，异频项因几何和有界而消失。有限三角和若收敛，合并同余频率后所有非恒定项都必须消失，所以序列从一开始就是常数。参见\hyperref[app:integration-fourier]{方法卷：Fourier 频率分离}。\end{remark}

\begin{example}[来源序号10，低年级组]\label{ex:fin12-low:9}\label{pr:fin12-low:9}

设 $g$ 为 $\mathbb{R}$ 上恒正的连续函数，对于正整数 $n$ 以及 $x_{0}, y_{0} \in \mathbb{R}$，考虑微分方程
\[
\begin{cases} y^{\prime}(x) = y^{\frac{1}{2n+1}}(x) g(x), \\ y\left(x_{0}\right) = y_{0}. \end{cases}
\]
证明：

（1）上述方程有定义在整个 $\mathbb{R}$ 上的解（称为全局解）；

（2）若 $y_{0} = 0$，则上述方程有无穷多个全局解；

（3）若 $y = y(x)$ 是上述方程的解，则 $y$ 在 $\mathbb{R}$ 上非负，或在 $\mathbb{R}$ 上非正。

\par\href{https://github.com/lzxcr/cmc-mathematics/blob/64916fa1fdbc3e444a652eca204d5669058ee321/src/problems/finals/12/10.md}{题目与来源条目}

\end{example}

\begin{solve}\label{ans:fin12-low:9}

正的 $g$ 可以吸收到递增的时间坐标中；再对 $|y|$ 取适当幂，分离变量后的方程成为线性的。记
\[
q=\frac{2n}{2n+1},\qquad G(x)=\int_{x_0}^xg(t)\,\mathrm dt,\qquad t_+=\max(t,0).
\]

（1）当 $y_0\ne0$ 时，显式函数
\[
y(x)=\operatorname{sgn}(y_0)\bigl(|y_0|^q+qG(x)\bigr)_+^{1/q}
\]
是全局解。由于 $1/q>1$，正部的这个幂是 $C^1$ 函数，在接点处导数为零，直接代入即能验证。$y_0=0$ 时可取零解。

（2）若 $y_0=0$，对任意等待时刻 $a\geq x_0$ 及符号 $\sigma=\pm1$，
\[
y_{a,\sigma}(x)=\sigma\left(q\int_a^xg(t)\,\mathrm dt\right)_+^{1/q}
\]
都是全局解：它在 $x\leq a$ 时保持零，从 $a$ 起向正或负方向离开，给出无穷多个解。

（3）在任意非零解区间内，方程使 $y>0$ 时严格递增，$y<0$ 时严格递减。因此非零连通分支不可能在某个有限右端点回到零，必一直延伸到正无穷。两个不同的非零分支因而不可能共存，所以全局解只能非负、非正或恒为零。

对非零解，$|y|$ 满足 $(|y|)'=g|y|^{1/(2n+1)}$，所以 $(|y|^q)'=qg$，这解释了显式公式的幂次。接点处正部幂的差商趋零，因为 $1/q>1$，而原方程右端也为零，故拼接没有破坏可导性。第三问中若某个非零分支含点 $x_1$，则此后 $|y|$ 严格递增，始终至少为 $|y(x_1)|>0$，不可能在有限右端点连续接回零；不同分支若都延伸到正无穷必相交，故至多存在一个非零分支，符号不能改变。

\end{solve}

\begin{remark}[方法与条件]不满足局部 Lipschitz 条件本身并不足以断言不唯一；本题还可算出 $\int_0^\varepsilon s^{-1/(2n+1)}\,\mathrm ds<\infty$，这才允许解在有限时间离开零点，并与零解拼接。有限离开时间和有限爆破时间由相同的倒数积分机制控制，见\hyperref[app:dynamics-blowup]{方法卷：爆破与等待解}。\end{remark}

\begin{example}[来源序号15，高年级组选做]\label{ex:fin12-high:13}\label{pr:fin12-high:13}

考虑求解常微分方程初值问题
\[
\begin{cases} y' = f(x,y), \\ y\left(x_0\right) = y_0 \end{cases}
\]
的下列数值格式：
$y_{n}+a_{1}y_{n-1}+a_{2}y_{n-2}=h\left(b_{0}f_{n}+b_{1}f_{n-1}+b_{2}f_{n-2}\right)$，
其中 $a_{1},a_{2},b_{0},b_{1},b_{2}$ 为常数，
$f_{j}=f\left(x_{j},y_{j}\right),j=n-2,n-1,n$。

（1）确定常数 $a_{1},a_{2},b_{0},b_{1},b_{2}$，使得上述数值格式具有尽可能高阶的精度；

（2）分析上一步得到的数值格式的稳定性与收敛性。

\par\href{https://github.com/lzxcr/cmc-mathematics/blob/64916fa1fdbc3e444a652eca204d5669058ee321/src/problems/finals/12/15.md}{题目与来源条目}

\end{example}

\begin{solve}\label{ans:fin12-high:13}

（1）阶数由格式对低次多项式的精确性决定。把精确解代入，记左端减右端为残量；它为 $O(h^{p+1})$ 时称格式为 $p$ 阶。在 $x_n$ 处展开后，$h^jy^{(j)}(x_n)$ 的系数为 $(-1)^jd_j/j!$，故阶条件为
\[
\begin{aligned}
d_0&=1+a_1+a_2=0,\\
d_1&=a_1+2a_2+b_0+b_1+b_2=0,\\
d_j&=a_1+2^ja_2+jb_1+j2^{j-1}b_2=0\quad(j\geq2).
\end{aligned}
\]
前五个条件 $d_0=\cdots=d_4=0$ 唯一确定
\[
a_1=0,\quad a_2=-1,\quad(b_0,b_1,b_2)=(1/3,4/3,1/3).
\]
消元过程可以逐步核对：$d_3-d_2=0$ 给出 $4a_2+b_1+8b_2=0$，$d_4-d_3=0$ 给出 $8a_2+b_1+20b_2=0$；相减得到 $a_2=-3b_2$，回代得 $b_1=4b_2$。再由 $d_2=0$ 得 $a_1=0$，由 $d_0=0$ 得 $a_2=-1$，最后 $d_1=0$ 给出 $b_0=1/3$。因此任何五阶或更高阶候选都必须先等于这一唯一解，再检验下一阶条件。

此时 $d_5=4/3\ne0$，故最高阶为 $4$，所得隐式 Simpson 多步格式为
\[
y_n-y_{n-2}=\frac h3(f_n+4f_{n-1}+f_{n-2}).
\]

（2）第一特征多项式 $\rho(\zeta)=\zeta^2-1$ 的根 $\pm1$ 均为单位圆上的单根，故格式零稳定；又 $\rho(1)=0$、$\rho'(1)=\sigma(1)=2$，其中 $\sigma(\zeta)=(\zeta^2+4\zeta+1)/3$，故相容。由 Dahlquist 等价定理，格式收敛；在通常的光滑性、Lipschitz 条件及四阶准确起步值下，具有四阶整体误差。

零稳定不等于绝对稳定。对测试方程 $y'=\lambda y$，令 $z=h\lambda<0$，放大因子满足 $\rho(\zeta)-z\sigma(\zeta)=0$。此多项式在 $\zeta=-1$ 的值为 $2z/3<0$，而在 $\zeta\to-\infty$ 时趋于正无穷，故有实根 $\zeta<-1$。所以该格式没有负实轴上的绝对稳定区间，更不是 A 稳定的。

收敛性的含义是固定有限时间区间上，当 $h\to0$ 且起步误差趋零时数值解趋于精确解。这里需假定精确解有足够连续导数，例如五阶连续导数，并且 $f$ 在解附近关于 $y$ 满足统一 Lipschitz 条件；$h$ 足够小时 $hb_0$ 乘该常数小于 $1$，隐式一步可局部唯一求解。零稳定允许单位圆上出现单根 $-1$，而绝对稳定要求固定 $z=h\lambda$ 下所有放大根都在单位圆内或满足边界单根条件，两者考察的问题不同，所以四阶收敛与负实轴绝对不稳定并不矛盾。

\end{solve}

\begin{remark}[方法与条件]这里的根 $-1$ 表示交错寄生模态：步长趋零、总时长固定时，它仍符合零稳定条件；步长固定、迭代次数趋于无穷时，它却破坏衰减。这解释了四阶收敛与缺乏负实轴绝对稳定区间为何能够同时成立。见\hyperref[app:dynamics-stability]{方法卷：数值稳定性}。\end{remark}

\section{决赛：第 13 届（2022 年）}

本节选入6道课程相关题；来源序号保留原题解编排，未选入题目不重新补号。

\begin{example}[来源序号1]\label{ex:fin13-high:1}\label{pr:fin13-high:1}

设函数 $y = y(x)$ 在 $x = 0$ 附近由方程 $y + 2y^2 + y^3 = e^{-x} + x - 1$ 所确定，且 $y = Ax^2 + Bx^3 + o\left(x^3\right)(x \rightarrow 0)$，则 $(A,B) = \underline{\qquad}$。

\par\href{https://github.com/lzxcr/cmc-mathematics/blob/64916fa1fdbc3e444a652eca204d5669058ee321/src/problems/finals/13/01.md}{题目与来源条目}

\end{example}

\begin{solve}\label{ans:fin13-high:1}

由 $y=O(x^2)$，原方程左端为 $y+O(x^4)$，右端为 $\frac12x^2-\frac16x^3+O(x^4)$。比较系数得
\[
(A,B)=\boxed{\left(\frac12,-\frac16\right)}.
\]

题设展开已保证 $y(0)=0$ 且 $y$ 至少为二阶小量，因此 $y^2=O(x^4)$、$y^3=O(x^6)$，对所求三阶系数都没有贡献。右端用 $e^{-x}=1-x+x^2/2-x^3/6+O(x^4)$，常数项与一次项恰被 $x-1$ 消去。两边相减得 $(A-1/2)x^2+(B+1/6)x^3+o(x^3)=0$，先除以 $x^2$ 取极限得 $A=1/2$，再消去二阶项并除以 $x^3$，得到 $B=-1/6$。

\end{solve}

\begin{example}[来源序号2]\label{ex:fin13-high:2}\label{pr:fin13-high:2}

极限 $\lim_{(x,y) \rightarrow (0,0)} \left( |x| + |y| \right)^{|xy|} = \underline{\qquad}$。

\par\href{https://github.com/lzxcr/cmc-mathematics/blob/64916fa1fdbc3e444a652eca204d5669058ee321/src/problems/finals/13/02.md}{题目与来源条目}

\end{example}

\begin{solve}\label{ans:fin13-high:2}

幂型极限先取对数。记 $r=|x|+|y|$，由 $|xy|\le r^2/4$，当 $0<r<1$ 时，
\[
\bigl||xy|\log r\bigr|\le\frac{r^2}{4}|\log r|\longrightarrow0.
\]
再用指数函数的连续性，所求极限为 $\boxed1$。

不等式 $4|xy|\le(|x|+|y|)^2$ 来自 $(|x|-|y|)^2\ge0$，因此估计对所有趋近路径同时成立。在坐标轴上 $|xy|=0$、底数仍为正，原式等于 $1$；离开坐标轴也可统一写成 $\exp(|xy|\ln r)$。令 $u=1/r$，则 $r^2|\ln r|=(\ln u)/u^2\to0$，故指数趋零。只检查固定直线路径不足以证明二维极限，这个径向上界则覆盖了全部路径。

\end{solve}

\begin{example}[来源序号4]\label{ex:fin13-high:4}\label{pr:fin13-high:4}

常微分方程 $y' - \cos x \sin^2 y = \frac{1}{2} \sin 2y$ 的通解为$\underline{\qquad}$。

\par\href{https://github.com/lzxcr/cmc-mathematics/blob/64916fa1fdbc3e444a652eca204d5669058ee321/src/problems/finals/13/04.md}{题目与来源条目}

\end{example}

\begin{solve}\label{ans:fin13-high:4}

右端含 $\sin^2y$ 与 $\sin y\cos y$，除以 $\sin^2y$ 后可用 $\frac{\mathrm d}{\mathrm dx}\cot y=-y'/\sin^2y$ 线性化。在 $\sin y\ne0$ 的区间令 $u=\cot y$，方程化为 $u'+u=-\cos x$，从而
\[
\boxed{\cot y=Ce^{-x}-\frac12(\sin x+\cos x)}.
\]
另有常值解 $\boxed{y\equiv k\pi\ (k\in\mathbb Z)}$。各非恒定解取连续的反余切分支。

除以 $\sin^2y$ 后原式为 $y'/\sin^2y-\cos x=\cot y$，所以 $u'=-y'/\sin^2y$ 给出 $u'+u=-\cos x$。乘 $e^x$ 并积分，$e^xu=C-\tfrac12e^x(\sin x+\cos x)$，于是得到通解。原方程的右端关于 $y$ 光滑，初值问题唯一；若某解在有限点取到 $k\pi$，就必须是常值解。因此其余解始终位于某个 $(k\pi,(k+1)\pi)$ 内，可在该区间选定唯一连续反余切分支，所得表达式在全实轴有定义。

\end{solve}

\begin{example}[来源序号7]\label{ex:fin13-high:7}\label{pr:fin13-high:7}

设 $\alpha > 0, f$ 在 $[0,+\infty)$ 上二阶可导，

\[
\lim_{x \rightarrow +\infty} \frac{|f(x)|}{x^\alpha} = 0, \qquad \lim_{x \rightarrow +\infty} x^\alpha |f''(x)| < +\infty.
\]

（1）若 $\alpha \in (0,1]$，证明：$\lim_{x \rightarrow +\infty} f'(x) = 0$。

（2）若 $\alpha > 1$，试构造满足题设条件的函数 $f$，使得 $\lim_{x \rightarrow +\infty} f'(x) = 0$ 不成立，且在 $[0,+\infty)$ 的任何子区间上，$f'$ 不恒等于常数。

\par\href{https://github.com/lzxcr/cmc-mathematics/blob/64916fa1fdbc3e444a652eca204d5669058ee321/src/problems/finals/13/07.md}{题目与来源条目}

\end{example}

\begin{solve}\label{ans:fin13-high:7}

（1）用长度为 $h$ 的割线近似导数：函数项除以 $h$，二阶余项乘以 $h$。因此按题设的增长尺度取 $h=sx^\alpha$，最后让 $s$ 趋零。令 $C=\sup_{t\ge x_0}t^\alpha|f''(t)|<\infty$。固定 $s\in(0,1)$，对充分大的 $x$ 取 $h=sx^\alpha\le x$。Taylor 公式给出
\[
|f'(x)|\le\frac{|f(x+h)|+|f(x)|}{h}
       +\frac h2\sup_{x\le t\le x+h}|f''(t)|
\le o(1)+\frac{Cs}{2}.
\]
这里固定 $s$ 后第一项趋于零，因为 $x+h\le2x$ 且 $f(t)=o(t^\alpha)$。先令 $x\to\infty$，再令 $s\downarrow0$，即得 $f'(x)\to0$。

（2）对任意 $\alpha>1$，取 $f(x)=x+e^{-x}$。它满足 $f=o(x^\alpha)$、$x^\alpha|f''(x)|\to0$，但 $f'(x)=1-e^{-x}\to1$；且 $f''(x)>0$，故 $f'$ 在任何非退化区间上均不恒定。

固定 $s$ 后，函数项可写成
\[
\frac{|f(x+h)|}{sx^\alpha}
=\frac1s\frac{|f(x+h)|}{(x+h)^\alpha}\left(\frac{x+h}{x}\right)^\alpha\to0,
\]
因为比值 $(x+h)/x$ 不超过 $2$。余项则由 $t\ge x$ 得 $|f''(t)|\le Cx^{-\alpha}$，乘 $h/2$ 恰为 $Cs/2$。因此结论是对每个固定 $s>0$ 有 $\limsup|f'|\le Cs/2$，然后再让 $s$ 趋零；不能在第一步就让 $s$ 随 $x$ 任意变化而失去函数项控制。

\end{solve}

\begin{example}[来源序号10，低年级组]\label{ex:fin13-low:9}\label{pr:fin13-low:9}

证明：

（1）函数方程 $e^{y} - y = e^{x} - x$ 唯一地确定了一个连续函数 $y = y(x) \in C(-\infty, 0]$ 满足 $y(0) = 0$ 以及 $y(x) \geq 0$。进一步，证明当 $x < 0$ 时成立 $y(x) > 0$。

（2）$y(x) \in C^{1}(-\infty, 0]$，且 $y'(0) = -1$。

（3）对于 $x < 0$ 成立 $x + y(x) < 0$。进一步，$y(x) \in C^{2}(-\infty, 0]$ 是凹函数，且 $y''(0) = -\frac{2}{3}$。

\par\href{https://github.com/lzxcr/cmc-mathematics/blob/64916fa1fdbc3e444a652eca204d5669058ee321/src/problems/finals/13/10.md}{题目与来源条目}

\end{example}

\begin{solve}\label{ans:fin13-low:9}

令 $H(t)=e^t-t$。它在 $(-\infty,0]$ 上严格递减，在 $[0,\infty)$ 上严格递增，两支值域均为 $[1,\infty)$，所以所求连续函数存在且唯一，且 $x<0$ 时 $y(x)>0$。

在 $x<0$ 时，$e^{y(x)}-1>0$，隐函数定理给出光滑性。原点处这个偏导数为零；应提取 $H(t)-H(0)$ 的二次零点，而不是反复除以消失的量。令
\[
\Phi(t)=t\sqrt{\frac{2(e^t-t-1)}{t^2}},\qquad \Phi(0)=0.
\]
根号内在零点解析延拓且值为 $1$，故 $\Phi$ 在原点附近解析，$\Phi(t)=t+t^2/6+O(t^3)$，$\Phi'(0)=1$。等式 $H(y)=H(x)$ 的异侧分支满足 $\Phi(y)=-\Phi(x)$，由反函数定理，$y$ 在原点也有解析延拓。比较展开得
\[
y(x)=-x-\frac{x^2}{3}+O(x^3),
\qquad y'(0)=-1,\quad y''(0)=-\frac23.
\]
对 $x<0$，写 $t=-x>0$。由 $H(t)-H(-t)=2(\sinh t-t)>0$，以及 $H$ 在正半轴单增，得 $0<y<t$，即 $x+y<0$。

最后隐式求导给出
\[
y'=\frac{e^x-1}{e^y-1},\qquad
 y''=\frac{(e^y-e^x)(e^{x+y}-1)}{(e^y-1)^3}<0.
\]
结合原点处的光滑性，得到 $y\in C^2(-\infty,0]$ 且严格凹。

展开的系数可由 $\Phi(y)=-\Phi(x)$ 直接核对：先得 $y=-x+O(x^2)$，设二阶项为 $Bx^2$，则左端为 $-x+(B+1/6)x^2+O(x^3)$，右端为 $-x-x^2/6+O(x^3)$，所以 $B=-1/3$。在内点，再次求导得到 $e^y(y')^2+(e^y-1)y''=e^x$，将 $y'=(e^x-1)/(e^y-1)$ 代入后，分子因式分解为 $(e^y-e^x)(e^{x+y}-1)$。因 $y>0>x$ 且 $x+y<0$，两个因子一正一负，严格凹性随之成立。

\end{solve}

\begin{remark}[方法与条件]局部坐标 $\Phi$ 将非退化极小点附近的 $H$ 化为 $H(0)+\Phi^2/2$，这是单变量 Morse 引理的具体实现。交换等值点的映射在这个坐标中就是反射 $\Phi\mapsto-\Phi$，因而临界点并不妨碍该分支的光滑性。\end{remark}

\begin{example}[来源序号12，低年级组]\label{ex:fin13-low:10}\label{pr:fin13-low:10}

设 $n$ 为正整数。证明微分方程 $y^{(2n)} - x^{2} y = |\sin x|$ 满足初值条件

\[
y(0) = y''(0) = \cdots = y^{(2n-2)}(0) = 1, \qquad y'(0) = y'''(0) = \cdots = y^{(2n-1)}(0) = 0
\]

的解为 $(-\infty, +\infty)$ 上的偶函数并且 $y(x) > 1, x \neq 0$。

\par\href{https://github.com/lzxcr/cmc-mathematics/blob/64916fa1fdbc3e444a652eca204d5669058ee321/src/problems/finals/13/12.md}{题目与来源条目}

\end{example}

\begin{solve}\label{ans:fin13-low:10}

线性微分方程的全局存在唯一性保证初值问题在 $\mathbb R$ 上有唯一解。函数 $y(-x)$ 满足同一方程及同一初值，故 $y$ 为偶函数。

要证明严格大于 $1$，先把高阶方程写成带非负核的积分方程。对 $x\ge0$，反复积分得到 Volterra 形式
\[
y(x)=\sum_{j=0}^{n-1}\frac{x^{2j}}{(2j)!}
+\frac1{(2n-1)!}\int_0^x(x-t)^{2n-1}\bigl(t^2y(t)+|\sin t|\bigr)\,\mathrm dt.
\]
若 $y$ 在正半轴有零点，取首个零点 $x_0$；在 $[0,x_0)$ 上 $y>0$，上式却给出 $y(x_0)\ge1$，矛盾。因此 $y(x)>0$。再代入上式，$x>0$ 时积分严格为正，于是 $y(x)>1$。偶性完成证明。

高阶线性方程可改写成 $2n$ 维一阶线性系统，系数与强迫项在任意有限区间连续，因此解可逐段唯一延拓到整个实轴，不会在有限位置失去存在性。反射 $x\mapsto-x$ 时，最高阶为偶数、$x^2$ 与 $|\sin x|$ 都为偶函数，奇数阶初值又全为零，所以反射解确实满足同一初值。积分表示中的初始多项式来自 Taylor 公式，核 $(x-t)^{2n-1}/(2n-1)!$ 非负；证明 $y>0$ 后，任意 $x>0$ 的内部区间上 $t^2y(t)>0$，故积分严格为正。

\end{solve}

\section{决赛：第 14 届（2023 年）}

本节选入6道课程相关题；来源序号保留原题解编排，未选入题目不重新补号。

\begin{example}[来源序号1]\label{ex:fin14-high:1}\label{pr:fin14-high:1}

由方程 $\begin{cases} x+y=t, \\ x^2-y^2=t \end{cases}$ 确定的曲线 $\begin{cases} x=x(t), \\ y=y(t) \end{cases}$ 在 $(1,0)$ 处的切线方程为$\underline{\qquad}$。

\par\href{https://github.com/lzxcr/cmc-mathematics/blob/64916fa1fdbc3e444a652eca204d5669058ee321/src/problems/finals/14/01.md}{题目与来源条目}

\end{example}

\begin{solve}\label{ans:fin14-high:1}

在 $(1,0)$ 附近有 $t=x+y\ne0$，故 $(x-y)(x+y)=x+y$ 化为 $x-y=1$。曲线局部就是直线，切线为 $\boxed{y=x-1}$。

由消元得到 $(x+y)(x-y-1)=0$，全体代数轨迹实际上还包含 $x+y=0$ 这一支；但点 $(1,0)$ 不在该支，故附近可以安全除以 $x+y$。在所需分支上联立 $x+y=t$、$x-y=1$，得 $x=(t+1)/2,y=(t-1)/2$，在 $t=1$ 时经过 $(1,0)$，且 $x'=y'=1/2$，切线斜率为 $1$，与消元所得直线一致。

\end{solve}

\begin{example}[来源序号2]\label{ex:fin14-high:2}\label{pr:fin14-high:2}

记 $L$ 为曲线 $\begin{cases} x^2+y^2+z^2=1, \\ x+2y+z=0 \end{cases}$，则曲线积分 $\int_L \left(2x^2+x+y^2+y\right) \mathrm{d} s=$$\underline{\qquad}$。

\par\href{https://github.com/lzxcr/cmc-mathematics/blob/64916fa1fdbc3e444a652eca204d5669058ee321/src/problems/finals/14/02.md}{题目与来源条目}

\end{example}

\begin{solve}\label{ans:fin14-high:2}

令 $\nu=(1,2,1)^{\mathrm T}/\sqrt6$ 为圆所在平面的单位法向量。单位大圆的二阶矩为
\[
\int_L uu^{\mathrm T}\,\mathrm ds=\pi(I-\nu\nu^{\mathrm T}),
\]
这是在该平面取正交基、令 $u=e_1\cos t+e_2\sin t$ 直接积分所得。线性项由中心对称性积分为零，故
\[
\int_L(2x^2+x+y^2+y)\,\mathrm ds
=\pi\left[2\left(1-\frac16\right)+1-\frac46\right]=\boxed{2\pi}.
\]

矩阵公式也可展开成一元积分。取平面内两个正交单位向量 $e_1,e_2$，令 $u(t)=e_1\cos t+e_2\sin t$，$0\le t\le2\pi$；此时 $|u'|=1$，所以 $ds=dt$。由于 $\int\cos^2t=\int\sin^2t=\pi$、$\int\sin t\cos t=0$，二阶矩等于 $\pi(e_1e_1^{\mathrm T}+e_2e_2^{\mathrm T})$。正交基补上 $\nu$ 后得到括号为 $I-\nu\nu^{\mathrm T}$，故 $\int_Lx^2ds=5\pi/6$、$\int_Ly^2ds=\pi/3$，加权后为 $2\pi$。

\end{solve}

\begin{example}[来源序号4]\label{ex:fin14-high:4}\label{pr:fin14-high:4}

设 $\mathbb{R}$ 上函数 $f$ 具有连续的一阶导数，且满足 $f(x)=\int_0^x (x-t)f'(t) \mathrm{d} t+x^2(x \in \mathbb{R})$，则 $f(x)=$$\underline{\qquad}$。

\par\href{https://github.com/lzxcr/cmc-mathematics/blob/64916fa1fdbc3e444a652eca204d5669058ee321/src/problems/finals/14/04.md}{题目与来源条目}

\end{example}

\begin{solve}\label{ans:fin14-high:4}

代入 $x=0$ 得 $f(0)=0$。对积分方程求导并用微积分基本定理，得 $f'(x)=f(x)+2x$。解这一阶线性初值问题可得
\[
\boxed{f(x)=2e^x-2x-2}.
\]

变上限求导时，上端点项为 $(x-x)f'(x)=0$，核对 $x$ 的偏导数为 $f'(t)$，所以积分的导数是 $\int_0^xf'(t)dt=f(x)-f(0)$。对 $f'-f=2x$ 乘积分因子 $e^{-x}$，积分得 $e^{-x}f(x)=2\int_0^xte^{-t}dt=2-2(x+1)e^{-x}$，即所示解。反向验证时，原方程左右差的导数为零，且在 $x=0$ 差为零，故所得微分方程解也确实满足原积分方程。

\end{solve}

\begin{example}[来源序号7]\label{ex:fin14-high:7}\label{pr:fin14-high:7}

设 $q_1, q_2, \dots, q_n, \dots$ 为有理数集 $\mathbb{Q}$ 的一个排列，满足 $|q_n| < n (\forall n \geq 1)$。又设 $r \in (0, 1), g(x) = \sum_{n=1}^{\infty} r^n (x - q_n)^{\frac{1}{3}}$，其中立方根取实值。证明：

（i）$g$ 在 $\mathbb{R}$ 上连续且严格单增，集合 $E = \{g(q) \mid q \in \mathbb{Q}\}$ 的闭包为 $\mathbb{R}$；

（ii）$\lim_{x \to t} \frac{g(x) - g(t)}{x - t} = \sum_{n=1}^{\infty} \frac{r^n}{3 (t - q_n)^{\frac{2}{3}}}$，这里规定 $\frac{1}{0} = +\infty$；

（iii）$g$ 的反函数 $f$ 处处可导，且 $f'$ 在 $E$ 上为零；

（iv）对任意 $a < b$，存在不同的 $\xi, \eta \in (a, b)$，使得 $f'(\xi) = f'(\eta) = \frac{f(b) - f(a)}{b - a}$。

\par\href{https://github.com/lzxcr/cmc-mathematics/blob/64916fa1fdbc3e444a652eca204d5669058ee321/src/problems/finals/14/07.md}{题目与来源条目}

\end{example}

\begin{solve}\label{ans:fin14-high:7}

构造的关键是先让 $g$ 在稠密集上具有无穷大差商，再由反函数把它变成零导数。

（i）在任意紧区间上，级数由 $\sum r^n(R+n)^{1/3}$ 控制，故局部一致收敛。各项严格递增，所以 $g$ 连续且严格递增；只保留差 $g(x)-g(0)$ 中第一项，就得 $g(x)\to\pm\infty$ 当 $x\to\pm\infty$。于是 $g$ 是 $\mathbb R$ 上的同胚，$E=g(\mathbb Q)$ 稠密。

（ii）设
\[
D(t)=\sum_{n\ge1}\frac{r^n}{3|t-q_n|^{2/3}}\in(0,\infty].
\]
每项差商非负，故对有限部分和取极限，再令项数增大，得到 $\liminf_{x\to t}(g(x)-g(t))/(x-t)\ge D(t)$。这已处理 $D(t)=\infty$ 的情形。

若 $D(t)<\infty$，则 $t\ne q_n$。对实立方根，置 $a=\sqrt[3]{x-q_n}$、$b=\sqrt[3]{t-q_n}$，有
\[
0<\frac{\sqrt[3]{x-q_n}-\sqrt[3]{t-q_n}}{x-t}
=\frac1{a^2+ab+b^2}\le\frac4{3b^2},
\]
因为 $a^2+ab+b^2=(a+b/2)^2+3b^2/4$。右端乘以 $r^n$ 后可求和，因此控制收敛定理允许逐项取极限，得到所述公式。

（iii）令 $s=g(t)$。反函数差商给出
\[
f'(s)=\frac1{D(t)},\qquad 1/\infty:=0.
\]
故 $f$ 处处可导，且 $f'$ 在 $E$ 上为零。

（iv）记 $k=(f(b)-f(a))/(b-a)>0$。中值定理给出 $\xi\in(a,b)$ 使 $f'(\xi)=k$。两个子区间的割线斜率至少一个不小于 $k$；在该子区间再次用中值定理，得到另一点 $c\ne\xi$ 使 $f'(c)\ge k$。若严格大于，利用 $E$ 的稠密性在同一子区间取 $d$ 使 $f'(d)=0$，再由导数的 Darboux 性质取到 $\eta\ne\xi$ 满足 $f'(\eta)=k$；若等于，直接取 $\eta=c$。

反函数导数公式在无穷大差商处需按定义理解：令 $s=g(t)$，当 $h\to0$ 时，因反函数连续，$u=f(s+h)\to t$，且
\[
\frac{f(s+h)-f(s)}h=\left(\frac{g(u)-g(t)}{u-t}\right)^{-1}.
\]
右侧是正数的倒数，若原差商趋于 $+\infty$ 就趋于 $0$。第四问中总割线斜率是两段割线斜率按长度的加权平均，所以至少一段不小于总斜率。在该段内选点 $c,d$，它们之间的区间不包含原先的 $\xi$，故 Darboux 定理产生的第二个等斜率点必与 $\xi$ 不同；此处只需导数的介值性，不假定 $f'$ 连续。

\end{solve}

\begin{remark}[方法与条件]这里构造的是严格递增、处处可导的函数，其导数却在稠密集上为零。因此“导数在稠密集上为零”不能代替“导数处处为零”；若另有导数连续性，结论才会改变。\end{remark}

\begin{example}[来源序号10，低年级组]\label{ex:fin14-low:9}\label{pr:fin14-low:9}

设 $f$ 是 $[0, 1]$ 上的可积函数，$\alpha > 0$ 是常数。证明：$\lim_{x \to 0^+} \frac{1}{x} \int_0^x f(t) \mathrm{d}t$ 存在当且仅当 $\lim_{x \to 0^+} \frac{1}{x^\alpha} \int_0^x t^{\alpha-1} f(t) \mathrm{d}t$ 存在。

\par\href{https://github.com/lzxcr/cmc-mathematics/blob/64916fa1fdbc3e444a652eca204d5669058ee321/src/problems/finals/14/10.md}{题目与来源条目}

\end{example}

\begin{solve}\label{ans:fin14-low:9}

两种平均未必来自 $f$ 的点态极限；应把已知平均的原函数作为分部积分中的整体。证明一个对称形式。设 $a,b>0$，
\[
F_a(x)=\int_0^xt^{a-1}f(t)\,\mathrm dt,
\qquad F_a(x)=Ax^a+o(x^a).
\]
在 $[\varepsilon,x]$ 上分部积分并令 $\varepsilon\downarrow0$，得到
\[
F_b(x)=x^{b-a}F_a(x)-(b-a)\int_0^xt^{b-a-1}F_a(t)\,\mathrm dt
=\frac abAx^b+o(x^b).
\]
这里边界项趋于零，且误差积分的绝对值不超过
$|b-a|\sup_{0<t\le x}|F_a(t)/t^a-A|\,x^b/b=o(x^b)$。

所以 $F_a(x)/x^a$ 有有限极限时，$F_b(x)/x^b$ 也有有限极限；交换 $a,b$ 得逆命题。取 $a=1,b=\alpha$，两极限的关系为后者等于前者除以 $\alpha$。若权函数在零点无界，积分按零点处的广义积分理解；上述恒等式同时证明所需积分的存在。

将 $F_a(t)=At^a+t^a\eta(t)$ 代入分部积分式，主项为 $Ax^b-(b-a)A\int_0^xt^{b-1}dt=(a/b)Ax^b$，而误差为 $x^b\eta(x)-(b-a)\int_0^xt^{b-1}\eta(t)dt=o(x^b)$。由于 $b>0$，$t^{b-1}$ 在零点可积，边界项 $\varepsilon^{b-a}F_a(\varepsilon)=O(\varepsilon^b)$ 也趋零。取 $a=1$ 得到一个方向，交换 $a,b$ 得到另一个方向，并同时确定两极限间的比例，不能只证明它们“同存在”而误认为数值相同。

\end{solve}

\begin{example}[来源序号12，低年级组]\label{ex:fin14-low:10}\label{pr:fin14-low:10}

求在 $\mathbb{R}$ 上处处可微且满足 $\left( \frac{\mathrm{d}y}{\mathrm{d}x} \right)^2 = 4 - y^2$ 的所有实函数 $y = y(x)$。

\par\href{https://github.com/lzxcr/cmc-mathematics/blob/64916fa1fdbc3e444a652eca204d5669058ee321/src/problems/finals/14/12.md}{题目与来源条目}

\end{example}

\begin{solve}\label{ans:fin14-low:10}

全部解由常值段与半个余弦波拼接而成：可在 $y=2$ 或 $y=-2$ 处停留任意长时间，然后用恰好长为 $\pi$ 的弧段到达另一端，继续停留或返回；也允许永远停留。每个从时刻 $a$、高度 $2\varepsilon$ 出发的运动段为
\[
y(x)=2\varepsilon\cos(x-a),\qquad a\le x\le a+\pi,\quad\varepsilon\in\{1,-1\}.
\]
相邻运动段的 $\varepsilon$ 交替，起点间隔至少为 $\pi$。两端可以延伸为常值半直线，也可有向两边无限延续的运动段；另有两个恒定解 $y\equiv\pm2$。

证明完备性：方程保证 $|y|\le2$。在开集 $\{|y|<2\}$ 的任一连通分支上，$y'$ 不为零，由 Darboux 性质知其符号不变。分离变量得到正弦弧；最大延拓到两端必分别取值 $-2,2$，且历时恰为 $\pi$。余下区间只能停在 $\pm2$。反之，上述拼接函数在接点两侧导数都为零，因而处处可微并满足方程。

在 $|y|<2$ 的一个分支上，令 $v=\arcsin(y/2)$，则 $v'=y'/\sqrt{4-y^2}$。方程使 $v'$ 恒为 $1$ 或恒为 $-1$，所以 $v$ 是斜率为 $\pm1$ 的仿射函数；其值域只能位于 $(-\pi/2,\pi/2)$，走完整个分支所需长度就是 $\pi$。分支端点若仍有 $|y|<2$，连续性会允许继续延伸，与分支极大性矛盾。每段运动至少占长度 $\pi$，因此接点不可能在有限位置无限聚集，所列拼接在每个有限区间内都能逐段核查可导性。

\end{solve}

\begin{remark}[方法与条件]分离变量只描述 $|y|<2$ 的运动段；在端点，右端 $\sqrt{4-y^2}$ 不满足局部 Lipschitz 条件，因此允许任意等待时间。若只写 $y=2\sin(x+C)$，会遗漏这些可微拼接解。\end{remark}

\section{决赛：第 15 届（2024 年）}

本节选入7道课程相关题；来源序号保留原题解编排，未选入题目不重新补号。

\begin{example}[来源序号2]\label{ex:fin15-high:2}\label{pr:fin15-high:2}

设 $f(x) = \mathrm{e}^{-x^{2}}$，则极限
$\lim_{n\rightarrow+\infty}\frac{1}{n}\sum_{k=1}^{n}(-1)^{k-1}f\left(\frac{k}{n}\right) = \underline{\qquad}$。

\par\href{https://github.com/lzxcr/cmc-mathematics/blob/64916fa1fdbc3e444a652eca204d5669058ee321/src/problems/finals/15/02.md}{题目与来源条目}

\end{example}

\begin{solve}\label{ans:fin15-high:2}

将相邻项配对。由 $f$ 在 $[0,1]$ 上一致连续，记连续模为 $\omega_f$，则
\[
\left|\frac1n\sum_{k=1}^n(-1)^{k-1}f(k/n)\right|
\le\frac12\omega_f(1/n)+\frac{\|f\|_\infty}{n}\longrightarrow0.
\]
所求极限为 $\boxed0$。

连续模定义为 $\omega_f(\delta)=\sup\{|f(u)-f(v)|:u,v\in[0,1],|u-v|\le\delta\}$，一致连续性给出 $\omega_f(\delta)\to0$。每一对相邻项 $f((2j-1)/n)-f(2j/n)$ 的绝对值不超过 $\omega_f(1/n)$，共至多 $n/2$ 对；若 $n$ 为奇数，余下一项除以 $n$ 后至多为 $\|f\|_\infty/n$。本题还可用 $|f'|\le2$ 取得 $\omega_f(1/n)\le2/n$，直接得到整体为 $O(1/n)$。

\end{solve}

\begin{example}[来源序号3]\label{ex:fin15-high:3}\label{pr:fin15-high:3}

设 $f$ 是 $[0,1]$ 上的连续函数，$\int_{0}^{1}x^{k}f(x)\mathrm{d}x = a_{k}$（$k = 0,1,2$），则 $\int_{0}^{1}\mathrm{d}x\int_{0}^{x}\mathrm{d}y\int_{0}^{y}f(z)\mathrm{d}z = \underline{\qquad}$。

\par\href{https://github.com/lzxcr/cmc-mathematics/blob/64916fa1fdbc3e444a652eca204d5669058ee321/src/problems/finals/15/03.md}{题目与来源条目}

\end{example}

\begin{solve}\label{ans:fin15-high:3}

交换积分次序，得到
\[
\int_0^1f(z)\left(\int_z^1\int_y^1\mathrm dx\,\mathrm dy\right)\mathrm dz
=\frac12\int_0^1(1-z)^2f(z)\,\mathrm dz
=\boxed{\frac{a_0-2a_1+a_2}{2}}.
\]

原积分区域是 $0\le z\le y\le x\le1$。固定 $z$ 后，$(x,y)$ 在 $z\le y\le x\le1$ 的直角三角形中，面积为 $(1-z)^2/2$，所以换序后成为 $f(z)$ 乘该面积的一重积分。$f$ 连续，整个区域有界，被积函数绝对可积，因此换序合法。将 $(1-z)^2$ 展开为 $1-2z+z^2$，再分别代入三个已知矩 $a_0,a_1,a_2$，即得答案。

\end{solve}

\begin{example}[来源序号4]\label{ex:fin15-high:4}\label{pr:fin15-high:4}

微分方程 $xy' + y = \cos x$ 定义在 $(-\infty,+\infty)$ 上的全局解为 $\underline{\qquad}$。

\par\href{https://github.com/lzxcr/cmc-mathematics/blob/64916fa1fdbc3e444a652eca204d5669058ee321/src/problems/finals/15/04.md}{题目与来源条目}

\end{example}

\begin{solve}\label{ans:fin15-high:4}

方程即 $(xy)'=\cos x$。全局解满足 $xy=\sin x+C$，代入 $x=0$ 知 $C=0$。因此唯一全局解为
\[
\boxed{y(x)=\begin{cases}\sin x/x,&x\ne0,\\1,&x=0.\end{cases}}
\]
该函数在原点可解析延拓，确实满足方程。

若先只在 $x\ne0$ 的区间解方程，会得到 $(\sin x+C)/x$，但非零常数会在零点产生极点，无法成为全局可微解。全局函数满足 $(xy)'=\cos x$，故积分常数在整条实轴上统一且由 $0\cdot y(0)=0$ 确定。连续性要求 $y(0)=\lim\sin x/x=1$；展开 $y(x)=1-x^2/6+O(x^4)$ 又给出 $y'(0)=0$。直接代入零点的原方程为 $y(0)=1$，与延拓值一致。

\end{solve}

\begin{example}[来源序号7]\label{ex:fin15-high:7}\label{pr:fin15-high:7}

设 $f,g$ 是 $[0,1]$ 上有正下界的 Riemann 可积函数，且 $\int_{0}^{1}f(x)\mathrm{d}x = \int_{0}^{1}g(x)\mathrm{d}x$，$\delta = \int_{0}^{1}|f(x) - g(x)|\mathrm{d}x > 0$。
（1）证明：$\int_{0}^{1}f(x)\mathrm{d}x < \int_{0}^{1}\frac{f^{2}(x)}{g(x)}\mathrm{d}x$。
（2）若进一步假设 $f,g$ 分段常值，$\int_{0}^{1}f(x)\mathrm{d}x = 1$，$\delta = \frac{3}{2}$，试给出在此种条件下 $\int_{0}^{1}\frac{f^{2}(x)}{g(x)}\mathrm{d}x$ 的下确界。

\par\href{https://github.com/lzxcr/cmc-mathematics/blob/64916fa1fdbc3e444a652eca204d5669058ee321/src/problems/finals/15/07.md}{题目与来源条目}

\end{example}

\begin{solve}\label{ans:fin15-high:7}

（1）记 $m=\int_0^1f=\int_0^1g>0$。恒等式与加权 Cauchy--Schwarz 不等式给出
\[
\int_0^1\frac{f^2}{g}-m=\int_0^1\frac{(f-g)^2}{g}
\ge\frac{\left(\int_0^1|f-g|\right)^2}{\int_0^1g}
=\frac{\delta^2}{m}>0.
\]
（2）整体 Cauchy--Schwarz 未利用 $f,g$ 总质量相同。把正、负差分开，两边各携带质量 $\delta/2$，即可保留这一约束。令 $E=\{f\ge g\}$、$B=\int_Eg$。因两函数均归一化且 $\delta=3/2$，有 $\int_Ef=B+3/4$，从而 $0<B<1/4$。在 $E$ 及其补集分别使用 Cauchy--Schwarz，不等式变为
\[
\int_0^1\frac{f^2}{g}
\ge\frac{(B+3/4)^2}{B}+\frac{(1/4-B)^2}{1-B}
=1+\frac9{16B(1-B)}>4.
\]
为说明下界最优，取 $0<B<1/4$，在 $[0,1/2)$ 与 $[1/2,1]$ 上分别令
\[
g=(2B,\,2-2B),\qquad f=(2B+3/2,\,1/2-2B).
\]
这些分段常值函数满足全部条件，且上式第一步取等。令 $B\uparrow1/4$，积分趋于 $4$。故下确界为 $\boxed4$，但不取得。

正负差的积分相等，因为 $\int(f-g)=0$，而它们之和是 $\delta$，所以各为 $3/4$。在补集上 $\int f=1-(B+3/4)=1/4-B>0$，严格正性保证 $B<1/4$；正差非零也保证 $B>0$。函数 $B(1-B)$ 在 $(0,1/4)$ 严格递增且小于 $3/16$，所以 $1+9/[16B(1-B)]>4$。构造中的每个固定 $B$ 都使两个函数有正下界，但当 $B\uparrow1/4$ 时 $f$ 的第二段值趋零，极限函数不再满足正下界条件，这就是下确界不取得的原因。

\end{solve}

\begin{remark}[方法与条件]当 $f,g$ 是概率密度时，$\int f^2/g-1$ 是 Pearson $\chi^2$ 散度，$\frac12\int|f-g|$ 是全变差距离。把空间按 $f\ge g$ 分成两块，是将分布压缩成二点分布；Cauchy--Schwarz 说明这种压缩只会降低散度。

这与条件平均降低凸函数积分的机制相同，见方法卷：条件期望与投影。\end{remark}

\begin{example}[来源序号10，低年级组]\label{ex:fin15-low:9}\label{pr:fin15-low:9}

设 $0 < a_{1} < a_{2}$，
$a_{n+1} = \ln\left(\mathrm{e}^{a_{n}} + \frac{a_{n} - a_{n-1}}{a_{n-1}}\right)$，$n = 2,3,\cdots$。
证明：数列 $\left\{a_{n}\right\}$ 单调递增且有界。

\par\href{https://github.com/lzxcr/cmc-mathematics/blob/64916fa1fdbc3e444a652eca204d5669058ee321/src/problems/finals/15/10.md}{题目与来源条目}

\end{example}

\begin{solve}\label{ans:fin15-low:9}

由递推式归纳得 $a_{n+1}>a_n>0$。有界性应控制总增量；利用 $\log(1+t)<t$ 及 $e^{a_n}>a_n$，把单步增量压到一个可望远镜求和的差：
\[
0<a_{n+1}-a_n
=\log\left(1+\frac{a_n-a_{n-1}}{e^{a_n}a_{n-1}}\right)
<\frac1{a_{n-1}}-\frac1{a_n}.
\]
求和立即得到 $a_{n+1}<a_2+1/a_1$。所以数列严格递增且有界。

归纳起点为 $a_2>a_1>0$。若 $a_n>a_{n-1}>0$，对数内的数严格大于 $e^{a_n}$，所以 $a_{n+1}>a_n$，递推同时始终有定义。对增量的界从 $n=2$ 加到 $N$，左端为 $a_{N+1}-a_2$，右端恰为 $1/a_1-1/a_N<1/a_1$。因此得到与 $N$ 无关的上界 $a_2+1/a_1$；仅知道单步增量趋零并不足以推出有界，关键在于它们被可望远镜求和的正项控制。

\end{solve}

\begin{example}[来源序号12，低年级组]\label{ex:fin15-low:10}\label{pr:fin15-low:10}

设 $f,g$ 为 $(-\infty,+\infty)$ 上的连续函数且当 $x \neq 0$ 时，$xf(x) > 0$。证明方程 $x'' + f\left(x'\right) + g(x) = 0$ 没有非常数的周期解 $x(\cdot)$。

\par\href{https://github.com/lzxcr/cmc-mathematics/blob/64916fa1fdbc3e444a652eca204d5669058ee321/src/problems/finals/15/12.md}{题目与来源条目}

\end{example}

\begin{solve}\label{ans:fin15-low:10}

令 $G(x)=\int_0^xg(s)\,\mathrm ds$。任一解的能量
\[
E(t)=\frac12x'(t)^2+G(x(t))
\]
满足
\[
E'(t)=x'(t)(x''(t)+g(x(t)))=-x'(t)f(x'(t))\le0,
\]
且等号仅在 $x'(t)=0$ 时成立。若 $x$ 为周期解，$E$ 也周期；单调的周期函数必为常数，因此 $x'\equiv0$。故不存在非常数周期解。

若 $x$ 为周期 $T$ 的可微函数，对 $x(t+T)=x(t)$ 求导得 $x'(t+T)=x'(t)$，所以能量也以 $T$ 为周期。积分一个周期可更直接看出矛盾：
\[
0=E(T)-E(0)=-\int_0^T x'(t)f(x'(t))\,dt.
\]
被积函数连续非负，且在 $x'\ne0$ 时严格正，因此若某点 $x'\ne0$，其附近一段积分就严格为正。只能有 $x'\equiv0$，故解为常数。整个论证不要求势能 $G$ 有正下界，也不要求 $f,g$ 可微。

\end{solve}

\begin{remark}[方法与条件]能量严格耗散排除了闭合轨道。这里不需要 $g$ 的符号条件，也不需要势能有下界，因为论证只比较一个假定周期内的能量变化。

构造能量的出发点是先把回复力 $g$ 积分成势能，使求导后只剩阻尼项；见\hyperref[app:dynamics-energy]{方法卷：能量与耗散}。\end{remark}

\begin{example}[来源序号16，高年级组选做]\label{ex:fin15-high:13}\label{pr:fin15-high:13}

设 $0 = x_{0} < x_{1} < \dots < x_{N} = 1$ 是区间 $[0,1]$ 的等距剖分点，即 $x_{j} = jh$，$j = 0,1,\dots,N$，$h = \frac{1}{N}$。令
$V_{h} = \left\{v \mid \text{对于任何 } 1 \leq j \leq N,\ v|_{I_{j}} \in P_{k}\left(I_{j}\right)\right\}$，
其中 $P_{k}\left(I_{j}\right)$ 是定义在区间 $I_{j} = \left[x_{j-1},x_{j}\right)$ 上次数不超过 $k$ 次的多项式全体。设函数 $u(x) \in C^{(k+2)}[0,1]$，$u(x)$ 在空间 $V_{h}$ 的 $L^{2}$ 投影 $u_{h} \in V_{h}$ 满足如下关系
$\|u - u_{h}\| \leq \|u - v\|, \quad \forall v \in V_{h}$，
其中范数为通常的 $L^{2}$ 范数。
（1）证明 $\|u - u_{h}\| \leq Ch^{k+1}$，并给出常数 $C$ 与 $u(x)$ 导数的依赖关系。
（2）假设函数 $\varphi(x) \in C^{(k+2)}[0,1]$，证明
$\left|\int_{0}^{1}\left(u(x) - u_{h}(x)\right)\varphi(x)\mathrm{d}x\right| \leq Ch^{2k+2}$，
并给出常数 $C$ 与 $u(x)$ 及 $\varphi(x)$ 导数的依赖关系。

\par\href{https://github.com/lzxcr/cmc-mathematics/blob/64916fa1fdbc3e444a652eca204d5669058ee321/src/problems/finals/15/16.md}{题目与来源条目}

\end{example}

\begin{solve}\label{ans:fin15-high:13}

记 $M_u=\|u^{(k+1)}\|_\infty$。在每个小区间 $I_j$ 上取 $u$ 在左端点的 $k$ 次 Taylor 多项式，拼成 $v_h\in V_h$。Taylor 余项估计给出
\[
\|u-v_h\|_2^2
\le\frac{M_u^2}{((k+1)!)^2}\,N\int_0^ht^{2k+2}\,\mathrm dt
=\frac{M_u^2h^{2k+2}}{((k+1)!)^2(2k+3)}.
\]
利用最佳逼近性质，得
\[
\|u-u_h\|_2\le\frac{M_u}{(k+1)!\sqrt{2k+3}}h^{k+1}.
\]
将 $v=u_h+tw_h$ 代入最小性，对实参数 $t$ 在零点求导，得 $\langle u-u_h,w_h\rangle=0$（$w_h\in V_h$）。对 $\varphi$ 作相同的局部 Taylor 逼近 $\varphi_h$，于是
\[
|\langle u-u_h,\varphi\rangle|
=|\langle u-u_h,\varphi-\varphi_h\rangle|
\le\frac{\|u^{(k+1)}\|_\infty\|\varphi^{(k+1)}\|_\infty}
{((k+1)!)^2(2k+3)}h^{2k+2}.
\]

第二问的额外阶数来自正交性，而不是再次提高对 $u$ 的逐点近似精度。令 $e=u-u_h$，最小性给出 $\|e-tw_h\|_2^2=\|e\|_2^2-2t\langle e,w_h\rangle+t^2\|w_h\|_2^2$ 在 $t=0$ 取最小，因此线性系数为零。取 $w_h=\varphi_h$，消掉可由离散空间精确表示的测试函数部分，再将两个各为 $O(h^{k+1})$ 的误差作 Cauchy--Schwarz 估计，就得到 $O(h^{2k+2})$。所列常数与网格数 $N$ 无关，只依赖固定的 $k$ 和相应导数范数。

\end{solve}

\begin{remark}[方法与条件]第二问利用正交性把误差同时作用于两个逼近余项，因此由 $h^{k+1}$ 提升到 $h^{2k+2}$。这是一种对偶性思想；题设的 $C^{k+2}$ 正则性实际上可放宽为 $C^{k+1}$。

最佳逼近与正交性的等价证明及其统一作用，见方法卷：Gram 矩阵与投影。\end{remark}

\section{决赛：第 16 届（2025 年）}

本节选入8道课程相关题；来源序号保留原题解编排，未选入题目不重新补号。

\begin{example}[来源序号2]\label{ex:fin16-high:2}\label{pr:fin16-high:2}

二元函数 $f(x,y) = x^2 + xy + y^2 - 3y$ 的极小值点为 $\underline{\qquad}$。

\par\href{https://github.com/lzxcr/cmc-mathematics/blob/64916fa1fdbc3e444a652eca204d5669058ee321/src/problems/finals/16/02.md}{题目与来源条目}

\end{example}

\begin{solve}\label{ans:fin16-high:2}

解驻点方程 $2x+y=0$、$x+2y=3$ 得 $(-1,2)$。Hessian 为 $\begin{pmatrix}2&1\\1&2\end{pmatrix}>0$，所以唯一极小值点为 $\boxed{(-1,2)}$。

正定性可由顺序主子式 $2>0$、$\det H=3>0$ 判断；也可以直接配方以同时得到全局结论。令 $u=x+1,v=y-2$，则
\[
f(x,y)=-3+u^2+uv+v^2=-3+(u+v/2)^2+3v^2/4.
\]
后两项非负，且同时为零当且仅当 $u=v=0$。所以 $(-1,2)$ 不仅是局部极小值点，也是唯一全局最小值点，最小值为 $-3$。

\end{solve}

\begin{example}[来源序号3]\label{ex:fin16-high:3}\label{pr:fin16-high:3}

二重积分 $\int_0^1 \mathrm{d}x \int_x^1 f(x,y) \mathrm{d}y$ 转换为极坐标形式为 $\underline{\qquad}$。

\par\href{https://github.com/lzxcr/cmc-mathematics/blob/64916fa1fdbc3e444a652eca204d5669058ee321/src/problems/finals/16/03.md}{题目与来源条目}

\end{example}

\begin{solve}\label{ans:fin16-high:3}

积分区域为 $0\le x\le y\le1$。以 $x=\rho\cos\theta$、$y=\rho\sin\theta$ 换元，得到
\[
\boxed{\int_{\pi/4}^{\pi/2}\mathrm d\theta
\int_0^{1/\sin\theta}f(\rho\cos\theta,\rho\sin\theta)\,\rho\,\mathrm d\rho}.
\]

区域在第一象限，且 $y\ge x$，所以极角从直线 $y=x$ 的 $\pi/4$ 到正 $y$ 轴的 $\pi/2$。沿固定极角的射线，从原点进入区域，直到上边界 $y=1$，即 $\rho\sin\theta=1$，所以径向上界为 $1/\sin\theta$。这一角度范围内 $x\le y$，故无需再另加 $x\le1$ 的限制。极坐标的 Jacobian 为 $\rho$，它必须与函数值一起保留在内积分中；换元在原函数可积的通常条件下成立。

\end{solve}

\begin{example}[来源序号4]\label{ex:fin16-high:4}\label{pr:fin16-high:4}

微分方程 $y' = 1 + x \mathrm{e}^{y}$ 的通解为 $\underline{\qquad}$。

\par\href{https://github.com/lzxcr/cmc-mathematics/blob/64916fa1fdbc3e444a652eca204d5669058ee321/src/problems/finals/16/04.md}{题目与来源条目}

\end{example}

\begin{solve}\label{ans:fin16-high:4}

乘以 $e^{-y}$ 后，非线性项化为 $x$，而 $e^{-y}y'$ 是 $-e^{-y}$ 的导数。令 $u=e^{-y}>0$，原方程化为 $u'+u=-x$。因此
\[
\boxed{y=-\log(1-x+Ce^{-x})},
\]
定义在 $1-x+Ce^{-x}>0$ 的任一连通区间上。

链式法则给出 $u'=-e^{-y}y'=-u-x$。乘以 $e^x$ 后，$(e^xu)'=-xe^x$，积分得 $e^xu=(1-x)e^x+C$，故 $u=1-x+Ce^{-x}$。由于 $u$ 最初定义为 $e^{-y}$，它必须严格为正，不能把对数参数为零或负的点纳入实解区间。反过来，在任何正值连通区间上由 $y=-\ln u$ 求导，$y'=-u'/u=1+x/u=1+xe^y$，确认没有引入额外解。

\end{solve}

\begin{example}[来源序号7]\label{ex:fin16-high:7}\label{pr:fin16-high:7}

设 $f : [a,b] \rightarrow [a,b]$ 连续，$x_0 \in [a,b]$，$x_{n+1} = f(x_n)\ (n \geq 0)$。证明：
（1）若 $\lim_{n \rightarrow +\infty} (x_{n+1} - x_n) = 0$，则数列 $\{x_n\}$ 收敛。
（2）若 $\lim_{n \rightarrow +\infty} (x_{n+2025} - x_n) = 0$，则 $\{x_n \mid n \geq 1\}$ 的聚点不超过 2025 个。

\par\href{https://github.com/lzxcr/cmc-mathematics/blob/64916fa1fdbc3e444a652eca204d5669058ee321/src/problems/finals/16/07.md}{题目与来源条目}

\end{example}

\begin{solve}\label{ans:fin16-high:7}

（1）设 $\alpha=\liminf x_n$、$\beta=\limsup x_n$。由有界性，它们均有限。若 $\alpha<\beta$，增量 $x_{n+1}-x_n\to0$ 意味着每个 $c\in(\alpha,\beta)$ 都是聚点：数列在 $c$ 两侧无穷次穿越，每次跨越的步长趋于零。

对任一聚点 $c$，选 $x_{n_j}\to c$，则 $x_{n_j+1}\to c$，连续性给出 $f(c)=c$。所以 $f$ 在 $(\alpha,\beta)$ 上是恒等映射。另一方面，足够晚的一次从 $\alpha$ 附近到 $\beta$ 附近的穿越必有某项落在这个开区间内；此后迭代恒定，与 $\alpha<\beta$ 矛盾。因此 $\alpha=\beta$，数列收敛。

（2）令 $F=f^{\circ2025}$。对每个 $j=0,\dots,2024$，子列 $y_k=x_{2025k+j}$ 满足 $y_{k+1}=F(y_k)$ 且 $y_{k+1}-y_k\to0$。由（1），每个子列都收敛，故原序列的聚点至多 $2025$ 个。

穿越论证可明确构造子列：固定 $c\in(\alpha,\beta)$，在任意晚的时刻之后都存在一项低于 $c$、另一项高于 $c$。取一次从低到高的穿越中首个满足 $x_{m+1}\ge c$ 的下标，则 $x_m<c\le x_{m+1}$，从而 $|x_m-c|\le|x_{m+1}-x_m|\to0$。所以 $c$ 的确是聚点。第二问中连续映射的有限次复合仍连续且保持 $[a,b]$；任何聚点子列又总有无穷多个下标属于同一个模 $2025$ 的剩余类，所以不会出现这 $2025$ 个子列极限之外的新聚点。

\end{solve}

\begin{remark}[方法与条件]第一问利用了区间的顺序：在两端聚点之间来回穿越，步长趋零迫使整个中间区间成为聚点。第二问还给出各余数子列的极限 $\ell_j$ 满足 $f(\ell_j)=\ell_{j+1}$（指标模 $2025$），所以聚点组成一个周期整除 $2025$ 的周期轨道。\end{remark}

\begin{example}[来源序号9，高年级组选做]\label{ex:fin16-high:8}\label{pr:fin16-high:8}

记 $S_r = \{(x,y) \in \mathbb{R}^2 \mid x^2 + y^2 = r^2\}\ (r \geq 0)$，$C(\mathbb{R}^2)$ 为 $\mathbb{R}^2$ 上连续函数全体。对于 $f \in C(\mathbb{R}^2)$，记 $E_f$ 为这样的 $\xi \in \mathbb{R}^2$ 的全体：$\xi$ 是 $f$ 在某个 $S_r$（$r \geq 0$）上的最小值点。
（1）设 $g \in C(\mathbb{R}^2)$，证明：$E_g$ 是闭集。
（2）设 $E \subseteq \mathbb{R}^2$ 为闭集，且 $E \cap S_r \neq \emptyset\ (\forall r \geq 0)$。证明：存在 $g \in C(\mathbb{R}^2)$，使得 $E_g = E$。
（3）证明：存在 $g \in C(\mathbb{R}^2)$，使得 $E_g$ 为不连通集。
（4）是否存在 $g \in C(\mathbb{R}^2)$，使得 $E_g$ 的任何连通子集均为单点集？

\par\href{https://github.com/lzxcr/cmc-mathematics/blob/64916fa1fdbc3e444a652eca204d5669058ee321/src/problems/finals/16/09.md}{题目与来源条目}

\end{example}

\begin{solve}\label{ans:fin16-high:8}

（1）设 $p_j\in E_g$ 且 $p_j\to p$。若 $p=0$，结论显然。若 $r=|p|>0$，对任意 $q\in S_r$，取 $q_j=|p_j|q/r$，则 $g(p_j)\le g(q_j)$。令 $j\to\infty$ 得 $g(p)\le g(q)$，故 $p\in E_g$。

（2）要使指定闭集恰为每个圆周的最小值集，只须找一个恰在该闭集上为零的非负连续函数。取 $g(p)=\operatorname{dist}(p,E)$。它连续，且每个圆周上的最小值为零，恰在 $E$ 中取得，故 $E_g=E$。

（3）可取 $E=([0,\infty)\times\{0\})\cup(\{0\}\times[1,\infty))$，再使用（2）。

（4）可以。需要一个完全不连通集合，其半径的连续像却覆盖区间；Cantor 集到区间的连续满射正提供这个结构。再用极角记住参数，避免压缩掉原集合的拓扑。令 $C$ 为标准 Cantor 集，$h:C\to[0,1]$ 为连续满射（把三进制数字 $0,2$ 换成二进制数字 $0,1$）。定义
\[
E_n=\{(n+h(t))(\cos t,\sin t):t\in C\},\qquad E=\bigcup_{n=0}^\infty E_n.
\]
每个参数映射连续且单射：非零点的极角恢复 $t$，原点只对应 $n=t=0$。因此 $E_n$ 同胚于 $C$，完全不连通。

$E_n$ 位于半径 $[n,n+1]$ 的闭圆环内。相邻两集在公共边界圆上的候选点分别具有极角 $1$ 与 $0$，故互不相交。这个紧集族局部有限，故 $E$ 闭，且各 $E_n$ 在 $E$ 中既开又闭。因此 $E$ 也完全不连通。由 $h$ 满射，$E$ 与每个 $S_r$ 相交，最后用（2）构造 $g$。

距离函数的连续性由 $|\operatorname{dist}(p,E)-\operatorname{dist}(q,E)|\le|p-q|$ 保证；闭性则保证距离为零当且仅当点属于 $E$，两者在第二问中缺一不可。第三问两部分的距离至少为 $1$，所以在并集的相对拓扑中确实分离。第四问每个 $E_n$ 都是紧集，且任何有界区域只与有限多个 $E_n$ 相交，因此并集闭；删去某个 $E_n$ 后余下的局部有限紧集并仍闭，故 $E_n$ 在 $E$ 中也开。任何连通子集若遇到两个不同 $E_n$ 就会被这些相对开闭集分开，只能落在一个 $E_n$ 内，而 Cantor 集的连通子集只有单点。

\end{solve}

\begin{remark}[方法与条件]连续像不保持完全不连通：Cantor 集可连续满射到区间。本题利用极角保留 Cantor 参数，使平面集合仍完全不连通，同时半径映射覆盖整个非负实轴。\end{remark}

\begin{example}[来源序号10，低年级组]\label{ex:fin16-low:9}\label{pr:fin16-low:9}

设 $x_0 > y_0 > 0$，对于 $n \geq 0$，有
$x_{n+1} = \frac{1}{2(n+1)} \sum_{k=0}^{n} (x_k + y_k)$，$y_{n+1} = \left[ \prod_{k=0}^{n} (x_k y_k) \right]^{\frac{1}{2n+2}}$。
证明：（1）数列 $\{x_n - y_n\}$ 收敛于零；
（2）数列 $\{n(x_n - y_n)\}$ 收敛于一个正数。

\par\href{https://github.com/lzxcr/cmc-mathematics/blob/64916fa1fdbc3e444a652eca204d5669058ee321/src/problems/finals/16/10.md}{题目与来源条目}

\end{example}

\begin{solve}\label{ans:fin16-low:9}

把相邻两式分别相减、相除，便能消去累积和与累积积，将问题还原为一步迭代：
\[
x_{n+1}=x_n-\frac{x_n-y_n}{2n+2},\qquad
y_{n+1}=y_n\left(\frac{x_n}{y_n}\right)^{1/(2n+2)}.
\]
加权算术--几何平均不等式归纳给出 $x_n>y_n>0$，$x_n$ 严格递减、$y_n$ 严格递增。故它们分别收敛于 $X\ge Y>0$。若 $X>Y$，则
\[
\log\frac{y_{n+1}}{y_n}=\frac1{2n+2}\log\frac{x_n}{y_n}
\]
的右边之和发散，而左边之和收敛于 $\log(Y/y_0)$，矛盾。所以 $X=Y$，$\delta_n=x_n-y_n\to0$。

为求速度，先证明一个较粗但可求和的衰减率，再控制 $n\delta_n$ 的乘法增量。记 $a=1/(2n+2)$、$t=\delta_n/y_n$。对 $0<a\le1/2$，Taylor 余项及凹性给出一致估计
\[
0\le1+at-(1+t)^a\le\frac a2t^2.
\]
因此存在常数 $C$，对 $n\ge1$ 有
\[
\frac n{n+1}\delta_n\le\delta_{n+1}
\le\frac n{n+1}\delta_n+\frac C{n+1}\delta_n^2.
\]
因 $\delta_n\to0$，充分大时 $\delta_{n+1}\le(1-1/(2n+2))\delta_n$，连乘得到 $\delta_n=O(n^{-1/2})$。再令 $b_n=n\delta_n>0$，上式给出
\[
1\le\frac{b_{n+1}}{b_n}\le1+\frac{C\delta_n}{n},
\qquad\sum_{n\ge1}\frac{\delta_n}{n}<\infty.
\]
故 $b_n$ 单调且有界，收敛于一个正数。这也得到更精确的 $\delta_n\sim c/n$。

一步递推式中的算术平均可写成 $x_{n+1}=(1-a)x_n+ay_n$，几何平均则为 $y_{n+1}=y_n^{1-a}x_n^a$，其中 $a=1/(2n+2)$。令 $t=(x_n-y_n)/y_n$，二者差为
\[
\delta_{n+1}=(1-2a)\delta_n+y_n[1+at-(1+t)^a].
\]
因 $-[(1+t)^a]''=a(1-a)(1+t)^{a-2}\le a$（$t\ge0$），Taylor 余项就在 $0$ 与 $at^2/2$ 之间，且 $y_n\ge y_0>0$，于是得到统一的二次误差界。最后 $b_n$ 的比例上界连乘不超过常数倍的 $\exp(C\sum\delta_n/n)<\infty$；又 $b_n\ge b_1>0$，所以极限确为正数。

\end{solve}

\begin{example}[来源序号12，低年级组]\label{ex:fin16-low:10}\label{pr:fin16-low:10}

在 $|x|,|y|<1$ 内求微分关系
\[
\frac{\mathrm dy}{\sqrt{1-y^4}}=\pm\frac{\mathrm dx}{\sqrt{1-x^4}}
\]
的全部局部实解，并核验代数曲线
\[
x^2+y^2+c^2x^2y^2=c^2+2xy\sqrt{1-c^4}\qquad(|c|\le1)
\]
的局部图像分支与这一微分关系的联系。再求
\[
\frac{\mathrm dy}{\sqrt{1+y^4}}=\pm\frac{\mathrm dx}{\sqrt{1+x^4}}
\]
的全部局部实解。这里局部解定义在非空开区间上，符号在每个连通解区间内固定。

\par\href{https://github.com/lzxcr/cmc-mathematics/blob/64916fa1fdbc3e444a652eca204d5669058ee321/src/problems/finals/16/12.md}{题目与来源条目}

\end{example}

\begin{solve}\label{ans:fin16-low:10}

分母只依赖各自变量，故先积分便能保留符号和任意积分常数；先平方消去根式则会混淆分支。设
\[
F_-(t)=\int_0^t\frac{\mathrm ds}{\sqrt{1-s^4}}\quad(|t|<1),
\qquad
F_+(t)=\int_0^t\frac{\mathrm ds}{\sqrt{1+s^4}}\quad(t\in\mathbb R).
\]
则全部局部实解分别是
\[
\boxed{F_-(y)=\sigma F_-(x)+C},\qquad
\boxed{F_+(y)=\sigma F_+(x)+C},\qquad\sigma\in\{1,-1\},
\]
其中 $C\in\mathbb R$。准确地说，取 $x$ 所在的非空开区间，使 $x$ 属于相应 $F_\pm$ 的定义域，且 $\sigma F_\pm(x)+C$ 属于其值域，再令
\[
y(x)=F_\pm^{-1}\bigl(\sigma F_\pm(x)+C\bigr).
\]
两个 $F_\pm$ 都严格递增且导数处处为正，故反函数定理保证这些函数确为解；反过来，对微分关系积分便得到上述形式。若将题面的 $\pm$ 理解为逐点任选，导数的 Darboux 性质及 $y'\ne0$ 仍迫使符号在连通区间内固定。

下面核验题设代数式。令 $d=\sqrt{1-c^4}$，$|c|\le1$。对
\[
x^2+y^2+c^2x^2y^2=c^2+2dxy
\]
隐式求导，并利用此关系整理，得
\[
\bigl(x(1+c^2y^2)-dy\bigr)^2=c^2(1-y^4),\qquad
\bigl(y(1+c^2x^2)-dx\bigr)^2=c^2(1-x^4).
\]
当 $c\ne0$ 且 $|x|,|y|<1$ 时，第二个平方恒等式保证对 $y$ 的偏导不为零，故隐函数定理给出局部图像。两式相除得到 $(y')^2=(1-y^4)/(1-x^4)$，即所求微分关系；$c=0$ 时直接得到 $y=x$。

将代数式右端的 $2dxy$ 换成 $-2dxy$，便得到原曲线经 $y\mapsto-y$ 反射后的解曲线；例如 $c=0$ 时得到 $y=-x$。第二个方程也有相应代数形式
\[
x^2+y^2-c^2x^2y^2=c^2\pm2xy\sqrt{1+c^4},
\]
其中 $c\in\mathbb R$。其局部图像分支可同样核验；完整定义域和全部积分常数仍以上面的 $F_+$ 形式为准。

代数核验中令 $A=x(1+c^2y^2)-dy$、$B=y(1+c^2x^2)-dx$。隐式求导给出 $A+By'=0$，而两个平方恒等式给出 $A^2=c^2(1-y^4)$、$B^2=c^2(1-x^4)$。当 $c\ne0$ 时 $B\ne0$，所以 $y'=-A/B$ 有定义，平方后得到所需比例；连续性和非零性使其符号在局部支上固定。积分形式中任意常数也需配合非空定义区间：即使 $F_+$ 定义在全实轴，它的值域也有有限端点，因为其导数在无穷远为 $O(t^{-2})$，不能把反函数的自变量取遍任意实数。

\end{solve}

\begin{remark}[方法与条件]$F_-$ 是双纽线椭圆积分，其逆函数的加法公式产生上述代数关系。单取题面所列代数族并不包含 $y=-x$；第二个微分关系还有 $y=1/x$（$x\ne0$）这样的局部解，也不能由上面有限参数 $c$ 的代数族表示。积分表示则保留了符号、定义域与全部可行的积分常数。核验一族曲线满足方程，与证明已经穷尽全部局部解，是两个不同的任务。\end{remark}

\begin{example}[来源序号16，高年级组选做]\label{ex:fin16-high:13}\label{pr:fin16-high:13}

设 $L_0(x)=1$、$L_1(x)=x$，且
\[
L_{n+1}(x)=\frac{2n+1}{n+1}xL_n(x)-\frac{n}{n+1}L_{n-1}(x)
\qquad(n\ge1).
\]
记 $\delta_{mn}$ 为 Kronecker 记号。已知 Legendre 多项式满足
\[
\begin{aligned}
\int_{-1}^1L_m(x)L_n(x)\,\mathrm dx&=\frac{2}{2n+1}\delta_{mn},\\
L'_{n-1}(x)&=xL'_n(x)-nL_n(x),\\
L'_{n+1}(x)&=xL'_n(x)+(n+1)L_n(x).
\end{aligned}
\]
（1）证明 $(x^2-1)L'_n(x)=n(xL_n(x)-L_{n-1}(x))$（$n\ge1$）。

（2）证明
\[
\int_{-1}^1(1-x^2)L'_i(x)L'_j(x)\,\mathrm dx
=\frac{2i(i+1)}{2i+1}\delta_{ij}\qquad(i,j\ge1).
\]
（3）设 $s_j$、$w_j$ 是 $n$ 点 Gauss--Legendre 求积的节点与权重，定义
\[
A_{ij}=L'_i(s_j)\quad(1\le i,j\le n),\qquad
W=\operatorname{diag}\bigl((1-s_j^2)w_j\bigr)_{j=1}^n.
\]
证明 $B=AWA^{\mathrm T}$ 是对角矩阵。

\par\href{https://github.com/lzxcr/cmc-mathematics/blob/64916fa1fdbc3e444a652eca204d5669058ee321/src/problems/finals/16/16.md}{题目与来源条目}

\end{example}

\begin{solve}\label{ans:fin16-high:13}

（1）将 $L'_{n-1}=xL'_n-nL_n$ 乘以 $x$，再减去 $L'_n=xL'_{n-1}+nL_{n-1}$，即得结论；$n=1$ 时直接验证。

（2）对（1）求导，并再次用已知微分关系，得到 Legendre 方程
\[
\bigl((x^2-1)L'_i\bigr)'=i(i+1)L_i.
\]
分部积分，利用端点处 $1-x^2=0$ 及正交性，便有
\[
\int_{-1}^1(1-x^2)L'_iL'_j\,\mathrm dx
=i(i+1)\int_{-1}^1L_iL_j\,\mathrm dx
=\frac{2i(i+1)}{2i+1}\delta_{ij}.
\]
（3）矩阵的每一行是一个导数多项式在全部节点处的取值，因此
\[
b_{ij}=\sum_{r=1}^nw_r(1-s_r^2)L'_i(s_r)L'_j(s_r).
\]
$n$ 点 Gauss 求积对次数不超过 $2n-1$ 的多项式精确。当 $i\ne j$ 时，多项式 $(1-x^2)L'_iL'_j$ 的次数为 $i+j\le2n-1$，故上式恰等于（2）中的积分，因而为零。于是 $B$ 为对角阵。

第二问中的分部积分为
\[
\int_{-1}^1(1-x^2)L'_iL'_j\,dx
=[(1-x^2)L'_iL_j]_{-1}^1-\int_{-1}^1[(1-x^2)L'_i]'L_j\,dx.
\]
第一项因端点因子为零而消失，第二项用微分方程变为 $i(i+1)\int L_iL_j$。第三问只需要处理非对角元；若误把求积精确性用于 $i=j=n$，被积多项式次数为 $2n$，已超出 $2n-1$ 的保证范围。对 $i\ne j$，最大可能和为 $n+(n-1)$，恰在保证范围内，所以结论是对角化，而无需额外声称最后一个对角元也等于连续积分值。

\end{solve}

\begin{remark}[方法与条件]关键是求积精度与多项式次数的配合：正交性本身不保证任意离散取样后仍然正交。这里所有非对角项都在精确度范围内，只有 $(n,n)$ 项达到 $2n$ 次，不能直接照搬积分值。

这最后一项也可以求出。令
\[
q(x)=(1-x^2)(L'_n(x))^2+n^2L_n(x)^2.
\]
两项的最高次项相消，故 $\deg q\le2n-1$；而节点是 $L_n$ 的零点，所以 $q(s_r)=(1-s_r^2)(L'_n(s_r))^2$。再次使用 Gauss 精确性与正交性，得到
\[
b_{nn}=\int_{-1}^1q(x)\,\mathrm dx
=\frac{2n(n+1)+2n^2}{2n+1}=2n.
\]
其余对角元为 $b_{ii}=2i(i+1)/(2i+1)$（$i<n$）。这体现了一个常用思路：在求积节点上不改变取值，却通过加上节点多项式的倍数降低次数，从而回到精确求积的范围。精确性的来源与一般函数的误差控制，见\hyperref[app:estimate-gauss]{方法卷：Gauss 求积}。

按 $A_{ij}=L'_i(s_j)$ 的约定，多项式对应行，故离散 Gram 矩阵是 $AWA^{\mathrm T}$；若把节点排在行上，则相应写成 $A^{\mathrm T}WA$。\end{remark}

\section{决赛：第 17 届（2026 年）}

本节选入7道课程相关题；来源序号保留原题解编排，未选入题目不重新补号。

\begin{example}[来源序号2]\label{ex:fin17-high:2}\label{pr:fin17-high:2}

设 $f(x,y)=x^2-y^2+9x^2y-3y^3$，$B_R=\{(x,y)\mid x^2+y^2\le R^2\}$，

\[
L_R:=\max_{(x,y)\in B_R}f(x,y),\qquad \ell_R:=\min_{(x,y)\in B_R}f(x,y),
\]

则

\[
\lim_{R\to+\infty}\frac{L_R+\ell_R}{R^2}=\underline{\hspace{2cm}}.
\]

\par\href{https://github.com/lzxcr/cmc-mathematics/blob/64916fa1fdbc3e444a652eca204d5669058ee321/src/problems/finals/17/02.md}{题目与来源条目}

\end{example}

\begin{solve}\label{ans:fin17-high:2}

缩放到单位圆盘，写
\[
f(Ru,Rv)=R^3H(u,v)+R^2Q(u,v),\qquad
H=9u^2v-3v^3,\quad Q=u^2-v^2.
\]
用极坐标 $u+iv=re^{i\theta}$，则 $H=3r^3\sin3\theta$、$Q=r^2\cos2\theta$。$H$ 的最大值为 $3$，其三个最大值点位于 $r=1$、$\theta=\pi/6,5\pi/6,3\pi/2$，相应 $Q$ 值为 $1/2,1/2,-1$。$H$ 的最小值为 $-3$，其最小值点上的 $Q$ 值为 $-1,1/2,1/2$。

对紧集上的连续函数，有一阶极值扰动公式
\[
\max(H+\varepsilon Q)=\max H+\varepsilon\max_{\operatorname{argmax}H}Q+o(\varepsilon)
\quad(\varepsilon\downarrow0).
\]
事实上，设 $M=\max H$、$q=\max_{\operatorname{argmax}H}Q$，取扰动后的最大值点 $z_\varepsilon$，便有 $q\le[\max(H+\varepsilon Q)-M]/\varepsilon\le Q(z_\varepsilon)$。比较函数值先得 $H(z_\varepsilon)\to M$；紧性保证其聚点都在 $\operatorname{argmax}H$，因此右端的上极限至多为 $q$。最小值同理。

取 $\varepsilon=1/R$，得到
\[
L_R=3R^3+\tfrac12R^2+o(R^2),\qquad
\ell_R=-3R^3-R^2+o(R^2).
\]
故所求极限为 $\boxed{-\tfrac12}$。

记 $\operatorname{argmax}H$ 为 $H$ 取最大值的点集。取其中使 $Q=q$ 的点 $z_0$，最大值性质给出 $H(z_\varepsilon)+\varepsilon Q(z_\varepsilon)\ge M+\varepsilon q$；结合 $H(z_\varepsilon)\le M$，得到正文的两边夹逼。又 $Q$ 在紧集上有界，所以 $0\le M-H(z_\varepsilon)\le\varepsilon(Q(z_\varepsilon)-q)=O(\varepsilon)$，聚点必落在主项最大值集上。最小值则选择主项最小值集上最小的 $Q=-1$。相加时 $R^3$ 项抵消，真正留下的是两个不同选择产生的 $R^2$ 系数 $1/2-1=-1/2$。

\end{solve}

\begin{remark}[方法与条件]主项决定极值点可能聚集的位置，下一阶项再在这个极值集合上作选择。因此不能先随意选一个主项极值点代入；主项有多个极值点时，次阶项决定最终答案。这是紧集上变分问题的一阶扰动原理。\end{remark}

\begin{example}[来源序号3]\label{ex:fin17-high:3}\label{pr:fin17-high:3}

设 $f(x)=\int_0^x\frac{1}{1+t^3}\,\mathrm{d}t$，$g(x):=f(x)-f\left(\frac{1}{x}\right)$，则 $g(2)=\underline{\hspace{2cm}}$。

\par\href{https://github.com/lzxcr/cmc-mathematics/blob/64916fa1fdbc3e444a652eca204d5669058ee321/src/problems/finals/17/03.md}{题目与来源条目}

\end{example}

\begin{solve}\label{ans:fin17-high:3}

记 $I=g(2)=\int_{1/2}^2(1+t^3)^{-1}\,\mathrm dt$。区间 $[1/2,2]$ 在倒数变换下保持不变；代换 $t=1/u$ 又得 $I=\int_{1/2}^2u(1+u^3)^{-1}\,\mathrm du$，两式相加：
\[
2I=\int_{1/2}^2\frac{\mathrm dt}{t^2-t+1}
=\frac2{\sqrt3}\left[\arctan\frac{2t-1}{\sqrt3}\right]_{1/2}^2
=\frac{2\pi}{3\sqrt3}.
\]
因此 $\boxed{g(2)=\pi/(3\sqrt3)}$。

倒数换元时 $dt=-du/u^2$，而 $(1+t^3)^{-1}=u^3/(1+u^3)$，反向积分后就得到 $u/(1+u^3)$。两式相加后分子为 $1+t$，利用 $1+t^3=(1+t)(t^2-t+1)$ 约去一次因子。再完成平方 $t^2-t+1=(t-1/2)^2+3/4$，原函数便是 $2\arctan((2t-1)/\sqrt3)/\sqrt3$。上端反正切为 $\pi/3$，下端为 $0$，最后还要除以两式相加产生的系数 $2$。

\end{solve}

\begin{example}[来源序号4]\label{ex:fin17-high:4}\label{pr:fin17-high:4}

微分方程

\[
tx''(t)-x'(t)-(4t-2)x(t)=0
\]

的通解为 $\underline{\hspace{2cm}}$。

\par\href{https://github.com/lzxcr/cmc-mathematics/blob/64916fa1fdbc3e444a652eca204d5669058ee321/src/problems/finals/17/04.md}{题目与来源条目}

\end{example}

\begin{solve}\label{ans:fin17-high:4}

试探 $x=e^{\lambda t}$ 时，左端除以 $e^{\lambda t}$ 为 $t(\lambda^2-4)+(2-\lambda)$，取 $\lambda=2$ 即恒为零。令 $x=e^{2t}v$，原方程化为
\[
tv''+(4t-1)v'=0.
\]
故 $v'=Cte^{-4t}$，积分并吸收常数得
\[
\boxed{x(t)=C_1e^{2t}+C_2(4t+1)e^{-2t}}.
\]
在不含 $0$ 的区间上，两解的 Wronskian 为 $-16t\ne0$，故构成基本解组。两解均光滑穿过 $0$；若要求跨过 $0$ 的二阶可微解，左右的函数值与二阶导数匹配也迫使同一对常数，所示形式仍完整。

在 $t\ne0$ 的区间，令 $w=v'$，则 $w'+(4-1/t)w=0$，解得 $w=Cte^{-4t}$；积分 $\int te^{-4t}dt=-(4t+1)e^{-4t}/16$，再把 $-1/16$ 吸收入常数，就得到第二个解。零点处不能除以 $t$，须另验匹配：两基本解在 $0$ 的值均为 $1$，一阶导数均为 $2$，二阶导数分别为 $4$ 与 $-12$。因此左右的 $x(0)=C_1+C_2$ 和 $x''(0)=4C_1-12C_2$ 相同就唯一确定同一对常数，说明二阶可微的跨零解也已完整包含。

\end{solve}

\begin{example}[来源序号7]\label{ex:fin17-high:7}\label{pr:fin17-high:7}

设 $f(x)$ 在 $(-\infty,+\infty)$ 内可导，数列 $\{a_n\}$ 收敛。定义数列 $\{x_n\}$：$x_{n+1}=f(x_n)+a_n$（$n\ge1$），其中 $x_1\in(-\infty,+\infty)$。

（1）若存在 $\alpha\in(0,1)$ 使得 $|f'(x)|\le\alpha$，证明：$\{x_n\}$ 收敛；

（2）若存在 $M>0$，使得 $|f'(x)|\le\frac{|x|+M}{|x|+M+1}$（$\forall x\in\mathbb R$），证明或举例否定：$\{x_n\}$ 收敛。

\par\href{https://github.com/lzxcr/cmc-mathematics/blob/64916fa1fdbc3e444a652eca204d5669058ee321/src/problems/finals/17/07.md}{题目与来源条目}

\end{example}

\begin{solve}\label{ans:fin17-high:7}

记 $a=\lim a_n$。

（1）映射 $F(x)=f(x)+a$ 是 $\mathbb R$ 上的压缩映射，故有唯一不动点 $\xi$。令 $e_n=|x_n-\xi|$、$d_n=|a_n-a|\to0$，则
\[
e_{n+1}\le\alpha e_n+d_n.
\]
$d_n$ 有界首先保证 $e_n$ 有界；再取上极限得到 $L\le\alpha L$，故 $L=0$，即 $x_n\to\xi$。

（2）结论仍成立。这里全局压缩常数可趋近 $1$，应先证明迭代有界，随后在该有界区间上恢复严格压缩。导数界给出
\[
|f(x)-f(0)|\le |x|-\log\left(1+\frac{|x|}{M+1}\right).
\]
具体地，在正、负半轴分别将 $\pm(f(x)-f(0))$ 与 $|x|-\log(1+|x|/(M+1))$ 比较，其差的单调性由中值定理和给定导数界得到；不要求 $f'$ 连续。

记 $C=\sup_n|f(0)+a_n|<\infty$，选择 $R\ge|x_1|$ 使 $\log(1+R/(M+1))\ge C$。因 $t-\log(1+t/(M+1))$ 在 $t\ge0$ 上递增，$|x_n|\le R$ 蕴含 $|x_{n+1}|\le R$，所以全部迭代都在 $[-R,R]$ 内。

极限映射 $F$ 也保持该闭区间，且在其中的 Lipschitz 常数不超过
\[
q=\frac{R+M}{R+M+1}<1.
\]
因此它在该区间有唯一不动点；将（1）的估计中的 $\alpha$ 换成 $q$，便得 $x_n$ 收敛。

第一问若 $D=\sup_nd_n$，反复展开得 $e_n\le\alpha^{n-1}e_1+D/(1-\alpha)$，所以取上极限前确有有限上界。也可固定 $\varepsilon>0$ 后取 $d_n\le\varepsilon$，得到尾部 $e_{N+k}\le\alpha^ke_N+\varepsilon/(1-\alpha)$，再依次令 $k\to\infty$、$\varepsilon\to0$。第二问设 $B(t)=t-\ln(1+t/(M+1))$，则 $B'(t)=(t+M)/(t+M+1)$；分别比较 $B(t)\pm[f(t)-f(0)]$ 以及 $B(t)\pm[f(-t)-f(0)]$，它们导数非负且在零点为零，便得到所用统一增长界。

\end{solve}

\begin{example}[来源序号9，低年级组]\label{ex:fin17-low:7}\label{pr:fin17-low:7}

记 $f_1(x)=|(x-1)(x-2)|$，定义 $f_{n+1}(x):=f_1(f_n(x))$，$E_n=\{x\mid f_n(x)=2\}$（$n\ge1$）。设 $a_n=\min E_n$，$b_n=\max E_n$。

（1）求 $a_n,b_n$ 的递推式；

（2）证明 $\lim_{n\to\infty}a_n$ 和 $\lim_{n\to\infty}b_n$ 存在并求其值。

\par\href{https://github.com/lzxcr/cmc-mathematics/blob/64916fa1fdbc3e444a652eca204d5669058ee321/src/problems/finals/17/09.md}{题目与来源条目}

\end{example}

\begin{solve}\label{ans:fin17-low:7}

记 $F(x)=|(x-1)(x-2)|$。由 $F(3-x)=F(x)$，所有迭代的原像集 $E_n$ 均关于 $3/2$ 对称，故 $a_n+b_n=3$。$E_1=\{0,3\}$。

又 $E_{n+1}=F^{-1}(E_n)$。$F\ge0$，故只需考虑 $E_n$ 中的非负目标值。对任意非负 $t$，$F(x)=t$ 的最大实根为
\[
G(t)=\frac{3+\sqrt{1+4t}}2,
\]
它随 $t$ 严格增加。因此最大根由 $t=b_n$ 产生，并有
\[
\boxed{b_{n+1}=\frac{3+\sqrt{1+4b_n}}2,\quad b_1=3},
\qquad
\boxed{a_{n+1}=\frac{3-\sqrt{13-4a_n}}2,\quad a_1=0}.
\]
这些集合非空且有限，所用极值均存在。

$G$ 将 $[3,2+\sqrt2]$ 映入自身，并在其中满足 $0<G'(t)\le1/\sqrt{13}<1$，其唯一不动点为 $2+\sqrt2$。故
\[
\boxed{b_n\to2+\sqrt2,\qquad a_n\to1-\sqrt2}.
\]

解 $F(x)=t$ 时，正号分支为 $x^2-3x+2=t$，较大根即 $G(t)$；负号分支若有根，则根只落在 $[1,2]$ 内，而 $G(t)\ge2$，所以它不可能给出更大的根。由 $G(3)>3$、$G(2+\sqrt2)=2+\sqrt2$ 及严格递增性，也可直接归纳 $b_n$ 递增且不超过 $2+\sqrt2$。极限方程 $b=G(b)$ 要求 $2b-3\ge0$，平方后有 $b=2\pm\sqrt2$，只有正号根属于已知区间；再由 $a_n=3-b_n$ 得到另一个极限。

\end{solve}

\begin{example}[来源序号10，高年级组选做]\label{ex:fin17-high:8}\label{pr:fin17-high:8}

设 $h$ 是 $[0,1]$ 上非负的 Lebesgue 可积函数，$A=\int_0^1h(x)\,\mathrm{d}x$。证明：

\[
\left(1+A^3\right)^{1/3}\le\int_0^1\left(1+h^3(x)\right)^{1/3}\mathrm{d}x\le1+A .
\]

\par\href{https://github.com/lzxcr/cmc-mathematics/blob/64916fa1fdbc3e444a652eca204d5669058ee321/src/problems/finals/17/10.md}{题目与来源条目}

\end{example}

\begin{solve}\label{ans:fin17-high:8}

对向量 $(1,h(x))\in\mathbb R^2$ 使用 $\ell^3$ 范数。积分三角不等式给出
\[
(1+A^3)^{1/3}
=\left\|\int_0^1(1,h(x))\,\mathrm dx\right\|_3
\le\int_0^1\|(1,h(x))\|_3\,\mathrm dx.
\]
另一方面，点态上 $(1+h^3)^{1/3}\le1+h$，积分即得上界 $1+A$。该估计也保证被积函数可积，无需假设 $h^3$ 可积。

函数 $t\mapsto(1+t^3)^{1/3}$ 在 $[0,\infty)$ 上严格凸，故左侧也可看作 Jensen 不等式，等号恰在 $h$ 几乎处处为常数时成立；右侧点态等号只在 $h=0$ 时成立，因而积分等号也恰在 $h=0$ 几乎处处时成立。

不使用向量范数也可验证左界。令 $\phi(t)=(1+t^3)^{1/3}$，则 $\phi'(t)=t^2/(1+t^3)^{2/3}$、$\phi''(t)=2t/(1+t^3)^{5/3}\ge0$，且导数严格递增，所以 $\phi$ 严格凸。它在 $A$ 处的切线下界为 $\phi(h)\ge\phi(A)+\phi'(A)(h-A)$，积分后线性项因 $\int(h-A)=0$ 消失。右界则来自 $(1+h)^3-(1+h^3)=3h+3h^2\ge0$，两边开立方并积分即可。

\end{solve}

\begin{example}[来源序号12，低年级组]\label{ex:fin17-low:10}\label{pr:fin17-low:10}

设常系数线性微分方程 $y''+py'+qy=f(x)$ 有两个互异的非零实特征值 $\lambda<\mu$，$f$ 为 $(-\infty,+\infty)$ 上的连续函数。证明该微分方程存在周期解的充分必要条件是 $f$ 为周期函数。

\par\href{https://github.com/lzxcr/cmc-mathematics/blob/64916fa1fdbc3e444a652eca204d5669058ee321/src/problems/finals/17/12.md}{题目与来源条目}

\end{example}

\begin{solve}\label{ans:fin17-low:10}

必要性直接来自 $f=y''+py'+qy$：周期解及其导数均具有同一个周期，所以 $f$ 周期。

反之，设 $f$ 的周期为 $T>0$。周期解等价于初值经过一个周期后回到自身，因此将方程写成一阶系统
\[
Y'=CY+F(t),\qquad C=\begin{pmatrix}0&1\\-q&-p\end{pmatrix},\quad
F(t)=\binom0{f(t)}.
\]
常数变易公式表明 $Y(T)=Y(0)$ 当且仅当
\[
(I-e^{TC})Y(0)=\int_0^Te^{(T-s)C}F(s)\,\mathrm ds.
\]
$e^{TC}$ 的特征值是 $e^{T\lambda},e^{T\mu}$，由于 $\lambda,\mu$ 均为非零实数，它们都不等于 $1$。所以上式唯一确定初值。对应解满足 $Y(T)=Y(0)$；而 $Y(t+T)$ 与 $Y(t)$ 满足同一初值问题，由唯一性知解为 $T$ 周期。

这还证明：对每个给定周期的连续强迫项，周期解唯一。

系统向量取 $Y=(y,y')^{\mathrm T}$。因为两个特征值互异，$C$ 可对角化，$e^{TC}$ 的两个特征值就是 $e^{T\lambda},e^{T\mu}$；它们不为 $1$，所以 $\det(I-e^{TC})=(1-e^{T\lambda})(1-e^{T\mu})\ne0$。解出初值后，线性系统在连续强迫项下全局存在唯一；平移函数 $Y(t+T)$ 的强迫项为 $F(t+T)=F(t)$，且在 $t=0$ 与 $Y(t)$ 初值相同，因此在整个实轴相同。该论证也覆盖常值强迫项，只需任取其一个正周期。

\end{solve}

\begin{remark}[方法与条件]一般线性周期系统的关键条件是 $1$ 不属于一个周期的演化矩阵的谱。若出现特征值 $1$，周期解的存在将取决于强迫项的相容条件；这正是共振现象。

这一周期初值映射及其非共振条件，见\hyperref[app:dynamics-flow]{方法卷：线性流}。\end{remark}

\part{非数学专业组高等数学真题选讲}

本编依据高校公开的试卷及参考答案核对数学表达，题意以规范化文字重新表述，解析由本选编重新推导。小题编号对应各节所链接原卷；未选入的小题不补号。原卷中的评分步骤不予照录。局部解与整体解、反常积分的次序、参数范围等易遗漏条件在解析中单独说明。本编共收录37道题，其中一题含多个小问仍计为一道；方法推广和反例均不另算真题。

\section{第十届初赛：2018年非数学类}
本节来源为华东师范大学公开的\href{https://math.ecnu.edu.cn/llin/pdf_files/nonmathans2018.pdf}{第十届初赛非数学类试卷及参考答案}。重点包括隐函数、定积分估计、空间区域与正项级数。

\begin{example}[原卷一（1）：相邻幂的差]\label{nm-2018-1a}
设 $0<\alpha<1$，求 $\displaystyle\lim_{n\to\infty}[(n+1)^\alpha-n^\alpha]$。
\end{example}
\begin{solve}
把数列看作函数 $g(t)=t^\alpha$ 在长度为1的区间上的增量。由拉格朗日中值定理，存在 $\xi_n\in(n,n+1)$ 使
\[
 (n+1)^\alpha-n^\alpha=\alpha\xi_n^{\alpha-1}.
\]
因为 $\alpha-1<0$，有 $0<\alpha\xi_n^{\alpha-1}\leq\alpha n^{\alpha-1}\to0$，故极限为0。更准确地说，$\xi_n/n\to1$，从而原差与 $\alpha n^{\alpha-1}$ 等价。这同时说明结论依赖参数范围：$\alpha=1$ 时差恒为1，$\alpha>1$ 时差趋于无穷大。

中值点虽然随 $n$ 改变，但始终满足 $1<\xi_n/n<1+1/n$，故夹逼得到比值趋于 $1$。若要直接得到两侧界，由 $t^{\alpha-1}$ 递减可写成 $\alpha(n+1)^{\alpha-1}\le(n+1)^\alpha-n^\alpha\le\alpha n^{\alpha-1}$；两边同趋零，且除以 $\alpha n^{\alpha-1}$ 后同趋 $1$，便同时得到极限与更细的渐近式。

\end{solve}
\begin{remark}
对两个都趋于无穷大的量直接分别取极限没有意义。作差之后应把它转化为导数控制的增量；只有在相对误差已经受控时，才可以在相减式中使用等价替换。
\end{remark}

\begin{example}[原卷一（2）：参数与隐函数]\label{nm-2018-1b}
曲线由 $x=t+\cos t$ 和 $e^y+ty+\sin t=1$ 确定，求 $t=0$ 处的切线。
\end{example}
\begin{solve}
当 $t=0$ 时，$x=1$、$e^y=1$，所以 $y=0$。记 $F(t,y)=e^y+ty+\sin t-1$，则 $F_y(0,0)=1\ne0$。隐函数定理保证在该点附近可把 $y$ 写成 $t$ 的可微函数。对等式求导得到
\[
 (e^y+t)y'(t)+y+\cos t=0,
 \qquad y'(0)=-1.
\]
另一方面 $x'(t)=1-\sin t$，故 $x'(0)=1$。因此 $dy/dx=-1$，切线方程为 $y=1-x$。这里求导的自变量起初是 $t$；只有除以非零的 $x'(0)$ 后，才得到平面曲线相对于 $x$ 的斜率。

隐式求导时，乘积 $ty(t)$ 的导数是 $y(t)+ty'(t)$，指数项的导数是 $e^{y(t)}y'(t)$，所以两项含 $y'$ 的部分应合并为 $(e^y+t)y'$。参数曲线在 $t=0$ 的速度向量为 $(x',y')=(1,-1)$，不为零，故切线也可写成 $(x,y)=(1,0)+s(1,-1)$；消去 $s$ 后正是 $x+y=1$。

\end{solve}

\begin{example}[原卷一（3）：反双曲函数型积分]\label{nm-2018-1c}
计算 $\displaystyle\int\frac{\ln(x+\sqrt{1+x^2})}{(1+x^2)^{3/2}}\,dx$。
\end{example}
\begin{solve}
对所有实数 $x$，$x+\sqrt{1+x^2}>0$。令 $u=\ln(x+\sqrt{1+x^2})$，则 $du=dx/\sqrt{1+x^2}$；又
\[
 d\left(\frac{x}{\sqrt{1+x^2}}\right)=\frac{dx}{(1+x^2)^{3/2}}.
\]
分部积分给出
\[
 \int\frac{\ln(x+\sqrt{1+x^2})}{(1+x^2)^{3/2}}\,dx
 =\frac{x\ln(x+\sqrt{1+x^2})}{\sqrt{1+x^2}}
 -\frac12\ln(1+x^2)+C.
\]
复核时，第一项的导数产生原被积函数及 $x/(1+x^2)$，后者恰被第二项导数抵消。也可令 $x=\sinh u$，将积分化为 $\int u\operatorname{sech}^2u\,du$；两种方法本质上相同。

分部积分具体取 $dv=(1+x^2)^{-3/2}dx$、$v=x/\sqrt{1+x^2}$，所以剩余积分为 $\int v\,du=\int x/(1+x^2)dx$。令 $w=1+x^2$，$dw=2x\,dx$，得到该项为 $\tfrac12\ln(1+x^2)$。由于 $\sqrt{1+x^2}>|x|$，原对数参数对每个实数都严格正，所得原函数在整条实轴有定义，积分常数只需一个。

\end{solve}

\begin{example}[原卷二：路径无关与分支常数]\label{nm-2018-2}
设 $f\in C^1(\mathbb R\setminus\{0\})$，$f(1)=0$。若
\[
 \int_L y[2-f(x^2-y^2)]\,dx+x f(x^2-y^2)\,dy
\]
在不穿过直线 $y=x$、$y=-x$ 的曲线中只依赖端点，求 $f$。
\end{example}
\begin{solve}
令 $t=x^2-y^2$、$P=y[2-f(t)]$、$Q=xf(t)$。路径无关推出每个连通分支内 $P_y=Q_x$，于是
\[
 2-f(t)+2y^2f'(t)=f(t)+2x^2f'(t),\qquad tf'(t)+f(t)=1.
\]
等价地，$(tf(t))'=1$。由于定义域有两个连通分支，积分常数必须分别保留。因此
\[
 f(t)=\begin{cases}1-1/t,&t>0,\\1+C_-/t,&t<0,\end{cases}
 \qquad C_-\in\mathbb R.
\]
条件 $f(1)=0$ 只决定正半轴的常数，不能决定负半轴。还需检验所得函数确实足够：在 $t$ 符号固定的每个扇形区域内，设相应常数为 $C$，则原微分形式为
\[
 d(xy)+C\frac{-y\,dx+x\,dy}{x^2-y^2}
 =d\left[xy+\frac C2\ln\left|\frac{x+y}{x-y}\right|\right].
\]
故每个连通分支上均存在势函数，曲线积分确实只依赖端点。

局部势函数的偏导必须满足混合偏导相等，所以先求 $P_y,Q_x$ 是必要条件；但一般区域上这并非全局路径无关的充分条件，因此还要给出最后的势函数验证。其对数部分满足
\[
d\log\left|\frac{x+y}{x-y}\right|
=\frac{-2y\,dx+2x\,dy}{x^2-y^2},
\]
恰好恢复所需系数。两条被删去的直线将平面分为四个扇形，但未知函数 $f(t)$ 的自变量只有正、负两个分支，所以只允许每个符号一个常数，而非四个独立常数。

\end{solve}
\begin{remark}[对参考答案的条件补充]
原参考答案给出 $f(t)=1-1/t$。它是合要求的一种取法，但在题面给定的 $\mathbb R\setminus\{0\}$ 上并不是唯一解。若另加负半轴上的初值或联结两侧的适当条件，才可进一步确定 $C_-$。穿过被删除直线的路径不属于题设允许的路径；不能用它把两个积分常数强行联系起来。
\end{remark}

\begin{example}[原卷三：积分乘积的上下界]\label{nm-2018-3}
设 $f$ 在 $[0,1]$ 上连续，且 $1\leq f(x)\leq3$。证明
\[
 1\leq\left(\int_0^1f(x)\,dx\right)
          \left(\int_0^1\frac{dx}{f(x)}\right)\leq\frac43.
\]
\end{example}
\begin{solve}
记两积分为 $A,B$。对 $\sqrt f$ 与 $1/\sqrt f$ 使用积分型柯西不等式，得 $1\leq AB$。又 $(f-1)(f-3)\leq0$，除以正数 $f$ 可得 $f+3/f\leq4$，积分后为 $A+3B\leq4$。由 $12AB\leq(A+3B)^2$，得到 $AB\leq4/3$。

下界取等当且仅当 $f$ 为常函数。若上界取等，必须同时有 $f(x)\in\{1,3\}$ 对所有 $x$ 成立，以及 $A=3B$。连续函数在区间上的像是区间，故前一个条件迫使 $f$ 恒为1或恒为3；这两种情况均不满足后一个条件。因此在连续函数类中，上界实际严格，但 $4/3$ 仍是最佳常数：取函数在前半段接近1、后半段接近3，在中间长度趋于0的区间内线性连接，则 $A\to2$、$B\to2/3$。

最后一步的均值不等式可直接展开：$(A+3B)^2-12AB=(A-3B)^2\ge0$。因此上界需要两步同时取等，一是积分中的连续非负函数 $4-f-3/f$ 恒为零，二是 $A=3B$。前者迫使 $f$ 只取 $1,3$，连续性又不允许在两值间跳变，这说明上界虽最佳却不能由连续函数取得。下界的 Cauchy 等号要求 $\sqrt f$ 与 $1/\sqrt f$ 成比例，等价于 $f$ 为正常数。

\end{solve}
\begin{remark}
最佳上界未必可以取得。题目只要求非严格不等式，进一步讨论取等条件时必须保留函数的连续性；把两值阶梯函数直接说成连续类中的等号例子是不正确的。
\end{remark}

\begin{example}[原卷四：两个球之间的积分]\label{nm-2018-4}
设 $\Omega$ 满足
\[
 x^2+y^2+(z-2)^2\geq4,\quad
 x^2+y^2+(z-1)^2\leq9,\quad z\geq0.
\]
求 $\displaystyle\iiint_\Omega(x^2+y^2)\,dx\,dy\,dz$。
\end{example}
\begin{solve}
固定高度 $z$，区域是同心圆环。两球在顶部 $z=4$ 内切，允许高度为 $0\leq z\leq4$。外、内半径的平方分别为
\[
 R^2=9-(z-1)^2=8+2z-z^2,
 \qquad r^2=4-(z-2)^2=4z-z^2.
\]
两者在该区间非负，且 $R^2-r^2=8-2z\geq0$。柱坐标的体积元为 $\rho\,d\rho\,d\theta\,dz$，故
\begin{align*}
 I&=\int_0^4\int_0^{2\pi}\int_r^R\rho^3\,d\rho\,d\theta\,dz\\
 &=\frac\pi2\int_0^4(8-2z)(8+6z-2z^2)\,dz
 =\frac{256\pi}{3}.
\end{align*}
这里积分函数贡献 $\rho^2$，坐标变换再贡献 $\rho$，两者不可混淆。若先对 $z$ 积分，通常需要处理多个根式分支；水平截面把几何关系直接转化为多项式积分。

径向积分为 $\int_r^R\rho^3d\rho=(R^4-r^4)/4$，角积分再乘 $2\pi$，所以整体系数是 $\pi/2$。两个平方半径的和为 $8+6z-2z^2$、差为 $8-2z$，相乘后为 $64+32z-28z^2+4z^3$。因此括号内从 $0$ 到 $4$ 的积分为 $256+256-1792/3+256=512/3$，乘 $\pi/2$ 得 $256\pi/3$。

\end{solve}

\begin{example}[原卷五：梯度界与函数增量]\label{nm-2018-5}
设 $f$ 在区域 $D$ 内可微，且 $f_x^2+f_y^2\leq M^2$。点 $A=(a,b)$、$B=(c,d)$ 之间的整条线段包含在 $D$ 中。证明
\[
 |f(B)-f(A)|\leq M\sqrt{(c-a)^2+(d-b)^2}.
\]
\end{example}
\begin{solve}
若 $A=B$，结论直接成立。否则令 $g(t)=f(A+t(B-A))$，$0\leq t\leq1$。多元可微性保证复合函数满足链式法则，因此
\[
 g'(t)=f_x(A+t(B-A))(c-a)+f_y(A+t(B-A))(d-b).
\]
由有限维柯西不等式，$|g'(t)|\leq M\|B-A\|$。对 $g$ 使用一元中值定理即得结论。这个证明不要求偏导连续；可微性及题设的逐点界已经足够。若要对区域内任意两点统一得到该估计，可以要求 $D$ 为凸集。仅连通而不凸时，仍可沿允许路径估计，但右端一般应改为 $M$ 乘路径长度。

中值定理在此给出某个 $t_*\in(0,1)$ 使 $f(B)-f(A)=g'(t_*)$，因为参数区间长度为 $1$。然后用 $|\nabla f\cdot(B-A)|\le\|\nabla f\|\|B-A\|$，题设第一因子至多为 $M$，第二因子正是两点距离。要求整条线段位于 $D$，是为了让复合函数在全部参数区间都有定义并满足中值定理条件，不能只检查两个端点。

\end{solve}

\begin{example}[原卷六：对数积分不等式]\label{nm-2018-6}
设 $f$ 在 $[0,1]$ 上连续且严格为正。证明 $\displaystyle\int_0^1\ln f(x)\,dx\leq\ln\int_0^1f(x)\,dx$。
\end{example}
\begin{solve}
设 $A=\int_0^1f(x)\,dx>0$。对任意 $u>0$，函数 $u-1-\ln u$ 在 $u=1$ 取得最小值0，故 $\ln u\leq u-1$。取 $u=f(x)/A$ 并积分，得到
\[
 \int_0^1\ln f(x)\,dx-\ln A
 \leq\frac1A\int_0^1f(x)\,dx-1=0.
\]
取等当且仅当 $f(x)/A=1$ 对所有 $x$ 成立，即 $f$ 为常函数。严格正的连续函数在紧区间上有正的最小值，所以这里的对数连续、可积，不需要额外讨论端点奇性。这是凹函数 Jensen 不等式在对数情形的一个完整证明。

基本不等式的证明是令 $q(u)=u-1-\ln u$，则 $q'(u)=(u-1)/u$，在 $(0,1)$ 负、在 $(1,\infty)$ 正，因此唯一最小点为 $u=1$，最小值为零。归一化为 $f/A$ 后其积分恰为 $1$，所以线性上界 $u-1$ 的积分为零。若整体取等，连续非负函数 $q(f/A)$ 的积分为零，必须处处为零，于是 $f/A\equiv1$；这使等号论证也有完整依据。

\end{solve}

\begin{example}[原卷七：均值不等式与求和换序]\label{nm-2018-7}
设 $a_n,b_n>0$，存在 $\delta>0$ 使 $b_{n+1}-b_n\geq\delta$，且 $\sum_{n=1}^\infty a_n$ 收敛。证明
\[
 \sum_{n=1}^\infty
 \frac{n(a_1\cdots a_n b_1\cdots b_n)^{1/n}}{b_nb_{n+1}}
\]
收敛。
\end{example}
\begin{solve}
算术—几何均值不等式给出
\[
 n(a_1b_1\cdots a_nb_n)^{1/n}\leq\sum_{j=1}^n a_jb_j.
\]
又因 $b_{n+1}-b_n\geq\delta$，有
\[
 \frac1{b_nb_{n+1}}\leq\frac1\delta\left(\frac1{b_n}-\frac1{b_{n+1}}\right).
\]
先在有限和中换序，避免预先假定无限和可交换。对任意 $N$，原级数的前 $N$ 项和不超过
\begin{align*}
 \frac1\delta\sum_{n=1}^N\sum_{j=1}^na_jb_j
 \left(\frac1{b_n}-\frac1{b_{n+1}}\right)
 &=\frac1\delta\sum_{j=1}^Na_jb_j
 \left(\frac1{b_j}-\frac1{b_{N+1}}\right)\\
 &\leq\frac1\delta\sum_{j=1}^Na_j.
\end{align*}
部分和单调且有统一上界，因此收敛。证明还得到和不超过 $\delta^{-1}\sum a_j$。它把乘积的困难转化为线性和，再利用差分结构望远镜相消。

换序的下标区域是 $1\le j\le n\le N$；固定 $j$ 后，$n$ 从 $j$ 到 $N$，所以内和为 $\sum_{n=j}^N(1/b_n-1/b_{n+1})=1/b_j-1/b_{N+1}$。乘上 $a_jb_j$ 后，第一项恰好化为 $a_j$，剩余项为非正，可舍去得到统一上界。由于原级数各项为正，部分和单调；单调性与有界性共同给出收敛，不能单凭每项小于某个趋零量就作判断。

\end{solve}

\section{第十一届初赛：2019年非数学类}
本节依据\href{https://math.ecnu.edu.cn/llin/pdf_files/nonmathans2019.pdf}{第十一届初赛非数学类试卷及参考答案}，选入8道题，保留原卷编号。

\begin{example}[原卷一（1）：含立方根的复合极限]\label{nm-2019-1a}
求
\[
 \lim_{x\to0}\frac{\ln\bigl(e^{\sin x}+\sqrt[3]{1-\cos x}\bigr)-\sin x}
 {\arctan\bigl(4\sqrt[3]{1-\cos x}\bigr)}.
\]
\end{example}
\begin{solve}
令 $r=\sqrt[3]{1-\cos x}$。当 $x\to0$ 时 $r\to0$，且在充分小的去心邻域内 $r>0$。将分子准确改写为 $\ln(1+r e^{-\sin x})$，原比值成为
\[
 \frac{\ln(1+r e^{-\sin x})}{r e^{-\sin x}}
 \cdot\frac{e^{-\sin x}}4\cdot
 \frac{4r}{\arctan(4r)}\longrightarrow\frac14.
\]
这个改写先消去两个主项之间的相减，随后每个因子都可以直接使用基本极限。没有必要把 $r$ 展成分数次幂，也不需要对立方根反复求导。

分子的改写利用 $\ln(e^{\sin x}+r)=\ln[e^{\sin x}(1+re^{-\sin x})]=\sin x+\ln(1+re^{-\sin x})$，所以相减是完全精确的。随后令 $u=re^{-\sin x}\to0$，用 $\ln(1+u)/u\to1$；令 $v=4r\to0$，用 $v/\arctan v\to1$，剩余中间因子趋于 $1/4$。所有分母在足够小的去心邻域内均非零，因此分解合法。

\end{solve}

\begin{example}[原卷一（3）：识别乘积导数]\label{nm-2019-1c}
计算 $\displaystyle\int_0^{\pi/2}e^x\frac{1+\sin x}{1+\cos x}\,dx$。
\end{example}
\begin{solve}
半角公式给出 $\tan(x/2)=\sin x/(1+\cos x)$，以及 $\frac12\sec^2(x/2)=1/(1+\cos x)$。因此
\[
 \frac{d}{dx}\left(e^x\tan\frac x2\right)
 =e^x\left(\tan\frac x2+\frac12\sec^2\frac x2\right)
 =e^x\frac{1+\sin x}{1+\cos x}.
\]
定积分为 $[e^x\tan(x/2)]_0^{\pi/2}=e^{\pi/2}$。积分区间内分母严格为正，因此这是普通定积分；无需把端点误判为奇点。

寻找原函数时可先设它为 $e^xv(x)$，则需满足 $v+v'=(1+\sin x)/(1+\cos x)$。观察右端恰好拆成 $\sin x/(1+\cos x)+1/(1+\cos x)$，分别对应 $\tan(x/2)$ 及其导数，这解释了乘积的选择。上限处 $\tan(\pi/4)=1$，下限处 $\tan0=0$，所以代端点后只剩 $e^{\pi/2}$。

\end{solve}

\begin{example}[原卷一（4）：局部原函数]\label{nm-2019-1d}
已知微分关系 $\displaystyle du=\frac{y\,dx-x\,dy}{3x^2-2xy+3y^2}$，求 $u$。
\end{example}
\begin{solve}
在 $y\ne0$ 的区域内令 $t=x/y$，则 $dt=(y\,dx-x\,dy)/y^2$，于是
\[
 du=\frac{dt}{3t^2-2t+3},\qquad
 u=\frac1{2\sqrt2}\arctan\frac{3x-y}{2\sqrt2\,y}+C.
\]
这里利用 $3t^2-2t+3=((3t-1)^2+8)/3$。公式只直接描述不跨越 $y=0$ 的局部坐标片；跨片时需调整反正切分支和常数。

原微分形式在去掉原点的平面上光滑，却没有单值的全局势函数。沿单位圆 $x=\cos\theta,y=\sin\theta$ 逆时针积分，分子等于 $-d\theta$，分母为 $3-\sin2\theta>0$，所以一周积分严格为负。闭路积分不为零排除了全局路径无关性。求出一个局部反正切表达式，并不等于证明它在穿孔平面上是全局函数。

令 $w=(3t-1)/(2\sqrt2)$，则 $dt=(2\sqrt2/3)dw$，且分母为 $8(1+w^2)/3$，所以原函数为 $(2\sqrt2)^{-1}\arctan w$。在 $y\ne0$ 的坐标片代回 $t=x/y$，反正切的自变量就是 $(3x-y)/(2\sqrt2\,y)$。全局障碍也与局部公式一致：绕原点一圈时角变量需要连续改变分支，反正切不能始终保持一个主值；因此只能在选定局部区域内把它作为单值势函数。

\end{solve}

\begin{example}[原卷一（5）：椭球面相切]\label{nm-2019-1e}
设 $a,b,c,\mu>0$。曲面 $xyz=\mu$ 与椭球面 $x^2/a^2+y^2/b^2+z^2/c^2=1$ 相切时，求 $\mu$。
\end{example}
\begin{solve}
记 $X=x/a,Y=y/b,Z=z/c$，椭球约束变为 $X^2+Y^2+Z^2=1$。由三个非负数的均值不等式，
\[
 |XYZ|\leq\left(\frac{X^2+Y^2+Z^2}{3}\right)^{3/2}
 =\frac1{3\sqrt3}.
\]
取等条件为 $|X|=|Y|=|Z|=1/\sqrt3$。相切也可由法向量平行检验：$\nabla(xyz)=(yz,xz,xy)$ 与 $(2x/a^2,2y/b^2,2z/c^2)$ 平行，乘各分量以相应坐标，得到 $x^2/a^2=y^2/b^2=z^2/c^2$。故只能有
\[
 \mu=\frac{abc}{3\sqrt3}.
\]
乘积为正的四个符号组合给出相切点。$\mu$ 较小时两曲面可能相交，但法向量不平行；相交条件不能代替相切条件。

相切点处 $xyz=\mu>0$，所以三个坐标均非零，两曲面的梯度也均非零。设 $(yz,xz,xy)=\lambda(2x/a^2,2y/b^2,2z/c^2)$，分别乘以 $x,y,z$ 得到同一个左端 $\mu$，故三项 $x^2/a^2,y^2/b^2,z^2/c^2$ 相等，再由和为 $1$ 得各为 $1/3$。反过来，在这些坐标值及正乘积符号组合下梯度确实成比例，所以求出的参数既必要又充分。

\end{solve}

\begin{example}[原卷二：球坐标中的非球形区域]\label{nm-2019-2}
设 $\Omega$ 位于第一卦限，且满足 $(x^2+y^2+z^2)^2\leq2xy$。求
\[
 I=\iiint_\Omega\frac{xyz}{x^2+y^2}\,dx\,dy\,dz.
\]
\end{example}
\begin{solve}
采用 $x=\rho\sin\varphi\cos\theta,y=\rho\sin\varphi\sin\theta,z=\rho\cos\varphi$，其中 $0\leq\theta,\varphi\leq\pi/2$。对 $\rho>0$，区域不等式化为
\[
 0\leq\rho\leq\sin\varphi\sqrt{2\cos\theta\sin\theta}.
\]
被积函数与体积元相乘成为 $\rho^3\sin\varphi\cos\varphi\cos\theta\sin\theta$。因而
\begin{align*}
 I&=\int_0^{\pi/2}\int_0^{\pi/2}
 \sin^5\varphi\cos\varphi\cos^3\theta\sin^3\theta\,d\varphi\,d\theta\\
 &=\left(\int_0^{\pi/2}\sin^5\varphi\cos\varphi\,d\varphi\right)
 \left(\int_0^{\pi/2}\cos^3\theta\sin^3\theta\,d\theta\right)
 =\frac16\cdot\frac1{12}=\frac1{72}.
\end{align*}
在 $x=y=0$ 上原式未定义，但这条轴与区域的交只有原点；修改单点函数值不影响积分。也可先删去 $\rho<\varepsilon$ 的小球再令 $\varepsilon\downarrow0$，以上非负且可积的表达式保证极限存在。

原不等式是 $\rho^4\le2\rho^2\sin^2\varphi\cos\theta\sin\theta$，对 $\rho>0$ 约去 $\rho^2$ 得径向上界。径向积分 $\int_0^{\rho_{\max}}\rho^3d\rho=\rho_{\max}^4/4$，而 $\rho_{\max}^4=4\sin^4\varphi\cos^2\theta\sin^2\theta$，其中的 $4$ 正好抵消。第一个角积分令 $u=\sin\varphi$ 得 $1/6$；第二个令 $v=\sin^2\theta$ 得 $\tfrac12\int_0^1v(1-v)dv=1/12$。

\end{solve}

\begin{example}[原卷三：导数被函数本身控制]\label{nm-2019-3}
设 $f$ 在 $[0,+\infty)$ 上可导，$f(0)=0$，且 $|f'(x)|\leq A|f(x)|$，其中 $A>0$。证明 $f\equiv0$。
\end{example}
\begin{solve}
令 $h(x)=e^{-2Ax}f(x)^2$。直接求导，并使用 $f f'\leq |f||f'|\leq A f^2$，得到
\[
 h'(x)=2e^{-2Ax}[f(x)f'(x)-Af(x)^2]\leq0.
\]
所以 $h$ 单调不增。另一方面 $h\geq0$ 且 $h(0)=0$，故对任意 $x\geq0$，$0\leq h(x)\leq h(0)=0$。于是 $f(x)^2=0$。

证明不需要除以 $f$，也不需要假定 $f'$ 连续。若先写 $f'/f$，就会在最关键的零点附近失去合法性；使用平方和积分因子可以保留所有零点。这也是一阶微分方程零初值解唯一性的一个基本模型。

积分因子 $e^{-2Ax}$ 的指数由平方的导数确定：$(f^2)'=2ff'\le2Af^2$，所以减去 $2Af^2$ 后导数非正。对任意有限的 $x>0$，中值定理已足以推出 $h(x)\le h(0)$，不需要先证明 $f'$ 黎曼可积。这个方法使用的是全区间微分不等式加零初值；若初值不是零，只能得到 $|f(x)|\le e^{Ax}|f(0)|$ 这样的增长界，并不会推出函数恒为零。

\end{solve}

\begin{example}[原卷四：球面上的指数积分]\label{nm-2019-4}
计算
\[
 I=\int_0^{2\pi}\int_0^\pi
 e^{\sin\theta(\cos\varphi-\sin\varphi)}\sin\theta\,d\theta\,d\varphi.
\]
\end{example}
\begin{solve}
把它看作单位球面 $S^2$ 上的第一型曲面积分 $\int_{S^2}e^{x-y}\,dS$。正交坐标变换保持球面和面积元不变，令新的极轴指向向量 $(1,-1,0)$，则 $x-y=\sqrt2\cos\alpha$。于是
\[
 I=2\pi\int_0^\pi e^{\sqrt2\cos\alpha}\sin\alpha\,d\alpha
 =2\pi\int_{-1}^1e^{\sqrt2 u}\,du
 =\sqrt2\pi(e^{\sqrt2}-e^{-\sqrt2}).
\]
这个思路可以推广到任意固定向量 $v$：$\int_{S^2}e^{v\cdot x}\,dS=4\pi\sinh\|v\|/\|v\|$；$v=0$ 时按连续延拓取球面积 $4\pi$。识别旋转不变性往往比逐层展开指数级数更短，也避免不必要的级数换序。

旋转后的单位极轴是 $v=(1,-1,0)/\sqrt2$，任一球面点与它的夹角为 $\alpha$，所以 $v\cdot(x,y,z)=\cos\alpha$，也就是 $x-y=\sqrt2\cos\alpha$。绕新极轴的方位角不再出现在被积函数中，其积分贡献 $2\pi$。最后令 $u=\cos\alpha$，$du=-\sin\alpha\,d\alpha$，上、下限分别从 $1$ 到 $-1$，反向后得到 $\int_{-1}^1e^{\sqrt2u}du=(e^{\sqrt2}-e^{-\sqrt2})/\sqrt2$。

\end{solve}

\begin{example}[原卷五：由系数恢复最小正根]\label{nm-2019-5}
设多项式 $f$ 的全部根都是正实数。在原点附近展开
\[
 -\frac{f'(x)}{f(x)}=\sum_{n=0}^\infty c_nx^n.
\]
证明 $c_n>0$，并求 $\displaystyle\lim_{n\to\infty}c_n^{-1/n}$。
\end{example}
\begin{solve}
设互异根为 $0<r_1<\cdots<r_m$，相应重数为正整数 $k_j$。因 $f(0)\ne0$，在 $|x|<r_1$ 内可对数求导并展开有限个几何级数：
\[
 -\frac{f'(x)}{f(x)}=\sum_{j=1}^m\frac{k_j}{r_j-x},
 \qquad c_n=\sum_{j=1}^m\frac{k_j}{r_j^{n+1}}>0.
\]
记次数 $d=\sum k_j$，有
\[
 \frac{k_1}{r_1^{n+1}}\leq c_n\leq\frac d{r_1^{n+1}}.
\]
取 $n$ 次方根，两端均趋于 $1/r_1$，故所求极限为 $r_1$。重数只进入不会影响根极限的常数因子。这里的系数渐近由离展开中心最近的极点决定；正系数又使上下界完全避免了相消问题。

这里按题意考察正次数的多项式。因式分解 $f(x)=C\prod_j(x-r_j)^{k_j}$ 的对数导数为 $f'/f=\sum_jk_j/(x-r_j)$；负号使展开后的系数全部为正。在 $|x|<r_1$ 内，$1/(r_j-x)=\sum_{n\ge0}x^n/r_j^{n+1}$。对上、下界取根后，常数因子满足 $k_1^{1/n},d^{1/n}\to1$，而 $r_1^{-(n+1)/n}\to1/r_1$，最后取倒数才得到最小正根 $r_1$。

\end{solve}

\section{第十二届初赛：2020年非数学类}
本节依据\href{https://math.ecnu.edu.cn/llin/pdf_files/nonmathans2020.pdf}{第十二届初赛非数学类试卷及参考答案}，选入9道题。涉及条件收敛的积分和级数时，解析保留求和、积分及取极限的次序。

\begin{example}[原卷一（1）：三阶消去]\label{nm-2020-1a}
求 $\displaystyle\lim_{x\to0}\frac{(x-\sin x)e^{-x^2}}{\sqrt{1-x^3}-1}$。
\end{example}
\begin{solve}
先把分母有理化，原式等于
\[
 -\frac{x-\sin x}{x^3}\,e^{-x^2}\bigl(\sqrt{1-x^3}+1\bigr).
\]
Taylor 公式给出 $\sin x=x-x^3/6+o(x^3)$，故三个因子分别趋于 $-1/6,1,2$，答案为 $-1/3$。根式在原点邻域有定义，且对足够小的非零 $x$，分母非零。采用三阶展开是因为一阶项在 $x-\sin x$ 中抵消；只写 $\sin x\sim x$ 无法确定剩余项的阶数。

有理化来自 $(\sqrt{1-x^3}-1)(\sqrt{1-x^3}+1)=-x^3$，负号必须保留。随后 $(x-\sin x)/x^3\to1/6$，无论从左还是从右趋近都成立；$e^{-x^2}\to1$、根式加一趋于 $2$，三者相乘便是 $-1/3$。这里两个三阶零点相除，既要匹配阶数，也要匹配首个非零项系数。

\end{solve}

\begin{example}[原卷一（2）：高阶零点与导数]\label{nm-2020-1b}
设 $f(x)=(x+1)^ne^{-x^2}$，$n$ 为正整数，求 $f^{(n)}(-1)$。
\end{example}
\begin{solve}
Leibniz 公式中，若只对 $(x+1)^n$ 求了 $k<n$ 次导数，该因子在 $x=-1$ 仍含正次幂而为零。唯一保留下来的项是把 $n$ 次导数全部作用在多项式因子上的项，所以
\[
 f^{(n)}(-1)=n!e^{-1}.
\]
也可令 $h=x+1$，则 $e^{-x^2}=e^{-1}+O(h)$，从而 $f(-1+h)=e^{-1}h^n+O(h^{n+1})$。这表明 $n$ 阶导数抽取的是局部 Taylor 展开中 $h^n$ 的系数乘 $n!$，不必展开高斯函数的所有高阶导数。

Leibniz 公式完整写成 $f^{(n)}=\sum_{k=0}^n\binom nk[(x+1)^n]^{(k)}(e^{-x^2})^{(n-k)}$。当 $k<n$ 时，多项式导数为 $n!(x+1)^{n-k}/(n-k)!$，在 $-1$ 处为零；当 $k=n$ 时，组合系数是 $1$、多项式导数是 $n!$，指数函数无需再求导，取值为 $e^{-1}$。因此只有这一项留下，所有高阶指数导数都无需计算。

\end{solve}

\begin{example}[原卷一（3）：隐函数切线]\label{nm-2020-1c}
由
\[
 \arctan\frac xy=\frac12\ln(x^2+y^2)-\frac12\ln2+\frac\pi4
\]
在 $(1,1)$ 附近确定的曲线，求该点切线。
\end{example}
\begin{solve}
将等式左减右记作 $F$，在 $(1,1)$ 附近 $y>0$，各项均光滑。计算 $F_y(1,1)=-1\ne0$，故存在唯一的局部可微隐函数。沿曲线求导可得
\[
 \frac{y-xy'}{x^2+y^2}=\frac{x+yy'}{x^2+y^2},
 \qquad (x+y)y'=y-x.
\]
代入 $(1,1)$ 得 $y'=0$，故切线为 $y=1$。隐函数定理的检验保证所求斜率确实来自一条局部曲线，而非单纯的形式运算。

对于 $\arctan(x/y)$，链式法则给出沿曲线的导数为 $(y-xy')/(x^2+y^2)$；对于 $\tfrac12\ln(x^2+y^2)$，导数为 $(x+yy')/(x^2+y^2)$。等式两边乘以非零的 $x^2+y^2$ 后得到一次方程，移项才有 $(x+y)y'=y-x$。在 $(1,1)$ 分母 $x+y=2$，分子为零，故切线水平。对应的偏导是 $F_x(1,1)=0,F_y(1,1)=-1$，也与 $y'=-F_x/F_y$ 一致。

\end{solve}

\begin{example}[原卷一（4）：条件收敛的二次积分]\label{nm-2020-1d}
已知 $\int_0^\infty(\sin x)/x\,dx=\pi/2$，求按先 $y$ 后 $x$ 的次序定义的反常积分
\[
 I=\int_0^\infty\frac{\sin x}{x}
 \left(\int_0^\infty\frac{\sin(x+y)}{x+y}\,dy\right)dx.
\]
\end{example}
\begin{solve}
在原点以连续延拓定义 $\sin x/x$，令 $F(x)=\int_0^x(\sin t)/t\,dt$，$A=\pi/2$。对每个固定 $x\geq0$，内积分等于 $A-F(x)$。先截断外层上限为 $R$，由普通微积分基本定理得到
\[
 I_R=\int_0^R F'(x)[A-F(x)]\,dx
 =AF(R)-\frac12F(R)^2.
\]
由于 $F(R)\to A$，有 $I= A^2/2=\pi^2/8$。全过程没有交换两个无限区间上的积分次序，因此不需要假定被积函数绝对可积。

内积分令 $t=x+y$ 后下限变为 $x$、上限仍为无穷，所以 $\int_0^\infty\sin(x+y)/(x+y)dy=\int_x^\infty\sin t/t\,dt=A-F(x)$。在有限区间上，$F'(x)[A-F(x)]$ 正是 $AF(x)-F(x)^2/2$ 的导数，而 $F(0)=0$，故下端点无贡献。令 $R\to\infty$ 时只对这个已求出的端点表达式取极限，就同时证明外层反常积分存在并算出它的值。

\end{solve}
\begin{remark}
把积分区域画成对称三角形后直接用“对称性减半”，必须另外证明有关二重反常积分及换序存在。本题已给出内外次序，利用尾积分的导数即可在该次序下严格计算。
\end{remark}

\begin{example}[原卷一（5）：函数与指数同时变化]\label{nm-2020-1e}
设 $f(x),g(x)\to a>0$（$x\to0$），且去心邻域内 $f(x)\ne g(x)$。求
\[
 \lim_{x\to0}\frac{f(x)^{g(x)}-g(x)^{g(x)}}{f(x)-g(x)}.
\]
\end{example}
\begin{solve}
由保号性，在足够小的去心邻域两底数均为正。固定 $x$，把 $g(x)$ 当作常数，对变量 $t$ 的函数 $t^{g(x)}$ 使用中值定理，得到某个介于 $f(x)$ 与 $g(x)$ 之间的 $\xi_x>0$，使比值等于
\[
 g(x)\xi_x^{g(x)-1}.
\]
夹逼定理给出 $\xi_x\to a$。由 $\xi_x^{g(x)-1}=\exp((g(x)-1)\ln\xi_x)$ 的连续性，所求极限为 $a\cdot a^{a-1}=a^a$。不需要假定 $f,g$ 可导，因为中值定理作用于辅助变量 $t$，并非作用于 $x$。

对每个固定的 $x$，辅助函数的导数是 $g(x)t^{g(x)-1}$，其中 $g(x)$ 不随辅助变量 $t$ 改变。因此没有产生 $g'(x)$ 或任何对数求导项。即使 $f,g$ 在零点附近振荡或不可导，只要二者都趋于同一个正数，介于它们之间的 $\xi_x$ 就趋于该数。正的共同极限同时保证实幂定义与对数连续性，因而每一步的条件都由题设满足。

\end{solve}

\begin{example}[原卷二：非线性递推线性化]\label{nm-2020-2}
设 $a_1=1$，$\displaystyle a_{n+1}=\frac{a_n}{(n+1)(a_n+1)}$，求 $\lim_{n\to\infty}n!a_n$。
\end{example}
\begin{solve}
归纳可得 $a_n>0$。取倒数消去分式递推，再除以阶乘，令 $b_n=1/(n!a_n)$，则
\[
 b_{n+1}=b_n+\frac1{n!},\qquad b_1=1.
\]
因此 $b_n=\sum_{k=0}^{n-1}1/k!\to e$，所求极限为 $1/e$。取倒数之前验证非零性、除以阶乘后检查初值，是这类递推题最容易遗漏的两个步骤。递推中的 $n+1$ 提示阶乘归一化，而 $a_n/(1+a_n)$ 的分式形状提示倒数变换。

取倒数的完整递推为 $1/a_{n+1}=(n+1)(1/a_n+1)$，再除以 $(n+1)!$ 得
\[
\frac1{(n+1)!a_{n+1}}=\frac1{n!a_n}+\frac1{n!}.
\]
从 $1$ 累加到 $n-1$ 后，$b_n=b_1+\sum_{j=1}^{n-1}1/j!=\sum_{j=0}^{n-1}1/j!$。初始项 $b_1=1$ 对应指数级数的 $j=0$ 项，不能漏掉。因极限 $e>0$，再取倒数得到 $n!a_n=1/b_n\to1/e$。

\end{solve}

\begin{example}[原卷三：两次中值定理的配合]\label{nm-2020-3}
设 $f\in C[0,1]$，在 $(0,1)$ 可导，且 $f(0)=0,f(1)=1$。证明存在 $x_0\in(0,1)$ 使 $f(x_0)=2-3x_0$；再证明存在互异的 $\xi,\eta\in(0,1)$ 使
\[
 [1+f'(\xi)][1+f'(\eta)]=4.
\]
\end{example}
\begin{solve}
设 $H(x)=f(x)+3x-2$。由 $H(0)=-2,H(1)=2$，介值定理给出所需 $x_0$。分别在 $[0,x_0]$ 和 $[x_0,1]$ 使用中值定理，得到
\[
 f'(\xi)=\frac{2-3x_0}{x_0},\quad
 f'(\eta)=\frac{3x_0-1}{1-x_0},\quad
 0<\xi<x_0<\eta<1.
\]
于是 $1+f'(\xi)=2(1-x_0)/x_0$，$1+f'(\eta)=2x_0/(1-x_0)$，相乘等于4。两个区间不相交的内部位置同时保证了 $\xi\ne\eta$。如果只分别声称“存在某个中值点”，却不控制其所在区间，就无法保证互异性。

辅助函数中系数 $3$ 并非任意：交点条件给 $f(x_0)=2-3x_0$，于是左段割线斜率加 $1$ 后为 $[f(x_0)+x_0]/x_0=2(1-x_0)/x_0$；右段则为 $[1-f(x_0)+1-x_0]/(1-x_0)=2x_0/(1-x_0)$。两个比值正好互为倒数再各乘 $2$，乘积自然固定为 $4$。这解释了为什么先找该特定直线与图像的交点，再分别使用中值定理。

\end{solve}

\begin{example}[原卷六：阶梯积分与因子和]\label{nm-2020-6}
对正整数 $n$ 定义
\[
 F(n)=\sum_{m=1}^n\int_0^m\cos\frac{2\pi n\lfloor x+1\rfloor}{m}\,dx.
\]
证明它等于 $n$ 的所有正因子之和，并计算 $F(2021)$。
\end{example}
\begin{solve}
在每个 $[k-1,k)$ 内，$\lfloor x+1\rfloor=k$。有限个端点不改变积分，故第 $m$ 项为
\[
 \sum_{k=1}^m\cos\frac{2\pi nk}{m}
 =\operatorname{Re}\sum_{k=1}^m q^k,
 \qquad q=e^{2\pi i n/m}.
\]
若 $m\mid n$，$q=1$，和为 $m$；否则 $q\ne1$ 而 $q^m=1$，几何级数之和为 $q(1-q^m)/(1-q)=0$。因此
\[
 F(n)=\sum_{m\mid n}m.
\]
因为 $2021=43\cdot47$ 且两因子均为素数，正因子为 $1,43,47,2021$，所以 $F(2021)=2112$。复数这里只是有限三角和的简写；也可以取几何级数恒等式的实部写成纯三角证明。

每个单位小区间上被积函数为常数，积分等于该常数乘长度 $1$，因此并不是用阶梯函数近似积分，而是精确拆分。几何和中 $q^m=e^{2\pi in}=1$ 对每个整数 $n$ 都成立；只有 $q=1$ 即 $n/m$ 为整数时，分母 $1-q$ 为零而不能套用分式公式，这一例外恰好筛出全部正因子。$43,47$ 的平方根均小于 $7$，检查不被 $2,3,5$ 整除即可确认它们是素数，故列出的四个因子完整。

\end{solve}

\begin{example}[原卷七：参数积分与级数]\label{nm-2020-7}
令 $u_n=\int_0^1(1+t^4)^{-n}\,dt$，$n\geq1$。证明 $u_n\to0$，证明 $\sum_{n\geq1}(-1)^nu_n$ 条件收敛，并证明 $p\geq1$ 时 $\sum_{n\geq1}u_n/n^p$ 收敛；求 $\sum_{n\geq1}u_n/n$。
\end{example}
\begin{solve}
首先 $0<u_{n+1}\leq u_n$。给定 $0<\varepsilon<1$，把积分拆成 $[0,\varepsilon]$ 与 $[\varepsilon,1]$，得到
\[
 0<u_n\leq\varepsilon+(1+\varepsilon^4)^{-n}.
\]
先令 $n\to\infty$，再令 $\varepsilon\downarrow0$，即知 $u_n\to0$。Leibniz 判别法推出交错级数收敛。另一方面，对 $n\geq2$，利用 $t^4\leq t$，有
\[
 u_n\geq\int_0^1(1+t)^{-n}\,dt
 =\frac{1-2^{1-n}}{n-1}.
\]
右端与 $1/n$ 等价，故 $\sum u_n$ 发散，交错级数仅条件收敛。

求带权和时，对 $t(1+t^4)^{-n}$ 求导并积分，有
\[
 2^{-n}=u_n-4n\int_0^1\frac{t^4}{(1+t^4)^{n+1}}\,dt
 =u_n-4n(u_n-u_{n+1}).
\]
因此有限部分和满足
\[
 \sum_{n=1}^N\frac{u_n}{n}
 =\sum_{n=1}^N\frac1{n2^n}+4(u_1-u_{N+1}).
\]
令 $N\to\infty$ 得 $S=\ln2+4u_1$，同时证明了收敛性。对 $1+t^4=(t^2+\sqrt2t+1)(t^2-\sqrt2t+1)$ 作部分分式，或配方积分，得到
\[
 u_1=\frac\pi{4\sqrt2}+\frac1{2\sqrt2}\ln(1+\sqrt2),
 \quad S=\ln2+\frac\pi{\sqrt2}+\sqrt2\ln(1+\sqrt2).
\]
当 $p\geq1$ 时，$0\leq u_n/n^p\leq u_n/n$，比较判别法完成最后一步。

最后的有理积分可展开如下，令 $D_\pm=t^2\pm\sqrt2t+1$，则
\[
\frac1{1+t^4}=
\frac{t/(2\sqrt2)+1/2}{D_+}
+\frac{-t/(2\sqrt2)+1/2}{D_-}.
\]
各分子再分成分母导数的一部分与常数项，积分得对数贡献 $[\ln(D_+/D_-)]_0^1/(4\sqrt2)=\ln(1+\sqrt2)/(2\sqrt2)$；剩余贡献为 $\tfrac14\int_0^1(D_+^{-1}+D_-^{-1})dt=\pi/(4\sqrt2)$。后者由 $D_\pm=(t\pm1/\sqrt2)^2+1/2$ 完成平方求出。这样得到 $u_1$ 的具体值，再代入已由有限和证明的 $S=\ln2+4u_1$，无需额外交换无限级数与积分。

\end{solve}
\begin{remark}
单独的 $u_n\to0$ 只保证交错判别的一部分条件，并不能推出 $\sum u_n$ 收敛。求和恒等式应先写成有限形式，让望远镜项保留 $u_{N+1}$，然后再使用已证明的极限。
\end{remark}

\section{第九届决赛：2018年非数学类}
本节来源为\href{https://math.ecnu.edu.cn/llin/pdf_files/NonMathFinalAns2018.pdf}{第九届决赛非数学类试卷及参考答案}，选入7道分析与微积分题。原卷中的线性代数题不在本节收录范围。

\begin{example}[原卷一（1）：三角函数之差]\label{nm-final9-1a}
求 $\displaystyle\lim_{x\to0}\frac{\tan x-\sin x}{x\ln(1+\sin^2x)}$。
\end{example}
\begin{solve}
利用准确恒等式 $\tan x-\sin x=\sin x(1-\cos x)/\cos x$，以及基本等价关系，得
\[
 \frac{\tan x-\sin x}{x\ln(1+\sin^2x)}
 =\frac{\sin x}{x}\cdot\frac{1-\cos x}{x^2}
 \cdot\frac{x^2}{\sin^2x}\cdot
 \frac{\sin^2x}{\ln(1+\sin^2x)}\cdot\frac1{\cos x}
 \longrightarrow\frac12.
\]
分子、分母均为三阶无穷小。此处先因式分解，再分别替换乘积中的因子，可以避免对 $\tan x$ 与 $\sin x$ 各自作一阶替换后错误地得到零。

五个因子的极限依次是 $1,1/2,1,1,1$。其中第四个使用 $u/\ln(1+u)\to1$，这里 $u=\sin^2x\to0$；第二个可由 $1-\cos x=2\sin^2(x/2)$ 得到极限 $1/2$。在足够小的去心邻域，$\sin x\ne0$、$\cos x\ne0$、$\ln(1+\sin^2x)>0$，所以逐因子分解没有引入零分母，左右极限一致。

\end{solve}

\begin{example}[原卷一（3）：全微分恢复函数]\label{nm-final9-1c}
已知 $df=ye^y\,dx+(1+y)xe^y\,dy$，且 $f(0,0)=0$，求 $f(x,y)$。
\end{example}
\begin{solve}
由 $f_x=ye^y$，对 $x$ 积分，得 $f(x,y)=xye^y+C(y)$。再对 $y$ 求导并与题设比较，得到 $C'(y)=0$，故 $C$ 为常数。初值决定该常数为0，所以 $f=xye^y$。直接求全微分可复核两个系数均正确。偏积分产生的是关于其余变量的函数，而不是在第一步就确定的常数。

将第一步结果对 $y$ 求导，$\partial_y(xye^y)=xe^y+xye^y=(1+y)xe^y$，故与题设相比只多出 $C'(y)$，必须 $C'(y)=0$。常数函数由 $f(0,0)=C=0$ 确定。在整个平面上 $xye^y$ 光滑，求其全微分恢复两项即可证明充分性；任意两个候选函数的两个偏导都相同，其差在连通平面上为常数，再由初值相同排除常数差，所以解唯一。

\end{solve}

\begin{example}[原卷一（4）：含整体积分的微分方程]\label{nm-final9-1d}
求可微函数 $u$，满足 $\displaystyle u'(t)=u(t)+\int_0^1u(s)\,ds$，以及 $u(0)=1$。
\end{example}
\begin{solve}
记 $C=\int_0^1u(s)\,ds$，它相对于变量 $t$ 是常数。解 $u'-u=C$ 并使用初值，得到 $u(t)=(1+C)e^t-C$。回代 $C$ 的定义，有
\[
 C=(1+C)(e-1)-C,
 \qquad C=\frac{e-1}{3-e}.
\]
由于 $e<3$，分母非零，故
\[
 u(t)=\frac{2e^t+1-e}{3-e}.
\]
最后代回原积分与导数等式即可验证。解普通微分方程只给出候选函数，整体积分约束才决定额外参数；漏掉回代会留下并不满足原题的一族函数。

线性方程乘以 $e^{-t}$ 后变为 $(e^{-t}u)'=Ce^{-t}$，利用初值求得 $u=1\cdot e^t+C(e^t-1)$。整体约束化为 $2C=(1+C)(e-1)$，移项得到 $(3-e)C=e-1$。代入最终函数后，$\int_0^1u=(e-1)/(3-e)$，而 $u'-u$ 也等于同一常数，且 $u(0)=1$，所以两个条件同时满足。

\end{solve}

\begin{example}[原卷二：割线斜率的介值性]\label{nm-final9-2}
设 $f$ 在 $(0,1)$ 连续，已知两条割线的斜率分别为 $\alpha<\beta$。证明对每个 $\lambda\in(\alpha,\beta)$，存在不同的 $s,t\in(0,1)$ 使 $(f(t)-f(s))/(t-s)=\lambda$。
\end{example}
\begin{solve}
把第一条割线端点按大小写成 $a_1<b_1$，第二条写成 $a_2<b_2$。对 $0\leq\tau\leq1$，令
\[
 a(\tau)=(1-\tau)a_1+\tau a_2,
 \quad b(\tau)=(1-\tau)b_1+\tau b_2.
\]
两端点始终在 $(0,1)$，且 $b(\tau)-a(\tau)=(1-\tau)(b_1-a_1)+\tau(b_2-a_2)>0$。因此
\[
 H(\tau)=\frac{f(b(\tau))-f(a(\tau))}{b(\tau)-a(\tau)}
\]
在 $[0,1]$ 连续，端值为 $\alpha,\beta$。介值定理给出 $\tau_0$，取 $s=a(\tau_0),t=b(\tau_0)$ 即可。这里不要求 $f$ 可导，割线斜率的连续变化来自端点空间中始终不碰撞的连线。

任意两条割线都可交换端点而不改变斜率，故先统一为左端点小于右端点不会损失一般性。插值过程中不仅分母为正，还满足 $b(\tau)-a(\tau)\ge\min\{b_1-a_1,b_2-a_2\}>0$，因此不会在极限过程中逼近端点碰撞。函数 $H$ 是连续函数之差除以连续且不为零的函数，故可直接应用普通介值定理；这一证明没有调用导数的 Darboux 性质。

\end{solve}

\begin{example}[原卷四：阶乘几何平均的增量]\label{nm-final9-4}
求 $\displaystyle\lim_{n\to\infty}\bigl(\sqrt[n+1]{(n+1)!}-\sqrt[n]{n!}\bigr)$。
\end{example}
\begin{solve}
记 $a_n=(n!)^{1/n}$、$L_n=\ln a_n$。先证明 $L_n=\ln n-1+o(1)$。由于 $\ln t$ 单调，可用积分与和的比较得到
\[
 \int_1^n\ln t\,dt\leq\ln(n!)\leq\int_1^n\ln t\,dt+\ln n.
\]
除以 $n$ 即得结论，因此 $a_n/n\to e^{-1}$。另一方面由阶乘递推准确得到
\[
 d_n=L_{n+1}-L_n=\frac{\ln(n+1)-L_n}{n+1},
 \qquad nd_n\to1.
\]
于是
\[
 a_{n+1}-a_n=\frac{a_n}{n}\cdot n d_n\cdot\frac{e^{d_n}-1}{d_n}
 \longrightarrow\frac1e.
\]
仅知道 $a_n\sim n/e$ 本身不足以推出相邻差趋于 $1/e$；例如在相似的主项上加入较小但快速振荡的误差，相邻差仍可能振荡。因此上面的对数递推是必要的额外控制。

积分上下界给出更明确的 $L_n=\ln n-1+O((\ln n)/n)$，因此 $\ln(n+1)-L_n=1+o(1)$，进而 $nd_n=\frac n{n+1}(1+o(1))\to1$，同时 $d_n\to0$。再由 $a_{n+1}=a_ne^{d_n}$ 得 $a_{n+1}-a_n=a_n(e^{d_n}-1)$，拆成正文三个因子后，每个都有独立确定的极限。这样相邻差中的误差由精确递推控制，不依赖对两个近似主项作相减。

\end{solve}

\begin{example}[原卷六：边值与梯度控制面积积分]\label{nm-final9-6}
设 $a>0$，$D=\{x^2+y^2\leq a^2\}$，$f\in C^1(D)$（导数连续延至边界），边界上 $f=a^2$，且 $\max_D(f_x^2+f_y^2)=a^2$。证明 $\displaystyle\iint_D f\,dx\,dy\leq4\pi a^4/3$。
\end{example}
\begin{solve}
由梯度界有 $\|\nabla f\|\leq a$。对半径为 $r>0$ 的点 $X$，沿同一射线取边界点 $Y=aX/r$。线段包含在圆盘内，故由一元中值定理和链式法则，
\[
 f(X)\leq f(Y)+a\|Y-X\|=a^2+a(a-r)=2a^2-ar.
\]
原点处可由连续性得到同样的上界。因此极坐标积分给出
\[
 \iint_D f\,dx\,dy\leq2\pi\int_0^a(2a^2-ar)r\,dr
 =\frac{4\pi a^4}{3}.
\]
也可对向量场 $(xf,yf)$ 使用平面散度公式，从边界积分中分离出 $\nabla f\cdot(x,y)$ 再估计。径向线段法更直接地揭示了边界值怎样控制内部函数值。

半径为 $r$ 的点到同射线边界点的距离为 $a-r$，所以梯度界把从边界向内的最大增长限制为 $a(a-r)$。极坐标面积元为 $r\,dr\,d\theta$，具体积分为
\[
2\pi\left[2a^2\frac{a^2}{2}-a\frac{a^3}{3}\right]
=2\pi\left(a^4-\frac{a^4}{3}\right)=\frac{4\pi a^4}{3}.
\]
原点方向不唯一，但任取径向方向的估计都有相同右端，或直接用连续性令 $r\downarrow0$，所以单点也不存在遗漏。

\end{solve}

\begin{example}[原卷七：对数阶判别]\label{nm-final9-7}
设 $0<a_n<1$，$\displaystyle q=\lim_{n\to\infty}\frac{\ln(1/a_n)}{\ln n}$ 存在，可以为 $+\infty$。证明 $q>1$ 时 $\sum a_n$ 收敛，$q<1$ 时发散；讨论 $q=1$。
\end{example}
\begin{solve}
若 $q>1$，取 $1<p<q$（$q=+\infty$ 时任选有限 $p>1$）。当 $n$ 足够大时，$\ln(1/a_n)>p\ln n$，所以 $a_n<n^{-p}$，级数收敛。若 $q<1$，取 $q<p<1$，则最终有 $a_n>n^{-p}$，级数发散。

当 $q=1$ 时不能仅凭该极限判断。对 $n\geq2$，取 $a_n=1/n$ 得发散例；对 $n\geq3$，取 $a_n=1/[n(\ln n)^2]$，则
\[
 \frac{\ln(1/a_n)}{\ln n}=1+2\frac{\ln\ln n}{\ln n}\to1,
\]
而积分判别法给出收敛。有限个初始项可任取 $(0,1)$ 中的数，不影响结论。临界点处，主导幂次相同，必须继续比较对数等更细的尺度。

对数比较时必须留意取负号后方向：$\ln(1/a_n)>p\ln n$ 等价于 $1/a_n>n^p$，即 $a_n<n^{-p}$；小于号则完全反向。由 $0<a_n<1$，对数商非负，所以有限的 $q$ 也非负。临界收敛例的积分为 $\int_3^Rdx/[x(\ln x)^2]=1/\ln3-1/\ln R$，有有限极限；而调和级数发散，两者对数商却都趋于 $1$，确实说明临界信息不足。

\end{solve}

\section{第八届决赛：2017年非数学类}
本节依据\href{https://math.ecnu.edu.cn/llin/pdf_files/nonmathfinalans2017.pdf}{第八届决赛非数学类试卷及参考答案}，选入4道分析题。

\begin{example}[原卷一（4）：有限和的整数部分]\label{nm-final8-1d}
求 $\displaystyle\sum_{n=1}^{100}\frac1{\sqrt n}$ 的整数部分。
\end{example}
\begin{solve}
函数 $x^{-1/2}$ 在正半轴严格递减。分别用左右端点矩形估计，得到
\[
 \int_1^{101}\frac{dx}{\sqrt x}
 <\sum_{n=1}^{100}\frac1{\sqrt n}
 <1+\int_1^{100}\frac{dx}{\sqrt x}=19.
\]
左端等于 $2(\sqrt{101}-1)>18$，因此该和位于 $(18,19)$，整数部分为18。严格不等号来自每个长度为1的区间上函数都非恒定。无需高精度数值求和，构造两个相邻整数之间的可靠界即可完成问题。

下界逐段使用 $\int_n^{n+1}x^{-1/2}dx<1/\sqrt n$，从 $n=1$ 加到 $100$ 得到 $\int_1^{101}$。上界则单独保留第一项 $1$，对 $n=2,\ldots,100$ 使用 $1/\sqrt n<\int_{n-1}^nx^{-1/2}dx$，得到 $1+\int_1^{100}$。两个方向的积分区间不同，正好把总和夹进长度为 $1$ 的区间；若两边都用同一组矩形，往往只能得到不足以确定整数部分的粗界。

\end{solve}

\begin{example}[原卷二：一个三角不等式]\label{nm-final8-2}
证明当 $0<x<\pi/2$ 时，$\displaystyle\frac4{\pi^2}<\frac1{x^2}-\frac1{\tan^2x}<\frac23$。
\end{example}
\begin{solve}
设 $F(x)=x^{-2}-\cot^2x$。其导数为
\[
 F'(x)=\frac{2(x^3\cos x-\sin^3x)}{x^3\sin^3x}.
\]
为证明分子为负，令 $H(x)=\sin x/(\cos x)^{1/3}-x$。直接求导后，写 $u=(\cos x)^{2/3}\in(0,1)$，有
\[
 H'(x)=\frac23u+\frac13u^{-2}-1>0.
\]
最后一个严格不等式由三个正数 $u,u,u^{-2}$ 的算术—几何均值不等式得到，等号只可能在 $u=1$。又 $H(0)=0$，因此 $\sin^3x>x^3\cos x$，从而 $F'<0$。

右端点极限为 $F(\pi/2-)=4/\pi^2$。左端点用 $\tan x=x+x^3/3+o(x^3)$ 可得
\[
 \frac1{\tan^2x}=\frac1{x^2}-\frac23+o(1),
 \qquad F(0+)=\frac23.
\]
严格单调性与两个端点极限即给出题设双边界。

左端极限中的倒数展开也应保留二阶修正：$\tan x=x(1+x^2/3+o(x^2))$，故
\[
\frac1{\tan^2x}=x^{-2}(1+x^2/3+o(x^2))^{-2}
=x^{-2}(1-2x^2/3+o(x^2)).
\]
这里使用 $(1+u)^{-2}=1-2u+o(u)$，乘回 $x^{-2}$ 后才留下常数 $-2/3$。严格递减性意味着每个内部点的值严格小于左端极限且严格大于右端极限，所以题目中的两侧都为严格不等式。

\end{solve}

\begin{example}[原卷四：由 Laplace 方程确定积分]\label{nm-final8-4}
设 $\Omega=\{x^2+y^2+z^2\leq1\}$，$f$ 具有连续二阶偏导数并连续延至边界，满足 $\Delta f=x^2+y^2+z^2$。求
\[
 I=\iiint_\Omega(xf_x+yf_y+zf_z)\,dV.
\]
\end{example}
\begin{solve}
记 $r^2=x^2+y^2+z^2$。在单位球面上 $r^2=1$，故 $\nabla f$ 与 $r^2\nabla f$ 的通量相等。对两者分别使用散度定理，得到
\[
 \iiint_\Omega\Delta f\,dV
 =\iiint_\Omega\operatorname{div}(r^2\nabla f)\,dV
 =2I+\iiint_\Omega r^2\Delta f\,dV.
\]
代入 $\Delta f=r^2$ 并用球坐标，
\[
 I=\frac12\iiint_\Omega(r^2-r^4)\,dV
 =2\pi\int_0^1(r^4-r^6)\,dr
 =\frac{4\pi}{35}.
\]
无须求出 $f$ 本身。边界上的等式把未知函数的两种通量联系起来，使题目所求的一阶导数组合出现；这是一种选择乘子并积分分部的思路。若试图直接解 Poisson 方程，反而会引入不必要的调和函数自由度。

乘积散度公式是 $\operatorname{div}(r^2\nabla f)=\nabla(r^2)\cdot\nabla f+r^2\Delta f$，而 $\nabla(r^2)=2(x,y,z)$，所以第一项恰好为所求被积函数的两倍。边界上两个向量场相同，因此各自的外通量相同，使用散度定理后未知边界量自动消去。球坐标中径向函数的体积分为 $4\pi\int_0^1(\cdot)r^2dr$，加上前面的 $1/2$，得到 $2\pi(1/5-1/7)=4\pi/35$。

\end{solve}

\begin{example}[原卷六：调和数与对数之差]\label{nm-final8-6}
设 $a_n=\sum_{k=1}^n1/k-\ln n$。证明 $a_n$ 有极限 $C$，并判断 $\sum_{n\geq1}(a_n-C)$ 的敛散性。
\end{example}
\begin{solve}
对 $t>0$，$t/(1+t)<\ln(1+t)<t$。于是
\[
 a_{n+1}-a_n=\frac1{n+1}-\ln\left(1+\frac1n\right)<0.
\]
积分比较又给出 $\sum_{k=1}^n1/k\geq\int_1^{n+1}dx/x$，从而 $a_n>0$。单调有界性保证极限存在。

设 $d_k=a_k-a_{k+1}=\ln(1+1/k)-1/(k+1)>0$。Taylor 展开或两次基本极限可得 $k^2d_k\to1/2$。由望远镜和，$a_n-C=\sum_{k=n}^\infty d_k$。给定 $0<\varepsilon<1/2$，足够大的 $n$ 满足
\[
 \left(\frac12-\varepsilon\right)\sum_{k=n}^\infty\frac1{k^2}
 \leq a_n-C\leq
 \left(\frac12+\varepsilon\right)\sum_{k=n}^\infty\frac1{k^2}.
\]
由积分比较，$n\sum_{k=n}^\infty k^{-2}\to1$，故 $n(a_n-C)\to1/2$。所以 $a_n-C\sim1/(2n)$，待判级数为正项发散级数。这说明趋向一个常数的误差，仍可能积累成发散的总量。

令 $u=1/k$，展开 $\ln(1+u)=u-u^2/2+O(u^3)$、$u/(1+u)=u-u^2+O(u^3)$，两式相减给出 $d_k=1/(2k^2)+O(k^{-3})$。平方倒数尾和的归一化可由
\[
\frac1n=\int_n^\infty x^{-2}dx\le\sum_{k=n}^\infty k^{-2}
\le\frac1{n^2}+\int_n^\infty x^{-2}dx
\]
夹逼得出。因而对 $d_k$ 取尾和后，误差的主要幂次由 $k^{-2}$ 变成 $n^{-1}$，这正是最终误差级数仍像调和级数一样发散的原因。

\end{solve}


\part{专题方法与证明工具}

本编各章主要改编自 Jacobi Lie《数学竞赛的方法与结构》的相应章节，原作者与 CC BY-SA 4.0 许可见编写说明。它们是解题方法说明，不计为历年真题。



\section{极限、上下极限与渐近记号}

求极限之前，先问对象是否真的收敛；求等价式之前，先问等价式将被怎样使用。上下极限描述尚未证明收敛的序列，渐近记号描述误差相对于指定尺度的大小。二者给后续的比较、展开和求和提供共同语言。

\subsection{上下极限与子列}

\label{app:limit-limsup}

对实数列 $(a_n)$，允许取扩展实数值，定义
\[
\limsup_na_n=\lim_{N\to\infty}\sup_{n\geq N}a_n,
\qquad
\liminf_na_n=\lim_{N\to\infty}\inf_{n\geq N}a_n.
\]
前一个尾部上界单调下降，后一个尾部下界单调上升，所以两极限总存在于扩展实数系中。它们相等且有限，当且仅当 $a_n$ 收敛到这个公共值。对有界序列，上极限是子列极限的最大值，下极限是最小值；紧性保证极端聚点能够被取到。

若 $L=\limsup a_n\in\mathbb R$，其准确含义有两部分：任意 $\varepsilon>0$ 下，最终 $a_n<L+\varepsilon$；同时有无穷多个 $n$ 使 $a_n>L-\varepsilon$。前者给最终统一上界，后者防止上界被进一步降低。证明极限时常先设上下极限为 $l,L$，把递推或不等式传到它们，再证明 $L\leq l$。

有界实数列满足
\[
\liminf a_n+\liminf b_n\leq\liminf(a_n+b_n)
\leq\limsup(a_n+b_n)\leq\limsup a_n+\limsup b_n.
\]
一般不能把最右不等号改成等号：$a_n=(-1)^n$、$b_n=-a_n$ 的两个最大聚值不会同时出现。若一个数列收敛，则可把它的极限移出另一个数列的上下极限。连续单调函数也可相应移出；递减函数则交换上下极限。

\textbf{应用：渐近压缩。} 若 $a_n\geq0$ 有界，$a_{n+1}\leq q a_n+\varepsilon_n$，其中 $0\leq q<1$、$\varepsilon_n\to0$，则对上极限取界得到 $L\leq qL$，从而 $a_n\to0$。若未预先知道有界，可先用 $\varepsilon_n$ 有界及几何求和补上这一步。这种论证不要求 $a_n$ 单调。

\subsection{大 O、小 o 与一致性}

\label{app:limit-notation}

在指定极限过程中，$f=O(g)$ 表示最终有 $|f|\leq C|g|$，常数 $C$ 不随趋近变量变化；$f=o(g)$ 表示 $f/g\to0$（分母须在所考虑区域内非零）；$f\sim g$ 表示 $f/g\to1$。$f\asymp g$ 通常表示二者非负且存在固定正数 $c,C$ 使 $cg\leq f\leq Cg$。每次都应明确是 $x\to0$、$n\to\infty$，还是另一个极限。

参数一致的 $r_\varepsilon(x)=O(\varepsilon^2)$（$x\in K$）表示同一个 $C$ 对所有 $x\in K$ 有效；一致小 o 则表示 $\sup_{x\in K}|r_\varepsilon(x)|/\varepsilon\to0$。若常数允许依赖参数，应写清其范围或记成 $O_a(\cdot)$。把固定参数结论用于随 $n$ 变化的参数，是误用渐近式的常见来源。

等价式可在乘除中替换，却不能任意穿过相减：$n+1\sim n$，但两者之差是 $1$，不是 $0$。存在抵消时，应改写成带余项的加法形式，例如
\[
\sqrt{1+x}=1+\frac x2-\frac{x^2}{8}+O(x^3).
\]
若要把差除以 $x^2$，只保留 $1+x/2+o(x)$ 就不够。展开到哪一阶，应由后续除法或累积的尺度反推。

逐项 $r_{n,k}=o(1)$ 不保证 $\sum_{k=1}^nr_{n,k}=o(1)$；即使所有项都等于 $1/n$，总和仍为 $1$。若有一致界 $|r_{n,k}|\leq Cn^{-2}$，才可立即得总和 $O(n^{-1})$。积分、乘积和行列式中的误差同样需要在操作层面控制。

\subsection{Cesàro 平均与 Stolz 的上下极限形式}

\label{app:limit-stolz}

若 $a_n\to L$，则 $n^{-1}\sum_{k=1}^na_k\to L$。证明把有限头部与最终落在 $L\pm\varepsilon$ 内的尾部分开，有限头部除以 $n$ 后消失。更一般地，正权 $w_k$ 满足 $W_n=\sum_{k\leq n}w_k\to\infty$ 时，$W_n^{-1}\sum_{k\leq n}w_ka_k\to L$。

\textbf{Stolz 的上下极限链。} 记
\[
r_n=\frac{a_n}{b_n},\qquad
u_n=\frac{a_{n+1}-a_n}{b_{n+1}-b_n}.
\]
在下列任一组条件下，都有
\[
\liminf u_n\leq\liminf r_n
\leq\limsup r_n\leq\limsup u_n,
\]
其中上下极限允许取 $\pm\infty$：

（1）无穷型：$b_n$ 严格递增到 $+\infty$，对实数列 $(a_n)$ 无其他要求；只需从 $b_n>0$ 的尾部起考虑商。

（2）零比零型：$b_n>0$ 严格递减到零，且 $a_n\to0$。

核心是正权平均。无穷型中，固定充分大的 $N$，则
\[
r_n=\frac{a_N}{b_n}
+\sum_{k=N}^{n-1}\frac{b_{k+1}-b_k}{b_n}u_k.
\]
有限头部 $a_N/b_n$ 趋于零，尾部正权之和趋于 $1$。因此 $u_k$ 的任意有限最终下界或上界，也是 $r_n$ 的下极限之下界或上极限之上界。让这些界逼近相应的上下极限，即得不等式链；用任意实数作界，也涵盖了无穷值。

零比零型则向尾端求和：由 $a_n,b_n\to0$，
\[
r_n=\sum_{k=n}^{\infty}\frac{b_k-b_{k+1}}{b_n}u_k.
\]
此式由有限望远镜求和后取极限得到，正权之和恰为 $1$，同样传递尾部的上下界。这里 $a_n\to0$ 不能删除，例如 $a_n=1+b_n$ 的差商恒为 $1$，而 $a_n/b_n\to+\infty$。

特别地，若 $u_n\to L$，不等式链便给出 $r_n\to L$，其中 $L$ 可以是扩展实数。这是通常的 Stolz 定理；无穷型并不要求 $a_n\to\infty$。

\textbf{何时可以反向。} 一般不能从商的收敛推出差商的收敛。例如 $a_n=nL+(-1)^n$、$b_n=n$ 时，$r_n\to L$，而 $u_n=L-2(-1)^n$ 振荡。以下两种附加条件分别控制尺度与振荡。

第一种是相邻尺度有固定分离：在上述任一种情形下，若最终有
\[
\frac{\max\{b_n,b_{n+1}\}}{|b_{n+1}-b_n|}\leq C,
\]
则 $r_n\to L\in\mathbb R$ 可以反推 $u_n\to L$。事实上，
\[
u_n-L=
\frac{b_{n+1}(r_{n+1}-L)-b_n(r_n-L)}{b_{n+1}-b_n},
\]
右侧两个系数都有界。无穷型的条件等价于增量 $(b_{n+1}-b_n)/b_n$ 最终有正下界；零比零型则是 $(b_n-b_{n+1})/b_n$ 最终有正下界。

若进一步有 $b_n/b_{n+1}\to0$（无穷型），或 $b_{n+1}/b_n\to0$（零比零型），且 $(r_n)$ 有界，则单个增量已支配相应尺度。由同一恒等式，分别有 $u_n-r_{n+1}\to0$ 或 $u_n-r_n\to0$，因而即使商不收敛，也有
\[
\liminf u_n=\liminf r_n,\qquad
\limsup u_n=\limsup r_n.
\]

第二种是单调性。若 $(v_n)$ 最终单调，且 $S_n=\sum_{k=1}^nv_k$ 满足 $S_n/n\to L\in\mathbb R$，则 $v_n\to L$。例如最终非减时，对充分大的 $n$，令 $m=\lfloor n/2\rfloor$，有
\[
\frac{S_n-S_m}{n-m}\leq v_n
\leq\frac{S_{2n}-S_n}{n}.
\]
两端均趋于 $L$；非增时不等号反向。这给出 Cesàro 平均的单调逆定理，也就是 $b_n=n$、差商最终单调时的 Stolz 逆向结论。平均会掩盖快速振荡，单调性正好排除了这一障碍。

\textbf{应用。} 对 $p>-1$，取 $a_n=\sum_{k=1}^nk^p$、$b_n=n^{p+1}$，差商趋于 $1/(p+1)$，所以 $\sum_{k\leq n}k^p\sim n^{p+1}/(p+1)$。若要求 $n^p/2$ 等修正项，则需\hyperref[app:taylor-em]{求和的 Euler--Maclaurin 公式}。Stolz 负责首阶比例，带余项的求和公式负责更精细的信息。

\subsection{L’Hospital 法则与连续平均}

\label{app:limit-lhospital}

Stolz 的连续对应同样是不等式链。统一考虑 $x\uparrow\omega$，其中 $\omega$ 可以有限，也可以是 $+\infty$。设 $f,g$ 在该端点的左邻域内可微，且最终 $g>0$。在下列任一组条件下，
\[
\liminf_{x\uparrow\omega}\frac{f'(x)}{g'(x)}
\leq\liminf_{x\uparrow\omega}\frac{f(x)}{g(x)}
\leq\limsup_{x\uparrow\omega}\frac{f(x)}{g(x)}
\leq\limsup_{x\uparrow\omega}\frac{f'(x)}{g'(x)}:
\]

（1）无穷型：$g'>0$，且 $g(x)\to+\infty$；不要求 $f$ 也趋于无穷。

（2）零比零型：$g'<0$，且 $f(x),g(x)\to0$。

上下极限允许取扩展实数值，导数比收敛时便得到通常的 L’Hospital 法则。其他单侧趋近方向或分母符号，可通过反向参数化或同时将 $f,g$ 变号处理。

证明只需 Cauchy 中值定理：对端点附近的 $x<y$，某个 $\xi\in(x,y)$ 满足
\[
\frac{f(y)-f(x)}{g(y)-g(x)}=\frac{f'(\xi)}{g'(\xi)}.
\]
因此导数比的任意最终下界或上界，也约束相应的增量比。无穷型先固定 $x$，令 $y\uparrow\omega$，利用
\[
\frac{f(y)}{g(y)}=
\frac{f(x)}{g(y)}+
\left(1-\frac{g(x)}{g(y)}\right)
\frac{f(y)-f(x)}{g(y)-g(x)}
\]
消去固定头部。零比零型则固定 $x$ 后令 $y\uparrow\omega$，增量比直接趋于 $f(x)/g(x)$。再让有限的上下界逼近导数比的上下极限，即得完整链。

两种定理的共同机制是正权平均：离散情形把 $\Delta a$ 写成 $u\,\Delta b$；若 $f,g\in C^1$，连续情形则把 $f'\,\mathrm dx$ 写成 $(f'/g')g'\,\mathrm dx$。无穷型的正权分别为 $\Delta b$ 与 $g'\,\mathrm dx$；零比零型向尾端求和、积分，正权分别为 $-\Delta b$ 与 $-g'\,\mathrm dx$。Cauchy 中值定理使连续版不必额外要求导数连续。

\textbf{逆向需要控制局部振荡。} 在 $x\to+\infty$ 时，$f(x)=Lx+\sin x$、$g(x)=x$ 满足 $f/g\to L$，而 $f'/g'=L+\cos x$ 不收敛。即使取指数分母也不够：令
\[
f(x)=Le^x+\sin(e^{2x}),\qquad g(x)=e^x,
\]
仍有 $f/g\to L$，但 $f'/g'=L+2e^x\cos(e^{2x})$ 无界振荡。离散情形的尺度分离控制的是一个完整步长上的差商，无法自动约束连续情形在任意小区间内的振荡。

\textbf{单调密度定理。} 设 $u:[0,+\infty)\to[0,+\infty)$ 局部可积且最终单调，
\[
F(x)=\int_0^xu(t)\,\mathrm dt\sim Cx^\rho,
\qquad C,\rho>0,\quad x\to+\infty.
\]
则 $u(x)\sim C\rho x^{\rho-1}$。这提供了一个严格的“等价式求导”条件，而不要求 $u$ 连续。

固定 $\lambda>1$。若 $u$ 最终非减，对充分大的 $x$ 有
\[
\frac{F(x)-F(x/\lambda)}{x-x/\lambda}
\leq u(x)\leq
\frac{F(\lambda x)-F(x)}{(\lambda-1)x}.
\]
除以 $Cx^{\rho-1}$ 并取上下极限，左、右两端分别趋于
\[
\frac{1-\lambda^{-\rho}}{1-\lambda^{-1}},
\qquad
\frac{\lambda^\rho-1}{\lambda-1}.
\]
再令 $\lambda\downarrow1$，两者均趋于 $\rho$。最终非增时不等号反向，结论相同。$\rho=1$ 正是平均值收敛加单调性推出函数值收敛；一般的 $\rho$ 则同时恢复增长阶与首项系数。这类借助附加条件从平均恢复局部量的结论，通常称为 Tauber 型结论。

\subsection{从乘积到对数}

\label{app:limit-products}

正数乘积适合取对数，因为乘法误差变成可累积的加法误差。若 $u_k\to0$，则在 $|u_k|\leq1/2$ 时有一致估计
\[
\log(1+u_k)=u_k-\frac{u_k^2}{2}+O(|u_k|^3).
\]
于是有限乘积满足 $\log\prod_{k\leq n}(1+u_k)=\sum_{k\leq n}\log(1+u_k)$。若 $\sum|u_k|^3<\infty$，三阶余项的累计有界；若要证明归一化乘积趋于非零常数，则还要证明对应累计余项收敛。

\textbf{应用。} 固定 $a>-1$，由 $\log(1+a/k)=a/k+O(k^{-2})$ 得
\[
\log\prod_{k=1}^n\left(1+\frac ak\right)-a\log n
=a(H_n-\log n)+\sum_{k=1}^n\left(\log(1+a/k)-a/k\right).
\]
右端收敛，所以乘积除以 $n^a$ 趋于一个正的有限常数。关键不只是 $\log(1+x)\sim x$，而是误差 $O(k^{-2})$ 可求和。用 \hyperref[app:special-gamma]{Gamma 函数}可进一步把常数识别为 $1/\Gamma(1+a)$。



\section{局部展开、插值与 Euler–Maclaurin 公式}

局部展开把函数分成可算的模型与可控的误差。只写出前几项还不是渐近证明；余项的阶数、所需导数及参数一致性，决定展开能否继续用于极限、积分或求和。

\subsection{Taylor 模型、抵消与余项}

\label{app:estimate-taylor}

局部模型的阶数由题设的正则性和所求精度共同决定。若 $f$ 在邻域内有连续的 $m+1$ 阶导数且 $|f^{(m+1)}|\leq M$，则
\[
f(a+h)=\sum_{j=0}^m\frac{f^{(j)}(a)}{j!}h^j+R_m(h),
\qquad |R_m(h)|\leq\frac{M|h|^{m+1}}{(m+1)!}.
\]
当参数、积分变量或展开中心也在变化时，要说明这个 $M$ 对它们一致。只有一点可微，通常只能用 $f(a+h)=f(a)+f'(a)h+o(h)$，不能无说明地补出二阶项。

抵消应发生在估计之前。若 $f\in C^2$，则
\[
f(x+h)+f(x-h)-2f(x)
=\int_{-h}^{h}(h-|t|)f''(x+t)\,\mathrm dt\qquad(h>0).
\]
于是二阶差分的量级和符号都由 $f''$ 控制。分别估计三个函数值会丢掉常数项与一次项的抵消。积分余项还允许只假设导数绝对连续，从而避免不必要的高阶光滑性。

同样，若 $f'$ 绝对连续且 $|f''|\leq M$ 几乎处处，就有 $|f(x+h)-f(x)-hf'(x)|\leq Mh^2/2$。使用更弱假设的意义在于找到真正控制误差的对象，并非一味追求定理的最大一般性。

\subsection{插值、有限差分与误差证书}

\label{app:estimate-interpolation}

对互异节点 $x_0,\ldots,x_m$，次数不超过 $m$ 的插值多项式唯一。若 $f\in C^{m+1}$，节点和 $x$ 都在同一区间内，且 $x$ 不是节点，则
\[
f(x)-p_m(x)=\frac{f^{(m+1)}(\xi)}{(m+1)!}\prod_{j=0}^m(x-x_j)
\]
对某个位于这些点之间的 $\xi$ 成立。证明是给误差减去适当倍数的节点乘积，使所得函数有 $m+2$ 个零点，再反复使用 Rolle 定理。它说明误差由两部分组成：函数的高阶变化，以及节点的几何分布。

等距节点下，定义 $\Delta_h f(x)=f(x+h)-f(x)$。若 $f\in C^m$，反复使用微积分基本定理可得
\[
\Delta_h^m f(x)=\int_{[0,h]^m}f^{(m)}(x+t_1+\cdots+t_m)\,\mathrm dt_1\cdots\mathrm dt_m
\qquad(h>0).
\]
因此 $m$ 阶差分消去所有次数小于 $m$ 的多项式，并继承 $m$ 阶导数的符号。遇到二项式系数交错求和，先把它识别为差分，通常比逐项展开更清楚。

更一般地，若线性误差泛函 $L$ 消去次数小于 $m$ 的多项式，且可与下式的积分交换，Taylor 积分余项给出
\[
L(f)=\int_a^b f^{(m)}(t)K(t)\,\mathrm dt,
\qquad K(t)=L\!\left(\frac{(x-t)_+^{m-1}}{(m-1)!}\right).
\]
这里 $L$ 作用于变量 $x$；在常见的有限次取值与积分组合中，交换由连续性直接保证。核 $K$ 的符号或 $L^1$ 范数就成为求积误差的证书。先检验多项式精确度，再研究核，是统一处理插值与数值积分的路线。

\subsection{Gauss 求积：精确性、正权与余项}

\label{app:estimate-gauss}

设 $n\ge1$，$s_1,\ldots,s_n\in(-1,1)$ 是 Legendre 多项式 $L_n$ 的互异零点。以这些点为节点，定义 Lagrange 基函数及求积公式
\[
\ell_j(x)=\prod_{k\ne j}\frac{x-s_k}{s_j-s_k},\qquad
w_j=\int_{-1}^1\ell_j(x)\,\mathrm dx,\qquad
Q_n(f)=\sum_{j=1}^nw_jf(s_j).
\]
插值直接保证 $Q_n$ 对次数小于 $n$ 的多项式精确；正交性把精确度提高到 $2n-1$。事实上，对 $\deg p\le2n-1$，作多项式除法 $p=qL_n+r$，其中 $\deg q,\deg r<n$。节点上 $L_n=0$，而 $L_n$ 与 $q$ 正交，故
\[
Q_n(p)=Q_n(r)=\int_{-1}^1r(x)\,\mathrm dx
=\int_{-1}^1p(x)\,\mathrm dx.
\]
次数限制不能继续放宽：$Q_n(L_n^2)=0$，但 $\int_{-1}^1L_n^2\,\mathrm dx>0$。因此使用离散正交关系之前，应先检查被积多项式的次数；\hyperref[pr:fin16-high:13]{第十六届决赛的 Legendre 矩阵题} 正是这一原则的应用。

正权也由精确性得到，而无需计算权重公式。因为 $\deg\ell_j^2=2n-2$，
\[
w_j=Q_n(\ell_j^2)=\int_{-1}^1\ell_j(x)^2\,\mathrm dx>0.
\]
又有 $\sum_jw_j=Q_n(1)=2$，故 $|Q_n(f)|\le2\|f\|_\infty$：正权保证求积不会把取值误差任意放大。

对一般函数，则用插值余项把精确度转成误差界。设 $f\in C^{2n}[-1,1]$，取次数不超过 $2n-1$ 的 Hermite 插值多项式 $H$，使 $H(s_j)=f(s_j)$、$H'(s_j)=f'(s_j)$。令 $\Pi_n(x)=\prod_{j=1}^n(x-s_j)$。反复使用 Rolle 定理，计入各节点的二重零点，得到
\[
f(x)-H(x)=\frac{f^{(2n)}(\xi_x)}{(2n)!}\Pi_n(x)^2
\qquad(x\ne s_j),
\]
其中 $\xi_x\in(-1,1)$。因 $Q_n(f)=Q_n(H)=\int_{-1}^1H$，于是
\[
\left|\int_{-1}^1f(x)\,\mathrm dx-Q_n(f)\right|
\le\frac{\|f^{(2n)}\|_\infty}{(2n)!}
\int_{-1}^1\Pi_n(x)^2\,\mathrm dx.
\]
若 $f^{(2n)}\ge0$，则误差 $\int f-Q_n(f)\ge0$；非正时方向相反。平方节点因子同时控制误差大小与符号。这里无需追踪随 $x$ 变化的中间点，只须对高阶导数作一致控制。

\subsection{三种余项与适用边界}

\label{app:taylor-forms}

\textbf{Peano 余项。} 若函数在 $a$ 附近具有直到 $m-1$ 阶导数，且 $f^{(m-1)}$ 在 $a$ 可微，则
\[
f(a+h)=\sum_{j=0}^m\frac{f^{(j)}(a)}{j!}h^j+o(h^m).
\]
它只给局部阶数，不给可计算的常数。适合单点极限，不应未经补充就对变化的展开中心宣称一致。

\textbf{Lagrange 与 Cauchy 余项。} 在足以应用反复中值定理的正则性下，例如 $f\in C^{m+1}$，存在介于 $a$ 与 $x$ 之间的 $\xi$，使
\[
R_m(x)=\frac{f^{(m+1)}(\xi)}{(m+1)!}(x-a)^{m+1}.
\]
另一种形式为 $R_m(x)=f^{(m+1)}(\xi)(x-a)(x-\xi)^m/m!$。中间点通常随 $x$ 改变；它不是能随意取极限或求导的新函数。

\textbf{积分余项。} 若 $f^{(m)}$ 绝对连续，则
\[
R_m(x)=\frac1{m!}\int_a^x(x-t)^mf^{(m+1)}(t)\,\mathrm dt.
\]
它比“存在某个中间点”提供更多结构：可读出符号，应用 Hölder 得到 $L^p$ 导数控制，并直接比较参数一致性。若 $f^{(m+1)}\in L^p$、$1<p\leq\infty$，令 $q$ 为共轭指数，对 $x>a$ 有
\[
|R_m(x)|\leq\frac{(x-a)^{m+1/q}}{m!(mq+1)^{1/q}}
\|f^{(m+1)}\|_{L^p(a,x)}
\]
（$p=\infty$ 时 $q=1$）。这展示了从点态导数界到积分导数界的自然推广。

\subsection{多变量与算子形式}

\label{app:taylor-multivariable}

对 $f:\mathbb R^d\to\mathbb R$，固定方向 $h$，把问题限制到 $g(t)=f(a+th)$。若 $f$ 在连线邻域内 $C^{m+1}$，一元积分余项给出
\[
f(a+h)=\sum_{j=0}^m\frac1{j!}D^jf(a)[h,\ldots,h]
+\frac1{m!}\int_0^1(1-t)^mD^{m+1}f(a+th)[h,\ldots,h]\,\mathrm dt.
\]
若高阶导数的多线性算子范数不超过 $M$，余项便不超过 $M\|h\|^{m+1}/(m+1)!$。这避免逐个控制大量混合偏导。Banach 空间中的 Fréchet 导数在相应连续性下有同样公式。

\textbf{应用：非退化极小点。} 若 $\nabla f(a)=0$ 且 Hessian $H=D^2f(a)$ 正定，存在 $c>0$ 使 $h^THh\geq c\|h\|^2$。二阶 Peano 余项在 $h\to0$ 时是 $o(\|h\|^2)$，故充分小时 $f(a+h)-f(a)\geq c\|h\|^2/4$。起作用的是正定主项对所有方向的统一下界，而非逐方向的正性检验。

\subsection{Euler–Maclaurin：由积分恢复离散修正}

\label{app:taylor-em}

令 Bernoulli 多项式由
\[
\frac{te^{xt}}{e^t-1}=\sum_{m=0}^\infty B_m(x)\frac{t^m}{m!}
\]
在 $t=0$ 附近定义，$B_m=B_m(0)$。最初几项为 $B_1(x)=x-1/2$、$B_2(x)=x^2-x+1/6$，而 $B_2=1/6$、$B_4=-1/30$。周期函数 $\widetilde B_m(x)=B_m(\{x\})$ 记录整数网格内的位置。

对整数 $a<b$，若 $f^{(2p-1)}$ 绝对连续且 $f^{(2p)}\in L^1[a,b]$，则
\[
\begin{aligned}
\sum_{k=a}^bf(k)
={}&\int_a^bf(x)\,\mathrm dx+\frac{f(a)+f(b)}2\\
&+\sum_{r=1}^p\frac{B_{2r}}{(2r)!}
\bigl(f^{(2r-1)}(b)-f^{(2r-1)}(a)\bigr)+R_p,\\
R_p={}&-\frac1{(2p)!}\int_a^b\widetilde B_{2p}(x)f^{(2p)}(x)\,\mathrm dx.
\end{aligned}
\]
最重要的部分是余项界
\[
|R_p|\leq\frac{\sup_{0\leq t\leq1}|B_{2p}(t)|}{(2p)!}
\int_a^b|f^{(2p)}(x)|\,\mathrm dx.
\]
公式不仅给出一串修正项，还说明截断后实际遗漏了多少。步长 $h$ 的版本由 $x=a+ht$ 缩放得到，每增加两阶导数就多出相应的 $h^2$ 因子。

证明起点是每个单位区间上的分部积分：离散梯形和减积分等于 $\int\widetilde B_1f'$。再利用 $B_m'=mB_{m-1}$ 反复分部积分，边界值给出导数修正，剩下最后一个积分。奇数 Bernoulli 数在 $m>1$ 时为零，解释了为何只有奇阶边界导数出现。

\textbf{应用一：调和数。} 对 $f(x)=1/x$ 应用尾部公式，因高阶导数在无穷远可积，得到
\[
H_n=\log n+\gamma+\frac1{2n}-\frac1{12n^2}+O(n^{-4}).
\]
常数 $\gamma=\lim_n(H_n-\log n)$ 可先用单调有界性证明存在；高精度系数由尾部导数给出。为了证明这里的 $O(n^{-4})$，可把公式用到 $p=2$ 并将四阶尾积分一起估计，不能只展示前两项便宣称余项阶数。

\textbf{应用二：幂函数的尾和。} 固定 $s>1$，对 $x^{-s}$ 得
\[
\sum_{k=n+1}^\infty k^{-s}
=\frac{n^{1-s}}{s-1}-\frac1{2n^s}
+\frac{s}{12n^{s+1}}+O_s(n^{-s-3}).
\]
它同时给出审敛、截断误差与加速求和方法。精度目标若仅是等价式，积分比较已经足够；需要相消后的首项时，边界修正往往决定答案。

\subsection{渐近级数不一定收敛}

\label{app:taylor-asymptotic-series}

渐近展开 $f(x)\sim\sum_{k\geq0}a_k\phi_k(x)$ 的含义是：对每个固定 $N$，截断后的误差是 $o(\phi_N(x))$，其中 $\phi_{k+1}=o(\phi_k)$。它不保证对固定 $x$ 把 $N\to\infty$ 后级数收敛。Euler--Maclaurin 与 Stirling 展开经常具有这种性质。

这一区别在应用中很实用：先固定所需精度和截断阶数，再让主变量趋于极限。若要让截断阶数也变化，必须另作一致估计。把“可展开到任意固定阶”误读为“等于无穷级数”，会丢掉整个误差控制的基础。



\section{级数、反常积分与振荡抵消}

审敛首先区分质量与抵消。正项对象没有相消，比较尺度就足够；一般项对象则须追踪部分和或原函数。函数项级数还多了一层要求：若希望交换求和与极限、积分或导数，收敛必须在适当范数中受控。

\subsection{正项级数：比较、比值与根值}

\label{app:series-positive}

对 $a_n\geq0$，部分和单调，所以收敛等价于部分和有界。若最终 $a_n\leq Cb_n$ 且 $\sum b_n<\infty$，则 $\sum a_n$ 收敛；若 $a_n\sim b_n>0$，两级数同敛散。模型 $\sum n^{-p}$ 在且仅在 $p>1$ 时收敛。

根值判别法给出 $L=\limsup_n\sqrt[n]{a_n}$：$L<1$ 时收敛，$L>1$ 时通项甚至不趋零；$L=1$ 时无结论。比值判别的常用充分条件为 $\limsup a_{n+1}/a_n<1$，或 $\liminf a_{n+1}/a_n>1$；需要各项最终为正。二者都在比较指数尺度，所以对 $n^{-p}$ 通常返回临界值 $1$，这时应换成幂次尺度。

例如若 $a_{n+1}/a_n=1-p/n+O(n^{-2})$，取对数并求和可得 $a_n\asymp n^{-p}$，事实上 $n^pa_n$ 趋于正数。于是 $p>1$ 收敛，$p\leq1$ 发散。这里 $O(n^{-2})$ 可求和，是精确结论的依据；仅有 $o(1/n)$ 在 $p=1$ 时不够，因为 $1/(n\log^2n)$ 与 $1/n$ 的比值主项相同，却敛散不同。

\subsection{单调分组与积分判别}

\label{app:series-integral}

若 $f:[1,\infty)\to[0,\infty)$ 单调递减，则每个单位区间上
\[
f(k+1)\leq\int_k^{k+1}f(x)\,\mathrm dx\leq f(k).
\]
求和可知 $\sum f(k)$ 与 $\int_1^\infty f$ 同敛散；收敛时，尾和位于两个相邻的尾积分之间。它既是审敛法，也是误差估计。

对递减 $a_n\geq0$，按 $[2^j,2^{j+1})$ 分组，每组和夹在 $2^ja_{2^{j+1}}$ 与 $2^ja_{2^j}$ 之间。故 Cauchy 凝聚判别给出 $\sum a_n$ 与 $\sum_j2^ja_{2^j}$ 同敛散。于是
\[
\sum_{n\geq2}\frac1{n(\log n)^p}
\quad\text{在且仅在 }p>1\text{ 时收敛}.
\]
重复取对数可以处理更细的临界尺度。单调性是这些分组比较的条件；缺少它时，稀疏尖峰可能让采样值与积分完全失配。

\subsection{一般项级数与 Abel 分部求和}

\label{app:series-dirichlet}

复数项级数满足 Cauchy 准则：收敛当且仅当任意足够远的有限尾和都一致小。绝对收敛保证这一点，并允许重排与换序；条件收敛不能任意重排。令 $A_k=\sum_{j=m}^ka_j$，离散分部积分为
\[
\sum_{k=m}^na_kb_k=A_nb_n+\sum_{k=m}^{n-1}A_k(b_k-b_{k+1}).
\]
若 $|A_k|\leq C$，则右侧由 $C(|b_n|+\sum|b_k-b_{k+1}|)$ 控制。它把通项的振荡放进部分和，把缓慢变化放进差分。

\textbf{Dirichlet 判别法。} 若 $a_n$ 的部分和有界，实数 $b_n$ 单调趋于零，则 $\sum a_nb_n$ 收敛。尾部原函数仍由某个固定常数控制，且 $\sum_{k=m}^{n-1}|b_k-b_{k+1}|=|b_m-b_n|$，所以尾和随 $m\to\infty$ 一致趋零。交错级数是 $a_n=(-1)^n$ 的特例。

\textbf{Abel 判别法。} 若 $\sum a_n$ 收敛，实数 $b_n$ 单调有界，则 $\sum a_nb_n$ 收敛。此时远处所有 $A_k$ 都一致小，而 $b_n$ 及其总变差有界。两判别的区别在于，Dirichlet 要求权重趋零，Abel 则利用原级数已经收敛。

\textbf{应用。} 对 $\theta\notin2\pi\mathbb Z$，几何和给出 $|\sum_{k=1}^Ne^{ik\theta}|\leq2/|1-e^{i\theta}|$。因此 $\sum e^{in\theta}/n^p$ 在 $p>0$ 时收敛；在 $p>1$ 时绝对收敛。若 $\theta$ 限制在远离 $2\pi\mathbb Z$ 的闭集，这个界一致，遂得到一致收敛；靠近共振点则不能沿用同一个常数。

\subsection{反常积分的 Abel–Dirichlet 判别}

\label{app:series-ad-integral}

若 $F(x)=\int_a^xf(t)\,\mathrm dt$ 有界，$g$ 连续可微、单调趋零，则分部积分给出
\[
\int_u^vf(x)g(x)\,\mathrm dx
=(F(v)-F(u))g(v)-\int_u^v(F(x)-F(u))g'(x)\,\mathrm dx.
\]
由 $\int_u^v|g'|=|g(u)-g(v)|$ 得统一尾界，所以积分收敛。这是 Dirichlet 积分判别；单调有界变差的权重也可用 Stieltjes 积分表达。若 $\int_a^\infty f$ 本已收敛，$g$ 仅须单调有界，这给出 Abel 判别。

例如 $\int_1^\infty\sin x/x^p\,\mathrm dx$ 对所有 $p>0$ 收敛，而绝对收敛要求 $p>1$。后一个必要性可在每周期内 $|\sin x|\geq1/2$ 的固定长度子区间上比较，不能仅用 $|\sin x|\leq1$ 得到。

对 $\int_1^\infty e^{ix^\alpha}x^{-\beta}\,\mathrm dx$，其中 $\alpha>0$，先令 $t=x^\alpha$，化为 $\alpha^{-1}\int_1^\infty e^{it}t^{(1-\beta)/\alpha-1}\,\mathrm dt$。Dirichlet 判别在 $\beta>1-\alpha$ 时给出收敛；绝对收敛仍要求 $\beta>1$。相位增长速度决定抵消的强度。

\subsection{一致收敛与逐项运算}

\label{app:series-uniform}

Weierstrass 判别：若 $|f_n(x)|\leq M_n$ 对所有 $x\in E$ 成立，且 $\sum M_n<\infty$，则级数绝对且一致收敛。紧区间上连续函数的一致极限连续，积分也可逐项交换。更一般的积分交换使用\hyperref[app:integration-exchange]{控制收敛与绝对可积性}，不必强求一致收敛。

逐项求导需要更强条件。若 $f_n\in C^1[a,b]$，$\sum f_n(x_0)$ 在一点收敛，且 $\sum f_n'$ 一致收敛，则由
\[
f_n(x)=f_n(x_0)+\int_{x_0}^xf_n'(t)\,\mathrm dt
\]
求和可知原级数一致收敛，和函数可微，导数等于导数级数。仅有原级数一致收敛不能推出导数级数收敛。

幂级数在收敛圆内的每个较小闭圆盘上一致绝对收敛，可以任意逐项求导；边界必须另作判断。若 $\sum a_n$ 收敛，Abel 极限定理给出 $\lim_{r\uparrow1}\sum a_nr^n=\sum a_n$。证明用分部求和把表达式写成部分和的加权平均，权重非负、总质量趋于一且有限头部权重消失。它描述从圆内逼近边界的合法方式，并不保证边界级数在所有角度都收敛。



\section{积分、收敛与正交方法}

\label{app:integration}

本章区分三个常被混用的问题：极限能否穿过积分，局部平均能否恢复函数，以及正交展开能否重建对象。相应工具依次是统一的积分控制、局部化和完备性。先辨认所需的收敛方式，再选择定理，通常能省去大量零散估计。

\subsection{极限交换与一致尾项}

\label{app:integration-exchange}

Lebesgue 控制收敛定理（也称支配收敛定理）的条件是：在固定测度空间上，可测函数 $f_n\to f$ 几乎处处，且存在与 $n$ 无关的 $g\in L^1$ 使 $|f_n|\leq g$。结论不只是积分收敛，还有 $\|f_n-f\|_1\to0$。对计数测度，它化为级数版本：若 $a_{n,k}\to a_k$ 且 $|a_{n,k}|\leq b_k$、$\sum_kb_k<\infty$，则 $\lim_n\sum_ka_{n,k}=\sum_ka_k$。

级数情形的证明展示了一般思路：先选 $K$ 使统一尾界 $\sum_{k>K}b_k$ 足够小，再对有限头部逐项取极限。这也是找不到单一支配函数时的可行替代：证明局部收敛，并独立证明尾部一致小。无界区间上，仅有每个紧区间的一致收敛不够；$\mathbf1_{[n,n+1]}$ 的积分始终为 $1$，质量只是逃向无穷远。

非负可测函数的递增极限可用单调收敛。在两个 $\sigma$ 有限测度空间的乘积上，非负可测函数可用 Tonelli 定理换序；有符号函数若绝对可积，则可用 Fubini 定理换序。本书使用的 Lebesgue 测度与可数集上的计数测度均满足这一条件。非负性防止抵消，绝对可积性则控制抵消后的总量。条件收敛时不可仅凭形式对称性交换求和次序。

参数积分求导也应归结到同一原则。若 $\partial_tf(t,x)$ 在参数邻域存在，且其绝对值由同一可积函数支配，同时某个参数处 $f$ 可积，则对差商用中值定理和支配收敛，得到 $\frac{\mathrm d}{\mathrm dt}\int f(t,x)\,\mathrm dx=\int\partial_tf(t,x)\,\mathrm dx$。关键控制对象是差商，而不仅是极限函数本身。

\subsection{单调收敛、Fatou 与一致可积性}

\label{app:integration-fatou}

三条基本定理可按逻辑组织。若非负可测函数 $f_n\uparrow f$ 几乎处处，则单调收敛定理给出 $\int f_n\uparrow\int f$，允许积分为无穷。它体现了非负质量的逐步累积，不需要一个预先可积的上界。

对任意非负可测 $f_n$，令 $g_n=\inf_{k\geq n}f_k$，则 $g_n\uparrow\liminf f_n$。由单调收敛及 $g_n\leq f_k$（$k\geq n$）得到 Fatou 引理
\[
\int\liminf_nf_n\leq\liminf_n\int f_n.
\]
它只保证下界；尖峰消失或质量逃向无穷远时，严格不等式确实可能发生。若有共同的可积下界，也可先减去下界再应用。

控制收敛可以从 Fatou 推出。若 $f_n\to f$、$|f_n|\leq g\in L^1$，则 $|f|\leq g$，对非负函数 $2g-|f_n-f|$ 使用 Fatou，便得 $\limsup\int|f_n-f|\leq0$。这解释了支配函数的作用：它排除在极限过程中丢失积分质量的可能。

在有限测度空间上，一个更灵活的条件是\textbf{一致可积性}：
\[
\lim_{M\to\infty}\sup_n\int_{\{|f_n|>M\}}|f_n|=0.
\]
若 $f_n\to f$ 依测度且 $\{f_n\}$ 一致可积，则 $f_n\to f$ 于 $L^1$，称为 Vitali 收敛定理。证明可先截断大值，使两侧尾部一致小；有界部分的依测度收敛在有限测度空间上给出 $L^1$ 收敛。极限函数的可积性与尾界由 Fatou 及依测度收敛的几乎处处收敛子列得到。

\textbf{可检查的充分条件。} 若存在 $p>1$ 使 $\sup_n\|f_n\|_p<\infty$，则
\[
\int_{\{|f_n|>M\}}|f_n|\leq M^{1-p}\|f_n\|_p^p\longrightarrow0
\]
一致成立。所以有限测度空间中的高阶矩界能排除集中尖峰。无限测度空间还要防止质量逃向远处；平移的区间示性函数满足上述大值尾条件，却不在 $L^1$ 中趋零。这也是定理中测度空间条件不可省的原因。

\subsection{移动网格与 Riemann 和}

\label{app:integration-riemann}

若 $f$ 在 $[a,b]$ 上 Riemann 可积，任意标记分割只要最大单元长度趋于零，Riemann 和都趋于 $\int_a^bf$。节点可以随 $n$ 移动，权重却必须对应单元长度。若 $f$ 连续，其连续模 $\omega_f$ 给出直接误差界
\[
\left|\sum_jf(\xi_j)|I_j|-\int_a^bf\right|
\leq(b-a)\omega_f\!\left(\max_j|I_j|\right).
\]
一般 Riemann 可积函数应用 Darboux 上下和；它未必能被阶梯函数一致逼近。

若题目中的权重不正好等于单元长度，可先写成标准和与误差两部分。对于有界 $f$，误差由 $\|f\|_\infty\sum_j|w_j-|I_j||$ 控制。这样把几何网格与代数权重分开检查，比把不规则求和直接称为“Riemann 和”更严谨。

含端点奇性的函数须再做截断：在 $[a+\delta,b-\delta]$ 上用通常结论，端点部分给出与网格无关的可积上界，最后令 $\delta\downarrow0$。可积性本身并不保证任意采样点都安全；采样点若极端靠近奇点，单个求和项就可能破坏收敛。

\subsection{局部化与近似恒等算子}

\label{app:integration-local}

连续函数可由短区间平均恢复：$h^{-1}\int_x^{x+h}f(t)\,\mathrm dt\to f(x)$。于是，对所有短区间成立的积分不等式可在连续点转成逐点不等式。只有局部可积性时，相应结论是 Lebesgue 微分定理所保证的几乎处处结论，不能任意指定点值。

同一思想可写成卷积。设 $K\geq0$ 可积，$\int_{\mathbb R}K=1$，令 $K_\varepsilon(t)=\varepsilon^{-1}K(t/\varepsilon)$。对有界函数 $f$ 的连续点 $x$，
\[
(K_\varepsilon*f)(x)-f(x)
=\int_{\mathbb R}K_\varepsilon(t)\bigl(f(x-t)-f(x)\bigr)\,\mathrm dt\longrightarrow0.
\]
证明将积分分为 $|t|<\delta$ 与其补集：前者由连续性小，后者由核的尾质量小。若 $f\in L^p$、$1\leq p<\infty$，利用平移在 $L^p$ 中连续，同样得到 $\|K_\varepsilon*f-f\|_p\to0$。

因此，“核越来越尖”只是图像；实际需要的是总质量归一、质量集中和足够的范数控制。若核允许变号，可用一致的 $L^1$ 范数界及绝对尾质量趋零替代非负性。没有这些条件，窄峰也可能放大误差。

\subsection{Hilbert 空间、投影与强收敛}

\label{app:integration-hilbert}

约定复内积对第一变量线性。在 Hilbert 空间 $H$ 中，任意闭线性子空间 $W$ 都有正交投影：每个 $v$ 唯一写成 $v=P_Wv+r$，$r\perp W$。其核心恒等式是
\[
\|v-w\|^2=\|r\|^2+\|P_Wv-w\|^2\qquad(w\in W).
\]
存在性可由极小化序列和平行四边形恒等式证明：接近最小距离的两个点必须相互接近，完备性给出极限，闭性确保极限仍在 $W$ 中。沿 $W$ 的任意方向求一次变分，便得正交性。有限维子空间自动闭；无限维时不能省略这个条件。

设 $e_1,e_2,\ldots$ 正交归一，$P_nv=\sum_{j\leq n}\langle v,e_j\rangle e_j$。投影恒等式给出 Bessel 不等式 $\sum_j|\langle v,e_j\rangle|^2\leq\|v\|^2$；只有当这族向量的线性张成稠密时，才进一步得到 $P_nv\to v$ 和 Parseval 等式。正交性负责计算，完备性负责重建。

若 $v_n$ 弱收敛到 $v$，且 $\|v_n\|\to\|v\|$，则
\[
\|v_n-v\|^2=\|v_n\|^2+\|v\|^2-2\operatorname{Re}\langle v_n,v\rangle\longrightarrow0.
\]
这给出由弱收敛升级为强收敛的标准路线。在 $L^2$ 中，先确认候选极限属于 $L^2$，再证明弱收敛或内积收敛，最后比较范数；不能把逐点收敛直接当作弱收敛。

\subsection{Fourier 正交性与加法计数}

\label{app:integration-fourier}

在 $[0,1]$ 上，整数频率满足 $\int_0^1e^{2\pi ikx}\,\mathrm dx=\mathbf1_{\{k=0\}}$。因此有限集 $S\subset\mathbb Z$ 的指数和 $F(x)=\sum_{n\in S}e^{2\pi inx}$ 有
\[
\int_0^1|F|^2=|S|,\qquad
\int_0^1|F|^4=\#\{(a,b,c,d)\in S^4:a+b=c+d\}.
\]
把乘积展开后，积分消去所有非零总频率，剩下的是有序四元组计数。若二元和除交换次序外表示唯一，结果为 $2|S|^2-|S|$：两种配对各给 $|S|^2$，四个元素相等的 $|S|$ 项重复了一次。这个量称为加法能量；它衡量频率集的加法结构，并可与\hyperref[app:estimate-moments]{矩插值}组合来估计 $\int|F|$。

对一般 $L^2$ 函数，三角系统的完备性使上述有限计算扩展为 Fourier 级数的均方收敛。均方收敛不自动意味着逐点收敛，更不意味着一致收敛；题目若要求点值，须另加正则性或采用适当的求和方式。

\subsection{振荡积分与反演的条件}

\label{app:integration-inversion}

若 $g\in L^1(\mathbb R)$，Riemann--Lebesgue 引理断言 $\int g(x)e^{itx}\,\mathrm dx\to0$。可先对区间示性函数直接积分，再用有限阶梯函数的 $L^1$ 稠密性逼近；误差统一由 $L^1$ 范数控制。这一证明也揭示条件：奇点附近有界不能代替整条实线可积。

若 $g$ 绝对连续，$g'\in L^1$ 且两端边界项消失，分部积分还能给出 $O(|t|^{-1})$ 的速度。对一般相位 $\phi$，在 $\phi'$ 不为零的区域写 $e^{it\phi}=(it\phi')^{-1}(e^{it\phi})'$；相位驻点附近则需单独缩放。振荡衰减来自抵消，先取绝对值会消灭这一信息。

采用 $\widehat f(\xi)=\int_{\mathbb R}f(x)e^{-ix\xi}\,\mathrm dx$ 的约定。一个简洁充分条件是 $f,\widehat f\in L^1$，此时反演式
\[
f(x)=\frac1{2\pi}\int_{\mathbb R}\widehat f(\xi)e^{ix\xi}\,\mathrm d\xi
\]
在 $f$ 的每个连续点成立，并在几乎处处意义下恢复 $f$。证明可先在频率端乘 Gaussian 衰减因子，使换序有绝对可积保证；空间端成为近似恒等核，最后用两边各自的收敛定理去掉因子。它把反演分成三个可检查步骤：正则化、换序、恢复原函数。



\section{渐近估计与尺度选择}

\label{app:estimates}

分析估计首先识别控制变化的尺度，再选择在这一尺度上有效的界。余项必须服从最终操作：求和要可累积，积分要一致可积，求极值要一致小。本章从振幅与导数的平衡出发，再讨论缩放、紧性、递推与积分集中。

\subsection{振幅与局部导数的两尺度估计}

\label{app:estimate-holder}

同一个量有两个上界时，应保留其中较小者。若 $Q\le A r^a$ 且 $Q\le B d r^{-c}$，其中 $a,c,A,B,d,r>0$，则
\[
Q\le\min\{A r^a,Bd r^{-c}\}
\le A^{c/(a+c)}B^{a/(a+c)}d^{a/(a+c)}.
\]
第二步是 $\min(U,V)\le U^\theta V^{1-\theta}$，取 $\theta=c/(a+c)$ 消去 $r$。等价地，解平衡方程 $A r^a=Bd r^{-c}$ 得到临界尺度。若两个界只在不同区域成立，就按这一尺度分区，分别使用各自有效的估计。

\textbf{振幅---导数准则。} 设 $a,b>0$，$u\in C[0,1]$，且 $u$ 在每个 $[\varepsilon,1]$ 上绝对连续。若存在常数 $A,B\ge0$，使
\[
|u(t)-u(0)|\le At^a,\qquad
|u'(t)|\le Bt^{a-b-1}\quad\text{几乎处处于 }(0,1),
\]
则存在与 $x,y$ 无关的常数 $C$，使
\[
|u(y)-u(x)|\le C|y-x|^\gamma,
\qquad \gamma=\min\left\{1,\frac a{b+1}\right\}.
\]

当 $a\ge b+1$ 时，$|u'|\le B$；先在 $0<x<y$ 上积分，再令 $x\downarrow0$，即得 Lipschitz 估计。设 $0<a<b+1$，记 $d=y-x$、$r=d^{1/(b+1)}$。对 $0\le x<y\le1$，有 $d\le r$。若 $x\le r$，则 $y\le2r$，振幅界给出
\[
|u(y)-u(x)|\le A(x^a+y^a)
\le A(1+2^a)r^a.
\]
若 $x>r$，则因 $a-b-1<0$，局部绝对连续性给出
\[
|u(y)-u(x)|\le B\int_x^yt^{a-b-1}\,\mathrm dt
\le Bd\,x^{a-b-1}\le Bd\,r^{a-b-1}.
\]
两者都是 $C d^{a/(b+1)}$。这里导数允许在零点附近不可积：靠近零点的一段已由振幅控制，无须跨过奇点积分。

\textbf{快速振荡的最优指数。} 对 $a,b>0$，令
\[
F_{a,b}(0)=0,\qquad F_{a,b}(x)=x^a\sin(x^{-b})\quad(0<x\le1).
\]
由 $|F_{a,b}(x)|\le x^a$ 及
\[
|F_{a,b}'(x)|
\le a x^{a-1}+b x^{a-b-1}
\le(a+b)x^{a-b-1},
\]
上述准则给出指数 $\gamma=\min\{1,a/(b+1)\}$，且端点指数本身可取。它也是最大的统一两点指数：若 $a<b+1$，取
\[
x_n=(2\pi n+\pi/2)^{-1/b},\qquad
y_n=(2\pi n+3\pi/2)^{-1/b}.
\]
两点的正弦值分别为 $1,-1$。于是
\[
|F_{a,b}(x_n)-F_{a,b}(y_n)|\asymp n^{-a/b},
\qquad |x_n-y_n|\asymp n^{-(b+1)/b},
\]
其中第二式由 $t\mapsto t^{-1/b}$ 的中值定理得到；符号 $\asymp$ 表示比值被两个正常数夹住。因此任意 $\beta>a/(b+1)$ 都使 Hölder 商趋于无穷。若 $a\ge b+1$，指数已经达到 $1$；区间上满足 $|u(y)-u(x)|\le C|y-x|^\beta$、$\beta>1$ 的函数必为常数，因为把 $[x,y]$ 等分为 $N$ 段后，右侧的总和是 $C|y-x|^\beta N^{1-\beta}\to0$。$F_{a,b}$ 非常数，故不能再提高指数。

这一区别属于两点几何：距零点约 $x$ 处，振幅为 $x^a$，一次振荡的距离为 $x^{b+1}$，两者决定统一指数。只与零点比较则有 $|F_{a,b}(x)|\le x^a$，并由正弦取值为 $1$ 的序列知这一幂次不能提高；这并不意味着任意两点都满足 $a$ 阶估计。上述结论限于 $[0,1]$；扩展到无界区间还须单独控制无穷远处的增长。

\subsection{导数积分与连续模}

\label{app:estimate-derivative}

对一般函数，直接对导数积分使用 Hölder 不等式是系统方法，但前提是函数的增量能由导数积分恢复。设 $u$ 在有限闭区间 $I$ 上绝对连续，$u'\in L^p(I)$，$1<p<\infty$。则对 $x<y$，
\[
\begin{aligned}
|u(y)-u(x)|
&=\left|\int_x^yu'(t)\,\mathrm dt\right|\\
&\le |y-x|^{1-1/p}
\left(\int_x^y|u'(t)|^p\,\mathrm dt\right)^{1/p}
\le\|u'\|_{L^p(I)}|y-x|^{1-1/p}.
\end{aligned}
\]
这就是一维 Morrey 估计，即 $W^{1,p}(I)$ 的绝对连续代表属于 $C^{0,1-1/p}(I)$。$p=\infty$ 时直接得到 Lipschitz 界；特别地，$p=2$ 给出指数 $1/2$。绝对连续性不能只换成“几乎处处可导”：Cantor 函数的导数几乎处处为零，却不是常数。

$p=1$ 时仍有连续模
\[
\omega_u(h):=\sup_{\substack{x,y\in I\\|x-y|\le h}}|u(x)-u(y)|
\le\sup_{\substack{J\subset I\ \text{为区间}\\|J|\le h}}
\int_J|u'(t)|\,\mathrm dt\longrightarrow0,
\]
最后一步是 Lebesgue 积分的绝对连续性；但一般不能得到正的幂次。例如 $u(0)=0$、$u(x)=1/\log(e/x)$（$0<x\le1$）的导数为 $1/[x\log^2(e/x)]\in L^1(0,1)$，且 $u(x)=\int_0^xu'(t)\,\mathrm dt$；然而对每个 $\alpha>0$，$u(x)/x^\alpha\to\infty$。

\textbf{何时应改用分区估计。} 对 $f(x)=x\sin(1/x)$，
\[
f'(x)=\sin(1/x)-\frac{\cos(1/x)}x,
\qquad
\int_0^1\frac{|\cos(1/x)|}{x}\,\mathrm dx
=\int_1^\infty\frac{|\cos t|}{t}\,\mathrm dt=\infty.
\]
最后一个积分在每段周期上的贡献与 $1/n$ 同阶，而 $\sin(1/x)$ 在 $(0,1)$ 上绝对可积，故 $f'\notin L^1(0,1)$，也不属于任何 $L^p(0,1)$（$p>1$）。因此不能在包含零点的区间上直接套用上述估计。对 $x>0$，在 $[x,y]$ 上积分完全合法；再与近零点的振幅界配合，上一节仍能给出最优的 $1/2$ 阶估计。导数的积分范数提供方便的充分条件，振幅与局部尺度则能保留它看不到的振荡抵消。

\subsection{缩放、紧性与极值扰动}

\label{app:estimate-scaling}

齐次结构提示先把变化的区域缩放到固定集合。例如 $F(Ru)=R^mP(u)+R^{m-1}Q(u)$，除以 $R^m$ 后只剩 $P+\varepsilon Q$，其中 $\varepsilon=R^{-1}$。这一步把“无穷远问题”变成紧集上的小扰动问题。

设 $K$ 是非空紧集，$P,Q$ 连续，$M=\max_KP$，$E=\{u\in K:P(u)=M\}$。则当 $\varepsilon\downarrow0$ 时，
\[
\max_K(P+\varepsilon Q)=M+\varepsilon\max_EQ+o(\varepsilon).
\]
证明体现两层选择。取极大点 $u_\varepsilon$，与 $E$ 中使 $Q$ 最大的点比较，利用 $Q$ 有界先得 $P(u_\varepsilon)\to M$；故 $u_\varepsilon$ 的每个聚点都在 $E$。又有
\[
\max_EQ\leq\frac{\max_K(P+\varepsilon Q)-M}{\varepsilon}
\leq Q(u_\varepsilon).
\]
右侧的上极限不超过 $\max_EQ$，结论随之成立。极小值同理，只需同时把两个最大值改为最小值。

若另有一致余项 $r_\varepsilon$，满足 $\sup_K|r_\varepsilon|=o(\varepsilon)$，结论不变。这里“一致”不能删去：极值点本身随参数移动，点态余项无法控制它。主项负责选出候选集，次项再在候选集中选点；这也解释了为何只比较最高次项不足以求出次阶极限。

\subsection{递推的坐标选择与累积误差}

\label{app:estimate-recursion}

设正数列已知趋于零，且
\[
x_{n+1}=x_n-cx_n^{1+\alpha}+o(x_n^{1+\alpha}),\qquad c,\alpha>0.
\]
增量在原坐标中不断缩小。选择 $y_n=x_n^{-\alpha}$，由 $(1-u)^{-\alpha}=1+\alpha u+o(u)$ 得
\[
y_{n+1}-y_n\longrightarrow\alpha c.
\]
将差分求和并取 Cesàro 平均，得到 $y_n/n\to\alpha c$，即
\[
x_n\sim(\alpha cn)^{-1/\alpha}.
\]
坐标选择可由连续模型 $x'=-cx^{1+\alpha}$ 预见，但严格证明仍是离散差分求和。应先独立证明正性、轨道留在所需区间以及 $x_n\to0$，再使用局部展开。

若需要下一阶，应继续研究变换后的增量。例如 $y_{n+1}-y_n=A+B/y_n+O(y_n^{-2})$ 且 $y_n\sim An$、$A>0$ 时，先求得 $y_n=An+O(\log n)$，代回后便有
\[
y_n=An+\frac BA\log n+O(1).
\]
原因是 $1/y_n-1/(An)=O(\log n/n^2)$ 可求和，而 $\sum_{k<n}1/k=\log n+O(1)$。局部误差必须经过求和审查：每步 $o(1)$ 只保证总误差 $o(n)$，不能自动得到有界总误差。

\subsection{Laplace 方法：质量集中与二次尺度}

\label{app:estimate-laplace}

设 $\phi\in C^2[a,b]$ 在内部点 $x_0$ 有唯一最大值，$\phi''(x_0)<0$，$a(x)$ 连续且 $a(x_0)\neq0$。则
\[
\int_a^ba(x)e^{n\phi(x)}\,\mathrm dx
\sim a(x_0)e^{n\phi(x_0)}
\sqrt{\frac{2\pi}{n|\phi''(x_0)|}}.
\]
证明先在 $x_0$ 的固定邻域外用紧性得到 $\phi\leq\phi(x_0)-\eta$，这部分指数级小；邻域内令 $x=x_0+u/\sqrt n$，二阶 Taylor 模型成为 $e^{-cu^2/2}$。由 $\phi(x)-\phi(x_0)\leq-c_1(x-x_0)^2$ 给出可积 Gaussian 支配，再控制缩放区间外的尾部，才可合法取极限。

这解释了 $n^{-1/2}$ 从何而来：指数中 $n(x-x_0)^2$ 保持常数量级的窗口宽度正是 $n^{-1/2}$。若最高点退化为 $\phi(x)=\phi(x_0)-c|x-x_0|^m+o(|x-x_0|^m)$，局部尺度相应变成 $n^{-1/m}$；若最大值落在端点且一阶导数非零，尺度通常为 $n^{-1}$。

\textbf{应用。} $\int_0^1[x(1-x)]^n\,\mathrm dx$ 的最大点为 $1/2$。对 $\phi=\log(x(1-x))$，有 $\phi(1/2)=-\log4$、$\phi''(1/2)=-8$。由于 $\phi$ 在端点不光滑，先固定 $0<\delta<1/2$，在 $[\delta,1-\delta]$ 上应用上述结论；两个端点区间的总贡献至多为 $2\delta[\delta(1-\delta)]^n$，相对于主项指数级小。因此原积分等价于 $4^{-n}\sqrt\pi/(2\sqrt n)$。也可用 Beta 函数与 Stirling 公式核对，见\hyperref[app:special-beta]{特殊函数}。两种方法分别从局部集中和精确恒等式解释同一个尺度。

\subsection{尾部渐近与 Watson 型展开}

\label{app:estimate-watson}

设 $a$ 在 $[0,\infty)$ 上可测且局部可积，$a(t)=\sum_{k=0}^mc_kt^k+O(t^{m+1})$ 在 $0\leq t\leq\delta$ 成立，并存在常数 $C>0$、$c\geq0$、$T\geq\delta$，使 $|a(t)|\leq Ce^{ct}$ 对所有 $t\geq T$ 成立。则当 $x\to+\infty$ 时
\[
\int_0^\infty e^{-xt}a(t)\,\mathrm dt
=\sum_{k=0}^m\frac{c_kk!}{x^{k+1}}+O(x^{-m-2}).
\]
证明将积分在 $\delta$ 处分开；近端余项的绝对值由 $C\int_0^\infty e^{-xt}t^{m+1}\,\mathrm dt$ 控制。有限中段 $[\delta,T]$ 由局部可积性控制，无穷尾段由上述指数增长界控制，两者都指数级小。分数幂 $t^{\alpha+k-1}$（$\alpha>0$）同理产生 $\Gamma(\alpha+k)x^{-\alpha-k}$。这是 Watson 引理的一种常用形式。

\textbf{应用。} 换元 $t=x+u$ 后，Gaussian 尾为
\[
\int_x^\infty e^{-t^2}\,\mathrm dt
=e^{-x^2}\int_0^\infty e^{-2xu-u^2}\,\mathrm du
=\frac{e^{-x^2}}{2x}\left(1-\frac1{2x^2}+O(x^{-4})\right).
\]
在近端展开 $e^{-u^2}$，远端用指数衰减控制。若只需上界，直接由 $t/x\geq1$ 得 $\int_x^\infty e^{-t^2}\,\mathrm dt\leq e^{-x^2}/(2x)$；无需使用完整展开。

\subsection{分区估计与一致参数}

\label{app:estimate-split}

遇到奇点、端点或移动峰，先把区域分成不同机制成立的部分。典型次序是：固定一个小邻域；在邻域内缩放并控制模型余项；在补集证明统一衰减；最后让邻域大小趋零。若把邻域大小过早设成 $n$ 的函数，容易把本来只对固定区域成立的估计误当作一致估计。

\textbf{反例提示。} $f_n(x)=n\mathbf1_{(0,1/n)}(x)$ 几乎处处趋零，积分却恒为 $1$。任何在每个固定 $x>0$ 处成立的展开，都看不到靠近零的主贡献。这正说明局部尺度与积分极限必须一起设计。实际计算中，可以先用图像或量纲猜尺度，但最后应写出支配、尾界或显式误差中的至少一种。



\section{不等式、对偶与矩方法}

常用不等式可按结构组织：凸性给出支撑界，内积给出正交分解，范数对偶把乘积变成可控的积。每次使用应同时检查指数匹配、等号条件，以及估计是否保留所需尺度。

\subsection{凸性、对偶界与等号结构}

\label{app:estimate-convexity}

可微凸函数满足切线不等式 $\phi(y)\geq\phi(x)+\phi'(x)(y-x)$。令 $x=\sum_jw_jx_j$，其中 $w_j\geq0$、$\sum_jw_j=1$，求和便得到 Jensen 不等式
\[
\phi\!\left(\sum_jw_jx_j\right)\leq\sum_jw_j\phi(x_j).
\]
严格凸时，所有正权对应的 $x_j$ 相等才可能取等。连续版本同样成立，但须保证相关期望存在。凸性把多变量不等式转化为一个局部支撑关系，并同时解释等号。

Young 不等式 $ab\leq a^p/p+b^q/q$（$a,b\geq0$，$p,q>1$，$1/p+1/q=1$）来自对 $a$ 的凸函数求极值。归一化后积分，便得到 Hölder 不等式。使用时应先选共轭指数来匹配已知范数，而非盲目把每项都估大。

积分或矩阵不等式中的“最佳常数”常来自对偶证书：构造一个已知非负的平方、正半定矩阵或支撑函数，使目标差恰等于该对象。求等号就是令证书消失；这比在一长串松弛中追踪等号更可靠。二次型形式见Gram 矩阵与正性。

\subsection{范数插值与矩方法}

\label{app:estimate-moments}

在同一测度空间上，若 $1\leq p<r<q<\infty$，$1/r=\theta/p+(1-\theta)/q$，且 $f\in L^p\cap L^q$，对 $|f|^{\theta r}|f|^{(1-\theta)r}$ 使用 Hölder 得
\[
\|f\|_r\leq\|f\|_p^\theta\|f\|_q^{1-\theta}.
\]
特别地，$p=1,r=2,q=4$ 时，非零函数满足 $\|f\|_1\geq\|f\|_2^3/\|f\|_4^2$。这把绝对值积分的下界转化为二阶、四阶矩的计算；振荡和中这两个矩常可由\hyperref[app:integration-fourier]{正交性计数}求出。

在概率空间上，若 $X\geq0$ 且 $0<\mathbb EX^2<\infty$，则对 $0<\eta<1$，
\[
\mathbb P(X>\eta\mathbb EX)\geq(1-\eta)^2\frac{(\mathbb EX)^2}{\mathbb EX^2}.
\]
设 $E=\{X>\eta\mathbb EX\}$。从 $(1-\eta)\mathbb EX\leq\mathbb E(X\mathbf1_E)$ 出发，用 Cauchy--Schwarz 即得。均值说明总质量有多少，二阶矩限制它集中在极少数点上的程度；两者共同迫使一个具有正概率的事件出现。

\subsection{内积不等式、加权配方与最优常数}

\label{app:inequality-cauchy}

Cauchy--Schwarz 来自 $\|u-tv\|^2\geq0$ 对标量 $t$ 的极小化。因此等号恰在两向量线性相关时成立，复数情形需相应使用共轭。加权形式
\[
\left|\sum_ka_kb_k\right|^2
\leq\left(\sum_kw_k|a_k|^2\right)
\left(\sum_k\frac{|b_k|^2}{w_k}\right),\qquad w_k>0,
\]
只是把权重的平方根分配到两个向量。选权重的目的，是使一侧恰好等于题设已知的量。

\textbf{应用。} 若 $x_i>0$、$\sum_i x_i=s$，则 $\sum_i a_i^2/x_i\geq(\sum_i|a_i|)^2/s$。对 $|a_i|/\sqrt{x_i}$ 与 $\sqrt{x_i}$ 使用上式即得；当所有 $a_i\neq0$ 时，最佳 $x_i$ 与 $|a_i|$ 成比例。这同时给出不等式、最佳常数与极值点，无需分别用多变量微分求解。

\subsection{Hölder、Minkowski 与对偶}

\label{app:inequality-holder}

对共轭指数 $p,q>1$，Hölder 给出 $\int|fg|\leq\|f\|_p\|g\|_q$。其等号条件是非零情形下 $|f|^p$ 与 $|g|^q$ 成比例；若还涉及复数积分的绝对值，须检查相位一致。由此可得
\[
\|f\|_p=\sup_{\|g\|_q\leq1}\left|\int f\overline g\right|,
\]
取 $g=|f|^{p-2}f/\|f\|_p^{p-1}$ 即达等。这个对偶式直接推出 Minkowski 三角不等式，并解释为何把困难的范数问题转成线性配对有利。

在总质量为 $1$ 的空间上，若 $1\leq p<q$，则 $\|f\|_p\leq\|f\|_q$。一般有限测度空间须补上因子 $\mu(E)^{1/p-1/q}$；在无限测度空间上没有这种无条件包含关系。尺度因子缺失往往说明归一化尚未完成。

\subsection{凸序、重排与主序}

\label{app:inequality-majorization}

将实向量的分量递减排列。若 $\sum_{i=1}^kx_i\geq\sum_{i=1}^ky_i$ 对 $k<n$ 成立，且总和相等，称 $x$ 主序于 $y$。它表示 $x$ 更分散。Karamata 不等式断言对凸函数 $\phi$ 有 $\sum_i\phi(x_i)\geq\sum_i\phi(y_i)$。

可把这看成反复平均的结果：用两个分量的保和混合替换原分量，凸函数和不会增加；一般主序关系可分解为有限次这种混合，这是主序定理的内容。竞赛中常不必引用完整定理，只要显式写出题目所需的一两次凸组合，应用 Jensen 即可。

另一类顺序信息是重排不等式：对同序实数列 $a_1\leq\cdots\leq a_n$、$b_1\leq\cdots\leq b_n$，同序配对的乘积和最大，逆序最小。交换一个逆序对使总和增加 $(a_j-a_i)(b_j-b_i)\geq0$，反复交换即证。这个证明说明，最优配对来自每次局部交换的符号，而非偶然的代数展开。

\subsection{积分型能量不等式}

\label{app:inequality-energy}

若绝对连续函数 $f$ 在 $[a,b]$ 满足 $f(a)=0$ 且 $f'\in L^2$，则
\[
|f(x)|^2=\left|\int_a^xf'(t)\,\mathrm dt\right|^2
\leq(x-a)\int_a^x|f'(t)|^2\,\mathrm dt.
\]
再积分即可得到一个 Poincaré 型界 $\|f\|_2\leq(b-a)\|f'\|_2/\sqrt2$。若两端都为零，最佳常数为 $(b-a)/\pi$，对应第一特征函数 $\sin(\pi(x-a)/(b-a))$；可由正弦级数的 Parseval 等式证明。

因此最佳不等式常与谱问题相连：最小 Rayleigh 商 $\int|f'|^2/\int|f|^2$ 就是相应边界条件下的最低特征值。备考中先用基本积分与 Cauchy--Schwarz 得到够用的界；题目追问最佳常数时，再寻找等号函数或相关谱结构。



\section{含参积分与积分变换}

参数能把难积分嵌入一个更容易微分或满足简单方程的函数族。积分变换则把同一思想系统化：先用一族核测试原函数，再让微分、卷积或缩放变成新变量上的代数运算。每一步都须保留定义域和交换运算的条件。

\subsection{连续性、求导与 Leibniz 公式}

\label{app:transform-parameter}

设 $F(t)=\int_Ef(t,x)\,\mathrm d\mu(x)$。若 $f(t,x)$ 对 $t$ 连续且在参数的每个紧邻域内由同一 $L^1$ 函数支配，控制收敛给出 $F$ 连续。若偏导存在且 $|\partial_tf(t,x)|\leq g(x)$、$g\in L^1$，对差商应用同样论证得到 $F'(t)=\int_E\partial_tf(t,x)\,\mathrm d\mu$。可反复求导时，每一阶都要有对应控制。

有限区间上，若 $f$ 与 $\partial_tf$ 连续，$a(t),b(t)$ 可微，则
\[
\frac{\mathrm d}{\mathrm dt}\int_{a(t)}^{b(t)}f(t,x)\,\mathrm dx
=f(t,b(t))b'(t)-f(t,a(t))a'(t)+\int_{a(t)}^{b(t)}\partial_tf(t,x)\,\mathrm dx.
\]
边界项来自积分区间本身的变化。对于反常端点，应先截断再证明所有项的极限存在，不能直接照抄有限区间公式。

\textbf{应用。} 固定 $a>0$，令 $F(b)=\int_0^\infty e^{-ax}\sin(bx)/x\,\mathrm dx$。求导后的被积函数由 $e^{-ax}$ 支配，所以
\[
F'(b)=\int_0^\infty e^{-ax}\cos(bx)\,\mathrm dx=\frac a{a^2+b^2}.
\]
由 $F(0)=0$ 得 $F(b)=\arctan(b/a)$。若再令 $a\downarrow0$，需要额外的统一振荡尾界；对 $b>0$，Dirichlet 分部积分给出尾部 $O(1/(bR))$，从而恢复 $\int_0^\infty\sin(bx)/x\,\mathrm dx=\pi/2$。最后一步并没有一个显然的全区间绝对可积支配函数。

\subsection{参数化差值与 Frullani 型积分}

\label{app:transform-frullani}

对 $a,b>0$，利用
\[
\frac{e^{-ax}-e^{-bx}}x=\int_a^be^{-tx}\,\mathrm dt
\]
可在有限参数区间内换序，得到
\[
\int_0^\infty\frac{e^{-ax}-e^{-bx}}x\,\mathrm dx
=\int_a^b\frac{\mathrm dt}{t}=\log\frac ba.
\]
当 $b>a$ 时可直接用 Tonelli；反向次序改变符号即可。这一表示让原来两项各自发散的差保持在同一个积分中。

一般地，若 $f$ 连续、在零与无穷处有有限极限，通过截断换元可得 Frullani 公式 $\int_0^\infty[f(ax)-f(bx)]\,\mathrm dx/x=(f(0)-f(\infty))\log(b/a)$。换元后仅余零端与无穷端的短区间积分，两端极限同时保证反常积分存在。一组便于直接换序的充分条件是 $f$ 局部绝对连续且 $f'\in L^1(0,\infty)$：写差为 $f(ax)-f(bx)=-\int_a^bxf'(tx)\,\mathrm dt$，绝对可积的二重积分保证换序。端点抵消与差值符号都必须保留。

\subsection{Laplace 变换与微分方程}

\label{app:transform-laplace}

对 $t\geq0$ 上的函数，定义 $\mathcal Lf(s)=\int_0^\infty e^{-st}f(t)\,\mathrm dt$。若 $f$ 局部可积且 $|f(t)|\leq Me^{ct}$，则 $\operatorname{Re}s>c$ 时绝对收敛。若 $f$ 绝对连续且相应边界项消失，分部积分给出
\[
\mathcal L(f')(s)=s\mathcal Lf(s)-f(0).
\]
因而常系数微分方程连同初值变成关于 $s$ 的代数方程。因果卷积 $(f*g)(t)=\int_0^tf(u)g(t-u)\,\mathrm du$ 在绝对可积条件下满足 $\mathcal L(f*g)=\mathcal Lf\,\mathcal Lg$。

\textbf{应用。} 对 $y'+ay=f$、$y(0)=y_0$，变换给出 $Y=(y_0+F)/(s+a)$。恢复后为
\[
y(t)=e^{-at}y_0+\int_0^te^{-a(t-u)}f(u)\,\mathrm du.
\]
这是积分因子公式；变换把它解释为初值响应与卷积响应。竞赛中若只需一阶方程，直接积分因子最短；面对高阶常系数、卷积项或分段外力时，变换的组织作用更明显。

Laplace 变换的唯一性和复反演公式需适当增长与正则性假设，可查专门教材。不能仅凭若干参数处的变换值就声称函数唯一；唯一性使用整个开半平面或足够丰富的参数集合。

\subsection{Fourier 变换、微分与卷积}

\label{app:transform-fourier}

约定 $\widehat f(\xi)=\int_{\mathbb R}f(x)e^{-ix\xi}\,\mathrm dx$，反演系数为 $1/(2\pi)$。在 Schwartz 空间（函数及各阶导数都比任意幂衰减得快）中，各项绝对可积，分部积分与求导给出
\[
\widehat{f'}(\xi)=i\xi\widehat f(\xi),\qquad
\widehat{xf}(\xi)=i\frac{\mathrm d}{\mathrm d\xi}\widehat f(\xi),\qquad
\widehat{f*g}=\widehat f\,\widehat g.
\]
平移变成相位，缩放改变频率尺度。一般 $L^1$ 情形仍有变换与卷积公式；$L^2$ 情形通过 Plancherel 定理延拓，满足 $\|\widehat f\|_2^2=2\pi\|f\|_2^2$。这一延拓是均方意义的，不能随意把原始积分当成逐点绝对收敛。

\textbf{Gaussian 核。} 令 $G_t(x)=(4\pi t)^{-1/2}e^{-x^2/(4t)}$，$t>0$。配方或参数微分可得 $\widehat G_t(\xi)=e^{-t\xi^2}$，故 $G_s*G_t=G_{s+t}$。因此 $u(t)=G_t*f$ 的 Fourier 变换满足 $\partial_t\widehat u=-\xi^2\widehat u$，在适当条件下 $u_t=u_{xx}$。复杂的空间扩散变成每个频率独立衰减的标量方程。

Gaussian 还可用于反演的正则化：先乘 $e^{-\varepsilon\xi^2}$，让换序绝对合法；空间端则卷积 $G_\varepsilon$，最后利用\hyperref[app:integration-local]{近似恒等算子}恢复原函数。这是一种可反复使用的“先控制，再取极限”的策略。

\subsection{Mellin 变换与乘法尺度}

\label{app:transform-mellin}

Mellin 变换定义为 $\mathcal Mf(s)=\int_0^\infty x^{s-1}f(x)\,\mathrm dx$，其收敛区域通常是一个竖直带。换元 $x=e^u$ 后，它成为 $u$ 变量上的双边 Laplace 变换；取 $s=\sigma+i\tau$ 时，又与 $e^{\sigma u}f(e^u)$ 的 Fourier 变换相连。

它把乘法缩放变成幂因子：$\mathcal M[f(ax)](s)=a^{-s}\mathcal Mf(s)$。因此涉及幂权、比例或 $x\mapsto1/x$ 对称的积分适合这一语言。$e^{-x}$ 的 Mellin 变换正是 $\Gamma(s)$；$1/(1+x)$ 在 $0<\operatorname{Re}s<1$ 中的变换为 $\pi/\sin\pi s$。

备考中未必需要完整 Mellin 反演理论。先识别可用的参数积分、写清竖直带及端点可积性，再用 Gamma--Beta 恒等式，往往就足以完成所需计算；这一视角还能预示为什么留数与幂次渐近会自然相遇。



\section{常用特殊函数与渐近联系}

特殊函数把反复出现的积分、微分方程或乘积封装成统一对象。使用它们的价值在于继承递推、变换与渐近公式，而非给一个未计算的积分换名字。本章只选择与竞赛常见结构直接相连的函数，并说明公式的定义域。

\subsection{Gamma 函数与阶乘}

\label{app:special-gamma}

对 $\operatorname{Re}z>0$，定义
\[
\Gamma(z)=\int_0^\infty t^{z-1}e^{-t}\,\mathrm dt,
\qquad t^{z-1}=e^{(z-1)\log t}.
\]
零端由 $t^{\operatorname{Re}z-1}$、无穷端由指数衰减保证可积；在右半平面的紧集上，加入任意固定次 $|\log t|$ 后仍可支配，所以可对 $z$ 求导。分部积分给出 $\Gamma(z+1)=z\Gamma(z)$，且 $\Gamma(1)=1$，因此 $\Gamma(n+1)=n!$。

由 Gaussian 二重积分转极坐标，$\int_{\mathbb R}e^{-x^2}\,\mathrm dx=\sqrt\pi$，换元 $t=x^2$ 得 $\Gamma(1/2)=\sqrt\pi$。递推便计算全部正半整数值。利用递推还可将 Gamma 亚纯延拓到复平面，极点为 $0,-1,-2,\ldots$，且留数为 $(-1)^n/n!$。

\textbf{应用。} 对 $a,p>0$、$q>-1$，换元 $t=ax^p$ 得
\[
\int_0^\infty x^qe^{-ax^p}\,\mathrm dx
=\frac1p a^{-(q+1)/p}\Gamma\!\left(\frac{q+1}{p}\right).
\]
先检查 $q>-1$ 这一零端条件，再代公式；无穷端由指数统一控制。

\subsection{Beta 函数与比例变量}

\label{app:special-beta}

对 $a,b>0$，定义 $B(a,b)=\int_0^1t^{a-1}(1-t)^{b-1}\,\mathrm dt$。在 $\Gamma(a)\Gamma(b)$ 的正二重积分中令 $x=ru,y=r(1-u)$，Jacobian 为 $r$，由 Tonelli 得
\[
B(a,b)=\frac{\Gamma(a)\Gamma(b)}{\Gamma(a+b)}.
\]
总量 $r=x+y$ 与比例 $u=x/(x+y)$ 的分离，也解释了概率中的一个结构：两个独立且尺度参数相同的 Gamma 随机变量，其和与比例相互独立；和仍服从 Gamma 分布，比例服从 Beta 分布。

换元 $t=x/(1+x)$ 或 $t=\sin^2\theta$ 分别得到
\[
\int_0^\infty\frac{x^{a-1}}{(1+x)^{a+b}}\,\mathrm dx=B(a,b),
\qquad
\int_0^{\pi/2}\sin^{2a-1}\theta\cos^{2b-1}\theta\,\mathrm d\theta=\frac12B(a,b).
\]
这统一了许多三角幂积分与有理幂积分。参数微分可生成对数因子，但仍需在参数紧集上验证端点支配。

当 $0<a<1$ 时，钥匙孔围道对 $z^{a-1}/(1+z)$ 的两岸相位给出 $B(a,1-a)=\pi/\sin\pi a$，于是
\[
\Gamma(a)\Gamma(1-a)=\frac\pi{\sin\pi a}.
\]
这称为反射公式，可经解析延拓推广。围道证明的关键是小圆弧 $O(\varepsilon^a)$、大圆弧 $O(R^{a-1})$ 都消失，恰对应 $0<a<1$ 的可积范围。

\subsection{Stirling 公式及其可控修正}

\label{app:special-stirling}

对正整数 $n$，有
\[
n!=\sqrt{2\pi n}\left(\frac ne\right)^n
\exp\!\left(\frac1{12n}+O(n^{-3})\right).
\]
主尺度可由 $\Gamma(n+1)=\int_0^\infty t^ne^{-t}\,\mathrm dt$ 的 Laplace 方法得到：令 $t=n(1+u/\sqrt n)$，指数在最大点的二次项为 $-u^2/2$，Gaussian 积分给出 $\sqrt{2\pi n}$。远离最大点的两段需用指数衰减估计。

更精细的对数形式来自对 $\log k$ 应用 Euler--Maclaurin：
\[
\log n!=\left(n+\frac12\right)\log n-n+\frac12\log(2\pi)
+\frac1{12n}-\frac1{360n^3}+O(n^{-5}).
\]
边界修正给出 $n^{-1},n^{-3}$ 系数，可积高阶导数控制余项；常数可由上述 Gaussian 极限或 Wallis 乘积确定。此处陈述固定阶展开，不声称 Bernoulli 修正级数对固定 $n$ 收敛。

\textbf{应用。} 中心二项式系数满足 $\binom{2n}{n}\sim4^n/\sqrt{\pi n}$。固定 $0<p<1$，考虑 $\binom{n}{k}p^k(1-p)^{n-k}$，令 $k=np+O(\sqrt n)$ 并展开对数，二次项产生 Gaussian 形状；这与中心极限定理的局部机制一致，但完整概率极限还需控制远处尾部。

固定实数 $a,b$ 时，还有 $\Gamma(n+a)/\Gamma(n+b)=n^{a-b}(1+O(n^{-1}))$。它常用于识别乘积、矩与积分的幂次尺度；若 $a,b$ 也随 $n$ 变化，则要重新检查一致性。

\subsection{调和数、对数导数与 Zeta 函数}

\label{app:special-zeta}

记 $\psi(z)=\Gamma'(z)/\Gamma(z)$。对递推取对数求导，得到 $\psi(z+1)-\psi(z)=1/z$，故 $\psi(n+1)=\psi(1)+H_n$，其中 $\psi(1)=-\gamma$。这说明调和数是阶乘连续延拓的对数导数。

对 $\operatorname{Re}s>1$，Riemann Zeta 函数定义为 $\zeta(s)=\sum_{n\geq1}n^{-s}$，级数在紧子区域上一致绝对收敛。将 $1/(e^t-1)=\sum_{n\geq1}e^{-nt}$ 积分，可得
\[
\Gamma(s)\zeta(s)=\int_0^\infty\frac{t^{s-1}}{e^t-1}\,\mathrm dt.
\]
实数 $s>1$ 时先用 Tonelli，复数情形再用绝对可积性。零端行为约为 $t^{s-2}$，所以同样要求 $\operatorname{Re}s>1$。

\textbf{应用。} 对 $m>0$，$\int_0^\infty t^m/(e^t-1)\,\mathrm dt=\Gamma(m+1)\zeta(m+1)$。特殊函数把参数依赖压缩成一个可递推、可展开的表达式。Zeta 的解析延拓及零点如何控制素数，见素数分布，其中深层结论会另行标明。

\subsection{Bessel 函数与对称微分方程}

\label{app:special-bessel}

二维极坐标中的 Laplacian 为 $\partial_r^2+r^{-1}\partial_r+r^{-2}\partial_\theta^2$。对 Helmholtz 方程 $\Delta u+u=0$ 分离变量，角频率 $m$ 对应径向方程
\[
r^2R''+rR'+(r^2-m^2)R=0.
\]
对非负整数 $m$，在原点正则的解可由幂级数递推定义为
\[
J_m(r)=\sum_{k=0}^\infty\frac{(-1)^k}{k!(k+m)!}\left(\frac r2\right)^{2k+m}.
\]
比值判别保证对所有复数 $r$ 收敛，逐项代入即可验证方程。它是旋转对称问题中正弦、余弦的径向对应物。

特别地，$J_0(r)=\pi^{-1}\int_0^\pi\cos(r\cos\theta)\,\mathrm d\theta$，可通过一致展开余弦并计算偶次三角积分验证。因而圆周上的平面波平均成为 $J_0$。备考中掌握定义、方程与这一积分表示已经足以识别常见情形；更完整的零点、正交性和大参数展开可查特殊函数教材，不必在单道计算题里从头重建。



\section{微分方程、离散迭代与数值方法}

微分方程题常在研究三种对象：轨道怎样变化，两条轨道怎样比较，以及离散近似怎样传播误差。显式求解只是其中一种工具。积分因子、能量、返回映射和不变量把这些问题组织到可控制的量上。

\subsection{积分因子与 Gronwall 估计}

\label{app:dynamics-factor}

对 $y'+a(t)y=f(t)$，令 $A(t)=\int_0^ta(s)\,\mathrm ds$。乘以 $e^{A(t)}$ 后左端成为一个导数，因此
\[
y(t)=e^{-A(t)}y(0)+\int_0^te^{-(A(t)-A(s))}f(s)\,\mathrm ds.
\]
第一项传递初值，第二项累计外力。若 $a>0$、$A(t)\to\infty$ 且 $f(t)/a(t)\to L$，对 $y-L$ 使用公式，将积分切成固定头部与小比例尾部，得到 $y(t)\to L$。头部随 $e^{-A(t)}$ 消失；尾部利用 $|f-La|\leq\varepsilon a$ 后，其核积分不超过 $1$。

Gronwall 不等式是同一机制的估计版。若 $u\geq0$ 连续且 $u(t)\leq c+\int_0^tb(s)u(s)\,\mathrm ds$，其中 $b\geq0$ 可积、$c\geq0$，则 $u(t)\leq c\exp(\int_0^tb)$。令右端为 $v(t)$，则 $v'\leq bv$，积分因子即给结论。它用于初值连续依赖、唯一性和数值误差传播。

\subsection{流映射、周期解与不动点}

\label{app:dynamics-flow}

若 $y'=f(t,y)$ 的右端连续且对 $y$ 局部 Lipschitz（常数对 $t$ 局部一致），局部解存在且唯一；在有限时间内轨道留在定义域的紧子集时可继续延拓。这些是标准存在唯一性与延拓定理。使用“流映射”前要确认唯一性，否则一个初值可能对应多条轨道。

若 $f(t+T,y)=f(t,y)$，把初值 $r$ 经过时间 $T$ 的值记为 $P(r)$，称返回映射。周期解恰对应 $P(r)=r$。若边界方向保证区间 $[a,b]$ 正向不变，连续依赖给出 $P:[a,b]\to[a,b]$ 连续，介值定理便给不动点。唯一性可另由两解差的严格收缩证明，不必强求存在统一的压缩常数。

\textbf{应用。} 对 $y'=\sin t-y^3$，区间 $[-2,2]$ 的两端都指向内部，所以有周期解。两解 $y,z$ 的差满足 $w'=-(y^2+yz+z^2)w$。若初值不同，系数在一周期内的积分严格为正，因此 $|w(T)|<|w(0)|$，两个周期解不可能并存。存在性用区间不变，唯一性用距离缩小，两个步骤的依据清楚分开。

\subsection{能量、Lyapunov 函数与 Riccati 坐标}

\label{app:dynamics-energy}

对 $x''+V'(x)=0$，乘以 $x'$ 得 $E=(x')^2/2+V(x)$ 守恒。若增加阻尼 $a(x)x'$、$a\geq0$，则 $E'=-a(x)(x')^2\leq0$。能量把二阶方程变成轨道可达区域与单调性的限制，往往比显式积分更有效。

更一般的 Lyapunov 函数是沿轨道不增加的标量函数。若能量有下界，$E'\leq0$ 保证能量有有限极限，却不自动保证状态趋于单点；还需紧性及对 $E'=0$ 集中不变轨道的分析。相关的 LaSalle 不变原理属于进一步理论，此处不将它当成未经说明的最后一步。

若 $x>0$ 满足 $x''=g(t)x$，令 $u=x'/x$，则 $u'=g-u^2$。这是 Riccati 方程，把增长率改作未知量。常数上下界 $0\leq a\leq g\leq b$ 提示以 $\sqrt a,\sqrt b$ 为比较尺度；是否能得到严格界，仍取决于边界条件和等号轨道。对正解做对数导数只在无零点区间合法。

\subsection{离散不动点与线性化}

\label{app:dynamics-iteration}

迭代 $x_{n+1}=F(x_n)$ 的极限若存在且 $F$ 连续，必须是不动点；反过来，不动点存在并不保证轨道收敛。应先找到不变区间，再用单调性、压缩性或能量控制轨道。

若 $F\in C^1$、$F(a)=a$ 且 $|F'(a)|<1$，在足够小邻域内有统一压缩常数 $q<1$，故 $|x_n-a|\leq q^n|x_0-a|$。若 $|F'(a)|>1$，则除恰好落到不动点的情形外，附近误差被放大。临界值 $|F'(a)|=1$ 没有统一结论，需要高阶项。

当 $u\downarrow0$ 时，若 $F(a+u)=a+u-cu^{1+\alpha}+o(u^{1+\alpha})$，其中 $c,\alpha>0$，且已证明正侧轨道趋向 $a$，则 $x_n-a\sim(\alpha cn)^{-1/\alpha}$，收敛从几何速度变成幂次速度，见\hyperref[app:estimate-recursion]{递推的坐标选择}。若 $F'(a)=-1$，先研究 $F\circ F$ 往往能消去交错符号。选择一步或两步映射应由线性项决定。

\subsection{爆破、等待解与倒数积分}

\label{app:dynamics-blowup}

对自治标量方程 $y'=g(y)$，在 $g>0$ 的区间，分离变量给出从 $y_0$ 到 $Y$ 所需时间 $\int_{y_0}^Y\mathrm du/g(u)$。若到无穷远的积分有限，解在有限时间爆破；若积分发散，不能经这条正向轨道在有限时间到达无穷。

例如 $y'=y^p$、$y(0)=y_0>0$、$p>1$ 的爆破时间为 $y_0^{1-p}/(p-1)$。对微分不等式 $y'\geq cy^p$，其中 $c>0$、$p>1$ 且 $y(0)=y_0>0$，有 $(y^{1-p})'\leq-c(p-1)$，所以正解存在时间不超过 $y_0^{1-p}/[c(p-1)]$；无需先解原方程。

零点处的机制相反。对 $y'=y^\alpha$、$0<\alpha<1$，因 $\int_0^\varepsilon u^{-\alpha}\,\mathrm du<\infty$，解可先等待任意时间 $T$，再沿 $y(t)=((1-\alpha)(t-T))^{1/(1-\alpha)}$ 离开零点。这给出同一零初值的多个解。缺乏 Lipschitz 条件只是提示，倒数积分有限才提供实际构造。

\subsection{Runge–Kutta 阶条件来自树形微分}

\label{app:dynamics-rk}

对自治系统 $y'=f(y)$，若 $f$ 在轨道邻域属于 $C^3$，精确解的一步展开为
\[
y(t+h)=y+hf+\frac{h^2}{2}f'f
+\frac{h^3}{6}\bigl(f''(f,f)+f'f'f\bigr)+O(h^4).
\]
这里 $f'f$ 是导数作用于向量，两个三阶项分别对应两个分支同时进入二阶导数，以及一条连续的导数链。它们是不同的独立结构，不能只在标量线性方程上检查其中一个。

设 RK 法 $K_i=f(y+h\sum_ja_{ij}K_j)$、$y_+=y+h\sum_ib_iK_i$，令 $c_i=\sum_ja_{ij}$。代入展开并匹配上式，三阶条件为
\[
\sum_ib_i=1,\qquad\sum_ib_ic_i=\frac12,
\qquad\sum_ib_ic_i^2=\frac13,
\qquad\sum_{i,j}b_ia_{ij}c_j=\frac16.
\]
前两项控制总量与一阶时间矩，第三项来自 $f''(f,f)$，第四项来自 $f'f'f$。非自治方程可加入 $t'=1$ 变成自治系统，因而自然得到内部节点一致性。更高阶条件由 rooted trees 系统组织；备考常只需低阶的直接展开。

阶 $p$ 意味着从精确初值出发的一步误差为 $O(h^{p+1})$。要得到固定总时间上的全局 $O(h^p)$，还需稳定的误差传播，通常借助 Lipschitz 条件和离散 Gronwall 估计。全局误差的阶数比一步误差低一阶，是因为约 $1/h$ 步的误差需要累积。

\subsection{零稳定与绝对稳定}

\label{app:dynamics-stability}

在测试方程 $y'=\lambda y$ 上，一步法变成 $y_{n+1}=R(h\lambda)y_n$。$|R(z)|<1$ 给出固定步长下的严格衰减；通常以 $|R(z)|\leq1$ 描述闭稳定域。显式 Euler 的 $R(z)=1+z$，负实轴上需 $-2<h\lambda<0$；隐式 Euler 的 $R(z)=1/(1-z)$ 在开左半平面严格衰减，且在闭左半平面满足 $|R(z)|\leq1$，所以是 A 稳定的。

线性多步法的零稳定性则研究步长趋零、总时间固定时初值误差是否一致受控。特征多项式的所有根须在闭单位圆内，单位圆上的根须为单根。这是根条件。它与固定步长、长时间的绝对稳定性不同：根 $-1$ 可以满足零稳定条件，却保留不衰减的交错模态。

一致性加零稳定性推出收敛，是 Dahlquist 等价定理，适用于标准的线性多步法及相应的正则初值问题，本章引用而不证明。竞赛中应先写清问的是哪种稳定性，再选取根条件或放大因子；两者不能仅凭名称相互替代。

\section{真题之间的联系：误差、边界与整体条件}\label{sec:supp-audit}
本章为本选编补写的专题讲解，沿用本书改编部分的 CC BY-SA 4.0 许可。以下命题、推广和反例均为方法说明，不另计为竞赛真题。它们从非数学类真题中的具体计算出发，说明哪些步骤能推广、哪些条件不能省去。阅读时宜先完成相应真题，再核对推广证明中的新假设。

\subsection{相减之前必须知道误差有多小}
设 $f(x)\sim g(x)$，其准确含义是 $f(x)=g(x)(1+\varepsilon(x))$，其中 $\varepsilon(x)\to0$。如果还要与 $h(x)$ 相减，则
\[
 f(x)-h(x)=g(x)-h(x)+g(x)\varepsilon(x).
\]
因此从 $f\sim g$ 推出 $f-h\sim g-h$，还须有 $g\varepsilon=o(g-h)$；仅有 $\varepsilon\to0$ 并不足够。这是等价无穷小通常能在乘除中替换、却不能无条件在加减中替换的原因。

例如，$\sin x\sim x$，但 $x-\sin x$ 的首项为 $x^3/6$，用 $x$ 替掉 $\sin x$ 会抹掉整个答案。例~\ref{nm-2020-1a} 的分母也是三阶量，所以必须把分子误差控制到 $o(x^3)$。相反，例~\ref{nm-2019-1a} 先使用对数恒等式把相减变成 $\ln(1+u)$，再应用基本极限，就不需要展开很多项。

\begin{proposition}[Taylor 展开的阶数匹配]\label{supp-taylor}
设 $g$ 在0附近具有 $m$ 阶导数，且在0具有 Peano 形式展开。若 $g(0)=g'(0)=\cdots=g^{(m-1)}(0)=0$、$g^{(m)}(0)\ne0$，则
\[
 g(x)=\frac{g^{(m)}(0)}{m!}x^m+o(x^m).
\]
若另一个函数 $h(x)=cx^m+o(x^m)$ 且 $c\ne0$，则 $g(x)/h(x)\to g^{(m)}(0)/(m!c)$。
\end{proposition}
\begin{proof}
Taylor 多项式中低于 $m$ 次的系数全部为零。分子分母同除以 $x^m$，余项均趋于零，分母趋于非零常数 $c$，商的极限法则即可使用。关键不是展开得越多越好，而是保留第一项未被消去的项，并验证余项相对于分母确实可以忽略。
\end{proof}

\begin{example}[方法例：两个相同主项的进一步相减]
求 $\displaystyle\lim_{x\to0}\frac{\tan x-\sin x-x^3/2}{x^5}$。
\end{example}
\begin{solve}
利用五阶展开
\[
 \tan x=x+\frac{x^3}{3}+\frac{2x^5}{15}+o(x^5),\qquad
 \sin x=x-\frac{x^3}{6}+\frac{x^5}{120}+o(x^5),
\]
三阶差正好为 $x^3/2$，被再次减去后，五阶系数为 $2/15-1/120=1/8$，所以极限为 $1/8$。本例的分母强制要求第五阶信息；仅知道 $\tan x-\sin x\sim x^3/2$ 仍不足以回答这个更精细的问题。
\end{solve}

数列中的同类陷阱更隐蔽。由 $a_n\sim cn$ 不能推出 $a_{n+1}-a_n\to c$。取 $a_n=cn+(-1)^n\sqrt n$，确有 $a_n/n\to c$，但相邻差中的振荡项无界。例~\ref{nm-final9-4} 之所以可以计算相邻根阶乘的差，是因为它还使用了准确的对数递推。主项估计与增量估计是两个不同层次的结论。

\subsection{变化参数下的中值定理}
例~\ref{nm-2020-1e} 的指数随 $x$ 变化，但在使用中值定理的那一刻，$x$ 已经固定。这个区分可以整理为一个常用原则。

\begin{proposition}[带参数差商]\label{supp-param-mvt}
设 $\Phi(t,p)$ 在 $(a,p_0)$ 的邻域内对第一变量可导，且 $\partial_t\Phi$ 在 $(a,p_0)$ 连续。若 $u(x),v(x)\to a$、$p(x)\to p_0$，并且 $u(x)\ne v(x)$，则
\[
 \frac{\Phi(u(x),p(x))-\Phi(v(x),p(x))}{u(x)-v(x)}
 \longrightarrow\partial_t\Phi(a,p_0).
\]
\end{proposition}
\begin{proof}
对每个足够接近极限点的 $x$，中值定理给出位于 $u(x)$ 与 $v(x)$ 之间的 $\xi_x$，使该差商等于 $\partial_t\Phi(\xi_x,p(x))$。由夹逼原理 $\xi_x\to a$，再用偏导数在两变量意义下的连续性即得结论。无需构造 $\xi_x$ 关于 $x$ 的连续选择，也无需让 $u,v,p$ 本身可导。
\end{proof}

当 $\Phi(t,p)=t^p$ 且 $a>0$ 时，偏导为 $pt^{p-1}$，在正底数区域连续，于是得到 $p_0a^{p_0-1}$。$a>0$ 不是装饰条件：如果底数趋于零，同时指数也变化，就可能出现依赖路径的行为。例如 $t^{p-1}$ 中的 $(p-1)\ln t$ 未必有统一极限。使用含对数或实指数的变换时，应先通过保号性确定共同的定义域。

\begin{example}[方法例：双参数指数差商]
设 $u(x),v(x)\to2$，$p(x)\to3$，且 $u(x)\ne v(x)$。求 $[u(x)^{p(x)}-v(x)^{p(x)}]/[u(x)-v(x)]$ 的极限。
\end{example}
\begin{solve}
由命题~\ref{supp-param-mvt}，极限为 $3\cdot2^2=12$。如果把指数换成 $p(x)+1$，相同论证给出 $4\cdot2^3=32$。两次计算都不需要知道三个函数接近极限的速度；连续偏导数已经提供了所需的局部一致控制。
\end{solve}

\subsection{从积分列到精确收敛阈值}
例~\ref{nm-2020-7} 只要求 $p\geq1$ 时 $\sum u_n/n^p$ 收敛。其真实阈值可以由一个更一般的积分渐近式得到，这也解释了为什么质量逐渐集中在原点附近。

\begin{theorem}[端点集中的幂尺度]\label{supp-boundary-scale}
设 $\alpha>0$，并令 $v_n=\int_0^1(1+t^\alpha)^{-n}\,dt$。则
\[
 n^{1/\alpha}v_n\longrightarrow
 C_\alpha:=\int_0^\infty e^{-s^\alpha}\,ds
 =\frac1\alpha\Gamma\left(\frac1\alpha\right)>0.
\]
这里 $\Gamma(q)=\int_0^\infty z^{q-1}e^{-z}\,dz$，$q>0$。
\end{theorem}
\begin{proof}
令 $s=n^{1/\alpha}t$，得到
\[
 n^{1/\alpha}v_n=\int_0^{n^{1/\alpha}}
 \left(1+\frac{s^\alpha}{n}\right)^{-n}\,ds.
\]
对固定 $s\geq0$，积分函数最终有定义，且趋于 $e^{-s^\alpha}$。为论证在不断扩大的区间上取极限，先把它在区间外补为零。对 $0\leq s\leq n^{1/\alpha}$，有 $0\leq s^\alpha/n\leq1$，而 $\ln(1+z)\geq z/2$ 对 $0\leq z\leq1$ 成立，故
\[
 0\leq\mathbf1_{[0,n^{1/\alpha}]}(s)
 \left(1+\frac{s^\alpha}{n}\right)^{-n}\leq e^{-s^\alpha/2}.
\]
右端在 $[0,\infty)$ 可积。因此控制收敛定理允许交换极限与积分，得到第一个等式。第二个等式来自 $z=s^\alpha$。若不用 Lebesgue 积分语言，也可先在 $[0,R]$ 上证明一致收敛，再用这个共同的尾积分界令 $R\to\infty$。
\end{proof}

对原题的 $\alpha=4$，有 $u_n\sim C_4n^{-1/4}$，从而对任意实数 $p$，
\[
 \sum_{n=1}^\infty\frac{u_n}{n^p}
 \text{ 收敛 }\quad\Longleftrightarrow\quad p+\frac14>1,
 \quad\Longleftrightarrow\quad p>\frac34.
\]
临界点 $p=3/4$ 的通项与正常数乘 $1/n$ 等价，故发散。这个推广没有改变原题答案，而是给出原题中 $p\geq1$ 条件的余量。

同理，$\sum(-1)^nv_n$ 对所有 $\alpha>0$ 都因单调趋零而收敛；它绝对收敛当且仅当 $1/\alpha>1$，即 $0<\alpha<1$。因此参数 $\alpha=1$ 是绝对收敛和条件收敛之间的分界。局部形状 $t^\alpha$ 决定有效积分宽度 $n^{-1/\alpha}$，而这个宽度又决定级数的幂次。

\subsection{积分不等式的等号、上确界与稳定性}
例~\ref{nm-2018-3} 可以推广为双侧有界正函数的统一估计。把积分区间看作总质量为1的概率空间，可以更清楚地看出条件的作用，但下面不需要概率论知识。

\begin{theorem}[正函数与倒数的积分乘积]\label{supp-kantorovich}
设 $f$ 在长度为1的区间上可积，$0<m\leq f\leq M$，记 $A=\int f$、$B=\int1/f$。则
\[
 1\leq AB\leq\frac{(M+m)^2}{4mM}.
\]
若 $m<M$ 且允许可测函数，上界可由在两个各占一半长度的集合上分别取值 $m,M$ 的函数取得；连续函数一般只能逼近这一上界。
\end{theorem}
\begin{proof}
下界来自 $1=(\int\sqrt f\,1/\sqrt f)^2\leq AB$。由 $(f-m)(f-M)\leq0$，除以 $f$ 并积分，得到
\[
 A+mMB\leq m+M.
\]
再用 $4mMAB\leq(A+mMB)^2$ 即得上界。取等要求 $f\in\{m,M\}$ 几乎处处，且 $A=mMB$。设取值 $m$ 的比例为 $\theta$，代入后得到 $(2\theta-1)(m-M)=0$，故 $\theta=1/2$。当要求连续且定义域为区间时，两值条件强制函数为常数，不能满足比例各半；但用越来越短的连接区间可以逼近两值函数。
\end{proof}

对长度为 $L>0$ 的区间，应先除以 $L$ 归一化。未归一化的结果是
\[
 L^2\leq\left(\int f\right)\left(\int\frac1f\right)
 \leq L^2\frac{(M+m)^2}{4mM}.
\]
漏掉 $L^2$ 是把概率型平均与普通积分混用造成的量纲错误。

例~\ref{nm-2018-6} 的对数不等式还能量化离等号有多远。设 $A=\int f$ 且 $m\leq f\leq M$。对 $\ln t$ 在 $A$ 处使用二阶 Taylor 公式，因二阶导数为 $-1/t^2$，逐点有
\[
 \ln f\leq\ln A+\frac{f-A}{A}-\frac{(f-A)^2}{2M^2}.
\]
积分时线性项消失，因此
\[
 \ln A-\int\ln f\geq\frac1{2M^2}\int(f-A)^2.
\]
右端衡量函数偏离常数的平方平均。由此不仅知道等号只能在常函数处成立，还知道只要平方偏差不小，对数不等式的差距就不会任意小。这里的统一上界 $M$ 用于把逐点 Taylor 余项化为统一常数。

\subsection{闭形式与全局势函数之间的障碍}
设 $\omega=P\,dx+Q\,dy$，其中 $P,Q\in C^1(D)$。若存在 $U$ 使 $dU=\omega$，则 $P_y=Q_x$。反过来，在单连通的平面区域上，这个偏导相等条件足以推出势函数存在；在有孔的区域上则还需要控制绕孔的闭路积分。这里的单连通是指任意闭曲线都可在区域内连续收缩为一点。

例~\ref{nm-2019-1d} 中的分母是正定二次型，仅在原点为零。令
\[
 \omega=\frac{y\,dx-x\,dy}{3x^2-2xy+3y^2}.
\]
它在穿孔平面上满足 $P_y=Q_x$，但单位圆逆时针积分为
\[
 \oint\omega=-\int_0^{2\pi}\frac{d\theta}{3-\sin2\theta}
 =-\frac\pi{\sqrt2}\ne0.
\]
为核验最后一个数值，可利用一般恒等式 $\int_0^{2\pi}d\theta/(a+b\sin\theta)=2\pi/\sqrt{a^2-b^2}$（$a>|b|$）。该恒等式由把积分区间切在半角正切的极点处、再令 $u=\tan(\theta/2)$ 配方积分得到；$2\theta$ 的代换只多覆盖一个周期，平均值不变。

因此问题的障碍不在局部求导，而在绕原点一周后反正切的分支发生变化。删去一条从原点出发的射线，可以选择连续的角度分支；不同坐标片上的势函数相差常数，沿完整一周却不能把这些常数拼成单值函数。例~\ref{nm-2018-2} 删除的是两条完整直线，余下四个扇形各自凸，因此每个分支都有全局势函数，但不同分支之间没有路径，不能通过端点积分确定彼此的常数。

\begin{example}[方法例：角度微分与半径微分]
在 $D=\mathbb R^2\setminus\{0\}$ 上比较
\[
 \omega_1=\frac{x\,dx+y\,dy}{x^2+y^2},\qquad
 \omega_2=\frac{-y\,dx+x\,dy}{x^2+y^2}.
\]
\end{example}
\begin{solve}
第一式等于 $d\bigl(\tfrac12\ln(x^2+y^2)\bigr)$，所以存在全局势函数。第二式在局部极坐标中等于 $d\theta$，单位圆积分为 $2\pi$，所以不存在单值全局势函数。两式都满足偏导相等条件；只有第二式受到角度绕转的影响。是否有孔并不自动否定任何势函数，而是说明局部条件不再足够，必须逐个微分形式检验周期积分。
\end{solve}

\subsection{约束极值应在切空间上判定}
例~\ref{nm-2019-1e} 可以由乘子法得到候选点，但乘子方程本身只是必要条件。为明确二阶判定的对象，设 $f,g\in C^2$，在点 $x_0$ 满足 $g(x_0)=0$、$\nabla g(x_0)\ne0$，且存在 $\lambda$ 使 $\nabla f(x_0)+\lambda\nabla g(x_0)=0$。定义 $L=f+\lambda g$，以及约束的切空间
\[
 T=\{v:\nabla g(x_0)\cdot v=0\}.
\]
若 $x_0$ 是约束局部极小点，则对任意 $v\in T$，有 $v^{\mathsf T}\nabla^2L(x_0)v\geq0$。证明可取约束面上一条光滑曲线 $\gamma$ 满足 $\gamma(0)=x_0,\gamma'(0)=v$。因为 $g\circ\gamma=0$，有 $f\circ\gamma=L\circ\gamma$，而 $\nabla L(x_0)=0$ 消掉二阶链式法则中的加速度项，故
\[
 (f\circ\gamma)''(0)=v^{\mathsf T}\nabla^2L(x_0)v.
\]
局部极小推出该值非负。隐函数定理保证每个切向量都可以用这样的局部曲线实现。

充分条件也只需在切空间上严格正定。用隐函数定理给约束面取局部坐标，受限函数的 Hessian 正是该二次型在坐标切向量上的表示；严格正定配合二阶 Taylor 余项即可推出严格局部极小。要求整个环境空间上的 Hessian 正定虽然足够，却比必要的条件强得多。

\begin{example}[方法例：环境 Hessian 不定而约束极小]
求 $f(x,y)=x^2-y^2$ 在约束 $y=0$ 下的局部极值，并解释 Hessian 的作用。
\end{example}
\begin{solve}
受限函数为 $f(x,0)=x^2$，原点是严格极小点。乘子方程给出原点处 $\lambda=0$，环境 Hessian 为 $\operatorname{diag}(2,-2)$，确实不定。但切空间只有向量 $(v,0)$，其上二次型等于 $2v^2>0$。负方向 $(0,v)$ 根本不允许沿约束移动，因此不能用它否定受限极小。这也解释了为什么不能把乘子变量与原变量一起当作普通极值问题求 Hessian 正定性。
\end{solve}

\subsection{尾和怎样改变幂次}
例~\ref{nm-final8-6} 中，相邻差是 $n^{-2}$ 量级，而到极限的误差是 $n^{-1}$ 量级。以下结论提供了统一解释。

\begin{proposition}[正项尾和渐近]\label{supp-tail}
设 $d_n\geq0$，且 $d_n\sim cn^{-p}$，其中 $c>0,p>1$。令 $R_n=\sum_{k=n}^\infty d_k$。则
\[
 R_n\sim\frac c{p-1}n^{1-p}.
\]
\end{proposition}
\begin{proof}
由极限定义，对任意 $\varepsilon\in(0,c)$，足够大的 $k$ 满足 $(c-\varepsilon)k^{-p}\leq d_k\leq(c+\varepsilon)k^{-p}$。积分比较给出
\[
 \int_n^\infty x^{-p}\,dx\leq\sum_{k=n}^\infty k^{-p}
 \leq n^{-p}+\int_n^\infty x^{-p}\,dx.
\]
两端同乘 $n^{p-1}$，趋于 $1/(p-1)$。把 $R_n$ 夹在这两个常数倍之间，再令 $\varepsilon\downarrow0$ 即得结论。
\end{proof}

由此可知 $\sum R_n$ 收敛当且仅当 $p>2$。也可不使用渐近式而利用非负项的换序：
\[
 \sum_{n=1}^\infty R_n
 =\sum_{n=1}^\infty\sum_{k=n}^\infty d_k
 =\sum_{k=1}^\infty k d_k,
\]
等式允许两边同时为无穷大。每个 $d_k$ 被计入前 $k$ 个尾和，因此多出一个权重 $k$。这与例~\ref{nm-2018-7} 把三角形求和区域换序的思路相同，但那里通过 $b_n$ 的差分抵消了累积权重。

\begin{example}[方法例：收敛数列的误差级数]
设 $a_n\downarrow L$，且 $a_n-a_{n+1}\sim3/n^{5/2}$。判断 $\sum(a_n-L)$ 是否收敛，并求 $n^{3/2}(a_n-L)$ 的极限。
\end{example}
\begin{solve}
由望远镜和，$a_n-L=\sum_{k=n}^\infty(a_k-a_{k+1})$。在命题~\ref{supp-tail} 中取 $c=3,p=5/2$，得到 $a_n-L\sim2n^{-3/2}$。所以该误差级数收敛，所求极限为2。若把差的幂次改为2，则误差只有 $n^{-1}$ 量级，级数会发散；不能把“原数列收敛”与“误差可求和”混为一谈。
\end{solve}

\subsection{零初值唯一性与变系数控制}
例~\ref{nm-2019-3} 的平方积分因子可以扩展到变系数情形。设 $f$ 可导，$a\geq0$ 连续，满足 $|f'(x)|\leq a(x)|f(x)|$，并在某个 $x_0$ 有 $f(x_0)=0$。在 $x\geq x_0$ 的区间上令
\[
 H(x)=\exp\left(-2\int_{x_0}^xa(t)\,dt\right)f(x)^2.
\]
链式法则给出 $H'\leq0$，故 $0\leq H(x)\leq H(x_0)=0$。向左的区间可以对反向变量 $s=x_0-x$ 使用同样论证，从而 $f$ 在包含 $x_0$ 的整个区间上恒为零。

同一思想能证明局部 Lipschitz 右端的一阶初值问题唯一。若 $u'=F(x,u)$、$v'=F(x,v)$，两解具有相同初值，且在考虑区域内 $|F(x,y)-F(x,z)|\leq A|y-z|$，则 $w=u-v$ 满足 $|w'|\leq A|w|$，故 $w\equiv0$。这里 Lipschitz 条件用于控制两个函数值之差，而仅有 $F$ 连续只能帮助存在性，不能保证唯一性。

\begin{example}[方法反例：连续右端仍可不唯一]
考察 $y'=2\sqrt{|y|}$、$y(0)=0$ 在 $x\geq0$ 上的解。
\end{example}
\begin{solve}
除了零解，对每个 $c\geq0$ 都可取
\[
 y_c(x)=\begin{cases}0,&0\leq x\leq c,\\(x-c)^2,&x\geq c.\end{cases}
\]
在连接点左右导数均为0，所以 $y_c$ 可导并满足方程。右端连续，但在 $y=0$ 附近 $2\sqrt{|y|}/|y|$ 无界，无法满足统一 Lipschitz 条件。分离变量时除以 $\sqrt y$ 会丢掉零解和等待一段时间再离开零点的解，因此解答必须另行核对被除去的情形。
\end{solve}

\subsection{利用边界恒等式选择散度乘子}
例~\ref{nm-final8-4} 的关键在于球面上 $r^2$ 恒定，并非必须先解偏微分方程。设 $B_R\subset\mathbb R^d$ 是半径 $R$ 的球，$f\in C^2(\overline{B_R})$，$h=\Delta f$。在边界上，向量场 $R^2\nabla f$ 和 $r^2\nabla f$ 的通量相等，散度定理给出
\[
 2\int_{B_R}x\cdot\nabla f\,dx
 =\int_{B_R}(R^2-r^2)h(x)\,dx.
\]
右端只依赖 Laplace 算子作用后的函数 $h$，与 $f$ 中可能加入的调和部分无关。这也说明原题为什么能在没有边界值的情况下求得特定积分：被询问的线性泛函恰好消掉了那部分自由度。

若 $h=r^m$ 且 $m\geq0$，记单位球面的面积为 $\sigma_{d-1}$，则径向积分得到
\[
 \int_{B_R}x\cdot\nabla f\,dx
 =\frac{\sigma_{d-1}}2\int_0^R(R^2-r^2)r^{m+d-1}\,dr
 =\frac{\sigma_{d-1}R^{m+d+2}}{(m+d)(m+d+2)}.
\]
对原题 $d=3,m=2,R=1$，$\sigma_2=4\pi$，恢复 $4\pi/35$。作为独立复核，可取特解 $f=r^4/20$，因为在三维中 $\Delta r^4=20r^2$；此时 $x\cdot\nabla f=r^4/5$，直接球坐标积分也给出相同结果。

选择辅助向量场时可以按结果反推：若目标含 $x\cdot\nabla f$，就考虑 $\operatorname{div}(\phi\nabla f)=\nabla\phi\cdot\nabla f+\phi\Delta f$，再取 $\phi=r^2$ 使 $\nabla\phi=2x$。边界值是否简单决定了这一步能否消掉未知边界项。对于一般区域，$r^2$ 不在边界上恒定，就不能照搬球域的结论。

\section{用反例核查极限交换与极值论证}\label{sec:supp-counterexamples}
竞赛解答中的错误往往发生在一行看似自然的变形中：把逐点极限放进积分号、把偏导存在当成可微、把驻点当成最值点，或者在一个条件收敛表达式中改变运算次序。下面的例子不属于真题，而是用来检验这些步骤的适用范围。它们与前面的证明工具配合使用，帮助判断某个条件究竟承担了什么工作。

\subsection{积分的质量可能集中到越来越窄的区域}
在 $[0,1]$ 上定义连续函数
\[
 f_n(x)=\begin{cases}
 n^2x,&0\leq x\leq1/n,\\
 n(2-nx),&1/n\leq x\leq2/n,\\
 0,&2/n\leq x\leq1,
 \end{cases}\qquad n\geq2.
\]
对 $x=0$，$f_n(0)=0$；对任意固定 $x>0$，当 $n>2/x$ 时 $f_n(x)=0$。所以 $f_n$ 逐点收敛于零函数。但是图形由底为 $2/n$、高为 $n$ 的三角形组成，或直接逐段积分可得
\[
 \int_0^1f_n(x)\,dx
 =\int_0^{1/n}n^2x\,dx+
   \int_{1/n}^{2/n}n(2-nx)\,dx=1.
\]
因此 $\lim_n\int f_n=1\ne0=\int\lim_nf_n$。问题不是函数不连续，因为每一个 $f_n$ 都连续；而是峰值无界，质量被压到越来越短的区间内。

一致收敛足以排除这种现象：若 $\sup_{[0,1]}|f_n-f|\to0$，则 $|\int f_n-\int f|\leq\sup|f_n-f|\to0$。控制收敛提供另一种排除方式：若存在与 $n$ 无关的可积函数 $g$ 使 $|f_n|\leq g$，则在逐点几乎处处收敛的条件下也能交换极限与积分。本例不存在这样的可积控制函数，否则会与已算出的积分极限矛盾。

定理~\ref{supp-boundary-scale} 中的缩放积分同样发生端点集中，但在那里重新选取变量以后，构造出了统一的可积界 $e^{-s^\alpha/2}$。这说明集中现象本身不妨碍计算，关键在于是否把它放到合适的尺度上并控制尾部。

\subsection{函数列的一致极限并不自动允许逐项求导}
令 $g_n(x)=\sin(nx)/n$，定义在整个实轴上。因为 $|g_n(x)|\leq1/n$，它一致收敛于零函数。但 $g_n'(x)=\cos(nx)$，在 $x=0$ 处恒为1，而极限函数的导数为0。因此即使函数列一致收敛，也不能直接推出导数列收敛到极限函数的导数。

一个可用的充分条件如下：设 $g_n\in C^1[a,b]$，某点 $x_0$ 上 $g_n(x_0)$ 收敛，且 $g_n'$ 一致收敛于连续函数 $h$。由微积分基本定理，
\[
 g_n(x)=g_n(x_0)+\int_{x_0}^xg_n'(t)\,dt.
\]
积分中的一致收敛给出极限函数
\[
 g(x)=c+\int_{x_0}^xh(t)\,dt,
 \qquad c=\lim_ng_n(x_0),
\]
且收敛在整个 $[a,b]$ 一致。因为 $h$ 连续，再次使用微积分基本定理可得 $g'=h$。注意假设的强条件作用于导数列，单点函数值用于确定积分常数；只控制导数而没有某点函数值，仍可能加入发散的常数项。

\begin{example}[方法例：为什么几何级数可以求导]
在 $|x|<1$ 内，从 $\sum_{n=0}^\infty x^n=1/(1-x)$ 推出 $\sum_{n=1}^\infty nx^{n-1}=1/(1-x)^2$。
\end{example}
\begin{solve}
固定 $0<r<1$，在 $[-r,r]$ 上有 $|nx^{n-1}|\leq nr^{n-1}$，而数项级数 $\sum nr^{n-1}$ 收敛，例如由比值判别法可知。Weierstrass 判别法保证导数级数在该闭区间一致收敛；原部分和在0处恒为1。于是上述充分条件允许逐项求导。因为每个 $|x|<1$ 都被某个这样的 $[-r,r]$ 包含，结论在整个开区间成立。不能把 $r=1$ 直接代入一致估计，也不能据此宣称端点处的导数级数收敛。
\end{solve}

\subsection{所有直线方向都良好仍不保证二维极限存在}
定义
\[
 F(x,y)=\begin{cases}\dfrac{x^2y}{x^4+y^2},&(x,y)\ne(0,0),\\0,&(x,y)=(0,0).\end{cases}
\]
沿任意非竖直过原点直线 $y=kx$，当 $k\ne0$ 时 $F(x,kx)=kx/(x^2+k^2)\to0$；当 $k=0$ 时函数恒为零。沿竖直线 $x=0$ 也恒为零。然而沿抛物线 $y=x^2$，对 $x\ne0$ 有 $F(x,x^2)=1/2$。因此二维极限不存在，函数在原点不连续。

这个例子说明检查有限条或全部直线路径，仍只能提供必要条件。证明多元极限存在应寻找对所有接近方式同时成立的估计，例如 $|F(x,y)-L|\leq\omega(\sqrt{x^2+y^2})$ 且 $\omega(r)\to0$。反证极限不存在时则只需找到两条路径给出不同极限，量词的方向完全不同。

偏导数也不足以保证可微。设 $G(x,y)=xy/\sqrt{x^2+y^2}$，并令 $G(0,0)=0$。沿坐标轴函数为零，故 $G_x(0,0)=G_y(0,0)=0$。若它在原点可微，线性主项只能为零，应有 $G=o(r)$。但沿 $x=y=t>0$，$G(t,t)=t/\sqrt2$、$r=\sqrt2t$，故 $G/r=1/2$，矛盾。验证可微性的对象是函数与线性近似之差相对于距离的比值，不能只分别检查两个偏导。

\subsection{驻点、边界和极值是否取得}
可导函数的内点局部极值满足导数为零，但这一必要条件既不保证局部极值，也不保证整体最值。函数 $x^3$ 在0处导数为零，却在任何邻域内同时取正值和负值；函数 $x$ 在开区间 $(0,1)$ 上没有驻点，仍有上确界1和下确界0，但两者都不取得。

在紧集上的连续函数一定取得最大值和最小值。实际解题时，可以先使用这一存在性结论，再把候选位置分为内点和边界；内点用梯度方程，边界用低一维约束或参数化分析。若定义域不是紧集，还须检查无穷远行为、缺失端点或奇异点附近的极限，不能仅凭求出的有限个驻点就宣布全局最值。

\begin{example}[方法例：圆盘上的极值候选应如何列出]
求 $f(x,y)=x^2+2y^2-x$ 在闭单位圆盘 $x^2+y^2\leq1$ 上的最小值和最大值。
\end{example}
\begin{solve}
内点驻点方程为 $2x-1=0,4y=0$，得到 $(1/2,0)$，函数值 $-1/4$。配方 $f=(x-1/2)^2+2y^2-1/4$ 直接证明它是全局最小值。边界上 $y^2=1-x^2$，故 $f=2-x^2-x=9/4-(x+1/2)^2$，其中 $-1\leq x\leq1$。边界最大值为 $9/4$，在 $x=-1/2,y=\pm\sqrt3/2$ 取得。由于内点唯一驻点的值较小，整体最大值即为 $9/4$。存在性、内点候选、边界比较各自承担不同任务，不能省去其中任一环节。
\end{solve}

\subsection{反常积分的两端不能用形式相消代替定义}
普通意义下，$\int_{-1}^1dx/x$ 需要分别考察0两侧的反常积分。两侧分别发散，故该积分不存在。对称截断却给出
\[
 \lim_{\varepsilon\downarrow0}
 \left(\int_{-1}^{-\varepsilon}\frac{dx}{x}
       +\int_{\varepsilon}^1\frac{dx}{x}\right)=0.
\]
这叫 Cauchy 主值，与普通反常积分是不同的定义。若让右侧截断为 $c\varepsilon$，$c>0$，同样的截断和变成 $-\ln c$，显示相消依赖两侧趋近速度。题目没有指定主值时，不能擅自把这种对称极限当成原积分。

绝对可积提供了更稳定的交换条件。若二重积分满足 $\iint|h(x,y)|\,dx\,dy<\infty$，Fubini 定理允许换序并保证积分值一致；若 $h\geq0$，Tonelli 定理允许换序，但结果可能为无穷大。对于变号且非绝对可积的函数，应按照题目给定的迭代积分或截断方式计算。例~\ref{nm-2020-1d} 先求内层尾积分，再对外层有限上限积分，正是为了在原定义内完成运算，而不是依赖未经验证的换序。


\part{多元微积分专题讲义}

本编选取本项目已有数学分析课程材料中的多元微分、重积分、曲线曲面积分及路径无关性内容，作为前述竞赛真题的前置方法讲义。它们包括定义、定理证明和方法例题，均不计作全国大学生数学竞赛真题。内容由项目原有 \texttt{math-14.tex}、\texttt{math-15.tex} 与 \texttt{math-16.tex} 的相关节选重组而来；本书重建编号与交叉引用，原文件保持不变。这里的项目课程材料不归属于 Jacobi Lie 的题解作品，其来源与权利状态沿用项目原稿。

多元函数的核心困难在于方向、维度和区域同时变化。先从可微性与链式法则理解局部线性近似，再用隐函数定理与极值条件处理约束；积分部分先把区域写清楚，再使用换序或换元。曲线曲面积分还须区分第一型与第二型，后者的方向决定符号。各章方法可对应第一届初赛的圆盘积分、第十五届初赛的约束极值，以及第七、八、九届决赛的重积分和曲面积分题。

\section{偏导数与全微分}\label{course-14-1}
本节用于核对多元极限、隐函数、驻点与约束极值题中的可微性和非退化条件。仅有偏导数存在不能替代全微分存在；用必要条件列出候选点以后，还要比较边界并验证极值是否取得。

对于一个多元函数 $u=f(x_1,x_2,\cdots,x_n)$, 用一元微分学的工具来研究它是一个很自然的问题, 因此引进偏导数以及方向导数等概念是水到渠成的事情. 在本节中, 我们主要引进这些概念, 并讨论它们的简单性质。

\subsection{偏导数}

\begin{definition}\label{course-definition-14-1-1}\quad 设函数 $u=f(x)=f(x_1,x_2,\cdots,x_n)$ 在区域 $D\subset\mathbb R^n$ 上有定义, $x_0=(x_1^0,x_2^0,\cdots,x_n^0)\in D$. 对于 $1\leqslant i\leqslant n$, 若一元函数
\[
  f(x_1^0,\cdots,x_{i-1}^0,x_i,x_{i+1}^0,\cdots,x_n^0)
\]
在 $x_i^0$ 处的导数, 即
\[
  \lim_{x_i\to x_i^0}
  \frac{f(x_1^0,\cdots,x_{i-1}^0,x_i,x_{i+1}^0,\cdots,x_n^0)
  -f(x_1^0,x_2^0,\cdots,x_n^0)}{x_i-x_i^0}
\]
存在, 则称 $f(x)$ 在 $x_0$ 处关于 $x_i$ 可偏导, 并称上述极限为 $f(x)$ 在 $x_0$ 处关于 $x_i$ 的偏导数, 记为 $\dfrac{\partial f(x_0)}{\partial x_i}, f'_{x_i}(x_0), \left.\dfrac{\partial u}{\partial x_i}\right|_{x_0}$ 等。

容易看出, 当一个 $n$ 元函数 $f(x)$ 关于每个分量都可偏导时, $f(x)$ 具有 $n$ 个偏导数. 若 $f(x)$ 在 $D$ 内的每一点关于 $x_i$ 都可偏导, 则将其偏导数记为 $\dfrac{\partial f(x)}{\partial x_i}, f'_{x_i}(x), \dfrac{\partial u}{\partial x_i}$ 等. 此时, 对每个固定的分量 $x_i$, $\dfrac{\partial f(x)}{\partial x_i}$ 仍然是 $D$ 中的一个 $n$ 元函数。

特别地, 若 $z=f(x,y)$ 是一个二元函数, 则我们用 $\dfrac{\partial f(x,y)}{\partial x}$ (或 $f'_x(x,y), \dfrac{\partial z}{\partial x}$) 与 $\dfrac{\partial f(x,y)}{\partial y}$ (或 $f'_y(x,y), \dfrac{\partial z}{\partial y}$) 来记它的两个偏导数. 同理, 对三元函数 $u=f(x,y,z)$, 可以类似地给出相应记号: $\dfrac{\partial f(x,y,z)}{\partial x}, \dfrac{\partial f(x,y,z)}{\partial y}, \dfrac{\partial f(x,y,z)}{\partial z}$, 或 $f'_x(x,y,z), f'_y(x,y,z), f'_z(x,y,z)$, 或 $\dfrac{\partial u}{\partial x}, \dfrac{\partial u}{\partial y}, \dfrac{\partial u}{\partial z}$。

对一个二元函数 $z=f(x,y)((x,y)\in D)$, 我们可以讨论其在 $(x_0,y_0)\in D$ 处偏导数的几何意义. 当函数具有较好的性质时, 该函数的图像是 $\mathbb R^3$ 中的一块曲面. $y=y_0$ 是一个过点 $(x_0,y_0,0)$ 且平行于 $Ozx$ 平面的平面, 它与曲面 $z=f(x,y)((x,y)\in D)$ 的交线为
\[
  L:\begin{cases}
    x=x,\\
    y=y_0,\\
    z=f(x,y_0).
  \end{cases}
\]
因此, $\dfrac{\partial f(x_0,y_0)}{\partial x}$ 是平面 $y=y_0$ 内的曲线 $L$ 在 $x=x_0$ 处的切线斜率 $k_1$. 同样, $\dfrac{\partial f(x_0,y_0)}{\partial y}$ 是平面 $x=x_0$ 与曲面 $z=f(x,y)((x,y)\in D)$ 的交线在 $y=y_0$ 处的切线斜率 $k_2$ 。

读者应该注意的是, 函数 $z=f(x,y)$ 在 $(x_0,y_0)$ 处的两个偏导数存在与否只与曲面与 $x=x_0$ 及 $y=y_0$ 的交线性质有关. 因此, 一个二元函数在点 $(x_0,y_0)$ 处即使两个偏导数都存在, 该函数在点 $(x_0,y_0)$ 处也未必连续. 例如, 设函数
\[
  f(x,y)=
  \begin{cases}
    0, & xy=0,\\
    1, & xy\ne0,
  \end{cases}
\]
容易看出 $f'_x(0,0)=f'_y(0,0)=0$, 但 $f(x,y)$ 在 $(0,0)$ 处不连续。

下面我们举两个例子:

\end{definition}
\begin{example}\label{course-example-14-1-1}\quad 设函数 $f(x,y)=\tan\dfrac{x^2}{y}$, 求 $\dfrac{\partial f(0,1)}{\partial x}$ 及 $\dfrac{\partial f(0,1)}{\partial y}$。

\end{example}
\begin{solve}\quad 由偏导数的定义有
\[
  \frac{\partial f(0,1)}{\partial x}
  =(\tan x^2)'\big|_{x=0}
  =2x\sec^2 x^2\big|_{x=0}=0,
\]
\[
  \frac{\partial f(0,1)}{\partial y}
  =(0)'\big|_{y=1}=0.
\]

\end{solve}
\begin{example}\label{course-example-14-1-2}\quad 设函数 $f(x,y)=\arcsin\dfrac{x}{\sqrt{x^2+y^2}}$, 求 $f'_x(x,y)$ 及 $f'_y(x,y)$。

\end{example}
\begin{solve}\quad 对这种计算题, 我们在求 $f'_x(x,y)$ 时将 $y$ 看成常数即可. 因此, 我们有
\[
\begin{aligned}
  f'_x(x,y)
  &=
  \frac{1}{\sqrt{1-\dfrac{x^2}{x^2+y^2}}}\cdot
  \frac{1}{x^2+y^2}
  \left(\sqrt{x^2+y^2}-\frac{x^2}{\sqrt{x^2+y^2}}\right)\\
  &=\frac{1}{\sqrt{y^2}}\cdot\frac{y^2}{x^2+y^2}
  =\frac{y\operatorname{sgn} y}{x^2+y^2};
\end{aligned}
\]
同理,
\[
  f'_y(x,y)
  =
  \frac{1}{\sqrt{1-\dfrac{x^2}{x^2+y^2}}}
  \left(-\frac{xy}{(x^2+y^2)^{3/2}}\right)
  =\frac{-x\operatorname{sgn} y}{x^2+y^2},
\]
其中 $\operatorname{sgn} y$ 是符号函数。

\end{solve}
\subsection{方向导数}

一个多元函数的方向导数是用来刻画定义域内从一点出发的射线上函数的变化情况的. 从一点出发的一条射线可以看成是一个方向, 而一个方向可由以原点为心的单位球面上的点所确定. 设 $n$ 维单位球面为 $x_1^2+x_2^2+\cdots+x_n^2=1$, 则每一个以原点为起点, 终点在该球面上的有向线段表示一个单位向量 $v$, 它可以记为 $v=(\cos\theta_1,\cos\theta_2,\cdots,\cos\theta_n)$, 其中 $\theta_i(1\leqslant i\leqslant n)$ 是该向量与 $x_i$ 轴正向的夹角, $\cos\theta_i(1\leqslant i\leqslant n)$ 也称为是 $v$ 的方向余弦. 特别地, 对于 $\mathbb R^2$ 内任一个单位向量 $v$, 记它与正实轴的夹角为 $\theta$, 则有 $v=(\cos\theta,\sin\theta)$. 有了这些准备, 我们可以给出以下定义。

\begin{definition}\label{course-definition-14-1-2}\quad 设函数 $u=f(x)$ 在区域 $D\subset\mathbb R^n$ 上有定义, $x_0\in D$,
\[
  v=(\cos\theta_1,\cos\theta_2,\cdots,\cos\theta_n)
\]
为一方向. 如果极限
\[
  \lim_{t\to0+0}\frac{f(x_0+tv)-f(x_0)}{t}
\]
存在, 则称该极限为 $f(x)$ 在 $x_0$ 处沿方向 $v$ 的方向导数, 记为 $\dfrac{\partial f(x_0)}{\partial v}$ 或 $\left.\dfrac{\partial u}{\partial v}\right|_{x_0}$。

\end{definition}
\noindent\keyterm{注 1}\quad 在上述定义中, 若记 $x_0+tv=x$, 则
\[
  \frac{\partial f(x_0)}{\partial v}
  =\lim_{x\to x_0}\frac{f(x)-f(x_0)}{(x-x_0)v}.
\]
从上述定义可以看出, 一个函数在某点处的方向导数不依赖于坐标系的选取。

\noindent\keyterm{注 2}\quad 设 $f(x)$ 在 $x_0=(x_1^0,x_2^0,\cdots,x_n^0)$ 处关于 $x_i$ 的偏导数 $\dfrac{\partial f(x_0)}{\partial x_i}(1\leqslant i\leqslant n)$ 存在并等于 $A$, 若记 $v_i$ 是第 $i$ 个分量为 $1$ 的单位向量, 则 $\dfrac{\partial f(x_0)}{\partial v_i}=A$. 但读者应该注意的是 $\dfrac{\partial f(x_0)}{\partial(-v_i)}=-A$。

\begin{example}\label{course-example-14-1-3}\quad 设函数 $f(x,y)=\sqrt{x^2+y^2}$, 试求 $f(x,y)$ 在 $(0,0)$ 处各个方向的方向导数。

\end{example}
\begin{solve}\quad 设 $v=(\cos\theta,\sin\theta)(0\leqslant\theta<2\pi)$, 则
\[
\begin{aligned}
  \frac{\partial f(0,0)}{\partial v}
  &=\lim_{t\to0+0}
  \frac{f(t\cos\theta,t\sin\theta)-f(0,0)}{t}\\
  &=\lim_{t\to0+0}
  \frac{\sqrt{t^2\cos^2\theta+t^2\sin^2\theta}}{t}\\
  &=\lim_{t\to0+0}\frac{t}{t}=1.
\end{aligned}
\]

从注 2 我们可知, 在上例中, $f(x,y)$ 在 $(0,0)$ 处两个偏导数都不存在. 当然, 我们也可从偏导数的定义直接证明这一结果。

\end{solve}
\begin{example}\label{course-example-14-1-4}\quad 设函数
\[
  f(x,y)=
  \begin{cases}
    1, & y=x^2,\ \text{且 } x\ne0,\\
    0, & \text{其他},
  \end{cases}
\]
即 $f(x,y)$ 在抛物线 $y=x^2$ 上除去原点外取值为 $1$, 其余各处取值为 $0$. 证明: 对任意方向 $v$, 有 $\dfrac{\partial f(0,0)}{\partial v}=0$, 但 $f(x,y)$ 在 $(0,0)$ 处不连续。

\end{example}
\begin{proof}\quad 记 $\mathbf 0=(0,0)$, 对任意方向 $v$, 当 $t$ 很小时, $\mathbf 0+tv$ 与抛物线不交, 从而 $f(tv)=0=f(0,0)$, 因此 $\dfrac{\partial f(0,0)}{\partial v}=0$. 由于
\[
  \lim_{k\to\infty}f\left(\frac1k,\frac1{k^2}\right)=1\ne f(0,0),
\]
所以 $f(x,y)$ 在 $(0,0)$ 处不连续。

\end{proof}
\subsection{全微分}

一元函数 $y=f(x)$ 的微分是函数关于自变量的一阶线性的近似, 由其几何意义, 实际上是用切线来局部近似曲线 $y=f(x)$. 对于一个二元函数, 由于它的图像一般是一块曲面, 因此需用切平面来局部近似它, 而这些都与多元函数的全微分有关. 对一个 $n$ 元函数, 我们有以下定义。

\begin{definition}\label{course-definition-14-1-3}\quad 设函数 $f(x)=f(x_1,x_2,\cdots,x_n)$ 在区域 $D\subset\mathbb R^n$ 上有定义, 且 $x_0=(x_1^0,x_2^0,\cdots,x_n^0)\in D$. 记 $\Delta x=(\Delta x_1,\Delta x_2,\cdots,\Delta x_n)$, 并称它为自变量的全增量. 再设
\[
  x=x_0+\Delta x=(x_1+\Delta x_1,x_2+\Delta x_2,\cdots,x_n+\Delta x_n)\in D.
\]
若存在仅依赖于 $x_0$ 的常数 $A_i(i=1,2,\cdots,n)$, 使得
\[
  \Delta f(x_0)=f(x_0+\Delta x)-f(x_0)
  =\sum_{i=1}^n A_i\Delta x_i+o(|\Delta x|),\quad |\Delta x|\to0,
\]
则称 $f(x)$ 在 $x_0$ 处可微, 并称 $\sum\limits_{i=1}^n A_i\Delta x_i$ 为 $f(x)$ 在 $x_0$ 处的全微分, 记为 $\mathrm df(x_0)$, 即
\[
  \mathrm df(x_0)=\sum_{i=1}^n A_i\Delta x_i.
\]
若 $f(x)$ 在 $D$ 内每一点处均可微, 则称 $f(x)$ 在 $D$ 内是可微函数。

当 $x_i(1\leqslant i\leqslant n)$ 是自变量时, 定义 $\mathrm dx_i=\Delta x_i$. 因此 $f(x)$ 在 $x_0$ 处的全微分可以记为 $\mathrm df(x_0)=\sum\limits_{i=1}^n A_i\mathrm dx_i$。

从全微分的定义可立即推出, 当 $f(x)$ 在 $x_0$ 处可微时, 它的全微分是函数改变量的线性主要部分. 本书中涉及一个多元函数的微分时, 一般情形下总是指上述的全微分, 因此今后我们将全微分简称为微分。

\end{definition}
\begin{theorem}\label{course-theorem-14-1-1}\quad 设函数 $f(x)$ 在区域 $D\subset\mathbb R^n$ 上有定义, 在 $x_0=(x_1^0,\cdots,x_n^0)\in D$ 处可微, 记其微分为 $\mathrm df(x_0)=\sum\limits_{i=1}^n A_i\mathrm dx_i$, 则
\[
\begin{gathered}
  (1)\quad f(x)\text{ 在 }x_0\text{ 处连续};\\
  (2)\quad \text{对于 } i(1\leqslant i\leqslant n), f(x)\text{ 关于 }x_i\text{ 可偏导, 并且有 }
  \frac{\partial f(x_0)}{\partial x_i}=A_i.
\end{gathered}
\]

\end{theorem}
\begin{proof}\quad (1) 记 $\Delta x=x-x_0$, 则 $x\to x_0$ 等价于 $\Delta x\to0$, 即它等价于: 对于 $\forall i(1\leqslant i\leqslant n)$, 有 $\Delta x_i\to0$. 由此我们有
\[
  \lim_{\Delta x\to0}[f(x_0+\Delta x)-f(x_0)]
  =
  \lim_{\Delta x\to0}\left[\sum_{i=1}^n A_i\Delta x_i+o(|\Delta x|)\right]=0,
\]
即 $f(x)$ 在 $x_0$ 处连续。

(2) 对任何固定的 $i(1\leqslant i\leqslant n)$, 当 $j=1,2,\cdots,n$ 且 $j\ne i$ 时, 令 $x_j=x_j^0$, 即 $\Delta x_j=0$. 此时我们有 $\Delta x=(0,\cdots,\Delta x_i,\cdots,0)$ 和 $|\Delta x|=|\Delta x_i|$, 因此有
\[
\begin{aligned}
&\lim_{\Delta x_i\to0}
  \frac{f(x_1^0,x_2^0,\cdots,x_{i-1}^0,x_i^0+\Delta x_i,x_{i+1}^0,\cdots,x_n^0)-f(x_0)}{\Delta x_i}\\
&\quad =
  \lim_{\Delta x_i\to0}\left(A_i+\frac{o(|\Delta x_i|)}{\Delta x_i}\right)=A_i,
\end{aligned}
\]
从而 $f(x)$ 关于 $x_i$ 可偏导, 并且有 $\dfrac{\partial f(x_0)}{\partial x_i}=A_i$. 证毕。

上述定理说明, 当函数 $f(x)$ 在 $x_0$ 处可微时, 必有
\[
  \mathrm df(x_0)=\sum_{i=1}^n\frac{\partial f(x_0)}{\partial x_i}\mathrm dx_i.
\]
特别地, 当 $z=f(x,y)$ 在 $(x_0,y_0)$ 处可微时, 我们有
\[
  \mathrm df(x_0,y_0)=\frac{\partial f(x_0,y_0)}{\partial x}\mathrm dx+
  \frac{\partial f(x_0,y_0)}{\partial y}\mathrm dy;
\]
而当 $u=f(x,y,z)$ 在 $(x_0,y_0,z_0)$ 处可微时, 则有
\[
  \mathrm df(x_0,y_0,z_0)=
  \frac{\partial f(x_0,y_0,z_0)}{\partial x}\mathrm dx+
  \frac{\partial f(x_0,y_0,z_0)}{\partial y}\mathrm dy+
  \frac{\partial f(x_0,y_0,z_0)}{\partial z}\mathrm dz.
\]

我们已经有了例子表明存在可偏导但不连续的多元函数, 因此这样的函数不可微. 这说明, 可偏导与可微在多元函数的情形不是等价的. 但是我们有以下的定理。

\end{proof}
\begin{theorem}\label{course-theorem-14-1-2}\quad 设函数 $f(x)$ 在区域 $D\subset\mathbb R^n$ 上有定义, $x_0=(x_1^0,x_2^0,\cdots,x_n^0)\in D$, 再设 $f(x)$ 在 $x_0$ 的邻域 $U(x_0,\delta_0)(\delta_0>0)$ 内存在各个偏导数, 并且这些偏导数在 $x_0$ 处连续, 则 $f(x)$ 在 $x_0$ 处可微。

\end{theorem}
\begin{proof}\quad 我们对 $n$ 作归纳法来证明定理。

当 $n=1$ 时, 定理结论显然成立。

假设当 $n=k$ 时定理成立, 即
\[
\begin{aligned}
&f(x_1^0+\Delta x_1,x_2^0+\Delta x_2,\cdots,x_k^0+\Delta x_k)
  -f(x_1^0,x_2^0,\cdots,x_k^0)\\
&\quad =
  \sum_{i=1}^k
  \frac{\partial f(x_1^0,x_2^0,\cdots,x_k^0)}{\partial x_i}\Delta x_i
  +o(1)\left(\sqrt{\sum_{i=1}^k(\Delta x_i)^2}\right)
\end{aligned}
\]
\[
  \left(\sqrt{\sum_{i=1}^k(\Delta x_i)^2}\to0\right).
\]
当 $n=k+1$ 时,
\[
\begin{aligned}
&\Delta f(x_1^0,\cdots,x_k^0,x_{k+1}^0)\\
&=f(x_1^0+\Delta x_1,\cdots,x_k^0+\Delta x_k,x_{k+1}^0+\Delta x_{k+1})
  -f(x_1^0,\cdots,x_k^0,x_{k+1}^0)\\
&=\big[f(x_1^0+\Delta x_1,\cdots,x_k^0+\Delta x_k,x_{k+1}^0+\Delta x_{k+1})
  -f(x_1^0+\Delta x_1,\cdots,x_k^0+\Delta x_k,x_{k+1}^0)\big]\\
&\quad+\big[f(x_1^0+\Delta x_1,\cdots,x_k^0+\Delta x_k,x_{k+1}^0)
  -f(x_1^0,\cdots,x_k^0,x_{k+1}^0)\big].
\end{aligned}
\]
由一元函数的拉格朗日微分中值定理及 $\dfrac{\partial f(x)}{\partial x_{k+1}}$ 在 $x_0$ 处的连续性知
\begin{equation}\label{course-eq-14-1-1}
\begin{aligned}
&f(x_1^0+\Delta x_1,\cdots,x_k^0+\Delta x_k,x_{k+1}^0+\Delta x_{k+1})
  -f(x_1^0+\Delta x_1,\cdots,x_k^0+\Delta x_k,x_{k+1}^0)\\
&=f'_{x_{k+1}}(x_1^0+\Delta x_1,\cdots,x_k^0+\Delta x_k,
  x_{k+1}^0+\theta\Delta x_{k+1})\Delta x_{k+1}\\
&=f'_{x_{k+1}}(x_1^0,\cdots,x_k^0,x_{k+1}^0)\Delta x_{k+1}
  +o(1)(|\Delta x_{k+1}|),
\end{aligned}
\end{equation}
其中 $0<\theta<1$. 由归纳法假设有
\begin{equation}\label{course-eq-14-1-2}
\begin{aligned}
&f(x_1^0+\Delta x_1,\cdots,x_k^0+\Delta x_k,x_{k+1}^0)
  -f(x_1^0,\cdots,x_k^0,x_{k+1}^0)\\
&=
\sum_{i=1}^k
\frac{\partial f(x_1^0,\cdots,x_k^0,x_{k+1}^0)}{\partial x_i}\Delta x_i
+o(1)\left(\sqrt{\sum_{i=1}^k(\Delta x_i)^2}\right)
\end{aligned}
\end{equation}
\[
  \left(\sqrt{\sum_{i=1}^k(\Delta x_i)^2}\to0\right).
\]
注意到当 $\sqrt{\sum\limits_{i=1}^{k+1}(\Delta x_i)^2}\to0$ 时, 有
\[
  \frac{
  o(1)(|\Delta x_{k+1}|)
  +o(1)\left(\sqrt{\sum\limits_{i=1}^k(\Delta x_i)^2}\right)}
  {\sqrt{\sum\limits_{i=1}^{k+1}(\Delta x_i)^2}}\to0,
\]
因此由式~\eqref{course-eq-14-1-1}, \eqref{course-eq-14-1-2}  即得
\[
\begin{aligned}
&\Delta f(x_1^0,\cdots,x_k^0,x_{k+1}^0)\\
&=
  f(x_1^0+\Delta x_1,\cdots,x_k^0+\Delta x_k,x_{k+1}^0+\Delta x_{k+1})
  -f(x_1^0,\cdots,x_k^0,x_{k+1}^0)\\
&=\sum_{i=1}^{k+1}
  \frac{\partial f(x_1^0,x_2^0,\cdots,x_k^0,x_{k+1}^0)}{\partial x_i}\Delta x_i
  +o(1)\left(\sqrt{\sum_{i=1}^{k+1}(\Delta x_i)^2}\right),
\end{aligned}
\]
即 $f(x_1,x_2,\cdots,x_{k+1})$ 在 $(x_1^0,x_2^0,\cdots,x_{k+1}^0)$ 处可微, 且
\[
  \mathrm df(x_1^0,x_2^0,\cdots,x_{k+1}^0)
  =\sum_{i=1}^{k+1}
  \frac{\partial f(x_1^0,x_2^0,\cdots,x_{k+1}^0)}{\partial x_i}\mathrm dx_i.
\]
这就证明了定理的结论对一切 $n\in\mathbb N$ 成立. 证毕。

\end{proof}
\begin{remark}\quad 若函数 $f(x)$ 在区域 $D\subset\mathbb R^n$ 上关于自变量的各个分量都具有连续偏导数, 通常我们称 $f(x)$ 在 $D$ 内是 $C^1$ 的, 记为 $f(x)\in C^1(D)$, 此时我们也称 $f(x)$ 在 $D$ 内连续可微。

\end{remark}
\begin{example}\label{course-example-14-1-5}\quad 证明函数
\[
  f(x,y)=
  \begin{cases}
    (x^2+y^2)\sin\dfrac1{x^2+y^2}, & x^2+y^2\ne0,\\
    0, & x^2+y^2=0
  \end{cases}
\]
在 $(0,0)$ 处可微, 但 $f'_x(x,y)$ 及 $f'_y(x,y)$ 在 $(0,0)$ 处不连续。

\end{example}
\begin{proof}\quad 显然 $f'_x(0,0)=f'_y(0,0)=0$. 由于
\[
  \Delta f(0,0)=f(\Delta x,\Delta y)
  =[(\Delta x)^2+(\Delta y)^2]\sin\frac1{(\Delta x)^2+(\Delta y)^2}
  =o(1)\sqrt{(\Delta x)^2+(\Delta y)^2},
\]
\[
  \left(\sqrt{(\Delta x)^2+(\Delta y)^2}\to0\right),
\]
因此 $f(x,y)$ 在 $(0,0)$ 处可微. 当 $(x,y)\ne(0,0)$ 时, 有
\[
  f'_x(x,y)=2x\sin\frac1{x^2+y^2}-\frac{2x}{x^2+y^2}\cos\frac1{x^2+y^2}.
\]
由于 $(x_k,y_k)=\left(\dfrac1{\sqrt{2k\pi}},0\right)\to(0,0)$ 时, $f'_x(x_k,y_k)\to\infty$, 因此 $f'_x(x,y)$ 在 $(0,0)$ 处不连续。

由函数的对称性, 同理 $f'_y(x,y)$ 在 $(0,0)$ 处也不连续。

对于可微函数, 它的方向导数与偏导数密切相关. 对此我们有下面的结论。

\end{proof}
\begin{theorem}\label{course-theorem-14-1-3}\quad 设函数 $f(x)$ 在区域 $D\subset\mathbb R^n$ 上有定义, 且在 $x_0\in D$ 处可微, 则 $f(x)$ 在 $x_0$ 处沿方向 $v=(\cos\theta_1,\cos\theta_2,\cdots,\cos\theta_n)$ 的方向导数为
\[
  \frac{\partial f(x_0)}{\partial v}
  =\sum_{i=1}^n\frac{\partial f(x_0)}{\partial x_i}\cos\theta_i.
\]

\end{theorem}
\begin{proof}\quad 由 $f(x)$ 的可微性及方向导数的定义, 我们有
\[
\begin{aligned}
  \lim_{t\to0+0}\frac{f(x_0+tv)-f(x_0)}{t}
  &=\lim_{t\to0+0}
  \frac{\sum\limits_{i=1}^n\dfrac{\partial f(x_0)}{\partial x_i}t\cos\theta_i+o(|t|)}{t}\\
  &=\sum_{i=1}^n\frac{\partial f(x_0)}{\partial x_i}\cos\theta_i.
\end{aligned}
\]
证毕。

在本小节的最后, 我们再举一个例子来说明多元函数微分的在近似计算中的应用。

\end{proof}
\begin{example}\label{course-example-14-1-6}\quad 求 $2.99^2\times1.02^3$ 的近似值。

\end{example}
\begin{solve}\quad 取函数 $f(x,y)=x^2y^3$, 显然 $f(x,y)$ 的两个偏导数在 $\mathbb R^2$ 上连续, 从而它在 $\mathbb R^2$ 处处可微. 取 $(x_0,y_0)=(3,1), \Delta x=-0.01, \Delta y=0.02$, 则
\[
\begin{aligned}
  2.99^2\times1.02^3
  &=f(2.99,1.02)=f(3-0.01,1+0.02)\\
  &\approx f(3,1)+f'_x(3,1)\Delta x+f'_y(3,1)\Delta y\\
  &=3^2\times1^3+f'_x(3,1)\cdot(-0.01)+f'_y(3,1)\cdot(0.02)\\
  &=9+6\cdot(-0.01)+27\cdot(0.02)=9.48.
\end{aligned}
\]

\end{solve}
\subsection{梯度}

与方向导数密切相关的概念是梯度. 在本节我们先介绍梯度的概念与基本的性质, 在本书后面的章节中我们还会多次讨论它的进一步的性质。

在上一小节中, 我们知道当 $f(x)$ 在 $x_0\in\mathbb R^n$ 处可微时, $f(x)$ 沿任意方向 $v=(\cos\theta_1,\cos\theta_2,\cdots,\cos\theta_n)$ 的方向导数为
\[
  \frac{\partial f(x_0)}{\partial v}=
  \sum_{i=1}^n\frac{\partial f(x_0)}{\partial x_i}\cos\theta_i.
\]
当 $f(x)$ 在 $x_0$ 处的 $n$ 个偏导数不全为零时, $f(x)$ 沿方向
\[
  v_0=
  \frac1{\sqrt{\sum\limits_{i=1}^n\left(\dfrac{\partial f(x_0)}{\partial x_i}\right)^2}}
  \left(\frac{\partial f(x_0)}{\partial x_1},\frac{\partial f(x_0)}{\partial x_2},\cdots,\frac{\partial f(x_0)}{\partial x_n}\right)
\]
的方向导数达到最大值. 因此向量
\[
  \left(\frac{\partial f(x_0)}{\partial x_1},\frac{\partial f(x_0)}{\partial x_2},\cdots,\frac{\partial f(x_0)}{\partial x_n}\right)
\]
是 $f(x)$ 在 $x_0$ 处方向导数达到最大的方向, 同时它的模就是该方向的方向导数. 由此我们引进下面的定义。

\begin{definition}\label{course-definition-14-1-4}\quad 设函数 $f(x)$ 在 $x_0$ 处可微, 则称向量
\[
  \left(\frac{\partial f(x_0)}{\partial x_1},\frac{\partial f(x_0)}{\partial x_2},\cdots,\frac{\partial f(x_0)}{\partial x_n}\right)
\]
为 $f(x)$ 在 $x_0$ 处的梯度, 记为 $\operatorname{grad}f(x_0)$, 即
\[
  \operatorname{grad}f(x_0)=
  \left(\frac{\partial f(x_0)}{\partial x_1},\frac{\partial f(x_0)}{\partial x_2},\cdots,\frac{\partial f(x_0)}{\partial x_n}\right).
\]

设 $f,g$ 均是可微函数, 则从定义容易推出梯度有以下简单性质:
\[
\begin{aligned}
  &(1)\quad \operatorname{grad}C=0,\ C\text{ 为常数};\\
  &(2)\quad \text{对于 } \forall\alpha,\beta\in\mathbb R,\ \text{有 }
  \operatorname{grad}(\alpha f+\beta g)=\alpha\operatorname{grad}f+\beta\operatorname{grad}g;\\
  &(3)\quad \operatorname{grad}(f\cdot g)=f\cdot\operatorname{grad}g+g\cdot\operatorname{grad}f;\\
  &(4)\quad \operatorname{grad}\frac fg=\frac1{g^2}(g\cdot\operatorname{grad}f-f\cdot\operatorname{grad}g)\quad(g\ne0).
\end{aligned}
\]
从梯度的定义可以看出, 当 $f(x)$ 在 $x_0$ 处可微时, $f(x)$ 沿方向 $v$ 的方向导数可以简单地记成
\[
  \frac{\partial f(x_0)}{\partial v}=\operatorname{grad}f(x_0)\cdot v.
\]
读者易于给出方向导数与梯度的关系的几何意义。

设函数 $u=f(x,y,z)$ 在 $(x_0,y_0,z_0)$ 处可微, $\operatorname{grad}f(x_0,y_0,z_0)$ 不是零向量. 记 $H=|\operatorname{grad}f(x_0,y_0,z_0)|$, $l_v$ 为从 $(x_0,y_0,z_0)$ 出发由方向 $v$ 确定的射线, 则对于 $\forall h\in(-H,H)$, 所有使得 $\dfrac{\partial f}{\partial v}=h$ 的 $l_v$ 的集合构成了顶点在 $(x_0,y_0,z_0)$, 母线是 $l_v$ 的锥面. 当 $h$ 从 $-H$ 连续变为 $H$ 时, 其母线 $l_v$ 与 $\operatorname{grad}f(x_0,y_0,z_0)$ 的夹角从 $\pi$ 变到 $0$. 特别地, 当它们的夹角为 $\dfrac\pi2$ 时, $l_v$ 的全体构成了过 $(x_0,y_0,z_0)$ 且以 $\operatorname{grad}f(x_0,y_0,z_0)$ 为法线的平面。

\end{definition}
\begin{example}\label{course-example-14-1-7}\quad 设 $k\geqslant2$ 为正整数, 函数 $f(x,y)$ 在极坐标 $(r,\theta)$ 下有表示式
\[
  f(x,y)=
  \begin{cases}
    r\sin k\theta, &(x,y)\ne(0,0),\\
    0, &(x,y)=(0,0).
  \end{cases}
\]
(1) 求 $f(x,y)$ 在 $(0,0)$ 处的方向导数; (2) $f(x,y)$ 在 $(0,0)$ 处是否连续? (3) $f(x,y)$ 在 $(0,0)$ 处是否可微?

\end{example}
\begin{solve}\quad (1) 对固定的 $\theta_0\in[0,2\pi)$, $\theta=\theta_0$ 即确定一个从原点出发的方向 $v_0=(\cos\theta_0,\sin\theta_0)$. 我们有
\[
  \frac{\partial f(0,0)}{\partial v_0}
  =\lim_{r\to0}\frac{r\sin k\theta_0-0}{r}
  =\sin k\theta_0.
\]
因此 $f(x,y)$ 在 $(0,0)$ 处沿方向 $(\cos\theta,\sin\theta)$ 的方向导数即为 $\sin k\theta$。

(2) 显然 $\lim\limits_{(x,y)\to(0,0)}f(x,y)=0$, 所以 $f(x,y)$ 在 $(0,0)$ 处连续。

(3) 当 $k\geqslant2$ 时, 由于 $\sin k\theta$ 至少在 $\theta=\dfrac{\pi}{2k}, \dfrac{5\pi}{2k}\in[0,2\pi)$ 处取最大值 $1$, 我们断言 $f(x,y)$ 在 $(0,0)$ 处不可微. 如若断言不真, 假设 $f(x,y)$ 在 $(0,0)$ 处可微, 则 $\operatorname{grad}f(0,0)$ 是非零向量, 因此它的方向导数只能在一个方向即梯度的方向达到最大值. 此矛盾证明了断言。

\end{solve}
\begin{remark}\quad 当 $k=2$ 且 $r\ne0$ 时,
\[
  r\sin2\theta=\frac{2r^2\sin\theta\cos\theta}{r}
  =\frac{2xy}{\sqrt{x^2+y^2}}.
\]
由此我们可将 $f(x,y)$ 写成如下形式:
\[
  f(x,y)=
  \begin{cases}
    \dfrac{2xy}{\sqrt{x^2+y^2}}, &x^2+y^2\ne0,\\
    0, &x^2+y^2=0.
  \end{cases}
\]
同理, 当 $k=3$ 且 $r\ne0$ 时,
\[
  f(x,y)=
  \begin{cases}
    \dfrac{3x^2y-y^3}{x^2+y^2}, &x^2+y^2\ne0,\\
    0, &x^2+y^2=0.
  \end{cases}
\]
请读者自己对上述函数在 $(x,y)$ 坐标下重新解例~\ref{course-example-14-1-7}。

\end{remark}
\subsection{向量函数的导数与全微分}

在本小节中, 我们将定义向量函数的导数与全微分. 由于涉及向量函数, 今后我们有时用行向量, 但有时也用列向量来记同一个向量, 不过从上下文中, 读者可以清楚地知道它何时是行向量, 何时是列向量。

\begin{definition}\label{course-definition-14-1-5}\quad 设向量函数 $u=f(x)=(f_1(x),f_2(x),\cdots,f_m(x))^\mathrm T$ 在区域 $D\subset\mathbb R^n$ 上有定义, $x_0\in D$, $\Delta x=(\Delta x_1,\cdots,\Delta x_n)^\mathrm T$ 为 $x$ 在 $x_0$ 处的全增量. 如果存在 $m\times n$ 矩阵
\[
  A=
  \begin{pmatrix}
    A_{11}&A_{12}&\cdots&A_{1n}\\
    A_{21}&A_{22}&\cdots&A_{2n}\\
    \cdots&\cdots&\cdots&\cdots\\
    A_{m1}&A_{m2}&\cdots&A_{mn}
  \end{pmatrix}
\]
使得当 $|\Delta x|\to0$ 时, 下式成立:
\begin{equation}
\begin{aligned}
  \Delta f(x_0)
  &=(\Delta f_1(x_0),\cdots,\Delta f_m(x_0))^\mathrm T\\
  &=A(\Delta x_1,\cdots,\Delta x_n)^\mathrm T
  +(\alpha_1(|\Delta x|),\cdots,\alpha_m(|\Delta x|))^\mathrm T,
\end{aligned}
\label{course-eq-14-1-3}
\end{equation}
其中 $A$ 中的元素仅依赖于 $x_0$ 而不依赖于 $\Delta x$, 对于 $\forall j(1\leqslant j\leqslant m)$, $\alpha_j(|\Delta x|)$ 依赖于 $\Delta x$, 并且满足 $\lim\limits_{|\Delta x|\to0}\dfrac{\alpha_j(|\Delta x|)}{|\Delta x|}=0$, 则称 $f(x)$ 在 $x_0$ 处可微或可导, 矩阵 $A$ 称为 $f(x)$ 在 $x_0$ 处的 Fréchet 导数 (简称导数), 记做 $f'(x_0)$ 或 $Df(x_0)$; $A\Delta x$ 称为 $f(x)$ 在 $x_0$ 处的全微分 (简称微分), 记做 $\mathrm df(x_0)$, 即
\[
  \mathrm df(x_0)=A\Delta x=f'(x_0)\Delta x=Df(x_0)\Delta x.
\]
在上述定义中, 若规定 $\mathrm dx=\Delta x$, 则有 $\mathrm df(x_0)=f'(x_0)\mathrm dx$。

式~\eqref{course-eq-14-1-3}  用矩阵可以简写成
\[
  \Delta f(x_0)=f(x_0+\Delta x)-f(x_0)=A\Delta x+\alpha(|\Delta x|),
\]
其中 $\alpha(|\Delta x|)=(\alpha_1(|\Delta x|),\cdots,\alpha_m(|\Delta x|))^\mathrm T$ 满足
\[
  \lim_{|\Delta x|\to0}\frac{\alpha(|\Delta x|)}{|\Delta x|}=0.
\]

下面我们考虑如何判断一个向量函数的可微性, 以及当它可微时, 如何求它的导数. 下面的定理对此给出了回答。

\end{definition}
\begin{theorem}\label{course-theorem-14-1-4}\quad 设 $D$ 是 $\mathbb R^n$ 中的区域, $x_0\in D$, 向量函数 $f(x)=(f_1(x),\cdots,f_m(x))^\mathrm T$ 在 $D$ 上有定义, 则 $f(x)$ 在 $x_0$ 处可微的充分必要条件是对于 $\forall j(1\leqslant j\leqslant m)$, $f_j(x)$ 在 $x_0$ 处可微. 记
\[
  A=
  \begin{pmatrix}
    \dfrac{\partial f_1(x_0)}{\partial x_1}&\dfrac{\partial f_1(x_0)}{\partial x_2}&\cdots&\dfrac{\partial f_1(x_0)}{\partial x_n}\\
    \dfrac{\partial f_2(x_0)}{\partial x_1}&\dfrac{\partial f_2(x_0)}{\partial x_2}&\cdots&\dfrac{\partial f_2(x_0)}{\partial x_n}\\
    \cdots&\cdots&\cdots&\cdots\\
    \dfrac{\partial f_m(x_0)}{\partial x_1}&\dfrac{\partial f_m(x_0)}{\partial x_2}&\cdots&\dfrac{\partial f_m(x_0)}{\partial x_n}
  \end{pmatrix},
\]
则当 $f(x)$ 在 $x_0$ 处可微时, 有
\[
  \mathrm df(x_0)=A\,\mathrm dx\quad\text{或}\quad f'(x_0)=A.
\]

\end{theorem}
\begin{proof}\quad \keyterm{必要性}\quad 设 $f(x)$ 在 $x_0$ 处可微, 从而存在 $m\times n$ 矩阵 $A=(A_{ji})$ 使得当 $|\Delta x|\to0$ 时, 有
\[
\begin{aligned}
  \Delta f(x_0)
  &=(\Delta f_1(x_0),\cdots,\Delta f_m(x_0))^\mathrm T\\
  &=\left(\sum_{i=1}^n A_{1i}\Delta x_i,\cdots,\sum_{i=1}^n A_{mi}\Delta x_i\right)^\mathrm T
  +(\alpha_1(|\Delta x|),\cdots,\alpha_m(|\Delta x|))^\mathrm T,
\end{aligned}
\]
其中 $\alpha_j(|\Delta x|)$ 依赖于 $\Delta x$, 且 $\lim\limits_{|\Delta x|\to0}\dfrac{\alpha_j(|\Delta x|)}{|\Delta x|}=0(j=1,\cdots,m)$. 比较上式两边向量的分量, 对 $j=1,\cdots,m$, 当 $|\Delta x|\to0$ 时, 有
\[
  \Delta f_j(x_0)=\sum_{i=1}^n A_{ji}\Delta x_i+\alpha_j(|\Delta x|).
\]
由多元函数可微的定义知, $f_j(x)$ 在 $x_0$ 处可微, 并且
\[
  A_{ji}=\frac{\partial f_j(x_0)}{\partial x_i}
  \quad (i=1,\cdots,n;\ j=1,\cdots,m).
\]

\keyterm{充分性}\quad 设对于 $\forall j(1\leqslant j\leqslant m)$, $f_j(x)$ 在 $x_0$ 处可微, 则有
\[
  \Delta f_j(x_0)=\sum_{i=1}^n\frac{\partial f_j(x_0)}{\partial x_i}\Delta x_i+\alpha_j(|\Delta x|),
\]
其中 $\alpha_j(|\Delta x|)$ 依赖于 $\Delta x$, 且 $\lim\limits_{|\Delta x|\to0}\dfrac{\alpha_j(|\Delta x|)}{|\Delta x|}=0$. 因此
\[
\begin{aligned}
  \Delta f(x_0)
  &=(\Delta f_1(x_0),\cdots,\Delta f_m(x_0))^\mathrm T\\
  &=
  \begin{pmatrix}
    \dfrac{\partial f_1(x_0)}{\partial x_1}&\cdots&\dfrac{\partial f_1(x_0)}{\partial x_n}\\
    \dfrac{\partial f_2(x_0)}{\partial x_1}&\cdots&\dfrac{\partial f_2(x_0)}{\partial x_n}\\
    \cdots&\cdots&\cdots\\
    \dfrac{\partial f_m(x_0)}{\partial x_1}&\cdots&\dfrac{\partial f_m(x_0)}{\partial x_n}
  \end{pmatrix}
  \begin{pmatrix}
    \Delta x_1\\
    \Delta x_2\\
    \vdots\\
    \Delta x_n
  \end{pmatrix}
  +
  \begin{pmatrix}
    \alpha_1(|\Delta x|)\\
    \alpha_2(|\Delta x|)\\
    \vdots\\
    \alpha_m(|\Delta x|)
  \end{pmatrix}\\
  &=A\Delta x+\alpha(|\Delta x|),
\end{aligned}
\]
其中 $\alpha(|\Delta x|)$ 满足 $\lim\limits_{|\Delta x|\to0}\dfrac{\alpha(|\Delta x|)}{|\Delta x|}=0$. 由定义知 $f(x)$ 在 $x_0$ 处可微, 且 $f'(x_0)=A$. 证毕。

当向量函数 $f(x)=(f_1,\cdots,f_m)^\mathrm T$ 中的每个分量函数在 $x_0$ 处均可微时, 矩阵 $A=f'(x_0)$ 称为 $f(x)$ 在 $x_0$ 处的雅可比 (Jacobi) 矩阵, 记为 $J_f(x_0)$. 特别地, 当 $f(x)$ 是 $n$ 维向量函数时, $A$ 是 $n\times n$ 矩阵, 此时 $J_f(x_0)$ 的行列式称为 $f(x)$ 在 $x_0$ 处的雅可比行列式, 记为
\[
  |f'(x_0)|\quad\text{或}\quad
  \left.\frac{\partial(f_1,f_2,\cdots,f_n)}{\partial(x_1,x_2,\cdots,x_n)}\right|_{x_0}.
\]
显然, 对于一个向量函数 $f(x)$, 若其分量函数在 $x_0$ 处具有各个偏导数时, 我们就可以形式地定义 $J_f(x_0)$, 但当 $f(x)$ 在 $x_0$ 处不可微, 仅仅存在所有的偏导数时, $J_f(x_0)$ 不能刻画 $f(x)$ 在 $x_0$ 附近的变化情况。

当 $f_j(x)(1\leqslant j\leqslant m)$ 的各个偏导数都在 $x_0$ 处连续时, $f_j(x)$ 在 $x_0$ 处可微, 因此 $f(x)$ 在 $x_0$ 处可微. 另外, 若对于 $\forall j(1\leqslant j\leqslant m)$, $f_j(x)$ 的各个偏导数在区域 $D$ 上连续, 我们称 $f(x)$ 在 $D$ 上是 $C^1$ 的, 记为 $f(x)\in C^1(D)$. 特别地, 我们称 $\mathbb R^n$ 中区域 $D$ 到 $\Omega$ 的变换 $y=f(x)$ 是 $C^1$ 的, 如果 $f(x)\in C^1(D)$, 并且 $f^{-1}(y)\in C^1(\Omega)$。

\end{proof}
\begin{remark}\quad 利用向量函数导数的记号, 对一个多元函数 $f(x)$, 当它可微时, 我们有
\[
  f'(x)=\left(\frac{\partial f}{\partial x_1},\cdots,\frac{\partial f}{\partial x_n}\right)=\operatorname{grad}f(x).
\]

\end{remark}
\begin{example}\label{course-example-14-1-8}\quad 设向量 $w=(r,\varphi,\theta)$, 向量函数
\[
  f(w)=(x(w),y(w),z(w))^\mathrm T
  =(r\sin\varphi\cos\theta,r\sin\varphi\sin\theta,r\cos\varphi)^\mathrm T,
\]
其中 $D=\{(r,\varphi,\theta):0<r<+\infty,0<\varphi<\pi,0<\theta<2\pi\}$; 试求 $f'(w)$ 以及 $f(w)$ 在 $w$ 处的雅可比行列式。

\end{example}
\begin{solve}\quad 由于 $x=r\sin\varphi\cos\theta$, $y=r\sin\varphi\sin\theta$, $z=r\cos\varphi$ 均在 $D$ 上具有连续偏导数, 因此 $f'(w)$ 存在. 由计算得
\[
\begin{aligned}
  f'(w)
  &=
  \begin{pmatrix}
    \dfrac{\partial x}{\partial r}&\dfrac{\partial x}{\partial\varphi}&\dfrac{\partial x}{\partial\theta}\\
    \dfrac{\partial y}{\partial r}&\dfrac{\partial y}{\partial\varphi}&\dfrac{\partial y}{\partial\theta}\\
    \dfrac{\partial z}{\partial r}&\dfrac{\partial z}{\partial\varphi}&\dfrac{\partial z}{\partial\theta}
  \end{pmatrix}\\
  &=
  \begin{pmatrix}
    \sin\varphi\cos\theta&r\cos\varphi\cos\theta&-r\sin\varphi\sin\theta\\
    \sin\varphi\sin\theta&r\cos\varphi\sin\theta&r\sin\varphi\cos\theta\\
    \cos\varphi&-r\sin\varphi&0
  \end{pmatrix}.
\end{aligned}
\]
对上述矩阵取行列式, 即得 $f(w)$ 在 $w=(r,\varphi,\theta)$ 处的雅可比行列式
\[
  \frac{\partial(x,y,z)}{\partial(r,\varphi,\theta)}=r^2\sin\varphi.
\]

\end{solve}


\section{多元函数求导法}\label{course-14-2}
本节用于核对多元极限、隐函数、驻点与约束极值题中的可微性和非退化条件。仅有偏导数存在不能替代全微分存在；用必要条件列出候选点以后，还要比较边界并验证极值是否取得。

\subsection{导数的四则运算}

在本小节中, 我们总假定涉及的多元 (向量) 函数存在各个偏导数, 然后来讨论多元函数求偏导数的一些公式. 读者应该注意的是, 一般情况下一个函数的导数是一个向量 (或矩阵)。

\begin{theorem}\label{course-theorem-14-2-1}\quad 设函数 $f(x),g(x)$ 在区域 $D\subset\mathbb R^n$ 上可导, 则对于 $\forall x\in D$, 有
\[
\begin{aligned}
  &(1)\quad (f(x)\pm g(x))'=f'(x)\pm g'(x);\\
  &(2)\quad (f(x)g(x))'=f(x)g'(x)+g(x)f'(x);\\
  &(3)\quad \left(\frac{f(x)}{g(x)}\right)'=
  \frac{g(x)f'(x)-f(x)g'(x)}{g^2(x)}\quad (g(x)\ne0).
\end{aligned}
\]
如果 $f(x):\mathbb R^n\to\mathbb R$ 是一个 $n$ 元可微函数, $g(x):\mathbb R^n\to\mathbb R^m$ 是一个可微的 $m$ 维列向量函数, 则有
\[
  (4)\quad (f(x)g(x))'=f(x)g'(x)+g(x)f'(x).
\]

\end{theorem}
\begin{proof}\quad 我们只证 (3). (1), (2) 和 (4) 的证明留给读者. 由定义有
\[
  \left(\frac{f(x)}{g(x)}\right)'=
  \left(
    \frac{\partial\left(\dfrac{f(x)}{g(x)}\right)}{\partial x_1},
    \frac{\partial\left(\dfrac{f(x)}{g(x)}\right)}{\partial x_2},
    \cdots,
    \frac{\partial\left(\dfrac{f(x)}{g(x)}\right)}{\partial x_n}
  \right).
\]
由于
\[
  \frac{\partial\left(\dfrac{f(x)}{g(x)}\right)}{\partial x_i}
  =
  \frac{g(x)\dfrac{\partial f(x)}{\partial x_i}-f(x)\dfrac{\partial g(x)}{\partial x_i}}{g^2(x)}
  \quad (i=1,2,\cdots,n),
\]
将其代入上式即知 (3) 成立. 证毕。

\end{proof}
\begin{remark}\quad 注意到 (4) 中等式左边 $(f(x)g(x))'$ 是一个 $m\times n$ 矩阵, $g(x)$ 是一个 $m\times1$ 矩阵, 而 $f'(x)$ 是一个 $1\times n$ 矩阵, 因此 $g(x)f'(x)$ 是一个 $m\times n$ 矩阵. 显然, $f(x)g'(x)$ 为一个 $m\times n$ 矩阵. 至于 (4) 的证明, 读者只要将 (4) 中等式左边矩阵中的每个分量与其右边矩阵中相对应的分量比较即可。

\end{remark}
\subsection{复合函数的求导法}

下面我们主要来考虑复合函数的求导法。

\begin{theorem}\label{course-theorem-14-2-2}\quad 设函数 $f(u)=f(u_1,u_2,\cdots,u_m)$ 在区域 $\Omega\subset\mathbb R^m$ 上有定义, 并且在 $u_0=(u_1^0,u_2^0,\cdots,u_m^0)^\mathrm T\in\Omega$ 处可微, 再设
\[
  u=u(x)=(u_1(x),u_2(x),\cdots,u_m(x))^\mathrm T
\]
在区域 $D\subset\mathbb R^n$ 上有定义, 在 $x_0=(x_1^0,x_2^0,\cdots,x_n^0)\in D$ 处可微, 并且 $u_0=u(x_0)$, 则 $f(u(x))$ 在 $x_0$ 处可微, 并且
\[
  \mathrm df(u(x_0))=f'(u(x_0))u'(x_0)\mathrm dx.
\]

\end{theorem}
\begin{proof}\quad 由 $u$ 在 $x_0$ 处可微, 我们有
\[
  \Delta u(x_0)=u(x_0+\Delta x)-u(x_0)
  =u'(x_0)\Delta x+\alpha(|\Delta x|),
\]
其中 $\alpha(|\Delta x|)$ 依赖于 $\Delta x$, 且满足
\[
  \frac{\alpha(|\Delta x|)}{|\Delta x|}\to0\quad(|\Delta x|\to0).
\]
再由 $f(u)$ 在 $u_0$ 处可微, 有
\[
  \Delta f(u_0)=f(u_0+\Delta u)-f(u_0)=f'(u_0)\Delta u+\beta(|\Delta u|),
\]
其中 $\beta(|\Delta u|)$ 依赖于 $\Delta u$, 且满足
\[
  \frac{\beta(|\Delta u|)}{|\Delta u|}\to0\quad(|\Delta u|\to0).
\]
我们规定 $\beta(0)=0$. 在这样的规定下, $\beta(|\Delta u|)$ 在 $\Delta u=0$ 时连续。

我们有
\[
\begin{aligned}
  \Delta f(u(x_0))
  &=f(u(x_0+\Delta x))-f(u(x_0))\\
  &=f'(u(x_0))(u'(x_0)\Delta x+\alpha(|\Delta x|))
  +\beta(|u'(x_0)\Delta x+\alpha(|\Delta x|)|)\\
  &=f'(u(x_0))u'(x_0)\Delta x+\gamma(|\Delta x|),
\end{aligned}
\]
其中
\[
\begin{aligned}
  \gamma(|\Delta x|)
  &=f'(u(x_0))\alpha(|\Delta x|)
  +\beta(|\Delta u(x_0)|)\\
  &=f'(u(x_0))\alpha(|\Delta x|)
  +\beta(|u'(x_0)\Delta x+\alpha(|\Delta x|)|).
\end{aligned}
\]
下面我们证明
\[
  \frac{\gamma(|\Delta x|)}{|\Delta x|}\to0\quad(|\Delta x|\to0).
\]
首先, 显然有
\begin{equation}
  \lim_{|\Delta x|\to0}
  \frac{|f'(u(x_0))\alpha(|\Delta x|)|}{|\Delta x|}
  \leqslant
  \lim_{|\Delta x|\to0}
  \frac{|f'(u(x_0))|\cdot|\alpha(|\Delta x|)|}{|\Delta x|}=0.
\label{course-eq-14-2-1}
\end{equation}
其次, 注意到当 $|\Delta x|$ 很小时, 有
\[
\begin{aligned}
  \frac{|\Delta u(x_0)|}{|\Delta x|}
  &=\frac{|u'(x_0)\Delta x+\alpha(|\Delta x|)|}{|\Delta x|}\\
  &\leqslant \frac1{|\Delta x|}\bigl(\|u'(x_0)\|\,|\Delta x|+|\alpha(|\Delta x|)|\bigr)
  \leqslant \|u'(x_0)\|+1,
\end{aligned}
\]
其中 $\|u'(x_0)\|$ 是矩阵 $u'(x_0)$ 的范数. 因此当 $|\Delta x|\to0$ 时, 必有 $|\Delta u(x_0)|\to0$. 我们有
\[
  \frac{|\beta(|\Delta u(x_0)|)|}{|\Delta x|}
  =
  \begin{cases}
    \dfrac{|\beta(|\Delta u(x_0)|)|}{|\Delta u(x_0)|}
    \cdot\dfrac{|\Delta u(x_0)|}{|\Delta x|}, &|\Delta u(x_0)|\ne0,\\[1.2em]
    0, &|\Delta u(x_0)|=0,
  \end{cases}
\]
从而有
\begin{equation}
  \frac{|\beta(|\Delta u(x_0)|)|}{|\Delta x|}\to0
  \quad(|\Delta x|\to0).
\label{course-eq-14-2-2}
\end{equation}
结合式~\eqref{course-eq-14-2-1}, \eqref{course-eq-14-2-2}  得
\[
  \frac{\gamma(|\Delta x|)}{|\Delta x|}\to0\quad(|\Delta x|\to0).
\]
我们最后有
\[
  \Delta f(u(x_0))=f'(u(x_0))u'(x_0)\Delta x+\gamma(|\Delta x|),
\]
其中 $\gamma(|\Delta x|)$ 依赖于 $\Delta x$ 且 $\dfrac{\gamma(|\Delta x|)}{|\Delta x|}\to0(|\Delta x|\to0)$. 由微分的定义知, $f(u(x))$ 在 $x_0$ 处可微, 并且
\[
  \mathrm df(u(x_0))=f'(u(x_0))u'(x_0)\mathrm dx.
\]
证毕。

从定理~\ref{course-theorem-14-2-2}  我们可以推出以下结论。

\end{proof}
\noindent\keyterm{推论 1}\quad 设向量函数 $f(u)=(f_1(u),f_2(u),\cdots,f_l(u))^\mathrm T$ 在区域 $\Omega\subset\mathbb R^m$ 上可微, $u(x)=(u_1(x),u_2(x),\cdots,u_m(x))^\mathrm T$ 在区域 $D\subset\mathbb R^n$ 上可微, 且 $u(D)\subset\Omega$, 则 $g=f(u(x))$ 在 $D$ 上可微, 并且
\[
  \mathrm df(u(x))=f'(u(x))u'(x)\mathrm dx.
\]
顺便我们有
\[
  [f(u(x))]'=f'(u(x))u'(x).
\]

\begin{proof}\quad 由定理~\ref{course-theorem-14-2-2}  知 $\mathrm df_j(u(x))=f'_j(u(x))u'(x)\mathrm dx\ (j=1,2,\cdots,l)$, 从而
\[
  \mathrm df(u(x))=
  \begin{pmatrix}
    f'_1(u(x))u'(x)\mathrm dx\\
    f'_2(u(x))u'(x)\mathrm dx\\
    \vdots\\
    f'_l(u(x))u'(x)\mathrm dx
  \end{pmatrix}
  =f'(u(x))u'(x)\mathrm dx.
\]
证毕。

\end{proof}
\begin{remark}\quad 读者应该注意到 $f'(u(x))u'(x)$ 是一个 $l\times n$ 矩阵。

\end{remark}
\noindent\keyterm{推论 2}\quad 设 $D$ 与 $\Omega$ 为 $\mathbb R^n$ 中的区域, $y=f(x)$ 是 $D$ 到 $\Omega$ 的 $C^1$ 变换, 则对于 $\forall x\in D$, 有
\begin{equation}
  (f^{-1})'(y)\cdot f'(x)=E,
\label{course-eq-14-2-3}
\end{equation}
其中 $y=f(x)$; 对于 $\forall y\in\Omega$, 有
\begin{equation}
  f'(x)\cdot(f^{-1})'(y)=E,
\label{course-eq-14-2-4}
\end{equation}
其中 $x=f^{-1}(y)$. 特别地, 当 $y=f(x)$ 时, 有
\begin{equation}
  (f^{-1})'(y)=[f'(x)]^{-1},
\label{course-eq-14-2-5}
\end{equation}
其中 $E$ 是 $n\times n$ 单位矩阵, $[f'(x)]^{-1}$ 为 $f'(x)$ 的逆矩阵。

\begin{proof}\quad 由 $C^1$ 变换的定义知 $f$ 与 $f^{-1}$ 都是可微向量函数, 因此 $(f^{-1}\circ f)(x)=x$ 两边对 $x$ 求导数即得式~\eqref{course-eq-14-2-3}, 而 $(f\circ f^{-1})(y)=y$ 两边对 $y$ 求导数即得式~\eqref{course-eq-14-2-4}. 在式~\eqref{course-eq-14-2-3}  的两边同时右乘 $[f'(x)]^{-1}$ 即得式~\eqref{course-eq-14-2-5}. 证毕。

\end{proof}
\begin{remark}\quad 在推论 2 中, 如果我们在式~\eqref{course-eq-14-2-4}  两边取行列式, 则有
\begin{equation}
  \frac{\partial(y_1,\cdots,y_n)}{\partial(x_1,\cdots,x_n)}
  \cdot
  \frac{\partial(x_1,\cdots,x_n)}{\partial(y_1,\cdots,y_n)}
  =1.
\label{course-eq-14-2-6}
\end{equation}
下面的推论 3 是多元复合函数求偏导数的基础。

\end{remark}
\noindent\keyterm{推论 3}\quad 设函数 $f(u)=f(u_1,u_2,\cdots,u_m)$ 在区域 $\Omega\subset\mathbb R^m$ 上有定义, 并且在 $u_0=(u_1^0,u_2^0,\cdots,u_m^0)^\mathrm T\in\Omega$ 处可微, 又设向量函数
\[
  u=u(x)=(u_1(x),u_2(x),\cdots,u_m(x))^\mathrm T
\]
在区域 $D\subset\mathbb R^n$ 上有定义, 在 $x_0=(x_1^0,x_2^0,\cdots,x_n^0)\in D$ 处可微, 并且 $u_0=u(x_0)$, 则对于 $\forall i(1\leqslant i\leqslant n)$, $f(u(x))$ 在 $x_0$ 处关于 $x_i$ 可偏导, 并且
\[
  \frac{\partial f(u(x_0))}{\partial x_i}
  =\sum_{j=1}^m\left(
  \frac{\partial f(u_0)}{\partial u_j}\cdot
  \frac{\partial u_j(x_0)}{\partial x_i}\right).
\]

\begin{proof}\quad 由定理~\ref{course-theorem-14-1-4}  的注有
\[
  \left.(f(u(x)))'\right|_{x=x_0}
  =\left(
  \frac{\partial f(u(x_0))}{\partial x_1},
  \frac{\partial f(u(x_0))}{\partial x_2},
  \cdots,
  \frac{\partial f(u(x_0))}{\partial x_n}
  \right).
\]
用矩阵形式表示 $f'(u(x_0))u'(x_0)$, 有
\[
\begin{aligned}
  f'(u(x_0))u'(x_0)
  &=
  \left(
  \frac{\partial f(u(x_0))}{\partial u_1},
  \frac{\partial f(u(x_0))}{\partial u_2},
  \cdots,
  \frac{\partial f(u(x_0))}{\partial u_m}
  \right)\\
  &\quad\cdot
  \begin{pmatrix}
    \dfrac{\partial u_1(x_0)}{\partial x_1}&\dfrac{\partial u_1(x_0)}{\partial x_2}&\cdots&\dfrac{\partial u_1(x_0)}{\partial x_n}\\
    \dfrac{\partial u_2(x_0)}{\partial x_1}&\dfrac{\partial u_2(x_0)}{\partial x_2}&\cdots&\dfrac{\partial u_2(x_0)}{\partial x_n}\\
    \vdots&\vdots&\ddots&\vdots\\
    \dfrac{\partial u_m(x_0)}{\partial x_1}&\dfrac{\partial u_m(x_0)}{\partial x_2}&\cdots&\dfrac{\partial u_m(x_0)}{\partial x_n}
  \end{pmatrix}\\
  &=
  \left(
  \sum_{j=1}^m\frac{\partial f(u_0)}{\partial u_j}\frac{\partial u_j(x_0)}{\partial x_1},
  \sum_{j=1}^m\frac{\partial f(u_0)}{\partial u_j}\frac{\partial u_j(x_0)}{\partial x_2},
  \cdots,
  \sum_{j=1}^m\frac{\partial f(u_0)}{\partial u_j}\frac{\partial u_j(x_0)}{\partial x_n}
  \right).
\end{aligned}
\]
由定理~\ref{course-theorem-14-2-2}  有 $\left.(f(u(x)))'\right|_{x=x_0}=f'(u(x_0))u'(x_0)$, 再比较它们的分量即得推论 3. 证毕。

\end{proof}
\begin{remark}\quad 利用定理~\ref{course-theorem-14-2-2}  类似的证明方法可以证明, 若将推论 3 中的条件 ``$u(x)$ 在 $x_0$ 处可微'' 减弱成 ``$u(x)$ 在 $x_0$ 处存在各个偏导数'', 则推论 3 的结论仍然成立. 请读者自己对此加以证明 (见本章习题). 但是, 在推论 3 中, 条件 ``$f(u)$ 在 $u_0$ 处可微'' 不能减弱成 ``$f(u)$ 在 $u_0$ 处存在各个偏导数''. 很容易找到例子来说明这一点, 例如取
\[
  f(u,v)=
  \begin{cases}
    1, &uv\ne0,\\
    0, &uv=0,
  \end{cases}
  \quad (u,v)\in\mathbb R^2,
\]
而令
\[
  \begin{cases}
    u=t,\\
    v=t
  \end{cases}
  \quad(-\infty<t<+\infty),
\]
则
\[
  g(t)=f(t,t)=
  \begin{cases}
    1, &t\ne0,\\
    0, &t=0
  \end{cases}
\]
在 $t=0$ 处不连续, 从而不可导。

如同一元函数, 我们将推论 3 中给出的复合函数求导公式称为链锁法则. 为了今后应用方便, 下面我们列出常见的二元复合函数的求导公式. 设 $z=f(u,v), u=u(x,y), v=v(x,y)$ 都为可微函数, 则
\[
  \frac{\partial z}{\partial x}
  =\frac{\partial f}{\partial u}\cdot\frac{\partial u}{\partial x}
  +\frac{\partial f}{\partial v}\cdot\frac{\partial v}{\partial x},
\]
\[
  \frac{\partial z}{\partial y}
  =\frac{\partial f}{\partial u}\cdot\frac{\partial u}{\partial y}
  +\frac{\partial f}{\partial v}\cdot\frac{\partial v}{\partial y}.
\]
有时, 若给出的函数具有形式 $z=f(u(x,y),v(x,y))$, 我们可以用以下记号:
\[
  \frac{\partial z}{\partial x}=f'_1\frac{\partial u}{\partial x}+f'_2\frac{\partial v}{\partial x},
  \qquad
  \frac{\partial z}{\partial y}=f'_1\frac{\partial u}{\partial y}+f'_2\frac{\partial v}{\partial y},
\]
其中 $f'_i(i=1,2)$ 指的是 $f$ 对其第 $i$ 个变量求偏导数. 对于具有多个中间变量的复合函数, 我们可以引进类似的记号, 这对抽象形式给出的函数的求导运算提供了简洁的记号。

\end{remark}
\begin{example}\label{course-example-14-2-1}\quad 设函数 $u=e^{x^2+y^2+z^2}; z=x^2\sin y$, 求 $\dfrac{\partial u}{\partial x}, \dfrac{\partial u}{\partial y}$。

\end{example}
\begin{solve}\quad
\[
\begin{aligned}
  \frac{\partial u}{\partial x}
  &=\frac{\partial e^{x^2+y^2+z^2}}{\partial x}\cdot1
  +\frac{\partial e^{x^2+y^2+z^2}}{\partial z}\cdot\frac{\partial z}{\partial x}\\
  &=2xe^{x^2+y^2+z^2}+e^{x^2+y^2+z^2}(2z)(2x\sin y)\\
  &=2x(1+2x^2\sin^2 y)e^{x^2+y^2+z^2},
\end{aligned}
\]
\[
\begin{aligned}
  \frac{\partial u}{\partial y}
  &=\frac{\partial e^{x^2+y^2+z^2}}{\partial y}\cdot1
  +\frac{\partial e^{x^2+y^2+z^2}}{\partial z}\cdot\frac{\partial z}{\partial y}\\
  &=2ye^{x^2+y^2+z^2}+e^{x^2+y^2+z^2}(2z)(x^2\cos y)\\
  &=(2y+x^4\sin2y)\cdot e^{x^2+y^2+z^2}.
\end{aligned}
\]

\end{solve}
\begin{example}\label{course-example-14-2-2}\quad 设函数 $u=f\left(xy,\dfrac yx,yz\right)$, 并设 $f$ 是可微函数, 试求 $\dfrac{\partial u}{\partial x}, \dfrac{\partial u}{\partial y}$ 和 $\dfrac{\partial u}{\partial z}$。

\end{example}
\begin{solve}\quad 对于这种抽象形式给出的函数, 我们采用上述约定的记号来求偏导数:
\[
  \frac{\partial u}{\partial x}
  =f'_1\frac{\partial(xy)}{\partial x}+f'_2\frac{\partial(y/x)}{\partial x}+f'_3\cdot0
  =yf'_1-\frac{y}{x^2}f'_2,
\]
\[
  \frac{\partial u}{\partial y}
  =f'_1\frac{\partial(xy)}{\partial y}
  +f'_2\frac{\partial(y/x)}{\partial y}
  +f'_3\frac{\partial(yz)}{\partial y}
  =xf'_1+\frac1x f'_2+zf'_3,
\]
\[
  \frac{\partial u}{\partial z}
  =f'_1\cdot0+f'_2\cdot0+f'_3\frac{\partial(yz)}{\partial z}
  =yf'_3.
\]

\end{solve}
\begin{example}\label{course-example-14-2-3}\quad 记向量 $v=(r,\theta_1,\cdots,\theta_{n-1})$, 并设向量函数
\[
  f(v)=
  \begin{pmatrix}
    x_1(v)\\
    x_2(v)\\
    \vdots\\
    x_{n-1}(v)\\
    x_n(v)
  \end{pmatrix}
  =
  \begin{pmatrix}
    r\cos\theta_1\\
    r\sin\theta_1\cos\theta_2\\
    \vdots\\
    r\sin\theta_1\sin\theta_2\cdots\cos\theta_{n-1}\\
    r\sin\theta_1\sin\theta_2\cdots\sin\theta_{n-1}
  \end{pmatrix},
\]
试求 $f'(v)$ 及其雅可比行列式。

\end{example}
\begin{solve}\quad 由定义有

\[
  f'(v)=
  \begin{pmatrix}
    \cos\theta_1&-r\sin\theta_1&\cdots&0\\
    \sin\theta_1\cos\theta_2&r\cos\theta_1\cos\theta_2&\cdots&0\\
    \vdots&\vdots&\ddots&\vdots\\
    \sin\theta_1\cdots\cos\theta_{n-1}&r\cos\theta_1\cdots\cos\theta_{n-1}&\cdots&-r\sin\theta_1\cdots\sin\theta_{n-1}\\
    \sin\theta_1\cdots\sin\theta_{n-1}&r\cos\theta_1\cdots\sin\theta_{n-1}&\cdots&r\sin\theta_1\cdots\cos\theta_{n-1}
  \end{pmatrix}.
\]
下面用归纳法来求 $f'(v)$ 的行列式 $|f'(v)|$. 当 $n=2$ 时, 有 $|f'(v)|=r$; 当 $n=3$ 时, 有 $|f'(v)|=r^2\sin\theta_1$. 我们猜测对于 $n\geqslant3$, 有
\[
  |f'(v)|=r^{n-1}\sin^{n-2}\theta_1\sin^{n-3}\theta_2\cdots\sin\theta_{n-2}.
\]
事实上, 已知 $n=3$ 时上式成立. 假设 $n=k-1$ 时上式成立, 现在我们来证明当 $n=k$ 时上式成立. 为此将变换 $f$ 拆成下列两个变换的复合:
\[
  f_1:\begin{cases}
    x_1=y_1,\\
    x_2=y_2\cos y_3,\\
    x_3=y_2\sin y_3\cos y_4,\\
    \cdots\\
    x_{n-1}=y_2\sin y_3\sin y_4\cdots\cos y_n,\\
    x_n=y_2\sin y_3\sin y_4\cdots\sin y_n,
  \end{cases}
  \qquad
  f_2:\begin{cases}
    y_1=r\cos\theta_1,\\
    y_2=r\sin\theta_1,\\
    y_3=\theta_2,\\
    \cdots\\
    y_{n-1}=\theta_{n-2},\\
    y_n=\theta_{n-1}.
  \end{cases}
\]
则有 $f=f_1\circ f_2$. 当 $n=k$ 时, 由归纳法假设得
\[
\begin{aligned}
  |f'(v)|
  &=
  \frac{\partial(x_1,x_2,\cdots,x_k)}
  {\partial(r,\theta_1,\cdots,\theta_{k-1})}\\
  &=
  \left.
  \frac{\partial(x_1,x_2,\cdots,x_k)}
  {\partial(y_1,y_2,\cdots,y_k)}
  \right|_{f_2(v)}
  \frac{\partial(y_1,y_2,\cdots,y_k)}
  {\partial(r,\theta_1,\cdots,\theta_{k-1})}\\
  &=
  \left.(y_2^{k-2}\sin^{k-3}y_3\sin^{k-4}y_4\cdots\sin y_{k-1})\right|_{f_2(v)}
  \cdot r\\
  &=r^{k-1}\sin^{k-2}\theta_1\sin^{k-3}\theta_2\cdots\sin\theta_{k-2}.
\end{aligned}
\]

\end{solve}
\begin{example}\label{course-example-14-2-4}\quad 设 $f(x)=(f_1(x),f_2(x),\cdots,f_m(x))$ 和 $g(x)=(g_1(x),g_2(x),\cdots,g_m(x))$ 为 $D\subset\mathbb R^n$ 到 $\mathbb R^m$ 的可微向量函数, 证明:
\[
  [f(x)(g(x))^\mathrm T]'=f(x)[(g(x))^\mathrm T]'+g(x)[(f(x))^\mathrm T]'.
\]

\end{example}
\begin{proof}\quad 令
\[
  F(x)=
  \begin{pmatrix}
    (f(x))^\mathrm T\\
    (g(x))^\mathrm T
  \end{pmatrix}
  =(f_1(x),\cdots,f_m(x),g_1(x),\cdots,g_m(x))^\mathrm T,
\]
则 $F(x)$ 是 $D\to\mathbb R^{2m}$ 的可微函数. 记 $y=(y_1,y_2,\cdots,y_{2m})^\mathrm T$, 并令
\[
  G(y)=\sum_{i=1}^m y_i y_{m+i},
\]
则 $G(y)$ 是 $\mathbb R^{2m}\to\mathbb R$ 的可微函数. 我们可以验证
\[
  f(x)(g(x))^\mathrm T=G(F(x)).
\]
由链锁法则则有
\[
\begin{aligned}
  [f(x)(g(x))^\mathrm T]'
  &=[G(F(x))]'\\
  &=(y_{m+1},\cdots,y_{2m},y_1,\cdots,y_m)\big|_{y=F(x)}
  \begin{pmatrix}
    [(f(x))^\mathrm T]'\\
    [(g(x))^\mathrm T]'
  \end{pmatrix}\\
  &=f(x)[(g(x))^\mathrm T]'+g(x)[(f(x))^\mathrm T]'.
\end{aligned}
\]

\end{proof}
\begin{example}\label{course-example-14-2-5}\quad 设 $f(x,y)$ 在 $\mathbb R^2$ 上具有连续偏导数, 试用极坐标来表示
\[
  |f'(x,y)|^2=
  \left(\frac{\partial f}{\partial x}\right)^2+
  \left(\frac{\partial f}{\partial y}\right)^2
  =|\operatorname{grad}f(x,y)|^2.
\]

\end{example}
\begin{solve}\quad 由
\[
  \begin{cases}
    x=r\cos\theta,\\
    y=r\sin\theta
  \end{cases}
\]
得
\[
  \frac{\partial f}{\partial x}
  =\frac{\partial f}{\partial r}\cdot\frac{\partial r}{\partial x}
  +\frac{\partial f}{\partial\theta}\cdot\frac{\partial\theta}{\partial x},
  \quad
  \frac{\partial f}{\partial y}
  =\frac{\partial f}{\partial r}\cdot\frac{\partial r}{\partial y}
  +\frac{\partial f}{\partial\theta}\cdot\frac{\partial\theta}{\partial y}.
\]
由 $r=\sqrt{x^2+y^2}$ 及 $\theta=\arctan\dfrac yx$ (此处分象限讨论, 我们以第一象限为例求之) 得
\[
  \frac{\partial r}{\partial x}=\frac{x}{\sqrt{x^2+y^2}}=\cos\theta,\quad
  \frac{\partial r}{\partial y}=\frac{y}{\sqrt{x^2+y^2}}=\sin\theta,
\]
\[
  \frac{\partial\theta}{\partial x}=-\frac{y}{x^2+y^2}=-\frac{\sin\theta}{r},\quad
  \frac{\partial\theta}{\partial y}=\frac{x}{x^2+y^2}=\frac{\cos\theta}{r}.
\]
因此
\[
  \frac{\partial f}{\partial x}
  =\cos\theta\frac{\partial f}{\partial r}
  -\frac{\sin\theta}{r}\cdot\frac{\partial f}{\partial\theta},
  \quad
  \frac{\partial f}{\partial y}
  =\sin\theta\frac{\partial f}{\partial r}
  +\frac{\cos\theta}{r}\cdot\frac{\partial f}{\partial\theta};
\]
从而有
\[
  |\operatorname{grad}f(x,y)|^2
  =
  \left(\frac{\partial f}{\partial x}\right)^2+
  \left(\frac{\partial f}{\partial y}\right)^2
  =
  \left(\frac{\partial f}{\partial r}\right)^2+
  \frac1{r^2}\left(\frac{\partial f}{\partial\theta}\right)^2.
\]

\end{solve}
\begin{example}\label{course-example-14-2-6}\quad 设函数 $u=f(x-\lambda t)$ ($\lambda$ 是常数), 其中 $f$ 是 $\mathbb R$ 中的一个可微函数 (即 $f(y)$ 是可微函数), 证明它满足微分方程
\begin{equation}
  \frac{\partial u}{\partial t}+\lambda\frac{\partial u}{\partial x}=0,
\label{course-eq-14-2-7}
\end{equation}
并证明方程 \eqref{course-eq-14-2-7}  的解一定具有形式 $u=g(x-\lambda t)$, 其中 $g$ 是任一个 $\mathbb R$ 中的可微函数。

\end{example}
\begin{proof}\quad 记 $y=x-\lambda t$, 则
\[
  \frac{\partial u}{\partial t}
  =\frac{\mathrm df}{\mathrm dy}\cdot\frac{\partial y}{\partial t}
  =f'(y)(-\lambda)=-\lambda f'(y),
  \quad
  \frac{\partial u}{\partial x}
  =\frac{\mathrm df}{\mathrm dy}\cdot\frac{\partial y}{\partial x}
  =\frac{\partial f}{\partial y}=f'(y),
\]
从而
\[
  \frac{\partial u}{\partial t}+\lambda\frac{\partial u}{\partial x}
  =-\lambda f'(y)+\lambda f'(y)=0.
\]
设可微函数 $u=u(x,t)$ 是方程 \eqref{course-eq-14-2-7}  的解, 作变量替换
\[
  \begin{cases}
    x=y+\lambda z,\\
    z=t,
  \end{cases}
\]
则 $u(x,t)=u(y+\lambda z,z)\triangleq u_1(y,z)$. 由
\[
  \frac{\partial u}{\partial t}
  =\frac{\partial u_1}{\partial y}\cdot\frac{\partial y}{\partial t}
  +\frac{\partial u_1}{\partial z}\cdot\frac{\partial z}{\partial t}
  =-\lambda\frac{\partial u_1}{\partial y}
  +\frac{\partial u_1}{\partial z},
\]
\[
  \frac{\partial u}{\partial x}
  =\frac{\partial u_1}{\partial y}\cdot\frac{\partial y}{\partial x}
  +\frac{\partial u_1}{\partial z}\cdot\frac{\partial z}{\partial x}
  =\frac{\partial u_1}{\partial y},
\]
有
\[
  \frac{\partial u}{\partial t}+\lambda\frac{\partial u}{\partial x}
  =\frac{\partial u_1}{\partial z}=0.
\]
这说明 $u_1(y,z)$ 与 $z$ 无关, 即 $u_1=g(y)$, 亦即 $u=g(x-\lambda t)$, 其中 $g$ 为 $\mathbb R$ 中的任一个可微函数。

\end{proof}
\subsection{高阶偏导数}

设 $f(x)=f(x_1,x_2,\cdots,x_n)$ 在区域 $D\subset\mathbb R^n$ 上具有各个偏导数. 由定义, 它的每个偏导数 $\dfrac{\partial f(x)}{\partial x_i}(i=1,2,\cdots,n)$ 是 $D$ 上的一个 $n$ 元函数. 若它们仍具有各个偏导数, 则称它们的偏导数为 $f(x)$ 的二阶偏导数. 类似地可以定义三阶以及更高阶的偏导数. 二阶及二阶以上的偏导数统称为高阶偏导数。

显然, 若一个函数 $f(x)$ 的各个偏导数都存在 $n$ 个偏导数, 则 $f(x)$ 具有 $n^2$ 个二阶偏导数. 当
\[
  \frac{\partial\left(\dfrac{\partial f(x)}{\partial x_i}\right)}{\partial x_k}
  \quad(1\leqslant i,k\leqslant n)
\]
存在时, 我们将其记为 $\dfrac{\partial^2 f(x)}{\partial x_k\partial x_i}$, $f''_{x_kx_i}(x)$ 或 $f''_{ki}(x)$。

\begin{example}\label{course-example-14-2-7}\quad 设函数
\[
  f(x,y)=
  \begin{cases}
    xy\dfrac{x^2-y^2}{x^2+y^2}, &x^2+y^2\ne0,\\
    0, &x^2+y^2=0,
  \end{cases}
\]
求 $f''_{xy}(0,0)$ 及 $f''_{yx}(0,0)$。

\end{example}
\begin{solve}\quad 由偏导数的定义得
\[
  f'_x(x,y)=
  \begin{cases}
    y\left[\dfrac{x^2-y^2}{x^2+y^2}+\dfrac{4x^2y^2}{(x^2+y^2)^2}\right], &x^2+y^2\ne0,\\
    0, &x^2+y^2=0,
  \end{cases}
\]
\[
  f'_y(x,y)=
  \begin{cases}
    x\left[\dfrac{x^2-y^2}{x^2+y^2}-\dfrac{4x^2y^2}{(x^2+y^2)^2}\right], &x^2+y^2\ne0,\\
    0, &x^2+y^2=0.
  \end{cases}
\]
特别地, 我们有
\[
  f'_x(0,y)=-y,\qquad f'_y(x,0)=x.
\]
由此得
\[
  f''_{yx}(0,0)=-1,\qquad f''_{xy}(0,0)=1.
\]

从上例可知, 一个函数的各个高阶偏导数都存在时, 若改变对自变量分量的求导顺序, 高阶偏导数的值也有可能改变. 但如果我们附加一些条件, 可使高阶偏导数的值对求导顺序无关. 事实上, 我们有下述结果。

\end{solve}
\begin{theorem}\label{course-theorem-14-2-3}\quad 设函数 $f(x)$ 在区域 $D\subset\mathbb R^n$ 上有定义, $x_0\in D$, 且对于 $1\leqslant j<k\leqslant n$, $f''_{kj}(x)$ 与 $f''_{jk}(x)$ 在 $U(x_0,\delta)(\delta>0)$ 内存在, 并且在 $x_0$ 处连续, 则有
\[
  f''_{kj}(x_0)=f''_{jk}(x_0).
\]

\end{theorem}
\begin{proof}\quad 记 $f(x)=f(x_1,x_2,\cdots,x_n)$. 为了使记号简便, 我们不妨设 $j=1,k=2$, 并记 $x=(x,y,x')$, $x_0=(x_0,y_0,x'_0)$ 其中 $x'=(x_3,x_4,\cdots,x_n)$, $x'_0=(x_3^0,x_4^0,\cdots,x_n^0)$. 对于充分小的 $\Delta x,\Delta y$, 观察:
\[
\begin{aligned}
  I(\Delta x,\Delta y)
  &\triangleq
  \frac{f(x_0+\Delta x,y_0+\Delta y,x'_0)-f(x_0+\Delta x,y_0,x'_0)}{\Delta x\Delta y}\\
  &\quad-\frac{f(x_0,y_0+\Delta y,x'_0)-f(x_0)}{\Delta x\Delta y}.
\end{aligned}
\]
记
\[
  g(x)=f(x,y_0+\Delta y,x'_0)-f(x,y_0,x'_0),
  \qquad
  h(y)=f(x_0+\Delta x,y,x'_0)-f(x_0,y,x'_0).
\]
一方面, 由一元函数的微分中值定理得
\[
\begin{aligned}
  I(\Delta x,\Delta y)
  &=\frac1{\Delta x\Delta y}
  \{[f(x_0+\Delta x,y_0+\Delta y,x'_0)-f(x_0+\Delta x,y_0,x'_0)]\\
  &\quad-[f(x_0,y_0+\Delta y,x'_0)-f(x_0)]\}
\end{aligned}
\]
\[
\begin{aligned}
  &=\frac{g(x_0+\Delta x)-g(x_0)}{\Delta x\Delta y}
  =\frac{g'(x_0+\theta_1\Delta x)}{\Delta y}\\
  &=\frac{f'_x(x_0+\theta_1\Delta x,y_0+\Delta y,x'_0)
  -f'_x(x_0+\theta_1\Delta x,y_0,x'_0)}{\Delta y}\\
  &=f''_{yx}(x_0+\theta_1\Delta x,y_0+\theta_2\Delta y,x'_0),
\end{aligned}
\]
其中 $0<\theta_1,\theta_2<1$. 另一方面, 我们有
\[
\begin{aligned}
  I(\Delta x,\Delta y)
  &=\frac1{\Delta x\Delta y}
  \{[f(x_0+\Delta x,y_0+\Delta y,x'_0)-f(x_0,y_0+\Delta y,x'_0)]\\
  &\quad-[f(x_0+\Delta x,y_0,x'_0)-f(x_0)]\}\\
  &=\frac{h(y_0+\Delta y)-h(y_0)}{\Delta x\Delta y}
  =\frac{h'(y_0+\theta_3\Delta y)}{\Delta x}\\
  &=\frac{f'_y(x_0+\Delta x,y_0+\theta_3\Delta y,x'_0)
  -f'_y(x_0,y_0+\theta_3\Delta y,x'_0)}{\Delta x}\\
  &=f''_{xy}(x_0+\theta_4\Delta x,y_0+\theta_3\Delta y,x'_0),
\end{aligned}
\]
其中 $0<\theta_3,\theta_4<1$. 由此得
\[
  f''_{yx}(x_0+\theta_1\Delta x,y_0+\theta_2\Delta y,x'_0)
  =f''_{xy}(x_0+\theta_4\Delta x,y_0+\theta_3\Delta y,x'_0).
\]
在上式中令 $\Delta x\to0,\Delta y\to0$, 并利用 $f''_{yx}(x)$ 与 $f''_{xy}(x)$ 在 $x_0$ 处的连续性即得
\[
  f''_{yx}(x_0)=f''_{xy}(x_0).
\]
证毕。

在今后, 我们碰到的大部分多元函数都具有各阶连续偏导数, 因此, 我们可以不考虑对自变量分量的求导顺序. 例如, 当 $y=f(x)(x\in D\subset\mathbb R^n)$ 具有各个二阶连续偏导数时, $f(x)$ 在 $x_0\in D$ 处具有 $\dfrac{n^2+n}{2}$ 个不同的二阶偏导数。

\end{proof}
\subsection{复合函数的高阶偏导数}

对于复合函数的高阶偏导数, 很难得出一般性的公式, 需要逐次求之. 例如, 设 $u=f(x,y),x=\varphi(s,t),y=\psi(s,t)$, 并假设所有涉及的函数均具有二阶连续偏导数, 则有
\[
  \frac{\partial u}{\partial s}
  =\frac{\partial f}{\partial x}\cdot\frac{\partial\varphi}{\partial s}
  +\frac{\partial f}{\partial y}\cdot\frac{\partial\psi}{\partial s},
\]
\[
\begin{aligned}
  \frac{\partial^2u}{\partial s^2}
  &=\frac{\partial^2 f}{\partial x^2}\left(\frac{\partial\varphi}{\partial s}\right)^2
  +\frac{\partial^2 f}{\partial x\partial y}\cdot
  \frac{\partial\varphi}{\partial s}\cdot\frac{\partial\psi}{\partial s}
  +\frac{\partial f}{\partial x}\cdot\frac{\partial^2\varphi}{\partial s^2}\\
  &\quad+\frac{\partial^2 f}{\partial x\partial y}\cdot
  \frac{\partial\varphi}{\partial s}\cdot\frac{\partial\psi}{\partial s}
  +\frac{\partial^2 f}{\partial y^2}\left(\frac{\partial\psi}{\partial s}\right)^2
  +\frac{\partial f}{\partial y}\cdot\frac{\partial^2\psi}{\partial s^2}\\
  &=\frac{\partial^2 f}{\partial x^2}\left(\frac{\partial\varphi}{\partial s}\right)^2
  +2\frac{\partial^2 f}{\partial x\partial y}\cdot
  \frac{\partial\varphi}{\partial s}\cdot\frac{\partial\psi}{\partial s}
  +\frac{\partial^2 f}{\partial y^2}\left(\frac{\partial\psi}{\partial s}\right)^2\\
  &\quad+\frac{\partial f}{\partial x}\cdot\frac{\partial^2\varphi}{\partial s^2}
  +\frac{\partial f}{\partial y}\cdot\frac{\partial^2\psi}{\partial s^2}.
\end{aligned}
\]

\begin{example}\label{course-example-14-2-8}\quad 设函数 $u=f\left(xy,\dfrac yx\right)$, 试求 $\dfrac{\partial^2u}{\partial x^2}$ 及 $\dfrac{\partial^2u}{\partial x\partial y}$ (假定其二阶偏导数均连续)。

\end{example}
\begin{solve}\quad 由
\[
  \frac{\partial u}{\partial x}=yf'_1-\frac y{x^2}f'_2
\]
得
\[
\begin{aligned}
  \frac{\partial^2u}{\partial x^2}
  &=y^2f''_{11}-\frac y{x^2}f''_{12}
  -\frac y{x^2}f''_{12}+\frac{y^2}{x^4}f''_{22}
  +\frac{2y}{x^3}f'_2\\
  &=y^2f''_{11}-\frac{2y}{x^2}f''_{12}
  +\frac{y^2}{x^4}f''_{22}+\frac{2y}{x^3}f'_2.
\end{aligned}
\]
\[
\begin{aligned}
  \frac{\partial^2u}{\partial x\partial y}
  &=f'_1+xyf''_{11}+\frac yx f''_{12}
  -\frac1{x^2}f'_2-\frac{xy}{x^2}f''_{12}
  -\frac y{x^3}f''_{22}\\
  &=xyf''_{11}-\frac y{x^3}f''_{22}+f'_1-\frac1{x^2}f'_2.
\end{aligned}
\]

\end{solve}
\begin{example}\label{course-example-14-2-9}\quad 设函数 $u=f(r),r=\sqrt{x^2+y^2+z^2}$, 且假设 $u$ 满足拉普拉斯方程
\[
  \frac{\partial^2u}{\partial x^2}
  +\frac{\partial^2u}{\partial y^2}
  +\frac{\partial^2u}{\partial z^2}=0,
\]
求 $f(r)$ 的表达式。

\end{example}
\begin{solve}\quad 由
\[
  \frac{\partial u}{\partial x}=f'(r)\frac xr
\]
得
\[
  \frac{\partial^2u}{\partial x^2}
  =f''(r)\frac{x^2}{r^2}+f'(r)\frac{r-\dfrac{x^2}{r}}{r^2}.
\]
同理,
\[
  \frac{\partial^2u}{\partial y^2}
  =f''(r)\frac{y^2}{r^2}+f'(r)\frac{r-\dfrac{y^2}{r}}{r^2},
\]
\[
  \frac{\partial^2u}{\partial z^2}
  =f''(r)\frac{z^2}{r^2}+f'(r)\frac{r-\dfrac{z^2}{r}}{r^2}.
\]
由于 $u$ 满足拉普拉斯方程, 因此有
\[
\begin{aligned}
  \frac{\partial^2u}{\partial x^2}
  +\frac{\partial^2u}{\partial y^2}
  +\frac{\partial^2u}{\partial z^2}
  &=f''(r)\frac{x^2+y^2+z^2}{r^2}
  +f'(r)\frac{3r^2-(x^2+y^2+z^2)}{r^3}\\
  &=f''(r)+\frac2r f'(r)=0.
\end{aligned}
\]
由此推出
\[
  (r^2f'(r))'=r^2f''(r)+2rf'(r)
  =r^2\left(f''(r)+\frac2r f'(r)\right)=0.
\]
从上式我们进一步推出: 存在常数 $C'$, 使得
\[
  f'(r)=\frac{C'}{r^2}.
\]
对上式两边求不定积分得
\[
  f(r)=-\frac{C'}r+C_1=\frac C{\sqrt{x^2+y^2+z^2}}+C_1,
\]
其中 $C=-C'$, $C_1$ 为常数。

\end{solve}
\subsection{一阶微分的形式不变性与高阶微分}

设 $D\subset\mathbb R^n$ 是区域, 函数 $f(x)$ 在 $x\in D$ 处可微, 即有
\[
  \mathrm df(x)=\sum_{i=1}^n\frac{\partial f(x)}{\partial x_i}\mathrm dx_i.
\]
若对于 $\forall i(1\leqslant i\leqslant n)$, $\dfrac{\partial f}{\partial x_i}$ 仍是可微函数, 则我们称
\[
  \sum_{i=1}^n\left(\sum_{k=1}^n
  \frac{\partial^2 f(x)}{\partial x_k\partial x_i}\mathrm dx_k\right)\mathrm dx_i
\]
为 $f(x)$ 的二阶微分, 并记为 $\mathrm d^2f(x)$, 即
\[
  \mathrm d^2f(x)=\sum_{i=1}^n\sum_{k=1}^n
  \frac{\partial^2 f(x)}{\partial x_k\partial x_i}\mathrm dx_k\mathrm dx_i.
\]
对于 $k\geqslant2(k\in\mathbb N)$, 我们可以归纳地给出多元函数 $f(x)$ 的 $k$ 阶微分
\[
  \mathrm d^kf(x)=\mathrm d(\mathrm d^{k-1}f(x)).
\]
若形式地记 $\mathrm df(x)=\left(\sum_{i=1}^n\mathrm dx_i\frac{\partial}{\partial x_i}\right)f(x)$, 当 $f(x)$ 的每个 $k$ 阶偏导数都连续时, 我们可以将 $\mathrm d^kf(x)$ 记为
\[
  \mathrm d^kf(x)=\left(\sum_{i=1}^n\mathrm dx_i\frac{\partial}{\partial x_i}\right)^kf(x).
\]
这样的记号对于高阶微分 (二阶及二阶以上的微分) 来说比较容易记忆. 特别地, 对于一个二元函数 $f(x,y)$, 若它各个 $k$ 阶偏导数都存在且连续, 则我们有
\[
\begin{aligned}
  \mathrm d^kf(x,y)
  &=\left(\mathrm dx\frac{\partial}{\partial x}
  +\mathrm dy\frac{\partial}{\partial y}\right)^kf(x,y)\\
  &=\sum_{j=0}^k C_k^j
  \frac{\partial^k f(x,y)}{\partial x^{k-j}\partial y^j}
  \mathrm dx^{k-j}\mathrm dy^j,
\end{aligned}
\]
其中 $C_k^j=\dfrac{k!}{j!(k-j)!}(j=1,2,\cdots,k),C_k^0=1$。

现设 $f(u)=f(u_1,u_2,\cdots,u_m)$ 在区域 $D\subset\mathbb R^m$ 上可微, 则 $f(u)$ 的微分 $\mathrm df(u)=f'(u)\mathrm du$. 注意到此时 $u$ 是自变量。

现在设 $u=(u_1(x),u_2(x),\cdots,u_m(x))^\mathrm T$ 是区域 $\Omega\subset\mathbb R^n$ 上的一个 $n$ 元可微向量函数, 并且 $u(\Omega)\subset D$, 则复合函数 $f(u(x))$ 在 $\Omega$ 处可微, 从而有
\[
  \mathrm df(u(x))=f'(u(x))u'(x)\mathrm dx.
\]
由于 $\mathrm du=(u'_1(x)\mathrm dx,u'_2(x)\mathrm dx,\cdots,u'_m(x)\mathrm dx)^\mathrm T=u'(x)\mathrm dx$, 因此, 当 $u$ 是中间变量时, 我们仍然有
\[
  \mathrm df(u)=f'(u(x))u'(x)\mathrm dx=f'(u)\mathrm du.
\]
以上讨论说明了对于多元函数的一阶微分仍具有形式不变性. 在一元微积分中, 我们已经知道了对于二阶以上的微分不再具有形式不变性, 因此多元函数情形也不可能具有高阶微分的形式不变性。

\begin{example}\label{course-example-14-2-10}\quad 设函数 $f\left(xy,\dfrac yx\right)$ 具有二阶连续偏导数, 求它的所有二阶偏导数。

\end{example}
\begin{solve}\quad 在例~\ref{course-example-14-2-8}  中我们曾经直接对该函数求过二阶偏函数, 现在我们利用 $f$ 的二阶微分来求之. 由
\[
  \mathrm df=f'_1\mathrm d(xy)+f'_2\mathrm d\left(\frac yx\right)
\]
得
\[
\begin{aligned}
  \mathrm d^2f
  &=\left(f''_{11}\mathrm d(xy)+f''_{12}\mathrm d\left(\frac yx\right)\right)\mathrm d(xy)
  +f'_1\mathrm d^2(xy)\\
  &\quad+\left(f''_{21}\mathrm d(xy)+f''_{22}\mathrm d\left(\frac yx\right)\right)
  \mathrm d\left(\frac yx\right)+f'_2\mathrm d^2\left(\frac yx\right)\\
  &=\left(y^2f''_{11}-\frac{2y^2}{x^2}f''_{12}
  +\frac{y^2}{x^4}f''_{22}+\frac{2y}{x^3}f'_2\right)\mathrm dx^2\\
  &\quad+2\left(xyf''_{11}-\frac y{x^3}f''_{22}
  +f'_1-\frac1{x^2}f'_2\right)\mathrm dx\mathrm dy\\
  &\quad+\left(x^2f''_{11}+2f''_{12}+\frac1{x^2}f''_{22}\right)\mathrm dy^2.
\end{aligned}
\]
因此
\[
  \frac{\partial^2 f}{\partial x^2}
  =y^2f''_{11}-\frac{2y^2}{x^2}f''_{12}
  +\frac{y^2}{x^4}f''_{22}+\frac{2y}{x^3}f'_2,
\]
\[
  \frac{\partial^2 f}{\partial x\partial y}
  =xyf''_{11}-\frac y{x^3}f''_{22}
  +f'_1-\frac1{x^2}f'_2,
\]
\[
  \frac{\partial^2 f}{\partial y^2}
  =x^2f''_{11}+2f''_{12}+\frac1{x^2}f''_{22}.
\]

\end{solve}


\section{泰勒公式}\label{course-14-3}
本节用于核对多元极限、隐函数、驻点与约束极值题中的可微性和非退化条件。仅有偏导数存在不能替代全微分存在；用必要条件列出候选点以后，还要比较边界并验证极值是否取得。

在利用高阶导数研究函数时, 泰勒公式是一个强有力的工具, 对于多元函数我们有以下定理。

\begin{theorem}[泰勒公式]\label{course-theorem-14-3-1}\quad 设函数 $f(x)$ 在 $x_0=(x_1^0,x_2^0,\cdots,x_n^0)\in\mathbb R^n$ 的邻域 $U(x_0,\delta_0)(\delta_0>0)$ 内具有 $K+1$ 阶连续偏导数, 则对于 $\forall x_0+h=(x_1^0+h_1,x_2^0+h_2,\cdots,x_n^0+h_n)\in U(x_0,\delta_0)$, 有
\begin{equation}
\begin{aligned}
  f(x_0+h)
  &=f(x_0)+\sum_{k=1}^K\frac1{k!}
  \left(\sum_{i=1}^n h_i\frac{\partial}{\partial x_i}\right)^k f(x_0)\\
  &\quad+\frac1{(K+1)!}
  \left(\sum_{i=1}^n h_i\frac{\partial}{\partial x_i}\right)^{K+1}
  f(x_0+\theta h),
\end{aligned}
\label{course-eq-14-3-1}
\end{equation}
其中 $0<\theta<1$。

\end{theorem}
\begin{proof}\quad 设 $h=(h_1,h_2,\cdots,h_n)$ 满足 $x_0+h\in U(x_0,\delta_0)$. 构造一元函数
\[
  \varphi(t)=f(x_0+th),\qquad t\in[0,1].
\]
由复合函数求导公式知, $\varphi(t)$ 具有 $K+1$ 阶连续导数. 因此, 由一元函数的泰勒公式有
\begin{equation}
  \varphi(t)=\varphi(0)+\sum_{k=1}^K
  \frac{\varphi^{(k)}(0)}{k!}t^k
  +\frac{\varphi^{(K+1)}(\theta t)}{(K+1)!}t^{K+1},
  \label{course-eq-14-3-2}
\end{equation}
其中 $0<\theta<1$. 注意到
\[
  \varphi'(t)=\sum_{i=1}^n\frac{\partial f(x_0+th)}{\partial x_i}h_i
  =\left(\sum_{i=1}^n h_i\frac{\partial}{\partial x_i}\right)f(x_0+th),
\]
\[
  \varphi''(t)=\sum_{j=1}^n\sum_{i=1}^n
  \frac{\partial^2 f(x_0+th)}{\partial x_i\partial x_j}h_ih_j
  =\left(\sum_{i=1}^n h_i\frac{\partial}{\partial x_i}\right)^2 f(x_0+th),
\]
一般地, 对于 $k=1,2,\cdots,K+1$, 有
\[
  \varphi^{(k)}(t)
  =\left(\sum_{i=1}^n h_i\frac{\partial}{\partial x_i}\right)^k f(x_0+th).
\]
因此, 将 $\varphi^{(k)}(t)(k=1,2,\cdots,K+1)$ 的表达式代入式~\eqref{course-eq-14-3-2}, 然后令 $t=1$ 即得所证. 证毕。

\end{proof}
\begin{remark}\quad 当 $f(x)$ 在区域 $D$ 内具有 $K+1$ 阶连续偏导数时, 定理~\ref{course-theorem-14-3-1}  中的余项
\[
  R_{K+1}(h)=\frac1{(K+1)!}
  \left(\sum_{i=1}^n h_i\frac{\partial}{\partial x_i}\right)^{K+1}f(x_0+\theta h)
  \quad(0<\theta<1)
\]
称为拉格朗日余项, 并称式~\eqref{course-eq-14-3-1}  为 $f(x)$ 带拉格朗日余项的泰勒公式。

\end{remark}
\noindent\keyterm{推论 1}\quad 设函数 $f(x)\in C^K(U(x_0,\delta_0))(\delta_0>0)$, 即 $f(x)$ 在邻域 $U(x_0,\delta_0)$ 内具有 $K(K\geqslant1)$ 阶连续偏导数, $x_0+h\in U(x_0,\delta_0)$, 则
\begin{equation}
  f(x_0+h)=f(x_0)+\sum_{k=1}^K\frac1{k!}
  \left(\sum_{i=1}^n h_i\frac{\partial}{\partial x_i}\right)^k f(x_0)
  +o(|h|^K)
  \quad(|h|\to0).
  \label{course-eq-14-3-3}
\end{equation}

\begin{proof}\quad 由定理~\ref{course-theorem-14-3-1}  我们有
\[
\begin{aligned}
  f(x_0+h)
  &=f(x_0)+\sum_{k=1}^K\frac1{k!}
  \left(\sum_{i=1}^n h_i\frac{\partial}{\partial x_i}\right)^kf(x_0)\\
  &\quad+\frac1{K!}
  \left(\sum_{i=1}^n h_i\frac{\partial}{\partial x_i}\right)^K
  (f(x_0+\theta h)-f(x_0)),
\end{aligned}
\]
其中 $0<\theta<1$。

记
\[
  R_K(h)=\frac1{K!}
  \left(\sum_{i=1}^n h_i\frac{\partial}{\partial x_i}\right)^K
  (f(x_0+\theta h)-f(x_0)).
\]
由 $f(x)$ 在 $U(x_0,\delta_0)$ 内具有 $K(K\geqslant1)$ 阶连续偏导数, 知
\[
  \lim_{|h|\to0}\frac{R_K(h)}{|h|^K}
  =\lim_{|h|\to0}\frac1{K!}
  \left(\sum_{i=1}^n\frac{h_i}{|h|}\frac{\partial}{\partial x_i}\right)^K
  (f(x_0+\theta h)-f(x_0))=0.
\]
上述极限为零是因为当 $|h|\to0$ 时, 和式中的每一项的系数都是有界的, 而 $f(x)$ 在 $x_0+\theta h$ 处的每个 $K$ 阶偏导数由于其连续性都趋于 $f(x)$ 在 $x_0$ 相应的 $K$ 阶偏导数. 证毕。

\end{proof}
\begin{remark}\quad 我们将推论 1 中的余项 $o(|h|^K)$ 称为皮亚诺余项, 并称式~\eqref{course-eq-14-3-3}  为 $f(x)$ 带皮亚诺余项的泰勒公式. 由推论 1 我们可知, 若 $f(x)$ 在 $x_0$ 处能展成
\[
  f(x_0+h)=P_K(h_1,\cdots,h_n)+o(|h|^K)
  \quad(|h|\to0),
\]
其中 $P_K$ 是 $n$ 元 $K$ 次多项式, 上式必定是 $f(x)$ 在 $U(x_0,\delta_0)$ 内的泰勒公式。

现设 $f(x)$ 在区域 $D$ 内具有 $K+1$ 阶连续偏导数, 从定理~\ref{course-theorem-14-3-1}  的证明中可以看出, 若要求 $f(x)(x\in D)$ 在 $x_0\in D$ 处的泰勒公式, 我们不必要求 $x$ 在 $x_0$ 的某个邻域内, 而只要求连接 $x_0$ 和 $x$ 的线段 $x_0+t(x-x_0)(t\in[0,1])$ 落在 $D$ 内即可。

在定理~\ref{course-theorem-14-3-1}  中取 $K=0$, 我们有以下多元函数的拉格朗日微分中值定理。

\end{remark}
\noindent\keyterm{推论 2}\quad 设 $f(x)$ 在区域 $D\subset\mathbb R^n$ 内具有连续偏导数, 且对于 $\forall t\in[0,1],x_0+t(x-x_0)\in D$, 则有
\[
\begin{aligned}
  f(x)-f(x_0)
  &=\sum_{i=1}^n
  \frac{\partial f(x_0+\theta(x-x_0))}{\partial x_i}(x_i-x_i^0)\\
  &=f'(x_0+\theta(x-x_0))\cdot(x-x_0)^\mathrm T,
\end{aligned}
\]
其中 $0<\theta<1$。

由上述推论 2, 我们容易推出下述结论。

\noindent\keyterm{推论 3}\quad 设函数 $f(x)$ 在区域 $D\subset\mathbb R^n$ 内的各个偏导数均为 0, 则 $f(x)$ 在 $D$ 内为常数函数。

请读者给出上述推论的证明。

以后我们还常常要用到泰勒公式中 $K=1$ 的情形. 在定理~\ref{course-theorem-14-3-1}  中, 令 $K=1$, 我们有
\[
  f(x_0+h)=f(x_0)+f'(x_0)\cdot h^\mathrm T
  +\frac12 h\cdot H_f(x_0)h^\mathrm T+o(|h|^2)
  \quad(|h|\to0),
\]
其中
\[
  H_f(x_0)=
  \begin{pmatrix}
    \dfrac{\partial^2 f(x_0)}{\partial x_1^2} &
    \dfrac{\partial^2 f(x_0)}{\partial x_1\partial x_2} &
    \cdots &
    \dfrac{\partial^2 f(x_0)}{\partial x_1\partial x_n}\\
    \dfrac{\partial^2 f(x_0)}{\partial x_2\partial x_1} &
    \dfrac{\partial^2 f(x_0)}{\partial x_2^2} &
    \cdots &
    \dfrac{\partial^2 f(x_0)}{\partial x_2\partial x_n}\\
    \vdots & \vdots & & \vdots\\
    \dfrac{\partial^2 f(x_0)}{\partial x_n\partial x_1} &
    \dfrac{\partial^2 f(x_0)}{\partial x_n\partial x_2} &
    \cdots &
    \dfrac{\partial^2 f(x_0)}{\partial x_n^2}
  \end{pmatrix}
\]
称为 $f(x)$ 在 $x_0$ 处的海色 (Hessi) 矩阵。

\begin{example}\label{course-example-14-3-1}\quad 求函数 $f(x,y)=\dfrac{x-y}{x+y}$ 在 $(1,0)$ 附近带皮亚诺余项的泰勒公式 (直到二次项)。

\end{example}
\begin{solve}\quad 对于这类函数, 一般可直接求偏导, 然后写出泰勒公式, 但我们下面利用一元函数已知的泰勒公式来求之. 由于
\[
\begin{aligned}
  f(x,y)
  &=\frac{(x-1)+1-y}{1+[(x-1)+y]}\\
  &=[(x-1)+1-y]\{1-(x-1)-y+[(x-1)+y]^2
  +o((x-1)^2+y^2)\}
\end{aligned}
\]
$((x-1)^2+y^2\to0)$, 将上式整理得
\[
  f(x,y)=1-2y+2(x-1)y+2y^2+o((x-1)^2+y^2),
  \quad ((x-1)^2+y^2\to0).
\]

\end{solve}
\begin{example}\label{course-example-14-3-2}\quad 求函数 $f(x,y)=xe^{x+y}$ 在原点 $(0,0)$ 处的所有四阶偏导数。

\end{example}
\begin{solve}\quad 由 $e^{x+y}$ 在 $(0,0)$ 处的泰勒公式有
\[
  f(x,y)=xe^{x+y}
  =x\left[1+\sum_{k=1}^K\frac{(x+y)^k}{k!}
  +o\left((\sqrt{x^2+y^2})^K\right)\right]
  \quad(\sqrt{x^2+y^2}\to0).
\]
令 $K=3$, 得
\[
\begin{aligned}
  f(x,y)
  &=x+x^2+xy+\frac12(x^3+2x^2y+xy^2)\\
  &\quad+\frac16(x^4+3x^3y+3x^2y^2+xy^3)
  +o\left((\sqrt{x^2+y^2})^4\right)
  \quad(\sqrt{x^2+y^2}\to0),
\end{aligned}
\]
于是有
\[
  \frac{C_4^4}{4!}\cdot\frac{\partial^4 f(0,0)}{\partial x^4}=\frac16,\quad
  \frac{C_4^3}{4!}\cdot\frac{\partial^4 f(0,0)}{\partial x^3\partial y}=\frac12,\quad
  \frac{C_4^2}{4!}\cdot\frac{\partial^4 f(0,0)}{\partial x^2\partial y^2}=\frac12,
\]
\[
  \frac{C_4^1}{4!}\cdot\frac{\partial^4 f(0,0)}{\partial x\partial y^3}=\frac16,\quad
  \frac{C_4^0}{4!}\cdot\frac{\partial^4 f(0,0)}{\partial y^4}=0,
\]
从而得
\[
  \frac{\partial^4 f(0,0)}{\partial x^4}=4,\quad
  \frac{\partial^4 f(0,0)}{\partial x^3\partial y}=3,\quad
  \frac{\partial^4 f(0,0)}{\partial x^2\partial y^2}=2,
\]
\[
  \frac{\partial^4 f(0,0)}{\partial x\partial y^3}=1,\quad
  \frac{\partial^4 f(0,0)}{\partial y^4}=0.
\]

\end{solve}


\section{隐函数存在定理}\label{course-14-4}
本节用于核对多元极限、隐函数、驻点与约束极值题中的可微性和非退化条件。仅有偏导数存在不能替代全微分存在；用必要条件列出候选点以后，还要比较边界并验证极值是否取得。

\subsection{单个方程的情形}

在讨论一般的隐函数存在问题之前, 我们先来讨论一下由一个二元方程确定隐函数的简单情形. 设 $F(x,y)=0$ 是区域 $D\subset\mathbb R^2$ 上的一个二元方程, 并且满足 $F(x_0,y_0)=0$, 其中 $(x_0,y_0)\in D$. 我们要研究的问题是: 何时方程 $F(x,y)=0$ 在点 $(x_0,y_0)$ 附近能唯一确定一个函数 $y=f(x)$, 使得 $f(x_0)=y_0$, 且 $F(x,f(x))=0$ 对 $x\in(x_0-\delta,x_0+\delta)(\delta>0)$ 成立? 我们称这样定义的函数 $y=f(x)$ 为由方程 $F(x,y)=0$ 所确定的隐函数. 显然这个问题等价于: 何时 $z=F(x,y)$ 的零点集合在 $(x_0,y_0)$ 附近是 $Oxy$ 平面内的一条曲线, 而且这条曲线是 $y=f(x)$ 的图像?

我们先来考查一个例子. $\mathbb R^2$ 中的单位圆周满足的方程为 $x^2+y^2=1$, 我们可将其写成 $x^2+y^2-1=0$. 现考虑函数 $z=F(x,y)=x^2+y^2-1$. 先看一个简单情况, 令 $(x_0,y_0)=(0,1)$, 我们来考查一下为什么 $F(x,y)=0$ 在该点的附近能确定一个隐函数 $y=\sqrt{1-x^2}$. 首先 $(x_0,y_0)=(0,1)$ 是 $F(x,y)=x^2+y^2-1$ 的零点, 并且平面 $x=0$ 与曲面 $z=F(x,y)$ 的交线是
\[
  \begin{cases}
    x=0,\\
    z=y^2-1.
  \end{cases}
\]
对于该曲线容易看出, 对于 $0<\delta_0<1$, 当 $y$ 从 $1-\delta_0$ 连续变化到 $1+\delta_0$ 时, $z$ 严格单调上升并且连续地从负变正. 因此, 对于充分小的 $\delta>0$, 对 $(-\delta,\delta)$ 中的每个 $x'$, 曲线
\[
  \begin{cases}
    x=x',\\
    z=y^2-1+x'^2
  \end{cases}
\]
都有此性质. 显然, 当具有上述性质时, $F(x,y)$ 在点 $(x_0,y_0)$ 附近的零点集合刚好是一个函数 $y=f(x)$ 的图像. 而什么时候能保证一个函数 $z=F(x,y)$ 在固定 $x$ 时关于 $y$ 是一个严格单调函数呢? 从多元函数偏导数的几何意义容易看出, 一个充分条件是 $F'_y(x,y)\ne0$。

对于满足 $F'_y(x,y)=0$ 的点 $(x_0,y_0)$, 读者容易发现此时未必有隐函数 $y=y(x)$ 存在, 如 $F(x,y)=x^2+y^2-1$ 在 $(1,0)$ 处。

下面的定理给出了隐函数存在的充分条件。

\begin{theorem}[隐函数存在定理]\label{course-theorem-14-4-1}\quad 设二元函数 $F(x,y)$ 在 $U((x_0,y_0),\delta)(\delta>0)$ 内满足以下条件:

(1) $F(x_0,y_0)=0$;

(2) $F(x,y),F'_y(x,y)$ 在 $U((x_0,y_0),\delta)$ 内连续;

(3) $F'_y(x_0,y_0)\ne0$,

则 $\exists\delta_0>0(0<\delta_0<\delta)$, 使得在 $U(x_0,\delta_0)$ 内存在唯一满足下述条件的连续函数 $y=f(x)$:

(a) $y_0=f(x_0)$;

(b) $F(x,f(x))=0,\forall x\in U(x_0,\delta_0)$;

(c) 如果 $F'_x(x,y)$ 在 $U((x_0,y_0),\delta)$ 内连续, 则 $f(x)$ 在 $U(x_0,\delta_0)$ 存在连续导数, 并且有
\[
  f'(x)=-\frac{F'_x(x,f(x))}{F'_y(x,f(x))}.
\]

\end{theorem}
\begin{proof}\quad 不妨设 $F'_y(x_0,y_0)>0$. 由 $F'_y(x,y)$ 的连续性及极限的保号性可知, 存在 $0<\delta_1,\delta_2<\delta$, 使得对于 $\forall(x,y)\in U(x_0,\delta_1)\times U(y_0,\delta_2)$, 有
\[
  F'_y(x,y)>0.
\]
特别地, 若固定 $x=x_0$, 则对于 $\forall y\in U(y_0,\delta_2)$, 有 $F'_y(x_0,y)>0$. 因此 $F(x_0,y)$ 在 $U(y_0,\delta_2)$ 内是 $y$ 的严格上升函数. 注意到 $F(x_0,y_0)=0$, 我们可知 $F(x_0,y_0-\delta_2)<0,F(x_0,y_0+\delta_2)>0$. 再由 $F(x,y)$ 在 $(x_0,y_0-\delta_2)$ 和 $(x_0,y_0+\delta_2)$ 处的连续性, 并利用极限的保号性, 我们可以找到 $\delta_0(0<\delta_0<\delta_1)$, 使得当 $x\in U(x_0,\delta_0)$ 时, 有
\[
  F(x,y_0-\delta_2)<0,\qquad F(x,y_0+\delta_2)>0.
\]
对任意给定的 $\tilde x\in U(x_0,\delta_0)$, 由 $F'_y(\tilde x,y)>0$ 推知, 当 $y\in U(y_0,\delta_2)$ 时, $F(\tilde x,y)$ 连续地从负数 $F(\tilde x,y_0-\delta_2)$ 严格上升到正数 $F(\tilde x,y_0+\delta_2)$ . 因此存在唯一的 $\tilde y\in U(y_0,\delta_2)$, 使得
\[
  F(\tilde x,\tilde y)=0.
\]
这说明了对于 $\forall\tilde x\in U(x_0,\delta_0)$, 有唯一确定的 $\tilde y$ 与之对应. 由函数的定义, 若令 $\tilde y=f(\tilde x)$, 则函数 $\tilde y=f(\tilde x)$ 在 $U(x_0,\delta_0)$ 内有定义, 并且满足 (a),(b). 该函数的唯一性由 $F(x,y)$ 在 $U(x_0,\delta_0)\times U(y_0,\delta_2)$ 内关于 $y$ 严格上升所保证。

下面证 $y=f(x)$ 在 $U(x_0,\delta_0)$ 内连续. 任意取定 $\bar x\in U(x_0,\delta_0)$, 记 $\bar y=f(\bar x)$. 对于充分小的 $\varepsilon>0$, 则有
\[
  F(\bar x,\bar y-\varepsilon)<0,\qquad F(\bar x,\bar y+\varepsilon)>0.
\]
由 $F(x,y)$ 在 $(\bar x,\bar y-\varepsilon)$ 及 $(\bar x,\bar y+\varepsilon)$ 处的连续性, 必存在充分小的 $\delta'(0<\delta'<\delta_0)$, 使得当 $x\in U(\bar x,\delta')\subset U(x_0,\delta_0)$ 时, 有
\[
  F(x,\bar y-\varepsilon)<0,\qquad F(x,\bar y+\varepsilon)>0.
\]
从上面关于隐函数存在性的证明过程知, 当 $x\in U(\bar x,\delta')$ 时, 必有 $f(x)\in(\bar y-\varepsilon,\bar y+\varepsilon)$, 即
\[
  |f(x)-f(\bar x)|<\varepsilon.
\]
这说明 $f(x)$ 在 $U(x_0,\delta_0)$ 内连续。

最后, 设 $F'_x(x,y)$ 在 $U((x_0,y_0),\delta)$ 内存在且连续, 我们证明 $y=f(x)$ 在 $U(x_0,\delta_0)$ 内具有连续导数. 任取 $\bar x\in U(x_0,\delta_0)$, 再取 $\Delta x$ 充分小, 使得 $\bar x+\Delta x\in U(x_0,\delta_0)$. 记
\[
  \bar y=f(\bar x),\qquad
  \Delta y=f(\bar x+\Delta x)-f(\bar x).
\]
由多元函数拉格朗日微分中值定理, $\exists0<\theta<1$, 使得下式成立:
\[
\begin{aligned}
  0&=F(\bar x+\Delta x,\bar y+\Delta y)-F(\bar x,\bar y)\\
  &=F'_x(\bar x+\theta\Delta x,\bar y+\theta\Delta y)\Delta x
  +F'_y(\bar x+\theta\Delta x,\bar y+\theta\Delta y)\Delta y.
\end{aligned}
\]
再由 $F'_y(x,y)\ne0((x,y)\in U((x_0,y_0),\delta))$, 得
\[
  \frac{\Delta y}{\Delta x}
  =-\frac{F'_x(\bar x+\theta\Delta x,\bar y+\theta\Delta y)}
  {F'_y(\bar x+\theta\Delta x,\bar y+\theta\Delta y)}.
\]
令 $\Delta x\to0$, 由 $F'_x(x,y)$ 及 $F'_y(x,y)$ 的连续性得
\[
  f'(\bar x)=-\frac{F'_x(\bar x,f(\bar x))}{F'_y(\bar x,f(\bar x))}.
\]
从上式可以看出 $f'(x)$ 在 $U(x_0,\delta_0)$ 内连续. 证毕。

关于隐函数存在定理, 我们要注意以下几点:

\end{proof}
\noindent\keyterm{注 1}\quad 隐函数存在定理是一个具有重要理论价值的定理, 这在许多后续课程中读者将会有所体验. 值得指出的是, 在定理的条件下, $F(x,y)=0$ 所确定的隐函数是局部存在的, 并且定理的结论只给出了存在性. 在很多情况下, 我们未必能找到解析式 $y=f(x)$ 来表示这个隐函数. 例如, 对于开普勒 (Kepler) 方程
\[
  y-x-\varepsilon\sin y=0\quad(0<\varepsilon<1),
\]
容易验证它在 $(-\infty,+\infty)$ 上能确定函数 $y=f(x)$, 但人们却无法写出 $y=f(x)$ 的解析式。

\noindent\keyterm{注 2}\quad 定理~\ref{course-theorem-14-4-1}  只是给了隐函数存在的一个充分条件, 若对定理证明过程加以分析, 读者容易举出例子来说明, 当 $F(x,y)$ 不满足定理的条件时, 有时也可以唯一确定一个隐函数 (见本章习题)。

在隐函数存在定理中, 将 $F(x,y)$ 的 $x\in U(x_0,\delta)$ 换成 $\mathbb R^n$ 中的点 $x\in U(x_0,\delta)(\delta>0)$ 时, 相应的隐函数存在和唯一性的结论不仅成立, 而且其证明与定理~\ref{course-theorem-14-4-1}  的证明相似. 下面不加证明地给出此结论。

\begin{theorem}[隐函数存在定理]\label{course-theorem-14-4-2}\quad 记 $x=(x_1,x_2,\cdots,x_n),x_0=(x_1^0,x_2^0,\cdots,x_n^0)\in\mathbb R^n$, 假设函数 $F(x,y)=F(x_1,x_2,\cdots,x_n,y)$ 在 $U(x_0,\delta)\times U(y_0,\delta)(\delta>0)$ 内有定义, 并且满足

(1) $F(x_0,y_0)=0$;

(2) $F(x,y)$ 和 $F'_y(x,y)$ 在 $U(x_0,\delta)\times U(y_0,\delta)$ 内连续;

(3) $F'_y(x_0,y_0)\ne0$,

则 $\exists\delta_0(0<\delta_0<\delta)$, 使得在 $U(x_0,\delta_0)$ 内存在唯一满足下述条件的连续函数 $y=f(x)$:

(a) $y_0=f(x_0)$;

(b) 对于 $\forall x\in U(x_0,\delta_0),F(x,f(x))=0$;

(c) 如果 $F(x,y)$ 在 $U(x_0,\delta)\times U(y_0,\delta)$ 内存在各个连续偏导数, 那么 $y=f(x)$ 在 $U(x_0,\delta_0)$ 具有各个连续偏导数, 并且对于 $i=1,2,\cdots,n$ 及 $\forall x\in U(x_0,\delta_0)$, 有
\[
  \frac{\partial f(x)}{\partial x_i}
  =-\frac{F'_{x_i}(x_1,x_2,\cdots,x_n,y)}
  {F'_y(x_1,x_2,\cdots,x_n,y)},\qquad
  \text{其中 }y=f(x).
\]

\end{theorem}
\begin{remark}\quad 在定理~\ref{course-theorem-14-4-1}  和定理~\ref{course-theorem-14-4-2}  中, 若 $F(x,y)$ (或 $F(x,y)$) 除了满足相应定理的条件外, 再假设它们具有各个高阶偏导数, 从隐函数的一阶导数或偏导数的表达式直接看出, 所确定的隐函数也具有相应的高阶导数或高阶偏导数。

\end{remark}
\begin{example}\label{course-example-14-4-1}\quad 设 $x=x(y,z),y=y(x,z),z=z(x,y)$ 都是由方程 $F(x,y,z)=0$ 确定的隐函数, 并且它们具有连续偏导数. 证明:
\[
  (1)\quad \frac{\partial x}{\partial y}\cdot
  \frac{\partial y}{\partial z}\cdot
  \frac{\partial z}{\partial x}=-1;
\]
\[
  (2)\quad \frac{\partial x}{\partial y}\cdot
  \frac{\partial x}{\partial z}\cdot
  \frac{\partial y}{\partial x}\cdot
  \frac{\partial y}{\partial z}\cdot
  \frac{\partial z}{\partial x}\cdot
  \frac{\partial z}{\partial y}=1.
\]

\end{example}
\begin{proof}\quad (1) 因为
\[
  \frac{\partial x}{\partial y}=-\frac{F'_y(x,y,z)}{F'_x(x,y,z)},\quad
  \frac{\partial y}{\partial z}=-\frac{F'_z(x,y,z)}{F'_y(x,y,z)},\quad
  \frac{\partial z}{\partial x}=-\frac{F'_x(x,y,z)}{F'_z(x,y,z)},
\]
所以
\[
  \frac{\partial x}{\partial y}\cdot
  \frac{\partial y}{\partial z}\cdot
  \frac{\partial z}{\partial x}=-1.
\]

(2) 由 (1), 又因为
\[
  \frac{\partial x}{\partial z}=-\frac{F'_z(x,y,z)}{F'_x(x,y,z)},\quad
  \frac{\partial y}{\partial x}=-\frac{F'_x(x,y,z)}{F'_y(x,y,z)},\quad
  \frac{\partial z}{\partial y}=-\frac{F'_y(x,y,z)}{F'_z(x,y,z)},
\]
所以
\[
  \frac{\partial x}{\partial y}\cdot
  \frac{\partial x}{\partial z}\cdot
  \frac{\partial y}{\partial x}\cdot
  \frac{\partial y}{\partial z}\cdot
  \frac{\partial z}{\partial x}\cdot
  \frac{\partial z}{\partial y}=1.
\]

\end{proof}
\begin{remark}\quad 在一元微分学中, 当 $y$ 是 $x$ 的函数时, $\dfrac{\mathrm dy}{\mathrm dx}$ 可以看成是 $\mathrm dy$ 与 $\mathrm dx$ 的商. 上例说明, 在多元微分学中, 当 $z=f(x,y)$ 可偏导时, $\dfrac{\partial z}{\partial x}$ 则是一个整体记号。

\end{remark}
\begin{example}\label{course-example-14-4-2}\quad 证明在 $(0,0,0)$ 的邻域内方程
\begin{equation}
  -2x+y-x^2+y^2+z+\sin z=0
  \label{course-eq-14-4-1}
\end{equation}
确定隐函数 $z=f(x,y)$, 并求 $f(x,y)$ 在 $(0,0)$ 处带皮亚诺余项的泰勒公式 (直到二次项)。

\end{example}
\begin{solve}\quad 令 $F(x,y,z)=-2x+y-x^2+y^2+z+\sin z$, 则 $F(x,y,z)$ 具有各阶连续偏导数且满足 $F(0,0,0)=0$ 和 $F'_z(0,0,0)=1+\cos0=2\ne0$. 因此 $F(x,y,z)=0$ 在 $(0,0,0)$ 的某一邻域内唯一确定隐函数 $z=f(x,y)$, 并且 $f(x,y)$ 在该邻域内具有各阶连续偏导数。

注意到 $f(0,0)=0$, 因此 $z=f(x,y)$ 的泰勒公式具有下述形式:
\begin{equation}
  f(x,y)=a_1x+a_2y+b_{11}x^2+b_{12}xy+b_{22}y^2+o(x^2+y^2)
  \label{course-eq-14-4-2}
\end{equation}
\[
  (\sqrt{x^2+y^2}\to0),
\]
其中 $a_i,b_{ij}(i,j=1,2)$ 为常数. 由于
\begin{equation}\label{course-eq-14-4-3}
  \sin z=z+o(|z|^2)
  =a_1x+a_2y+b_{11}x^2+b_{12}xy+b_{22}y^2+o(x^2+y^2)
\end{equation}
\[
  (\sqrt{x^2+y^2}\to0),
\]
将式~\eqref{course-eq-14-4-2},\eqref{course-eq-14-4-3}  代入式~\eqref{course-eq-14-4-1}  得
\[
  -2x+y-x^2+y^2+2a_1x+2a_2y+2b_{11}x^2+2b_{12}xy+2b_{22}y^2+o(x^2+y^2)=0.
\]
由此推出
\[
  (2a_1-2)x+(2a_2+1)y+(2b_{11}-1)x^2
  +2b_{12}xy+(2b_{22}+1)y^2+o(x^2+y^2)=0,
\]
从而有
\[
  a_1=1,\qquad a_2=-\frac12,\qquad b_{11}=\frac12,\qquad
  b_{12}=0,\qquad b_{22}=-\frac12.
\]
这样我们就求出了如下 $f(x,y)$ 的带皮亚诺余项的泰勒公式:
\[
  f(x,y)=x-\frac12y+\frac12x^2-\frac12y^2+o(x^2+y^2)
  \quad(\sqrt{x^2+y^2}\to0).
\]

\end{solve}
\subsection{方程组的情形}

我们下面来考查两个方程所组成的方程组
\[
  \begin{cases}
    F(x,u,v)=0,\\
    G(x,u,v)=0
  \end{cases}
\]
何时能在 $(x_0,u_0,v_0)$ 的邻域内确定 $u,v$ 是 $x$ 的函数. 对一个方程 $F(x,y)=0$ 的情形, $F'_y(x_0,y_0)\ne0$ 起着重要作用. 因此, 当我们考虑向量函数
\[
  \begin{pmatrix}
    z_1\\
    z_2
  \end{pmatrix}
  =
  \begin{pmatrix}
    F(x,u,v)\\
    G(x,u,v)
  \end{pmatrix}
\]
时, 若将 $x$ 看做常量, 它在 $(x_0,u_0,v_0)$ 处关于 $(u,v)$ 的导数将起着隐函数存在定理中 $F'_y(x_0,y_0)$ 的作用. 注意到此时该导数是
\[
  \left.
  \begin{pmatrix}
    F'_u & F'_v\\
    G'_u & G'_v
  \end{pmatrix}
  \right|_{(x_0,u_0,v_0)}.
\]
由于要确定两个隐函数, 因此该雅可比矩阵应该是非退化的, 即
\[
  \left.
  \begin{vmatrix}
    F_u & F_v\\
    G_u & G_v
  \end{vmatrix}
  \right|_{(x_0,u_0,v_0)}\ne0.
\]
对于上述问题, 我们有以下一般性的定理。

\begin{theorem}[隐函数组存在定理]\label{course-theorem-14-4-3}\quad 设向量函数
\[
  F(x,u)=(F_1(x,u),F_2(x,u),\cdots,F_m(x,u))
\]
在 $U(x_0,\delta)\times U(u_0,\delta)(\delta>0)$ 内有定义, 其中 $x=(x_1,x_2,\cdots,x_n)\in U(x_0,\delta)$, $u=(u_1,u_2,\cdots,u_m)\in U(u_0,\delta)$, 并且满足以下条件:

(1) $F_j(x_0,u_0)=0,\quad j=1,2,\cdots,m$;

(2) 对于 $j=1,2,\cdots,m$, $F_j(x,u)$ 以及它的各个偏导数在 $U(x_0,\delta)\times U(u_0,\delta)$ 内连续;

(3)
\[
  \left.
  \frac{\partial(F_1,F_2,\cdots,F_m)}
  {\partial(u_1,u_2,\cdots,u_m)}
  \right|_{(x_0,u_0)}
  \ne0,
\]

则 $\exists\delta_0(0<\delta_0<\delta)$, 使得在 $U(x_0,\delta_0)$ 内存在唯一的 $m$ 维 $n$ 元向量函数
\[
  f(x)=(f_1(x),f_2(x),\cdots,f_m(x)),
\]
满足

(a) $u_0=(f_1(x_0),f_2(x_0),\cdots,f_m(x_0))$;

(b) 对于 $\forall j(1\leqslant j\leqslant m)$ 及 $\forall x\in U(x_0,\delta_0)$, 有
\[
  F_j(x,f_1(x),\cdots,f_m(x))=0;
\]

(c) $f(x)$ 的每个分量函数 $f_j(x)(j=1,2,\cdots,m)$ 在 $U(x_0,\delta_0)$ 内存在连续偏导数, 记
\[
  A=
  \begin{pmatrix}
    \dfrac{\partial F_1(x,u)}{\partial x_1} &
    \dfrac{\partial F_1(x,u)}{\partial x_2} &
    \cdots &
    \dfrac{\partial F_1(x,u)}{\partial x_n}\\
    \dfrac{\partial F_2(x,u)}{\partial x_1} &
    \dfrac{\partial F_2(x,u)}{\partial x_2} &
    \cdots &
    \dfrac{\partial F_2(x,u)}{\partial x_n}\\
    \vdots & \vdots & & \vdots\\
    \dfrac{\partial F_m(x,u)}{\partial x_1} &
    \dfrac{\partial F_m(x,u)}{\partial x_2} &
    \cdots &
    \dfrac{\partial F_m(x,u)}{\partial x_n}
  \end{pmatrix},
\]
\[
  B=
  \begin{pmatrix}
    \dfrac{\partial F_1(x,u)}{\partial u_1} &
    \dfrac{\partial F_1(x,u)}{\partial u_2} &
    \cdots &
    \dfrac{\partial F_1(x,u)}{\partial u_m}\\
    \dfrac{\partial F_2(x,u)}{\partial u_1} &
    \dfrac{\partial F_2(x,u)}{\partial u_2} &
    \cdots &
    \dfrac{\partial F_2(x,u)}{\partial u_m}\\
    \vdots & \vdots & & \vdots\\
    \dfrac{\partial F_m(x,u)}{\partial u_1} &
    \dfrac{\partial F_m(x,u)}{\partial u_2} &
    \cdots &
    \dfrac{\partial F_m(x,u)}{\partial u_m}
  \end{pmatrix},
\]
那么
\[
  f'(x)=-B^{-1}\cdot A,
\]
其中矩阵 $B^{-1}$ 为 $B$ 的逆矩阵。

\end{theorem}
\begin{proof}\quad 为了证明本定理, 可以对 $m$ 作数学归纳法. 对 $m=2$ 的证明是本质的, 而对 $m\geqslant3$ 的证明则只具有形式的复杂性. 在此我们只给出 $m=2$ 的证明。

由于
\[
  \left.
  \frac{\partial(F_1,F_2)}{\partial(u_1,u_2)}
  \right|_{(x_0,u_0)}
  =
  \left.
  \begin{vmatrix}
    \dfrac{\partial F_1}{\partial u_1} &
    \dfrac{\partial F_1}{\partial u_2}\\
    \dfrac{\partial F_2}{\partial u_1} &
    \dfrac{\partial F_2}{\partial u_2}
  \end{vmatrix}
  \right|_{(x_0,u_0)}
  \ne0,
\]
不失一般性, 我们假定 $\dfrac{\partial F_2(x_0,u_0)}{\partial u_2}\ne0$. 记 $u_0=(u_1^0,u_2^0)$. 由定理~\ref{course-theorem-14-4-2}  知, $\exists\delta_1(0<\delta_1<\delta)$, 使得方程 $F_2(x,u_1,u_2)=0$ 在 $U(x_0,\delta_1)\times U(u_1^0,\delta_1)$ 内唯一确定一个函数
\[
  u_2=g(x,u_1),
\]
满足
\[
  F_2(x,u_1,g(x,u_1))=0,\qquad u_2^0=g(x_0,u_1^0),
\]
且 $g(x,u_1)$ 在该邻域内具有各个连续偏导数. 将 $u_2=g(x,u_1)$ 代入方程 $F_1(x,u_1,u_2)=0$, 得到一个关于 $(x,u_1)$ 的方程
\[
  H(x,u_1)\triangleq F_1(x,u_1,g(x,u_1))=0.
\]
由于
\[
\begin{aligned}
  \frac{\partial H}{\partial u_1}
  &=\frac{\partial F_1}{\partial u_1}
  +\frac{\partial F_1}{\partial u_2}\cdot
  \frac{\partial g(x,u_1)}{\partial u_1}\\
  &=\frac{\partial F_1}{\partial u_1}
  +\frac{\partial F_1}{\partial u_2}
  \left(-\frac{\partial F_2}{\partial u_1}\bigg/
  \frac{\partial F_2}{\partial u_2}\right)\\
  &=\frac1{\dfrac{\partial F_2}{\partial u_2}}\cdot
  \frac{\partial(F_1,F_2)}{\partial(u_1,u_2)},
\end{aligned}
\]
我们有
\[
  \frac{\partial H(x_0,u_1^0)}{\partial u_1}
  =\frac1{\dfrac{\partial F_2(x_0,u_1^0,u_2^0)}{\partial u_2}}\cdot
  \left.
  \frac{\partial(F_1,F_2)}{\partial(u_1,u_2)}
  \right|_{(x_0,u_1^0,u_2^0)}
  \ne0.
\]
因此 $\exists\delta_0(0<\delta_0<\delta_1)$, 使得 $H(x,u_1)=0$ 在 $U(x_0,\delta_0)$ 内唯一确定一个函数 $u_1=f_1(x)$, 满足 $u_1^0=f_1(x_0)$, 且
\[
  F_1(x,f_1(x),g(x,f_1(x)))=0.
\]
现记
\[
  u_1=f_1(x),\qquad u_2=f_2(x)=g(x,f_1(x)),
\]
则在 $U(x_0,\delta_0)$ 内, 有
\[
  \begin{cases}
    F_1(x,f_1(x),f_2(x))=0,\\
    F_2(x,f_1(x),f_2(x))=0
  \end{cases}
\]
和 $u_0=(f_1(x_0),f_2(x_0))$. 再由 $F(x,u)$ 具有各个连续偏导数, 我们可以推出 $f_1(x),f_2(x)$ 在 $U(x_0,\delta_0)$ 内具有各个连续偏导数 (从而 $f_1(x)$ 与 $f_2(x)$ 在该邻域内连续)。

对于 $\forall i(1\leqslant i\leqslant n)$, 利用复合函数的求导法则有
\[
  \begin{cases}
    \dfrac{\partial F_1}{\partial x_i}
    +\dfrac{\partial F_1}{\partial u_1}\cdot\dfrac{\partial f_1}{\partial x_i}
    +\dfrac{\partial F_1}{\partial u_2}\cdot\dfrac{\partial f_2}{\partial x_i}=0,\\[1.2ex]
    \dfrac{\partial F_2}{\partial x_i}
    +\dfrac{\partial F_2}{\partial u_1}\cdot\dfrac{\partial f_1}{\partial x_i}
    +\dfrac{\partial F_2}{\partial u_2}\cdot\dfrac{\partial f_2}{\partial x_i}=0.
  \end{cases}
\]
将上式写成向量形式有
\[
  \begin{pmatrix}
    \dfrac{\partial F_1}{\partial u_1} &
    \dfrac{\partial F_1}{\partial u_2}\\
    \dfrac{\partial F_2}{\partial u_1} &
    \dfrac{\partial F_2}{\partial u_2}
  \end{pmatrix}
  \begin{pmatrix}
    \dfrac{\partial f_1}{\partial x_i}\\
    \dfrac{\partial f_2}{\partial x_i}
  \end{pmatrix}
  =-
  \begin{pmatrix}
    \dfrac{\partial F_1}{\partial x_i}\\
    \dfrac{\partial F_2}{\partial x_i}
  \end{pmatrix},
\]
因此有

\[
  \begin{pmatrix}
    \dfrac{\partial F_1}{\partial u_1} &
    \dfrac{\partial F_1}{\partial u_2}\\
    \dfrac{\partial F_2}{\partial u_1} &
    \dfrac{\partial F_2}{\partial u_2}
  \end{pmatrix}
  \begin{pmatrix}
    \dfrac{\partial f_1}{\partial x_1} &
    \dfrac{\partial f_1}{\partial x_2} &
    \cdots &
    \dfrac{\partial f_1}{\partial x_n}\\
    \dfrac{\partial f_2}{\partial x_1} &
    \dfrac{\partial f_2}{\partial x_2} &
    \cdots &
    \dfrac{\partial f_2}{\partial x_n}
  \end{pmatrix}
  =-
  \begin{pmatrix}
    \dfrac{\partial F_1}{\partial x_1} &
    \dfrac{\partial F_1}{\partial x_2} &
    \cdots &
    \dfrac{\partial F_1}{\partial x_n}\\
    \dfrac{\partial F_2}{\partial x_1} &
    \dfrac{\partial F_2}{\partial x_2} &
    \cdots &
    \dfrac{\partial F_2}{\partial x_n}
  \end{pmatrix}.
\]
由此得
\[
  \begin{pmatrix}
    \dfrac{\partial f_1}{\partial x_1} &
    \dfrac{\partial f_1}{\partial x_2} &
    \cdots &
    \dfrac{\partial f_1}{\partial x_n}\\
    \dfrac{\partial f_2}{\partial x_1} &
    \dfrac{\partial f_2}{\partial x_2} &
    \cdots &
    \dfrac{\partial f_2}{\partial x_n}
  \end{pmatrix}
  =-
  \begin{pmatrix}
    \dfrac{\partial F_1}{\partial u_1} &
    \dfrac{\partial F_1}{\partial u_2}\\
    \dfrac{\partial F_2}{\partial u_1} &
    \dfrac{\partial F_2}{\partial u_2}
  \end{pmatrix}^{-1}
  \begin{pmatrix}
    \dfrac{\partial F_1}{\partial x_1} &
    \dfrac{\partial F_1}{\partial x_2} &
    \cdots &
    \dfrac{\partial F_1}{\partial x_n}\\
    \dfrac{\partial F_2}{\partial x_1} &
    \dfrac{\partial F_2}{\partial x_2} &
    \cdots &
    \dfrac{\partial F_2}{\partial x_n}
  \end{pmatrix}.
\]
现证所确定的隐函数组是唯一的. 事实上, 倘若由
\[
  F(x,u_1,u_2)=0
\]
在 $U(x_0,\delta_0)$ 内确定另一组函数
\[
  \begin{cases}
    u_1=\tilde f_1(x),\\
    u_2=\tilde f_2(x),
  \end{cases}
\]
并且它们仍然满足定理结论 (a), (b) 和 (c), 则容易看出 $\tilde f_1$ 与 $f_1$, $\tilde f_2$ 与 $f_2$ 在该邻域内具有相同的各个偏导数, 从而 $f_1(x)-\tilde f_1(x)$, $f_2(x)-\tilde f_2(x)$ 在 $U(x_0,\delta_0)$ 内均为常数函数. 但它们在 $x_0$ 处有 $f_1(x_0)=\tilde f_1(x_0)$, $f_2(x_0)=\tilde f_2(x_0)$, 因此在 $U(x_0,\delta_0)$ 内有 $f_1(x)=\tilde f_1(x)$, $f_2(x)=\tilde f_2(x)$. 证毕。

\end{proof}
\begin{example}\label{course-example-14-4-3}\quad 证明方程组
\begin{equation}
  \begin{cases}
    x^2+y^2+u^2-v^2=2,\\
    u+v+x+y=0
  \end{cases}
  \label{course-eq-14-4-4}
\end{equation}
在 $(x_0,y_0,u_0,v_0)=(1,1,-1,-1)$ 的某个邻域 $U((1,1,-1,-1),\delta_0)(\delta_0>0)$ 内确定隐函数组
\[
  \begin{cases}
    u=u(x,y),\\
    v=v(x,y),
  \end{cases}
\]
并在该邻域内求 $\dfrac{\partial u}{\partial x}$, $\dfrac{\partial v}{\partial y}$。

\end{example}
\begin{solve}\quad 显然点 $x_0=(1,1,-1,-1)$ 满足方程组 \eqref{course-eq-14-4-4}. 记
\[
  F(x,y,u,v)=x^2+y^2+u^2-v^2-2,
  \qquad
  G(x,y,u,v)=u+v+x+y,
\]
则 $F,G$ 都具有连续偏导数, 且
\[
  \left.
  \frac{\partial(F,G)}{\partial(u,v)}
  \right|_{x_0}
  =
  \left.
  \begin{vmatrix}
    2u & -2v\\
    1 & 1
  \end{vmatrix}
  \right|_{x_0}
  =-2-2=-4\ne0.
\]
因此, $\exists\delta_0>0$, 使得方程组 \eqref{course-eq-14-4-4}  在 $U((1,1,-1,-1),\delta_0)$ 内唯一确定隐函数组
\[
  \begin{cases}
    u=u(x,y),\\
    v=v(x,y),
  \end{cases}
\]
并且 $u(x,y)$ 与 $v(x,y)$ 在 $U((1,1),\delta_0)$ 内存在连续偏导数. 下面我们来求偏导数. 对方程组 \eqref{course-eq-14-4-4}  微分得
\begin{equation}
  \begin{cases}
    2x\,dx+2y\,dy+2u\,du-2v\,dv=0,\\
    du+dv+dx+dy=0.
  \end{cases}
  \label{course-eq-14-4-5}
\end{equation}
令 $dy=0$, 得
\[
  \begin{cases}
    u\,du-v\,dv=-x\,dx,\\
    du+dv=-dx.
  \end{cases}
\]
从中解出 $du=-\dfrac{x+v}{u+v}\,dx$, 从而
\[
  \frac{\partial u}{\partial x}=-\frac{x+v}{u+v}.
\]
在方程组 \eqref{course-eq-14-4-5}  中令 $dx=0$, 即得
\[
  \begin{cases}
    u\,du-v\,dv=-y\,dy,\\
    du+dv=-dy,
  \end{cases}
\]
从而可求出
\[
  \frac{\partial v}{\partial y}=\frac{y-u}{u+v}.
\]

\end{solve}
\subsection{逆映射存在定理}

对于一个可微映射, 如何确定它的逆映射的存在性呢? 隐函数组存在定理可以给出一个连续可微映射局部存在逆映射的充分条件。

\begin{theorem}[逆映射存在定理]\label{course-theorem-14-4-4}\quad 设
\[
  y=(y_1,y_2,\cdots,y_n)=(f_1(x),f_2(x),\cdots,f_n(x))=f(x)
\]
是区域 $D\subset\mathbb R^n$ 到区域 $\Omega\subset\mathbb R^n$ 的一个 $C^1$ 映射, 并且在 $x_0\in D$ 处有
\[
  \left.
  \frac{\partial(f_1,f_2,\cdots,f_n)}
       {\partial(x_1,x_2,\cdots,x_n)}
  \right|_{x_0}
  \ne0.
\]
记 $y_0=f(x_0)$, 则存在 $x_0$ 的邻域 $U(x_0,\delta_0)\subset D$, 使得映射 $y=f(x)$ 是 $U(x_0,\delta_0)$ 到 $f(U(x_0,\delta_0))$ 的 $C^1$ 同胚映射 (变换), 其中 $f(U(x_0,\delta_0))$ 是包含 $y_0$ 的一个区域。

\end{theorem}
\begin{proof}\quad 对于 $j=1,2,\cdots,n$, 记 $F_j(x,y)=y_j-f_j(x)$, 并考虑下面 $n$ 个方程组成的方程组
\[
  \begin{cases}
    F_1(x,y)=0,\\
    F_2(x,y)=0,\\
    \cdots\\
    F_n(x,y)=0.
  \end{cases}
\]
由定理条件知它们满足

\noindent (1) 对任意的 $j(1\leqslant j\leqslant n)$, $F_j(x_0,y_0)=0$;

\noindent (2) $\exists\delta'>0$, 对任意的 $j(1\leqslant j\leqslant n)$, 使得 $F_j(x,y)$ 在 $U((x_0,y_0),\delta')$ 内具有连续偏导数 (因此 $F_j(x,y)$ 在该邻域连续), 并且
\[
  \left.
  \frac{\partial(F_1,F_2,\cdots,F_n)}
       {\partial(x_1,x_2,\cdots,x_n)}
  \right|_{(x_0,y_0)}
  =
  (-1)^n
  \left.
  \frac{\partial(f_1,f_2,\cdots,f_n)}
       {\partial(x_1,x_2,\cdots,x_n)}
  \right|_{x_0}
  \ne0,
\]
因此由定理~\ref{course-theorem-14-4-3}  知, 在 $y_0$ 的某个邻域 $U(y_0,\delta_1)(0<\delta_1<\delta')$ 内存在一个 $n$ 维向量函数
\[
  x=g(y)=(g_1(y),\cdots,g_n(y)),
\]
满足 $x_0=(g_1(y_0),\cdots,g_n(y_0))$ 和
\[
  \begin{cases}
    y_1-f_1(g_1(y),\cdots,g_n(y))=0,\\
    \cdots\\
    y_n-f_n(g_1(y),\cdots,g_n(y))=0,
  \end{cases}
\]
并且 $(g_1(y),\cdots,g_n(y))$ 在 $U(y_0,\delta_1)$ 内具有各个连续偏导数. 上面的 $n$ 个等式同时表明, 在 $U(y_0,\delta_1)$ 内 $f(g(y))\equiv y$, 即 $x=g(y)$ 是 $y=f(x)$ 的逆映射. 我们从定理~\ref{course-theorem-14-4-3}  推知 $f'(x)g'(y)=E$ ($n$ 阶单位矩阵), 因此由 $|f'(x_0)|\ne0$ 有 $g'(y_0)=[f'(x_0)]^{-1}$, 从而 $g(y)$ 在 $y_0$ 处的雅可比行列式也不为零. 如果将 $f(x)$ 与 $g(y)$ 的地位对换, 则在 $x_0$ 的某个邻域 $U(x_0,\delta_1')(0<\delta_1'<\delta')$ 内存在 $g(y)$ 的逆映射, 显然该逆映射是 $f(x)$。

现证存在 $x_0$ 的邻域 $U(x_0,\delta_0)$, 使得 $f(x)$ 是 $U(x_0,\delta_0)$ 到区域 $f(U(x_0,\delta_0))$ 的 $C^1$ 同胚映射. 首先取 $\delta_2(0<\delta_2<\delta_1)$ 充分小, 使得对于 $\forall x\in U(x_0,\delta_2)$, $f(x)$ 在 $x$ 处的雅可比行列式不等于零. 下面我们证明, 对于 $\forall x'\in U(x_0,\delta_2)$ 及任何包含 $x'$ 的开集 $U_{x'}$, $f(x')$ 是 $f(U_{x'})$ 的内点. 由于我们证明该事实只是利用 $f(x)$ 是 $C^1$ 映射, 并且它在 $x'$ 处的雅可比行列式不为零, 容易看出, 我们只要对 $x_0$ 证明上述结论即可. 任取包含 $x_0$ 的开集 $U_{x_0}$, 则由 $g(y)$ 在 $y_0$ 处的连续性, 存在 $y_0$ 的邻域 $U(y_0,\delta_3)$ (当然是开集), 使得 $g(U(y_0,\delta_3))\subset U_{x_0}$. 因此有
\[
  U(y_0,\delta_3)=f(g(U(y_0,\delta_3)))\subset f(U_{x_0}).
\]
再注意到连续映射将连通集映成连通集, 因此 $f(U(x_0,\delta_2))$ 是一个包含 $y_0$ 的区域. 同理, $\exists\delta_2'(0<\delta_2'<\delta_1')$, 使得 $g(U(y_0,\delta_2'))$ 是包含 $x_0$ 的一个区域. 因此取 $\delta_0(0<\delta_0<\delta_2')$ 充分小, 使得 $f(U(x_0,\delta_0))\subset U(y_0,\delta_2')$, 则 $f(U(x_0,\delta_0))$ 是 $\mathbb R^n$ 中包含 $y_0$ 的区域. 由于在 $U(x_0,\delta_0)$ 满足 $g(f(x))=x$, 并且 $f$ 在 $U(x_0,\delta_0)$ 内是 $C^1$ 映射, $g(y)=f^{-1}(y)$ 在 $f(U(x_0,\delta_0))$ 内是 $C^1$ 映射, 因此 $f(x)$ 是 $U(x_0,\delta_0)$ 到 $f(U(x_0,\delta_0))$ 的 $C^1$ 同胚映射 (变换). 证毕。

从逆映射存在定理的证明可知, 如果一个 $n$ 维 $n$ 元向量函数
\[
  y=f(x)=(f_1(x),f_2(x),\cdots,f_n(x))
\]
在区域 $D\subset\mathbb R^n$ 内具有连续偏导数, 并且对于 $\forall x\in D$, $f'(x)$ 非退化, 即它的雅可比行列式
\[
  \left.
  \frac{\partial(f_1,f_2,\cdots,f_n)}
       {\partial(x_1,x_2,\cdots,x_n)}
  \right|_x
  \ne0,
\]
则局部总存在逆映射, 因此 $f$ 必将开集映成开集. 另外, 由于此时 $y=f(x)$ 是连续函数, 因此推出 $f(D)$ 为一个区域。

\end{proof}
\noindent\keyterm{注 1}\quad 当一个 $n$ 维 $n$ 元向量函数 $y=f(x)$ 在区域 $D\subset\mathbb R^n$ 内具有连续偏导数, 并且对于 $\forall x\in D$ 有 $|f'(x)|\ne0$ 时, 我们仅仅能推出 $y=f(x)$ 是局部同胚映射, 但不能推出它是 $D$ 到 $f(D)$ 的整体同胚映射. 例如:
\[
  \begin{pmatrix}
    u\\
    v
  \end{pmatrix}
  =
  \begin{pmatrix}
    e^x\cos y\\
    e^x\sin y
  \end{pmatrix},
  \qquad (x,y)\in\mathbb R^2
\]
是 $\mathbb R^2$ 到 $\mathbb R^2\setminus\{(0,0)\}$ 的映射, 对于 $\forall(x,y)\in\mathbb R^2$, 有
\[
  \frac{\partial(u,v)}{\partial(x,y)}
  =
  \begin{vmatrix}
    e^x\cos y & -e^x\sin y\\
    e^x\sin y & e^x\cos y
  \end{vmatrix}
  =e^x\ne0,
\]
因此该映射局部总是同胚映射. 由 $\cos y$ 及 $\sin y$ 的周期性, 上述映射显然不是 $\mathbb R^2\to\mathbb R^2\setminus\{(0,0)\}$ 的同胚映射。

\noindent\keyterm{注 2}\quad 对于逆映射存在定理, 我们容易举例说明, 如果仅仅要求存在逆映射, 而不要求逆映射具有连续偏导数, 雅可比行列式不为零的条件不是必要条件. 例如,
\[
  \begin{pmatrix}
    u\\
    v
  \end{pmatrix}
  =
  \begin{pmatrix}
    x^3\\
    y
  \end{pmatrix}
\]
是 $\mathbb R^2\to\mathbb R^2$ 的同胚映射. 事实上, 它的逆映射为
\[
  \begin{pmatrix}
    x\\
    y
  \end{pmatrix}
  =
  \begin{pmatrix}
    u^{1/3}\\
    v
  \end{pmatrix},
\]
但在 $(0,0)$ 处有
\[
  \left.
  \frac{\partial(u,v)}{\partial(x,y)}
  \right|_{(0,0)}
  =0.
\]
此例同时说明了当 $\dfrac{\partial x}{\partial u}$ (或 $\dfrac{\partial y}{\partial u}$) 不是处处存在时, 仍然具有连续的逆映射。



\section{多元函数的极值}\label{course-14-5}
本节用于核对多元极限、隐函数、驻点与约束极值题中的可微性和非退化条件。仅有偏导数存在不能替代全微分存在；用必要条件列出候选点以后，还要比较边界并验证极值是否取得。

\subsection{通常极值问题}

在一元函数中, 我们利用导数讨论了函数的极值. 在本节中, 我们来讨论多元函数的极值问题. 对于多元函数来说, 极值的定义与一元函数是相似的。

\begin{definition}\label{course-definition-14-5-1}\quad 设函数 $u=f(x)$ 在区域 $D$ 内有定义, $x_0\in D$, 若存在 $x_0$ 的邻域 $U(x_0,\delta_0)\subset D(\delta_0>0)$, 使得当 $x\in U(x_0,\delta_0)$ 时, 有 $f(x)\leqslant f(x_0)$ (或 $f(x)\geqslant f(x_0)$), 则称 $f(x)$ 在 $x_0$ 处取极大值 (或极小值), $x_0$ 称为 $f(x)$ 的极大值点 (或极小值点). 如果上述不等式为严格不等式, 则称 $x_0$ 为 $f(x)$ 的严格极大值点 (或严格极小值点)。

对于一个二元可微函数 $z=f(x,y)((x,y)\in D)$, 一般来说它的图像是一块曲面. 从几何上来看, 若 $(x_0,y_0)\in D$ 是 $z=f(x,y)$ 的极值点, 则 $(x_0,y_0,f(x_0,y_0))$ 是曲面上局部的最高点或最低点. 显然, 对于曲面上的任何一条过 $(x_0,y_0,f(x_0,y_0))$ 的曲线 $\Gamma$, 该点也是 $\Gamma$ 在它附近的最高点或最低点. 当 $\Gamma$ 在该点具有切线时, 它的切线应该平行于 $Oxy$ 平面. 所以, 当 $z=f(x,y)$ 在 $(x_0,y_0)$ 处可微时, 有
\[
  \frac{\partial f(x_0,y_0)}{\partial x}
  =
  \frac{\partial f(x_0,y_0)}{\partial y}
  =0,
\]
即 $f'(x_0,y_0)=\operatorname{grad} f(x_0)=0$。

\end{definition}
\begin{theorem}\label{course-theorem-14-5-1}\quad 设函数 $f(x)$ 在区域 $D\subset\mathbb R^n$ 内有定义, $x_0=(x_1^0,x_2^0,\cdots,x_n^0)\in D$, 再设 $f(x)$ 在 $x_0$ 处取极值并且 $f(x)$ 在该点关于 $x_i(1\leqslant i\leqslant n)$ 可偏导, 则有
\[
  \frac{\partial f(x_0)}{\partial x_i}=0.
\]
特别地, 若 $f(x)$ 在该点可微, 则 $f'(x_0)=0$。

\end{theorem}
\begin{proof}\quad 在定理的条件下, 一元函数
\[
  F_i(x_i)=f(x_1^0,\cdots,x_{i-1}^0,x_i,x_{i+1}^0,\cdots,x_n^0)
  \qquad (1\leqslant i\leqslant n)
\]
在 $x_i^0$ 处取极值, 因此
\[
  F_i'(x_i^0)=\frac{\partial f(x_0)}{\partial x_i}=0
  \qquad (1\leqslant i\leqslant n).
\]
证毕。

如同一元函数, 对于一个 $n$ 元可微函数 $f(x)$, 若 $f'(x_0)=0$, 则称 $x_0$ 为 $f(x)$ 的一个驻点或临界点. 因此, 对可微函数 $f(x)$ 而言, $f(x)$ 的极值点一定是驻点, 但反命题一般不真。

例如, $f(x,y)=y^2-x^2$ 在 $(0,0)$ 处显然不取极值. 该函数沿两条坐标轴的变化方向相反，因而称其图像为马鞍面. 当一个可微函数 $f(x)$ 的驻点不是极值点时, 该驻点也称为 $f(x)$ 的一个鞍点。

在研究一元函数的极值问题中, 我们曾经利用它的高阶导数来判别驻点是否为极值点. 对于多元函数, 我们也可以利用高阶偏导数来讨论其极值问题. 由于多元函数高阶偏导数的复杂性, 我们只考虑涉及二阶偏导数的情形。

\end{proof}
\begin{theorem}\label{course-theorem-14-5-2}\quad 设函数 $f(x)=f(x_1,x_2,\cdots,x_n)$ 在区域 $D\subset\mathbb R^n$ 内具有二阶连续偏导数, 且 $f'(x_0)=0(x_0\in D)$, 再设 $f(x)$ 在 $x_0$ 处的海色矩阵 $H_f(x_0)$ 为满秩矩阵, 则

\noindent (1) 当 $H_f(x_0)$ 正定时, $f(x)$ 在 $x_0$ 处取极小值;

\noindent (2) 当 $H_f(x_0)$ 负定时, $f(x)$ 在 $x_0$ 处取极大值;

\noindent (3) 当 $H_f(x_0)$ 不定时, $f(x)$ 在 $x_0$ 处不取极值。

\end{theorem}
\begin{proof}\quad 在 $x_0$ 的充分小邻域 $U(x_0,\delta_0)(\delta_0>0)$ 内, $f(x)$ 有泰勒公式
\[
\begin{aligned}
  f(x)
  &=f(x_0)+f'(x_0)(x-x_0)^{\mathrm T}
    +\frac12(x-x_0)H_f(x_0)(x-x_0)^{\mathrm T}
    +o(|x-x_0|^2)\\
  &=f(x_0)+\frac12
    \frac{x-x_0}{|x-x_0|}
    H_f(x_0)
    \frac{(x-x_0)^{\mathrm T}}{|x-x_0|}
    |x-x_0|^2
    +o(|x-x_0|^2)
    \qquad (|x-x_0|\to0).
\end{aligned}
\]
对于 $\forall x\in U_0(x_0,\delta_0)$, 记 $x'=\dfrac{x-x_0}{|x-x_0|}$, 则 $|x'|=1$. 由于 $x'H_f(x_0)x'^{\mathrm T}$ 是 $S(0;1)=\{x':|x'|=1\}$ 上的一个连续函数, 注意到 $S(0;1)$ 是紧集, $x'H_f(x_0)x'^{\mathrm T}$ 在 $S(0;1)$ 上取到最大值 $M$ 和最小值 $m$。

\noindent (1) 当 $H_f(x_0)$ 正定时, 对任意非零向量 $x-x_0$ 均有
\[
  (x-x_0)H_f(x_0)(x-x_0)^{\mathrm T}>0,
\]
因此有 $m>0$, 从而存在 $0<\delta_1<\delta_0$, 使得当 $x\in U_0(x_0,\delta_1)$ 时, 有
\[
\begin{aligned}
  f(x)
  &=f(x_0)+\frac12
    \frac{x-x_0}{|x-x_0|}
    H_f(x_0)
    \frac{(x-x_0)^{\mathrm T}}{|x-x_0|}
    |x-x_0|^2+o(|x-x_0|^2)\\
  &>f(x_0)+\frac m4|x-x_0|^2>f(x_0).
\end{aligned}
\]
这就证明了 $f(x)$ 在 $x_0$ 处取极小值。

\noindent (2) 当 $H_f(x_0)$ 负定时, 有 $M<0$. 因此, 存在 $0<\delta_2<\delta_0$, 当 $x\in U_0(x_0,\delta_2)$ 时, 有
\[
\begin{aligned}
  f(x)
  &=f(x_0)+\frac12
    \frac{x-x_0}{|x-x_0|}
    H_f(x_0)
    \frac{(x-x_0)^{\mathrm T}}{|x-x_0|}
    |x-x_0|^2+o(|x-x_0|^2)\\
  &<f(x_0)+\frac M4|x-x_0|^2<f(x_0).
\end{aligned}
\]
这就证明了 $f(x)$ 在 $x_0$ 取极大值。

\noindent (3) 当 $H_f(x_0)$ 不定时, 有 $m<0$, $M>0$. 任取 $x_1\in\mathbb R^n$, 使得
\[
  \frac{x_1-x_0}{|x_1-x_0|}
  H_f(x_0)
  \frac{(x_1-x_0)^{\mathrm T}}{|x_1-x_0|}
  =m,
\]
则对充分小的 $t>0$, 有
\[
  f\left(x_0+t\frac{x_1-x_0}{|x_1-x_0|}\right)
  =f(x_0)+\frac12mt^2+o(t^2)
  \leqslant f(x_0)+\frac14mt^2<f(x_0).
\]
再取 $x_2\in\mathbb R^n$, 使得
\[
  \frac{x_2-x_0}{|x_2-x_0|}
  H_f(x_0)
  \frac{(x_2-x_0)^{\mathrm T}}{|x_2-x_0|}
  =M,
\]
则对充分小的 $t>0$, 有
\[
  f\left(x_0+t\frac{x_2-x_0}{|x_2-x_0|}\right)
  =f(x_0)+\frac12Mt^2+o(t^2)
  \geqslant f(x_0)+\frac14Mt^2>f(x_0).
\]
这说明 $f(x)$ 在 $x_0$ 处不取极值. 证毕。

在代数课程中, 有关于 $H_f(x_0)$ 正定、负定和不定的判定法则, 在这里我们就不再赘述了. 我们经常用到 $n=2$ 的情况, 下面给出定理~\ref{course-theorem-14-5-2}  在 $\mathbb R^2$ 情形下的特殊形式。

\end{proof}
\noindent\keyterm{推论}\quad 设函数 $f(x,y)$ 在区域 $D\subset\mathbb R^2$ 内具有二阶连续偏导数, $(x_0,y_0)\in D$, 再设 $f'(x_0,y_0)=0$, 且记
\[
  H_f(x_0,y_0)
  =
  \left.
  \begin{pmatrix}
    \dfrac{\partial^2 f}{\partial x^2} &
    \dfrac{\partial^2 f}{\partial x\partial y}\\
    \dfrac{\partial^2 f}{\partial x\partial y} &
    \dfrac{\partial^2 f}{\partial y^2}
  \end{pmatrix}
  \right|_{(x_0,y_0)}
  \triangleq
  \begin{pmatrix}
    A & B\\
    B & C
  \end{pmatrix},
\]
则

\noindent (1) 当 $A>0$,
\[
  \begin{vmatrix}
    A & B\\
    B & C
  \end{vmatrix}
  =AC-B^2>0
\]
时, $f(x,y)$ 在 $(x_0,y_0)$ 处取极小值;

\noindent (2) 当 $A<0$,
\[
  \begin{vmatrix}
    A & B\\
    B & C
  \end{vmatrix}
  =AC-B^2>0
\]
时, $f(x,y)$ 在 $(x_0,y_0)$ 处取极大值;

\noindent (3) 当
\[
  \begin{vmatrix}
    A & B\\
    B & C
  \end{vmatrix}
  =AC-B^2<0
\]
时, $f(x,y)$ 在 $(x_0,y_0)$ 处不取极值。

\begin{example}\label{course-example-14-5-1}\quad 试求函数 $f(x,y)=x^2+3xy+3y^2-6x-3y$ 的极值。

\end{example}
\begin{solve}\quad 由
\[
  f_x'(x,y)=2x+3y-6=0,\qquad
  f_y'(x,y)=3x+6y-3=0
\]
得唯一驻点 $(9,-4)$, 再由于
\[
  f_{xx}''(x,y)=2,\qquad f_{xy}''(x,y)=3,\qquad f_{yy}''(x,y)=6,
\]
从而
\[
  f_{xx}''(9,-4)=2>0,
\]
\[
  \begin{vmatrix}
    f_{xx}''(9,-4) & f_{xy}''(9,-4)\\
    f_{xy}''(9,-4) & f_{yy}''(9,-4)
  \end{vmatrix}
  =
  \begin{vmatrix}
    2 & 3\\
    3 & 6
  \end{vmatrix}
  =3>0,
\]
因此 $f(x,y)$ 在 $(9,-4)$ 处取极小值 $-21$. 因为 $f(x,y)$ 在 $\mathbb R^2$ 处处可微, 所以 $(9,-4)$ 是 $f(x,y)$ 的唯一极值点。

\end{solve}
\begin{example}\label{course-example-14-5-2}\quad 某工厂生产一批长方体无盖盒子, 要求其体积为 $1\mathrm{m}^3$, 盒子的底的厚度是侧面厚度的三倍. 问: 如何设定盒子的长、宽、高才能使得用料最省?

\end{example}
\begin{solve}\quad 设盒子的长、宽、高分别为 $x,y,z$, 并令
\[
  F(x,y,z)=3xy+2yz+2xz.
\]
由于盒子的体积为 $1\mathrm{m}^3$, 即 $xyz=1$, 因此, 若令
\[
  S(x,y)=3xy+2\left(\frac1x+\frac1y\right),
\]
则本题的最值问题转化成求 $S(x,y)$ 的最小值问题. 由
\[
  \frac{\partial S(x,y)}{\partial x}=3y-\frac2{x^2}=0,\qquad
  \frac{\partial S(x,y)}{\partial y}=3x-\frac2{y^2}=0
\]
解得
\[
  x=y=\sqrt[3]{\frac23}.
\]
由于该实际问题总是存在最小值, 且 $\left(\sqrt[3]{\dfrac23},\sqrt[3]{\dfrac23}\right)$ 是 $S(x,y)$ 唯一的驻点, 因此它是 $S(x,y)$ 的极小值点, 且 $S(x,y)$ 必在该点取最小值. 所以当
\[
  x=y=\sqrt[3]{\frac23},\qquad z=\sqrt[3]{\frac94}
\]
时可使用料最省。

我们也可以由 $S(x,y)$ 的海色矩阵的正定性来判断该点为极小值点. 事实上, 在 $\left(\sqrt[3]{\dfrac23},\sqrt[3]{\dfrac23}\right)$ 处, $S(x,y)$ 的海色矩阵为
\[
  H_S\left(\sqrt[3]{\frac23},\sqrt[3]{\frac23}\right)
  =
  \begin{pmatrix}
    6 & 3\\
    3 & 6
  \end{pmatrix}.
\]
容易看出它是正定的, 从而 $\left(\sqrt[3]{\dfrac23},\sqrt[3]{\dfrac23}\right)$ 为 $S(x,y)$ 的极小值点。

\end{solve}
\begin{example}\label{course-example-14-5-3}\quad 设 $(x_k,y_k)(k=1,2,\cdots,K)$ 为 $\mathbb R^2$ 内 $K$ 个两两不同的点, 且它们不同在 $Oxy$ 平面内的一条垂直于 $x$ 轴的直线上. 证明: 存在唯一直线 $L_0:y=ax+b$, 使得函数
\[
  f(s,t)=\sum_{k=1}^K(sx_k+t-y_k)^2
\]
在 $(a,b)$ 处达到最小。

\end{example}
\begin{proof}\quad 显然 $f(s,t)=\sum\limits_{k=1}^K(sx_k+t-y_k)^2$ 在 $\mathbb R^2$ 中处处可微. 由
\[
  \begin{cases}
    \dfrac{\partial f}{\partial s}
    =2\displaystyle\sum_{k=1}^K x_k(sx_k+t-y_k)=0,\\[1.2ex]
    \dfrac{\partial f}{\partial t}
    =2\displaystyle\sum_{k=1}^K (sx_k+t-y_k)=0
  \end{cases}
\]
得
\[
  \begin{cases}
    s\displaystyle\sum_{k=1}^K x_k^2
    +t\displaystyle\sum_{k=1}^K x_k
    =\displaystyle\sum_{k=1}^K x_ky_k,\\[1.2ex]
    s\displaystyle\sum_{k=1}^K x_k
    +Kt
    =\displaystyle\sum_{k=1}^K y_k.
  \end{cases}
\]
由此求出上述方程组的解为
\[
  s=
  \frac{
    K\displaystyle\sum_{k=1}^K x_ky_k
    -\left(\displaystyle\sum_{k=1}^K x_k\right)
     \left(\displaystyle\sum_{k=1}^K y_k\right)}
  {
    K\displaystyle\sum_{k=1}^K x_k^2
    -\left(\displaystyle\sum_{k=1}^K x_k\right)^2}
  \triangleq a,
\]
\[
  t=
  \frac{
    \left(\displaystyle\sum_{k=1}^K x_k^2\right)
    \left(\displaystyle\sum_{k=1}^K y_k\right)
    -\left(\displaystyle\sum_{k=1}^K x_k\right)
     \left(\displaystyle\sum_{k=1}^K x_ky_k\right)}
  {
    K\displaystyle\sum_{k=1}^K x_k^2
    -\left(\displaystyle\sum_{k=1}^K x_k\right)^2}
  \triangleq b.
\]
由于该问题中的最小值总是存在的, 且 $f(s,t)$ 只有唯一的驻点 $(a,b)$, 因此 $(a,b)$ 必为 $f(s,t)$ 的最小值点。

在上述例子中, 已知平面 $\mathbb R^2$ 内 $K$ 个不落在一条垂直于 $x$ 轴的直线上的点, 求满足定理条件的唯一直线 $y=ax+b$ 的方法称为是最小二乘法。

\end{proof}
\subsection{条件极值问题}

为了导出条件极值问题, 我们先来看一个例子. 设 $\Gamma\subset\mathbb R^3$ 是一条曲线, 现在要求它到原点的最短距离. 为了求解此问题, 我们可令 $f(x,y,z)=\sqrt{x^2+y^2+z^2}$, 该函数表示了 $\mathbb R^3$ 中任一点到原点的距离. 于是该问题转化为: 当 $(x,y,z)$ 在 $\Gamma$ 上变化时, 求 $f(x,y,z)$ 的最小值. 实际上, 这是 $(x,y,z)$ 在 $\Gamma$ 上变化的条件下, 求 $f(x,y,z)$ 的极小值问题. 像这样具有外加约束条件的极值问题称为条件极值问题. 为了求解这一条件极值问题, 一个自然的想法是: 求出 $\Gamma$ 的参数方程然后代入 $f(x,y,z)$, 从而将此问题转化成上一小节讨论的极值问题. 但在很多情形下, 这绝非易事. 因此我们需用其他方法来求解条件极值问题。

首先, 设 $\Gamma$ 是一条平面曲线, 它由方程 $\varphi(x,y)=0$ 所确定. 我们假定 $\varphi(x,y)$ 在区域 $D\subset\mathbb R^2$ 内具有连续偏导数, 且对于 $\forall(x,y)\in D$,
\[
  \left.
  \left(\frac{\partial\varphi}{\partial x},\frac{\partial\varphi}{\partial y}\right)
  \right|_{(x,y)}
\]
的秩为 $1$, 再设 $(x_0,y_0)\in\Gamma$ 且是函数 $f(x,y)=\sqrt{x^2+y^2}$ 的一个极值点. 在理论上 (由隐函数存在定理) 我们可以从 $\varphi(x,y)=0$ 在 $(x_0,y_0)$ 附近解出 $y=y(x)$ 或 $x=x(y)$. 现在不妨设
\[
  \frac{\partial\varphi(x_0,y_0)}{\partial y}\ne0,
\]
从而 $\exists\delta>0$, 使得在 $U(x_0,\delta)$ 内存在 $y=y(x)$, 满足 $y_0=y(x_0)$ 且 $\varphi(x,y(x))=0$. 将 $y=y(x)$ 代入 $f(x,y)$ 得 $f(x,y(x))$, 则 $x_0$ 是 $f(x,y(x))$ 的一个通常极值点, 从而有
\[
  \frac{\partial f(x_0,y_0)}{\partial x}
  +
  \frac{\partial f(x_0,y_0)}{\partial y}y'(x_0)=0.
\]
由
\[
  y'(x_0)=-\frac{\varphi_x'(x_0,y_0)}{\varphi_y'(x_0,y_0)}
\]
知
\[
  \frac{\dfrac{\partial f(x_0,y_0)}{\partial x}}
       {\dfrac{\partial\varphi(x_0,y_0)}{\partial x}}
  =
  \frac{\dfrac{\partial f(x_0,y_0)}{\partial y}}
       {\dfrac{\partial\varphi(x_0,y_0)}{\partial y}},
\]
即 $\exists\lambda\in\mathbb R$, 使得
\[
  \begin{cases}
    \dfrac{\partial f(x_0,y_0)}{\partial x}
    +\lambda\dfrac{\partial\varphi(x_0,y_0)}{\partial x}=0,\\[1.2ex]
    \dfrac{\partial f(x_0,y_0)}{\partial y}
    +\lambda\dfrac{\partial\varphi(x_0,y_0)}{\partial y}=0.
  \end{cases}
\]
再设 $\Gamma$ 是由下述方程组确定的一条空间曲线:
\[
  \begin{cases}
    \varphi_1(x,y,z)=0,\\
    \varphi_2(x,y,z)=0,
  \end{cases}
\]
并假定 $\varphi_1,\varphi_2$ 在区域 $D\subset\mathbb R^3$ 内具有连续偏导数, 且对于 $\forall(x,y,z)\in D$,
\[
  \left.
  \begin{pmatrix}
    \dfrac{\partial\varphi_1}{\partial x} &
    \dfrac{\partial\varphi_1}{\partial y} &
    \dfrac{\partial\varphi_1}{\partial z}\\
    \dfrac{\partial\varphi_2}{\partial x} &
    \dfrac{\partial\varphi_2}{\partial y} &
    \dfrac{\partial\varphi_2}{\partial z}
  \end{pmatrix}
  \right|_{(x,y,z)}
\]
的秩为 $2$. 由此条件, 理论上从方程组我们仍然可解出曲线的参数方程, 我们将其代入 $f(x,y,z)$ 就可以推出在 $f(x,y,z)$ 的极值点 $(x_0,y_0,z_0)$ 处, 存在常数 $\lambda_1,\lambda_2$ 成立下述方程组:
\[
  \begin{cases}
    \dfrac{\partial f(x_0,y_0,z_0)}{\partial x}
    +\lambda_1\dfrac{\partial\varphi_1(x_0,y_0,z_0)}{\partial x}
    +\lambda_2\dfrac{\partial\varphi_2(x_0,y_0,z_0)}{\partial x}=0,\\[1.2ex]
    \dfrac{\partial f(x_0,y_0,z_0)}{\partial y}
    +\lambda_1\dfrac{\partial\varphi_1(x_0,y_0,z_0)}{\partial y}
    +\lambda_2\dfrac{\partial\varphi_2(x_0,y_0,z_0)}{\partial y}=0,\\[1.2ex]
    \dfrac{\partial f(x_0,y_0,z_0)}{\partial z}
    +\lambda_1\dfrac{\partial\varphi_1(x_0,y_0,z_0)}{\partial z}
    +\lambda_2\dfrac{\partial\varphi_2(x_0,y_0,z_0)}{\partial z}=0.
  \end{cases}
\]
对于一般情形, 我们可以证明下面的定理。

\begin{theorem}\label{course-theorem-14-5-3}\quad 设函数 $f(x)$, $\varphi(x)=(\varphi_1(x),\varphi_2(x),\cdots,\varphi_m(x))$ 在区域 $D\subset\mathbb R^n(m<n)$ 内具有各个连续偏导数, 再设 $x_0=(x_1^0,x_2^0,\cdots,x_n^0)\in D$ 为 $f(x)$ 在约束条件
\begin{equation}
  \begin{cases}
    \varphi_1(x)=0,\\
    \varphi_2(x)=0,\\
    \cdots\\
    \varphi_m(x)=0
  \end{cases}
  \label{course-eq-14-5-1}
\end{equation}
下的极值点, 并且 $\varphi'(x_0)$ 的秩为 $m$, 则存在常数 $\lambda_1,\lambda_2,\cdots,\lambda_m\in\mathbb R$, 使得在 $x_0$ 处成立下述等式:
\begin{equation}
  \begin{cases}
    \dfrac{\partial f(x_0)}{\partial x_i}
    +\displaystyle\sum_{j=1}^m\lambda_j
    \dfrac{\partial\varphi_j(x_0)}{\partial x_i}=0,
    \qquad i=1,2,\cdots,n,\\[1.2ex]
    \varphi_j(x_0)=0,\qquad j=1,2,\cdots,m.
  \end{cases}
\label{course-eq-14-5-2}
\end{equation}

\end{theorem}
\begin{proof}\quad 由于 $\varphi'(x_0)$ 的秩为 $m$, 我们不妨设行列式
\[
  \left.
  \frac{\partial(\varphi_1,\varphi_2,\cdots,\varphi_m)}
       {\partial(x_{n-m+1},x_{n-m+2},\cdots,x_n)}
  \right|_{x_0}
  \ne0.
\]
因此, 在 $x_0$ 的某个邻域内唯一确定一组具有各个连续偏导数的隐函数
\begin{equation}
  \begin{cases}
    x_{n-m+1}=g_1(x_1,x_2,\cdots,x_{n-m}),\\
    x_{n-m+2}=g_2(x_1,x_2,\cdots,x_{n-m}),\\
    \cdots\\
    x_n=g_m(x_1,x_2,\cdots,x_{n-m}),
  \end{cases}
  \label{course-eq-14-5-3}
\end{equation}
满足
\[
  x_j^0=g_j(x_1^0,x_2^0,\cdots,x_{n-m}^0)
  \qquad (j=n-m+1,n-m+2,\cdots,n)
\]
及对 $k=1,2,\cdots,m$, 有
\[
  \varphi_k(x_1,\cdots,x_{n-m},g_1(x_1,x_2,\cdots,x_{n-m}),\cdots,g_m(x_1,x_2,\cdots,x_{n-m}))=0.
\]
将隐函数组 \eqref{course-eq-14-5-3}  代入 $f(x)$, 得
\[
  f(x_1,\cdots,x_{n-m},g_1(x_1,x_2,\cdots,x_{n-m}),\cdots,g_m(x_1,x_2,\cdots,x_{n-m})).
\]
因此, $x_0$ 是条件极值点转化为 $(x_1^0,x_2^0,\cdots,x_{n-m}^0)$ 为上述函数的通常极值点。

令 $x_0'=(x_1^0,x_2^0,\cdots,x_{n-m}^0)$, 则对 $i=1,2,\cdots,n-m$, 有
\[
  \frac{\partial f(x_0)}{\partial x_i}
  +\frac{\partial f(x_0)}{\partial x_{n-m+1}}\cdot
  \frac{\partial g_1(x_0')}{\partial x_i}
  +\cdots+
  \frac{\partial f(x_0)}{\partial x_n}\cdot
  \frac{\partial g_m(x_0')}{\partial x_i}
  =0.
\]
令 $g(x')=(g_1(x'),g_2(x'),\cdots,g_m(x'))^{\mathrm T}$, 其中 $x'=(x_1,x_2,\cdots,x_{n-m})$. 将上述 $n-m$ 个等式写成向量形式, 有
\begin{equation}
  \left(
    \frac{\partial f(x_0)}{\partial x_1},
    \cdots,
    \frac{\partial f(x_0)}{\partial x_{n-m}}
  \right)
  +
  \left(
    \frac{\partial f(x_0)}{\partial x_{n-m+1}},
    \cdots,
    \frac{\partial f(x_0)}{\partial x_n}
  \right)
  g'(x_0')=0.
  \label{course-eq-14-5-4}
\end{equation}
由于
\begin{equation}
\begin{aligned}
  g'(x_0')
  &=
  -
  \begin{pmatrix}
    \dfrac{\partial\varphi_1(x_0)}{\partial x_{n-m+1}} &
    \dfrac{\partial\varphi_1(x_0)}{\partial x_{n-m+2}} &
    \cdots &
    \dfrac{\partial\varphi_1(x_0)}{\partial x_n}\\
    \dfrac{\partial\varphi_2(x_0)}{\partial x_{n-m+1}} &
    \dfrac{\partial\varphi_2(x_0)}{\partial x_{n-m+2}} &
    \cdots &
    \dfrac{\partial\varphi_2(x_0)}{\partial x_n}\\
    \vdots & \vdots & & \vdots\\
    \dfrac{\partial\varphi_m(x_0)}{\partial x_{n-m+1}} &
    \dfrac{\partial\varphi_m(x_0)}{\partial x_{n-m+2}} &
    \cdots &
    \dfrac{\partial\varphi_m(x_0)}{\partial x_n}
  \end{pmatrix}^{-1}\\
  &\quad\cdot
  \begin{pmatrix}
    \dfrac{\partial\varphi_1(x_0)}{\partial x_1} &
    \dfrac{\partial\varphi_1(x_0)}{\partial x_2} &
    \cdots &
    \dfrac{\partial\varphi_1(x_0)}{\partial x_{n-m}}\\
    \dfrac{\partial\varphi_2(x_0)}{\partial x_1} &
    \dfrac{\partial\varphi_2(x_0)}{\partial x_2} &
    \cdots &
    \dfrac{\partial\varphi_2(x_0)}{\partial x_{n-m}}\\
    \vdots & \vdots & & \vdots\\
    \dfrac{\partial\varphi_m(x_0)}{\partial x_1} &
    \dfrac{\partial\varphi_m(x_0)}{\partial x_2} &
    \cdots &
    \dfrac{\partial\varphi_m(x_0)}{\partial x_{n-m}}
  \end{pmatrix}
  \triangleq -A^{-1}B.
\end{aligned}
\label{course-eq-14-5-5}
\end{equation}
注意到
\[
  -
  \left(
    \frac{\partial f(x_0)}{\partial x_{n-m+1}},
    \frac{\partial f(x_0)}{\partial x_{n-m+2}},
    \cdots,
    \frac{\partial f(x_0)}{\partial x_n}
  \right)\cdot A^{-1}
\]
是一个 $m$ 维行向量, 我们可以将其记为
\begin{equation}
  -
  \left(
    \frac{\partial f(x_0)}{\partial x_{n-m+1}},
    \frac{\partial f(x_0)}{\partial x_{n-m+2}},
    \cdots,
    \frac{\partial f(x_0)}{\partial x_n}
  \right)\cdot A^{-1}
  =(\lambda_1,\lambda_2,\cdots,\lambda_m).
  \label{course-eq-14-5-6}
\end{equation}
将式~\eqref{course-eq-14-5-5}  与 \eqref{course-eq-14-5-6}  代入式~\eqref{course-eq-14-5-4}  得
\begin{equation}
\begin{aligned}
  &\left(
    \frac{\partial f(x_0)}{\partial x_1},
    \frac{\partial f(x_0)}{\partial x_2},
    \cdots,
    \frac{\partial f(x_0)}{\partial x_{n-m}}
  \right)\\
  &\quad+
  (\lambda_1,\lambda_2,\cdots,\lambda_m)
  \begin{pmatrix}
    \dfrac{\partial\varphi_1(x_0)}{\partial x_1} &
    \dfrac{\partial\varphi_1(x_0)}{\partial x_2} &
    \cdots &
    \dfrac{\partial\varphi_1(x_0)}{\partial x_{n-m}}\\
    \dfrac{\partial\varphi_2(x_0)}{\partial x_1} &
    \dfrac{\partial\varphi_2(x_0)}{\partial x_2} &
    \cdots &
    \dfrac{\partial\varphi_2(x_0)}{\partial x_{n-m}}\\
    \vdots & \vdots & & \vdots\\
    \dfrac{\partial\varphi_m(x_0)}{\partial x_1} &
    \dfrac{\partial\varphi_m(x_0)}{\partial x_2} &
    \cdots &
    \dfrac{\partial\varphi_m(x_0)}{\partial x_{n-m}}
  \end{pmatrix}
  =0.
\end{aligned}
\label{course-eq-14-5-7}
\end{equation}
另外, 我们可以将式~\eqref{course-eq-14-5-6}  改写成下述形式
\begin{equation}
  \left(
    \frac{\partial f(x_0)}{\partial x_{n-m+1}},
    \frac{\partial f(x_0)}{\partial x_{n-m+2}},
    \cdots,
    \frac{\partial f(x_0)}{\partial x_n}
  \right)
  +
  (\lambda_1,\lambda_2,\cdots,\lambda_m)
  A=0.
  \label{course-eq-14-5-8}
\end{equation}
最后将向量方程 \eqref{course-eq-14-5-7}  与 \eqref{course-eq-14-5-8}  写成分量形式的方程, 再加上约束条件 \eqref{course-eq-14-5-1}, 即得式~\eqref{course-eq-14-5-2}. 证毕。

\end{proof}
\begin{remark}\quad 若构造函数
\[
  F(x_1,\cdots,x_n,\lambda_1,\cdots,\lambda_m)
  =
  f(x)+\sum_{j=1}^m\lambda_j\varphi_j(x),
\]
则上述条件极值的必要条件形式上化为 $F$ 的通常极值的必要条件
\begin{equation}
  \begin{cases}
    \dfrac{\partial F(x_0)}{\partial x_i}=0
    \qquad (i=1,2,\cdots,n),\\[1.2ex]
    \dfrac{\partial F(x_0)}{\partial\lambda_j}=0
    \qquad (j=1,2,\cdots,m).
  \end{cases}
\label{course-eq-14-5-9}
\end{equation}
上述求条件极值点的必要条件的方法称为拉格朗日乘数法。

由于在拉格朗日乘数法中, 条件极值的驻点方程可借助
\[
  F(x)=f(x)+\sum_{j=1}^m\lambda_j\varphi_j(x)
\]
的驻点方程来表达，并不意味着对乘子也取通常极值。而当 $x_0$ 是驻点时, $\lambda_j(j=1,2,\cdots,m)$ 是 $m$ 个常数, 因此我们利用 $F(x)$ 在 $x_0$ 处的海色矩阵 $H_F(x_0)$ 来判定 $x_0$ 是否为极值点. 可以证明: 当 $H_F(x_0)$ 正定时, 条件极值点为极小值点; 当 $H_F(x_0)$ 负定时, 条件极值为极大值点。

值得提出的是, 当 $H_F(x_0)$ 不定时, 条件极值仍可以取到. 例如, 设函数 $f(x,y)=x^2-y^2$, $\varphi(x,y)=y=0$, 则 $\varphi(x,y)=0$ 确定直线方程 $y=0$, 显然 $f(x,0)=x^2$ 在 $(0,0)$ 处取极小值 $0$. 容易算出
\[
  F(x,y)=x^2-y^2+\lambda y
\]
的海色矩阵在 $(0,0)$ 处为
\[
  \begin{pmatrix}
    2 & 0\\
    0 & -2
  \end{pmatrix},
\]
它是一个不定矩阵。

\end{remark}
\begin{example}\label{course-example-14-5-4}\quad 求椭球面 $16x^2+4y^2+9z^2=144$ 内接且各面均平行于坐标平面的长方体的最大体积。

\end{example}
\begin{solve}\quad 设 $(x,y,z)$ 是该椭球面内接长方体在第一卦限内的顶点, 由对称性, 该长方体的体积为
\[
  V=f(x,y,z)=8xyz.
\]
因此所求的最大体积转化为求 $f(x,y,z)$ 在约束条件
\[
  g(x,y,z)=16x^2+4y^2+9z^2-144=0
\]
下的最大值. 由拉格朗日乘数法, 我们作函数
\[
  F(x,y,z,\lambda)=f(x,y,z)+\lambda g(x,y,z).
\]
求 $F(x,y,z,\lambda)$ 的各个偏导数得方程组
\[
  \begin{cases}
    8yz+32\lambda x=0,\\
    8xz+8\lambda y=0,\\
    8xy+18\lambda z=0,\\
    16x^2+4y^2+9z^2-144=0.
  \end{cases}
  \label{course-eq-lagrange-box}
\]
由上面方程组易推出
\begin{equation}
  xyz=-12\lambda.
  \label{course-eq-14-5-14}
\end{equation}
将式~\eqref{course-eq-14-5-14}  代入式 (14.5.10) 得
\[
  -32\lambda(3-x^2)=0.
\]
因此得 $\lambda=0$ 或 $x=\sqrt3$. 当 $\lambda=0$ 时得 $xyz=0$. 显然不合题意. 因此必有 $\lambda\ne0$, 从而有 $x=\sqrt3$。

同理将式~\eqref{course-eq-14-5-14}  分别代入式 (14.5.11) 和 (14.5.12) 可得 $y=2\sqrt3$ 和 $z=4\sqrt3/3$。

由于在该问题中最大值总是存在的, 因此所求最大值为
\[
  V_{\max}=8\times\sqrt3\times2\sqrt3\times\frac{4\sqrt3}{3}=64\sqrt3.
\]

\end{solve}
\begin{example}\label{course-example-14-5-5}\quad 设 $n$ 个正数之和为 $c$, 试求它们乘积的最大值, 并证明对任何正数 $a_1,a_2,\cdots,a_n$, 有
\[
  \sqrt[n]{a_1a_2\cdots a_n}
  \leqslant
  \frac{a_1+a_2+\cdots+a_n}{n}.
\]

\end{example}
\begin{solve}\quad 在 $E=\{(x_1,x_2,\cdots,x_n):x_i\geqslant0,\ i=1,2,\cdots,n\}$ 上定义函数
\[
  f(x_1,x_2,\cdots,x_n)=\prod_{i=1}^n x_i.
\]
依题意, 我们要求 $f(x_1,x_2,\cdots,x_n)$ 在约束条件
\[
  \sum_{i=1}^n x_i=c
\]
下的最大值。

作辅助函数
\[
  F(x_1,x_2,\cdots,x_n,\lambda)
  =
  \prod_{i=1}^n x_i
  +\lambda\left(\sum_{i=1}^n x_i-c\right).
\]
求 $F$ 的 $n$ 个偏导数及 $\dfrac{\partial F}{\partial\lambda}$, 并令其为零, 得方程组
\[
  \begin{cases}
    x_2x_3\cdots x_n+\lambda=0,\\
    x_1x_3\cdots x_n+\lambda=0,\\
    \cdots\\
    x_1x_2\cdots x_{n-1}+\lambda=0,\\
    \displaystyle\sum_{i=1}^n x_i=c,
  \end{cases}
\]
解出
\[
  x_1=x_2=\cdots=x_n=\frac cn,\qquad
  \lambda=-\left(\frac cn\right)^{n-1}.
\]
我们知道 $f(x_1,x_2,\cdots,x_n)$ 在紧集 $E$ 上取到最大值. 另外, 容易看出, 若 $(x_1^0,x_2^0,\cdots,x_n^0)$ 是 $f(x_1,x_2,\cdots,x_n)$ 的最大值点, 必有 $x_i^0\ne0(i=1,2,\cdots,n)$. 由于
\[
  \left(\frac cn,\frac cn,\cdots,\frac cn,-\left(\frac cn\right)^{n-1}\right)
\]
是 $F(x_1,x_2,\cdots,x_n,\lambda)$ 的唯一驻点, 因此它为 $F(x_1,x_2,\cdots,x_n,\lambda)$ 的极大值点, 从而 $f(x_1,x_2,\cdots,x_n)$ 在该点取最大值
\[
  \frac{c^n}{n^n}.
\]
现对 $n$ 个正数 $a_1,a_2,\cdots,a_n$, 令 $c=\displaystyle\sum_{i=1}^n a_i$, 并利用上述函数 $f(x)$ 所证结论, 则对任何满足 $\displaystyle\sum_{i=1}^n x_i=c$ 的正数 $x_i$ 均有
\[
  \prod_{i=1}^n x_i\leqslant\frac{c^n}{n^n}.
\]
取 $x_i=a_i(i=1,2,\cdots,n)$, 则有
\[
  \sqrt[n]{a_1a_2\cdots a_n}
  \leqslant
  \frac{a_1+a_2+\cdots+a_n}{n}.
\]

\end{solve}


\section{化重积分为累次积分}\label{course-15-3}
本节用于把竞赛中的积分区域转写为可计算的积分限。换序保持区域不变，换元须同时改变区域和面积、体积元；奇异边界处应通过截断或可积性论证解释极限。

我们在前面介绍了 $\mathbb R^n$ 中可求体积的有界闭区域上的重积分理论. 在本节及下节中, 我们将讨论重积分的计算问题. 我们主要介绍二重积分与三重积分的计算。

\subsection{化二重积分为累次积分}

在讨论一般可求体积的有界闭区域上的二重积分的计算之前, 我们先来讨论一下矩形区域上的二重积分。

\begin{theorem}\label{course-theorem-15-3-1}\quad 设函数 $f(x,y)$ 在 $D=[a,b]\times[c,d]$ 上可积, 且对于 $\forall x\in[a,b]$, $I(x)=\int_c^d f(x,y)\mathrm dy$ 存在, 则定积分 $\int_a^b I(x)\mathrm dx$ 存在, 并且
\[
  \iint_D f(x,y)\mathrm dx\mathrm dy
  =\int_a^b I(x)\mathrm dx
  \triangleq\int_a^b\mathrm dx\int_c^d f(x,y)\mathrm dy.
\]

\end{theorem}
\begin{proof}\quad 对 $[a,b]$ 作分割
\[
  \Delta_x:a=x_0<x_1<\cdots<x_J=b,
\]
并记 $\Delta x_j=x_j-x_{j-1}(j=1,2,\cdots,J)$; 对 $[c,d]$ 作分割
\[
  \Delta_y:c=y_0<y_1<\cdots<y_K=d,
\]
并记 $\Delta y_k=y_k-y_{k-1}(k=1,2,\cdots,K)$. 过点 $(x_j,0)(j=1,2,\cdots,J)$ 且平行于 $y$ 轴的直线以及过点 $(0,y_k)(k=1,2,\cdots,K)$ 且平行于 $x$ 轴的直线将 $D$ 分成了一些小矩形, 它们组成了 $D$ 的一个分割:
\[
  \Delta=\{\Delta D_{jk}:1\leqslant j\leqslant J,\ 1\leqslant k\leqslant K\}.
\]
显然 $\lambda(\Delta)\to0$ 的充分必要条件是 $\lambda(\Delta_x)\to0$ 和 $\lambda(\Delta_y)\to0$. 在 $[x_{j-1},x_j]$ 上任取 $\xi_j$, 在 $[y_{k-1},y_k]$ 上任取 $\eta_k$, 则 $(\xi_j,\eta_k)\in\Delta D_{jk}(1\leqslant j\leqslant J,1\leqslant k\leqslant K)$。

记
\[
  m_{jk}=\inf_{(x,y)\in\Delta D_{jk}}f(x,y),\quad
  M_{jk}=\sup_{(x,y)\in\Delta D_{jk}}f(x,y),
\]
则有
\[
  \sum_{j=1}^J\sum_{k=1}^K m_{jk}\Delta x_j\Delta y_k
  \leqslant \sum_{j=1}^J\sum_{k=1}^K f(\xi_j,\eta_k)\Delta x_j\Delta y_k
  \leqslant \sum_{j=1}^J\sum_{k=1}^K M_{jk}\Delta x_j\Delta y_k,
\]
从而有
\[
\begin{aligned}
  \sum_{j=1}^J\sum_{k=1}^K m_{jk}\Delta x_j\Delta y_k
  &\leqslant \sum_{j=1}^J\left(\int_c^d f(\xi_j,y)\mathrm dy\right)\Delta x_j\\
  &=\sum_{j=1}^J I(\xi_j)\Delta x_j
  \leqslant \sum_{j=1}^J\sum_{k=1}^K M_{jk}\Delta x_j\Delta y_k.
\end{aligned}
\]
由于
\[
  \lim_{\lambda(\Delta)\to0}\sum_{j=1}^J\sum_{k=1}^K m_{jk}\Delta x_j\Delta y_k
  =
  \lim_{\lambda(\Delta)\to0}\sum_{j=1}^J\sum_{k=1}^K M_{jk}\Delta x_j\Delta y_k
  =
  \iint_D f(x,y)\mathrm dx\mathrm dy,
\]
从定积分的定义知, $I(x)$ 在 $[a,b]$ 上可积, 并且
\[
  \int_a^b I(x)\mathrm dx
  =
  \lim_{\lambda(\Delta_x)\to0}\sum_{j=1}^J I(\xi_j)\Delta x_j
  =
  \iint_D f(x,y)\mathrm dx\mathrm dy.
\]
证毕。

\end{proof}
\noindent\keyterm{注 1}\quad 若 $f(x,y)$ 在 $D=[a,b]\times[c,d]$ 上可积, 且对于 $\forall y\in[c,d]$, $J(y)=\int_a^b f(x,y)\mathrm dx$ 存在, 则类似地可以证明 $\int_c^d J(y)\mathrm dy$ 存在, 并且
\[
  \iint_D f(x,y)\mathrm dx\mathrm dy
  =
  \int_c^d\mathrm dy\int_a^b f(x,y)\mathrm dx.
\]
今后我们称形如 $\int_a^b\mathrm dx\int_c^d f(x,y)\mathrm dy$ 或 $\int_c^d\mathrm dy\int_a^b f(x,y)\mathrm dx$ 的积分为累次积分。

\noindent\keyterm{注 2}\quad 在矩形区域上的二重积分相当于一个二元函数的极限, 而累次积分相当于一个二元函数的累次极限. 因此, 类似于极限与累次极限的关系, 当一个二重积分存在时, 我们不能保证该积分能用累次积分求出. 这样的例子可以容易举出. 例如设 $D=[0,1]\times[0,1]$, 当 $(x,y)\in D$ 时, 令
\[
  f(x,y)=
  \begin{cases}
    \dfrac1k, & x=\dfrac1k(k\in\mathbb N),\ y\text{ 为有理数},\\
    0, & \text{其他}.
  \end{cases}
\]
我们来说明 $f(x,y)$ 在 $D$ 上的二重积分不能通过累次积分求出. 给定 $\varepsilon>0$，先取 $K$ 使 $1/K<\varepsilon$。在 $x=1,1/2,\ldots,1/K$ 的有限条竖线外取很窄的条带，总面积小于 $\varepsilon$；条带中 $|f|\le1$，条带外 $|f|\le1/K$。于是上积分不超过 $2\varepsilon$，下积分为零，故 $f$ 黎曼可积且二重积分为零。 但当 $x=\dfrac1k$ 时, $f\left(\dfrac1k,y\right)=\dfrac1k\mathrm D(y)$, 其中 $\mathrm D(y)$ 是 $[0,1]$ 上的狄利克雷函数, 从而它在 $[0,1]$ 上不可积。

\begin{example}\label{course-example-15-3-1}\quad 设 $D=[0,1]\times[0,1]$, 计算
\[
  I=\iint_D x(x-y)^2\mathrm dx\mathrm dy.
\]

\end{example}
\begin{solve}\quad 由于被积函数 $x(x-y)^2$ 在 $D$ 上连续, 因此
\[
\begin{aligned}
  I
  &=\int_0^1 x\mathrm dx\int_0^1(x-y)^2\mathrm dy
    =\int_0^1 x\left[-\frac13(x-y)^3\right]_0^1\mathrm dx\\
  &=\frac13\int_0^1x[x^3-(x-1)^3]\mathrm dx
    =\frac1{12}.
\end{aligned}
\]

下面我们讨论 $\mathbb R^2$ 中两类特殊有界闭区域上的二重积分问题. 在此我们先来复习一下 X 型区域与 Y 型区域的概念. 若区域
\[
  D=\{(x,y):a\leqslant x\leqslant b,\ \varphi_1(x)\leqslant y\leqslant\varphi_2(x)\},
\]
其中 $\varphi_1(x)$ 与 $\varphi_2(x)$ 是区间 $[a,b]$ 上的连续函数, 且对于 $\forall x\in(a,b)$, 有 $\varphi_1(x)<\varphi_2(x)$, 则称 $D$ 为 X 型区域。

若
\[
  D=\{(x,y):c\leqslant y\leqslant d,\ \psi_1(y)\leqslant x\leqslant\psi_2(y)\},
\]
其中 $\psi_1(y)$ 与 $\psi_2(y)$ 是区间 $[c,d]$ 上的连续函数, 且对于 $\forall y\in(c,d)$, 有 $\psi_1(y)<\psi_2(y)$, 则称 $D$ 为 Y 型区域。

显然任何 X 型区域与 Y 型区域都是可求面积的. 对于此类区域上的二重积分, 我们有以下结论。

\end{solve}
\begin{theorem}\label{course-theorem-15-3-2}\quad 设 $D=\{(x,y):a\leqslant x\leqslant b,\varphi_1(x)\leqslant y\leqslant\varphi_2(x)\}$ 是 X 型区域, 函数 $f(x,y)$ 在 $D$ 上可积, 且对于 $\forall x\in[a,b]$,
\[
  I(x)=\int_{\varphi_1(x)}^{\varphi_2(x)}f(x,y)\mathrm dy
\]
存在, 则定积分 $\int_a^b I(x)\mathrm dx$ 存在, 且
\begin{equation}
  \iint_D f(x,y)\mathrm dx\mathrm dy
  =\int_a^b I(x)\mathrm dx
  =\int_a^b\mathrm dx\int_{\varphi_1(x)}^{\varphi_2(x)}f(x,y)\mathrm dy.
  \label{course-eq-15-3-1}
\end{equation}

\end{theorem}
\begin{proof}\quad 由于函数 $\varphi_1(x)$ 与 $\varphi_2(x)$ 在 $[a,b]$ 上连续, 我们取
\[
  c=\min_{x\in[a,b]}\varphi_1(x)-1,\qquad
  d=\max_{x\in[a,b]}\varphi_2(x)+1,
\]
则 $D\subset D_1=[a,b]\times[c,d]$. 记
\[
  \tilde f(x,y)=
  \begin{cases}
    f(x,y), & (x,y)\in D,\\
    0, & (x,y)\in D_1\setminus D.
  \end{cases}
\]
由重积分的性质知, $\tilde f(x,y)$ 在 $D_1$ 上可积, 且对于 $\forall x\in[a,b]$,
\[
  \int_c^d\tilde f(x,y)\mathrm dy
  =\int_{\varphi_1(x)}^{\varphi_2(x)}f(x,y)\mathrm dy.
\]
因此
\[
  \iint_D f(x,y)\mathrm dx\mathrm dy
  =\iint_{D_1}\tilde f(x,y)\mathrm dx\mathrm dy
  =\int_a^b\mathrm dx\int_{\varphi_1(x)}^{\varphi_2(x)}f(x,y)\mathrm dy.
\]
证毕。

\end{proof}
\begin{remark}\quad 当 $f(x,y)$ 在 Y 型区域
\[
  D=\{(x,y):c\leqslant y\leqslant d,\ \psi_1(y)\leqslant x\leqslant\psi_2(y)\}
\]
上可积, 且对于 $\forall y\in[c,d]$, $J(y)=\int_{\psi_1(y)}^{\psi_2(y)}f(x,y)\mathrm dx$ 存在时, 则定积分 $\int_c^dJ(y)\mathrm dy$ 存在, 且
\begin{equation}
  \iint_D f(x,y)\mathrm dx\mathrm dy
  =\int_c^dJ(y)\mathrm dy
  =\int_c^d\mathrm dy\int_{\psi_1(y)}^{\psi_2(y)}f(x,y)\mathrm dx.
  \label{course-eq-15-3-2}
\end{equation}

通常我们称 \eqref{course-eq-15-3-1}  与 \eqref{course-eq-15-3-2}  右边的积分为重积分 $\iint_D f(x,y)\mathrm dx\mathrm dy$ 的累次积分. 另外, 当 $D$ 是一般有界闭区域时, 若添加有限条光滑曲线可以将 $D$ 分成有限个 X 型或 Y 型区域, 我们也可以通过累次积分来计算 $D$ 上的二重积分。

\end{remark}
\begin{example}\label{course-example-15-3-2}\quad 设有界闭区域 $D\subset\mathbb R^2$ 由直线 $y=0$ 及曲线 $y=x^3,x+y=2$ 所围成, $f(x,y)$ 是 $D$ 上的连续函数, 试用两种不同的积分顺序将二重积分 $\iint_D f(x,y)\mathrm dx\mathrm dy$ 化为累次积分。

\end{example}
\begin{solve}\quad 由题给积分区域的不等式，我们有
\[
\begin{aligned}
  \iint_D f(x,y)\mathrm dx\mathrm dy
  &=\int_0^1\mathrm dy\int_{y^{1/3}}^{2-y}f(x,y)\mathrm dx\\
  &=\int_0^1\mathrm dx\int_0^{x^3}f(x,y)\mathrm dy
    +\int_1^2\mathrm dx\int_0^{2-x}f(x,y)\mathrm dy.
\end{aligned}
\]

从上例中可以看出, 选择不同的积分顺序, 计算的难易程度可能会有所差别。

\end{solve}
\begin{example}\label{course-example-15-3-3}\quad 计算累次积分
\[
  \int_0^1\mathrm dx\int_x^{\sqrt x}\frac{\sin y}{y}\mathrm dy.
\]

\end{example}
\begin{solve}\quad 由于 $\dfrac{\sin y}{y}$ 的原函数无法求出, 累次积分不能直接计算. 为此考虑函数
\[
  f(x,y)=
  \begin{cases}
    \dfrac{\sin y}{y}, & y\ne0,\\
    1, & y=0
  \end{cases}
\]
在平面区域 $D=\{(x,y):0\leqslant x\leqslant1,\ x\leqslant y\leqslant\sqrt x\}$ 上的二重积分。

由于 $f(x,y)$ 在 $D$ 上连续, 因此该二重积分可以化为累次积分. 我们有
\[
\begin{aligned}
  \int_0^1\mathrm dx\int_x^{\sqrt x}\frac{\sin y}{y}\mathrm dy
  &=\iint_D\frac{\sin y}{y}\mathrm dx\mathrm dy
    =\int_0^1\mathrm dy\int_{y^2}^{y}\frac{\sin y}{y}\mathrm dx\\
  &=\int_0^1\frac{(y-y^2)\sin y}{y}\mathrm dy
    =1-\sin1.
\end{aligned}
\]
从例~\ref{course-example-15-3-3}  可以看出, 对于一些二重积分, 其中一个累次积分很难计算出二重积分的值, 而另一个累次积分却能容易地算出二重积分的值。

\end{solve}
\begin{example}\label{course-example-15-3-4}\quad 计算 $\mathbb R^3$ 中以柱面 $x^2+y^2=a^2$ 与 $x^2+z^2=a^2(a>0)$ 为边界且含有原点的立体体积 $V$。

\end{example}
\begin{solve}\quad 由对称性, 该立体在八个卦限中都具有相等的体积. 在第一卦限中, 该立体是以区域
\[
  D=\{(x,y):x^2+y^2\leqslant a^2,\ x\geqslant0,\ y\geqslant0\}
\]
为底, 以曲面 $z=\sqrt{a^2-x^2}$ 为顶的曲顶柱体。
由二重积分的几何意义知
\[
  V=8\iint_D z\mathrm dx\mathrm dy.
\]
化二重积分为累次积分即得
\[
\begin{aligned}
  V
  &=8\iint_D z\mathrm dx\mathrm dy
    =8\int_0^a\mathrm dx\int_0^{\sqrt{a^2-x^2}}\sqrt{a^2-x^2}\mathrm dy\\
  &=8\int_0^a(a^2-x^2)\mathrm dx
    =\frac{16}{3}a^3.
\end{aligned}
\]

\end{solve}
\subsection{化三重积分为累次积分}

对于三重积分, 当其积分区域具有特殊形状时, 我们可以将三重积分化为累次积分。

设 $\Omega\subset\mathbb R^2$ 是可求面积的有界闭区域, $\varphi(x,y),\psi(x,y),(x,y)\in\Omega$ 是 $\Omega$ 上的连续函数, 且对于 $\forall(x,y)\in\Omega^\circ$, 有 $\varphi(x,y)<\psi(x,y)$, 记
\[
  D=\{(x,y,z):(x,y)\in\Omega,\ \varphi(x,y)\leqslant z\leqslant\psi(x,y)\}
\]
, 则容易推出 $D$ 是可求体积的闭区域(见本章习题). 再设函数 $f(x,y,z)$ 在 $D$ 上连续, 则
\[
  I(x,y)=\int_{\varphi(x,y)}^{\psi(x,y)}f(x,y,z)\mathrm dz
\]
在 $\Omega$ 上可积, 且有
\[
  \iiint_D f(x,y,z)\mathrm dV
  =\iint_\Omega I(x,y)\mathrm dx\mathrm dy
  \triangleq \iint_\Omega\mathrm dx\mathrm dy
    \int_{\varphi(x,y)}^{\psi(x,y)}f(x,y,z)\mathrm dz.
\]
若三重积分 $\iiint_D f(x,y,z)\mathrm dV$ 的积分区域 $D$ 具有如下两种形式:
\[
  D=\{(x,y,z):(y,z)\in\Omega,\ \varphi(y,z)\leqslant x\leqslant\psi(y,z)\},
\]
\[
  D=\{(x,y,z):(x,z)\in\Omega,\ \varphi(x,z)\leqslant y\leqslant\psi(x,z)\},
\]
类似地, 我们也可以将其化为累次积分。

现在设空间有界闭区域 $D\subset\mathbb R^3$ 的边界是由有限块分片光滑曲面所组成, 并且当 $(x,y,z)\in D$ 时, 有 $a\leqslant z\leqslant b$, 又设对于 $\forall z\in[a,b]$, 过点 $(0,0,z)$ 且与 $Oxy$ 平面平行的平面截 $D$ 所得到的截面是可求面积的平面闭区域 $D_z$, 则当函数 $f(x,y,z)$ 在 $D$ 上连续时, 有
\[
  \iiint_D f(x,y,z)\mathrm dV
  =\int_a^b\left(\iint_{D_z}f(x,y,z)\mathrm dx\mathrm dy\right)\mathrm dz
  \triangleq \int_a^b\mathrm dz\iint_{D_z}f(x,y,z)\mathrm dx\mathrm dy.
\]
当用平行于 $Ozx$ 平面或 $Oyz$ 平面的平面截 $D$ 所得的截面为可求面积的平面闭区域时, 我们也有类似的公式。

\begin{example}\label{course-example-15-3-5}\quad 计算三重积分
\[
  I=\iiint_D\frac{\mathrm dx\mathrm dy\mathrm dz}{(1+x+y+z)^3},
\]
其中 $D$ 是由平面 $x=0,y=0,z=0$ 及 $x+y+z=1$ 所围成的闭区域。

\end{example}
\begin{solve}\quad 记 $\Omega\subset\mathbb R^2$ 为由直线 $x=0,y=0$ 及 $x+y=1$ 所围的闭区域, 则
\[
  D=\{(x,y,z):0\leqslant z\leqslant1-x-y,\ (x,y)\in\Omega\}.
\]
于是有
\[
\begin{aligned}
  I
  &=\iint_\Omega\mathrm dx\mathrm dy\int_0^{1-x-y}\frac{\mathrm dz}{(1+x+y+z)^3}\\
  &=\int_0^1\mathrm dx\int_0^{1-x}\mathrm dy\int_0^{1-x-y}\frac{\mathrm dz}{(1+x+y+z)^3}\\
  &=\int_0^1\mathrm dx\int_0^{1-x}
    \left.\frac{-1}{2(1+x+y+z)^2}\right|_0^{1-x-y}\mathrm dy\\
  &=\int_0^1\mathrm dx\int_0^{1-x}\frac12\left[\frac1{(1+x+y)^2}-\frac14\right]\mathrm dy\\
  &=\frac12\int_0^1\left.\left(\frac{-1}{1+x+y}-\frac y4\right)\right|_0^{1-x}\mathrm dx\\
  &=\frac12\int_0^1\left(\frac1{1+x}-\frac{3-x}{4}\right)\mathrm dx
    =\frac12\ln2-\frac5{16}.
\end{aligned}
\]

\end{solve}
\begin{example}\label{course-example-15-3-6}\quad 计算三重积分
\[
  I=\iiint_D (x+y+z)^2\mathrm dx\mathrm dy\mathrm dz,
\]
其中
\[
  D=\left\{(x,y,z):\frac{x^2}{a^2}+\frac{y^2}{b^2}+\frac{z^2}{c^2}\leqslant1,\ a,b,c>0\right\}.
\]

\end{example}
\begin{solve}\quad 由积分区域的对称性以及被积函数的对称性容易看出
\[
  \iiint_D xy\mathrm dx\mathrm dy\mathrm dz
  =\iiint_D xz\mathrm dx\mathrm dy\mathrm dz
  =\iiint_D yz\mathrm dx\mathrm dy\mathrm dz=0,
\]
因此
\[
  I=\iiint_D (x^2+y^2+z^2)\mathrm dx\mathrm dy\mathrm dz.
\]
对任意 $x\in[-a,a]$, 过点 $(x,0,0)$ 且平行于 $Oyz$ 平面的平面与 $D$ 的交为
\[
  D_x=\left\{(y,z):\frac{y^2}{b^2\left(1-\frac{x^2}{a^2}\right)}
  +\frac{z^2}{c^2\left(1-\frac{x^2}{a^2}\right)}\leqslant1\right\},
\]
其面积为 $\pi bc\left(1-\dfrac{x^2}{a^2}\right)$. 因此有
\[
\begin{aligned}
  \iiint_D x^2\mathrm dx\mathrm dy\mathrm dz
  &=\int_{-a}^a x^2\mathrm dx\iint_{D_x}\mathrm dy\mathrm dz\\
  &=\int_{-a}^a x^2\pi bc\left(1-\frac{x^2}{a^2}\right)\mathrm dx
    =\frac4{15}\pi a^3bc.
\end{aligned}
\]
同理可计算出
\[
  \iiint_D y^2\mathrm dx\mathrm dy\mathrm dz=\frac4{15}\pi ab^3c,\qquad
  \iiint_D z^2\mathrm dx\mathrm dy\mathrm dz=\frac4{15}\pi abc^3.
\]
最后我们有
\[
  I=\frac4{15}\pi abc(a^2+b^2+c^2).
\]

\end{solve}


\section{重积分的变量替换}\label{course-15-4}
本节用于把竞赛中的积分区域转写为可计算的积分限。换序保持区域不变，换元须同时改变区域和面积、体积元；奇异边界处应通过截断或可积性论证解释极限。

\subsection{重积分的变量替换公式}

在本节中, 我们主要证明二重积分的变量替换公式. 对于 $n(n>2)$ 重积分, 我们将不加证明地给出其变量替换公式。

设
\[
  T(u,v):
  \begin{cases}
    x=x(u,v),\\
    y=y(u,v)
  \end{cases}
\]
是可求面积的有界闭区域 $D$ 到可求面积的有界闭区域 $\Omega$ 的一个 $C^1$ 同胚映射, 并且假定对于 $\forall(u,v)\in D$ 有
\[
  \frac{\partial(x,y)}{\partial(u,v)}\ne0.
\]
由此假定与偏导数的连续性知, $\dfrac{\partial(x,y)}{\partial(u,v)}$ 在 $D$ 上不改变符号. 注意到向量函数的导数定义, $T(u,v)$ 在 $(u,v)\in D$ 处导数的行列式即为 $\dfrac{\partial(x,y)}{\partial(u,v)}$。

在一元微积分中, 当 $f'(x)$ 在区间 $[a,b]$ 上连续且不变号时, $f(x)$ 在 $[a,b]$ 上单调, 从而
\[
  \int_a^b|f'(x)|\mathrm dx
  =\left|\int_a^bf'(x)\mathrm dx\right|
  =|f(b)-f(a)|.
\]
这说明 $\int_a^b|f'(x)|\mathrm dx$ 是函数 $y=f(x)$ 的值域的“长度”. 因此, 我们有理由猜测
\[
  \iint_D\left|\frac{\partial(x,y)}{\partial(u,v)}\right|\mathrm du\mathrm dv
\]
应该是 $\Omega$ 的面积。

再从微元法的观点来看, 在点 $(u_0,v_0)\in D$ 处给一增量 $(\mathrm du,\mathrm dv)$, 则以 $(u_0,v_0)$, $(u_0+\mathrm du,v_0)$, $(u_0+\mathrm du,v_0+\mathrm dv)$, $(u_0,v_0+\mathrm dv)$ 为顶点的小矩形的面积为 $\mathrm du\mathrm dv$. 而这个小矩形在变换
\[
  T(u,v):
  \begin{cases}
    x=x(u,v),\\
    y=y(u,v)
  \end{cases}
\]
下近似地映成一个平行四边形, 其四个顶点分别为
\[
\begin{aligned}
  P_1&=(x(u_0,v_0),y(u_0,v_0)),\\
  P_2&=(x(u_0+\mathrm du,v_0),y(u_0+\mathrm du,v_0)),\\
  P_3&=(x(u_0+\mathrm du,v_0+\mathrm dv),y(u_0+\mathrm du,v_0+\mathrm dv)),\\
  P_4&=(x(u_0,v_0+\mathrm dv),y(u_0,v_0+\mathrm dv)).
\end{aligned}
\]
由于 $T(u,v)$ 是 $C^1$ 映射, 因此有
\[
  |\overrightarrow{P_1P_2}|\approx
  \left|\frac{\partial x(u_0,v_0)}{\partial u}\boldsymbol i
  +\frac{\partial y(u_0,v_0)}{\partial u}\boldsymbol j\right|\mathrm du,
\]
\[
  |\overrightarrow{P_1P_4}|\approx
  \left|\frac{\partial x(u_0,v_0)}{\partial v}\boldsymbol i
  +\frac{\partial y(u_0,v_0)}{\partial v}\boldsymbol j\right|\mathrm dv.
\]
该近似平行四边形的面积约为
\[
  |\overrightarrow{P_1P_2}\times\overrightarrow{P_1P_4}|
  \approx
  \left|\frac{\partial(x,y)}{\partial(u,v)}\right|\mathrm du\mathrm dv.
\]
这也说明了 $\dfrac{\partial(x,y)}{\partial(u,v)}$ 起着一元函数 $f(x)$ 的导数相似的作用. 下面我们将给出上述结论的严格证明, 并推出相应的二重积分的变量替换公式。

\begin{theorem}\label{course-theorem-15-4-1}\quad 设 $D\subset\mathbb R^2$ 是有界闭区域, $\partial D$ 由有限条分段光滑曲线所组成, 变换
\[
  T(u,v):
  \begin{cases}
    x=x(u,v),\\
    y=y(u,v)
  \end{cases}
\]
是 $D$ 到有界闭区域 $\Omega$ 的 $C^1$ 同胚映射, 并且在 $D$ 上处处有
\[
  \frac{\partial(x,y)}{\partial(u,v)}\ne0,
\]
再设函数 $f(x,y)$ 在 $\Omega$ 上可积, 则
\[
  \iint_\Omega f(x,y)\mathrm dx\mathrm dy
  =\iint_D f(x(u,v),y(u,v))
  \left|\frac{\partial(x,y)}{\partial(u,v)}\right|\mathrm du\mathrm dv.
\]

为了证明定理~\ref{course-theorem-15-4-1}, 我们先做一些准备工作. 设 $D_0$ 为 $D$ 内的一个正方形, 它的四个顶点为 $(u_0,v_0)$, $(u_0+\delta,v_0)$, $(u_0+\delta,v_0+\delta)$, $(u_0,v_0+\delta)$, 其中 $\delta>0$. 同胚映射 $T$ 将 $D_0$ 映成了区域 $\Omega$ 内的一个由四条光滑曲线组成的一个曲边四边形 $T(D_0)$. 如果我们进一步假定变换 $T$ 的四个偏导数 $x'_u,x'_v,y'_u,y'_v$ 中至少有一个在 $D_0$ 上恒大(小)于零, 取 $D_0$ 的某一条对角线将 $D_0$ 分成两个三角形, 则不难看出变换 $T$ 将这两个三角形映成两个 X 型区域(或两个 Y 型区域). 下面我们先证明一个引理。

\end{theorem}
\begin{lemma}\label{course-lemma-15-4-2}\quad 假设变换
\[
  T(u,v):
  \begin{cases}
    x=x(u,v),\\
    y=y(u,v)
  \end{cases}
\]
在 $D$ 内具有二阶连续偏导数, $D_0$ 为 $D$ 内的一个正方形并且 $T(D_0)$ 可以分解成有限个由分段光滑曲线所围成的 X 型区域(或 Y 型区域), 则 $T(D_0)$ 的面积
\[
  \sigma(T(D_0))=\iint_{D_0}\left|\frac{\partial(x,y)}{\partial(u,v)}\right|\mathrm du\mathrm dv.
\]

\end{lemma}
\begin{proof}\quad 在定积分的几何应用中, 我们知道, 对于 $T(D_0)$ 的面积, 有以下的计算公式:
\[
  \sigma(T(D_0))=-\int_{\partial T(D_0)}y\mathrm dx
  =\int_{\partial T(D_0)}x\mathrm dy
  =\frac12\int_{\partial T(D_0)}x\mathrm dy-y\mathrm dx,
\]
其中 $\partial T(D_0)$ 的方向取正向. 曲线 $\partial T(D_0)$ 的方程可以由参数 $(u,v)$ 在 $D_0$ 四条边上的变化来描述. 先假设 $T$ 将 $\partial D_0$ 的正向映成了 $\partial T(D_0)$ 的正向, 这时我们有
\[
\begin{aligned}
  -\int_{\partial T(D_0)}y\mathrm dx
  &=-\int_{u_0}^{u_0+\delta}y(u,v_0)\frac{\partial x(u,v_0)}{\partial u}\mathrm du
    -\int_{v_0}^{v_0+\delta}y(u_0+\delta,v)\frac{\partial x(u_0+\delta,v)}{\partial v}\mathrm dv\\
  &\quad-\int_{u_0+\delta}^{u_0}y(u,v_0+\delta)\frac{\partial x(u,v_0+\delta)}{\partial u}\mathrm du
    -\int_{v_0+\delta}^{v_0}y(u_0,v)\frac{\partial x(u_0,v)}{\partial v}\mathrm dv\\
  &=\int_{u_0}^{u_0+\delta}
    \left[y(u,v_0+\delta)\frac{\partial x(u,v_0+\delta)}{\partial u}
    -y(u,v_0)\frac{\partial x(u,v_0)}{\partial u}\right]\mathrm du\\
  &\quad-\int_{v_0}^{v_0+\delta}
    \left[y(u_0+\delta,v)\frac{\partial x(u_0+\delta,v)}{\partial v}
    -y(u_0,v)\frac{\partial x(u_0,v)}{\partial v}\right]\mathrm dv\\
  &=\iint_{D_0}\left[
    \frac{\partial}{\partial v}\left(y\frac{\partial x}{\partial u}\right)
    -\frac{\partial}{\partial u}\left(y\frac{\partial x}{\partial v}\right)
    \right]\mathrm du\mathrm dv\\
  &=\iint_{D_0}\frac{\partial(x,y)}{\partial(u,v)}\mathrm du\mathrm dv.
\end{aligned}
\]
当 $T$ 将 $\partial D_0$ 的正向映成了 $\partial T(D_0)$ 的负向时, 由上述方法同样可以证明
\[
  \sigma(T(D_0))=-\iint_{D_0}\frac{\partial(x,y)}{\partial(u,v)}\mathrm du\mathrm dv
  =\iint_{D_0}\left|\frac{\partial(x,y)}{\partial(u,v)}\right|\mathrm du\mathrm dv.
\]
证毕。

\end{proof}
\begin{remark}\quad 在引理~\ref{course-lemma-15-4-2}  中, 我们对 $x(u,v),y(u,v)$ 假定了它们具有二阶连续偏导数. 当 $x(u,v)$ 与 $y(u,v)$ 具有一阶连续偏导数时, 利用具有二阶连续偏导数的函数对它们进行逼近, 则可以证明此时引理的结论仍成立. 由于这种逼近过程过于复杂, 我们在此不作介绍, 但我们还是假定引理在此稍弱的条件下仍成立。

\end{remark}
\begin{proof}\quad 由于函数 $f(x,y)$ 在 $\Omega$ 上可积, 从而
\[
  I=\iint_\Omega f(x,y)\mathrm dx\mathrm dy
\]
存在. 由于 $T$ 是 $D$ 到 $\Omega$ 的同胚映射, 因此 $f(x(u,v),y(u,v))$ 在 $D$ 上可积(见本章习题). 再注意到 $\dfrac{\partial(x,y)}{\partial(u,v)}$ 在 $D$ 上连续, 从而
\[
  \iint_D f(x(u,v),y(u,v))
  \left|\frac{\partial(x,y)}{\partial(u,v)}\right|\mathrm du\mathrm dv
\]
存在. 下面只要证明它等于 $I$ 即可。

取 $\delta>0$ 充分小, 并在 $Ouv$ 平面上作平行于两坐标轴的两个直线族, 使得它们的间距为 $\delta$. 这两个直线族对 $D$ 产生了一个分割:
\[
  \Delta_D=\{\Delta D_1,\Delta D_2,\cdots,\Delta D_K\},
\]
即每个 $\Delta D_k(k=1,2,\cdots,K)$ 都是一个小正方形与 $D$ 的交集. 利用变换 $T$, 我们得到 $\Omega$ 的一个分割
\[
  \Delta_\Omega=\{T(\Delta D_1),T(\Delta D_2),\cdots,T(\Delta D_K)\},
\]
由于 $T$ 在 $D$ 上的一致连续性, 当 $\delta\to0$ 时, 有 $\lambda(\Delta_D)\to0$ 和 $\lambda(\Delta_\Omega)\to0$。

现将 $\Delta_D$ 中的小闭集分成两类: 记 $\Delta'_D=\{\Delta D'_1,\Delta D'_2,\cdots,\Delta D'_J\}$ 为 $\Delta(D)$ 中所有完全落在 $D$ 内部的正方形的全体, 而记 $\Delta''_D=\{\Delta D''_1,\Delta D''_2,\cdots,\Delta D''_L\}$ 为 $\Delta_D$ 中所有含有 $\partial D$ 的点的集合. 显然, $\Delta'_\Omega=\{T(\Delta D'_1),T(\Delta D'_2),\cdots,T(\Delta D'_J)\}$ 中的每个集合均落在 $\Omega$ 的内部, 而 $\Delta''_\Omega=\{T(\Delta D''_1),T(\Delta D''_2),\cdots,T(\Delta D''_L)\}$ 中的每个集合均与 $\partial\Omega$ 的交非空。

由假设 $D$ 由有限条分段光滑曲线所围成, 从而 $D$ 可求面积. 由此我们有
\[
  \lim_{\delta\to0}\sigma\left(\bigcup_{l=1}^L\Delta D''_l\right)=0.
\]
由于变换 $T$ 是 $C^1$ 同胚映射, $\partial\Omega$ 也是由有限条分段光滑曲线所组成. 因此对于 $\forall\varepsilon>0$, 存在有限个长方体组成的简单集 $A$, 使得 $\partial\Omega\subset A^\circ$ 且 $\sigma(A)<\varepsilon$. 因为 $\partial\Omega$ 与 $\partial A$ 均为紧集, 从而有
\[
  \rho=\min_{x\in\partial\Omega,\ y\in\partial A}|x-y|>0.
\]
再由 $T$ 在 $\overline D$ 的一致连续性, 当 $\delta<\dfrac\rho2$ 时, 必有 $T(\Delta D''_l)\subset A^\circ(l=1,2,\cdots,L)$. 这说明了
\[
  \lim_{\delta\to0}\sigma\left(\bigcup_{l=1}^LT(\Delta D''_l)\right)=0.
\]
注意到在 $D$ 上处处有 $\left|\dfrac{\partial(x,y)}{\partial(u,v)}\right|\ne0$, 对每个 $\Delta D'_j(j=1,2,\cdots,J)$, 由有限覆盖定理不难看出, 若将其等分成 $k^2(k\in\mathbb N)$ 个小正方形, 则当 $k$ 充分大时, 在每个小正方形上 $x'_u,x'_v,y'_u$ 及 $y'_v$ 中至少有一个不取零. 因此, $T$ 将该小正方形映成一个曲边小四边形, 这个小曲边四边形必是两个由分段光滑曲线围成的 X 型或 Y 型区域的并. 由引理~\ref{course-lemma-15-4-2}  有
\[
  \sigma(T(\Delta D'_j))
  =\iint_{\Delta D'_j}\left|\frac{\partial(x,y)}{\partial(u,v)}\right|\mathrm du\mathrm dv.
\]
由重积分第一中值定理, 在 $\Delta D'_j(j=1,2,\cdots,J)$ 上存在 $(u'_j,v'_j)$, 使得
\[
  \sigma(T(\Delta D'_j))
  =\left|\frac{\partial(x,y)}{\partial(u,v)}\right|_{(u'_j,v'_j)}
  \sigma(\Delta D'_j).
\]

记 $x'_j=x(u'_j,v'_j), y'_j=y(u'_j,v'_j)(j=1,2,\cdots,J)$. 在 $\Delta D''_l$ 内任取 $(u''_l,v''_l)$, 在 $T(\Delta D''_l)$ 内任取 $(x''_l,y''_l)(l=1,2,\cdots,L)$, 作和式
\begin{equation}
  \sum_{j=1}^J f(x'_j,y'_j)\sigma(T(\Delta D'_j))
  +\sum_{l=1}^L f(x''_l,y''_l)\sigma(T(\Delta D''_l)),
  \label{course-eq-15-4-1}
\end{equation}
则当 $\delta\to0$ 时, 上述和式趋于
\[
  \iint_\Omega f(x,y)\mathrm dx\mathrm dy.
\]
另外, 我们作和式
\begin{equation}
\begin{aligned}
  &\sum_{j=1}^J f(x(u'_j,v'_j),y(u'_j,v'_j))
  \left|\frac{\partial(x,y)}{\partial(u,v)}\right|_{(u'_j,v'_j)}
  \sigma(\Delta D'_j)\\
  &\quad+\sum_{l=1}^L f(x(u''_l,v''_l),y(u''_l,v''_l))
  \left|\frac{\partial(x,y)}{\partial(u,v)}\right|_{(u''_l,v''_l)}
  \sigma(\Delta D''_l),
\end{aligned}
  \label{course-eq-15-4-2}
\end{equation}
则当 $\delta\to0$, 上述和式趋于
\[
  \iint_D f(x(u,v),y(u,v))
  \left|\frac{\partial(x,y)}{\partial(u,v)}\right|\mathrm du\mathrm dv.
\]
由于 $f(x(u,v),y(u,v))$ 和
\[
  \left|f(x(u,v),y(u,v))\left|\frac{\partial(x,y)}{\partial(u,v)}\right|\right|
\]
在 $D$ 上可积, 从而存在 $M>0$, 使得当 $(x,y)\in\Omega$ 时, 有 $|f(x,y)|\leqslant M$, 以及当 $(u,v)\in D$ 时, 有
\[
  \left|f(x(u,v),y(u,v))\left|\frac{\partial(x,y)}{\partial(u,v)}\right|\right|\leqslant M.
\]
由此我们有
\[
\begin{aligned}
  &\left|
  \sum_{j=1}^J f(x'_j,y'_j)\sigma(T(\Delta D'_j))
  +\sum_{l=1}^L f(x''_l,y''_l)\sigma(T(\Delta D''_l))\right.\\
  &\quad-\left[
  \sum_{j=1}^J f(x(u'_j,v'_j),y(u'_j,v'_j))
  \left|\frac{\partial(x,y)}{\partial(u,v)}\right|_{(u'_j,v'_j)}
  \sigma(\Delta D'_j)\right.\\
  &\qquad\left.\left.
  +\sum_{l=1}^L f(x(u''_l,v''_l),y(u''_l,v''_l))
  \left|\frac{\partial(x,y)}{\partial(u,v)}\right|_{(u''_l,v''_l)}
  \sigma(\Delta D''_l)\right]\right|\\
  &=\left|
  \sum_{l=1}^L f(x''_l,y''_l)\sigma(T(\Delta D''_l))
  -\sum_{l=1}^L f(x(u''_l,v''_l),y(u''_l,v''_l))
  \left|\frac{\partial(x,y)}{\partial(u,v)}\right|_{(u''_l,v''_l)}
  \sigma(\Delta D''_l)\right|\\
  &\leqslant M\sigma\left(\bigcup_{l=1}^L T(\Delta D''_l)\right)
  +M\sigma\left(\bigcup_{l=1}^L\Delta D''_l\right)\to0\quad(\delta\to0).
\end{aligned}
\]
因此有
\[
  \iint_\Omega f(x,y)\mathrm dx\mathrm dy
  =\iint_D f(x(u,v),y(u,v))
  \left|\frac{\partial(x,y)}{\partial(u,v)}\right|\mathrm du\mathrm dv.
\]
证毕。

对于 $n$ 重积分, 我们不加证明地给出下述定理。

\end{proof}
\begin{theorem}\label{course-theorem-15-4-3}\quad 设变换
\[
  \boldsymbol x(\boldsymbol u)=(x_1(\boldsymbol u),x_2(\boldsymbol u),\cdots,x_n(\boldsymbol u))
\]
是可求体积的有界闭区域 $D\subset\mathbb R^n$ 到可求体积的有界闭区域 $\Omega\subset\mathbb R^n$ 的同胚映射, 它的各个偏导数在包含 $D$ 的区域上连续, 并且在 $D$ 上
\[
  \frac{\partial(x_1,\cdots,x_n)}{\partial(u_1,\cdots,u_n)}\ne0,
\]
再设函数 $f(\boldsymbol x)$ 在 $\Omega$ 上可积, 则有
\[
\begin{aligned}
  &\int\!\!\int\cdots\int_\Omega f(x_1,\cdots,x_n)\mathrm dx_1\cdots\mathrm dx_n\\
  &\quad=\int\!\!\int\cdots\int_D
  f(x_1(u_1,\cdots,u_n),\cdots,x_n(u_1,\cdots,u_n))
  \left|\frac{\partial(x_1,\cdots,x_n)}{\partial(u_1,\cdots,u_n)}\right|
  \mathrm du_1\cdots\mathrm du_n.
\end{aligned}
\]

\end{theorem}
\subsection{利用变量替换计算重积分}

在本小节中, 我们举例说明如何用变量替换公式来计算重积分. 在二重积分变量替换中, 极坐标变换
\[
  \begin{cases}
    x=r\cos\theta,\\
    y=r\sin\theta
  \end{cases}
  \quad(0\leqslant r<+\infty,0\leqslant\theta\leqslant2\pi)
\]
是非常重要的变换之一. 我们前面已经知道, 它的雅可比行列式为
\[
  \left.\frac{\partial(x,y)}{\partial(r,\theta)}\right|_{(r,\theta)}=r.
\]
注意到极坐标变换在正实轴与原点处不是一一映射的, 但它是 $\{(r,\theta):0<r<+\infty,0<\theta<2\pi\}$ 到 $\mathbb R^2\setminus\{(x,y):x\geqslant0\}$ 的 $C^1$ 同胚映射。

\begin{example}\label{course-example-15-4-1}\quad 计算闭单位圆盘 $\Omega=\{(x,y):x^2+y^2\leqslant1\}$ 的面积。

\end{example}
\begin{solve}\quad 极坐标变换
\[
  \begin{cases}
    x=r\cos\theta,\\
    y=r\sin\theta
  \end{cases}
\]
把 $D=\{(r,\theta):0\leqslant r\leqslant1,0\leqslant\theta\leqslant2\pi\}$ 映成 $\Omega$, 且 $\dfrac{\partial(x,y)}{\partial(r,\theta)}=r$. 因此 $\Omega$ 的面积为
\[
  \sigma(\Omega)=\iint_\Omega\mathrm dx\mathrm dy
  =\iint_D r\mathrm dr\mathrm d\theta
  =\int_0^1 r\mathrm dr\int_0^{2\pi}\mathrm d\theta=\pi.
\]

细心的读者可以发现, 尽管极坐标变换不是 $D$ 到 $\Omega$ 的同胚. 但在上例中应用极坐标变换仍然能正确计算出 $\sigma(\Omega)$. 我们可以用下面的方法来说明上述计算的合理性. 对于 $\forall\varepsilon>0$, 记 $\Omega_\varepsilon=\Omega\setminus\{(x,y):\varepsilon<x<1,-\varepsilon x<y<\varepsilon x\}$, 则极坐标变换是 $D_\varepsilon=\{(r,\theta):\varepsilon\leqslant r\leqslant1,2\pi-\arctan\varepsilon\leqslant\theta\leqslant\arctan\varepsilon\}$ 到 $\Omega_\varepsilon$ 的 $C^1$ 变换. 因此
\[
\begin{aligned}
  \sigma(\Omega)
  &=\lim_{\varepsilon\to0+0}\iint_{\Omega_\varepsilon}\mathrm dx\mathrm dy
  =\lim_{\varepsilon\to0+0}\iint_{D_\varepsilon}r\mathrm dr\mathrm d\theta\\
  &=\lim_{\varepsilon\to0+0}\int_\varepsilon^1 r\mathrm dr
  \int_{\arctan\varepsilon}^{2\pi-\arctan\varepsilon}\mathrm d\theta
  =\int_0^1 r\mathrm dr\int_0^{2\pi}\mathrm d\theta=\pi.
\end{aligned}
\]

因此, 若今后我们遇到一个映射在挖掉有限条光滑曲线后的区域之间是 $C^1$ 同胚映射时, 我们可以在原来的闭区域上直接应用定理~\ref{course-theorem-15-4-1}。

\end{solve}
\begin{example}\label{course-example-15-4-2}\quad 计算二重积分
\[
  \iint_D(x^2+y^2)\mathrm dx\mathrm dy,
\]
其中 $D$ 是由曲线 $(x^2+y^2)^2=a^2(x^2-y^2)$ 所围且落在第一、四象限的部分, 这里 $a>0$ 为常数。

\end{example}
\begin{solve}\quad 利用极坐标变换
\[
  \begin{cases}
    x=r\cos\theta,\\
    y=r\sin\theta
  \end{cases}
\]
可把曲线方程 $(x^2+y^2)^2=a^2(x^2-y^2)$ 化为
\[
  r^4=a^2r^2(\cos^2\theta-\sin^2\theta)=a^2r^2\cos2\theta.
\]
注意到 $x>0$ 时, 有
\[
  r=a\sqrt{\cos2\theta},\quad -\frac\pi4\leqslant\theta\leqslant\frac\pi4,
\]
再注意到 $\dfrac{\partial(x,y)}{\partial(r,\theta)}=r$, 我们有
\[
\begin{aligned}
  \iint_D(x^2+y^2)\mathrm dx\mathrm dy
  &=\int_{-\frac\pi4}^{\frac\pi4}\mathrm d\theta
  \int_0^{a\sqrt{\cos2\theta}}r^2\cdot r\mathrm dr\\
  &=2\int_0^{\frac\pi4}\frac14a^4\cos^22\theta\mathrm d\theta
  =\frac\pi8a^4.
\end{aligned}
\]

下面我们介绍利用一些特殊的变量替换来计算重积分的例子。

\end{solve}
\begin{example}\label{course-example-15-4-3}\quad 计算二重积分
\[
  \iint_D(x^2+y^2)\mathrm dx\mathrm dy,
\]
其中 $D$ 是由曲线 $x^2-y^2=1,x^2-y^2=9,xy=2,xy=4$ 所围成的闭区域。

\end{example}
\begin{solve}\quad 作变换 $T$:
\[
  \begin{cases}
    u=x^2-y^2,\\
    v=2xy,
  \end{cases}
\]
则 $T^{-1}$ 将矩形
\[
  \Omega=\{(u,v):1\leqslant u\leqslant9,4\leqslant v\leqslant8\}
  \]
一一地映成 $D$. 由
\[
  \frac{\partial(x,y)}{\partial(u,v)}
  =\frac1{\dfrac{\partial(u,v)}{\partial(x,y)}}
  =\frac1{\left|\begin{matrix}2x&-2y\\2y&2x\end{matrix}\right|}
  =\frac1{4(x^2+y^2)}
  =\frac1{4\sqrt{u^2+v^2}},
\]
得
\[
\begin{aligned}
  \iint_D(x^2+y^2)\mathrm dx\mathrm dy
  &=\iint_\Omega\sqrt{u^2+v^2}\frac1{4\sqrt{u^2+v^2}}\mathrm du\mathrm dv\\
  &=\frac14\iint_\Omega\mathrm du\mathrm dv=8.
\end{aligned}
\]

\end{solve}
\begin{example}\label{course-example-15-4-4}\quad 求椭球体
\[
  \frac{x^2}{a^2}+\frac{y^2}{b^2}+\frac{z^2}{c^2}\leqslant1\quad(a,b,c>0)
\]
的体积。

\end{example}
\begin{solve}\quad 取
\[
  D=\left\{(x,y):\frac{x^2}{a^2}+\frac{y^2}{b^2}\leqslant1\right\},
\]
则由二重积分的几何意义知, 所求体积为
\[
  V=2\iint_D c\sqrt{1-\frac{x^2}{a^2}-\frac{y^2}{b^2}}\mathrm dx\mathrm dy.
\]
作广义极坐标变换
\[
  \begin{cases}
    x=ar\cos\theta,\\
    y=br\sin\theta,
  \end{cases}
\]
则 $\dfrac{\partial(x,y)}{\partial(r,\theta)}=abr$, 且
\[
  \Omega=\{(r,\theta):0\leqslant r<1,0\leqslant\theta<2\pi\}
\]
经变换一一地映成 $D\setminus\{(0,0)\}$, 从而
\[
\begin{aligned}
  V&=2\iint_\Omega c\sqrt{1-r^2}\,abr\mathrm dr\mathrm d\theta
  =2\int_0^{2\pi}\mathrm d\theta\int_0^1 c\sqrt{1-r^2}\,abr\mathrm dr\\
  &=4\pi abc\cdot\frac12\cdot\frac23(1-r^2)^{\frac32}\Big|_0^1
  =\frac43\pi abc.
\end{aligned}
\]

\end{solve}
\begin{example}\label{course-example-15-4-5}\quad 计算二重积分
\[
  \iint_D\sqrt{\frac{xy}{x+y}}\mathrm dx\mathrm dy,
\]
其中 $D$ 是由直线 $x=0,y=0,x+y=1$ 所围成的闭区域。

\end{example}
\begin{solve}\quad 首先我们注意到
\[
  \lim_{D\ni(x,y)\to(0,0)}\frac{xy}{x+y}=0,
\]
因此该二重积分存在. 显然, 对于积分区域 $D$, 该二重积分易于化成累次积分. 但由于被积函数比较复杂, 该积分仍不易求出. 为此我们利用下述变换
\[
  \begin{cases}
    x=r\cos^2\theta,\\
    y=r\sin^2\theta.
  \end{cases}
\]
该变换将
\[
  \Omega=\left\{(r,\theta):0<r<1,0<\theta<\frac\pi2\right\}
\]
一一地变到 $D$ 的内部, 并且
\[
  \frac{\partial(x,y)}{\partial(r,\theta)}
  =\left|\begin{matrix}
    \cos^2\theta&-r\sin2\theta\\
    \sin^2\theta&r\sin2\theta
  \end{matrix}\right|
  =r\sin2\theta.
\]
因此有
\[
\begin{aligned}
  \iint_D\sqrt{\frac{xy}{x+y}}\mathrm dx\mathrm dy
  &=\frac12\iint_\Omega r^{\frac12}\sin2\theta\cdot r\sin2\theta\,\mathrm dr\mathrm d\theta\\
  &=\frac12\int_0^{\frac\pi2}\sin^22\theta\,\mathrm d\theta\int_0^1 r^{\frac32}\mathrm dr
  =\frac\pi{20}.
\end{aligned}
\]

对于三重积分的计算, 下面的两种变换是经常用到的。

\end{solve}
\noindent\keyterm{柱坐标变换}
\[
  \begin{cases}
    x=r\cos\theta,\\
    y=r\sin\theta,\\
    z=z,
  \end{cases}
\]
其中 $0\leqslant r<+\infty,0\leqslant\theta<2\pi,-\infty<z<+\infty$. 这种变换本质上是将空间一点投影到 $Oxy$ 平面内, 然后取它投影点的极坐标. 因此它将
\[
  \{(r,\theta,z):0\leqslant r<+\infty,0\leqslant\theta<2\pi,-\infty<z<+\infty\}
\]
映成 $\mathbb R^3$, 其雅可比行列式为
\[
  \frac{\partial(x,y,z)}{\partial(r,\theta,z)}=r.
\]

\noindent\keyterm{球坐标变换}
\[
  \begin{cases}
    x=r\sin\varphi\cos\theta,\\
    y=r\sin\varphi\sin\theta,\\
    z=r\cos\varphi,
  \end{cases}
\]
其中 $0\leqslant r<+\infty,0\leqslant\varphi\leqslant\pi,0\leqslant\theta<2\pi$. 它的雅可比行列式
\[
  \frac{\partial(x,y,z)}{\partial(r,\varphi,\theta)}=r^2\sin\varphi.
\]
显然, 在球坐标下, $r=a$ 对应于球面 $x^2+y^2+z^2=a^2$, $\varphi=b$ 对应于锥面 $z=\cot b\sqrt{x^2+y^2}$, $\theta=c$ 为以 $z$ 轴为边的半平面。

\begin{example}\label{course-example-15-4-6}\quad 设空间立体 $D$ 由曲面 $z=25-x^2-y^2$ 与平面 $z=0$ 所围成, 试求它的体积 $V$。

\end{example}
\begin{solve}\quad 由三重积分的几何意义得
\[
  V=\iiint_D\mathrm dx\mathrm dy\mathrm dz.
\]
取柱坐标变换
\[
  \begin{cases}
    x=r\cos\theta,\\
    y=r\sin\theta,\\
    z=z,
  \end{cases}
\]
则它将
\[
  \Omega=\{(r,\theta,z):0\leqslant r\leqslant5,0\leqslant\theta<2\pi,0\leqslant z\leqslant25-r^2\}
\]
映成 $D$. 因此
\[
\begin{aligned}
  V&=\iiint_\Omega\frac{\partial(x,y,z)}{\partial(r,\theta,z)}
  \mathrm dr\mathrm d\theta\mathrm dz
  =\int_0^{2\pi}\mathrm d\theta\int_0^5 r\mathrm dr\int_0^{25-r^2}\mathrm dz\\
  &=2\pi\int_0^5 r(25-r^2)\mathrm dr=\frac{625}2\pi.
\end{aligned}
\]

\end{solve}
\begin{remark}\quad 上例中的柱坐标变换在 $\Omega$ 和 $D$ 内分别挖掉一个光滑曲面后是 $C^1$ 同胚, 利用处理极坐标变换相似的方法, 容易验证上述的重积分的变量替换公式仍然成立。

\end{remark}
\begin{example}\label{course-example-15-4-7}\quad 设 $\mathbb R^3$ 中均匀物体 $D$ 由球面 $x^2+y^2+z^2=1$ 和锥面 $z=\sqrt{x^2+y^2}$ 所围成, 求它的质心坐标。

\end{example}
\begin{solve}\quad 由对称性, 可设其质心坐标为 $(0,0,z_0)$. 锥面方程 $z=\sqrt{x^2+y^2}$ 在球坐标系下为 $\varphi=\dfrac\pi4$. 因此 $D$ 的体积
\[
\begin{aligned}
  V&=\iiint_D\mathrm dV
  =\int_0^{2\pi}\mathrm d\theta\int_0^{\frac\pi4}\mathrm d\varphi\int_0^1 r^2\sin\varphi\mathrm dr\\
  &=2\pi\int_0^{\frac\pi4}\sin\varphi\frac{r^3}3\Big|_0^1\mathrm d\varphi
  =\frac{2\pi}3\left(1-\frac{\sqrt2}2\right).
\end{aligned}
\]
下面还要计算物体到 $Oxy$ 平面的力矩 $M_{xy}$. 利用微元法, 我们有
\[
\begin{aligned}
  M_{xy}&=\iiint_D z\mathrm dV
  =\int_0^{2\pi}\mathrm d\theta\int_0^{\frac\pi4}\mathrm d\varphi\int_0^1
  r\cos\varphi\cdot r^2\sin\varphi\mathrm dr\\
  &=2\pi\int_0^{\frac\pi4}\sin\varphi\cos\varphi\mathrm d\varphi\int_0^1r^3\mathrm dr
  =\frac\pi8.
\end{aligned}
\]
因此有
\[
  z_0=\frac{M_{xy}}V=\frac{3\left(1+\dfrac{\sqrt2}2\right)}8.
\]

\end{solve}


\section{第一型曲线积分}\label{course-16-1}
本节用于曲线、曲面积分题中的参数化、方向和积分公式。应用 Green、Gauss 或 Stokes 公式之前，应逐项检查光滑性、闭合性、定向以及被积场的奇点；缺失边界必须补上并单独计算。

\subsection{第一型曲线积分的定义}

在定积分的应用中, 我们学习过曲线弧长的定义. 现在我们简要回顾一下. 设 $\Gamma$ 是平面内或空间 $\mathbb R^3$ 中的一条连续的简单曲线, 以 $A,B$ 为其端点 (当 $A=B$ 时, 则 $\Gamma$ 是一条约当曲线). 在 $\Gamma$ 上从 $A$ 到 $B$ 依次取点 $A=A_0,A_1,\cdots,A_K=B$, 它们构成了 $\Gamma$ 的一个分割 $T$. 记 $\overline{A_{k-1}A_k}$ 为连接 $A_{k-1},A_k$ 的线段, 其长为 $s_k(k=1,2,\cdots,K)$, 则当
\[
  \sup_T\sum_{k=1}^K s_k<+\infty
\]
时, 其中上确界是对 $\Gamma$ 所有的分割 $T$ 取的, 我们称该曲线是可求长的, 并称该上确界为 $\Gamma$ 的弧长。

现设光滑曲线 $\Gamma$ 由参数方程
\[
  \begin{cases}
    x=x(t),\\
    y=y(t),\\
    z=z(t)
  \end{cases}
  \quad(\alpha\leqslant t\leqslant \beta)
\]
给出, 则其弧长 $L$ 可由下述公式计算:
\[
  L=\int_\alpha^\beta \sqrt{(x'(t))^2+(y'(t))^2+(z'(t))^2}\mathrm dt.
\]

为了引入第一型曲线积分, 我们先来研究以下物质曲线的质量问题. 设 $\Gamma$ 是一条物质曲线, 它可被看成落在 $\mathbb R^3$ 中的一条曲线, 并在它上面定义了一个密度函数. 现在设 $\Gamma$ 以 $A$ 为起点, 以 $B$ 为终点, 在 $\Gamma$ 上任一点 $(x,y,z)$ 处该物质曲线的线密度为 $\rho(x,y,z)$, 我们来求 $\Gamma$ 的质量. 当线密度函数为常数函数, 即 $\rho(x,y,z)\equiv\rho_0$ 时, 我们可以求出该物质曲线的质量为 $\rho_0L$, 其中 $L$ 为该曲线的弧长. 当 $\rho(x,y,z)$ 在 $\Gamma$ 上不是常数函数时, 依照定积分的思想, 我们可以在 $\Gamma$ 上依次取点 $A=A_0,A_1,\cdots,A_K=B$, 并称其为 $\Gamma$ 的一个分割. 此分割将 $\Gamma$ 分成 $K$ 个小弧段. 在以 $A_{k-1},A_k$ 为端点的弧段 $\widehat{A_{k-1}A_k}$ 上任取一点 $(\xi_k,\eta_k,\zeta_k)(k=1,2,\cdots,K)$, 作和式 $\sum\limits_{k=1}^K \rho(\xi_k,\eta_k,\zeta_k)\Delta s_k$, 其中 $\Delta s_k(k=1,2,\cdots,K)$ 是 $\widehat{A_{k-1}A_k}$ 的弧长. 如同求曲边梯形的面积, 当分割越来越细时, 若该和式的极限趋向于常数 $I$, 则我们认为 $I$ 即为该物质曲线的质量. 由此我们引进下面的定义。

\begin{definition}\label{course-definition-16-1-1}\quad 设 $\Gamma\subset\mathbb R^3$ 是可求长的简单曲线, 它以 $A$ 为起点, $B$ 为终点, $f(x,y,z)$ 是定义在 $\Gamma$ 上的函数. 对于 $\Gamma$ 的任一分割 $T$: $A=A_0,A_1,\cdots,A_K=B$, 记 $\Delta s_k(k=1,2,\cdots,K)$ 为弧段 $\widehat{A_{k-1}A_k}$ 的弧长 (即 $\Gamma$ 介于 $A_{k-1},A_k$ 之间部分的长度) 及 $\lambda(T)=\max\limits_{1\leqslant k\leqslant K}\Delta s_k$. 在 $\widehat{A_{k-1}A_k}$ 上任取一点 $(\xi_k,\eta_k,\zeta_k)(k=1,2,\cdots,K)$, 作和式
\[
  \sum_{k=1}^K f(\xi_k,\eta_k,\zeta_k)\Delta s_k.
\]
若存在 $I\in\mathbb R$, 使得对于 $\forall\varepsilon>0$, $\exists\delta>0$, 对 $\Gamma$ 的任意分割 $T$ 及每个小弧段上任意选取的 $(\xi_k,\eta_k,\zeta_k)$, 当 $\lambda(T)<\delta$ 时, 有
\[
  \left|\sum_{k=1}^K f(\xi_k,\eta_k,\zeta_k)\Delta s_k-I\right|<\varepsilon,
\]
即
\[
  \lim_{\lambda(T)\to0}\sum_{k=1}^K f(\xi_k,\eta_k,\zeta_k)\Delta s_k=I,
\]
则称 $f(x,y,z)$ 沿 $\Gamma$ 的第一型曲线积分存在, 并称 $I$ 为 $f(x,y,z)$ 在 $\Gamma$ 上的第一型曲线积分, 记为
\[
  I=\int_\Gamma f(x,y,z)\mathrm ds
  \quad\text{或}\quad
  I=\int_{\widehat{AB}} f(x,y,z)\mathrm ds,
\]
其中 $f(x,y,z)$ 称为被积函数, $\Gamma$ 称为积分曲线。

从定义~\ref{course-definition-16-1-1}  可以看出, 第一型曲线积分只与积分曲线 $\Gamma$ 的弧长及被积函数 $f(x,y,z)$ 的取值有关, 而与 $\Gamma$ 的方向无关. 设 $\Gamma=\widehat{AB}$, 即 $\Gamma$ 是以 $A$ 为起点, $B$ 为终点的一条曲线. 记 $\Gamma^-=\widehat{BA}$ 是 $\Gamma$ 的反向曲线, 即将 $\Gamma$ 的起点与终点对调后所得到的曲线, 则当 $f(x,y,z)$ 在 $\Gamma$ 上的第一型曲线积分存在时, 有
\[
  \int_\Gamma f(x,y,z)\mathrm ds=\int_{\Gamma^-} f(x,y,z)\mathrm ds.
\]

从定义~\ref{course-definition-16-1-1}  还可以推知第一型曲线积分具有下述性质:

(1) 若对于 $\forall(x,y,z)\in\Gamma$, 有 $f(x,y,z)\equiv1$, 则 $\int_\Gamma f(x,y,z)\mathrm ds=\int_\Gamma \mathrm ds$ 即为 $\Gamma$ 的弧长。

(2) 若 $\Gamma$ 落在 $Oxy$ 平面上, 且对于 $\forall(x,y)\in\Gamma$ 有 $f(x,y)>0$, 则 $\int_\Gamma f(x,y)\mathrm ds$ 表示空间柱面 $S=\{(x,y,z):(x,y)\in\Gamma,0\leqslant z\leqslant f(x,y)\}$ 的面积 （曲面面积的定义见第~\ref{course-16-3} 节）。

(3) 第一型曲线积分具有关于被积函数的线性性, 即设 $\Gamma$ 是可求长的简单曲线, 函数 $f(x,y,z)$ 与 $g(x,y,z)$ 在 $\Gamma$ 上的第一型曲线积分存在, $\alpha,\beta\in\mathbb R$, 则 $\alpha f(x,y,z)+\beta g(x,y,z)$ 在 $\Gamma$ 上的第一型曲线积分存在, 并且
\[
  \int_\Gamma [\alpha f(x,y,z)+\beta g(x,y,z)]\mathrm ds
  =\alpha\int_\Gamma f(x,y,z)\mathrm ds
  +\beta\int_\Gamma g(x,y,z)\mathrm ds.
\]

(4) 第一型曲线积分具有关于积分曲线的可加性, 即设 $\Gamma_1=\widehat{AB}$, $\Gamma_2=\widehat{BC}$ 均为可求长的简单曲线, 且 $\Gamma=\Gamma_1+\Gamma_2=\widehat{AC}$ 也为可求长的简单曲线, 则函数 $f(x,y,z)$ 在 $\Gamma$ 上的第一型曲线积分存在的充分必要条件是 $f(x,y,z)$ 在 $\Gamma_1$ 和 $\Gamma_2$ 上的第一型曲线积分均存在, 且当第一型曲线积分存在时, 有
\[
  \int_{\Gamma_1} f(x,y,z)\mathrm ds+\int_{\Gamma_2} f(x,y,z)\mathrm ds
  =\int_\Gamma f(x,y,z)\mathrm ds.
\]
以上性质的证明请读者自己给出。

\end{definition}
\subsection{第一型曲线积分的存在性与计算公式}

对于第一型曲线积分的存在性, 我们有下面的结论。

\begin{theorem}\label{course-theorem-16-1-1}\quad 设 $\Gamma=\widehat{AB}\subset\mathbb R^3$ 是可求长的简单曲线, 函数 $f(x,y,z)$ 在 $\Gamma$ 上连续, 则第一型曲线积分 $\int_\Gamma f(x,y,z)\mathrm ds$ 存在。

\end{theorem}
\begin{proof}\quad 设 $\Gamma$ 以弧长 $s$ 为参数的参数方程为
\[
  \Gamma:
  \begin{cases}
    x=x(s),\\
    y=y(s),\\
    z=z(s)
  \end{cases}
  \quad(0\leqslant s\leqslant L),
\]
其中 $L$ 为 $\Gamma$ 的弧长. 对于 $\Gamma=\widehat{AB}$ 的一个分割 $T$: $A=A_0,A_1,\cdots,A_K=B$, 则对应于闭区间 $[0,L]$ 的一个分割 $\Delta$: $0=s_0<s_1<\cdots<s_K=L$. 在 $\widehat{A_{k-1}A_k}$ 上任取一点 $(\xi_k,\eta_k,\zeta_k)$, 则必存在 $s'_k\in[s_{k-1},s_k]$, 使得 $(\xi_k,\eta_k,\zeta_k)=(x(s'_k),y(s'_k),z(s'_k))(k=1,2,\cdots,K)$. 因此
\[
  \sum_{k=1}^K f(\xi_k,\eta_k,\zeta_k)\Delta s_k
  =\sum_{k=1}^K f(x(s'_k),y(s'_k),z(s'_k))\Delta s_k.
\]
上述等式右端是连续函数 $f(x(s),y(s),z(s))$ 在区间 $[0,L]$ 上的黎曼和. 注意到 $\lambda(T)\to0$ 的充分必要条件是 $\lambda(\Delta)\to0$, 因此 $\int_\Gamma f(x,y,z)\mathrm ds$ 存在, 并且有
\[
  \int_\Gamma f(x,y,z)\mathrm ds=\int_0^L f(x(s),y(s),z(s))\mathrm ds.
\]
证毕。

在假定 $\Gamma\subset\mathbb R^3$ 是光滑曲线时, 我们有以下关于第一型曲线积分的计算公式。

\end{proof}
\begin{theorem}\label{course-theorem-16-1-2}\quad 设 $\Gamma$ 是光滑曲线, 其参数方程为
\[
  \begin{cases}
    x=x(t),\\
    y=y(t),\\
    z=z(t)
  \end{cases}
  \quad(\alpha\leqslant t\leqslant \beta),
\]
再设函数 $f(x,y,z)$ 在 $\Gamma$ 上连续, 则有
\[
  \int_\Gamma f(x,y,z)\mathrm ds
  =\int_\alpha^\beta f(x(t),y(t),z(t))
  \sqrt{[x'(t)]^2+[y'(t)]^2+[z'(t)]^2}\mathrm dt.
\]

\end{theorem}
\begin{proof}\quad 由于 $\Gamma$ 是光滑曲线, 设其弧长为 $L$. 作 $[\alpha,\beta]\to[0,L]$ 的变换如下: 对于 $\forall t\in[\alpha,\beta]$, 记 $s(t)$ 为 $\Gamma$ 介于 $(x(\alpha),y(\alpha),z(\alpha))$ 与 $(x(t),y(t),z(t))$ 之间部分的弧长, 即
\[
  s=s(t)=\int_\alpha^t \sqrt{[x'(u)]^2+[y'(u)]^2+[z'(u)]^2}\mathrm du.
\]
注意到
\[
  \frac{\mathrm ds}{\mathrm dt}
  =\sqrt{[x'(t)]^2+[y'(t)]^2+[z'(t)]^2}>0,
\]
记 $t=t(s)$ 为 $s(t)$ 的反函数, 因此由定积分换元公式得
\[
\begin{aligned}
  \int_\Gamma f(x,y,z)\mathrm ds
  &=\int_0^L f(x(t(s)),y(t(s)),z(t(s)))\mathrm ds\\
  &=\int_\alpha^\beta f(x(t),y(t),z(t))
    \sqrt{[x'(t)]^2+[y'(t)]^2+[z'(t)]^2}\mathrm dt.
\end{aligned}
\]
证毕。

\end{proof}
\begin{example}\label{course-example-16-1-1}\quad 计算第一型曲线积分 $\int_\Gamma \sqrt{1-x^2-y^2}\mathrm ds$, 其中 $\Gamma$ 是曲线 $x^2+y^2=x$。

\end{example}
\noindent\keyterm{解法 1}\quad 取曲线 $\Gamma$ 的参数方程
\[
  \begin{cases}
    x=\dfrac12+\dfrac12\cos t,\\
    y=\dfrac12\sin t
  \end{cases}
  \quad(0\leqslant t\leqslant 2\pi),
\]
则
\[
  \mathrm ds=\sqrt{\frac14\sin^2t+\frac14\cos^2t}\mathrm dt
  =\frac12\mathrm dt.
\]
于是
\[
\begin{aligned}
  \int_\Gamma \sqrt{1-x^2-y^2}\mathrm ds
  &=\int_\Gamma \sqrt{1-x}\mathrm ds
    =\frac12\int_0^{2\pi}\sqrt{\frac{1-\cos t}{2}}\mathrm dt\\
  &=\frac12\int_0^{2\pi}\left|\sin\frac t2\right|\mathrm dt
    =\int_0^\pi\sin\frac t2\mathrm dt=2.
\end{aligned}
\]

\noindent\keyterm{解法 2}\quad 由于被积函数及积分曲线的对称性, 我们只要考虑积分曲线落在第一象限的部分即可. 此时曲线 $\Gamma$ 可写成如下的参数方程
\[
  \begin{cases}
    x=x,\\
    y=\sqrt{x-x^2}
  \end{cases}
  \quad(0\leqslant x\leqslant1),
\]
从而 $\mathrm ds=\dfrac{\mathrm dx}{2\sqrt{x-x^2}}$, 因此
\[
  \int_\Gamma \sqrt{1-x^2-y^2}\mathrm ds
  =2\int_0^1\sqrt{1-x}\frac{\mathrm dx}{2\sqrt{x-x^2}}
  =\int_0^1\frac{\mathrm dx}{\sqrt x}=2.
\]

\begin{remark}\quad 在上例中, 解法 2 中的计算更简单些, 但最后的积分是一个瑕积分. 另外, 由第一型曲线积分的性质, 上例中的积分值是柱面 $x^2+y^2=x$ 被单位球面所截得的在 $Oxy$ 平面上方那一块的面积。

\end{remark}
\begin{example}\label{course-example-16-1-2}\quad 计算第一型曲线积分 $\int_\Gamma (x+y)^2\mathrm ds$, 其中 $\Gamma$ 是单位圆盘上半部分的边界, 即 $\Gamma=\Gamma_1\cup\Gamma_2$, 其中
\[
  \Gamma_1=\{(x,y):x^2+y^2=1,y\geqslant0\},\quad
  \Gamma_2=\{(x,y):-1\leqslant x\leqslant1,y=0\}.
\]

\end{example}
\begin{solve}\quad 由第一型曲线积分的性质, 我们有
\[
  \int_\Gamma (x+y)^2\mathrm ds
  =\int_{\Gamma_1}(x+y)^2\mathrm ds+\int_{\Gamma_2}(x+y)^2\mathrm ds.
\]
由于
\[
  \int_{\Gamma_1}(x+y)^2\mathrm ds
  =\int_{\Gamma_1}(x^2+y^2)\mathrm ds+2\int_{\Gamma_1}xy\mathrm ds
  =\int_{\Gamma_1}\mathrm ds+0=\pi,
\]
其中 $\int_{\Gamma_1}xy\mathrm ds=0$ 由对称性得到, 另外我们有
\[
  \int_{\Gamma_2}(x+y)^2\mathrm ds
  =\int_{-1}^1 x^2\mathrm dx=\frac23,
\]
所以
\[
  \int_\Gamma (x+y)^2\mathrm ds=\frac23+\pi.
\]

\end{solve}


\section{第二型曲线积分}\label{course-16-2}
本节用于曲线、曲面积分题中的参数化、方向和积分公式。应用 Green、Gauss 或 Stokes 公式之前，应逐项检查光滑性、闭合性、定向以及被积场的奇点；缺失边界必须补上并单独计算。

\subsection{第二型曲线积分的定义}

第二型曲线积分的物理模型是变力做功. 设空间有一个力场, 如地球对空间中任一点都存在引力, 该引力在空间的分布就可称为地球的引力场. 从数学的观点来看一个力场即是在 $\mathbb R^3$ 中定义了一个向量函数
\[
  \boldsymbol F(x,y,z)=(P(x,y,z),Q(x,y,z),R(x,y,z)).
\]
假定在它的作用下, 一个质点沿空间曲线 $\Gamma$ 从点 $A$ 移到了点 $B$. 现在我们来求该力场对质点所做的功。

首先注意到, 若 $\boldsymbol F$ 是常力, 即它的大小与方向均是常数, 并且在它的作用下, 质点沿连接点 $A$ 和 $B$ 的直线段 $\overline{AB}$ 从 $A$ 移到 $B$, 则它所做的功为
\[
  W=\boldsymbol F\cdot\overrightarrow{AB}.
\]
为了求解上述的一般问题, 我们用熟知的分割、求和、取极限的积分技巧. 记 $\Gamma=\widehat{AB}$, 沿 $\Gamma$ 依次取分点
\[
  A=A_0,A_1,\cdots,A_K=B.
\]
它们构成 $\Gamma$ 的一个分割 $T$, 记 $\lambda(T)=\max\limits_{1\leqslant k\leqslant K}\{\operatorname{diam}(\widehat{A_{k-1}A_k})\}$. 在每一段小弧 $\widehat{A_{k-1}A_k}(k=1,2,\cdots,K)$ 上, 以线段 $\overline{A_{k-1}A_k}$ 代替 $\widehat{A_{k-1}A_k}$, 并且将 $\boldsymbol F(x,y,z)$ 在 $\widehat{A_{k-1}A_k}$ 上取常力 $\boldsymbol F(\xi_k,\eta_k,\zeta_k)((\xi_k,\eta_k,\zeta_k)\in\widehat{A_{k-1}A_k})$, 则在变力 $\boldsymbol F$ 作用下质点沿 $\Gamma$ 从 $A_{k-1}$ 移动到 $A_k$ 所做的功
\[
  \Delta W_k\approx \boldsymbol F(\xi_k,\eta_k,\zeta_k)\cdot\overrightarrow{A_{k-1}A_k}.
\]
因此
\[
  W=\sum_{k=1}^K\Delta W_k
  \approx \sum_{k=1}^K \boldsymbol F(\xi_k,\eta_k,\zeta_k)\cdot\overrightarrow{A_{k-1}A_k}.
\]
当 $\lambda(T)\to0$ 时, 若上式右端和式的极限存在, 则
\[
  W=\lim_{\lambda(T)\to0}\sum_{k=1}^K
  \boldsymbol F(\xi_k,\eta_k,\zeta_k)\cdot\overrightarrow{A_{k-1}A_k}.
\]
记 $A_k=(x_k,y_k,z_k)(k=1,2,\cdots,K)$, 则
\[
  \overrightarrow{A_{k-1}A_k}
  =(x_k-x_{k-1},y_k-y_{k-1},z_k-z_{k-1})
  \triangleq(\Delta x_k,\Delta y_k,\Delta z_k).
\]
于是我们有
\[
\begin{aligned}
  &\sum_{k=1}^K \boldsymbol F(\xi_k,\eta_k,\zeta_k)\cdot\overrightarrow{A_{k-1}A_k}\\
  &\quad=\sum_{k=1}^K\bigl(P(\xi_k,\eta_k,\zeta_k)\Delta x_k
    +Q(\xi_k,\eta_k,\zeta_k)\Delta y_k
    +R(\xi_k,\eta_k,\zeta_k)\Delta z_k\bigr).
\end{aligned}
\]
从上述变力做功的问题中, 我们得到一个新的和式极限. 由此和式极限我们可以引进一类新的积分. 另外要注意, 我们记空间一条曲线 $\Gamma=\widehat{AB}$ 是为了强调 $\Gamma$ 是以 $A$ 为起点, $B$ 为终点的. 当 $A=B$ 时, 我们则必须在 $\Gamma$ 上取定一个方向。

\begin{definition}\label{course-definition-16-2-1}\quad 设 $\Gamma=\widehat{AB}$ 为空间连续曲线, 向量函数
\[
  \boldsymbol F(x,y,z)=(P(x,y,z),Q(x,y,z),R(x,y,z))
\]
在 $\widehat{AB}$ 上有定义, 记在 $\widehat{AB}$ 上依次取分点
\[
  A=A_0,A_1,\cdots,A_K=B
\]
组成的分割为 $T$,\allowbreak $\overrightarrow{A_{k-1}A_k}=(\Delta x_k,\Delta y_k,\Delta z_k)(k=1,2,\cdots,K)$, $\lambda(T)=\max\limits_{1\leqslant k\leqslant K}\{\operatorname{diam}(\widehat{A_{k-1}A_k})\}$. 在 $\widehat{A_{k-1}A_k}$ 上任取点 $(\xi_k,\eta_k,\zeta_k)(k=1,2,\cdots,K)$. 若极限
\[
  \lim_{\lambda(T)\to0}\sum_{k=1}^K
  [P(\xi_k,\eta_k,\zeta_k)\Delta x_k
  +Q(\xi_k,\eta_k,\zeta_k)\Delta y_k
  +R(\xi_k,\eta_k,\zeta_k)\Delta z_k]
\]
存在并且等于 $I$, 其中常数 $I$ 不依赖 $\widehat{AB}$ 的分割 $T$ 以及 $(\xi_k,\eta_k,\zeta_k)$ 的选取, 则称该极限值 $I$ 为 $\boldsymbol F(x,y,z)$ 在 $\Gamma$ 上的第二型曲线积分, 记为
\[
  I=\int_\Gamma P(x,y,z)\mathrm dx+Q(x,y,z)\mathrm dy+R(x,y,z)\mathrm dz,
\]
或者
\[
  I=\int_{\widehat{AB}}P\mathrm dx+Q\mathrm dy+R\mathrm dz,
\]
其中 $\boldsymbol F(x,y,z)=(P,Q,R)$ 称为被积函数, $\Gamma$ 称为积分曲线。

从上述定义可以看出, 第二型曲线积分是考虑一个向量函数在曲线上沿着一个方向的积分, 因此它与曲线的方向有关. 显然, 第一型曲线积分是考虑一个多元函数在曲线上的积分. 当然, 仅从纯数学的观点, 第二型曲线积分也可仅仅定义一个多元函数 $P(x,y,z)$ 的下述积分 $\int_{\widehat{AB}}P\mathrm dx$, $\int_{\widehat{AB}}P\mathrm dy$, $\int_{\widehat{AB}}P\mathrm dz$. 但是这些积分不能反映出第二型曲线积分的物理背景。

从定义我们可以看出第二型曲线积分具有下述性质:

(1) 若 $\int_{\widehat{AB}}P\mathrm dx+Q\mathrm dy+R\mathrm dz$ 存在, 则 $\int_{\widehat{BA}}P\mathrm dx+Q\mathrm dy+R\mathrm dz$ 存在, 并且
\[
  \int_{\widehat{BA}}P\mathrm dx+Q\mathrm dy+R\mathrm dz
  =-\int_{\widehat{AB}}P\mathrm dx+Q\mathrm dy+R\mathrm dz.
\]
特别地, 当 $\Gamma$ 是一条简单闭曲线时, 我们有
\[
  \int_{\Gamma^-}P\mathrm dx+Q\mathrm dy+R\mathrm dz
  =-\int_\Gamma P\mathrm dx+Q\mathrm dy+R\mathrm dz.
\]

(2) 设向量函数 $\boldsymbol f=(P_1,Q_1,R_1)$ 和 $\boldsymbol g=(P_2,Q_2,R_2)$ 在空间曲线 $\widehat{AB}$ 上的第二型曲线积分都存在, $\alpha,\beta\in\mathbb R$, 则 $\alpha\boldsymbol f+\beta\boldsymbol g$ 在 $\widehat{AB}$ 上的第二型曲线积分存在, 并且
\[
\begin{aligned}
  &\int_{\widehat{AB}}(\alpha P_1\mathrm dx+\alpha Q_1\mathrm dy+\alpha R_1\mathrm dz)
  +(\beta P_2\mathrm dx+\beta Q_2\mathrm dy+\beta R_2\mathrm dz)\\
  &\quad=\alpha\int_{\widehat{AB}}P_1\mathrm dx+Q_1\mathrm dy+R_1\mathrm dz
  +\beta\int_{\widehat{AB}}P_2\mathrm dx+Q_2\mathrm dy+R_2\mathrm dz.
\end{aligned}
\]

(3) 设 $\widehat{AB}=\widehat{AC}\cup\widehat{CB}$, 即 $\widehat{AB}$ 是由两条连续曲线 $\widehat{AC},\widehat{CB}$ 组成, 并且 $\widehat{AC}$ 与 $\widehat{CB}$ 仅在端点 $C$ 处相交, 则向量函数 $\boldsymbol F=(P,Q,R)$ 在 $\widehat{AB}$ 上的第二型曲线积分存在的充分必要条件是它在 $\widehat{AC}$ 及 $\widehat{CB}$ 上的第二型曲线积分均存在, 并且第二型曲线积分存在时, 成立等式
\[
  \int_{\widehat{AB}}P\mathrm dx+Q\mathrm dy+R\mathrm dz
  =\int_{\widehat{AC}}P\mathrm dx+Q\mathrm dy+R\mathrm dz
  +\int_{\widehat{CB}}P\mathrm dx+Q\mathrm dy+R\mathrm dz.
\]

\end{definition}
\subsection{第二型曲线积分的存在性与计算公式}

在这里, 我们不准备讨论关于非常广泛的积分曲线和被积函数的第二型曲线积分存在性问题, 只讨论积分曲线是光滑曲线, 而被积函数是连续的情形。

\begin{theorem}\label{course-theorem-16-2-1}\quad 设 $\Gamma=\widehat{AB}$ 是以 $A$ 为起点, $B$ 为终点的光滑曲线, 其参数方程为
\[
  \begin{cases}
    x=x(t),\\
    y=y(t),\\
    z=z(t)
  \end{cases}
  \quad(\alpha\leqslant t\leqslant \beta)
\]
且当 $t$ 从 $\alpha$ 连续变为 $\beta$ 时, 对应曲线上的点从 $A$ 连续变为 $B$, 再假定函数\allowbreak $P(x,y,z),Q(x,y,z),R(x,y,z)$ 在 $\Gamma$ 上连续, 则向量函数 $\boldsymbol F=(P,Q,R)$ 在 $\widehat{AB}$ 上的第二型曲线积分存在, 并且
\[
\begin{aligned}
  \int_{\widehat{AB}}P\mathrm dx+Q\mathrm dy+R\mathrm dz
  &=\int_\alpha^\beta P(x(t),y(t),z(t))x'(t)\mathrm dt
    +\int_\alpha^\beta Q(x(t),y(t),z(t))y'(t)\mathrm dt\\
  &\quad+\int_\alpha^\beta R(x(t),y(t),z(t))z'(t)\mathrm dt.
\end{aligned}
\]

\end{theorem}
\begin{proof}\quad 我们用微元法证之. 由光滑曲线的定义知 $x'(t),y'(t),z'(t)$ 均是 $t$ 的连续函数, 且对于 $\forall t\in(\alpha,\beta)$, 有
\[
  x'^2(t)+y'^2(t)+z'^2(t)\ne0.
\]
对于 $\forall[t,t+\mathrm dt]\subset[\alpha,\beta]$, 记 $A_t=(x(t),y(t),z(t))\in\Gamma$, $A_{t+\mathrm dt}=(x(t+\mathrm dt),y(t+\mathrm dt),z(t+\mathrm dt))\in\Gamma$, 由于
\[
\begin{aligned}
  \overrightarrow{A_tA_{t+\mathrm dt}}
  &\approx (x'(t),y'(t),z'(t))\mathrm dt\\
  &=\frac{(x'(t),y'(t),z'(t))}
    {\sqrt{x'^2(t)+y'^2(t)+z'^2(t)}}
    \sqrt{x'^2(t)+y'^2(t)+z'^2(t)}\mathrm dt\\
  &=\frac{(x'(t),y'(t),z'(t))}
    {\sqrt{x'^2(t)+y'^2(t)+z'^2(t)}}\mathrm ds\\
  &=(\cos\alpha,\cos\beta,\cos\gamma)\mathrm ds.
\end{aligned}
\]
其中 $(\cos\alpha,\cos\beta,\cos\gamma)$ 是曲线在点 $(x(t),y(t),z(t))$ 处的单位切向量. 因此有
\[
\begin{aligned}
  P\mathrm dx+Q\mathrm dy+R\mathrm dz
  &=(P,Q,R)(\mathrm dx,\mathrm dy,\mathrm dz)\\
  &=(P,Q,R)(\cos\alpha,\cos\beta,\cos\gamma)\mathrm ds\\
  &=(P\cos\alpha+Q\cos\beta+R\cos\gamma)\mathrm ds.
\end{aligned}
\]
由假设, $\Gamma$ 是光滑曲线, 从而 $\cos\alpha,\cos\beta,\cos\gamma$ 是 $t$ 的连续函数. 由第一型曲线积分的存在性定理知 $\int_{\widehat{AB}}(P\cos\alpha+Q\cos\beta+R\cos\gamma)\mathrm ds$ 存在, 从而 $\int_{\widehat{AB}}P\mathrm dx+Q\mathrm dy+R\mathrm dz$ 存在, 并且有公式
\[
\begin{aligned}
  \int_{\widehat{AB}}P\mathrm dx+Q\mathrm dy+R\mathrm dz
  &=\int_\alpha^\beta P(x(t),y(t),z(t))x'(t)\mathrm dt
    +\int_\alpha^\beta Q(x(t),y(t),z(t))y'(t)\mathrm dt\\
  &\quad+\int_\alpha^\beta R(x(t),y(t),z(t))z'(t)\mathrm dt.
\end{aligned}
\]
证毕。

\end{proof}
\begin{remark}\quad 在上面定理的证明过程中, 我们清楚地知道, 第二型曲线积分与第一型曲线积分有如下的关系:
\[
  \int_{\widehat{AB}}P\mathrm dx+Q\mathrm dy+R\mathrm dz
  =\int_{\widehat{AB}}(P,Q,R)\cdot(\cos\alpha,\cos\beta,\cos\gamma)\mathrm ds,
\]
其中 $(\cos\alpha,\cos\beta,\cos\gamma)$ 为 $\widehat{AB}$ 在 $(x,y,z)$ 处的单位切向量。

到现在为止, 我们讨论了平面及空间的曲线积分问题. 值得指出的是, 在 $\mathbb R^n(n>3)$ 中也可以讨论相应的曲线积分问题. 事实上, 我们已经有了 $\mathbb R^n$ 中曲线的定义, 因此可以用完全平行于 $\mathbb R^3$ 中的理论来讨论可求长曲线、光滑曲线等问题, 从而可以定义两类不同的积分. 由于这些理论与我们已经研究的 $\mathbb R^3$ 中的相应理论有很大的相似性, 我们在这里就不深入讨论了。

\end{remark}
\begin{example}\label{course-example-16-2-1}\quad 计算第二型曲线积分 $\int_{\widehat{AB}}(-y)\mathrm dx+x\mathrm dy$, 其中 $\widehat{AB}$ 为单位圆周 $x^2+y^2=1$ 的上半部分, 方向为从点 $A(1,0)$ 到点 $B(-1,0)$。

\end{example}
\noindent\keyterm{解法 1}\quad 令
\[
  \begin{cases}
    x=\cos t,\\
    y=\sin t
  \end{cases}
  \quad(0\leqslant t\leqslant\pi),
\]
则 $t$ 从 $0$ 连续变化到 $\pi$ 时, 曲线上的点从 $A$ 连续变化到 $B$. 因此
\[
\begin{aligned}
  \int_{\widehat{AB}}(-y)\mathrm dx+x\mathrm dy
  &=\int_0^\pi[-\sin t\cdot(-\sin t)+\cos t\cdot\cos t]\mathrm dt\\
  &=\int_0^\pi\mathrm dt=\pi.
\end{aligned}
\]

\noindent\keyterm{解法 2}\quad 由 $x^2+y^2=1$ 知 $x\mathrm dx+y\mathrm dy=0$, 得 $\mathrm dy=-\dfrac{x}{y}\mathrm dx$, 从而有
\[
\begin{aligned}
  \int_{\widehat{AB}}(-y)\mathrm dx+x\mathrm dy
  &=\int_1^{-1}(-y)\mathrm dx+\int_1^{-1}x\left(-\frac{x}{y}\right)\mathrm dx\\
  &=\int_{-1}^1\left(\frac{x^2+y^2}{y}\right)\mathrm dx
    =\int_{-1}^1\frac{\mathrm dx}{\sqrt{1-x^2}}
    =\arcsin x\bigg|_{-1}^1=\pi.
\end{aligned}
\]

\begin{example}\label{course-example-16-2-2}\quad 计算第二型曲线积分 $\int_\Gamma (x^2-y^2)\mathrm dx-2xy\mathrm dy$, 其中 $\Gamma$ 是从点 $(0,0)$ 沿曲线 $y=x^\alpha(\alpha>0)$ 到点 $(1,1)$ 的部分。

\end{example}
\begin{solve}\quad 将曲线
\[
  \begin{cases}
    x=x,\\
    y=x^\alpha
  \end{cases}
  \quad(x\in[0,1])
\]
代入得
\[
\begin{aligned}
  \int_\Gamma (x^2-y^2)\mathrm dx-2xy\mathrm dy
  &=\int_0^1(x^2-x^{2\alpha})\mathrm dx
    -\int_0^1 2xx^\alpha(\alpha x^{\alpha-1})\mathrm dx\\
  &=\int_0^1[x^2-(2\alpha+1)x^{2\alpha}]\mathrm dx
    =-\frac23.
\end{aligned}
\]
此题说明, 该第二型曲线积分与 $\alpha$ 的选取无关, 而只与这些曲线的端点有关。

\end{solve}
\begin{example}\label{course-example-16-2-3}\quad 计算第二型曲线积分 $\int_\Gamma x\mathrm dx+y\mathrm dy+z\mathrm dz$, 其中 $\Gamma$ 为球面 $x^2+y^2+z^2=1$ 与平面 $x+y+z=0$ 的交线, 从 $z$ 轴看去取逆时针方向。

\end{example}
\noindent\keyterm{解法 1}\quad 作正交变换将 $\Gamma$ 变到 $O\xi\eta$ 平面:
\[
  \begin{cases}
    \xi=\dfrac1{\sqrt2}(x-y),\\
    \eta=\dfrac1{\sqrt6}(x+y-2z),\\
    \zeta=\dfrac1{\sqrt3}(x+y+z).
  \end{cases}
\]
然后利用 $\xi=\cos t,\eta=\sin t$ 代入上述 $\Gamma$ 的方程, 即得 $\Gamma$ 的参数方程
\[
  \begin{cases}
    x=\dfrac1{\sqrt2}\cos t+\dfrac1{\sqrt6}\sin t,\\
    y=-\dfrac1{\sqrt2}\cos t+\dfrac1{\sqrt6}\sin t,\\
    z=-\dfrac2{\sqrt6}\sin t,
  \end{cases}
\]
从而
\[
\begin{aligned}
  \int_\Gamma x\mathrm dx+y\mathrm dy+z\mathrm dz
  &=\int_0^{2\pi}\bigg[
    \left(\frac1{\sqrt2}\cos t+\frac1{\sqrt6}\sin t\right)
    \left(-\frac1{\sqrt2}\sin t+\frac1{\sqrt6}\cos t\right)\\
  &\quad+
    \left(-\frac1{\sqrt2}\cos t+\frac1{\sqrt6}\sin t\right)
    \left(\frac1{\sqrt2}\sin t+\frac1{\sqrt6}\cos t\right)\\
  &\quad+
    \left(-\frac2{\sqrt6}\sin t\right)
    \left(-\frac2{\sqrt6}\cos t\right)
  \bigg]\mathrm dt\\
  &=\int_0^{2\pi}0\mathrm dt=0.
\end{aligned}
\]

\noindent\keyterm{解法 2}\quad 由于球面 $x^2+y^2+z^2=1$ 上每个点 $(x,y,z)$ 在 $\mathbb R^3$ 中表示的向量即是球面在该点的单位法向量 $\boldsymbol n$, 而曲线 $\Gamma$ 每个点处的单位切向量 $\boldsymbol v$ 与球面在该点的法向量正交, 因此有
\[
  \int_\Gamma x\mathrm dx+y\mathrm dy+z\mathrm dz
  =\int_\Gamma \boldsymbol n\cdot\boldsymbol v\mathrm ds=0.
\]



\section{第一型曲面积分}\label{course-16-3}
本节用于曲线、曲面积分题中的参数化、方向和积分公式。应用 Green、Gauss 或 Stokes 公式之前，应逐项检查光滑性、闭合性、定向以及被积场的奇点；缺失边界必须补上并单独计算。

\subsection{曲面的面积}

在介绍第一型曲面积分之前, 我们先来讨论一下曲面的面积问题. 如何利用积分的思想来求曲面的面积呢? 受求曲线弧长的启发, 我们能否用曲面内接多边形的面积来逼近曲面的面积? 对此施瓦茨曾经有一个著名的例子, 他考虑了由圆柱面的内接全等等腰三角形构成的多面体面积的逼近: 设圆柱面 $S$ 为 $\{(x,y,z):x^2+y^2=1,0\leqslant z\leqslant1\}$. 将 $S$ 的高进行了 $K$ 等分, 从而在 $S$ 上得到 $K+1$ 个平行于 $Oxy$ 平面的圆周 $\Gamma_k(k=1,2,\cdots,K+1)$. 对每个圆周 $\Gamma_k(k=1,2,\cdots,K+1)$ 再 $L$ 等分, 使得 $\Gamma_{k+1}$ 上的分点在 $Oxy$ 平面的投影落在 $\Gamma_k$ 上的分点在 $Oxy$ 平面的投影的中间, 然后将每个 $\Gamma_k$ 上相邻两个分点连接, 并且它们分别与 $\Gamma_{k+1}$ 上的位于它上方中间的点连接, 这样我们就得到了一个等腰三角形. 所有的这些三角形组成了 $S$ 的一个内接多边形. 这个多边形的面积用初等数学的知识即可求出. 容易证明, 适当选取 $K,L$ 的比例, 如 $K=L^3$, 则可以使得多边形的面积趋于 $\infty$. 在上述逼近中, 若取 $K=L$, 则容易推出多边形的面积趋于 $2\pi$. 因此用内接多边形的面积来逼近曲面面积的方法是行不通的。

仔细分析一下施瓦茨的例子, 当 $K$ 与 $L$ 的比很大时, 即三角形的高与底边长的比趋于零时, 每个小三角形所在平面几乎与柱面的切平面垂直, 因此产生很多由于弯曲而增加的面积. 如果该三角形所在平面越来越“平行”于切平面, 则多边形面积逼近将会成功. 注意到对于圆柱面, 我们若将其摊平到平面, 则可以计算出它的面积为
\[
  \text{底周长}\times\text{高}=2\pi.
\]

为了求曲面面积, 我们的想法是用每一点处切平面上的部分来构造曲面的外切多边形进行逼近. 由于要用到切平面, 我们只考虑光滑曲面。

设曲面 $S$ 由参数方程
\begin{equation}
  \begin{cases}
    x=x(u,v),\\
    y=y(u,v),\\
    z=z(u,v),
  \end{cases}
  \quad (u,v)\in D
  \label{course-eq-16-3-1}
\end{equation}
定义, 其中 $D$ 是可求面积的有界闭区域, 函数 $x(u,v),y(u,v),z(u,v)$ 均在 $D$ 上具有连续偏导数. 作为 $D$ 到 $\mathbb R^3$ 的映射, \eqref{course-eq-16-3-1}  是单的且其雅可比矩阵
\[
  \begin{pmatrix}
    x'_u & y'_u & z'_u\\
    x'_v & y'_v & z'_v
  \end{pmatrix}
\]
在每一点 $(u,v)$ 处的秩均为 $2$. 在本节中, 我们总假定曲面 $S$ 为满足上述条件的光滑曲面。

现记 $\boldsymbol r(u,v)=(x(u,v),y(u,v),z(u,v))$ 及
\begin{equation}
  A=\frac{\partial(y,z)}{\partial(u,v)},\quad
  B=\frac{\partial(z,x)}{\partial(u,v)},\quad
  C=\frac{\partial(x,y)}{\partial(u,v)}.
  \label{course-eq-16-3-2}
\end{equation}
则在任一点 $(u,v)$ 处行列式 $A,B,C$ 中至少有一个不为零, 并且它们均是 $(u,v)$ 的连续函数, 因此法向量 $\boldsymbol n=(A,B,C)$ 在曲面 $S$ 上连续变化。

下面我们利用微元法来求曲面的面积. 给定以 $(u_0,v_0),(u_0+\mathrm du,v_0),(u_0+\mathrm du,v_0+\mathrm dv),(u_0,v_0+\mathrm dv)$ 为顶点的一个小矩形 $D_0$. 当 $(u,v)$ 在 $D_0$ 上变化时, 在 $S$ 上可得到一小块曲面, 它可近似看成一个小平行四边形, 其两邻边分别为
\[
  \boldsymbol r(u_0+\mathrm du,v_0)-\boldsymbol r(u_0,v_0)
  \approx \boldsymbol r'_u(u_0,v_0)\mathrm du
\]
与
\[
  \boldsymbol r(u_0,v_0+\mathrm dv)-\boldsymbol r(u_0,v_0)
  \approx \boldsymbol r'_v(u_0,v_0)\mathrm dv.
\]
因此它们张成的平行四边形的面积, 即面积微元为
\[
  \mathrm dS(u_0,v_0)=|\boldsymbol r'_u(u_0,v_0)\times\boldsymbol r'_v(u_0,v_0)|\mathrm du\mathrm dv,
\]
从而曲面 $S$ 的面积 (仍用 $S$ 表示其面积) 为
\[
  S=\iint_D |\boldsymbol r'_u(u,v)\times\boldsymbol r'_v(u,v)|\mathrm du\mathrm dv.
\]
现在我们进一步计算 $|\boldsymbol r'_u(u,v)\times\boldsymbol r'_v(u,v)|$. 由
\[
  |\boldsymbol r'_u(u,v)\times\boldsymbol r'_v(u,v)|
  =|(x'_u,y'_u,z'_u)\times(x'_v,y'_v,z'_v)|
  =\sqrt{A^2+B^2+C^2},
\]
若记
\begin{equation}
  \begin{cases}
    E=\boldsymbol r'_u\boldsymbol r'_u=x_u'^2+y_u'^2+z_u'^2,\\
    F=\boldsymbol r'_u\boldsymbol r'_v=x'_ux'_v+y'_uy'_v+z'_uz'_v,\\
    G=\boldsymbol r'_v\boldsymbol r'_v=x_v'^2+y_v'^2+z_v'^2,
  \end{cases}
  \label{course-eq-16-3-3}
\end{equation}
则可知
\[
  EG-F^2=A^2+B^2+C^2.
\]
因此, 我们有下面的定理。

\begin{theorem}\label{course-theorem-16-3-1}\quad 设光滑曲面 $S$ 由参数方程 \eqref{course-eq-16-3-1}  给出, 则它的面积为
\[
  S=\iint_D \sqrt{A^2+B^2+C^2}\mathrm du\mathrm dv
  =\iint_D \sqrt{EG-F^2}\mathrm du\mathrm dv,
\]
其中 $A,B,C$ 和 $E,F,G$ 的定义分别见式~\eqref{course-eq-16-3-2}  和 \eqref{course-eq-16-3-3}。

特别地, 当光滑曲面 $S$ 是由 $z=f(x,y)((x,y)\in D)$ 给出时, 则有
\[
  S=\iint_D \sqrt{1+f_x'^2(x,y)+f_y'^2(x,y)}\mathrm dx\mathrm dy.
\]

\end{theorem}
\begin{example}\label{course-example-16-3-1}\quad 求单位圆柱面 $x^2+y^2=1$ 在平面 $z=0$ 与 $z=1$ 之间的面积。

\end{example}
\begin{solve}\quad 设所求面积为 $S$, 则 $S$ 为曲面 $y=f(z,x)=\sqrt{1-x^2}((z,x)\in D,\text{ 其中 }D=[0,1]\times[-1,1])$ 的面积的两倍. 因此有
\[
\begin{aligned}
  S&=2\iint_D\sqrt{1+f_x'^2+f_z'^2}\mathrm dz\mathrm dx\\
  &=2\int_0^1\mathrm dz\int_{-1}^1\sqrt{1+\left(\frac{-x}{\sqrt{1-x^2}}\right)^2}\mathrm dx\\
  &=2\int_{-1}^1\frac{\mathrm dx}{\sqrt{1-x^2}}=2\pi.
\end{aligned}
\]

\end{solve}
\begin{example}\label{course-example-16-3-2}\quad 计算球面 $x^2+y^2+z^2=1$ 被柱面 $\left(x-\dfrac12\right)^2+y^2=\dfrac14$ 所截在柱面内的部分的面积。

\end{example}
\begin{solve}\quad 由对称性, 所求面积 $S$ 为所截得曲面在 $Oxy$ 平面上方部分面积的两倍. 设 $z=f(x,y)=\sqrt{1-x^2-y^2}((x,y)\in D)$, 其中
\[
  D=\left\{(x,y):\left(x-\frac12\right)^2+y^2\leqslant\frac14\right\},
\]
则
\[
  S=2\iint_D\sqrt{1+f_x'^2+f_y'^2}\mathrm dx\mathrm dy
  =2\iint_D\frac{\mathrm dx\mathrm dy}{\sqrt{1-x^2-y^2}}.
\]
利用极坐标变换得
\[
\begin{aligned}
  S&=4\int_0^{\frac\pi2}\mathrm d\theta\int_0^{\cos\theta}
  \frac{r\mathrm dr}{\sqrt{1-r^2}}
  =4\int_0^{\frac\pi2}\left(-\sqrt{1-r^2}\right)\bigg|_0^{\cos\theta}\mathrm d\theta\\
  &=4\int_0^{\frac\pi2}(1-\sin\theta)\mathrm d\theta=2\pi-4.
\end{aligned}
\]

\end{solve}
\subsection{第一型曲面积分的定义}

由于曲面的复杂性, 我们在这里只讨论光滑曲面上的第一型曲面积分. 具体来说, 在本小节中我们所指的曲面 $S$ 将由参数方程 \eqref{course-eq-16-3-1}  定义, 并且它们满足所假定的光滑条件. 特别地, 我们总假定闭区域 $D\subset\mathbb R^2$ 的边界由有限条光滑曲线组成. 设
\[
  \Delta=\{\Delta D_1,\Delta D_2,\cdots,\Delta D_K\}
\]
为 $D$ 的一个分割, 且假定每个 $\Delta D_k(k=1,2,\cdots,K)$ 均为由有限条光滑曲线所围成的区域. 在此假定下, $\Delta S_k=\{\boldsymbol r(u,v):(u,v)\in\Delta D_k\}(k=1,2,\cdots,K)$ 是一块光滑小曲面, $T=\{\Delta S_1,\Delta S_2,\cdots,\Delta S_K\}$ 是 $S$ 的一个分割. 记 $\lambda(\Delta)=\max\limits_{1\leqslant k\leqslant K}\operatorname{diam}(\Delta D_k)$ 与 $\lambda(T)=\max\limits_{1\leqslant k\leqslant K}\operatorname{diam}(\Delta S_k)$. 由于 $\boldsymbol r(u,v)$ 在 $D$ 上一致连续且其逆映射在 $S$ 上一致连续, 因此 $\lambda(\Delta)\to0$ 的充分必要条件是 $\lambda(T)\to0$。

对于 $S$ 的分割 $T=\{\Delta S_1,\Delta S_2,\cdots,\Delta S_K\}$, 以下我们总是假定每个 $\Delta S_k$ 均是 $S$ 上可求面积的一小块曲面, 并且存在 $D$ 的可求面积的闭子区域 $\Delta D_k$, 使得 $\boldsymbol r(u,v)$ 是 $\Delta D_k$ 到 $\Delta S_k$ 的一一对应, 从而 $\Delta=\{\Delta D_1,\Delta D_2,\cdots,\Delta D_K\}$ 为 $D$ 的一个分割。

\begin{definition}\label{course-definition-16-3-1}\quad 设 $S\subset\mathbb R^3$ 是光滑曲面, 函数 $f(x,y,z)$ 在 $S$ 上有定义, 又设\allowbreak $T=\{\Delta S_1,\Delta S_2,\cdots,\Delta S_K\}$ 是 $S$ 的一个分割, 且仍用 $\Delta S_k(k=1,2,\cdots,K)$ 来记每块小曲面的面积, 并记
\[
  \lambda(T)=\max_{1\leqslant k\leqslant K}\operatorname{diam}(\Delta S_k).
\]
在 $\Delta S_k$ 上任取一点 $(\xi_k,\eta_k,\zeta_k)$, 作和式 $\sum\limits_{k=1}^K f(\xi_k,\eta_k,\zeta_k)\Delta S_k$. 若存在常数 $I$, 使得对于 $\forall\varepsilon>0$, $\exists\delta>0$, 对于 $S$ 的任意分割 $T$ 及任取的 $(\xi_k,\eta_k,\zeta_k)\in\Delta S_k(k=1,2,\cdots,K)$, 当 $\lambda(T)<\delta$ 时, 有
\[
  \left|\sum_{k=1}^K f(\xi_k,\eta_k,\zeta_k)\Delta S_k-I\right|<\varepsilon,
\]
即
\[
  \lim_{\lambda(T)\to0}\sum_{k=1}^K f(\xi_k,\eta_k,\zeta_k)\Delta S_k=I,
\]
则称 $f(x,y,z)$ 在 $S$ 上的第一型曲面积分存在, 并称 $I$ 为 $f(x,y,z)$ 在 $S$ 上的第一型曲面积分, 记为
\[
  I=\iint_S f(x,y,z)\mathrm dS,
\]
其中 $f(x,y,z)$ 称为被积函数, $S$ 称为积分曲面。

从定义可以推出第一型曲面积分具有以下性质:

(1) 若在 $S$ 上函数 $f(x,y,z)\equiv1$, 则 $\iint_S f(x,y,z)\mathrm dS=\iint_S\mathrm dS$ 即为曲面 $S$ 的面积。

(2) 若曲面 $S\subset\mathbb R^2$, 则函数 $f(x,y)$ 在 $S$ 上的曲面积分即为二重积分 $\iint_S f(x,y)\mathrm dx\mathrm dy$。

另外, 第一型曲面积分同样具有关于被积函数的线性性以及积分曲面的可加性等性质, 请读者自己给出上述性质的精确描述。

\end{definition}
\subsection{第一型曲面积分的存在性与计算公式}

对于光滑曲面上连续函数的第一型曲面积分, 我们有以下的存在性定理以及计算公式。

\begin{theorem}\label{course-theorem-16-3-2}\quad 设 $S\subset\mathbb R^3$ 是光滑曲面, 其参数方程为
\[
  \begin{cases}
    x=x(u,v),\\
    y=y(u,v),\\
    z=z(u,v),
  \end{cases}
  \quad (u,v)\in D,
\]
其中 $D$ 是由有限条光滑曲线围成的有界闭区域, 又设函数 $f(x,y,z)$ 在 $S$ 上连续, 则 $f(x,y,z)$ 在 $S$ 上的第一型曲面积分存在, 并且
\begin{equation}
  \iint_S f(x,y,z)\mathrm dS
  =\iint_D f(x(u,v),y(u,v),z(u,v))\sqrt{EG-F^2}\mathrm du\mathrm dv,
  \label{course-eq-16-3-4}
\end{equation}
其中 $E,G,F$ 由式~\eqref{course-eq-16-3-3}  定义。

\end{theorem}
\begin{proof}\quad 考查曲面 $S$ 的任一分割 $T=\{\Delta S_1,\Delta S_2,\cdots,\Delta S_K\}$, 记其相应于 $D$ 的分割为 $\Delta=\{\Delta D_1,\Delta D_2,\cdots,\Delta D_K\}$. 在每个 $\Delta S_k(k=1,2,\cdots,K)$ 上任取点 $(\xi_k,\zeta_k,\eta_k)$, 则存在 $(u_k,v_k)\in\Delta D_k$, 使得
\[
  (\xi_k,\zeta_k,\eta_k)=(x(u_k,v_k),y(u_k,v_k),z(u_k,v_k)).
\]
因为
\[
  \Delta S_k=\iint_{\Delta D_k}\sqrt{EG-F^2}\mathrm du\mathrm dv,
\]
所以由重积分第一中值定理得
\[
  \Delta S_k=\sqrt{EG-F^2}\bigg|_{(u'_k,v'_k)}\Delta\sigma_k,
\]
其中 $(u'_k,v'_k)\in\Delta D_k,\Delta\sigma_k$ 为 $\Delta D_k$ 的面积. 因此
\[
\begin{aligned}
  &\sum_{k=1}^K f(\xi_k,\zeta_k,\eta_k)\Delta S_k\\
  &\quad=\sum_{k=1}^K f(x(u_k,v_k),y(u_k,v_k),z(u_k,v_k))
    \sqrt{EG-F^2}\bigg|_{(u'_k,v'_k)}\Delta\sigma_k\\
  &\quad=\sum_{k=1}^K f(x(u_k,v_k),y(u_k,v_k),z(u_k,v_k))
    \sqrt{EG-F^2}\bigg|_{(u_k,v_k)}\Delta\sigma_k\\
  &\qquad+\bigg[
    \sum_{k=1}^K f(x(u_k,v_k),y(u_k,v_k),z(u_k,v_k))
    \sqrt{EG-F^2}\bigg|_{(u'_k,v'_k)}\Delta\sigma_k\\
  &\qquad\quad-
    \sum_{k=1}^K f(x(u_k,v_k),y(u_k,v_k),z(u_k,v_k))
    \sqrt{EG-F^2}\bigg|_{(u_k,v_k)}\Delta\sigma_k
  \bigg]\\
  &\quad\triangleq I_1+I_2.
\end{aligned}
\]
显然
\[
  \lim_{\lambda(T)\to0}I_1=\lim_{\lambda(\Delta)\to0}I_1
  =\iint_D f(x(u,v),y(u,v),z(u,v))\sqrt{EG-F^2}\mathrm du\mathrm dv.
\]
记 $M$ 为 $|f(x,y,z)|$ 在 $S$ 上的最大值, $M_k,m_k$ 分别为连续函数 $\sqrt{EG-F^2}$ 在 $\Delta D_k(k=1,2,\cdots,K)$ 上的最大、最小值, 则有
\[
  |I_2|\leqslant M\sum_{k=1}^K(M_k-m_k)\Delta\sigma_k\to0
  \quad(\lambda(\Delta)\to0).
\]
因此
\[
\begin{aligned}
  \lim_{\lambda(T)\to0}\sum_{k=1}^K f(\xi_k,\zeta_k,\eta_k)\Delta S_k
  &=\lim_{\lambda(T)\to0}(I_1+I_2)
    =\lim_{\lambda(\Delta)\to0}(I_1+I_2)\\
  &=\iint_D f(x(u,v),y(u,v),z(u,v))\sqrt{EG-F^2}\mathrm du\mathrm dv.
\end{aligned}
\]
这就证明了 $f(x,y,z)$ 在 $S$ 上的第一型曲面积分存在, 并且
\[
  \iint_S f(x,y,z)\mathrm dS
  =\iint_D f(x(u,v),y(u,v),z(u,v))\sqrt{EG-F^2}\mathrm du\mathrm dv.
\]
证毕。

当曲面 $S$ 由方程 $z=g(x,y)((x,y)\in D)$ 给出时, 由定理~\ref{course-theorem-16-3-2}  有
\[
  \iint_S f(x,y,z)\mathrm dS
  =\iint_D f(x,y,g(x,y))\sqrt{1+g_x'^2(x,y)+g_y'^2(x,y)}\mathrm dx\mathrm dy.
\]

\end{proof}
\begin{example}\label{course-example-16-3-3}\quad 求第一型曲面积分 $I=\iint_S x^2z\mathrm dS$, 其中 $S$ 为以下的曲面:

(1) 上半单位球面 $z=\sqrt{1-x^2-y^2}$;

(2) 单位球面 $x^2+y^2+z^2=1$。

\end{example}
\begin{solve}\quad (1) 取 $D=\{(x,y):x^2+y^2\leqslant1\}$, 则
\[
\begin{aligned}
  I&=\iint_S x^2z\mathrm dS\\
  &=\iint_D x^2\cdot\sqrt{1-x^2-y^2}
    \sqrt{1+\frac{x^2}{1-x^2-y^2}+\frac{y^2}{1-x^2-y^2}}\mathrm dx\mathrm dy\\
  &=\iint_D x^2\mathrm dx\mathrm dy
    =\frac12\iint_D(x^2+y^2)\mathrm dx\mathrm dy\\
  &=\frac12\int_0^{2\pi}\mathrm d\theta\int_0^1r^3\mathrm dr=\frac\pi4.
\end{aligned}
\]

(2) 分别记
\[
  S_1:z=\sqrt{1-x^2-y^2},\quad (x,y)\in D,
\]
\[
  S_2:z=-\sqrt{1-x^2-y^2},\quad (x,y)\in D,
\]
则由积分曲面及被积函数的对称性得
\[
  \iint_{S_1}x^2z\mathrm dS=-\iint_{S_2}x^2z\mathrm dS,
\]
所以
\[
  \iint_S x^2z\mathrm dS=0.
\]

\end{solve}
\begin{example}\label{course-example-16-3-4}\quad 求均匀物质曲面 $S:z=f(x,y)=2-(x^2+y^2)(z\geqslant0)$ 的质心坐标。

\end{example}
\begin{solve}\quad 设其质心坐标为 $(x_0,y_0,z_0)$, 由对称性有 $x_0=y_0=0$, 并且由微元法有
\[
  z_0=\frac{M_{xy}}{\sigma}=\frac{\iint_S z\mathrm dS}{\iint_S\mathrm dS}
\]
(其中 $M_{xy}$ 是曲面 $S$ 到 $Oxy$ 平面的力矩, $\sigma$ 是曲面 $S$ 的面积). 取 $D=\{(x,y):x^2+y^2\leqslant2\}$. 由 $\sqrt{1+f_x'^2+f_y'^2}=\sqrt{1+4x^2+4y^2}$ 有
\[
\begin{aligned}
  \iint_S\mathrm dS
  &=\iint_D\sqrt{1+4x^2+4y^2}\mathrm dx\mathrm dy
    =\int_0^{2\pi}\mathrm d\theta\int_0^{\sqrt2}\sqrt{1+4r^2}r\mathrm dr\\
  &=2\pi\cdot\frac1{12}(1+4r^2)^{\frac32}\bigg|_0^{\sqrt2}
    =\frac{13}3\pi,
\end{aligned}
\]
\[
\begin{aligned}
  \iint_S z\mathrm dS
  &=\iint_D(2-x^2-y^2)\sqrt{1+4x^2+4y^2}\mathrm dx\mathrm dy\\
  &=\int_0^{2\pi}\mathrm d\theta\int_0^{\sqrt2}r(2-r^2)\sqrt{1+4r^2}\mathrm dr
    =\frac{37}{10}\pi,
\end{aligned}
\]
所以
\[
  z_0=\frac{111}{130}.
\]

\end{solve}


\section{第二型曲面积分}\label{course-16-4}
本节用于曲线、曲面积分题中的参数化、方向和积分公式。应用 Green、Gauss 或 Stokes 公式之前，应逐项检查光滑性、闭合性、定向以及被积场的奇点；缺失边界必须补上并单独计算。

\subsection{曲面的侧}

在日常生活中, 我们常见的曲面大都具有不同的两侧, 如一个自来水管具有内侧与外侧, 当水管破裂时, 管内的水将从内侧流向外侧。

是不是所有的曲面都能分出两侧呢? 著名的麦比乌斯 (Möbius) 带就是一个单侧曲面. 该曲面的构造非常简单: 取一条长方形纸带, 记为 $ABCD$, 将纸带扭转后, 让 $A$ 与 $C$, $B$ 与 $D$ 重合后粘起来, 使纸带构成一个环带 $S$, 该环带则称为麦比乌斯带. 原长方形的边 $AB$ 与 $DC$ 构成该环带 $S$ 的边界. 用颜色来涂该环带, 可以不越过其边界, 而将它全部涂遍颜色. 由此我们无法分出该环带不同的侧 。

如何用精确的语言来描述一个曲面 $S$ 是单侧还是双侧的呢? 我们有以下定义。

\begin{definition}\label{course-definition-16-4-1}\quad 设 $S$ 是光滑曲面 ($S$ 可以是封闭的, 则此时无边界, 否则它具有边界). 对 $S$ 内部的任一点 $P_0$, 取定 $S$ 在 $P_0$ 处的一个法向的朝向. 若一个动点 $P$ 从 $P_0$ 出发, 不越过 $S$ 的边界, 沿 $S$ 上任何路径 $\Gamma$ 运动回到 $P_0$ 处, $S$ 在 $P$ 的法向从 $P_0$ 处选定的朝向出发连续地沿路径 $\Gamma$ 变化回到 $P_0$ 处时, 总是保持原先在 $P_0$ 处选定的朝向, 则称 $S$ 为双侧曲面, 否则称 $S$ 为单侧曲面。

由定义, 一个单侧曲面 $S$ 上必定存在点 $P_0\in S$, 以及 $S$ 上的一条经过 $P_0$ 的闭路径 $\Gamma$, 若在 $P_0$ 处选择 $S$ 的一个法向的朝向, 然后让 $P$ 从 $P_0$ 出发沿 $\Gamma$ 运动, $P$ 处的法向连续地变化, 当 $P$ 经过 $\Gamma$ 再回到 $P_0$ 时, 法向的朝向与原来确定的朝向相反。

现设 $S$ 是一双侧曲面, 从定义可以看出在 $S$ 上任取一点 $P_0$, 然后取定 $S$ 在 $P_0$ 处一个法向的朝向, 则该曲面一侧法向的朝向即由该处法向的朝向所取定. 这个事实请读者自己来验证. 利用此事实, 我们可以容易地描述由一些参数方程表示的双侧曲面的定侧问题。

例如, 若方程
\[
  \begin{cases}
    x=x(u,v),\\
    y=y(u,v),\\
    z=z(u,v),
  \end{cases}\quad (u,v)\in D
\]
定义了一个光滑的双侧曲面 $S$, 令
\begin{equation}
  A=\frac{\partial(y,z)}{\partial(u,v)},\quad
  B=\frac{\partial(z,x)}{\partial(u,v)},\quad
  C=\frac{\partial(x,y)}{\partial(u,v)},
  \label{course-eq-16-4-1}
\end{equation}
则 $\boldsymbol n=\pm(A,B,C)$ 即为 $S$ 的两个不同朝向的法向量. 记 $\boldsymbol n_1=(A,B,C)$, $\boldsymbol n_2=(-A,-B,-C)$. 当 $P_0\in S$, 且在 $P_0$ 处 $C>0$ 时, 我们将由 $\boldsymbol n_1$ 确定的一侧称为 $S$ 的上侧, $\boldsymbol n_2$ 确定的一侧称为 $S$ 下侧. 类似地, 当 $A\ne0$ 时, 我们可以说 $S$ 有前侧和后侧, 而当 $B\ne0$ 时, $S$ 有左侧和右侧。

若光滑曲面 $S$ 的方程为
\[
  z=f(x,y),\quad (x,y)\in D,
\]
则 $S$ 上点 $(x,y,z)$ 处的法向量为 $\boldsymbol n=\pm(-f_x',-f_y',1)$. 因此, 若 $\boldsymbol n$ 取 ``+'' 号, 便得到了 $S$ 的上侧; 若 $\boldsymbol n$ 取 ``-'' 号, 则得到 $S$ 的下侧。

\end{definition}
\subsection{第二型曲面积分的定义}

第二型曲面积分的物理背景是流量的计算问题. 设区域 $D\subset\mathbb R^3$ 内有某一流体在流动 (如风的运动, 海水的运动等), 对于 $\forall(x,y,z)\in D$, 其速度函数为
\[
  \boldsymbol v(x,y,z)=(P(x,y,z),Q(x,y,z),R(x,y,z)),
\]
简记为 $\boldsymbol v=(P,Q,R)$. 若不考虑流体的种类, 上述速度函数也可认为是 $D$ 内的一个流速场. 进一步假设 $\boldsymbol v$ 仅仅依赖于 $D$ 内点的位置, 而不依赖于时间的变化, 这样的流速场也称为稳定流速场。

设 $S$ 是 $D$ 内的一个光滑双侧曲面. 我们的问题是: 如何来求单位时间内流体流过 $S$ 的质量? 为了使问题简单化, 我们可设该流体的密度为 1. 因此, 我们只要计算该流体通过 $S$ 的流量即可。

若 $S$ 是平面的一部分, 其面积仍由 $S$ 表示, 而 $\boldsymbol v$ 是常向量函数, 即 $P=C_1,Q=C_2,R=C_3$ 均为常数函数, 记 $S$ 的单位法向量为 $\boldsymbol n$, 则不难看出单位时间内该流体流过 $S$ 的流量
\[
  W=(\boldsymbol v\cdot\boldsymbol n)S,
\]
即流量 $W$ 为 $S$ 的面积与 $\boldsymbol v$ 在 $\boldsymbol n$ 上的投影的乘积. 有了上述的计算公式, 利用积分思想就不难推导出在一般情形下流量的计算公式了。

我们下面仍用微元法来求解此问题. 在 $S$ 上任取一小块 $\Delta S$, 其面积仍然用 $\Delta S$ 记之. 任取 $(\xi,\eta,\zeta)\in\Delta S$, 设 $\Delta S$ 在 $(\xi,\eta,\zeta)$ 处的单位法向量为 $\boldsymbol n(\xi,\eta,\zeta)$, 则可以认为 $\Delta S$ 近似地是法向量为 $\boldsymbol n(\xi,\eta,\zeta)$, 面积为 $\Delta S$ 的一小块平面. 再假设 $\boldsymbol v$ 在 $\Delta S$ 上也近似为 $\boldsymbol v(\xi,\eta,\zeta)$, 则由上面的分析, 我们便得到了流量微元
\[
  \mathrm dW=\boldsymbol v(x,y,z)\cdot\boldsymbol n(x,y,z)\mathrm dS.
\]
因此, 整个流量便是
\[
  W=\iint_S\boldsymbol v\cdot\boldsymbol n\,\mathrm dS
  =\iint_S(P\cos\alpha+Q\cos\beta+R\cos\gamma)\mathrm dS,
\]
其中 $\cos\alpha,\cos\beta,\cos\gamma$ 分别是 $\boldsymbol n$ 的方向余弦. 由于上述积分中含有曲面法向量的方向余弦, 因此它不是第一型曲面积分. 它其实是下面我们要讨论的第二型曲面积分。

对于第二型曲面积分所涉及的曲面 $S$, 我们假定它是分片光滑的双侧曲面. 这类曲面由有限块光滑曲面组成, 任何两块曲面均在边界相交, 并且不存在三块曲面具有相同一段弧作为边界. 具有三块曲面共边等非流形接合的对象不在这里的讨论范围内。

现在我们给出分片光滑双侧曲面上的第二型曲面积分的定义。

\begin{definition}\label{course-definition-16-4-2}\quad 设 $S\subset\mathbb R^3$ 是分片光滑双侧曲面, 若它有边界, 则其边界由有限条分段光滑曲线组成. 给定 $S$ 的一侧, $S$ 上每点处的单位法向量记为 $\boldsymbol n=(\cos\alpha,\cos\beta,\cos\gamma)$, 向量函数 $\boldsymbol f(x,y,z)=(P(x,y,z),Q(x,y,z),R(x,y,z))$ 在 $S$ 上有定义. 对 $S$ 作任意分割 $T=\{\Delta S_1,\Delta S_2,\cdots,\Delta S_K\}$, 其中 $\Delta S_k$ 是以光滑曲线为边界的小光滑曲面, 并且仍用 $\Delta S_k(k=1,2,\cdots,K)$ 来记其面积. 记 $\lambda(T)=\max_{1\leqslant k\leqslant K}\operatorname{diam}(\Delta S_k)$. 在 $\Delta S_k$ 上任取一点 $(\xi_k,\eta_k,\zeta_k)(k=1,2,\cdots,K)$. 若存在不依赖于分割 $T$ 以及 $(\xi_k,\eta_k,\zeta_k)(k=1,2,\cdots,K)$ 的选取的常数 $I$, 使得
\begin{equation}
\begin{aligned}
  \lim_{\lambda(T)\to0}\sum_{k=1}^K
  (&P(\xi_k,\eta_k,\zeta_k)\cos\alpha
  +Q(\xi_k,\eta_k,\zeta_k)\cos\beta\\
  &+R(\xi_k,\eta_k,\zeta_k)\cos\gamma)\Delta S_k=I,
\end{aligned}
\label{course-eq-16-4-2}
\end{equation}
则称 $I$ 是 $\boldsymbol f=(P,Q,R)$ 在 $S$ 上的第二型曲面积分, 记为
\[
  I=\iint_S P(x,y,z)\mathrm dy\mathrm dz
    +Q(x,y,z)\mathrm dz\mathrm dx
    +R(x,y,z)\mathrm dx\mathrm dy,
\]
其中 $\boldsymbol f=(P,Q,R)$ 称为被积函数, $S$ 称为积分曲面。

\end{definition}
\noindent\keyterm{注 1}\quad 我们来解释一下第二型曲面积分记号的合理性. 由定义, 一个第二型曲面积分是式~\eqref{course-eq-16-4-2}  定义的极限. 当小曲面 $\Delta S_k$ 的直径很小时, $\Delta S_k$ 近似于 $(\xi_k,\eta_k,\zeta_k)$ 处一小块切平面, 因此它的面积与 $\cos\alpha$ 的乘积近似为该小块切平面在 $Oyz$ 平面上的投影 $\Delta\sigma^k_{yz}$ 的有向面积, 即当 $\cos\alpha>0$ 时, 取其投影的面积为正, 否则为负. 因此, 若仍由 $\Delta\sigma^k_{yz}$ 记其有向面积, 则
\[
\begin{aligned}
  &\lim_{\lambda(T)\to0}\sum_{k=1}^K
    P(\xi_k,\eta_k,\zeta_k)\cos\alpha\Delta S_k\\
  &\quad=\lim_{\lambda(T)\to0}\sum_{k=1}^K
    P(\xi_k,\eta_k,\zeta_k)\Delta\sigma^k_{yz}.
\end{aligned}
\]
上式右边可以自然地记成 $\iint_S P\,\mathrm dy\mathrm dz$. 同理我们可解释记号 $\iint_S Q\,\mathrm dz\mathrm dx$ 和 $\iint_S R\,\mathrm dx\mathrm dy$ 的合理性。

对于第二型曲面积分, 读者不难看出它们仍然具有关于被积函数的线性性以及积分曲面的相加性. 值得注意的是, 当曲面相加时要指定曲面的侧. 特别需要指出的是: 双侧曲面上同一个向量函数在两个侧上的第二型曲面积分的绝对值相等, 但相差一个负号。

\subsection{第二型曲面积分的存在性与计算公式}

设光滑曲面 $S$ 由参数方程
\begin{equation}
  \begin{cases}
    x=x(u,v),\\
    y=y(u,v),\\
    z=z(u,v),
  \end{cases}\quad (u,v)\in D
  \label{course-eq-16-4-3}
\end{equation}
给出, 则 $S$ 上每一点处的方向余弦为
\[
  \cos\alpha=\pm\frac{A}{\sqrt{A^2+B^2+C^2}},\quad
  \cos\beta=\pm\frac{B}{\sqrt{A^2+B^2+C^2}},
\]
\[
  \cos\gamma=\pm\frac{C}{\sqrt{A^2+B^2+C^2}},
\]
其中 $A,B,C$ 由式~\eqref{course-eq-16-4-1}  定义. 值得指出的是, 它们都是 $(u,v)\in D$ 的连续函数。

\begin{theorem}\label{course-theorem-16-4-1}\quad 设 $S$ 为光滑双侧曲面, 其参数方程由 \eqref{course-eq-16-4-3}  给出. 选定 $S$ 的一侧, 记其单位法向量为
\[
  \boldsymbol n=\left(
    \frac{A}{\sqrt{A^2+B^2+C^2}},
    \frac{B}{\sqrt{A^2+B^2+C^2}},
    \frac{C}{\sqrt{A^2+B^2+C^2}}
  \right),
\]
再设向量函数 $\boldsymbol f(x,y,z)=(P,Q,R)$ 在 $S$ 上连续, 则 $\boldsymbol f$ 在 $S$ 上的第二型曲面积分存在, 并且
\[
  \iint_S P\,\mathrm dy\mathrm dz+Q\,\mathrm dz\mathrm dx+R\,\mathrm dx\mathrm dy
  =\iint_D(PA+QB+RC)\,\mathrm du\mathrm dv.
\]

\end{theorem}
\begin{proof}\quad 由 $S$ 的光滑性及式~\eqref{course-eq-16-4-1}  知 $\boldsymbol n$ 在 $S$ 上是连续向量函数. 由
\begin{equation}
  \iint_S P\,\mathrm dy\mathrm dz+Q\,\mathrm dz\mathrm dx+R\,\mathrm dx\mathrm dy
  =\iint_S(P\cos\alpha+Q\cos\beta+R\cos\gamma)\mathrm dS,
  \label{course-eq-16-4-4}
\end{equation}
并注意到上式的右边是连续函数 $P\cos\alpha+Q\cos\beta+R\cos\gamma$ 的第一型曲面积分, 因此左边的第二型曲面积分存在. 定理的可积性部分得证。

由公式~\eqref{course-eq-16-3-4}  得
\[
  \iint_S(P\cos\alpha+Q\cos\beta+R\cos\gamma)\mathrm dS
  =\iint_D(PA+QB+RC)\,\mathrm du\mathrm dv.
\]
证毕。

特别地, 若光滑曲面 $S$ 的方程为
\[
  z=f(x,y),\quad (x,y)\in D,
\]
并取定其上侧, 其中 $D\subset\mathbb R^2$ 是由分段光滑曲线所围成的有界闭区域. 由于 $S$ 上的单位法向量平行于 $(-f_x',-f_y',1)$, 容易推出
\[
  \iint_S P(x,y,z)\mathrm dy\mathrm dz+Q(x,y,z)\mathrm dz\mathrm dx+R(x,y,z)\mathrm dx\mathrm dy
\]
\[
  =\iint_D[-P(x,y,f(x,y))f_x'(x,y)-Q(x,y,f(x,y))f_y'(x,y)
  +R(x,y,f(x,y))]\mathrm dx\mathrm dy.
\]

\end{proof}
\begin{example}\label{course-example-16-4-2}\quad 计算第二型曲面积分
\[
  I=\iint_S x\mathrm dy\mathrm dz+y\mathrm dz\mathrm dx+z\mathrm dx\mathrm dy,
\]
其中 $S$ 为由三个坐标平面及平面 $x+y+z=1$ 所围四面体的外侧。

\end{example}
\noindent\keyterm{解法 1}\quad 记 $S$ 落在 $Oxy,Oyz,Ozx$ 平面上的部分分别记为 $S_z,S_x$ 及 $S_y$, 落在平面 $x+y+z=1$ 的部分记为 $S_1$ . 在 $S_z$ 上, $z=0$, $\mathrm dy\mathrm dz=\mathrm dz\mathrm dx=0$, 从而
\[
  \iint_{S_z}x\mathrm dy\mathrm dz+y\mathrm dz\mathrm dx+z\mathrm dx\mathrm dy=0.
\]
同理, 在 $S_y$ 与 $S_x$ 上的积分都为零. 因此
\[
  I=\iint_{S_1}x\mathrm dy\mathrm dz+y\mathrm dz\mathrm dx+z\mathrm dx\mathrm dy.
\]
记 $D=\{(x,y):x\geqslant0,y\geqslant0,x+y\leqslant1\}$, 则由对称性有
\[
  I=3\iint_D(1-x-y)\mathrm dx\mathrm dy
  =3\int_0^1\mathrm dx\int_0^{1-x}(1-x-y)\mathrm dy=\frac12.
\]

\noindent\keyterm{解法 2}\quad 与解法 1 同理得到
\[
  I=\iint_{S_1}x\mathrm dy\mathrm dz+y\mathrm dz\mathrm dx+z\mathrm dx\mathrm dy.
\]
由于在 $S_1$ 上点 $(x,y,z)$ 处的方向余弦 $(\cos\alpha,\cos\beta,\cos\gamma)$ 即为 $S_1$ 的单位法向量, 且
\[
  \cos\alpha=\cos\beta=\cos\gamma=\frac1{\sqrt3},
\]
而 $S_1$ 是一个三角形, 其面积为 $\frac{\sqrt3}{2}$, 因此
\[
  I=\iint_{S_1}\frac{x+y+z}{\sqrt3}\mathrm dS
  =\frac1{\sqrt3}\iint_{S_1}\mathrm dS=\frac12.
\]



\section{各类积分之间的联系}\label{course-16-5}
本节用于曲线、曲面积分题中的参数化、方向和积分公式。应用 Green、Gauss 或 Stokes 公式之前，应逐项检查光滑性、闭合性、定向以及被积场的奇点；缺失边界必须补上并单独计算。

设 $f(x)\in C^1[a,b]$. 众所周知, 牛顿-莱布尼茨公式
\[
  \int_a^b\mathrm df(x)=f(b)-f(a)
\]
是一元微积分中最重要的公式. 换种观点来看该公式, 如果我们定义一个函数 $f(x)$ 在 $a,b(a<b)$ 两点上的积分为 $f(b)-f(a)$, 则上述公式建立了区间 $[a,b]$ 上 $f'(x)$ 的积分与 $f(x)$ 在端点处的积分的联系. 因此, 一个自然的问题是: 是否平面或空间中一个几何体的边界与内部上的某些积分总存在联系? 本节中我们将讨论这个问题。

\subsection{格林公式}

格林 (Green) 公式给出的是平面区域上函数的二重积分与其相关的函数在边界上的第二型曲线积分之间的联系. 在这里, 我们只讨论 $\mathbb R^2$ 中的有界闭区域 $D$, 其边界是有限条约当曲线, 且每条约当曲线都是分段光滑的. 允许区域有有限多个孔洞。

对于平面 $\mathbb R^2$ 内一条约当曲线 $\Gamma$, 以 $\Gamma$ 为边界的区域有两个: 一个是有界区域, 该区域也称为 $\Gamma$ 的内部; 而另一个是无界区域, 该区域也称为 $\Gamma$ 的外部. 例如有一个外边界和若干内边界时，$D$ 位于 $\Gamma_0$ 的内部, 而在 $\Gamma_k(k=1,2,\cdots,K)$ 的外部. 一个区域 $D$ 称为单连通区域, 如果位于 $D$ 内的任一条约当曲线的内部总包含在 $D$ 内; 否则称 $D$ 是多连通区域. 在我们假定的情形下, 若 $D$ 是一条约当曲线的内部, 则它是单连通的; 若 $D$ 由 $K(K>1)$ 条约当曲线围成, 则 $D$ 是多连通的, 此时也称 $D$ 是 $K$ 连通的。

一条约当曲线具有两个互为相反的方向. 现设 $D$ 是由有限条约当曲线所围成的区域. 这里规定 $\partial D$ 的定向, 即若一个人沿着 $D$ 的边界曲线上的一个方向前进时, 曲线所围的区域总在他的左边, 则称该方向为 $\partial D$ 的正向. 相对 $D$ 来说, $\Gamma_0$ 的正向是逆时针方向, 而 $\Gamma_k(k=1,2,\cdots,K)$ 的正向是顺时针方向。

有了上述准备, 我们可以证明下述定理。

\begin{theorem}[格林公式]\label{course-theorem-16-5-1}\quad 设 $D\subset\mathbb R^2$ 是有界闭区域, 其边界由有限条分段光滑的约当曲线组成. 若函数 $P(x,y),Q(x,y)$ 在 $D$ 上具有连续偏导数, 则有
\[
  \int_{\partial D}P\mathrm dx+Q\mathrm dy
  =\iint_D\left(\frac{\partial Q}{\partial x}-\frac{\partial P}{\partial y}\right)\mathrm dx\mathrm dy,
\]
其中 $\partial D$ 的方向为正向。

\end{theorem}
\begin{proof}\quad 我们先证 $D$ 是单连通区域的情形. 证明分以下几步进行:

(1) 设 $D$ 是 $Y$ 型区域, 即
\[
  D=\{(x,y):\varphi(y)\leqslant x\leqslant\psi(y),c\leqslant y\leqslant d\},
\]
下面根据区域的上下边界方程证明
\[
  \int_{\partial D}Q(x,y)\mathrm dy=\iint_D\frac{\partial Q}{\partial x}\mathrm dx\mathrm dy.
\]
我们有
\[
  \iint_D\frac{\partial Q}{\partial x}\mathrm dx\mathrm dy
  =\int_c^d\mathrm dy\int_{\varphi(y)}^{\psi(y)}
  \frac{\partial Q}{\partial x}\mathrm dx
  =\int_c^d[Q(\psi(y),y)-Q(\varphi(y),y)]\mathrm dy.
\]
由
\[
\begin{aligned}
  \int_{\partial D}Q(x,y)\mathrm dy
  &=\int_{\widehat{AB}}Q(x,y)\mathrm dy+\int_{\widehat{BC}}Q(x,y)\mathrm dy
    +\int_{\widehat{CE}}Q(x,y)\mathrm dy+\int_{\widehat{EA}}Q(x,y)\mathrm dy,\\
  \int_{\widehat{AB}}Q(x,y)\mathrm dy&=\int_{\widehat{CE}}Q(x,y)\mathrm dy=0,\\
  \int_{\widehat{BC}}Q(x,y)\mathrm dy&=\int_c^d Q(\psi(y),y)\mathrm dy,
\end{aligned}
\]
及
\[
  \int_{\widehat{EA}}Q(x,y)\mathrm dy
  =\int_d^cQ(\varphi(y),y)\mathrm dy
  =-\int_c^dQ(\varphi(y),y)\mathrm dy,
\]
因此有
\[
  \int_{\partial D}Q(x,y)\mathrm dy=\iint_D\frac{\partial Q}{\partial x}\mathrm dx\mathrm dy.
\]

(2) 若 $D$ 是 $X$ 型区域, 与 (1) 同理可证
\[
  \int_{\partial D}P(x,y)\mathrm dx=-\iint_D\frac{\partial P}{\partial y}\mathrm dx\mathrm dy.
\]

(3) 若 $D$ 既是 $X$ 型又是 $Y$ 型区域, 则成立
\[
  \int_{\partial D}P\mathrm dx+Q\mathrm dy
  =\iint_D\left(\frac{\partial Q}{\partial x}-\frac{\partial P}{\partial y}\right)\mathrm dx\mathrm dy.
\]

若区域 $D$ 加上有限条光滑曲线后, 可分成有限个既是 $X$ 型又是 $Y$ 型区域, 则可在每个小区域上应用上述结果. 将每个小区域边界上的曲线积分与区域上的重积分分别相加, 由于添加的曲线必然是两个小区域的公共边界, 因此在曲线积分求和中出现两次, 但因方向相反而相互抵消, 所以总有
\[
  \int_{\partial D}P\mathrm dx+Q\mathrm dy
  =\iint_D\left(\frac{\partial Q}{\partial x}-\frac{\partial P}{\partial y}\right)\mathrm dx\mathrm dy.
\]
对一般的单连通区域, 可以用分成有限个既是 $X$ 型又是 $Y$ 型区域的区域逼近方法证之, 但这方面内容在数学分析课程中难以介绍, 因此我们略去该证明。

当 $D$ 是 $K(K>1)$ 连通时, 我们同样可以在区域 $D$ 内添加若干条光滑曲线, 将 $D$ 分成有限个单连通区域 . 在每个小单连通区域上由上所证成立格林公式, 然后将每个小区域边界上的曲线积分与区域上的重积分分别相加, 注意到所添加的光滑曲线仍然是两个小区域的公共边界, 因此和式中在所添加曲线上的曲线积分将相互抵消, 所以格林公式对 $K$ 连通区域仍然成立. 证毕。

在上述格林公式中, 我们假定了函数 $P$ 和 $Q$ 在 $D$ 上具有连续偏导数, 这就要求 $P,Q$ 在包含 $D$ 的某个开集内具有连续偏导数. 另外, 对于有限条封闭曲线 $\Gamma$ 所围的区域 $D$, 为了强调曲线的封闭性, 今后我们常常用
\[
  \oint_{\partial D}P\mathrm dx+Q\mathrm dy
\]
来表示沿 $\partial D$ 关于 $D$ 的正向的积分。

对于格林公式, 我们还可以有以下形式. 设光滑曲线 $\Gamma$ 由
\[
  \begin{cases}
    x=x(t),\\
    y=y(t),
  \end{cases}\quad (t\in[\alpha,\beta])
\]
给出, $t$ 的增加方向恰好是 $\Gamma$ 的正向. 记该曲线在 $(x,y)$ 处的单位切向量为 $\boldsymbol v$, 则有
\[
  \mathrm dx=x'(t)\mathrm dt
  =\frac{x'(t)}{\sqrt{x'^2(t)+y'^2(t)}}\sqrt{x'^2(t)+y'^2(t)}\mathrm dt
  =\cos(\boldsymbol v,x)\mathrm ds,
\]
其中 $\cos(\boldsymbol v,x)$ 为切向量 $\boldsymbol v$ 与 $x$ 轴正向的夹角的余弦. 同理, 有
\[
  \mathrm dy=y'(t)\mathrm dt=\cos(\boldsymbol v,y)\mathrm ds,
\]
其中 $\cos(\boldsymbol v,y)$ 为切向量 $\boldsymbol v$ 与 $y$ 轴正向的夹角的余弦. 设函数 $P(x,y),Q(x,y)$ 在 $\Gamma$ 上连续, 则
\[
  \int_\Gamma P\mathrm dx+Q\mathrm dy
  =\int_\Gamma[P\cos(\boldsymbol v,x)+Q\cos(\boldsymbol v,y)]\mathrm ds.
\]
现取 $\Gamma$ 在 $(x,y)$ 处的单位法向量为 $\boldsymbol n$, 使得 $\boldsymbol n$ 与 $\boldsymbol v$ 成右手系. 这样, 在封闭曲线 $\Gamma$ 围成区域 $D$ 时, 若取 $\Gamma$ 的正向, 则选取的法向量指向 $D$ 的外部 , 从而有
\[
  \begin{cases}
    \cos(\boldsymbol n,x)=\cos(\boldsymbol v,y),\\
    \cos(\boldsymbol n,y)=-\cos(\boldsymbol v,x),
  \end{cases}
\]
其中 $\cos(\boldsymbol n,x)$ 与 $\cos(\boldsymbol n,y)$ 分别为法向 $\boldsymbol n$ 与 $x$ 轴, $y$ 轴正向的夹角的余弦. 因此
\[
\begin{aligned}
  &\int_\Gamma[P\cos(\boldsymbol n,x)+Q\cos(\boldsymbol n,y)]\mathrm ds\\
  &=\int_\Gamma[-Q\cos(\boldsymbol v,x)+P\cos(\boldsymbol v,y)]\mathrm ds\\
  &=\int_\Gamma(-Q)\mathrm dx+P\mathrm dy
  =\iint_D\left(\frac{\partial P}{\partial x}+\frac{\partial Q}{\partial y}\right)\mathrm dx\mathrm dy.
\end{aligned}
\]

另外, 在定积分的应用中, 我们曾用
\[
  \int_\Gamma x\mathrm dy=-\int_\Gamma y\mathrm dx
  =\frac12\int_\Gamma x\mathrm dy-y\mathrm dx
\]
来表示封闭曲线 $\Gamma$ 所围区域 $D$ 的面积. 现在利用格林公式可以清楚看到上述三个积分都等于 $\iint_D\mathrm dx\mathrm dy$, 即 $D$ 的面积。

\end{proof}
\begin{example}\label{course-example-16-5-1}\quad 设 $\Gamma=\widehat{AB}\subset\mathbb R^2$ 是从点 $A(0,0)$ 到点 $B(0,1)$ 的任一条光滑曲线, 试计算第二型曲线积分
\[
  \int_\Gamma 2xy\mathrm dx+(x^2-y^2)\mathrm dy.
\]

\end{example}
\begin{solve}\quad 由于 $\Gamma$ 是任一条光滑曲线, 因此直接计算将无法求出积分值. 注意到若令
\[
  P(x,y)=2xy,\quad Q(x,y)=x^2-y^2,
\]
则有
\[
  \frac{\partial Q}{\partial x}=2x=\frac{\partial P}{\partial y}.
\]
因此对于 $\forall(x,y)\in\mathbb R^2$, 有 $\frac{\partial Q}{\partial x}-\frac{\partial P}{\partial y}=0$。

为了应用格林公式, 设实轴上从点 $(0,0)$ 到点 $(0,1)$ 的线段为 $\Gamma_2$. 我们再取一条光滑曲线 $\Gamma_1$, 使得 $\Gamma_1$ 和 $\Gamma_2$ 都以点 $(0,0),(0,1)$ 为端点, 并且 $\Gamma\cup\Gamma_1^-$ 和 $\Gamma_1\cup\Gamma_2^-$ 均为约当区域 . 设 $\Gamma\cup\Gamma_1^-$ 所围的有界闭区域为 $D_1$, 而 $\Gamma_1\cup\Gamma_2^-$ 所围的有界闭区域为 $D_2$, 由格林公式有
\[
\begin{aligned}
  \int_\Gamma P\mathrm dx+Q\mathrm dy
  &=\int_{\Gamma_1}P\mathrm dx+Q\mathrm dy+\int_{\Gamma\cup\Gamma_1^-}P\mathrm dx+Q\mathrm dy\\
  &=\int_{\Gamma_1}P\mathrm dx+Q\mathrm dy+\iint_{D_1}0\mathrm dx\mathrm dy\\
  &=\int_{\Gamma_2}P\mathrm dx+Q\mathrm dy+\int_{\Gamma_1\cup\Gamma_2^-}P\mathrm dx+Q\mathrm dy\\
  &=\int_{\Gamma_2}P\mathrm dx+Q\mathrm dy+\iint_{D_2}0\mathrm dx\mathrm dy\\
  &=\int_0^1 0\mathrm dx=0.
\end{aligned}
\]

\end{solve}
\begin{example}\label{course-example-16-5-2}\quad 计算第二型曲线积分
\[
  \oint_\Gamma\frac{x\mathrm dy-y\mathrm dx}{x^2+y^2},
\]
其中 $\Gamma\subset\mathbb R^2$ 为任一条光滑约当曲线, 且 $(0,0)\notin\Gamma$。

\end{example}
\begin{solve}\quad 记
\[
  P(x,y)=\frac{-y}{x^2+y^2},\quad Q(x,y)=\frac{x}{x^2+y^2},
\]
则对任意的 $(x,y)\ne(0,0)$, 有
\[
  \frac{\partial Q}{\partial x}=\frac{y^2-x^2}{(x^2+y^2)^2}=\frac{\partial P}{\partial y}.
\]
因此, 设 $\Gamma$ 所围成的区域为 $D$, 则当 $(0,0)\notin D$ 时, 由格林公式有
\[
  \oint_\Gamma\frac{x\mathrm dy-y\mathrm dx}{x^2+y^2}
  =\iint_D0\mathrm dx\mathrm dy=0.
\]
当 $(0,0)\in D$ 时, 格林公式条件不满足. 对于 $R>0$, 设
\[
  \Gamma_R=\{(x,y):x^2+y^2=R^2\}.
\]
取 $R$ 充分大, 使得 $\Gamma$ 落在 $\Gamma_R$ 的内部. 此时 $\Gamma$ 与 $\Gamma_R$ 围成一个二连通区域 $D_R$, 在此区域上应用格林公式, 有
\[
  0=\iint_{D_R}\left(\frac{\partial Q}{\partial x}-\frac{\partial P}{\partial y}\right)\mathrm dx\mathrm dy
  =\int_{\Gamma_R}P\mathrm dx+Q\mathrm dy+\int_{\Gamma^-}P\mathrm dx+Q\mathrm dy.
\]
因此, 我们有
\[
\begin{aligned}
  \int_\Gamma P\mathrm dx+Q\mathrm dy
  &=\int_{\Gamma_R}\frac{x\mathrm dy-y\mathrm dx}{x^2+y^2}
    =\frac1{R^2}\int_{\Gamma_R}x\mathrm dy-y\mathrm dx\\
  &=\frac2{R^2}\iint_{D_R}\mathrm dx\mathrm dy=2\pi.
\end{aligned}
\]

\end{solve}
\subsection{高斯公式}

在本小节中, 设 $D\subset\mathbb R^3$ 是有界闭区域, $\partial D$ 由有限块光滑双侧曲面所组成. 选定 $\partial D$ 的侧为 $D$ 的外侧, 即 $D$ 的外法向量确定的侧 (此时也称外侧为 $\partial D$ 的正向). 高斯 (Gauss) 公式将给出函数在 $D$ 上的三重积分与相关函数在 $\partial D$ 上的曲面积分的关系。

\begin{theorem}[高斯公式]\label{course-theorem-16-5-2}\quad 设 $D\subset\mathbb R^3$ 是有界闭区域, $\partial D$ 是 $D$ 的外侧, 函数\allowbreak $P(x,y,z),Q(x,y,z),R(x,y,z)$ 在 $D$ 上具有连续偏导数, 则有
\[
  \iint_{\partial D}P\mathrm dy\mathrm dz+Q\mathrm dz\mathrm dx+R\mathrm dx\mathrm dy
  =\iiint_D\left(\frac{\partial P}{\partial x}+\frac{\partial Q}{\partial y}
  +\frac{\partial R}{\partial z}\right)\mathrm dx\mathrm dy\mathrm dz.
\]

\end{theorem}
\begin{proof}\quad 高斯公式的证明与格林公式的证明具有相似性, 基本思想是: 先设 $D$ 是一些特殊的区域, 然后再用区域逼近来证明一般区域上的高斯公式. 同样, 区域逼近这一步在这里也不作介绍。

(1) 设 $D$ 是如下的闭区域:
\[
  D=\{(x,y,z):(y,z)\in\Omega,\varphi(y,z)\leqslant x\leqslant\psi(y,z)\}
\]
(通常也称为 $X$ 型区域), 其中 $\Omega\subset\mathbb R^2$ 是由有限多条光滑闭曲线所围的有界闭区域, $\varphi(y,z)$ 与 $\psi(y,z)$ 是 $\Omega$ 上的连续函数. 我们来证明
\[
  \iint_{\partial D}P(x,y,z)\mathrm dy\mathrm dz
  =\iiint_D\frac{\partial P}{\partial x}\mathrm dx\mathrm dy\mathrm dz.
\]
其证明方法是将上式两边的积分都化为 $\Omega$ 上的二重积分. 记
\[
\begin{aligned}
  S_1&=\{(x,y,z):x=\psi(y,z),(y,z)\in\Omega\},\quad \text{取前侧};\\
  S_2&=\{(x,y,z):x=\varphi(y,z),(y,z)\in\Omega\},\quad \text{取后侧};\\
  S_3&=\{(x,y,z):\varphi(y,z)\leqslant x\leqslant\psi(y,z),(y,z)\in\Omega\},
    \quad \text{取 $S_3$ 的侧为 $S$ 的外侧在 $S_3$ 的限制}.
\end{aligned}
\]
由于
\[
  \iint_{S_1}P(x,y,z)\mathrm dy\mathrm dz
  =\iint_\Omega P(\psi(y,z),y,z)\mathrm dy\mathrm dz,
\]
\[
  \iint_{S_2}P(x,y,z)\mathrm dy\mathrm dz
  =-\iint_\Omega P(\varphi(y,z),y,z)\mathrm dy\mathrm dz,\quad
  \iint_{S_3}P(x,y,z)\mathrm dy\mathrm dz=0,
\]
因此
\[
  \iint_{\partial D}P(x,y,z)\mathrm dy\mathrm dz
  =\iint_\Omega[P(\psi(y,z),y,z)-P(\varphi(y,z),y,z)]\mathrm dy\mathrm dz.
\]
而
\[
\begin{aligned}
  \iiint_D\frac{\partial P}{\partial x}\mathrm dx\mathrm dy\mathrm dz
  &=\iint_\Omega\mathrm dy\mathrm dz
    \int_{\varphi(y,z)}^{\psi(y,z)}\frac{\partial P}{\partial x}\mathrm dx\\
  &=\iint_\Omega[P(\psi(y,z),y,z)-P(\varphi(y,z),y,z)]\mathrm dy\mathrm dz,
\end{aligned}
\]
所以
\[
  \iint_{\partial D}P(x,y,z)\mathrm dy\mathrm dz
  =\iiint_D\frac{\partial P}{\partial x}\mathrm dx\mathrm dy\mathrm dz.
\]

(2) 若 $D=\{(x,y,z):(x,y)\in\Omega,\varphi(x,y)\leqslant z\leqslant\psi(x,y)\}$ (称之为 $Z$ 型区域, 其中 $\Omega,\varphi,\psi$ 所满足的条件与 (1) 中的相同), 则我们可以证明
\[
  \iint_{\partial D}R(x,y,z)\mathrm dx\mathrm dy
  =\iiint_D\frac{\partial R}{\partial z}\mathrm dx\mathrm dy\mathrm dz.
\]

(3) 当 $D=\{(x,y,z):(z,x)\in\Omega,\varphi(z,x)\leqslant y\leqslant\psi(z,x)\}$ (也称为 $Y$ 型区域, 其中 $\Omega,\varphi,\psi$ 所满足的条件与 (1) 中的相同) 时, 我们可以证明
\[
  \iint_{\partial D}Q(x,y,z)\mathrm dz\mathrm dx
  =\iiint_D\frac{\partial Q}{\partial y}\mathrm dx\mathrm dy\mathrm dz.
\]

(4) 显然, 当 $D$ 刚好同时是这三种区域时, 定理结论成立. 对于一般的中间无 ``洞'' 的有界闭区域 $D$, 当区域 $D$ 可以用有限块光滑曲面将其分成一些小闭区域, 且使得每个小闭区域同时是 $X$ 型, $Y$ 型与 $Z$ 型区域时 (如四面体区域), 高斯公式在每个小区域上成立. 再将每个小区域上相应积分相加, 注意到所添加曲面必是两个小区域的公共边界, 因此在第二型曲面积分中出现两次, 但因它们的侧刚好相反, 从而相互抵消. 所以对这类区域, 高斯公式依然成立. 若 $D$ 不能分成有限块同时是上述的 $X$ 型, $Y$ 型和 $Z$ 型区域, 我们不加证明地指出, 利用分成同时是上述三种区域的区域逼近方法, 可以证明对于 $D$ 仍成立高斯公式。

(5) 若 $D$ 是中间有 ``洞'' 的情形. 此时, 我们亦可添加若干块分片光滑的曲面, 将 $D$ 分成有限个中间无 ``洞'' 的小区域. 由上所证, 高斯公式在这些小区域上成立. 注意到所添加曲面必是两个小区域的公共边界, 因此对这类相当广泛的区域, 高斯公式依然成立. 证毕。

\end{proof}
\begin{remark}\quad 对于区域有 ``洞'' 的情形, 区域 $D$ 内的 ``洞'' $D_1$ 的边界曲面的侧相对于区域 $D$ 来说, 应该是 $\partial D_1$ 的内侧 。

\end{remark}
\begin{example}\label{course-example-16-5-3}\quad 计算第二型曲面积分
\[
  I=\iint_S(\sin yz+x)\mathrm dy\mathrm dz+(e^{xz}+y)\mathrm dz\mathrm dx+(xy+z)\mathrm dx\mathrm dy,
\]
其中 $S$ 是单位球面上半部分的上侧, 即
\[
  S=\{(x,y,z):x^2+y^2+z^2=1,z\geqslant0\}.
\]

\end{example}
\begin{solve}\quad 直接计算似乎比较复杂, 因此利用高斯公式来计算. 由于 $S$ 不封闭, 若我们取 $Oxy$ 平面上的单位圆盘 $S_1=\{(x,y):x^2+y^2\leqslant1\}$, 并取 $S_1$ 的下侧, 则 $S\cup S_1$ 构成了上半单位球体 $D$ 的边界, 且取外侧. 由高斯公式得
\[
\begin{aligned}
  &\iint_{S\cup S_1}(\sin yz+x)\mathrm dy\mathrm dz+(e^{xz}+y)\mathrm dz\mathrm dx+(xy+z)\mathrm dx\mathrm dy\\
  &=3\iiint_D\mathrm dx\mathrm dy\mathrm dz
    =3\cdot\frac12\cdot\frac{4\pi}{3}=2\pi,
\end{aligned}
\]
而在 $S_1$ 上 $\mathrm dy\mathrm dz=0,\mathrm dz\mathrm dx=0$, 又
\[
  \iint_{S_1}(xy+z)\mathrm dx\mathrm dy
  =-\iint_{x^2+y^2\leqslant1}xy\mathrm dx\mathrm dy=0,
\]
所以 $I=2\pi$。

\end{solve}
\begin{example}\label{course-example-16-5-4}\quad 设 $S\subset\mathbb R^3$ 为封闭光滑曲面, 取外侧, 以它为边界的有界闭区域是 $D$. 已知点 $(\xi,\eta,\zeta)\in\mathbb R^3$ 不在 $S$ 上. 计算高斯积分
\[
  I=\iint_S\frac{\cos(\boldsymbol r,\boldsymbol n)}{r^2}\mathrm dS,
\]
其中 $\boldsymbol r=(x-\xi,y-\eta,z-\zeta)$, $r=|\boldsymbol r|$, $\boldsymbol n$ 是 $S$ 的单位外法向量。

\end{example}
\begin{solve}\quad 由于
\[
  \cos(\boldsymbol r,\boldsymbol n)
  =\frac{x-\xi}{r}\cos(\boldsymbol n,x)
   +\frac{y-\eta}{r}\cos(\boldsymbol n,y)
   +\frac{z-\zeta}{r}\cos(\boldsymbol n,z),
\]
因此
\[
  I=\iint_S\frac{x-\xi}{r^3}\mathrm dy\mathrm dz
    +\frac{y-\eta}{r^3}\mathrm dz\mathrm dx
    +\frac{z-\zeta}{r^3}\mathrm dx\mathrm dy.
\]
因为
\[
  \frac{\partial}{\partial x}\left(\frac{x-\xi}{r^3}\right)=\frac1{r^3}-\frac{3(x-\xi)^2}{r^5},\quad
  \frac{\partial}{\partial y}\left(\frac{y-\eta}{r^3}\right)=\frac1{r^3}-\frac{3(y-\eta)^2}{r^5},
\]
\[
  \frac{\partial}{\partial z}\left(\frac{z-\zeta}{r^3}\right)=\frac1{r^3}-\frac{3(z-\zeta)^2}{r^5},
\]
所以
\[
  \frac{\partial}{\partial x}\left(\frac{x-\xi}{r^3}\right)
  +\frac{\partial}{\partial y}\left(\frac{y-\eta}{r^3}\right)
  +\frac{\partial}{\partial z}\left(\frac{z-\zeta}{r^3}\right)=0.
\]
因此, 当 $(\xi,\eta,\zeta)\notin D$ 时, 由高斯公式得 $I=0$。

当 $(\xi,\eta,\zeta)\in D$ 时, 我们可取 $\varepsilon$ 充分小, 使得球面
\[
  S_\varepsilon=\{(x,y,z):(x-\xi)^2+(y-\eta)^2+(z-\zeta)^2=\varepsilon^2\}
\]
完全落在 $D$ 的内部. 取 $S_\varepsilon$ 的内侧 $S_\varepsilon^-$, 设闭区域 $D_\varepsilon$ 以 $S\cup S_\varepsilon^-$ 为边界, 则
\[
  \iint_{S\cup S_\varepsilon^-}\frac{\cos(\boldsymbol r,\boldsymbol n)}{r^2}\mathrm dS
  =\iiint_{D_\varepsilon}0\mathrm dx\mathrm dy\mathrm dz=0.
\]
注意到在 $S_\varepsilon$ 上, $\boldsymbol r$ 与 $\boldsymbol n$ 平行, 所以
\[
\begin{aligned}
  \iint_S\frac{\cos(\boldsymbol r,\boldsymbol n)}{r^2}\mathrm dS
  &=-\iint_{S_\varepsilon^-}\frac{\cos(\boldsymbol r,\boldsymbol n)}{r^2}\mathrm dS
    =\iint_{S_\varepsilon}\frac{\cos(\boldsymbol r,\boldsymbol n)}{r^2}\mathrm dS\\
  &=\iint_{S_\varepsilon}\frac{\mathrm dS}{\varepsilon^2}
    =\frac1{\varepsilon^2}4\pi\varepsilon^2=4\pi.
\end{aligned}
\]

\end{solve}
\subsection{斯托克斯公式}

斯托克斯 (Stokes) 公式是格林公式的直接推广, 它揭示了函数在空间曲面 $S$ 上的第二型曲面积分与相关函数在 $\partial S$ 上的第二型曲线积分之间的关系. 为此我们先讨论一下空间曲面与其边界的定向问题。

设 $S$ 是空间分片光滑的双侧曲面, 其边界由有限条分段光滑的空间闭曲线组成. 给定 $S$ 的一侧, 设其由法向量 $\boldsymbol n$ 所确定. 若一个人站立与该法向量保持一致并且沿 $\partial S$ 前进时 $S$ 在其左边, 则称他的前进方向为 $\partial S$ 关于该曲面确定的侧的正向。

$\Gamma_0$ 的逆时针方向及 $\Gamma_k(k=1,2,\cdots,K)$ 的顺时针方向构成了 $S$ 关于 $\boldsymbol n$ 确定的侧的正向. 显然, 这样的定向与平面区域的边界定向是保持一致的。

\begin{theorem}[斯托克斯公式]\label{course-theorem-16-5-3}\quad 设 $S\subset\mathbb R^3$ 是光滑双侧曲面, $\partial S$ 由有限多条分段光滑曲线组成, 给定 $S$ 的一侧并取 $\partial S$ 关于该侧为正向的方向, 再设函数 $P(x,y,z),Q(x,y,z),R(x,y,z)$ 在包含 $S$ 的某个区域内具有连续偏导数, 则
\[
\begin{aligned}
  \int_{\partial S}P\mathrm dx+Q\mathrm dy+R\mathrm dz
  &=\iint_S\left(\frac{\partial R}{\partial y}-\frac{\partial Q}{\partial z}\right)\mathrm dy\mathrm dz\\
  &\quad+\left(\frac{\partial P}{\partial z}-\frac{\partial R}{\partial x}\right)\mathrm dz\mathrm dx
    +\left(\frac{\partial Q}{\partial x}-\frac{\partial P}{\partial y}\right)\mathrm dx\mathrm dy.
\end{aligned}
\]

\end{theorem}
\begin{proof}\quad 若 $\partial S$ 由有限条闭曲线所组成, 我们可以在 $S$ 上添加若干条曲线将 $S$ 分成有限块小曲面, 而每块小曲面的边界只有一条闭曲线. 如同格林公式的证明, 若在小曲面上斯托克斯公式成立, 则 $S$ 上也成立. 因此不妨设 $S$ 的边界是一条分段光滑的闭曲线. 另外, 我们也只对几种简单的曲面进行证明, 对于一般的曲面, 可以由简单曲面逼近来证明, 但这个过程非常复杂, 在此就不作介绍了。

具体地说, 我们先设 $S$ 由方程
\begin{equation}
  z=f(x,y),\quad (x,y)\in D
  \label{course-eq-16-5-1}
\end{equation}
给出, 其中 $D\subset\mathbb R^2$ 是由一条分段光滑的约当曲线所围成的区域, $\partial D$ 的定向由 $S$ 在 $D$ 的投影所确定。

一方面, 由第二型曲线积分的计算公式以及格林公式得
\[
\begin{aligned}
  \int_{\partial S}P(x,y,z)\mathrm dx
  &=\int_{\partial D}P(x,y,f(x,y))\mathrm dx\\
  &=-\iint_D\left(\frac{\partial P}{\partial y}
    +\frac{\partial P}{\partial z}\cdot\frac{\partial f}{\partial y}\right)\mathrm dx\mathrm dy.
\end{aligned}
\]
另一方面, 有
\[
\begin{aligned}
  \iint_S\frac{\partial P}{\partial z}\mathrm dz\mathrm dx
  -\frac{\partial P}{\partial y}\mathrm dx\mathrm dy
  &=\iint_D\left[\frac{\partial P}{\partial z}\left(-\frac{\partial f}{\partial y}\right)
    -\frac{\partial P}{\partial y}\right]\mathrm dx\mathrm dy\\
  &=-\iint_D\left(\frac{\partial P}{\partial y}
    +\frac{\partial P}{\partial z}\cdot\frac{\partial f}{\partial y}\right)\mathrm dx\mathrm dy.
\end{aligned}
\]
因此
\[
  \int_{\partial S}P\mathrm dx
  =\iint_S\frac{\partial P}{\partial z}\mathrm dz\mathrm dx
    -\frac{\partial P}{\partial y}\mathrm dx\mathrm dy.
\]
同样的计算可以得出
\[
  \int_{\partial S}Q\mathrm dy
  =\iint_S\left(-\frac{\partial Q}{\partial z}\right)\mathrm dy\mathrm dz
    +\frac{\partial Q}{\partial x}\mathrm dx\mathrm dy.
\]

同理, 对于曲面 $S$ 由方程 $x=g(y,z)((y,z)\in D)$ 给出的情形, 其中 $S$ 与 $D$ 的光滑性假设如 \eqref{course-eq-16-5-1}, 则有
\[
  \int_{\partial S}Q\mathrm dy+\int_{\partial S}R\mathrm dz
  =\iint_S\left(\frac{\partial R}{\partial y}-\frac{\partial Q}{\partial z}\right)\mathrm dy\mathrm dz
    +\frac{\partial Q}{\partial x}\mathrm dx\mathrm dy-\frac{\partial R}{\partial x}\mathrm dz\mathrm dx;
\]
而对于曲面 $S$ 由方程 $y=h(z,x)((z,x)\in D)$ 给出的情形, 则有
\[
  \int_{\partial S}P\mathrm dx+\int_{\partial S}R\mathrm dz
  =\iint_S\left(\frac{\partial P}{\partial z}-\frac{\partial R}{\partial x}\right)\mathrm dz\mathrm dx
    +\frac{\partial R}{\partial y}\mathrm dy\mathrm dz-\frac{\partial P}{\partial y}\mathrm dx\mathrm dy.
\]
因此, 若 $S$ 是同时能表示成上述三种形式的曲面或用光滑曲线可以将 $S$ 分成有限块这样的曲面, 则有
\[
\begin{aligned}
  \int_{\partial S}P\mathrm dx+Q\mathrm dy+R\mathrm dz
  &=\iint_S\left(\frac{\partial R}{\partial y}-\frac{\partial Q}{\partial z}\right)\mathrm dy\mathrm dz\\
  &\quad+\left(\frac{\partial P}{\partial z}-\frac{\partial R}{\partial x}\right)\mathrm dz\mathrm dx
    +\left(\frac{\partial Q}{\partial x}-\frac{\partial P}{\partial y}\right)\mathrm dx\mathrm dy.
\end{aligned}
\]
证毕。

在斯托克斯公式中, 曲面积分的被积函数似乎不太容易记忆, 我们可以用下述方式来记该公式:
\[
  \int_{\partial S}P\mathrm dx+Q\mathrm dy+R\mathrm dz
  =\iint_S
  \begin{vmatrix}
    \cos\alpha&\cos\beta&\cos\gamma\\
    \frac{\partial}{\partial x}&\frac{\partial}{\partial y}&\frac{\partial}{\partial z}\\
    P&Q&R
  \end{vmatrix}\mathrm dS
\]
\[
  =\iint_S
  \begin{vmatrix}
    \mathrm dy\mathrm dz&\mathrm dz\mathrm dx&\mathrm dx\mathrm dy\\
    \frac{\partial}{\partial x}&\frac{\partial}{\partial y}&\frac{\partial}{\partial z}\\
    P&Q&R
  \end{vmatrix},
\]
其中 $(\cos\alpha,\cos\beta,\cos\gamma)$ 为与 $\partial S$ 的正向保持一致的法向量的方向余弦. 另外, 由于此时 $\partial S$ 为封闭曲线, 我们也用 $\oint_{\partial S}P\mathrm dx+Q\mathrm dy+R\mathrm dz$ 来记 $\int_{\partial S}P\mathrm dx+Q\mathrm dy+R\mathrm dz$。

\end{proof}
\begin{example}\label{course-example-16-5-5}\quad 计算第二型曲线积分
\[
  I=\oint_{\Gamma_h}(y^2-z^2)\mathrm dx+(z^2-x^2)\mathrm dy+(x^2-y^2)\mathrm dz,
\]
其中 $\Gamma_h$ 是平面 $x+y+z=h(-1<h<1)$ 与球面 $x^2+y^2+z^2=1$ 的交线, 从 $z$ 轴看去取逆时针方向。

\end{example}
\begin{solve}\quad 设平面 $x+y+z=h(-1<h<1)$ 被圆周 $\Gamma_h$ 所围的部分为 $S_h$, 则 $S_h$ 是半径为 $\sqrt{1-h^2}$ 的圆盘. 取 $S_h$ 的法向量向上, 即选取法向量的方向余弦为
\[
  (\cos\alpha,\cos\beta,\cos\gamma)=\left(\frac1{\sqrt3},\frac1{\sqrt3},\frac1{\sqrt3}\right).
\]
利用斯托克斯公式, 有
\[
  I=\iint_{S_h}
  \begin{vmatrix}
    \frac1{\sqrt3}&\frac1{\sqrt3}&\frac1{\sqrt3}\\
    \frac{\partial}{\partial x}&\frac{\partial}{\partial y}&\frac{\partial}{\partial z}\\
    y^2-z^2&z^2-x^2&x^2-y^2
  \end{vmatrix}\mathrm dS
\]
\[
  =-\frac4{\sqrt3}\iint_{S_h}(x+y+z)\mathrm dS
  =-\frac4{\sqrt3}\iint_{S_h}h\mathrm dS
  =-\frac{4\pi}{\sqrt3}h(1-h^2).
\]

\end{solve}


\section{曲线积分与路径的无关性}\label{course-16-7}
本节用于曲线、曲面积分题中的参数化、方向和积分公式。应用 Green、Gauss 或 Stokes 公式之前，应逐项检查光滑性、闭合性、定向以及被积场的奇点；缺失边界必须补上并单独计算。

设 $D\subset\mathbb R^3$ 是一个区域, $u(x,y,z)$ 是 $D$ 内的 $C^1$ 函数, 则 $\mathrm du=u_x'\mathrm dx+u_y'\mathrm dy+u_z'\mathrm dz$ 是 $u(x,y,z)$ 在 $D$ 内的全微分. 给定 $A_0\in D$, 则对于 $\forall A\in D$, $A$ 与 $A_0$ 之间均可用 $D$ 内无穷多条分段光滑的曲线连接. 如果我们写出其中任意一条曲线 $\Gamma$ (以 $A_0$ 为起点, 以 $A$ 为终点) 的参数方程, 然后计算可知
\[
  \int_\Gamma u_x'\mathrm dx+u_y'\mathrm dy+u_z'\mathrm dz=u(A)-u(A_0).
\]
换句话说, $u_x'\mathrm dx+u_y'\mathrm dy+u_z'\mathrm dz$ 在 $D$ 内沿任何光滑曲线上的积分, 只依赖于该曲线的起点与终点, 而与曲线的选取无关. 对于这种现象, 我们引入以下概念。

\begin{definition}\label{course-definition-16-7-1}\quad 设 $D\subset\mathbb R^3$ 为一个区域, 函数 $P(x,y,z),Q(x,y,z),R(x,y,z)$ 在 $D$ 内连续. 如果对于 $\forall A_0,A\in D$ 及 $D$ 内任意一条以 $A_0$ 为起点, $A$ 为终点的分段光滑的曲线 $\Gamma$, 曲线积分
\[
  \int_\Gamma P\mathrm dx+Q\mathrm dy+R\mathrm dz
\]
的值只与 $A_0,A$ 有关, 而与曲线 $\Gamma$ 的选取无关, 则称曲线积分 $\int_\Gamma P\mathrm dx+Q\mathrm dy+R\mathrm dz$ 在 $D$ 内与路径无关。

从第二型曲线积分的物理背景可以看出, 该曲线积分与路径无关是指 $D$ 内力场 $\boldsymbol F=(P,Q,R)$ 将单位质量的粒子从点 $A_0$ 移动到点 $A$ 所做的功不依赖于粒子所走的路线, 而只依赖粒子运动的起点与终点. 从直观上看, 重力做功应具有该性质, 因为重力对质点所做的功只依赖于空间这两点之间的垂直距离。

现在, 我们来研究对于区域 $D$ 内给定的三个函数 $P,Q,R$, 何时曲线积分 $\int_\Gamma P\mathrm dx+Q\mathrm dy+R\mathrm dz$ 在 $D$ 内与路径无关. 设 $P(x,y,z),Q(x,y,z),R(x,y,z)$ 是空间 $\mathbb R^3$ 中区域 $D$ 内的三个函数, 若存在 $D$ 内的可微函数 $u=u(x,y,z)$, 使得
\[
  \mathrm du=P\mathrm dx+Q\mathrm dy+R\mathrm dz,
\]
则称 $u$ 为 $P\mathrm dx+Q\mathrm dy+R\mathrm dz$ 的一个原函数. 因此, $P\mathrm dx+Q\mathrm dy+R\mathrm dz$ 有原函数 $u$ 的等价定义是: 存在可微函数 $u$, 使得 $P\mathrm dx+Q\mathrm dy+R\mathrm dz$ 是 $u$ 的全微分。

在本节中, 我们将主要证明以下两个定理。

\end{definition}
\begin{theorem}\label{course-theorem-16-7-1}\quad 设 $D$ 是 $\mathbb R^3$ 中的区域, 函数 $P(x,y,z),Q(x,y,z),R(x,y,z)$ 在 $D$ 内连续, 则以下三个命题等价:

(1) 曲线积分 $\int_\Gamma P\mathrm dx+Q\mathrm dy+R\mathrm dz$ 在 $D$ 内与路径无关;

(2) 对 $D$ 内任何分段光滑的约当曲线 $\Gamma$, 有
\[
  \oint_\Gamma P\mathrm dx+Q\mathrm dy+R\mathrm dz=0;
\]

(3) $P\mathrm dx+Q\mathrm dy+R\mathrm dz$ 在 $D$ 内存在原函数。

\end{theorem}
\begin{proof}\quad 先证 $(1)\Rightarrow(2)$. 设 $\Gamma$ 是 $D$ 内任一条分段光滑的约当曲线. 在 $\Gamma$ 上取定两点 $A_0,A$, 则 $\Gamma$ 构成了从点 $A_0$ 到点 $A$ 的两条曲线 $\Gamma_1,\Gamma_2$, 其中 $\Gamma_1$ 取曲线 $\Gamma$ 的正向, $\Gamma_2$ 则取 $\Gamma$ 的负向 . 由 (1) 我们有
\[
  0=\int_{\Gamma_1}P\mathrm dx+Q\mathrm dy+R\mathrm dz
  -\int_{\Gamma_2}P\mathrm dx+Q\mathrm dy+R\mathrm dz
\]
\[
  =\int_{\Gamma_1+\Gamma_2^-}P\mathrm dx+Q\mathrm dy+R\mathrm dz
  =\int_\Gamma P\mathrm dx+Q\mathrm dy+R\mathrm dz,
\]
因此 (2) 成立。

再证 $(2)\Rightarrow(1)$. 对于 $\forall A_0,A\in D$, 任取两条以 $A_0$ 为起点, $A$ 为终点的光滑曲线 $\Gamma_1,\Gamma_2$. 若 $\Gamma_1\cup\Gamma_2^-$ 组成一条约当曲线, 令 $\Gamma=\Gamma_1\cup\Gamma_2^-$, 则我们有
\[
  \int_{\Gamma_1}P\mathrm dx+Q\mathrm dy+R\mathrm dz
  -\int_{\Gamma_2}P\mathrm dx+Q\mathrm dy+R\mathrm dz
\]
\[
  =\int_{\Gamma_1\cup\Gamma_2^-}P\mathrm dx+Q\mathrm dy+R\mathrm dz
  =\int_\Gamma P\mathrm dx+Q\mathrm dy+R\mathrm dz=0.
\]
当 $\Gamma_1\cup\Gamma_2^-$ 不是约当曲线时, 在 $D$ 内选取一条从点 $A_0$ 到点 $A$ 的光滑曲线 $\Gamma_3$, 使得 $\Gamma_1\cup\Gamma_3^-,\Gamma_2\cup\Gamma_3^-$ 均是约当曲线. 由上面所证有
\[
  \int_{\Gamma_1}P\mathrm dx+Q\mathrm dy+R\mathrm dz
  =\int_{\Gamma_3}P\mathrm dx+Q\mathrm dy+R\mathrm dz
  =\int_{\Gamma_2}P\mathrm dx+Q\mathrm dy+R\mathrm dz,
\]
因此 (1) 成立。

现证 $(1)\Rightarrow(3)$. 取定 $A_0=(x_0,y_0,z_0)\in D$, 任取 $A=(x,y,z)\in D$, 则可记
\[
  u(x,y,z)=\int_{\Gamma_0}P(\xi,\eta,\zeta)\mathrm d\xi
  +Q(\xi,\eta,\zeta)\mathrm d\eta
  +R(\xi,\eta,\zeta)\mathrm d\zeta,
\]
其中 $\Gamma_0$ 为 $D$ 中一条从点 $A_0$ 到点 $A$ 的固定的光滑曲线。

由于 $A\in D$, 因此 $A$ 是 $D$ 的内点. 对于充分小的 $\Delta x$, 有 $A_{\Delta x}=(x+\Delta x,y,z)\in D$, 于是线段 $\overline{AA_{\Delta x}}\subset D$. 记 $\Gamma_{\Delta x}=\Gamma_0\cup\overline{AA_{\Delta x}}$, 则
\[
  u(x+\Delta x,y,z)-u(x,y,z)
  =
  \int_{\Gamma_{\Delta x}}P\mathrm d\xi+Q\mathrm d\eta+R\mathrm d\zeta
  -\int_{\Gamma_0}P\mathrm d\xi+Q\mathrm d\eta+R\mathrm d\zeta
\]
\[
  =\int_{\overline{AA_{\Delta x}}}P\mathrm d\xi+Q\mathrm d\eta+R\mathrm d\zeta
  =\int_x^{x+\Delta x}P(\xi,\eta,\zeta)\mathrm d\xi
  =P(x+\theta\Delta x,y,z)\Delta x\quad(0<\theta<1).
\]
因此
\[
  \frac{\partial u}{\partial x}
  =\lim_{\Delta x\to0}\frac{u(x+\Delta x,y,z)-u(x,y,z)}{\Delta x}
  =\lim_{\Delta x\to0}P(x+\theta\Delta x,y,z)=P(x,y,z).
\]
同理可证 $\dfrac{\partial u}{\partial y}=Q$ 及 $\dfrac{\partial u}{\partial z}=R$。

由 $P,Q,R$ 的连续性知, $\mathrm du=P\mathrm dx+Q\mathrm dy+R\mathrm dz$ 在 $D$ 内成立, 因此 $u(x,y,z)$ 是 $P\mathrm dx+Q\mathrm dy+R\mathrm dz$ 的一个原函数。

最后证 $(3)\Rightarrow(1)$. 设 $u(x,y,z)$ 满足 $\mathrm du=P\mathrm dx+Q\mathrm dy+R\mathrm dz$, 再设 $A_0=(x_0,y_0,z_0),A_1=(x_1,y_1,z_1)\in D$ 是任意两点, 任取一条从点 $A_0$ 到点 $A$ 的分段光滑曲线
\[
  \Gamma:\begin{cases}
    x=x(t),\\
    y=y(t),\\
    z=z(t),
  \end{cases}\quad t\in[\alpha,\beta],
\]
且 $(x(\alpha),y(\alpha),z(\alpha))=(x_0,y_0,z_0),(x(\beta),y(\beta),z(\beta))=(x_1,y_1,z_1)$, 则
\[
  \int_\Gamma P\mathrm dx+Q\mathrm dy+R\mathrm dz
\]
\[
  =\int_\alpha^\beta
  [P(x(t),y(t),z(t))x'(t)+Q(x(t),y(t),z(t))y'(t)
  +R(x(t),y(t),z(t))z'(t)]\mathrm dt
\]
\[
  =\int_\alpha^\beta\frac{\mathrm du(x(t),y(t),z(t))}{\mathrm dt}\mathrm dt
  =u(x(\beta),y(\beta),z(\beta))-u(x(\alpha),y(\alpha),z(\alpha))
\]
\[
  =u(x_1,y_1,z_1)-u(x_0,y_0,z_0).
\]
这说明了曲线积分 $\int_\Gamma P\mathrm dx+Q\mathrm dy+R\mathrm dz$ 在 $D$ 内与路径无关. 证毕。

当曲线积分 $\int_\Gamma P\mathrm dx+Q\mathrm dy+R\mathrm dz$ 在 $D$ 内与路径无关时, 对于 $A_0=(x_0,y_0,z_0),A_1=(x_1,y_1,z_1)\in D$, 我们用
\[
  \int_{(x_0,y_0,z_0)}^{(x_1,y_1,z_1)}P\mathrm dx+Q\mathrm dy+R\mathrm dz
\]
来表示 $D$ 内从点 $A_0$ 到点 $A_1$ 的任一条路径的曲线积分。

为了得到进一步的判定曲线积分与路径无关的准则, 我们需要对空间区域作一些限制. 设 $D\subset\mathbb R^3$ 是一个区域, 如果 $D$ 内任何一条约当曲线 $\Gamma$ 总能在 $D$ 内连续地缩成 $D$ 内的一点, 则称 $D$ 是一个单连通区域. 因此对于 $\mathbb R^3$ 中的单连通区域 $D$, $D$ 内部可以有“洞”, 如同心球面围成的区域就是单连通的. 而对于形如轮胎的环面所围区域就不是单连通的了. 空间的单连通区域 $D$ 的一个特点是: 对于 $D$ 中任何一条分段光滑的约当曲线 $\Gamma$, 我们总可以在 $D$ 内找一个分片光滑的双侧曲面 $S$, 使得 $\Gamma=\partial S$. 上述事实的证明超出了数学分析课程的范围, 在这里我们就不介绍了. 但是在下面定理的证明中我们将不加证明地利用该事实。

\end{proof}
\begin{theorem}\label{course-theorem-16-7-2}\quad 设 $D\subset\mathbb R^3$ 是单连通区域, 函数 $P(x,y,z),Q(x,y,z),R(x,y,z)$ 在 $D$ 内具有连续偏导数, 则曲线积分 $\int_\Gamma P\mathrm dx+Q\mathrm dy+R\mathrm dz$ 在 $D$ 内与路径无关的充分必要条件是: 在 $D$ 内有
\[
  \frac{\partial R}{\partial y}\equiv\frac{\partial Q}{\partial z},
  \quad
  \frac{\partial P}{\partial z}\equiv\frac{\partial R}{\partial x},
  \quad
  \frac{\partial Q}{\partial x}\equiv\frac{\partial P}{\partial y},
\]
即向量
\[
  \begin{vmatrix}
    \boldsymbol i&\boldsymbol j&\boldsymbol k\\
    \frac{\partial}{\partial x}&\frac{\partial}{\partial y}&\frac{\partial}{\partial z}\\
    P&Q&R
  \end{vmatrix}
  \equiv0.
\]

\end{theorem}
\begin{proof}\quad 必要性\quad 设 $\int_\Gamma P\mathrm dx+Q\mathrm dy+R\mathrm dz$ 在 $D$ 内与路径无关, 则由定理~\ref{course-theorem-16-7-1}  知, 存在定义在 $D$ 内的可微函数 $u(x,y,z)$, 使得
\[
  P\mathrm dx+Q\mathrm dy+R\mathrm dz=u_x'\mathrm dx+u_y'\mathrm dy+u_z'\mathrm dz.
\]
因此, 在 $D$ 内恒有
\[
  \frac{\partial R}{\partial y}
  =\frac{\partial^2u}{\partial y\partial z}
  =\frac{\partial^2u}{\partial z\partial y}
  =\frac{\partial Q}{\partial z},
  \quad
  \frac{\partial P}{\partial z}
  =\frac{\partial^2u}{\partial z\partial x}
  =\frac{\partial^2u}{\partial x\partial z}
  =\frac{\partial R}{\partial x},
\]
\[
  \frac{\partial Q}{\partial x}
  =\frac{\partial^2u}{\partial x\partial y}
  =\frac{\partial^2u}{\partial y\partial x}
  =\frac{\partial P}{\partial y},
\]
即必要性成立。

充分性\quad 任取 $D$ 内一条分段光滑的约当曲线 $\Gamma$, 作 $D$ 内一个分片光滑的双侧曲面 $S$, 使得 $\partial S=\Gamma$, 且选取 $S$ 的侧使得 $\Gamma$ 定向为正向. 由斯托克斯公式我们有
\[
  \oint_\Gamma P\mathrm dx+Q\mathrm dy+R\mathrm dz
  =
  \iint_S
  \begin{vmatrix}
    \mathrm dy\mathrm dz&\mathrm dz\mathrm dx&\mathrm dx\mathrm dy\\
    \frac{\partial}{\partial x}&\frac{\partial}{\partial y}&\frac{\partial}{\partial z}\\
    P&Q&R
  \end{vmatrix}
  =0.
\]
由定理~\ref{course-theorem-16-7-1}  知, $\int_\Gamma P\mathrm dx+Q\mathrm dy+R\mathrm dz$ 在 $D$ 内与路径无关. 证毕。

由于我们在很多时候要讨论平面区域内的曲线积分与路径的关系, 我们不加证明地给出下面定理。

\end{proof}
\begin{theorem}\label{course-theorem-16-7-3}\quad 设 $D\subset\mathbb R^2$ 是一个区域, 函数 $P(x,y),Q(x,y)$ 在 $D$ 内连续, 则以下断言等价:

(1) 曲线积分 $\int_\Gamma P\mathrm dx+Q\mathrm dy$ 在 $D$ 内与路径无关;

(2) 对 $D$ 内任何分段光滑的约当曲线 $\Gamma$, 有 $\oint_\Gamma P\mathrm dx+Q\mathrm dy=0$;

(3) 在 $D$ 内存在可微函数 $u(x,y)$, 使得 $\mathrm du=P\mathrm dx+Q\mathrm dy$。

当 $D$ 是单连通区域且 $P,Q$ 具有连续偏导数时, 则下述断言与上述三个断言等价:
\[
  (4)\quad \text{在 }D\text{ 内成立 }\frac{\partial Q}{\partial x}\equiv\frac{\partial P}{\partial y}.
\]
此定理的证明与定理~\ref{course-theorem-16-7-1}  和 16.7.2 证明方法相同, 不过在证明 (4) 时, 要用到格林公式。

\end{theorem}
\begin{example}\label{course-example-16-7-1}\quad 记 $E=\{(x,y):x\in[0,1],y=0\}$, 证明曲线积分
\[
  \int_\Gamma\frac{x\mathrm dy-y\mathrm dx}{x^2+y^2}
  -\frac{(x-1)\mathrm dy-y\mathrm dx}{(x-1)^2+y^2}
\]
在 $D=\mathbb R^2\setminus E$ 内与路径无关。

\end{example}
\begin{proof}\quad 设 $\Gamma_0\subset D$ 是任一条分段光滑的约当曲线, 记
\[
  I=\int_{\Gamma_0}\frac{x\mathrm dy-y\mathrm dx}{x^2+y^2}
  -\frac{(x-1)\mathrm dy-y\mathrm dx}{(x-1)^2+y^2},
\]
则
\[
  I=\int_{\Gamma_0}\frac{x\mathrm dy-y\mathrm dx}{x^2+y^2}
  -\int_{\Gamma_0}\frac{(x-1)\mathrm dy-y\mathrm dx}{(x-1)^2+y^2}.
\]
由于
\[
  \frac{\partial}{\partial x}\left(\frac{x}{x^2+y^2}\right)
  =\frac{y^2-x^2}{(x^2+y^2)^2}
  =\frac{\partial}{\partial y}\left(\frac{-y}{x^2+y^2}\right)
\]
及
\[
  \frac{\partial}{\partial x}\left(\frac{x-1}{(x-1)^2+y^2}\right)
  =\frac{y^2-(x-1)^2}{[(x-1)^2+y^2]^2}
  =\frac{\partial}{\partial y}\left(\frac{-y}{(x-1)^2+y^2}\right),
\]
因此, 对任意的光滑约当曲线 $\Gamma_0\subset D$, 设 $\Gamma_0$ 所围区域为 $D_0$, 若 $D_0\cap E=\varnothing$, 则由定理~\ref{course-theorem-16-7-3}  知 $I=0$. 当 $E\subset D_0$ 时, 点 $(0,0)$ 及 $(1,0)$ 均在 $D_0$ 内, 从而由例~\ref{course-example-16-5-2}  知
\[
  \int_{\Gamma_0}\frac{x\mathrm dy-y\mathrm dx}{x^2+y^2}=2\pi,
  \quad
  \int_{\Gamma_0}\frac{(x-1)\mathrm dy-y\mathrm dx}{(x-1)^2+y^2}=2\pi.
\]
因此同样有 $I=0$. 由定理~\ref{course-theorem-16-7-3}  知, 曲线积分
\[
  \int_\Gamma\frac{x\mathrm dy-y\mathrm dx}{x^2+y^2}
  -\frac{(x-1)\mathrm dy-y\mathrm dx}{(x-1)^2+y^2}
\]
在 $D$ 内与路径无关。

\end{proof}
\begin{remark}\quad 显然, 若在 $\mathbb R^2\setminus\{(0,0),(0,1)\}$ 内考虑曲线积分
\[
  \int_\Gamma\frac{x\mathrm dy-y\mathrm dx}{x^2+y^2}
  -\frac{(x-1)\mathrm dy-y\mathrm dx}{(x-1)^2+y^2},
\]
则它在 $D$ 内与路径有关. 另外, 我们容易看出, 在 $\mathbb R^2\setminus\{(-\infty,0]\cup[1,+\infty)\}$ 内该曲线积分与路径无关. 以后从复变函数课程中, 我们还将讨论上述的积分问题, 在那里我们可以更加清楚地知道这些曲线积分与路径无关的本质。

\end{remark}
\begin{example}\label{course-example-16-7-2}\quad 设 $f(r)$ 是定义在 $(0,+\infty)$ 内的一元连续函数, 其中 $r=\sqrt{x^2+y^2+z^2}$. 证明: 当区域 $D\subset\mathbb R^3$ 不含原点时, 曲线积分
\[
  \int_\Gamma f(r)(x\mathrm dx+y\mathrm dy+z\mathrm dz)
\]
在 $D$ 内与路径无关。

\end{example}
\begin{proof}\quad 由于 $rf(r)$ 连续, 从而它存在原函数 $u(r)$. 取 $u(\sqrt{x^2+y^2+z^2})$, 则当 $D$ 不含原点时, 有
\[
  \mathrm du
  =\frac{\partial u}{\partial r}\cdot\frac{\partial r}{\partial x}\mathrm dx
  +\frac{\partial u}{\partial r}\cdot\frac{\partial r}{\partial y}\mathrm dy
  +\frac{\partial u}{\partial r}\cdot\frac{\partial r}{\partial z}\mathrm dz
\]
\[
  =f(r)r\left(\frac{x}{r}\mathrm dx+\frac{y}{r}\mathrm dy+\frac{z}{r}\mathrm dz\right)
  =f(r)(x\mathrm dx+y\mathrm dy+z\mathrm dz).
\]
由定理~\ref{course-theorem-16-7-1}  知 $\int_\Gamma f(r)(x\mathrm dx+y\mathrm dy+z\mathrm dz)$ 在 $D$ 内与路径无关。

\end{proof}
\begin{remark}\quad 选择不同的 $f(r)$, 我们可以得到许多在 $\mathbb R^3$ 中不含原点的区域内的函数 $P,Q,R$, 使得 $\int_\Gamma P\mathrm dx+Q\mathrm dy+R\mathrm dz$ 在 $D$ 内与路径无关。

最后我们介绍一下如何来求原函数, 即给定微分 $P\mathrm dx+Q\mathrm dy+R\mathrm dz$, 当 $\int_\Gamma P\mathrm dx+Q\mathrm dy+R\mathrm dz$ 在区域 $D$ 内与路径无关时, 如何找到在 $D$ 内的函数 $u(x,y,z)$, 使得 $\mathrm du=P\mathrm dx+Q\mathrm dy+R\mathrm dz$。

第一种方法是不定积分法: 由于
\[
  \frac{\partial u}{\partial x}=P,\quad
  \frac{\partial u}{\partial y}=Q,\quad
  \frac{\partial u}{\partial z}=R,
\]
将 $y,z$ 视为常数, 对 $\dfrac{\partial u}{\partial x}=P$ 关于 $x$ 求不定积分得
\[
  u(x,y,z)=\int P\mathrm dx+\varphi(y,z),
\]
其中 $\varphi(y,z)$ 是与 $x$ 无关的 $(y,z)$ 的待定函数. 因此我们有
\[
  \frac{\partial\varphi}{\partial y}
  =Q-\frac{\partial}{\partial y}\int P\mathrm dx.
\]
在上面等式两边对 $y$ 求不定积分得
\[
  \varphi(y,z)=\int\left(Q-\frac{\partial}{\partial y}\int P\mathrm dx\right)\mathrm dy+\psi(z),
\]
其中 $\psi(z)$ 是与 $y$ 无关的 $z$ 的待定函数. 因此我们有
\[
  \frac{\partial\psi}{\partial z}
  =
  R-\frac{\partial}{\partial z}\int P\mathrm dx
  -\frac{\partial}{\partial z}
  \left(\int\left(Q-\frac{\partial}{\partial y}\int P\mathrm dx\right)\mathrm dy\right).
\]
在上面等式两边对 $z$ 求不定积分得
\[
  \psi(z)=\int\left[
  R-\frac{\partial}{\partial z}\int P\mathrm dx
  -\frac{\partial}{\partial z}
  \left(\int\left(Q-\frac{\partial}{\partial y}\int P\mathrm dx\right)\mathrm dy\right)
  \right]\mathrm dz.
\]
最后我们有
\[
  u(x,y,z)=\int P\mathrm dx
  +\int\left(Q-\frac{\partial}{\partial y}\int P\mathrm dx\right)\mathrm dy
\]
\[
  \quad+\int\left[
  R-\frac{\partial}{\partial z}\int P\mathrm dx
  -\frac{\partial}{\partial z}
  \left(\int\left(Q-\frac{\partial}{\partial y}\int P\mathrm dx\right)\mathrm dy\right)
  \right]\mathrm dz.
\]

第二种方法是曲线积分法: 给定 $D$ 中一点 $A_0=(x_0,y_0,z_0)$, 对于任意一点 $A=(x,y,z)\in D$, 在 $D$ 内选择一些连接 $A_0,A$ 且易于求曲线积分的路径, 然后计算曲线积分即得。

\end{remark}
\begin{example}\label{course-example-16-7-3}\quad 在 $\mathbb R^3$ 内求 $u(x,y,z)$, 使得
\[
  \mathrm du(x,y,z)=(yze^{xyz}+3x^2)\mathrm dx+(xze^{xyz}+\sin y)\mathrm dy+(xye^{xyz}+2z)\mathrm dz.
\]

\end{example}
\begin{solve}\quad 记
\[
  P(x,y,z)=yze^{xyz}+3x^2,\quad Q(x,y,z)=xze^{xyz}+\sin y,
\]
\[
  R(x,y,z)=xye^{xyz}+2z.
\]
容易验证它们处处成立
\[
  \frac{\partial R}{\partial y}=\frac{\partial Q}{\partial z},\quad
  \frac{\partial P}{\partial z}=\frac{\partial R}{\partial x},\quad
  \frac{\partial Q}{\partial x}=\frac{\partial P}{\partial y},
\]
因此 $\int_\Gamma P\mathrm dx+Q\mathrm dy+R\mathrm dz$ 在 $\mathbb R^3$ 内与路径无关。

下面我们用两种方法来求 $u(x,y,z)$。

方法 1\quad 由 $\dfrac{\partial u}{\partial x}=yze^{xyz}+3x^2$, 两边对 $x$ 求不定积分得
\begin{equation}
  u(x,y,z)=e^{xyz}+x^3+\varphi(y,z),
  \label{course-eq-16-7-1}
\end{equation}
其中 $\varphi(y,z)$ 为与 $x$ 无关的 $(y,z)$ 的待定函数. 上式两边对 $y$ 求偏导数得
\begin{equation}
  Q(x,y,z)=xze^{xyz}+\sin y
  =\frac{\partial u}{\partial y}
  =xze^{xyz}+\frac{\partial\varphi}{\partial y},
  \label{course-eq-16-7-2}
\end{equation}
因此 $\dfrac{\partial\varphi}{\partial y}=\sin y$. 该式两边对 $y$ 求不定积分得
\begin{equation}
  \varphi(y,z)=-\cos y+\psi(z),
  \label{course-eq-16-7-3}
\end{equation}
其中 $\psi(z)$ 为与 $x,y$ 无关的 $z$ 的待定函数. 将式~\eqref{course-eq-16-7-3}  代入式~\eqref{course-eq-16-7-1}, 再对 $z$ 求偏导数得
\[
  R(x,y,z)=xye^{xyz}+2z=\frac{\partial u}{\partial z}=xye^{xyz}+\psi'(z),
\]
从而 $\psi'(z)=2z$. 此式两边求不定积分得 $\psi(z)=z^2+C$, 其中 $C$ 为常数. 因此
\[
  u(x,y,z)=e^{xyz}+x^3-\cos y+z^2+C.
\]

方法 2\quad 设 $\mathbb R^3$ 中任意一点的坐标为 $(\xi,\eta,\zeta)$, 并设 $\Gamma_1,\Gamma_2,\Gamma_3$ 分别是连接 $(0,0,0)$ 与 $(\xi,0,0)$, $(\xi,0,0)$ 与 $(\xi,\eta,0)$, $(\xi,\eta,0)$ 与 $(\xi,\eta,\zeta)$ 的线段, 则 $\Gamma_1\cup\Gamma_2\cup\Gamma_3$ 是连接 $(0,0,0)$ 与 $(\xi,\eta,\zeta)$ 的一条分段光滑曲线。

由于曲线积分与路径无关, 我们得到 $P\mathrm dx+Q\mathrm dy+R\mathrm dz$ 的一个原函数如下:
\[
  u_1(\xi,\eta,\zeta)
  =\int_{(0,0,0)}^{(\xi,\eta,\zeta)}P\mathrm dx+Q\mathrm dy+R\mathrm dz
\]
\[
  =\int_0^\xi P(x,0,0)\mathrm dx
  +\int_0^\eta Q(\xi,y,0)\mathrm dy
  +\int_0^\zeta R(\xi,\eta,z)\mathrm dz
\]
\[
  =\int_0^\xi 3x^2\mathrm dx+\int_0^\eta\sin y\mathrm dy
  +\int_0^\zeta(\xi\eta e^{\xi\eta z}+2z)\mathrm dz
\]
\[
  =\xi^3-\cos\eta+1+e^{\xi\eta\zeta}-1+\zeta^2
  =e^{\xi\eta\zeta}+\xi^3-\cos\eta+\zeta^2.
\]
将 $(\xi,\eta,\zeta)$ 换回 $(x,y,z)$, 并注意到两个原函数可以相差一个常数, 我们有
\[
  u(x,y,z)=e^{xyz}+x^3-\cos y+z^2+C.
\]

\end{solve}


\section{来源索引与编目说明}

\label{sec:sources}
公开题解的来源为\href{https://github.com/lzxcr/cmc-mathematics}{cmc-mathematics 仓库}；编目规则见\href{https://github.com/lzxcr/cmc-mathematics/blob/64916fa1fdbc3e444a652eca204d5669058ee321/docs/COVERAGE.md}{收录与编号}，许可详见\href{https://github.com/lzxcr/cmc-mathematics/blob/64916fa1fdbc3e444a652eca204d5669058ee321/LICENSE.md}{作者许可声明}。初赛覆盖数学A类及分设A、B类之前的数学专业卷，不能据此声称同时收齐数学B类。本选编保留对应源码中的来源标识，以便在原库逐题追溯。

组别与课程范围参考\href{https://ist.yctu.edu.cn/2025/0926/c3316a140974/page.htm}{盐城师范学院转发的第十七届竞赛通知}；非数学类初赛为高等数学，决赛另有线性代数，数学类初赛还含高等代数与解析几何。这是该届通知的范围说明，不作为以后届次的规则保证。检索到的\href{https://www.cmathc.org.cn/st/560.html}{第1至17届初赛索引}只作检索入口，该网站自述为网络资料整理站，并非据此认定为主办方官网。

高校原卷入口包括华东师范大学的\href{https://math.ecnu.edu.cn/llin/20161022.HTM}{第八届}、\href{https://math.ecnu.edu.cn/llin/20181027.HTM}{第十届}、\href{https://math.ecnu.edu.cn/llin/20191026.HTM}{第十一届}与\href{https://math.ecnu.edu.cn/llin/20201128.HTM}{第十二届}资料页。第八届资料页另列非数学类第三题解答更正，使用原卷答案时应同时查阅勘误。网页栏目日期可能与实际决赛年份不一致，届次与年份应结合试卷标题核验。

下表列出本书数学专业组部分实际选入的来源序号。一个大题含若干小问时仍作为一道题；同题在高低年级卷重复出现时不重复计数。页码链接指向各场次首道题及其紧随的解答，供核对收录范围。
\begin{longtable}{p{0.24\textwidth}p{0.42\textwidth}p{0.22\textwidth}}
\toprule 场次 & 来源序号 & 题目与解析入口\\\midrule\endhead


初赛第 1 届（2009 年） & 4, 5, 6, 7 & 第~\pageref{ex:pre01-math:4} 页；第~\pageref{ans:pre01-math:4} 页\\

初赛第 2 届（2010 年） & 1, 3, 4, 7, 8 & 第~\pageref{ex:pre02-math:1} 页；第~\pageref{ans:pre02-math:1} 页\\

初赛第 3 届（2011 年） & 2, 4, 5, 7 & 第~\pageref{ex:pre03-math:2} 页；第~\pageref{ans:pre03-math:2} 页\\

初赛第 4 届（2012 年） & 2, 3, 5, 6 & 第~\pageref{ex:pre04-math:2} 页；第~\pageref{ans:pre04-math:2} 页\\

初赛第 5 届（2013 年） & 3, 4, 5 & 第~\pageref{ex:pre05-math:3} 页；第~\pageref{ans:pre05-math:3} 页\\

初赛第 6 届（2014 年） & 2, 4, 6 & 第~\pageref{ex:pre06-math:2} 页；第~\pageref{ans:pre06-math:2} 页\\

初赛第 7 届（2015 年） & 4, 5, 6 & 第~\pageref{ex:pre07-math:4} 页；第~\pageref{ans:pre07-math:4} 页\\

初赛第 8 届（2016 年） & 4, 5, 6 & 第~\pageref{ex:pre08-math:4} 页；第~\pageref{ans:pre08-math:4} 页\\

初赛第 9 届（2017 年） & 2, 5, 6 & 第~\pageref{ex:pre09-math:2} 页；第~\pageref{ans:pre09-math:2} 页\\

初赛第 10 届（2018 年） & 4, 5, 6 & 第~\pageref{ex:pre10-math:4} 页；第~\pageref{ans:pre10-math:4} 页\\

初赛第 11 届（2019 年） & 2, 5, 6 & 第~\pageref{ex:ch10-11th-A:2} 页；第~\pageref{ans:ch10-11th-A:2} 页\\

初赛第 12 届（2020 年） & 2, 5, 6 & 第~\pageref{ex:ch09-12th-A:2} 页；第~\pageref{ans:ch09-12th-A:2} 页\\

初赛第 13 届（2021 年） & 2, 5, 6 & 第~\pageref{ex:ch07-13th-A:2} 页；第~\pageref{ans:ch07-13th-A:2} 页\\

初赛第 13 届补赛（2021 年） & 2, 5, 6 & 第~\pageref{ex:ch08-13th-A-makeup:2} 页；第~\pageref{ans:ch08-13th-A-makeup:2} 页\\

初赛第 14 届（2022 年） & 2, 5, 6 & 第~\pageref{ex:ch04-14th-A:2} 页；第~\pageref{ans:ch04-14th-A:2} 页\\

初赛第 14 届补赛（2022 年） & 2, 5, 6 & 第~\pageref{ex:ch05-14th-A-makeup:2} 页；第~\pageref{ans:ch05-14th-A-makeup:2} 页\\

初赛第 14 届第二次补赛（2022 年） & 2, 5, 6 & 第~\pageref{ex:ch06-14th-A-makeup2:2} 页；第~\pageref{ans:ch06-14th-A-makeup2:2} 页\\

初赛第 15 届（2023 年） & 2, 5, 6 & 第~\pageref{ex:ch03-15th-A:2} 页；第~\pageref{ans:ch03-15th-A:2} 页\\

初赛第 16 届（2024 年） & 2, 5, 6 & 第~\pageref{ex:ch02-16th-A:2} 页；第~\pageref{ans:ch02-16th-A:2} 页\\

初赛第 17 届（2025 年） & 2, 5, 6 & 第~\pageref{ex:ch01-17th-A:2} 页；第~\pageref{ans:ch01-17th-A:2} 页\\

决赛第 1 届（2010 年） & 1, 2, 3, 5, 6, 7, 8 & 第~\pageref{ex:fin01-all:1} 页；第~\pageref{ans:fin01-all:1} 页\\

决赛第 2 届（2011 年） & 2, 3, 5 & 第~\pageref{ex:fin02-all:2} 页；第~\pageref{ans:fin02-all:2} 页\\

决赛第 3 届（2012 年） & 2, 3, 4, 7 & 第~\pageref{ex:fin03-all:2} 页；第~\pageref{ans:fin03-all:2} 页\\

决赛第 4 届（2013 年） & 1, 2, 4, 6 & 第~\pageref{ex:fin04-all:1} 页；第~\pageref{ans:fin04-all:1} 页\\

决赛第 5 届（2014 年） & 3, 4, 5, 12 & 第~\pageref{ex:fin05-hi:3} 页；第~\pageref{ans:fin05-hi:3} 页\\

决赛第 6 届（2015 年） & 2, 3, 7, 8, 9 & 第~\pageref{ex:fin06-lo:2} 页；第~\pageref{ans:fin06-lo:2} 页\\

决赛第 7 届（2016 年） & 2, 3, 7, 8, 10, 13 & 第~\pageref{ex:fin07-hi:2} 页；第~\pageref{ans:fin07-hi:2} 页\\

决赛第 8 届（2017 年） & 2, 3, 7, 8, 10 & 第~\pageref{ex:fin08-hi:2} 页；第~\pageref{ans:fin08-hi:2} 页\\

决赛第 9 届（2018 年） & 2, 3, 7, 10 & 第~\pageref{ex:fin09-hi:2} 页；第~\pageref{ans:fin09-hi:2} 页\\

决赛第 10 届（2019 年） & 2, 3, 7, 8, 10, 13 & 第~\pageref{ex:fin10-hi:2} 页；第~\pageref{ans:fin10-hi:2} 页\\

决赛第 11 届（2021 年） & 1, 2, 7, 10 & 第~\pageref{ex:fin11-high:1} 页；第~\pageref{ans:fin11-high:1} 页\\

决赛第 12 届（2021 年） & 1, 2, 7, 8, 10, 15 & 第~\pageref{ex:fin12-high:1} 页；第~\pageref{ans:fin12-high:1} 页\\

决赛第 13 届（2022 年） & 1, 2, 4, 7, 10, 12 & 第~\pageref{ex:fin13-high:1} 页；第~\pageref{ans:fin13-high:1} 页\\

决赛第 14 届（2023 年） & 1, 2, 4, 7, 10, 12 & 第~\pageref{ex:fin14-high:1} 页；第~\pageref{ans:fin14-high:1} 页\\

决赛第 15 届（2024 年） & 2, 3, 4, 7, 10, 12, 16 & 第~\pageref{ex:fin15-high:2} 页；第~\pageref{ans:fin15-high:2} 页\\

决赛第 16 届（2025 年） & 2, 3, 4, 7, 9, 10, 12, 16 & 第~\pageref{ex:fin16-high:2} 页；第~\pageref{ans:fin16-high:2} 页\\

决赛第 17 届（2026 年） & 2, 3, 4, 7, 9, 10, 12 & 第~\pageref{ex:fin17-high:2} 页；第~\pageref{ans:fin17-high:2} 页\\

\bottomrule\end{longtable}\end{document}